% !TeX program = pdflatex
% Topologia ogólna — rozdziały 01–25 stabilne; rozdział 26 w audycie.
\documentclass[11pt,a4paper,oneside]{report}
\usepackage[utf8]{inputenc}
\usepackage[T1]{fontenc}
\usepackage{lmodern}
\usepackage[polish]{babel}
\usepackage[margin=2.5cm,headheight=14pt]{geometry}
\usepackage{amsmath,amssymb,amsthm,mathtools}
\usepackage{enumitem}
\usepackage{tikz}
\usetikzlibrary{arrows.meta,calc,positioning}
\usepackage{microtype}
\usepackage{placeins}
\usepackage[unicode]{hyperref}
\hypersetup{colorlinks=true,linkcolor=blue!45!black,urlcolor=blue!45!black}
\usepackage{fancyhdr}
\theoremstyle{plain}
\newtheorem{twierdzenie}{Twierdzenie}[section]
\newtheorem{lemat}[twierdzenie]{Lemat}
\newtheorem{stwierdzenie}[twierdzenie]{Stwierdzenie}
\newtheorem{wniosek}[twierdzenie]{Wniosek}
\theoremstyle{definition}
\newtheorem{definicja}[twierdzenie]{Definicja}
\newtheorem{przyklad}[twierdzenie]{Przykład}
\theoremstyle{remark}
\newtheorem{uwaga}[twierdzenie]{Uwaga}
\renewcommand{\proofname}{Dowód}
\newcommand{\R}{\mathbb{R}}
\newcommand{\Q}{\mathbb{Q}}
\setlength{\parindent}{0pt}
\setlength{\parskip}{4pt}
\setlength{\emergencystretch}{2em}
\pagestyle{fancy}
\fancyhf{}
\fancyhead[L]{\small Topologia ogólna}
\fancyhead[R]{\small\ifnum\value{chapter}>0 Rozdział \thechapter\else Spis treści\fi}
\fancyfoot[C]{\thepage}
\begin{document}
\hypersetup{pageanchor=false}
\begin{titlepage}
\centering
\vspace*{4cm}
{\Huge\bfseries Topologia ogólna\par}
\vspace{.7cm}
{\Large Przestrzenie, zbieżność, zwartość i konstrukcje topologiczne\par}
\vspace{1.3cm}
{\large Wersja 1.26 — rozdziały 01–25 stabilne; rozdział 26 w audycie\par}
\vfill
{\normalsize Biblioteka matematyczna\par}
\end{titlepage}
\hypersetup{pageanchor=true}
\pagenumbering{roman}
\tableofcontents
\clearpage
\pagenumbering{arabic}
\chapter{Przestrzenie topologiczne}
\label{ch:topog-przestrzenie-topologiczne}

\noindent\textbf{Cel rozdziału.}
W analizie pojęcia granicy i ciągłości poznajemy najpierw za pomocą odległości.
Zaczynamy od elementarnego wprowadzenia do przestrzeni metrycznych:
metryki, kul i sfer oraz przykładów na prostej i w przestrzeni euklidesowej.
Następnie sprawdzamy, że wiele argumentów nie wymaga liczbowej
wartości odległości: wystarcza wiedzieć, które zbiory są otwarte.
Precyzujemy aksjomaty topologii, wyprowadzamy pierwsze konsekwencje,
konstruujemy topologie z baz i podbaz oraz porównujemy różne topologie
na tym samym zbiorze. Zaawansowane własności przestrzeni metrycznych
(zbieżność, zupełność, metryzowalność) omówimy w dalszych rozdziałach.
Nie zakładamy żadnych aksjomatów oddzielania, przeliczalności
ani niepustości przestrzeni.

\section{Przestrzenie metryczne: definicje i podstawowe przykłady}
\label{sec:topog-motywacja}

Zanim zdefiniujemy topologię, wprowadzimy znany z geometrii
sposób mierzenia odległości. Posłuży nam on jako punkt wyjścia,
ale nie będziemy zakładać istnienia odległości w dowolnej
przestrzeni topologicznej.

\begin{definicja}[Metryka i przestrzeń metryczna]
\label{def:topog-metryka}
Niech $X$ będzie zbiorem. \emph{Metryką} na $X$ nazywamy
odwzorowanie $d\colon X\times X\to[0,\infty)$ takie, że
dla wszystkich $x,y,z\in X$ zachodzą:
\begin{enumerate}[label=(M\arabic*)]
\item $d(x,y)=0$ wtedy i tylko wtedy, gdy $x=y$;
\item $d(x,y)=d(y,x)$ (symetria);
\item $d(x,z)\le d(x,y)+d(y,z)$ (nierówność trójkąta).
\end{enumerate}
Parę $(X,d)$ nazywamy \emph{przestrzenią metryczną}.
Liczba $d(x,y)$ jest \emph{odległością} między $x$ a $y$.
Dopuszczamy $X=\varnothing$; jedyna funkcja
$\varnothing\times\varnothing\to[0,\infty)$ jest wtedy metryką.
\end{definicja}

\begin{przyklad}[Prosta rzeczywista i przestrzeń euklidesowa]
\label{ex:topog-metryka-euklidesowa}
Na $\R$ funkcja $d(x,y)=|x-y|$ jest metryką.
Warunek (M1) wynika z własności wartości bezwzględnej,
(M2) z równości $|x-y|=|y-x|$, a (M3) z nierówności
$|x-z|\le |x-y|+|y-z|$.

Dla $n\ge1$ na $\R^n$ kładziemy
$d_2(x,y)=\sqrt{\sum_{j=1}^{n}(x_j-y_j)^2}=\|x-y\|_2$.
Symetria i warunek (M1) są natychmiastowe.
Uzasadnijmy używaną tu nierówność Cauchy'ego--Schwarza.
Dla $v\ne0$, przyjmując
$t=\bigl(\sum_j u_jv_j\bigr)/\sum_j v_j^2$, z nierówności
$0\le\sum_j(u_j-tv_j)^2$ otrzymujemy
$(\sum_j u_jv_j)^2\le(\sum_j u_j^2)(\sum_j v_j^2)$;
dla $v=0$ nierówność jest oczywista.
Zatem $\|u+v\|_2^2=\|u\|_2^2+2\sum_j u_jv_j+\|v\|_2^2
\le(\|u\|_2+\|v\|_2)^2$.
Po wzięciu pierwiastków i podstawieniu
$u=x-y$, $v=y-z$ otrzymujemy (M3).
Dla $n=1$ wzór ten sprowadza się do wartości bezwzględnej.
\end{przyklad}

\begin{przyklad}[Metryka maksimum]
\label{ex:topog-metryka-maksimum}
Na $\R^n$, $n\ge1$, określmy
$d_\infty(x,y)=\max_{1\le j\le n}|x_j-y_j|$.
Warunki (M1) i (M2) wynikają z własności wartości
bezwzględnej. Dla każdych $x,y,z\in\R^n$ i każdego $j$
mamy $|x_j-z_j|\le|x_j-y_j|+|y_j-z_j|
\le d_\infty(x,y)+d_\infty(y,z)$.
Biorąc maksimum po $j$, dostajemy (M3).
W szczególności na $\R^2$ kula tej metryki ma kształt
kwadratu o bokach równoległych do osi współrzędnych.
\end{przyklad}

\begin{przyklad}[Metryka dyskretna]
\label{ex:topog-metryka-dyskretna}
Na dowolnym zbiorze $X$ określmy
$d_{\mathrm{dys}}(x,y)=0$, gdy $x=y$, oraz
$d_{\mathrm{dys}}(x,y)=1$, gdy $x\ne y$.
Symetria i (M1) są oczywiste. Jeśli $x\ne z$, to
przynajmniej jedna z nierówności $x\ne y$ lub $y\ne z$
jest prawdziwa. Stąd
$1=d_{\mathrm{dys}}(x,z)\le
 d_{\mathrm{dys}}(x,y)+d_{\mathrm{dys}}(y,z)$.
Dla $x=z$ nierówność trójkąta także zachodzi.
Metryka dyskretna nie wymaga żadnej struktury algebraicznej
ani porządku na $X$.
\end{przyklad}

\begin{definicja}[Kule i sfera metryczna]
\label{def:topog-kule-sfera}
Niech $(X,d)$ będzie przestrzenią metryczną,
$x\in X$ i $r>0$. Definiujemy:
\begin{align*}
B_d(x,r)&=\{y\in X:d(x,y)<r\},
&&\text{\emph{kula otwarta}},\\
\overline B_d(x,r)&=\{y\in X:d(x,y)\le r\},
&&\text{\emph{kula domknięta}},\\
S_d(x,r)&=\{y\in X:d(x,y)=r\},
&&\text{\emph{sfera metryczna}}.
\end{align*}
Punkt $x$ nazywamy \emph{środkiem}, a $r$ \emph{promieniem}.
W każdym z tych zbiorów bierzemy pod uwagę wyłącznie
punkty należące do $X$.
\end{definicja}

\begin{przyklad}[Kule na prostej i płaszczyźnie]
\label{ex:topog-kule-euklidesowe}
Dla $X=\R$ i $d(x,y)=|x-y|$ mamy
$B_d(x,r)=(x-r,x+r)$,
$\overline B_d(x,r)=[x-r,x+r]$ oraz
$S_d(x,r)=\{x-r,x+r\}$.
Na płaszczyźnie euklidesowej $B_{d_2}(x,r)$ jest
wnętrzem dysku o środku $x$ i promieniu $r$,
$\overline B_{d_2}(x,r)$ jest całym dyskiem,
a $S_{d_2}(x,r)$ jest okręgiem.
W $\R^3$ sfera $S_{d_2}(x,r)$ jest powierzchnią kuli,
nie jej wnętrzem. Dla metryki maksimum na $\R^2$
kula $B_{d_\infty}(x,r)$ jest otwartym kwadratem,
a $S_{d_\infty}(x,r)$ jego brzegiem geometrycznym.
\end{przyklad}

\begin{figure}[ht]
\centering
\begin{tikzpicture}[>=Latex, font=\small]
  \begin{scope}[shift={(0,0)}]
    \fill[blue!10] (0,0) circle (0.83);
    \draw[blue!65!black,densely dashed,thick] (0,0) circle (0.83);
    \fill (0,0) circle (1.2pt);
    \draw[->,thin] (0,0)--(.58,.58) node[pos=.61,above left] {$r$};
    \node at (0,-1.20) {$B_d(x,r)$};
    \node[font=\scriptsize,align=center] at (0,-1.60) {bez okręgu};
  \end{scope}
  \begin{scope}[shift={(3.4,0)}]
    \fill[blue!10] (0,0) circle (0.83);
    \draw[blue!65!black,thick] (0,0) circle (0.83);
    \fill (0,0) circle (1.2pt);
    \draw[->,thin] (0,0)--(.58,.58);
    \node at (0,-1.20) {$\overline B_d(x,r)$};
    \node[font=\scriptsize] at (0,-1.60) {z okręgiem};
  \end{scope}
  \begin{scope}[shift={(6.8,0)}]
    \draw[blue!65!black,thick] (0,0) circle (0.83);
    % Środek służy tylko jako odniesienie: nie jest punktem sfery.
    \draw[gray!75,fill=white] (0,0) circle (1.2pt);
    \draw[->,thin] (0,0)--(.58,.58);
    \node at (0,-1.20) {$S_d(x,r)$};
    \node[font=\scriptsize] at (0,-1.60) {tylko okrąg};
  \end{scope}
\end{tikzpicture}
\caption{Kula otwarta, kula domknięta i sfera
w płaszczyźnie z metryką euklidesową. Na rysunku
linia przerywana nie należy do kuli otwartej, a zaznaczony
środek nie należy do sfery. Dla ogólnej metryki kształt tych
zbiorów może być inny.}
\label{fig:topog-kule-sfera}
\end{figure}

\begin{definicja}[Zbiór otwarty w przestrzeni metrycznej]
\label{def:topog-otwarty-metrycznie}
Niech $(X,d)$ będzie przestrzenią metryczną.
Podzbiór $U\subseteq X$ nazywamy \emph{otwartym
w sensie metrycznym}, jeśli dla każdego $x\in U$
istnieje $r>0$ takie, że $B_d(x,r)\subseteq U$.
W szczególności $\varnothing$ jest otwarty:
nie ma w nim punktów wymagających sprawdzenia.
Podzbiór $F\subseteq X$ nazywamy \emph{domkniętym
w sensie metrycznym}, jeśli $X\setminus F$ jest otwarty.
\end{definicja}

\begin{stwierdzenie}[Otwartość i domkniętość kul oraz sfer]
\label{prop:topog-kule-otwarte-sfery-domkniete}
Niech $(X,d)$ będzie przestrzenią metryczną,
$x\in X$ i $r>0$. Wówczas kula $B_d(x,r)$ jest otwarta
w sensie metrycznym, a zbiory $\overline B_d(x,r)$
i $S_d(x,r)$ są domknięte w sensie metrycznym.
\end{stwierdzenie}
\begin{proof}
Z nierówności trójkąta i symetrii otrzymujemy dla
wszystkich $y,z\in X$
$|d(x,y)-d(x,z)|\le d(y,z)$.
Jeżeli $y\in B_d(x,r)$, to $\varepsilon=r-d(x,y)>0$.
Dla każdego $z\in B_d(y,\varepsilon)$ zachodzi
$d(x,z)\le d(x,y)+d(y,z)<r$, więc
$B_d(y,\varepsilon)\subseteq B_d(x,r)$.

Jeżeli $y\notin\overline B_d(x,r)$, to
$\varepsilon=d(x,y)-r>0$. Dla każdego
$z\in B_d(y,\varepsilon)$ mamy
$d(x,z)\ge d(x,y)-d(y,z)>r$.
Dopełnienie kuli domkniętej jest zatem otwarte.

Wreszcie, jeśli $y\notin S_d(x,r)$, to
$\varepsilon=|d(x,y)-r|>0$. Każdy
$z\in B_d(y,\varepsilon)$ spełnia
$|d(x,z)-d(x,y)|<\varepsilon$,
więc $d(x,z)$ leży po tej samej stronie $r$ co $d(x,y)$.
W szczególności $z\notin S_d(x,r)$.
Zatem dopełnienie sfery jest otwarte.
\end{proof}

\begin{uwaga}[Kula domknięta nie musi być domknięciem kuli otwartej]
\label{rem:topog-kula-domknieta-a-domkniecie}
Nazwy „kula domknięta” i „sfera” odnoszą się do
warunków określających odległość od środka. W ogólnej przestrzeni
metrycznej $\overline B_d(x,r)$ nie musi być domknięciem
$B_d(x,r)$, a sfera nie musi być brzegiem kuli otwartej.
Przykładowo, dla metryki dyskretnej na zbiorze $X$
mającym co najmniej dwa punkty i promienia $r=1$ mamy
$B_d(x,1)=\{x\}$, $\overline B_d(x,1)=X$ i
$S_d(x,1)=X\setminus\{x\}$.
Singleton $\{x\}$ jest otwarty, a także domknięty,
bo $X\setminus\{x\}$ jest sumą otwartych kul o promieniu
$1/2$ ze środkami w pozostałych punktach. Wobec tego
$\{x\}$ jest już zbiorem domkniętym, podczas gdy
kula domknięta $\overline B_d(x,1)$ jest całym $X$.
To rozróżnienie oraz pojęcie brzegu dokładniej omówimy
w rozdziale 02.
\end{uwaga}

\section{Od metryki do aksjomatów topologii}
\label{sec:topog-odleglosc-otwarte}

W przestrzeni metrycznej zadana metryka pozwala określić
zbiory otwarte. Okazuje się, że rodzina tych zbiorów
spełnia proste aksjomaty, które można przyjąć bez
odwoływania się do odległości.

\begin{stwierdzenie}[Zbiory otwarte w przestrzeni metrycznej]
\label{prop:topog-topologia-metryczna}
Niech $(X,d)$ będzie dowolną przestrzenią metryczną, a $\tau_d$
rodziną wszystkich podzbiorów $X$ otwartych w sensie
Definicji~\ref{def:topog-otwarty-metrycznie}.
Wówczas:
\begin{enumerate}[label=(\roman*)]
\item $\varnothing,X\in\tau_d$;
\item suma dowolnej rodziny zbiorów z $\tau_d$ należy do $\tau_d$;
\item przecięcie dowolnej skończonej rodziny zbiorów z $\tau_d$
      należy do $\tau_d$.
\end{enumerate}
\end{stwierdzenie}
\begin{proof}
Zbiór pusty jest otwarty na mocy definicji. Jeżeli $x\in X$,
to $B_d(x,1)\subseteq X$, więc $X$ również jest otwarty;
dotyczy to także $X=\varnothing$.

Niech $(U_i)_{i\in I}$ będzie dowolną rodziną zbiorów otwartych.
Jeżeli $x\in\bigcup_{i\in I}U_i$, to $x\in U_j$ dla pewnego $j\in I$.
Istnieje $r>0$ takie, że $B_d(x,r)\subseteq U_j\subseteq
\bigcup_{i\in I}U_i$. A zatem suma jest otwarta.
Gdy $I=\varnothing$, suma jest zbiorem pustym.

Dla dwóch zbiorów otwartych $U,V$ i punktu $x\in U\cap V$
wybieramy $r,s>0$ takie, że $B_d(x,r)\subseteq U$ i
$B_d(x,s)\subseteq V$. Wówczas
$B_d(x,\min\{r,s\})\subseteq U\cap V$.
Indukcja daje tezę dla dowolnej niepustej skończonej rodziny;
przecięciem rodziny pustej jest $X$, które jest otwarte.
\end{proof}

\begin{uwaga}
\label{rem:topog-nie-metryczna}
Aksjomaty z powyższego stwierdzenia mają sens na dowolnym zbiorze,
nawet gdy nie jest na nim zadana metryka. To prowadzi do podstawowej
definicji topologii. W dalszych rozdziałach zobaczymy także topologie,
których nie można otrzymać z żadnej metryki.
\end{uwaga}

\section{Aksjomaty przestrzeni topologicznej}
\label{sec:topog-aksjomaty}

\begin{definicja}[Topologia i przestrzeń topologiczna]
\label{def:topog-topologia}
Niech $X$ będzie dowolnym zbiorem. \emph{Topologią na $X$}
nazywamy rodzinę $\tau\subseteq\mathcal P(X)$ spełniającą:
\begin{enumerate}[label=(T\arabic*)]
\item $\varnothing\in\tau$ i $X\in\tau$;
\item dla każdej rodziny $(U_i)_{i\in I}$ elementów $\tau$
      zbiór $\bigcup_{i\in I}U_i$ należy do $\tau$;
\item dla każdych $U,V\in\tau$ zbiór $U\cap V$ należy do $\tau$.
\end{enumerate}
Parę $(X,\tau)$ nazywamy \emph{przestrzenią topologiczną},
a elementy $\tau$ \emph{zbiorami otwartymi}. Zbiór $F\subseteq X$
nazywamy \emph{domkniętym}, jeśli $X\setminus F\in\tau$.
Jeżeli topologia nie budzi wątpliwości, piszemy krótko $X$
zamiast $(X,\tau)$.
\end{definicja}

Warunek (T3) jest równoważny zamknięciu na \emph{skończone}
przecięcia: wykonujemy kolejne przecięcia dwóch zbiorów,
a przecięcie pustej rodziny to $X$. W (T2) indeks $I$
może być dowolny, także pusty. Nie wymagamy, by przecięcie
nieskończenie wielu zbiorów otwartych było otwarte.

\begin{przyklad}[Przecięcie nieskończenie wielu zbiorów otwartych]
\label{ex:topog-nieskonczone-przeciecie}
W $\R$ z topologią pochodzącą od metryki $d(x,y)=|x-y|$
zbiory $U_n=(-1/n,1/n)$, $n\geq1$, są otwarte, lecz
$\bigcap_{n\geq1}U_n=\{0\}$ nie jest otwarty.
Istotnie, dla każdego $r>0$ punkt $r/2$ należy do
$B_d(0,r)$, ale nie do $\{0\}$.
Ponadto $x\neq0$ nie należy do $U_n$ dla dostatecznie
dużego $n$, dlatego przecięcie jest dokładnie $\{0\}$.
\end{przyklad}

\begin{stwierdzenie}[Aksjomaty zbiorów domkniętych]
\label{prop:topog-zbiory-domkniete}
Niech $(X,\tau)$ będzie przestrzenią topologiczną, a
$\mathcal F=\{F\subseteq X:X\setminus F\in\tau\}$.
Wówczas $\varnothing,X\in\mathcal F$, przecięcie dowolnej
rodziny elementów $\mathcal F$ należy do $\mathcal F$
i suma dwóch elementów $\mathcal F$ należy do $\mathcal F$.
Odwrotnie, każda rodzina $\mathcal F\subseteq\mathcal P(X)$
spełniająca te trzy warunki wyznacza dokładnie jedną topologię,
której rodziną zbiorów domkniętych jest $\mathcal F$.
\end{stwierdzenie}
\begin{proof}
Pierwsze trzy własności wynikają z praw de Morgana:
$X\setminus\bigcap_{i\in I}F_i
=\bigcup_{i\in I}(X\setminus F_i)$ i
$X\setminus(F\cup G)=(X\setminus F)\cap(X\setminus G)$.
Pierwsza równość jest prawdziwa również dla pustej rodziny.

Odwrotnie, dla zadanej $\mathcal F$ definiujemy
$\tau=\{X\setminus F:F\in\mathcal F\}$. Z warunku
$\varnothing,X\in\mathcal F$ otrzymujemy (T1),
z zamknięcia $\mathcal F$ na dowolne przecięcia otrzymujemy
(T2), a z jej zamknięcia na sumy dwóch zbiorów otrzymujemy
(T3). Ponieważ $\varnothing\in\mathcal F$, rodzina
$\mathcal F$ jest także zamknięta na sumy skończone.
Rodzina zbiorów domkniętych tej topologii jest
równa $\mathcal F$. Każda topologia o tej rodzinie
zbiorów domkniętych musi składać się z ich dopełnień,
więc jest jednoznacznie wyznaczona.
\end{proof}

\section{Pierwsze przykłady i kontrprzykłady}
\label{sec:topog-przyklady}

\begin{przyklad}[Topologia dyskretna i antydyskretna]
\label{ex:topog-dyskretna-antydyskretna}
Dla dowolnego zbioru $X$ rodzina $\mathcal P(X)$ jest topologią,
ponieważ wszystkie występujące w aksjomatach zbiory są
podzbiorami $X$. Jest to \emph{topologia dyskretna}.
Jest ona topologią pochodzącą od metryki dyskretnej
z Przykładu~\ref{ex:topog-metryka-dyskretna}: dla każdego $x\in X$
mamy $B_{d_{\mathrm{dys}}}(x,1/2)=\{x\}$, a więc
każdy podzbiór $U\subseteq X$ jest otwarty, także gdy $U=\varnothing$.
Rodzina $\{\varnothing,X\}$ także spełnia (T1)--(T3):
sumy i przecięcia jej elementów należą do tej rodziny.
Jest to \emph{topologia antydyskretna}.
Jeżeli $X=\varnothing$ lub $X$ jest jednoelementowy,
obie topologie są identyczne; jeżeli $X$ ma co najmniej
dwa elementy, są różne.
\end{przyklad}

\begin{przyklad}[Topologia koskończona]
\label{ex:topog-kofinita}
Niech $X$ będzie dowolnym zbiorem. Ustalmy
\[
\tau_{\mathrm{ks}}=\{\varnothing\}\cup
\{U\subseteq X:X\setminus U\text{ jest skończony}\}.
\]
Zbiory $\varnothing$ i $X$ należą do $\tau_{\mathrm{ks}}$.
Jeżeli rodzina elementów $\tau_{\mathrm{ks}}$ zawiera
niepusty składnik $U$, to dopełnienie jej sumy jest zawarte
w skończonym zbiorze $X\setminus U$, a zatem jest skończone.
Jeżeli wszystkie składniki są puste albo rodzina jest pusta,
ich suma jest pusta.
Dla niepustych $U,V\in\tau_{\mathrm{ks}}$ dopełnienie
$U\cap V$ jest sumą dwóch zbiorów skończonych.
Przecięcie z $\varnothing$ również jest otwarte.
Są to dokładnie aksjomaty topologii.

Gdy $X$ jest skończony, $\tau_{\mathrm{ks}}=\mathcal P(X)$.
Gdy $X$ jest nieskończony, dwa niepuste zbiory otwarte
nie mogą być rozłączne: ich dopełnienia są skończone,
a zbiory dopełniające nie mogłyby razem pokryć
nieskończonego $X$.
\end{przyklad}

\begin{przyklad}[Topologia koprzeliczalna]
\label{ex:topog-koprzeliczalna}
Niech $X$ będzie dowolnym zbiorem. Definiujemy
\[
\tau_{\mathrm{kp}}=\{\varnothing\}\cup
\{X\setminus C:C\subseteq X,\ C\text{ jest co najwyżej przeliczalny}\}.
\]
Argument z poprzedniego przykładu działa po zastąpieniu
słowa „skończony” słowami „co najwyżej przeliczalny”:
podzbiór zbioru co najwyżej przeliczalnego jest
co najwyżej przeliczalny, a suma dwóch zbiorów
co najwyżej przeliczalnych jest co najwyżej przeliczalna.
Zatem jest to topologia. Jeśli $X$ jest co najwyżej
przeliczalny, otrzymujemy topologię dyskretną.
Jeżeli $X$ jest nieprzeliczalny, to każde dwa niepuste
zbiory otwarte mają niepuste przecięcie, ponieważ suma
dwóch co najwyżej przeliczalnych dopełnień nie może pokryć $X$.
\end{przyklad}

\begin{przyklad}[Przestrzeń Sierpińskiego]
\label{ex:topog-sierpinski}
Na $X=\{0,1\}$ określmy
$\tau_{\mathrm{Si}}=\{\varnothing,\{1\},X\}$. Łatwo bezpośrednio
sprawdzić, że suma i przecięcie dowolnych elementów
tej rodziny należą do $\tau_{\mathrm{Si}}$: jedynymi możliwościami
są $\varnothing,\{1\},X$. Jest to topologia różna
zarówno od dyskretnej, jak i antydyskretnej.
Zbiór $\{1\}$ jest otwarty, lecz nie jest domknięty,
ponieważ jego dopełnienie $\{0\}$ nie jest otwarte.
Natomiast $\{0\}$ jest domknięty, lecz nieotwarty.
\end{przyklad}

\begin{przyklad}[Topologia wyznaczona przez podział]
\label{ex:topog-podzial}
Niech $\mathcal A$ będzie podziałem zbioru $X$:
każdy $A\in\mathcal A$ jest niepusty, różne elementy
$\mathcal A$ są rozłączne, a $\bigcup\mathcal A=X$.
Rodzina $\tau_{\mathcal A}=\{\bigcup\mathcal E:
\mathcal E\subseteq\mathcal A\}$ jest topologią.
Zbiór pusty powstaje jako suma pustej podrodziny,
a $X$ jako suma całego podziału. Suma dowolnych sum
bloków znów jest sumą bloków. Przecięcie dwóch takich
sum jest sumą bloków należących do obu wybranych
podrodzin: każdy punkt należy do dokładnie jednego bloku.
Gdy $\mathcal A=\{X\}$ i $X\neq\varnothing$, mamy
topologię antydyskretną; gdy
$\mathcal A=\{\{x\}:x\in X\}$, mamy topologię dyskretną.
\end{przyklad}

\begin{uwaga}[Przestrzeń pusta]
\label{rem:topog-pusta}
Na $X=\varnothing$ istnieje tylko jedna topologia,
$\tau=\{\varnothing\}$. W aksjomacie (T1) zbiory
$\varnothing$ i $X$ są wtedy tym samym zbiorem.
Uwzględnianie przypadku pustego zapobiega
niepotrzebnym wyjątkom w ogólnych konstrukcjach.
\end{uwaga}

\section{Otoczenia punktów}
\label{sec:topog-otoczenia}

\begin{definicja}[Otoczenie i otoczenie otwarte]
\label{def:topog-otoczenie}
Niech $(X,\tau)$ będzie przestrzenią topologiczną i $x\in X$.
Zbiór $N\subseteq X$ nazywamy \emph{otoczeniem punktu $x$},
jeśli istnieje $U\in\tau$ takie, że $x\in U\subseteq N$.
Jeżeli dodatkowo $N\in\tau$, mówimy o
\emph{otoczeniu otwartym}. Rodzinę wszystkich otoczeń
punktu $x$ oznaczamy przez $\mathcal N(x)$.
\end{definicja}

Nie każde otoczenie jest otwarte: w standardowej topologii
na $\R$ zbiór $[-1,1]$ jest otoczeniem zera, ponieważ
$0\in(-1,1)\subseteq[-1,1]$, lecz sam nie jest otwarty.

\begin{stwierdzenie}[Podstawowe własności otoczeń]
\label{prop:topog-otoczenia}
Dla dowolnej przestrzeni topologicznej $(X,\tau)$
i każdego $x\in X$ zachodzą następujące własności:
\begin{enumerate}[label=(\roman*)]
\item $X\in\mathcal N(x)$, a każde $N\in\mathcal N(x)$ zawiera $x$;
\item jeżeli $N\in\mathcal N(x)$ oraz $N\subseteq M\subseteq X$,
      to $M\in\mathcal N(x)$;
\item jeżeli $N,M\in\mathcal N(x)$, to
      $N\cap M\in\mathcal N(x)$;
\item jeżeli $N\in\mathcal N(x)$, to istnieje $M\in\mathcal N(x)$
      takie, że $N\in\mathcal N(y)$ dla każdego $y\in M$.
\end{enumerate}
Zbiór $U\subseteq X$ jest otwarty wtedy i tylko wtedy, gdy
$U\in\mathcal N(x)$ dla każdego $x\in U$.
\end{stwierdzenie}
\begin{proof}
Ponieważ $X$ jest otwarty, jest otoczeniem każdego swojego punktu.
Pozostała część (i) wynika wprost z definicji.
Jeżeli $x\in V\subseteq N\subseteq M$ i $V$ jest otwarty,
to $M$ również jest otoczeniem $x$, co daje (ii).
Dla (iii) wybieramy otwarte $V,W$ takie, że
$x\in V\subseteq N$ i $x\in W\subseteq M$.
Zbiór $V\cap W$ jest otwarty i
$x\in V\cap W\subseteq N\cap M$.
W (iv) możemy wziąć $M=V$, gdzie $V$ jest dowolnym
otwartym zbiorem spełniającym $x\in V\subseteq N$:
dla każdego $y\in V$ ten sam $V$ świadczy,
że $N$ jest otoczeniem $y$.

Jeżeli $U$ jest otwarty, jest otoczeniem każdego punktu,
który zawiera. Odwrotnie, załóżmy, że $U\in\mathcal N(x)$
dla każdego $x\in U$. Każdy taki punkt należy do pewnego
zbioru otwartego $V$ zawartego w $U$. Wobec tego
\[
U=\bigcup\{V\in\tau:V\subseteq U\},
\]
więc $U$ jest otwarty. Równość pozostaje prawdziwa także dla
$U=\varnothing$: rodzina po prawej stronie jest wtedy
jednoelementowa i składa się ze zbioru pustego.
\end{proof}

\begin{uwaga}
\label{rem:topog-otoczenia-topologia}
Własność (iv) odzwierciedla fakt, że punkt nie musi być
jedynym punktem posiadającym dane otoczenie: pewien
jego mniejszy fragment jest otoczeniem każdego swojego punktu.
Pełne odwrotne twierdzenie o opisywaniu topologii za pomocą
rodzin otoczeń jest użyteczne przy przejściu do przestrzeni
jednostajnych; na obecnym etapie wystarczy powyższa charakteryzacja.
\end{uwaga}

\section{Bazy topologii}
\label{sec:topog-bazy}

\begin{definicja}[Baza topologii]
\label{def:topog-baza}
Niech $(X,\tau)$ będzie przestrzenią topologiczną.
Rodzinę $\mathcal B\subseteq\tau$ nazywamy
\emph{bazą topologii $\tau$}, jeśli dla każdego
$U\in\tau$ oraz $x\in U$ istnieje $B\in\mathcal B$
takie, że $x\in B\subseteq U$.
\end{definicja}

Równoważnie, każdy zbiór otwarty jest sumą pewnej podrodziny
$\mathcal B$. Istotnie, wybierając wszystkie elementy bazy
zawarte w $U$, otrzymujemy sumę równą $U$;
w drugą stronę, jeśli $U$ jest taką sumą, każdy
jego punkt należy do któregoś składnika.

\begin{twierdzenie}[Kryterium bycia bazą]
\label{thm:topog-kryterium-bazy}
Niech $X$ będzie zbiorem, a $\mathcal B\subseteq\mathcal P(X)$.
Istnieje topologia na $X$, dla której $\mathcal B$ jest bazą,
wtedy i tylko wtedy, gdy:
\begin{enumerate}[label=(B\arabic*)]
\item dla każdego $x\in X$ istnieje $B\in\mathcal B$
      takie, że $x\in B$;
\item jeżeli $B_1,B_2\in\mathcal B$ oraz
      $x\in B_1\cap B_2$, to istnieje $B_3\in\mathcal B$
      takie, że $x\in B_3\subseteq B_1\cap B_2$.
\end{enumerate}
Jeżeli te warunki są spełnione, odpowiednia topologia
jest wyznaczona jednoznacznie i ma postać
\[
\tau_{\mathcal B}=
\left\{\bigcup\mathcal E:\mathcal E\subseteq\mathcal B\right\}.
\]
\end{twierdzenie}
\begin{proof}
Jeżeli $\mathcal B$ jest bazą topologii $\tau$, to dla
$U=X$ jej definicja daje (B1). Jeżeli
$x\in B_1\cap B_2$, zbiór $B_1\cap B_2$ jest otwarty,
więc z definicji bazy istnieje $B_3$ jak w (B2).

Załóżmy odwrotnie (B1)--(B2) i zdefiniujmy
$\tau_{\mathcal B}$ podanym wzorem. Suma pustej
podrodziny daje $\varnothing$; z (B1) wynika
$X=\bigcup\mathcal B\in\tau_{\mathcal B}$.
Suma dowolnej rodziny elementów $\tau_{\mathcal B}$
jest sumą wszystkich występujących w ich zapisie
elementów bazy, więc należy do $\tau_{\mathcal B}$.

Niech $U=\bigcup\mathcal E$ oraz $V=\bigcup\mathcal D$,
gdzie $\mathcal E,\mathcal D\subseteq\mathcal B$.
Dla każdego $x\in U\cap V$ istnieją
$B_1\in\mathcal E$ i $B_2\in\mathcal D$
takie, że $x\in B_1\cap B_2$. Z (B2) wynika więc
istnienie $B\in\mathcal B$ spełniającego
$x\in B\subseteq B_1\cap B_2\subseteq U\cap V$.
W konsekwencji
\[
U\cap V=\bigcup\{B\in\mathcal B:B\subseteq U\cap V\},
\]
co dowodzi, że $U\cap V\in\tau_{\mathcal B}$.
Ten zapis obejmuje również przypadek $U\cap V=\varnothing$
i nie wymaga wyboru po jednym elemencie bazy dla każdego punktu.
Otrzymujemy topologię. Wszystkie elementy
$\mathcal B$ są w niej otwarte, a każdy zbiór otwarty
jest sumą elementów $\mathcal B$, więc $\mathcal B$
jest jej bazą. Jednoznaczność wynika z równoważnej
charakteryzacji bazy podanej przed twierdzeniem.
\end{proof}

\begin{stwierdzenie}[Kryterium otwartości przez bazę]
\label{prop:topog-otwartosc-przez-baze}
Niech $\mathcal B$ będzie bazą topologii $\tau$ na $X$.
Dla podzbioru $U\subseteq X$ następujące warunki są równoważne:
\begin{enumerate}[label=(\roman*)]
\item $U\in\tau$;
\item dla każdego $x\in U$ istnieje $B\in\mathcal B$
      takie, że $x\in B\subseteq U$.
\end{enumerate}
\end{stwierdzenie}
\begin{proof}
Implikacja (i)$\Rightarrow$(ii) jest definicją bazy.
W drugą stronę $U$ jest sumą wszystkich elementów
$\mathcal B$ zawartych w $U$, a więc zbiorem otwartym.
Gdy $U=\varnothing$, warunek (ii) jest spełniony,
a $U$ jest otwarty z aksjomatu (T1).
\end{proof}

\begin{przyklad}[Kule metryczne jako baza]
\label{ex:topog-kule-baza}
W dowolnej przestrzeni metrycznej $(X,d)$ rodzina
$\mathcal B_d=\{B_d(x,r):x\in X,\ r>0\}$
jest bazą topologii metrycznej. Istotnie, z
Definicji~\ref{def:topog-otwarty-metrycznie} wynika,
że dla każdego otwartego $U\subseteq X$ i $x\in U$
istnieje $B_d(x,r)$ zawarta w $U$. Każda kula należy
przy tym do topologii metrycznej na mocy
Stwierdzenia~\ref{prop:topog-kule-otwarte-sfery-domkniete}.
Są to dokładnie warunki z Definicji~\ref{def:topog-baza}.
Jeżeli $X=\varnothing$, baza jest pusta i warunek
jej definicji zachodzi bez konieczności wyboru punktu.
\end{przyklad}

\begin{przyklad}[Baza o końcach wymiernych]
\label{ex:topog-wymierna-baza}
W zwykłej topologii na $\R$ rodzina
$\mathcal B_{\mathbb Q}=\{(p,q):p,q\in\Q,\ p<q\}$
jest bazą. Niech $U$ będzie otwarty i $x\in U$.
Istnieje $\varepsilon>0$ takie, że
$(x-\varepsilon,x+\varepsilon)\subseteq U$.
Z gęstości $\Q$ w $\R$ wybieramy liczby wymierne
$p,q$ tak, aby $x-\varepsilon<p<x<q<x+\varepsilon$.
Wtedy $x\in(p,q)\subseteq U$. Rodzina jest przeliczalna: jest obrazem przeliczalnego
zbioru $\{(p,q)\in\Q\times\Q:p<q\}$ przy odwzorowaniu,
które parze liczb $(p,q)$ przyporządkowuje przedział $(p,q)$.
Własność tę wykorzystamy przy aksjomatach przeliczalności.
\end{przyklad}

\begin{przyklad}[Topologia dolnej granicy]
\label{ex:topog-sorgenfrey}
Niech $X=\R$ i
$\mathcal B_{\mathrm S}=\{[a,b):a,b\in\R,\ a<b\}$.
Rodzina pokrywa $\R$, gdyż $x\in[x,x+1)$.
Jeżeli $x\in[a,b)\cap[c,d)$, to
$x<\min\{b,d\}$ i
$[x,\min\{b,d\})\subseteq[a,b)\cap[c,d)$.
Kryterium bazy wyznacza topologię
$\tau_{\mathrm{So}}$, zwaną \emph{topologią dolnej granicy}
(topologią prostej Sorgenfreya).
Każdy zbiór otwarty w zwykłej topologii na $\R$
jest otwarty w $\tau_{\mathrm{So}}$: dla $x\in U$
można wybrać $\varepsilon>0$ z
$(x-\varepsilon,x+\varepsilon)\subseteq U$, skąd
$x\in[x,x+\varepsilon)\subseteq U$.
Natomiast $[0,1)$ jest otwarty w $\tau_{\mathrm{So}}$,
ale nie jest otwarty w topologii zwykłej,
bo żadna kula wokół zera nie mieści się w $[0,1)$.
Topologia dolnej granicy jest zatem ściśle silniejsza
od topologii zwykłej.
\end{przyklad}

\section{Podbazy i topologie generowane}
\label{sec:topog-podbazy}

\begin{definicja}[Topologia generowana przez rodzinę zbiorów]
\label{def:topog-topologia-generowana}
Niech $X$ będzie zbiorem, a $\mathcal S\subseteq\mathcal P(X)$.
\emph{Topologią generowaną przez $\mathcal S$} nazywamy
najmniejszą względem zawierania topologię na $X$,
która zawiera wszystkie elementy $\mathcal S$.
Rodzinę $\mathcal S$ nazywamy \emph{podbazą} tej topologii.
\end{definicja}

Istnienie najmniejszej topologii nie jest założeniem:
wynika z następnego twierdzenia. W definicji
\emph{nie} zakładamy, że $\mathcal S$ pokrywa $X$.

\begin{twierdzenie}[Konstrukcja topologii z podbazy]
\label{thm:topog-podbaza-konstrukcja}
Niech $X$ będzie dowolnym zbiorem i
$\mathcal S\subseteq\mathcal P(X)$. Niech $\mathcal B$
będzie rodziną wszystkich przecięć skończonych podrodzin
$\mathcal S$, przy czym przecięcie pustej podrodziny
jest równe $X$. Wówczas $\mathcal B$ jest bazą pewnej
topologii, a topologia ta jest dokładnie topologią
generowaną przez $\mathcal S$. Każdy jej zbiór otwarty
jest sumą skończonych przecięć elementów $\mathcal S$.
\end{twierdzenie}
\begin{proof}
Ponieważ $X\in\mathcal B$, każdy punkt jest zawarty
w pewnym elemencie $\mathcal B$. Przecięcie dwóch
skończonych przecięć elementów $\mathcal S$ jest
znów takim skończonym przecięciem. Zatem $\mathcal B$
spełnia (B1)--(B2) z
Twierdzenia~\ref{thm:topog-kryterium-bazy}.
Oznaczmy wyznaczoną topologię przez $\tau$.
Każdy $S\in\mathcal S$ należy do $\mathcal B$ jako
przecięcie jednoelementowej rodziny, zatem $\mathcal S
\subseteq\tau$.

Niech $\rho$ będzie dowolną topologią na $X$
zawierającą $\mathcal S$. Z (T1) i (T3) wynika,
że każde skończone przecięcie elementów $\mathcal S$
należy do $\rho$, a z (T2), że należą do niej także
wszystkie sumy takich przecięć. Zatem $\tau\subseteq\rho$.
Dowodzi to minimalności i jednoznaczności $\tau$.
\end{proof}

\begin{przyklad}[Podbaza nie musi być bazą]
\label{ex:topog-podbaza-niebaza}
Niech $X=\{a,b,c\}$ i
$\mathcal S=\{\{a,b\},\{b,c\}\}$.
Przecięcie obu elementów $\mathcal S$ jest równe
$\{b\}$, ale w $\mathcal S$ nie ma zbioru
zawierającego $b$ i zawartego w $\{b\}$.
Zatem $\mathcal S$ nie spełnia (B2), choć pokrywa $X$.
Topologia przez nią generowana to
\[
\tau=\{\varnothing,\{b\},\{a,b\},\{b,c\},X\}.
\]
Istotnie, przecięcia skończone dają tylko $X$,
$\{a,b\}$, $\{b,c\}$ i $\{b\}$,
a sumy tych zbiorów, wraz z sumą pustą,
dają dokładnie wypisaną rodzinę.
\end{przyklad}

\begin{przyklad}[Pusta podbaza]
\label{ex:topog-pusta-podbaza}
Gdy $\mathcal S=\varnothing$, jedynym skończonym
przecięciem zbiorów z $\mathcal S$ jest $X$
(przecięcie puste). Sumy podrodzin $\{X\}$
dają $\varnothing$ oraz $X$, a więc topologię
antydyskretną. Jest to zgodne z tym, że żadnych
zbiorów nie nakazano uczynić otwartymi.
\end{przyklad}

\section{Porównywanie topologii na ustalonym zbiorze}
\label{sec:topog-porownywanie}

\begin{definicja}[Topologia słabsza i silniejsza]
\label{def:topog-silniejsza}
Niech $\tau_1$ i $\tau_2$ będą topologiami na tym
samym zbiorze $X$. Mówimy, że $\tau_1$ jest
\emph{słabsza} (albo \emph{grubsza}) od $\tau_2$,
a $\tau_2$ jest \emph{silniejsza} (albo \emph{drobniejsza})
od $\tau_1$, jeżeli $\tau_1\subseteq\tau_2$.
Jeśli zawieranie jest właściwe, mówimy odpowiednio
o topologii \emph{ściśle słabszej} i \emph{ściśle silniejszej}.
\end{definicja}

To porównanie dotyczy liczby i układu zbiorów otwartych,
a nie wartości odległości. Topologia dyskretna jest
najsilniejsza, a antydyskretna najsłabsza ze wszystkich
topologii na $X$.

\begin{stwierdzenie}[Kryterium porównania przez bazy]
\label{prop:topog-porownanie-bazy}
Niech $\tau_1,\tau_2$ będą topologiami na $X$,
a $\mathcal B_1$ bazą $\tau_1$.
Wówczas $\tau_1\subseteq\tau_2$ wtedy i tylko wtedy,
gdy każdy zbiór $B\in\mathcal B_1$ należy do $\tau_2$.
Jeżeli $\mathcal B_2$ jest bazą $\tau_2$, jest to
równoważne warunkowi: dla wszystkich $B\in\mathcal B_1$
i $x\in B$ istnieje $C\in\mathcal B_2$ takie,
że $x\in C\subseteq B$.
\end{stwierdzenie}
\begin{proof}
Jeżeli $\tau_1\subseteq\tau_2$, każdy element
$\mathcal B_1\subseteq\tau_1$ należy do $\tau_2$.
Odwrotnie, jeśli wszystkie elementy $\mathcal B_1$
są otwarte w $\tau_2$, to dowolny $U\in\tau_1$
jako ich suma również należy do $\tau_2$.
Ostatnia równoważność wynika z kryterium otwartości
przez bazę $\mathcal B_2$.
\end{proof}

\begin{przyklad}[Topologie nieporównywalne]
\label{ex:topog-nieporownywalne}
Na $X=\{a,b\}$ określmy
$\tau_a=\{\varnothing,\{a\},X\}$ i
$\tau_b=\{\varnothing,\{b\},X\}$.
Obie rodziny są topologiami, ale
$\{a\}\in\tau_a\setminus\tau_b$ i
$\{b\}\in\tau_b\setminus\tau_a$.
Żadna z tych topologii nie jest słabsza od drugiej.
Na zbiorze dwuelementowym istnieją dokładnie cztery
topologie: antydyskretna, dyskretna oraz $\tau_a,\tau_b$.
Wynika to z tego, że $\varnothing$ i $X$ są obowiązkowe,
a każdy z dwóch pozostałych podzbiorów można dołączyć
lub pominąć, nie naruszając aksjomatów.
\end{przyklad}

\begin{figure}[ht]
\centering
\begin{tikzpicture}[>=Latex, every node/.style={font=\small},
 box/.style={draw,rounded corners,align=center,inner sep=4pt}]
\node[box] (top) at (0,2.6) {$\mathcal P(X)$};
\node[box] (left) at (-2.5,1.3) {$\tau_a$};
\node[box] (right) at (2.5,1.3) {$\tau_b$};
\node[box] (bottom) at (0,0) {$\{\varnothing,X\}$};
\draw (bottom)--(left)--(top)--(right)--(bottom);
\end{tikzpicture}
\caption{Cztery topologie na $X=\{a,b\}$ uporządkowane
przez zawieranie; kierunek ku górze oznacza topologię silniejszą.}
\label{fig:topog-cztery-topologie}
\end{figure}

\begin{twierdzenie}[Przecięcia topologii i topologia generowana]
\label{thm:topog-przeciecia-topologii}
Niech $X$ będzie zbiorem i $(\tau_i)_{i\in I}$
dowolną rodziną topologii na $X$. Wówczas
$\bigcap_{i\in I}\tau_i$ jest topologią na $X$,
przy czym dla $I=\varnothing$ przyjmujemy
$\bigcap_{i\in\varnothing}\tau_i=\mathcal P(X)$.
W szczególności dla każdego
$\mathcal S\subseteq\mathcal P(X)$ topologia
generowana przez $\mathcal S$ jest przecięciem
wszystkich topologii na $X$ zawierających $\mathcal S$.
\end{twierdzenie}
\begin{proof}
Dla $I\neq\varnothing$ zbiory $\varnothing,X$
należą do każdej $\tau_i$, więc także do ich przecięcia.
Jeżeli $(U_j)_{j\in J}$ jest rodziną zbiorów należących
do wszystkich $\tau_i$, to dla każdego $i\in I$
zbiór $\bigcup_{j\in J}U_j$ należy do $\tau_i$;
zatem należy do ich przecięcia. Analogicznie,
jeżeli $U,V$ należą do wszystkich $\tau_i$, to
$U\cap V$ należy do wszystkich $\tau_i$.
Dla $I=\varnothing$ otrzymujemy topologię dyskretną.

Rodzina wszystkich topologii zawierających $\mathcal S$
nie jest pusta, ponieważ zawiera $\mathcal P(X)$.
Jej przecięcie jest topologią zawierającą $\mathcal S$
i zawartą w każdej innej topologii o tej własności.
Zatem jest to najmniejsza topologia zawierająca
$\mathcal S$, zgodnie z Definicją~\ref{def:topog-topologia-generowana}.
\end{proof}

\begin{przyklad}[Porównanie trzech topologii na prostej]
\label{ex:topog-porownanie-r}
Na $\R$ topologia koskończona jest ściśle słabsza od
zwykłej topologii euklidesowej. Istotnie, dla każdego
$a\in\R$ zbiór $\R\setminus\{a\}$ jest otwarty:
jeśli $x\ne a$, to
$B_d(x,|x-a|/2)\subseteq\R\setminus\{a\}$.
Dopełnienie dowolnego zbioru skończonego jest zatem
przecięciem skończenie wielu takich zbiorów otwartych,
a więc jest otwarte. Jednak $(-1,1)$ nie jest otwarty
w topologii koskończonej, ponieważ jego dopełnienie
jest nieskończone.

Topologia euklidesowa i topologia koprzeliczalna
na $\R$ są nieporównywalne. Zbiór
$\R\setminus\Q$ jest koprzeliczalnie otwarty,
ale nie jest euklidesowo otwarty: każdy przedział
wokół liczby niewymiernej zawiera liczbę wymierną.
Z kolei $(-1,1)$ jest euklidesowo otwarty,
ale jego dopełnienie nie jest przeliczalne,
więc nie jest koprzeliczalnie otwarty.
\end{przyklad}

\section{Podsumowanie i dalsze powiązania}
\label{sec:topog-podsumowanie}

W przestrzeni metrycznej odległość umożliwia definiowanie
kul i sfer; z kul wyprowadzamy pojęcie zbioru otwartego.
Przestrzeń topologiczna to zbiór wyposażony w rodzinę zbiorów
otwartych, zamkniętą na dowolne sumy i skończone przecięcia.
Te same informacje można opisywać za pomocą zbiorów
domkniętych, otoczeń, baz albo podbaz. Znajomość bazy
pozwala sprawdzać otwartość lokalnie, w otoczeniu każdego
punktu. Podbaza daje z kolei sposób wygenerowania
najmniejszej topologii zawierającej zadaną rodzinę zbiorów.

W następnym rozdziale przejdziemy do operacji na podzbiorach
przestrzeni topologicznej: wnętrza, domknięcia, brzegu i zbioru
pochodnego. Następnie zbadamy ciągłość odwzorowań, dla której
istotna będzie właśnie rodzina zbiorów otwartych,
a nie wybrana metryka.

\chapter{Operacje topologiczne na zbiorach}
\label{ch:topog-operacje-zbiorowe}

\noindent\textbf{Cel rozdziału.}
Aksjomaty topologii określają, które podzbiory przestrzeni są otwarte.
Na tej podstawie możemy badać dowolny podzbiór, nawet jeżeli sam nie jest
ani otwarty, ani domknięty. Jego wnętrze opisuje punkty, wokół których
zbiór zawiera całe otwarte otoczenie; domknięcie dodaje punkty, do których
zbiór dochodzi w sensie topologicznym; brzeg odróżnia punkty, w których
każde otoczenie spotyka zarówno zbiór, jak i jego dopełnienie.
Następnie rozróżnimy punkty skupienia i punkty izolowane, wprowadzimy
gęstość i zbiory nigdziegęste oraz pokażemy, że domknięcie można
scharakteryzować bezpośrednio przez cztery aksjomaty Kuratowskiego.
Nie zakładamy metryzowalności, aksjomatu oddzielania $T_1$ ani
niepustości przestrzeni.

W całym rozdziale $(X,\tau)$ oznacza przestrzeń topologiczną, zaś
$A,B\subseteq X$ jej podzbiory. Dopełnienie $A$ w $X$ oznaczamy przez
$X\setminus A$. Otoczenie punktu może nie być otwarte, lecz zgodnie
z definicją zawiera wtedy pewne otwarte otoczenie tego punktu.

\section{Wnętrze i domknięcie}
\label{sec:topog-wnetrze-domkniecie}

\begin{definicja}[Wnętrze i domknięcie zbioru]
\label{def:topog-wnetrze-domkniecie}
Niech $(X,\tau)$ będzie przestrzenią topologiczną i niech $A\subseteq X$.
\emph{Wnętrzem} zbioru $A$ nazywamy zbiór
\[
\operatorname{int}_X A
 =\bigcup\{U\in\tau:U\subseteq A\}.
\]
\emph{Domknięciem} zbioru $A$ nazywamy zbiór
\[
\overline A^{\,X}
 =\bigcap\{F\subseteq X:F\text{ jest domknięty w }X,
                    \ A\subseteq F\}.
\]
Gdy przestrzeń $X$ jest jednoznaczna, piszemy $\operatorname{int}A$
oraz $\overline A$. Rodzina w drugim wzorze jest niepusta,
ponieważ zawiera $X$. Suma zbiorów po prawej stronie pierwszego wzoru
jest określona również dla $A=\varnothing$ i wtedy wynosi $\varnothing$.
\end{definicja}

Intuicyjnie $\operatorname{int}A$ to największy zbiór otwarty
zawarty w $A$, a $\overline A$ to najmniejszy zbiór domknięty
zawierający $A$. Oba określenia wynikają z następującego stwierdzenia.

\begin{stwierdzenie}[Własności ekstremalne]
\label{prop:topog-ekstremalne}
Dla każdej przestrzeni topologicznej $X$ i każdego $A\subseteq X$
zachodzą następujące własności:
\begin{enumerate}[label=(\roman*)]
\item $\operatorname{int}A$ jest otwarty, $\operatorname{int}A\subseteq A$,
      a każdy otwarty $U\subseteq A$ jest zawarty w $\operatorname{int}A$;
\item $\overline A$ jest domknięty, $A\subseteq\overline A$,
      a każdy domknięty $F\supseteq A$ zawiera $\overline A$;
\item $A$ jest otwarty wtedy i tylko wtedy, gdy
      $\operatorname{int}A=A$;
\item $A$ jest domknięty wtedy i tylko wtedy, gdy $\overline A=A$.
\end{enumerate}
\end{stwierdzenie}
\begin{proof}
Wnętrze $A$ jest sumą rodziny zbiorów otwartych,
więc jest otwarte na mocy aksjomatów topologii. Każdy składnik tej
sumy jest zawarty w $A$ i tym samym cała suma jest zawarta w $A$.
Jeżeli $U\subseteq A$ jest otwarty, to jest jednym z sumowanych zbiorów,
zatem $U\subseteq\operatorname{int}A$. Dowodzi to (i).

Domknięcie $A$ jest przecięciem rodziny zbiorów domkniętych,
więc jest domknięte. Ponieważ $A$ jest zawarty w każdym z tych zbiorów,
mamy $A\subseteq\overline A$. Każdy domknięty $F\supseteq A$
występuje w rodzinie przecinanych zbiorów, dlatego
$\overline A\subseteq F$. To dowodzi (ii).

Jeżeli $A$ jest otwarty, to w (i) możemy przyjąć $U=A$,
co razem z $\operatorname{int}A\subseteq A$ daje równość.
Jeżeli $A=\operatorname{int}A$, to $A$ jest otwarty, ponieważ
wnętrze jest otwarte. Analogicznie (ii) dowodzi (iv).
\end{proof}

\begin{twierdzenie}[Charakteryzacja punktowa]
\label{thm:topog-punktowa-charakteryzacja}
Niech $X$ będzie przestrzenią topologiczną, $A\subseteq X$ i $x\in X$.
Wówczas:
\begin{enumerate}[label=(\roman*)]
\item $x\in\operatorname{int}A$ wtedy i tylko wtedy, gdy istnieje
      otwarty zbiór $U$ taki, że $x\in U\subseteq A$;
\item $x\in\overline A$ wtedy i tylko wtedy, gdy dla każdego
      otwartego $U\ni x$ zachodzi $U\cap A\ne\varnothing$.
\end{enumerate}
Oba warunki pozostają równoważne po zastąpieniu otwartych
zbiorów $U\ni x$ dowolnymi otoczeniami punktu $x$.
\end{twierdzenie}
\begin{proof}
Warunek (i) wynika wprost z definicji wnętrza jako sumy
wszystkich otwartych podzbiorów $A$: punkt należy do tej sumy
dokładnie wtedy, gdy należy do przynajmniej jednego jej składnika.

Załóżmy, że $x\notin\overline A$. Zbiór $X\setminus\overline A$
jest wtedy otwartym otoczeniem punktu $x$ i jest rozłączny z $A$,
ponieważ $A\subseteq\overline A$. Zatem warunek (ii) nie zachodzi.
Odwrotnie, niech istnieje otwarty $U\ni x$ rozłączny z $A$.
Wówczas $X\setminus U$ jest domknięty i zawiera $A$,
więc $\overline A\subseteq X\setminus U$. Stąd
$x\notin\overline A$. Otrzymujemy równoważność przez
zanegowanie obu warunków.

Każdy zbiór otwarty zawierający $x$ jest jego otoczeniem.
Jeżeli $N$ jest otoczeniem $x$, to istnieje otwarty
$U$ taki, że $x\in U\subseteq N$.
Dlatego oba kryteria zachowują swoją treść również przy
kwantyfikowaniu po wszystkich otoczeniach.
\end{proof}

\begin{przyklad}[Przedziały w prostej rzeczywistej]
\label{ex:topog-wnetrze-domkniecie-przedzialu}
Rozważmy $\R$ z topologią euklidesową i liczby $a<b$.
Dla $A=[a,b)$ otrzymujemy
$\operatorname{int}A=(a,b)$ oraz $\overline A=[a,b]$.
Rzeczywiście, każdy punkt $x\in(a,b)$ ma przedział
$(x-\varepsilon,x+\varepsilon)\subseteq A$ dla dostatecznie
małego $\varepsilon>0$. Punkt $a$ nie jest wewnętrzny,
bo każde jego otwarte otoczenie zawiera liczby mniejsze od $a$.
Każdy punkt $[a,b]$ należy do $\overline A$: dla $b$
wynika to z faktu, że każde otwarte otoczenie $b$
zawiera liczby z $[a,b)$; dla pozostałych punktów
wynika z zawierania $A\subseteq\overline A$.
Punkt spoza $[a,b]$ ma otwarte otoczenie rozłączne z $A$.

Dla $A=\varnothing$ i $A=\R$ mamy odpowiednio
$\operatorname{int}\varnothing=\overline\varnothing=\varnothing$
oraz $\operatorname{int}\R=\overline\R=\R$.
\end{przyklad}

\begin{twierdzenie}[Rachunek wnętrz i domknięć]
\label{thm:topog-rachunek-wnetrz-domkniec}
Niech $X$ będzie przestrzenią topologiczną, a $A,B\subseteq X$.
Wówczas:
\begin{enumerate}[label=(\roman*)]
\item jeśli $A\subseteq B$, to
      $\operatorname{int}A\subseteq\operatorname{int}B$
      i $\overline A\subseteq\overline B$;
\item $\operatorname{int}(\operatorname{int}A)=\operatorname{int}A$
      i $\overline{\overline A}=\overline A$;
\item $\operatorname{int}(A\cap B)
      =\operatorname{int}A\cap\operatorname{int}B$;
\item $\overline{A\cup B}=\overline A\cup\overline B$;
\item $\operatorname{int}(A\cup B)\supseteq
      \operatorname{int}A\cup\operatorname{int}B$ oraz
      $\overline{A\cap B}\subseteq\overline A\cap\overline B$.
\end{enumerate}
Nie należy w ogólności zastępować obu zawierań w (v)
równościami.
\end{twierdzenie}
\begin{proof}
Jeżeli $A\subseteq B$, to każdy otwarty zbiór zawarty w $A$
jest także zawarty w $B$; po wzięciu ich sum otrzymujemy
pierwsze zawieranie w (i). Z kolei $\overline B$ jest
domknięty i zawiera $A$, więc z minimalności $\overline A$
wynika $\overline A\subseteq\overline B$.

Równości z (ii) wynikają z faktu, że $\operatorname{int}A$
jest otwarty, zaś $\overline A$ domknięty, oraz z (iii)--(iv)
poprzedniego stwierdzenia.

Zbiór $\operatorname{int}A\cap\operatorname{int}B$ jest otwarty
i zawiera się w $A\cap B$, stąd jest zawarty w
$\operatorname{int}(A\cap B)$. W odwrotną stronę,
z monotoniczności i zawierań $A\cap B\subseteq A,B$
wynika $\operatorname{int}(A\cap B)\subseteq
\operatorname{int}A\cap\operatorname{int}B$.

Zbiór $\overline A\cup\overline B$ jest domknięty
jako suma dwóch zbiorów domkniętych i zawiera $A\cup B$.
Stąd $\overline{A\cup B}\subseteq\overline A\cup\overline B$.
Monotoniczność daje oba zawierania
$\overline A,\overline B\subseteq\overline{A\cup B}$,
a zatem również ich sumy. Otrzymujemy (iv).

W (v) zawierania wynikają bezpośrednio z monotoniczności
zastosowanej do $A,B\subseteq A\cup B$ oraz
$A\cap B\subseteq A,B$.
\end{proof}

\begin{przyklad}[Ścisłe zawierania]
\label{ex:topog-scisle-zawierania}
W $\R$ z topologią euklidesową weźmy
$A=[0,1]$ i $B=[1,2]$. Wówczas
$\operatorname{int}A\cup\operatorname{int}B
=(0,1)\cup(1,2)$, podczas gdy
$\operatorname{int}(A\cup B)=(0,2)$.
Punkt $1$ należy tylko do drugiego ze zbiorów.

Dla $A=(0,1)$ i $B=(1,2)$ mamy $A\cap B=\varnothing$,
a więc $\overline{A\cap B}=\varnothing$,
natomiast $\overline A\cap\overline B=\{1\}$.
W ten sposób obie nierówności z poprzedniego
twierdzenia mogą być ścisłe.
\end{przyklad}

\begin{stwierdzenie}[Zależności między dopełnieniem, wnętrzem i domknięciem]
\label{prop:topog-dualnosc-wnetrze-domkniecie}
Niech $X$ będzie przestrzenią topologiczną i $A\subseteq X$.
Wtedy
\[
X\setminus\operatorname{int}A=\overline{X\setminus A},
\qquad
X\setminus\overline A=\operatorname{int}(X\setminus A).
\]
\end{stwierdzenie}
\begin{proof}
Zbiór $X\setminus\operatorname{int}A$ jest domknięty,
gdyż $\operatorname{int}A$ jest otwarty. Zawiera
$X\setminus A$, ponieważ $\operatorname{int}A\subseteq A$.
Dlatego $\overline{X\setminus A}\subseteq
X\setminus\operatorname{int}A$.
Odwrotnie, niech $F$ będzie dowolnym domkniętym
zbiorem zawierającym $X\setminus A$. Wówczas
$X\setminus F$ jest otwarty i zawarty w $A$,
więc $X\setminus F\subseteq\operatorname{int}A$.
Po przejściu do dopełnień otrzymujemy
$X\setminus\operatorname{int}A\subseteq F$.
Przecinając wszystkie takie $F$, dostajemy pierwszą równość.
Druga powstaje z pierwszej po zastąpieniu $A$ przez
$X\setminus A$ i ponownym przejściu do dopełnień.
\end{proof}

\section{Brzeg i zewnętrze zbioru}
\label{sec:topog-brzeg-zewnetrze}

\begin{definicja}[Brzeg i zewnętrze]
\label{def:topog-brzeg-zewnetrze}
Niech $X$ będzie przestrzenią topologiczną i $A\subseteq X$.
\emph{Brzegiem} $A$ nazywamy
$\partial_X A:=\overline A\setminus\operatorname{int}A$,
a \emph{zewnętrzem} $A$ zbiór
$\operatorname{ext}_X A:=\operatorname{int}(X\setminus A)$.
Jeśli przestrzeń jest ustalona, pomijamy indeks $X$.
\end{definicja}

\begin{twierdzenie}[Charakteryzacja brzegu]
\label{thm:topog-charakteryzacja-brzegu}
Niech $X$ będzie przestrzenią topologiczną,
$A\subseteq X$ i $x\in X$. Wówczas:
\begin{enumerate}[label=(\roman*)]
\item $\partial A=\overline A\cap\overline{X\setminus A}$;
\item $x\in\partial A$ wtedy i tylko wtedy, gdy każde
      otwarte otoczenie $U$ punktu $x$ spełnia jednocześnie
      $U\cap A\ne\varnothing$ i
      $U\cap(X\setminus A)\ne\varnothing$;
\item $\partial A$ jest domknięty i
      $\partial A=\partial(X\setminus A)$;
\item $X$ jest sumą rozłączną trzech zbiorów
      $\operatorname{int}A$, $\partial A$ i
      $\operatorname{ext}A$;
\item $\overline A=\operatorname{int}A\cup\partial A$.
\end{enumerate}
\end{twierdzenie}
\begin{proof}
Z twierdzenia o zależności domknięcia i wnętrza od dopełnienia
wynika $X\setminus\operatorname{int}A=\overline{X\setminus A}$.
Dlatego $\partial A=\overline A\cap(X\setminus\operatorname{int}A)
=\overline A\cap\overline{X\setminus A}$, co daje (i).
Równoważność z (ii) wynika z zastosowania punktowej
charakteryzacji domknięcia do $A$ oraz do $X\setminus A$.

W (iii) domkniętość brzegu wynika z (i), bo jest on
przecięciem dwóch zbiorów domkniętych.
Ten sam wzór jest symetryczny względem $A$ i jego
dopełnienia, stąd $\partial A=\partial(X\setminus A)$.

Zbiory $\operatorname{int}A$ i $\operatorname{ext}A$
są rozłączne, bo należą odpowiednio do $A$ i do jego dopełnienia.
Ponadto $\partial A$ nie spotyka żadnego z nich:
z definicji jest rozłączny z $\operatorname{int}A$,
a z (i) i równości
$\operatorname{ext}A=X\setminus\overline A$
wynika rozłączność z $\operatorname{ext}A$.
Niech $x\in X$ nie należy ani do $\operatorname{int}A$,
ani do $\operatorname{ext}A$. Drugi warunek daje
$x\in\overline A$, a pierwszy wtedy $x\in\partial A$.
Dowodzi to (iv). Ponieważ $\operatorname{int}A\subseteq\overline A$,
rozłączna dekompozycja $\overline A$ na
$\operatorname{int}A$ i $\overline A\setminus\operatorname{int}A$
daje (v).
\end{proof}

\begin{przyklad}[Brzeg przedziału i zbioru liczb wymiernych]
\label{ex:topog-brzeg-przyklady}
W $\R$ z topologią euklidesową dla $A=(a,b)$, gdzie $a<b$,
mamy $\operatorname{int}A=(a,b)$, $\overline A=[a,b]$,
$\partial A=\{a,b\}$ i
$\operatorname{ext}A=(-\infty,a)\cup(b,\infty)$.

Dla $A=\Q$ każdy niepusty przedział otwarty zawiera
liczbę wymierną i niewymierną. Zatem
$\overline\Q=\overline{\R\setminus\Q}=\R$, więc
$\partial\Q=\R$. Zarazem $\operatorname{int}\Q=\varnothing$
i $\operatorname{ext}\Q=\varnothing$.
Z tych samych powodów brzeg zbioru liczb niewymiernych
jest całą prostą rzeczywistą.
\end{przyklad}

\begin{figure}[ht]
\centering
\begin{tikzpicture}[x=1.15cm,y=.92cm, font=\small]
\foreach \y in {0,1,2,3} {
  \draw[gray!55] (-3,\y)--(3,\y);
}
\node[left] at (-3.25,3) {$A=[a,b)$};
\node[left] at (-3.25,2) {$\operatorname{int}A$};
\node[left] at (-3.25,1) {$\partial A$};
\node[left] at (-3.25,0) {$\operatorname{ext}A$};
\draw[blue!65!black,ultra thick] (-1.5,3)--(1.5,3);
\fill[blue!65!black] (-1.5,3) circle (2pt);
\draw[blue!65!black,fill=white,thick] (1.5,3) circle (2.4pt);
\draw[blue!65!black,ultra thick] (-1.5,2)--(1.5,2);
\draw[blue!65!black,fill=white,thick] (-1.5,2) circle (2.4pt);
\draw[blue!65!black,fill=white,thick] (1.5,2) circle (2.4pt);
\fill[red!65!black] (-1.5,1) circle (2.5pt);
\fill[red!65!black] (1.5,1) circle (2.5pt);
\draw[teal!70!black,ultra thick,<-] (-3,0)--(-1.5,0);
\draw[teal!70!black,ultra thick,->] (1.5,0)--(3,0);
\draw[teal!70!black,fill=white,thick] (-1.5,0) circle (2.4pt);
\draw[teal!70!black,fill=white,thick] (1.5,0) circle (2.4pt);
\node[below] at (-1.5,-.22) {$a$};
\node[below] at (1.5,-.22) {$b$};
\end{tikzpicture}
\caption{Wnętrze, brzeg i zewnętrze przedziału $A=[a,b)$
na prostej rzeczywistej. Przy końcach przedziałów pełna
kropka oznacza przynależność do zbioru w danym wierszu.}
\label{fig:topog-rozklad-przedzialu}
\end{figure}

\begin{przyklad}[Brzeg w przestrzeni Sierpińskiego]
\label{ex:topog-brzeg-sierpinski}
Niech $X=\{0,1\}$ i
$\tau=\{\varnothing,\{1\},X\}$. Dla $A=\{1\}$ mamy
$\operatorname{int}A=\{1\}$ oraz $\overline A=X$,
ponieważ jedynym domkniętym zbiorem zawierającym $1$ jest $X$.
Wobec tego $\partial A=\{0\}$.
Dla $B=\{0\}$ zachodzi $\operatorname{int}B=\varnothing$
i $\overline B=\{0\}$, więc również $\partial B=\{0\}$.
Pokazuje to, że brzeg nie musi być zawarty w samym zbiorze:
w tym przypadku $\partial\{1\}$ jest od niego rozłączny,
chociaż zawsze leży w jego domknięciu.
\end{przyklad}

\begin{stwierdzenie}[Brzeg sumy zbiorów]
\label{prop:topog-brzeg-sumy}
Dla dowolnej przestrzeni topologicznej $X$ i dowolnych
$A,B\subseteq X$ zachodzi
$\partial(A\cup B)\subseteq\partial A\cup\partial B$.
\end{stwierdzenie}
\begin{proof}
Niech $x\notin\partial A\cup\partial B$.
Z rozłącznego podziału przestrzeni wynika, że
$x$ należy do wnętrza albo do zewnętrza każdego ze zbiorów
$A$ i $B$. Jeżeli należy do wnętrza przynajmniej jednego
z nich, to należy do $\operatorname{int}(A\cup B)$,
a więc nie leży na brzegu sumy. Jeśli zaś należy do obu
zewnętrz, to otwarty zbiór
$\operatorname{ext}A\cap\operatorname{ext}B$
jest otoczeniem $x$ rozłącznym z $A\cup B$.
Stąd $x\in\operatorname{ext}(A\cup B)$ i znów
$x\notin\partial(A\cup B)$. Dowiedliśmy żądanego zawierania.
\end{proof}

\section{Punkty skupienia i punkty izolowane}
\label{sec:topog-punkty-skupienia}

\begin{definicja}[Zbiór pochodny i punkt izolowany]
\label{def:topog-zbior-pochodny}
Niech $X$ będzie przestrzenią topologiczną i $A\subseteq X$.
Punkt $x\in X$ nazywamy \emph{punktem skupienia} zbioru $A$,
jeżeli każde otwarte otoczenie $U\ni x$ zawiera punkt
zbioru $A$ różny od $x$, czyli
$U\cap(A\setminus\{x\})\ne\varnothing$.
Zbiór wszystkich punktów skupienia $A$ oznaczamy przez
$A'$ i nazywamy \emph{zbiorem pochodnym} zbioru $A$.

Punkt $x\in A$ nazywamy \emph{punktem izolowanym} w $A$,
jeżeli istnieje otwarty $U\ni x$ taki, że
$U\cap A=\{x\}$. Zbiór punktów izolowanych $A$
oznaczamy przez $\operatorname{Iso}(A)$.
\end{definicja}

Punkt skupienia nie musi należeć do samego zbioru $A$.
Z kolei punkt izolowany musi należeć do $A$ z definicji.
Pojęcie punktu skupienia wymaga wyraźnego wyłączenia
samego punktu $x$ z przecięcia: w przeciwnym razie
otrzymalibyśmy charakteryzację domknięcia, a nie zbioru pochodnego.

\begin{twierdzenie}[Domknięcie i zbiór pochodny]
\label{thm:topog-domkniecie-pochodny}
Dla dowolnej przestrzeni topologicznej $X$ i dowolnego $A\subseteq X$
zachodzą równości
\[
\overline A=A\cup A',\qquad
\operatorname{Iso}(A)=A\setminus A'.
\]
\end{twierdzenie}
\begin{proof}
Każdy $x\in A$ należy do $\overline A$.
Jeżeli $x\in A'$, to każde otwarte otoczenie $x$
spotyka $A\setminus\{x\}$, a zatem także $A$;
z charakteryzacji punktowej wynika $x\in\overline A$.
Stąd $A\cup A'\subseteq\overline A$.

Odwrotnie, niech $x\in\overline A\setminus A$.
Każde otwarte otoczenie $x$ spotyka $A$, a ponieważ
$x\notin A$, jego przecięcie z $A\setminus\{x\}=A$
jest niepuste. Zatem $x\in A'$, skąd pierwsza równość.

Dla $x\in A$ równość
$U\cap(A\setminus\{x\})=\varnothing$
jest równoważna warunkowi $U\cap A=\{x\}$.
Punkt $x\in A$ nie jest więc punktem skupienia dokładnie
wtedy, gdy jest izolowany w $A$. Daje to drugą równość.
\end{proof}

\begin{przyklad}[Zbiór wyrazów ciągu i jego punkt skupienia]
\label{ex:topog-pochodny-ciag}
W $\R$ z topologią euklidesową niech
$A=\{1/n:n\in\mathbb N,\ n\ge1\}$.
Każdy otwarty przedział wokół $0$ zawiera $1/n$
dla dostatecznie dużych $n$, więc $0\in A'$.
Dla ustalonego $n$ punkt $1/n$ jest izolowany w $A$:
można wybrać dostatecznie mały przedział wokół $1/n$
niezawierający żadnego innego elementu $A$
(dla $n=1$ wybieramy przedział niewchodzący w $1/2$,
a dla $n\ge2$ korzystamy z luk między kolejnymi wyrazami).
Każdy $x\notin A\cup\{0\}$ ma otwarty przedział
rozłączny z $A$; dla $x<0$ lub $x>1$ jest to oczywiste,
a dla $x\in(0,1)\setminus A$ wynika z istnienia dwóch
kolejnych wyrazów $1/(n+1)<x<1/n$.
Stąd $A'=\{0\}$ oraz $\overline A=A\cup\{0\}$.
\end{przyklad}

\begin{stwierdzenie}[Monotoniczność zbioru pochodnego]
\label{prop:topog-pochodny-monotoniczny}
Niech $X$ będzie przestrzenią topologiczną oraz $A,B\subseteq X$.
Jeśli $A\subseteq B$, to $A'\subseteq B'$.
Ponadto $(A\cup B)'=A'\cup B'$.
\end{stwierdzenie}
\begin{proof}
Jeżeli $x\in A'$, każde jego otwarte otoczenie spotyka
$A\setminus\{x\}\subseteq B\setminus\{x\}$,
więc $x\in B'$. To daje monotoniczność i zawieranie
$A'\cup B'\subseteq(A\cup B)'$.

Jeżeli $x\notin A'\cup B'$, to istnieją otwarte otoczenia
$U\ni x$ oraz $V\ni x$, dla których
$U\cap(A\setminus\{x\})=\varnothing$ i
$V\cap(B\setminus\{x\})=\varnothing$.
Otwarte otoczenie $U\cap V$ nie zawiera punktu
$(A\cup B)\setminus\{x\}$, toteż
$x\notin(A\cup B)'$. Otrzymujemy odwrotne zawieranie.
\end{proof}

\begin{twierdzenie}[Domkniętość zbioru pochodnego w przestrzeni $T_1$]
\label{thm:topog-pochodny-t1}
Niech $X$ będzie przestrzenią topologiczną spełniającą
aksjomat $T_1$, czyli każdy singleton $\{x\}$ jest domknięty.
Wówczas dla każdego $A\subseteq X$ zbiór pochodny $A'$ jest domknięty.
\end{twierdzenie}
\begin{proof}
Pokażemy, że $X\setminus A'$ jest otwarty.
Niech $x\notin A'$. Istnieje otwarty zbiór $U\ni x$
taki, że $U\cap(A\setminus\{x\})=\varnothing$.
Oczywiście sam $x$ nie należy do $A'$.
Niech $y\in U\setminus\{x\}$. Ponieważ $X$ jest $T_1$,
zbiór $X\setminus\{x\}$ jest otwarty.
Zatem $V=U\cap(X\setminus\{x\})$ jest otwartym
otoczeniem $y$, rozłącznym z $A$:
jedyny punkt $A$, który mógł należeć do $U$, to $x$,
a z $V$ usunęliśmy właśnie $x$.
Wobec tego $y\notin A'$. Otrzymujemy $U\subseteq X\setminus A'$.
Każdy punkt dopełnienia ma zatem otwarte otoczenie
zawarte w dopełnieniu, co dowodzi jego otwartości.
\end{proof}

\begin{przyklad}[Dlaczego założenia $T_1$ nie można pominąć]
\label{ex:topog-pochodny-nie-domkniety}
Niech $X=\{a,b\}$ z topologią antydyskretną
$\tau=\{\varnothing,X\}$ i weźmy $A=\{a\}$.
Dla $b$ jedyne otwarte otoczenie to $X$, które spotyka
$A\setminus\{b\}=\{a\}$, więc $b\in A'$.
Dla $a$ mamy $A\setminus\{a\}=\varnothing$, zatem
$a\notin A'$. Otrzymujemy $A'=\{b\}$, który nie jest
domknięty w topologii antydyskretnej.
Zatem domkniętość zbioru pochodnego nie wynika
z samych aksjomatów topologii.
\end{przyklad}

\begin{definicja}[Zbiór bez punktów izolowanych i zbiór doskonały]
\label{def:topog-doskonaly}
Niech $X$ będzie przestrzenią topologiczną, $A\subseteq X$.
Mówimy, że $A$ \emph{nie ma punktów izolowanych}, jeśli
$\operatorname{Iso}(A)=\varnothing$, równoważnie $A\subseteq A'$.
Zbiór $A$ nazywamy \emph{doskonałym}, jeśli jest domknięty
w $X$ i nie ma punktów izolowanych.
\end{definicja}

\begin{uwaga}
Dla dowolnej przestrzeni topologicznej $X$ i $A\subseteq X$
zbiór $A$ jest doskonały wtedy i tylko wtedy, gdy $A=A'$.
Istotnie, jeżeli $A$ jest domknięty, to
$A'\subseteq\overline A=A$, a brak punktów izolowanych
jest równoważny $A\subseteq A'$. Odwrotnie, jeśli $A=A'$,
to $\overline A=A\cup A'=A$, zatem $A$ jest domknięty;
ponadto $\operatorname{Iso}(A)=A\setminus A'=\varnothing$.
Żaden z tych argumentów nie wymaga aksjomatu $T_1$.
\end{uwaga}

\section{Gęstość i zbiory nigdziegęste}
\label{sec:topog-gestosc}

\begin{definicja}[Zbiór gęsty i zbiór nigdziegęsty]
\label{def:topog-gesty-nigdziegesty}
Niech $X$ będzie przestrzenią topologiczną i $A\subseteq X$.
Zbiór $A$ nazywamy \emph{gęstym w $X$}, jeżeli
$\overline A=X$.
Zbiór $A$ nazywamy \emph{nigdziegęstym w $X$},
jeżeli $\operatorname{int}(\overline A)=\varnothing$.
\end{definicja}

\begin{twierdzenie}[Kryteria gęstości i nigdziegęstości]
\label{thm:topog-gesty-kryteria}
Niech $X$ będzie przestrzenią topologiczną, $A\subseteq X$.
Wówczas:
\begin{enumerate}[label=(\roman*)]
\item $A$ jest gęsty w $X$ wtedy i tylko wtedy, gdy
      każdy niepusty zbiór otwarty w $X$ spotyka $A$;
\item $A$ jest nigdziegęsty w $X$ wtedy i tylko wtedy, gdy
      każdy niepusty zbiór otwarty $U\subseteq X$ zawiera
      niepusty zbiór otwarty $V\subseteq U$ rozłączny z $A$;
\item $A$ jest nigdziegęsty w $X$ wtedy i tylko wtedy, gdy
      $X\setminus\overline A$ jest otwarty i gęsty w $X$.
\end{enumerate}
\end{twierdzenie}
\begin{proof}
Załóżmy $\overline A=X$ i niech $U$ będzie niepusty
i otwarty. Wybierzmy $x\in U$.
Z punktowej charakteryzacji domknięcia wynika $U\cap A\ne\varnothing$.
Odwrotnie, jeśli każdy niepusty otwarty zbiór spotyka $A$,
to każde otwarte otoczenie dowolnego $x\in X$ spotyka $A$,
a zatem $x\in\overline A$. Dla $X=\varnothing$
oba warunki również zachodzą, ponieważ $\overline\varnothing=\varnothing$.

Jeśli $\operatorname{int}\overline A=\varnothing$,
to żaden niepusty otwarty $U$ nie może być zawarty
w $\overline A$. Wobec tego
$V=U\setminus\overline A$ jest niepusty, otwarty,
zawarty w $U$ i rozłączny z $A$.
Odwrotnie, załóżmy własność z (ii) i przypuśćmy,
że $\operatorname{int}\overline A\ne\varnothing$.
Stosując założenie do otwartego
$U=\operatorname{int}\overline A$, otrzymalibyśmy
niepusty otwarty $V\subseteq\overline A$ rozłączny z $A$.
Każdy punkt $V$ miałby wtedy otwarte otoczenie $V$
rozłączne z $A$, co przeczy $V\subseteq\overline A$.

Zbiór $X\setminus\overline A$ jest otwarty.
Na podstawie (i) jest gęsty dokładnie wtedy, gdy
każdy niepusty zbiór otwarty $U$ spotyka
$X\setminus\overline A$, czyli gdy żaden taki $U$
nie jest zawarty w $\overline A$. Ostatni warunek
jest równoważny $\operatorname{int}\overline A=\varnothing$.
\end{proof}

\begin{przyklad}[Zbiór o pustym wnętrzu nie musi być nigdziegęsty]
\label{ex:topog-q-nigdziegesty}
Na prostej rzeczywistej $\Q$ ma puste wnętrze,
ponieważ każdy niepusty otwarty przedział zawiera
liczbę niewymierną. Jednocześnie $\overline\Q=\R$,
więc $\operatorname{int}\overline\Q=\R\ne\varnothing$.
Zbiór $\Q$ jest zatem gęsty, ale nie jest nigdziegęsty.

Natomiast zbiór $A=\{0\}\cup\{1/n:n\ge1\}$
jest domknięty, ponieważ jego dopełnienie jest otwarte:
dla punktu spoza $A$ można dobrać przedział rozłączny
z $A$ (tak jak w przykładzie o zbiorze pochodnym).
Ma puste wnętrze, ponieważ $A\subseteq\Q$,
a wcześniej wykazaliśmy, że
$\operatorname{int}\Q=\varnothing$.
Dlatego $A$ jest nigdziegęsty.
\end{przyklad}

\begin{stwierdzenie}[Podzbiory i skończone sumy]
\label{prop:topog-nigdziegesty-sumy}
Niech $X$ będzie przestrzenią topologiczną.
Każdy podzbiór zbioru nigdziegęstego w $X$
jest nigdziegęsty w $X$.
Suma skończonej rodziny zbiorów nigdziegęstych
w $X$ jest nigdziegęsta w $X$.
\end{stwierdzenie}
\begin{proof}
Jeśli $B\subseteq A$ i $A$ jest nigdziegęsty,
to $\overline B\subseteq\overline A$,
zatem $\operatorname{int}\overline B
\subseteq\operatorname{int}\overline A=\varnothing$.

Dla dwóch zbiorów $A,B$ nigdziegęstych
weźmy dowolny niepusty zbiór otwarty $U$.
Ponieważ $A$ jest nigdziegęsty, istnieje niepusty
otwarty $V\subseteq U$ rozłączny z $\overline A$;
możemy przyjąć $V=U\setminus\overline A$.
Analogicznie $W=V\setminus\overline B$ jest
niepusty, otwarty i rozłączny z
$\overline A\cup\overline B=\overline{A\cup B}$.
Stąd $\operatorname{int}\overline{A\cup B}=\varnothing$.
Przez indukcję otrzymujemy tezę dla dowolnej
niepustej rodziny skończonej. Suma rodziny pustej
jest $\varnothing$ i również jest nigdziegęsta.
\end{proof}

\begin{przyklad}[Operacje w topologii koskończonej]
\label{ex:topog-operacje-koskonczona}
Niech $X$ będzie nieskończonym zbiorem z topologią
koskończoną. Jeżeli $A\subseteq X$ jest skończony,
to $A$ jest domknięty, a więc $\overline A=A$.
Jeśli $A$ jest nieskończony, każde niepuste otwarte
$U$ ma skończone dopełnienie i musi spotykać $A$,
więc $\overline A=X$.

Wnętrze $A$ jest równe $A$, jeśli $A$ jest otwarty
(czyli pusty albo ma skończone dopełnienie);
w każdym innym przypadku jest puste.
Jeżeli $A$ jest nieskończony i $X\setminus A$
również jest nieskończony, to
$\operatorname{int}A=\varnothing$ i $\overline A=X$,
a zatem $\partial A=X$.
Dla skończonego $A$ zbiór pochodny jest pusty:
każdy punkt $x$ ma otwarte otoczenie rozłączne
z $A\setminus\{x\}$, mianowicie
$X\setminus(A\setminus\{x\})$.
Jeśli $A$ jest nieskończony, to $A'=X$,
gdyż każde niepuste otwarte otoczenie spotyka
nieskończony zbiór $A\setminus\{x\}$.
\end{przyklad}

\begin{przyklad}[Wnętrze i domknięcie dla topologii wyznaczonej przez podział]
\label{ex:topog-operacje-podzial}
Niech $\mathcal A$ będzie podziałem zbioru $X$
i niech $\tau_{\mathcal A}$ składa się z sum bloków
podziału. Dla $E\subseteq X$ zachodzą wzory
\[
\operatorname{int}E
 =\bigcup\{C\in\mathcal A:C\subseteq E\},\qquad
\overline E
 =\bigcup\{C\in\mathcal A:C\cap E\ne\varnothing\}.
\]
Pierwszy wynika stąd, że każdy otwarty podzbiór $E$
jest sumą bloków w całości zawartych w $E$.
Każdy blok jest jednocześnie otwarty i domknięty:
jego dopełnienie jest sumą pozostałych bloków.
Jeżeli domknięty zbiór $F$ zawiera $E$,
to zawiera cały blok spotykający $E$,
ponieważ $F$ również jest sumą bloków.
Wynika stąd drugi wzór. W szczególności
brzeg $E$ jest sumą bloków, które spotykają
zarówno $E$, jak i jego dopełnienie.
\end{przyklad}

\section{Aksjomaty domknięcia Kuratowskiego}
\label{sec:topog-kuratowski}

Dotychczas definiowaliśmy domknięcie za pomocą zbiorów domkniętych.
Można odwrócić konstrukcję: wskazać własności operatora
domknięcia i odtworzyć z nich całą topologię.

\begin{twierdzenie}[Aksjomaty domknięcia Kuratowskiego]
\label{thm:topog-kuratowski}
Niech $X$ będzie zbiorem. Odwzorowanie
$c\colon\mathcal P(X)\to\mathcal P(X)$ jest operatorem
domknięcia pewnej jednoznacznie wyznaczonej topologii
na $X$ wtedy i tylko wtedy, gdy dla wszystkich $A,B\subseteq X$
zachodzą warunki
\begin{enumerate}[label=(K\arabic*)]
\item $c(\varnothing)=\varnothing$;
\item $A\subseteq c(A)$;
\item $c(A\cup B)=c(A)\cup c(B)$;
\item $c(c(A))=c(A)$.
\end{enumerate}
W takim przypadku zbiory domknięte tej topologii
to dokładnie zbiory $F\subseteq X$ spełniające $c(F)=F$.
\end{twierdzenie}
\begin{proof}
Jeżeli $c(A)=\overline A$ dla ustalonej topologii,
to (K1) wynika z domkniętości zbioru pustego,
(K2) z zawierania $A\subseteq\overline A$,
(K3) z twierdzenia o domknięciu sumy dwóch zbiorów,
a (K4) z domkniętości $\overline A$.

Załóżmy teraz, że odwzorowanie $c$ spełnia (K1)--(K4).
Najpierw wyprowadzimy jego monotoniczność.
Jeżeli $A\subseteq B$, to $B=A\cup B$, więc
$c(B)=c(A)\cup c(B)$ na mocy (K3), a zatem
$c(A)\subseteq c(B)$.
Z (K2) zastosowanego do $X$ wynika $X\subseteq c(X)$,
a ponieważ $c(X)\subseteq X$, otrzymujemy $c(X)=X$.

Niech $\mathcal F=\{F\subseteq X:c(F)=F\}$.
Z (K1) i poprzedniej równości wynika
$\varnothing,X\in\mathcal F$.
Jeżeli $F,G\in\mathcal F$, to
$c(F\cup G)=c(F)\cup c(G)=F\cup G$,
a zatem $F\cup G\in\mathcal F$.
Niech teraz $(F_i)_{i\in I}$ będzie dowolną rodziną
zbiorów z $\mathcal F$ i oznaczmy
$F=\bigcap_{i\in I}F_i$; dla $I=\varnothing$ przyjmujemy $F=X$.
Jeśli $I\ne\varnothing$, to dla każdego $i\in I$
mamy z monotoniczności
$c(F)\subseteq c(F_i)=F_i$.
Wobec tego $c(F)\subseteq\bigcap_{i\in I}F_i=F$,
a z (K2) wynika odwrotne zawieranie.
Dla $I=\varnothing$ już wiemy, że $c(F)=c(X)=X$.
Zatem $F\in\mathcal F$ w obu przypadkach.

Rodzina $\tau=\{X\setminus F:F\in\mathcal F\}$
jest więc topologią: zawiera $\varnothing,X$,
jest zamknięta na dowolne sumy (z praw de Morgana)
i na skończone przecięcia.
Pozostaje sprawdzić, że jej operator domknięcia
jest dokładnie równy $c$.
Na mocy (K4) zbiór $c(A)$ należy do $\mathcal F$,
a (K2) daje $A\subseteq c(A)$. Jeśli $F\in\mathcal F$
i $A\subseteq F$, to monotoniczność daje
$c(A)\subseteq c(F)=F$.
Wobec tego $c(A)$ jest najmniejszym domkniętym
zbiorem zawierającym $A$, a więc jego domknięciem
w skonstruowanej topologii.

Jednoznaczność wynika z charakteryzacji zbiorów domkniętych:
w każdej topologii zbiór jest domknięty wtedy
i tylko wtedy, gdy jest równy własnemu domknięciu.
Dwie topologie mające ten sam operator $c$ muszą więc
mieć te same zbiory domknięte, a tym samym te same otwarte.
\end{proof}

\begin{przyklad}[Dwa skrajne operatory domknięcia]
\label{ex:topog-operatory-skr}
Na dowolnym zbiorze $X$ operator $c_{\mathrm d}(A)=A$
spełnia aksjomaty (K1)--(K4) i wyznacza topologię dyskretną.
Operator
\[
c_{\mathrm a}(A)=
\begin{cases}
\varnothing,&A=\varnothing,\\
X,&A\ne\varnothing
\end{cases}
\]
również spełnia te aksjomaty: (K3) wynika z faktu,
że $A\cup B$ jest niepusty dokładnie wtedy,
gdy przynajmniej jeden z $A,B$ jest niepusty.
Wyznacza on topologię antydyskretną.
Oba przykłady pozostają poprawne dla $X=\varnothing$.
\end{przyklad}

\section{Podsumowanie i powiązania}
\label{sec:topog-operacje-podsumowanie}

Dla dowolnego podzbioru przestrzeni topologicznej
wnętrze jest jego największym otwartym podzbiorem,
a domknięcie najmniejszym domkniętym nadzbiorem.
Brzeg jest przecięciem domknięć zbioru i jego dopełnienia;
zewnętrze jest wnętrzem dopełnienia. Zbiór pochodny
odróżnia punkty skupienia od punktów izolowanych,
przy czym $\overline A=A\cup A'$ bez jakichkolwiek
aksjomatów oddzielania, natomiast domkniętość
$A'$ wymaga dodatkowych założeń, np. $T_1$.
Gęstość oraz nigdziegęstość można badać przez
otwarte otoczenia, a aksjomaty Kuratowskiego
pozwalają odtworzyć topologię z samej operacji domknięcia.

W następnym rozdziale użyjemy tych narzędzi do badania
ciągłości odwzorowań i homeomorfizmów, a w dalszych
rozdziałach do opisu zbieżności, zwartości i oddzielania.

\chapter{Odwzorowania ciągłe}
\label{ch:topog-odwzorowania-ciagle}

\noindent\textbf{Cel rozdziału.}
W rozdziałach poprzednich opisaliśmy przestrzeń przez zbiory otwarte
oraz przez operacje wnętrza, domknięcia i brzegu.
Teraz zbadamy odwzorowania zachowujące tę strukturę.
Wskażemy równoważne kryteria ciągłości, ich postać lokalną i metryczną,
wyprowadzimy kryteria za pomocą baz i podbaz oraz wykażemy
własności składania. Następnie rozróżnimy odwzorowania ciągłe,
otwarte, domknięte, homeomorfizmy i zanurzenia topologiczne.
Szczególnie ważne będą kontrprzykłady: ciągła bijekcja nie musi mieć
ciągłej odwrotności, a ciągłe odwzorowanie nie musi być ani otwarte,
ani domknięte. Nie zakładamy metryzowalności ani żadnych aksjomatów
oddzielania, chyba że wyraźnie je wymieniono.

\section{Ciągłość jako zachowanie zbiorów otwartych}
\label{sec:topog-ciaglosc-definicja}

\begin{definicja}[Odwzorowanie ciągłe]
\label{def:topog-odwzorowanie-ciagle}
Niech $(X,\tau_X)$ i $(Y,\tau_Y)$ będą przestrzeniami topologicznymi.
Odwzorowanie $f\colon X\to Y$ nazywamy \emph{ciągłym}, jeżeli dla
każdego zbioru otwartego $V\subseteq Y$ przeciwobraz
$f^{-1}(V)=\{x\in X:f(x)\in V\}$ jest otwarty w $X$.
Symbol $f^{-1}(V)$ oznacza tu przeciwobraz zbioru, nawet jeśli $f$
nie jest bijekcją. W definicji dopuszczamy $X=\varnothing$
i $Y=\varnothing$, o ile istnieje rozważane odwzorowanie.
\end{definicja}

\begin{uwaga}[Dlaczego przeciwobrazy, a nie obrazy?]
Dla dowolnego odwzorowania $f\colon X\to Y$ i rodziny
$(V_i)_{i\in I}$ podzbiorów $Y$ zachodzą równości
\begin{align*}
f^{-1}\Bigl(\bigcup_{i\in I}V_i\Bigr)
 &=\bigcup_{i\in I}f^{-1}(V_i), &
f^{-1}\Bigl(\bigcap_{i\in I}V_i\Bigr)
 &=\bigcap_{i\in I}f^{-1}(V_i).
\end{align*}
Ponadto $f^{-1}(\varnothing)=\varnothing$ i $f^{-1}(Y)=X$.
Prawa te wynikają bezpośrednio z definicji przeciwobrazu i działają
również dla rodzin pustych. Dzięki temu warunek ciągłości jest
zgodny z aksjomatami topologii. Dla obrazów ogólnego odwzorowania
nie zachodzi analogiczna równość dla przecięć.
\end{uwaga}

\begin{przyklad}[Odwzorowania stałe, identycznościowe i zbiory puste]
\label{ex:topog-ciaglosc-elementarna}
Dla każdej przestrzeni topologicznej $X$ odwzorowanie
$\operatorname{id}_X\colon X\to X$, $\operatorname{id}_X(x)=x$,
jest ciągłe: przeciwobrazem zbioru otwartego jest ten sam zbiór.
Jeżeli $Y$ jest przestrzenią topologiczną, $y_0\in Y$, to
$f\colon X\to Y$, $f(x)=y_0$, jest ciągłe, bo dla dowolnego
$V\subseteq Y$ jego przeciwobrazem jest $X$ albo $\varnothing$.
Każde odwzorowanie ze zbioru pustego do $Y$ jest ciągłe,
ponieważ wszystkie przeciwobrazy są puste. Odwzorowanie
$X\to\varnothing$ istnieje tylko wtedy, gdy $X=\varnothing$.
\end{przyklad}

\begin{przyklad}[Zależność ciągłości od topologii]
\label{ex:topog-ciaglosc-skrajne}
Jeżeli $X$ ma topologię dyskretną, to każde odwzorowanie
$f\colon X\to Y$ jest ciągłe dla dowolnej przestrzeni $Y$:
wszystkie podzbiory $X$ są otwarte.
Jeżeli $Y$ ma topologię antydyskretną, to każde
$f\colon X\to Y$ jest ciągłe, ponieważ jedyne zbiory otwarte
w $Y$ to $\varnothing$ i $Y$.

Niech $S=\{0,1\}$. Odwzorowanie identycznościowe
$\operatorname{id}\colon(S,\mathcal P(S))\to
(S,\{\varnothing,S\})$ jest ciągłą bijekcją. Odwrotne
odwzorowanie nie jest ciągłe: przeciwobraz otwartego zbioru
$\{0\}$ przestrzeni dyskretnej nie jest otwarty w przestrzeni
antydyskretnej.
\end{przyklad}

\begin{twierdzenie}[Równoważne globalne kryteria ciągłości]
\label{thm:topog-ciaglosc-kryteria-globalne}
Niech $X,Y$ będą przestrzeniami topologicznymi, a
$f\colon X\to Y$ odwzorowaniem. Następujące warunki są
równoważne:
\begin{enumerate}[label=(\roman*)]
\item $f$ jest ciągłe;
\item przeciwobraz każdego zbioru domkniętego w $Y$
      jest domknięty w $X$;
\item dla każdego $A\subseteq X$ zachodzi
      $f(\overline A^{\,X})\subseteq\overline{f(A)}^{\,Y}$;
\item dla każdego $B\subseteq Y$ zachodzi
      $\overline{f^{-1}(B)}^{\,X}\subseteq
        f^{-1}(\overline B^{\,Y})$.
\end{enumerate}
\end{twierdzenie}
\begin{proof}
Załóżmy (i). Jeżeli $F\subseteq Y$ jest domknięty, to
$Y\setminus F$ jest otwarty. Z równości
$X\setminus f^{-1}(F)=f^{-1}(Y\setminus F)$ wynika,
że $f^{-1}(F)$ jest domknięty. Otrzymujemy (ii).

Załóżmy (ii), ustalmy $A\subseteq X$ i oznaczmy
$F=\overline{f(A)}^{\,Y}$. Zbiór $F$ jest domknięty,
więc $f^{-1}(F)$ jest domknięty w $X$. Ponieważ
$A\subseteq f^{-1}(F)$, dostajemy
$\overline A^{\,X}\subseteq f^{-1}(F)$, co daje (iii).

Załóżmy (iii) i wybierzmy $B\subseteq Y$.
Stosując (iii) do $A=f^{-1}(B)$, otrzymujemy
\[
f(\overline{f^{-1}(B)}^{\,X})
 \subseteq\overline{f(f^{-1}(B))}^{\,Y}
 \subseteq\overline B^{\,Y}.
\]
Przeciwobraz ostatniego zawierania daje (iv).

Załóżmy (iv) i niech $F\subseteq Y$ będzie domknięty.
Ponieważ $\overline F^{\,Y}=F$, mamy
$\overline{f^{-1}(F)}^{\,X}\subseteq f^{-1}(F)$.
Drugie zawieranie wynika z ekstensywności domknięcia,
więc $f^{-1}(F)$ jest domknięty. Otrzymujemy (ii),
a z (ii) wynika (i), ponieważ dla otwartego $V\subseteq Y$
zbiór $f^{-1}(V)$ jest dopełnieniem domkniętego
$f^{-1}(Y\setminus V)$. Dowodzi to równoważności.
\end{proof}

\begin{stwierdzenie}[Ciągłość i wnętrze]
\label{prop:topog-ciaglosc-wnetrze}
Jeżeli $f\colon X\to Y$ jest ciągłe, to dla każdego
$B\subseteq Y$ zachodzi
$f^{-1}(\operatorname{int}_Y B)
 \subseteq\operatorname{int}_X(f^{-1}(B))$.
\end{stwierdzenie}
\begin{proof}
Zbiór $\operatorname{int}_Y B$ jest otwarty i zawarty w $B$.
Z ciągłości jego przeciwobraz jest otwarty i zawarty w
$f^{-1}(B)$. Ponieważ wnętrze jest największym otwartym
podzbiorem, dostajemy tezę.
\end{proof}

\section{Ciągłość w punkcie i otoczenia}
\label{sec:topog-ciaglosc-lokalna}

\begin{definicja}[Ciągłość w punkcie]
\label{def:topog-ciaglosc-w-punkcie}
Niech $f\colon X\to Y$ będzie odwzorowaniem przestrzeni
topologicznych, a $x\in X$. Mówimy, że $f$ jest
\emph{ciągłe w punkcie $x$}, jeśli dla każdego otwartego
zbioru $V\subseteq Y$ zawierającego $f(x)$ istnieje
otwarty zbiór $U\subseteq X$ taki, że
$x\in U$ i $f(U)\subseteq V$.
\end{definicja}

\begin{twierdzenie}[Kryterium punktowe i otoczeniowe]
\label{thm:topog-ciaglosc-punktowa-globalna}
Dla odwzorowania $f\colon X\to Y$ przestrzeni topologicznych
następujące warunki są równoważne:
\begin{enumerate}[label=(\roman*)]
\item $f$ jest ciągłe;
\item $f$ jest ciągłe w każdym punkcie $x\in X$;
\item dla każdego $x\in X$ i każdego otoczenia $N$
      punktu $f(x)$ zbiór $f^{-1}(N)$ jest otoczeniem $x$.
\end{enumerate}
\end{twierdzenie}
\begin{proof}
Załóżmy (i), ustalmy $x\in X$ i otwarty $V\ni f(x)$.
Przeciwobraz $U=f^{-1}(V)$ jest otwarty, zawiera $x$
i spełnia $f(U)\subseteq V$. Mamy (ii).

Załóżmy (ii). Niech $N$ będzie otoczeniem $f(x)$.
Z definicji otoczenia istnieje otwarty zbiór $V$ taki,
że $f(x)\in V\subseteq N$. Ciągłość w $x$ daje otwarty
$U\ni x$ z $f(U)\subseteq V$. Wtedy
$x\in U\subseteq f^{-1}(N)$, a zatem $f^{-1}(N)$
jest otoczeniem $x$. Otrzymujemy (iii).

Załóżmy (iii) i niech $V\subseteq Y$ będzie otwarty.
Dla każdego $x\in f^{-1}(V)$ zbiór $V$ jest otoczeniem
$f(x)$, więc $f^{-1}(V)$ jest otoczeniem $x$.
Z kryterium otwartości przez otoczenia wynika,
że $f^{-1}(V)$ jest otwarty. Jeśli $f^{-1}(V)=\varnothing$,
jest otwarty na mocy aksjomatów topologii.
Zatem zachodzi (i).
\end{proof}

\begin{twierdzenie}[Metryczne kryterium $\varepsilon$--$\delta$]
\label{thm:topog-ciaglosc-eps-delta}
Niech $(X,d_X)$ i $(Y,d_Y)$ będą przestrzeniami
metrycznymi, wyposażonymi w topologie wyznaczone przez metryki.
Dla $f\colon X\to Y$ i $x\in X$ następujące warunki
są równoważne:
\begin{enumerate}[label=(\roman*)]
\item $f$ jest ciągłe w $x$ w sensie topologicznym;
\item dla każdego $\varepsilon>0$ istnieje $\delta>0$
      takie, że dla wszystkich $z\in X$ z
      $d_X(z,x)<\delta$ zachodzi
      $d_Y(f(z),f(x))<\varepsilon$.
\end{enumerate}
\end{twierdzenie}
\begin{proof}
Załóżmy (i) i ustalmy $\varepsilon>0$.
Kula $B_{d_Y}(f(x),\varepsilon)$ jest otwarta,
więc ciągłość w $x$ daje otwarty $U\ni x$ taki, że
$f(U)\subseteq B_{d_Y}(f(x),\varepsilon)$.
Ponieważ kule metryczne tworzą bazę topologii $X$,
istnieje $\delta>0$ spełniające $B_{d_X}(x,\delta)\subseteq U$.
To jest dokładnie (ii).

Załóżmy (ii). Jeśli $V\subseteq Y$ jest otwarty i
$f(x)\in V$, to istnieje $\varepsilon>0$ takie, że
$B_{d_Y}(f(x),\varepsilon)\subseteq V$.
Dla wybranego z (ii) $\delta>0$ mamy
$f(B_{d_X}(x,\delta))\subseteq V$, więc $f$
jest ciągłe w $x$ w sensie topologicznym.
\end{proof}

\begin{przyklad}[Moduł na prostej]
\label{ex:topog-ciaglosc-modul}
Odwzorowanie $f\colon\R\to\R$, $f(x)=|x|$,
jest ciągłe dla topologii euklidesowych. Dla dowolnych
$x,z\in\R$ z nierówności trójkąta wynikają
$|z|\le|z-x|+|x|$ oraz $|x|\le|x-z|+|z|$.
Zatem $\bigl||z|-|x|\bigr|\le|z-x|$.
Dla danego $\varepsilon>0$ wystarczy przyjąć
$\delta=\varepsilon$ w kryterium $\varepsilon$--$\delta$.
\end{przyklad}

\begin{przyklad}[Funkcja nieciągła]
\label{ex:topog-funkcja-nieciagla}
Na $\R$ z topologią euklidesową rozważmy
$f\colon\R\to\R$ dane wzorem
$f(x)=0$ dla $x\le0$ i $f(x)=1$ dla $x>0$.
Zbiór $V=(1/2,3/2)$ jest otwarty w $\R$, natomiast
$f^{-1}(V)=(0,\infty)$ jest otwarty; to samo
nie wystarcza jeszcze do wykazania ciągłości.
Dla otwartego $W=(-1/2,1/2)$ mamy
$f^{-1}(W)=(-\infty,0]$, który nie jest otwarty.
Zatem $f$ nie jest ciągłe. Bezpośrednio w punkcie
$0$ każde otwarte otoczenie zawiera dodatnie liczby,
na których $f$ ma wartość $1$, więc nie można zawrzeć
jego obrazu w $W$.
\end{przyklad}

\section{Ciągłość sprawdzana na bazach i podbazach}
\label{sec:topog-ciaglosc-bazy}

\begin{twierdzenie}[Kryterium bazy i podbazy w przeciwdziedzinie]
\label{thm:topog-ciaglosc-bazy-podbazy}
Niech $f\colon X\to Y$ będzie odwzorowaniem przestrzeni
topologicznych.
\begin{enumerate}[label=(\roman*)]
\item Jeżeli $\mathcal B$ jest bazą topologii na $Y$,
      to $f$ jest ciągłe wtedy i tylko wtedy, gdy
      $f^{-1}(B)$ jest otwarty w $X$ dla każdego
      $B\in\mathcal B$.
\item Jeżeli $\mathcal S$ jest podbazą topologii na $Y$,
      to $f$ jest ciągłe wtedy i tylko wtedy, gdy
      $f^{-1}(S)$ jest otwarty w $X$ dla każdego
      $S\in\mathcal S$.
\end{enumerate}
\end{twierdzenie}
\begin{proof}
Jeśli $f$ jest ciągłe, obie tezy są oczywiste,
gdyż elementy bazy i podbazy są otwarte.

Załóżmy, że każdy $f^{-1}(B)$ dla $B\in\mathcal B$
jest otwarty. Każdy otwarty $V\subseteq Y$ jest sumą
wszystkich elementów bazy zawartych w $V$. Zatem
\[
f^{-1}(V)=\bigcup_{\substack{B\in\mathcal B\\B\subseteq V}}
 f^{-1}(B),
\]
a więc jest otwarty. Dotyczy to również $V=\varnothing$,
niezależnie od tego, czy $\varnothing\in\mathcal B$.
Dowodzi to (i).

Jeśli przeciwobrazy wszystkich $S\in\mathcal S$
są otwarte, to przeciwobraz każdego skończonego
przecięcia elementów $\mathcal S$ jest otwarty.
Przecięcie pustej rodziny to $Y$, którego przeciwobrazem
jest $X$. Takie skończone przecięcia tworzą bazę
z definicji topologii generowanej przez podbazę.
Z punktu (i) wynika ciągłość, co kończy dowód (ii).
\end{proof}

\begin{przyklad}[Topologia prostej Sorgenfreya]
\label{ex:topog-ciaglosc-sorgenfrey}
Niech $\tau_e$ oznacza topologię euklidesową na $\R$,
a $\tau_{\mathrm{So}}$ topologię generowaną przez
przedziały $[a,b)$, $a<b$. Z rozdziału~1 wiemy, że
$\tau_e\subsetneq\tau_{\mathrm{So}}$.
Odwzorowanie identycznościowe
$\operatorname{id}\colon(\R,\tau_{\mathrm{So}})
\to(\R,\tau_e)$ jest ciągłe, ponieważ każdy zbiór
otwarty w $\tau_e$ jest otwarty także w $\tau_{\mathrm{So}}$.
Odwzorowanie identycznościowe w przeciwnym kierunku
nie jest ciągłe: przeciwobraz $[0,1)$ byłby tym samym
zbiorem, a nie jest on euklidesowo otwarty.
\end{przyklad}

\begin{stwierdzenie}[Porównanie topologii przez identyczność]
\label{prop:topog-porownanie-przez-ciaglosc}
Niech $\tau_1,\tau_2$ będą topologiami na tym samym zbiorze
$X$. Odwzorowanie identycznościowe
$\operatorname{id}\colon(X,\tau_1)\to(X,\tau_2)$
jest ciągłe wtedy i tylko wtedy, gdy $\tau_2\subseteq\tau_1$.
\end{stwierdzenie}
\begin{proof}
Dla każdego $V\subseteq X$ zachodzi
$\operatorname{id}^{-1}(V)=V$. Warunek ciągłości
mówi więc dokładnie, że każdy zbiór należący do
$\tau_2$ należy do $\tau_1$.
\end{proof}

\begin{przyklad}[Przestrzeń Sierpińskiego wykrywa otwartość]
\label{ex:topog-sierpinski-klasyfikator}
Niech $S=\{0,1\}$ ma topologię
$\{\varnothing,\{1\},S\}$, a $A\subseteq X$.
Określmy odwzorowanie $\chi_A\colon X\to S$ wzorem
$\chi_A(x)=1$ dla $x\in A$ i $\chi_A(x)=0$ dla $x\notin A$.
Ponieważ $\chi_A^{-1}(\{1\})=A$,
a przeciwobrazy dwóch pozostałych otwartych zbiorów
są zawsze otwarte, odwzorowanie $\chi_A$ jest ciągłe
dokładnie wtedy, gdy $A$ jest otwarty w $X$.
W szczególności przykład działa także dla $X=\varnothing$.
\end{przyklad}

\section{Składanie, ograniczenia i sklejanie}
\label{sec:topog-ciaglosc-konstrukcje}

\begin{twierdzenie}[Składanie odwzorowań ciągłych]
\label{thm:topog-skladanie-ciaglych}
Niech $X,Y,Z$ będą przestrzeniami topologicznymi, a
$f\colon X\to Y$, $g\colon Y\to Z$ odwzorowaniami
ciągłymi. Wtedy $g\circ f\colon X\to Z$ jest ciągłe.
\end{twierdzenie}
\begin{proof}
Jeżeli $W\subseteq Z$ jest otwarty, to $g^{-1}(W)$
jest otwarty w $Y$, a więc $f^{-1}(g^{-1}(W))$
jest otwarty w $X$. Ponieważ
$(g\circ f)^{-1}(W)=f^{-1}(g^{-1}(W))$,
zachodzi definicja ciągłości złożenia.
\end{proof}

\begin{definicja}[Topologia podprzestrzeni]
\label{def:topog-topologia-podprzestrzeni-wstep}
Jeśli $(X,\tau)$ jest przestrzenią topologiczną,
a $A\subseteq X$, to na $A$ określamy rodzinę
$\tau_A=\{A\cap U:U\in\tau\}$. Jest ona topologią:
zawiera $\varnothing$ i $A$, zaś przecięcia ze zbiorem
$A$ zachowują dowolne sumy i skończone przecięcia.
Nazywamy ją \emph{topologią podprzestrzeni}.
Konstrukcję tę zbadamy szczegółowo w rozdziale~4.
\end{definicja}

\begin{stwierdzenie}[Ograniczanie dziedziny i przeciwdziedziny]
\label{prop:topog-ograniczanie-ciaglosci}
Niech $f\colon X\to Y$ będzie odwzorowaniem ciągłym.
\begin{enumerate}[label=(\roman*)]
\item Dla każdego $A\subseteq X$ ograniczenie
      $f|_A\colon A\to Y$ jest ciągłe, gdy $A$
      ma topologię podprzestrzeni.
\item Jeżeli $B\subseteq Y$ i $f(X)\subseteq B$,
      to odwzorowanie $\widetilde f\colon X\to B$,
      $\widetilde f(x)=f(x)$, jest ciągłe,
      gdy $B$ ma topologię podprzestrzeni.
\end{enumerate}
\end{stwierdzenie}
\begin{proof}
Dla dowolnego otwartego $V\subseteq Y$ mamy
$(f|_A)^{-1}(V)=A\cap f^{-1}(V)$, co jest otwarte
w podprzestrzeni $A$. Dowodzi to (i).

Jeżeli $W\subseteq B$ jest otwarty w $B$,
to ma postać $B\cap V$ z $V$ otwartym w $Y$.
Ponieważ $f(X)\subseteq B$, zachodzi
$\widetilde f^{-1}(W)=f^{-1}(V)$, zbiór otwarty
w $X$. Otrzymujemy (ii).
\end{proof}

\begin{twierdzenie}[Lokalność ciągłości na otwartym pokryciu]
\label{thm:topog-ciaglosc-lokalna-pokrycie}
Niech $(U_i)_{i\in I}$ będzie rodziną otwartych
podzbiorów przestrzeni topologicznej $X$ taką, że
$X=\bigcup_{i\in I}U_i$. Niech $f\colon X\to Y$ będzie
odwzorowaniem do przestrzeni topologicznej $Y$.
Wówczas $f$ jest ciągłe wtedy i tylko wtedy, gdy
dla każdego $i\in I$ ograniczenie
$f|_{U_i}\colon U_i\to Y$ jest ciągłe.
\end{twierdzenie}
\begin{proof}
Jeśli $f$ jest ciągłe, ciągłość wszystkich ograniczeń
wynika z poprzedniego stwierdzenia.
Odwrotnie, niech $V\subseteq Y$ będzie otwarty.
Dla każdego $i$ zbiór
$U_i\cap f^{-1}(V)=(f|_{U_i})^{-1}(V)$ jest otwarty
w $U_i$. Ponieważ $U_i$ jest otwarty w $X$,
każdy zbiór otwarty w $U_i$ jest także otwarty w $X$.
Zatem $f^{-1}(V)=\bigcup_{i\in I}(U_i\cap f^{-1}(V))$
jest otwarty w $X$. Dla $I=\varnothing$ pokrycie
wymusza $X=\varnothing$ i ten sam argument działa.
\end{proof}

\begin{twierdzenie}[Lemat o sklejaniu na pokryciu skończonym domkniętym]
\label{thm:topog-sklejanie-domkniete}
Niech $X$ i $Y$ będą przestrzeniami topologicznymi,
a $F_1,\ldots,F_m$ domkniętymi podzbiorami $X$, gdzie
$m\ge1$ i $X=F_1\cup\cdots\cup F_m$.
Jeżeli $f\colon X\to Y$ jest odwzorowaniem takim, że
$f|_{F_j}\colon F_j\to Y$ jest ciągłe dla każdego
$j=1,\ldots,m$, to $f$ jest ciągłe.
\end{twierdzenie}
\begin{proof}
Niech $C\subseteq Y$ będzie domknięty. Dla każdego
$j$ zbiór $F_j\cap f^{-1}(C)$ jest domknięty
w podprzestrzeni $F_j$. Jego dopełnienie w $F_j$ jest
otwarte w $F_j$, czyli ma postać $F_j\cap U_j$
dla pewnego otwartego $U_j\subseteq X$.
Zatem $F_j\cap f^{-1}(C)=F_j\cap E_j$,
gdzie $E_j=X\setminus U_j$ jest domknięty w $X$.
Ponieważ $F_j$ jest domknięty w $X$, zbiór
$F_j\cap E_j$ jest domknięty w $X$.
Równość $X=\bigcup_{j=1}^{m}F_j$ daje
\[
f^{-1}(C)=\bigcup_{j=1}^{m}(F_j\cap f^{-1}(C)),
\]
więc $f^{-1}(C)$ jest domknięty jako skończona suma
zbiorów domkniętych. Ciągłość wynika z kryterium
przeciwobrazów zbiorów domkniętych.
\end{proof}

\begin{przyklad}[Dlaczego skończoność pokrycia domkniętego jest istotna?]
\label{ex:topog-sklejanie-nieskonczone}
Niech $X=Y=\R$ mają topologię euklidesową i niech
$f(x)=1$ dla $x=0$ oraz $f(x)=0$ dla $x\ne0$.
Funkcja nie jest ciągła, ponieważ
$f^{-1}((1/2,3/2))=\{0\}$ nie jest otwarty.
Tymczasem każdy singleton $F_x=\{x\}$ jest domknięty
w $\R$, rodzina $(F_x)_{x\in\R}$ pokrywa $\R$,
a wszystkie ograniczenia $f|_{F_x}$ są stałe,
więc ciągłe. Nie można więc bez dodatkowych założeń
zastąpić skończonej rodziny domkniętej w lemacie
rodziną dowolną.
\end{przyklad}

\section{Odwzorowania otwarte i domknięte}
\label{sec:topog-otwarte-domkniete-odwzorowania}

\begin{definicja}[Odwzorowanie otwarte i domknięte]
\label{def:topog-otwarte-domkniete-odwzorowanie}
Niech $f\colon X\to Y$ będzie odwzorowaniem przestrzeni
topologicznych. Nazywamy je \emph{otwartym}, jeśli
$f(U)$ jest otwarty w $Y$ dla każdego otwartego
$U\subseteq X$. Nazywamy je \emph{domkniętym},
jeśli $f(F)$ jest domknięty w $Y$ dla każdego
domkniętego $F\subseteq X$.
Żadna z tych definicji nie zawiera założenia ciągłości
ani surjektywności $f$.
\end{definicja}

\begin{stwierdzenie}[Kryteria przez wnętrze i domknięcie]
\label{prop:topog-otwarte-domkniete-kryteria}
Dla dowolnego odwzorowania $f\colon X\to Y$
przestrzeni topologicznych zachodzą równoważności:
\begin{enumerate}[label=(\roman*)]
\item $f$ jest otwarte wtedy i tylko wtedy, gdy
      $f(\operatorname{int}_X A)
        \subseteq\operatorname{int}_Y f(A)$
      dla każdego $A\subseteq X$;
\item $f$ jest domknięte wtedy i tylko wtedy, gdy
      $\overline{f(A)}^{\,Y}\subseteq f(\overline A^{\,X})$
      dla każdego $A\subseteq X$.
\end{enumerate}
\end{stwierdzenie}
\begin{proof}
Jeśli $f$ jest otwarte, to $f(\operatorname{int}A)$
jest otwarty w $Y$ i zawarty w $f(A)$, więc
jest zawarty we wnętrzu $f(A)$. Odwrotnie,
dla otwartego $U\subseteq X$ mamy
$\operatorname{int}_X U=U$, a założenie daje
$f(U)\subseteq\operatorname{int}_Y f(U)$.
Odwrotne zawieranie zawsze zachodzi, zatem
$f(U)$ jest otwarty. Dowodzi to (i).

Jeśli $f$ jest domknięte, to $f(\overline A^{\,X})$
jest domknięty w $Y$ i zawiera $f(A)$, więc
$\overline{f(A)}^{\,Y}\subseteq f(\overline A^{\,X})$.
Odwrotnie, dla dowolnego domkniętego $F\subseteq X$
założenie daje
$\overline{f(F)}^{\,Y}\subseteq f(\overline F^{\,X})=f(F)$.
Druga inkluzja zachodzi zawsze, więc $f(F)$ jest
domknięty. Dowodzi to (ii).
\end{proof}

\begin{przyklad}[Ciągłość nie wymusza otwartości ani domkniętości]
\label{ex:topog-ciaglosc-a-otwartosc-domknietosc}
Niech $S=\{0,1\}$ i
$f\colon(S,\mathcal P(S))\to(S,\{\varnothing,S\})$
będzie odwzorowaniem identycznościowym.
Jak już wykazaliśmy, $f$ jest ciągłe. Nie jest ono
otwarte: zbiór $\{0\}$ jest otwarty w dziedzinie,
a jego obraz nie jest otwarty w przeciwdziedzinie.
Nie jest też domknięte: zbiór $\{0\}$ jest domknięty
w dziedzinie, a jego obraz nie jest domknięty
w topologii antydyskretnej.

W drugą stronę, identyczność
$g\colon(S,\{\varnothing,S\})\to(S,\mathcal P(S))$
jest otwarta i domknięta (obrazami zbiorów otwartych,
odpowiednio domkniętych, są $\varnothing$ i $S$),
ale nie jest ciągła, gdyż przeciwobraz $\{0\}$ nie
jest otwarty w dziedzinie.
\end{przyklad}

\begin{stwierdzenie}[Składanie odwzorowań otwartych i domkniętych]
\label{prop:topog-skladanie-otwartych-domknietych}
Jeżeli $f\colon X\to Y$ i $g\colon Y\to Z$ są
odwzorowaniami otwartymi, to $g\circ f$ jest otwarte.
Jeżeli oba są domknięte, to $g\circ f$ jest domknięte.
\end{stwierdzenie}
\begin{proof}
Dla otwartego $U\subseteq X$ zbiór $f(U)$ jest otwarty
w $Y$, a zatem $g(f(U))$ jest otwarty w $Z$.
Dla domkniętego $F\subseteq X$ analogicznie
$f(F)$ jest domknięty w $Y$, a $g(f(F))$ w $Z$.
\end{proof}

\section{Homeomorfizmy i zanurzenia topologiczne}
\label{sec:topog-homeomorfizmy}

\begin{definicja}[Homeomorfizm]
\label{def:topog-homeomorfizm}
Niech $X,Y$ będą przestrzeniami topologicznymi.
\emph{Homeomorfizmem} $f\colon X\to Y$ nazywamy
odwzorowanie bijektywne, ciągłe i takie, że
odwzorowanie odwrotne $f^{-1}\colon Y\to X$ jest ciągłe.
Przestrzenie, między którymi istnieje homeomorfizm,
nazywamy \emph{homeomorficznymi} i piszemy $X\cong Y$.
\end{definicja}

\begin{twierdzenie}[Ciągła bijekcja: kryteria homeomorfizmu]
\label{thm:topog-homeomorfizm-kryteria}
Niech $f\colon X\to Y$ będzie ciągłą bijekcją
przestrzeni topologicznych. Następujące warunki są
równoważne:
\begin{enumerate}[label=(\roman*)]
\item $f$ jest homeomorfizmem;
\item $f$ jest odwzorowaniem otwartym;
\item $f$ jest odwzorowaniem domkniętym.
\end{enumerate}
\end{twierdzenie}
\begin{proof}
Niech $h=f^{-1}\colon Y\to X$ będzie odwrotnością
w sensie odwzorowania, istniejącą dzięki bijektywności $f$.
Dla dowolnego $U\subseteq X$ zachodzi
$h^{-1}(U)=f(U)$, dlatego $h$ jest ciągłe
dokładnie wtedy, gdy $f$ jest otwarte. To dowodzi
równoważności (i) i (ii).

Podobnie, stosując kryterium przeciwobrazów zbiorów
domkniętych do $h$, otrzymujemy, że $h$ jest ciągłe
dokładnie wtedy, gdy $f(F)$ jest domknięte dla każdego
domkniętego $F\subseteq X$. To jest warunek (iii).
\end{proof}

\begin{przyklad}[Ciągła bijekcja niebędąca homeomorfizmem]
\label{ex:topog-bijekcja-niehomeomorfizm}
Identyczność
$f\colon(S,\mathcal P(S))\to(S,\{\varnothing,S\})$
dla dwuelementowego $S$ jest ciągłą bijekcją,
ale nie jest homeomorfizmem: jej odwrotność nie jest
ciągła. W świetle poprzedniego twierdzenia wynika to
również z tego, że $f$ nie jest otwarte.
Przykład nie wymaga ani metryki, ani nieskończonego zbioru.
\end{przyklad}

\begin{stwierdzenie}[Zachowanie operacji zbiorowych]
\label{prop:topog-homeomorfizm-operacje}
Niech $f\colon X\to Y$ będzie homeomorfizmem przestrzeni
topologicznych. Wówczas dla każdego $A\subseteq X$
zachodzą równości:
\begin{align*}
f(\operatorname{int}_X A)&=\operatorname{int}_Y f(A),\\
f(\overline A^{\,X})&=\overline{f(A)}^{\,Y},\\
f(\partial_X A)&=\partial_Y f(A),\\
f(A')&=(f(A))'.
\end{align*}
W ostatniej równości symbole $A'$ i $(f(A))'$ oznaczają
zbiory pochodne obliczone odpowiednio w $X$ i w $Y$.
W szczególności homeomorfizmy zachowują gęstość,
domkniętość, otwartość i własność bycia punktem izolowanym
podzbioru.
\end{stwierdzenie}
\begin{proof}
Homeomorfizm $f$ i jego odwrotność są otwarte.
Zatem $f(\operatorname{int}_X A)$ jest otwarty
i zawarty w $f(A)$, co daje
$f(\operatorname{int}_X A)\subseteq\operatorname{int}_Y f(A)$.
Tę samą inkluzję dla $f^{-1}$ i $f(A)$
przekształcamy w odwrotną, otrzymując równość.

Ciągłość $f$ daje
$f(\overline A^{\,X})\subseteq\overline{f(A)}^{\,Y}$,
a ciągłość $f^{-1}$ daje inkluzję odwrotną.
Ponieważ $f$ jest bijekcją, mamy
$f(X\setminus A)=Y\setminus f(A)$
i $f(B\cap C)=f(B)\cap f(C)$.
Ze wzoru
$\partial_X A=\overline A^{\,X}\cap
\overline{X\setminus A}^{\,X}$
wynika więc równość dla brzegów.

Dla $x\in X$ warunek $x\in A'$ oznacza,
że każde otwarte otoczenie $U\ni x$
spotyka $A\setminus\{x\}$.
Homeomorfizm przeprowadza otoczenia $x$
bijektywnie na otoczenia $f(x)$, przy czym
\[
f\bigl((U\setminus\{x\})\cap A\bigr)
=(f(U)\setminus\{f(x)\})\cap f(A).
\]
Daje to równoważność
$x\in A'\Longleftrightarrow f(x)\in(f(A))'$.
Zachowanie gęstości, otwartości i domkniętości
wynika z równości dotyczących domknięcia
i wnętrza, natomiast punkt $x\in A$
jest izolowany w $A$ dokładnie wtedy,
gdy $x\notin A'$.
\end{proof}

\begin{definicja}[Zanurzenie topologiczne]
\label{def:topog-zanurzenie-topologiczne}
Niech $X,Y$ będą przestrzeniami topologicznymi.
Odwzorowanie $f\colon X\to Y$ nazywamy
\emph{zanurzeniem topologicznym}, jeżeli odwzorowanie
$f\colon X\to f(X)$ (z topologią podprzestrzeni
odziedziczoną z $Y$) jest homeomorfizmem.
W szczególności zanurzenie topologiczne jest ciągłe
i różnowartościowe, ale nie musi być surjektywne
na całe $Y$.
\end{definicja}

\begin{stwierdzenie}[Kryterium zanurzenia]
\label{prop:topog-zanurzenie-kryterium}
Niech $f\colon X\to Y$ będzie odwzorowaniem
przestrzeni topologicznych. Jest ono zanurzeniem
topologicznym wtedy i tylko wtedy, gdy jest różnowartościowe,
ciągłe i dla każdego otwartego $U\subseteq X$ zbiór
$f(U)$ jest otwarty w podprzestrzeni $f(X)\subseteq Y$.
\end{stwierdzenie}
\begin{proof}
Jeśli $f$ jest zanurzeniem, to jako homeomorfizm
$X\to f(X)$ jest różnowartościowe, ciągłe i otwarte
do $f(X)$, więc ma podane własności.
Odwrotnie, jeżeli $f$ jest różnowartościowe,
to $f\colon X\to f(X)$ jest bijekcją.
Jest ono ciągłe, ponieważ dla otwartego
$W=f(X)\cap V$ w $f(X)$, z $V$ otwartym w $Y$,
mamy $f^{-1}(W)=f^{-1}(V)$, otwarty w $X$.
Z założenia jest także otwarte do $f(X)$.
Kryterium homeomorfizmu dowodzi tezy.
\end{proof}

\begin{przyklad}[Naturalne włożenie i kontrprzykład]
\label{ex:topog-zanurzenie-przyklady}
Włożenie podprzestrzeni $A\hookrightarrow Y$ jest zanurzeniem.
Ciągła różnowartościowa tożsamość z
Przykładu~\ref{ex:topog-bijekcja-niehomeomorfizm}
nie jest nim, gdyż jej odwrotność nie jest ciągła.
\end{przyklad}

\section{Podsumowanie i zależności}
\label{sec:topog-ciaglosc-podsumowanie}
Ciągłość opisują przeciwobrazy zbiorów otwartych
i domkniętych, zachowanie domknięcia, otoczenia,
a w przestrzeniach metrycznych warunek
$\varepsilon$--$\delta$; wystarcza baza lub podbaza.
Zachowują ją składanie, ograniczanie i sklejanie
na pokryciach otwartych lub skończonych domkniętych.
Otwartość i domkniętość nie gwarantują ciągłości.
Homeomorfizm ma ciągłą odwrotność,
zanurzenie jest homeomorfizmem na obraz.
Dalej rozpatrujemy podprzestrzenie, produkty,
sumy rozłączne, topologie początkowe i końcowe,
zwartość, spójność oraz zbieżność.

\chapter{Konstrukcje przestrzeni topologicznych}
\label{ch:topog-konstrukcje}

\noindent\textbf{Cel rozdziału.}
Dotychczas rozpatrywaliśmy gotowe przestrzenie topologiczne i odwzorowania
między nimi. Teraz pokażemy, jak z danych przestrzeni otrzymywać nowe.
Konstrukcje będą określane przez wymaganie ciągłości naturalnych odwzorowań.
Najpierw wprowadzimy topologię podprzestrzeni, następnie iloczynu i sumy
rozłącznej, a na końcu topologie początkowe i końcowe, w szczególności
topologię ilorazową. Wyprowadzimy jawne kryteria otwartości i własności
uniwersalne, nie zakładając metryzowalności ani aksjomatów oddzielania.

\section{Topologia podprzestrzeni}
\label{sec:topog-podprzestrzen}

\begin{definicja}[Topologia podprzestrzeni]
\label{def:topog-podprzestrzen}
Niech $(X,\tau)$ będzie przestrzenią topologiczną i niech $A\subseteq X$.
\emph{Topologią podprzestrzeni} na $A$ nazywamy rodzinę
\[
\tau|_A=\{A\cap U:U\in\tau\}.
\]
Elementy $\tau|_A$ nazywamy zbiorami otwartymi \emph{w $A$}.
Dotyczy to również $A=\varnothing$.
\end{definicja}

\begin{stwierdzenie}[Aksjomaty topologii podprzestrzeni]
\label{prop:topog-podprzestrzen-topologia}
Dla każdej przestrzeni topologicznej $(X,\tau)$ i każdego $A\subseteq X$
rodzina $\tau|_A$ jest topologią na $A$.
\end{stwierdzenie}
\begin{proof}
Z równości $A\cap\varnothing=\varnothing$ i $A\cap X=A$
wynika obecność $\varnothing$ i $A$. Dla dowolnej rodziny $(U_i)_{i\in I}$
zbiorów otwartych w $X$ mamy
\[
\bigcup_{i\in I}(A\cap U_i)=A\cap\bigcup_{i\in I}U_i,
\]
a dla $U,V\in\tau$ zachodzi
$(A\cap U)\cap(A\cap V)=A\cap(U\cap V)$.
Prawa strona każdej równości należy do $\tau|_A$, ponieważ $\tau$
jest topologią. Argument działa także dla pustej rodziny.
\end{proof}

\begin{przyklad}[Otwarty względnie, lecz nie w całej przestrzeni]
\label{ex:topog-otwarty-wzglednie}
Niech $X=\mathbb R$ z topologią euklidesową i $A=[0,2]$.
Zbiór $[0,1)$ jest otwarty w $A$, ponieważ
$[0,1)=A\cap(-1,1)$, ale nie jest otwarty w $\mathbb R$.
Podobnie $\{0\}$ jest otwarty w podprzestrzeni
$\{0\}\cup[2,3]$, chociaż nie jest otwarty w $\mathbb R$.
\end{przyklad}

\begin{figure}[htbp]
\centering
\begin{tikzpicture}[>=Stealth,scale=1]
 \draw[->] (-2.25,0)--(4.15,0) node[right] {$\mathbb R$};
 \foreach \x in {-2,-1,0,1,2,3} \draw (\x,.08)--(\x,-.08);
 \draw[line width=2.8pt,blue!55] (0,0)--(2,0);
 \fill[blue!70!black] (0,0) circle (2.3pt);
 \fill[blue!70!black] (2,0) circle (2.3pt);
 \draw[line width=2.3pt,red!75!black] (0,.34)--(1,.34);
 \fill[red!75!black] (0,.34) circle (2.7pt);
 \draw[fill=white,draw=red!75!black,thick] (1,.34) circle (2.7pt);
 \node[blue!65!black,above] at (2.35,.05) {$A=[0,2]$};
 \node[red!75!black,above] at (.48,.42) {$A\cap(-1,1)=[0,1)$};
 \node[below] at (0,-.11) {$0$};
 \node[below] at (1,-.11) {$1$};
 \node[below] at (2,-.11) {$2$};
\end{tikzpicture}
\caption{Zbiór $[0,1)$ jest otwarty w $A=[0,2]$, a nie w $\mathbb R$.}
\label{fig:topog-podprzestrzen-przedzial}
\end{figure}

\begin{stwierdzenie}[Domkniętość i domknięcie względne]
\label{prop:topog-podprzestrzen-domkniecie}
Niech $A\subseteq X$, gdzie $X$ jest przestrzenią topologiczną.
Dla każdego $B\subseteq A$ zachodzą równoważność i równość:
\begin{enumerate}[label=(\roman*)]
 \item $B$ jest domknięty w $A$ wtedy i tylko wtedy, gdy
       $B=A\cap F$ dla pewnego zbioru domkniętego $F\subseteq X$;
 \item $\overline B^{\,A}=A\cap\overline B^{\,X}$.
\end{enumerate}
\end{stwierdzenie}
\begin{proof}
Zbiór $B$ jest domknięty w $A$ dokładnie wtedy, gdy
$A\setminus B=A\cap U$ dla pewnego otwartego $U\subseteq X$.
Biorąc dopełnienia w $A$, otrzymujemy
$B=A\cap(X\setminus U)$, co dowodzi (i).

Niech $C=A\cap\overline B^{\,X}$. Jest on domknięty
w $A$ i zawiera $B$, więc $\overline B^{\,A}\subseteq C$.
Jeśli $D\subseteq A$ jest dowolnym domkniętym podzbiorem $A$
zawierającym $B$, to na mocy (i) istnieje domknięty $F\subseteq X$
taki, że $D=A\cap F$. Z $B\subseteq F$ wynika
$\overline B^{\,X}\subseteq F$, a więc $C\subseteq D$.
Stąd $C$ jest najmniejszym domkniętym podzbiorem $A$
zawierającym $B$, co dowodzi (ii).
\end{proof}

\begin{stwierdzenie}[Wnętrze i brzeg względny]
\label{prop:topog-podprzestrzen-wnetrze}
Dla przestrzeni topologicznej $X$ oraz zbiorów $B\subseteq A\subseteq X$
mamy
\begin{align*}
 \operatorname{int}_A B
 &=A\setminus\overline{A\setminus B}^{\,X},\\
 \partial_A B
 &=A\cap\overline B^{\,X}\cap\overline{A\setminus B}^{\,X}.
\end{align*}
\end{stwierdzenie}
\begin{proof}
W dowolnej przestrzeni wnętrze zbioru jest dopełnieniem
domknięcia jego dopełnienia. Stosując tę własność w $A$ oraz
wzór na domknięcie względne, dostajemy
\[
\operatorname{int}_A B
 =A\setminus\overline{A\setminus B}^{\,A}
 =A\setminus(A\cap\overline{A\setminus B}^{\,X})
 =A\setminus\overline{A\setminus B}^{\,X}.
\]
Brzeg jest przecięciem domknięć zbioru i jego dopełnienia
w rozważanej przestrzeni; zastosowanie poprzedniego
stwierdzenia dwukrotnie daje drugą równość.
\end{proof}

\begin{uwaga}[Wnętrze względne nie jest przecięciem wnętrz]
W ogólności $\operatorname{int}_A B$ nie musi równać się
$A\cap\operatorname{int}_X B$. Dla $X=\mathbb R$,
$A=B=[0,1]$ mamy $\operatorname{int}_A B=A$, podczas gdy
$A\cap\operatorname{int}_{\mathbb R}B=(0,1)$.
\end{uwaga}

\begin{stwierdzenie}[Dziedziczenie topologii i włożenie]
\label{prop:topog-podprzestrzen-przechodniosc}
Niech $B\subseteq A\subseteq X$, gdzie $X$ jest przestrzenią
topologiczną. Topologia na $B$ odziedziczona najpierw z $X$
przez $A$, a następnie z $A$, jest równa topologii
odziedziczonej bezpośrednio z $X$. Ponadto kanoniczne
włożenie $\iota\colon A\to X$, $\iota(a)=a$, jest ciągłe.
\end{stwierdzenie}
\begin{proof}
Zbiór otwarty w $B$ otrzymany dwuetapowo ma postać
$B\cap(A\cap U)=B\cap U$, gdzie $U$ jest otwarty w $X$;
są to dokładnie zbiory topologii odziedziczonej bezpośrednio.
Dla otwartego $U\subseteq X$ przeciwobraz
$\iota^{-1}(U)=A\cap U$ jest otwarty w $A$, zatem
$\iota$ jest ciągłe.
\end{proof}

\begin{twierdzenie}[Ciągłość odwzorowania do podprzestrzeni]
\label{thm:topog-do-podprzestrzeni}
Niech $X,Y$ będą przestrzeniami topologicznymi,
$A\subseteq Y$, a $f\colon X\to A$ dowolnym odwzorowaniem.
Jest ono ciągłe do $A$ z topologią podprzestrzeni wtedy
i tylko wtedy, gdy złożenie $\iota\circ f\colon X\to Y$
z kanonicznym włożeniem $\iota\colon A\to Y$ jest ciągłe.
\end{twierdzenie}
\begin{proof}
Jeśli $f$ jest ciągłe, złożenie z ciągłym włożeniem
jest ciągłe. Odwrotnie, każdy otwarty $W\subseteq A$
ma postać $A\cap V$ dla otwartego $V\subseteq Y$.
Ponieważ wartości $f$ należą do $A$, zachodzi
$f^{-1}(W)=(\iota\circ f)^{-1}(V)$, a zatem
przeciwobraz $W$ jest otwarty.
\end{proof}

\section{Topologia iloczynu}
\label{sec:topog-iloczyn}

\begin{definicja}[Iloczyn przestrzeni topologicznych]
\label{def:topog-iloczyn}
Niech $I$ będzie zbiorem indeksów, a $(X_i,\tau_i)_{i\in I}$
rodziną przestrzeni topologicznych. Na zbiorze iloczynowym
$P=\prod_{i\in I}X_i$ rozważamy odwzorowania
$\pi_i\colon P\to X_i$, $\pi_i((x_j)_{j\in I})=x_i$.
\emph{Topologią iloczynu} na $P$ nazywamy topologię
generowaną przez rodzinę
\[
\mathcal S=\{\pi_i^{-1}(U):i\in I,\ U\in\tau_i\}.
\]
Jest ona zatem najmniejszą topologią, w której wszystkie
rzutowania $\pi_i$ są ciągłe. Dla $I=\varnothing$
przyjmujemy $P=\{()\}$, a topologia iloczynu jest
jedyną topologią na zbiorze jednoelementowym.
\end{definicja}

\begin{twierdzenie}[Baza topologii iloczynu]
\label{thm:topog-iloczyn-baza}
Niech $(X_i)_{i\in I}$ będzie rodziną przestrzeni
topologicznych, a $P=\prod_{i\in I}X_i$.
Bazę topologii iloczynu stanowią wszystkie zbiory
\[
\prod_{i\in I}U_i,\qquad
U_i\subseteq X_i\text{ otwarty},\qquad
U_i=X_i\text{ dla wszystkich poza skończenie wieloma }i.
\]
Dla $I=\varnothing$ zbiór bazowy jest równy $P$.
\end{twierdzenie}
\begin{proof}
Z definicji topologii generowanej przez podbazę jej
bazą są skończone przecięcia elementów $\mathcal S$,
włącznie z przecięciem pustym, które wynosi $P$.
Dla skończonego zbioru indeksów $F\subseteq I$ oraz
otwartych $U_i\subseteq X_i$ mamy
\[
\bigcap_{i\in F}\pi_i^{-1}(U_i)=
\prod_{i\in I}V_i,\qquad
V_i=\begin{cases}U_i,&i\in F,\\X_i,&i\notin F.\end{cases}
\]
Gdy ten sam indeks pojawia się kilkakrotnie, należy
zastąpić odpowiadające mu zbiory otwarte ich skończonym
przecięciem. Odwrotnie, każdy iloczyn z tezy jest takim
skończonym przecięciem przeciwobrazów.
\end{proof}

\begin{przyklad}[Prostokąty w $\mathbb R^2$]
\label{ex:topog-iloczyn-prostokaty}
Dla $X_1=X_2=\mathbb R$ z topologią euklidesową bazą
topologii iloczynu na $\mathbb R^2$ są prostokąty otwarte
$(a,b)\times(c,d)$. Topologia ta pokrywa się z topologią
metryki euklidesowej. Każdy punkt prostokąta ma kulę
euklidesową w nim zawartą, a każda kula euklidesowa
wokół danego punktu zawiera prostokąt otwarty zawierający
ten punkt. W pierwszą stronę wybieramy promień mniejszy
od odległości obu współrzędnych punktu do brzegów
odpowiednich przedziałów; w drugą wystarcza prostokąt
o półbokach $\varepsilon/3$ dla kuli o promieniu
$\varepsilon$, ponieważ odległość punktów w takim
prostokącie nie przekracza $\sqrt 2\,\varepsilon/3$.
\end{przyklad}

\begin{figure}[htbp]
\centering
\begin{tikzpicture}[>=Stealth,scale=.85]
 \draw[->] (-.3,0)--(5.3,0) node[right] {$X_1=\mathbb R$};
 \draw[->] (0,-.3)--(0,3.65) node[above] {$X_2=\mathbb R$};
 \fill[blue!13] (1.1,.7) rectangle (4.0,2.7);
 \draw[blue!70!black,thick,dashed] (1.1,.7) rectangle (4.0,2.7);
 \fill (2.6,1.75) circle (2.2pt);
 \node[above right] at (2.6,1.75) {$(x_1,x_2)$};
 \node[below] at (2.5,.7) {$U_1\times U_2$};
 \draw[blue!70!black,thick] (1.1,-.14)--(4,-.14);
 \draw[blue!70!black,thick] (-.14,.7)--(-.14,2.7);
 \node[below] at (1.1,-.14) {$a$};
 \node[below] at (4,-.14) {$b$};
 \node[left] at (-.14,.7) {$c$};
 \node[left] at (-.14,2.7) {$d$};
\end{tikzpicture}
\caption{Bazowe otoczenie punktu w topologii iloczynu dwóch prostych.}
\label{fig:topog-iloczyn-prostokat}
\end{figure}

\begin{twierdzenie}[Własność uniwersalna iloczynu]
\label{thm:topog-iloczyn-uniwersalnosc}
Niech $(X_i)_{i\in I}$ będzie rodziną przestrzeni
topologicznych, $P=\prod_{i\in I}X_i$ z topologią
iloczynu, a $Z$ dowolną przestrzenią topologiczną.
Odwzorowanie $f\colon Z\to P$ jest ciągłe wtedy i tylko
wtedy, gdy dla każdego $i\in I$ odwzorowanie
$\pi_i\circ f\colon Z\to X_i$ jest ciągłe.
\end{twierdzenie}
\begin{proof}
Jeśli $f$ jest ciągłe, to każde złożenie z ciągłym
rzutowaniem jest ciągłe. Odwrotnie, niech wszystkie
$\pi_i\circ f$ będą ciągłe. Dla dowolnego elementu
podbazy $\pi_i^{-1}(U)$ mamy
\[
f^{-1}(\pi_i^{-1}(U))=(\pi_i\circ f)^{-1}(U),
\]
a prawa strona jest otwarta w $Z$. Z kryterium ciągłości
przez podbazę wynika ciągłość $f$. Dotyczy to również
pustej rodziny indeksów, gdy $P$ jest jednoelementowe.
\end{proof}

\begin{stwierdzenie}[Rzutowania są odwzorowaniami otwartymi]
\label{prop:topog-rzutowania-otwarte}
Niech $P=\prod_{i\in I}X_i$ będzie iloczynem przestrzeni
topologicznych. Dla każdego $j\in I$ rzutowanie
$\pi_j\colon P\to X_j$ jest odwzorowaniem otwartym.
\end{stwierdzenie}
\begin{proof}
Niech $W=\prod_{i\in I}U_i$ będzie bazowym zbiorem
otwartym, z $U_i=X_i$ poza skończonym zbiorem indeksów.
Jeśli $W=\varnothing$, to $\pi_j(W)=\varnothing$ jest
otwarty. W przeciwnym przypadku istnieje punkt
$(x_i)_{i\in I}\in W$. Dla dowolnego $y\in U_j$
rodzina, która ma współrzędną $y$ w miejscu $j$, a
pozostałe współrzędne $x_i$, należy do $W$. Zatem
$\pi_j(W)=U_j$, zbiór otwarty. Każdy zbiór otwarty
w $P$ jest sumą takich zbiorów bazowych, a obraz
zachowuje dowolne sumy. Stąd rzutowanie jest otwarte.
\end{proof}

\begin{stwierdzenie}[Domknięcie prostokąta]
\label{prop:topog-iloczyn-domkniecie}
Niech $X,Y$ będą przestrzeniami topologicznymi oraz
$A\subseteq X$, $B\subseteq Y$. W topologii iloczynu
na $X\times Y$ zachodzi
\[
\overline{A\times B}^{\,X\times Y}
 =\overline A^{\,X}\times\overline B^{\,Y}.
\]
\end{stwierdzenie}
\begin{proof}
Jeśli $A=\varnothing$ lub $B=\varnothing$, obie strony
są puste. W przeciwnym przypadku punkt $(x,y)$ należy
do domknięcia $A\times B$ wtedy i tylko wtedy, gdy każdy
bazowy prostokąt $U\times V$ zawierający $(x,y)$ przecina
$A\times B$. To jest równoważne temu, że każde otwarte
otoczenie $U$ punktu $x$ przecina $A$ i każde otwarte
otoczenie $V$ punktu $y$ przecina $B$.
Z definicji domknięcia dostajemy
$x\in\overline A^{\,X}$ i $y\in\overline B^{\,Y}$.
\end{proof}

\begin{uwaga}[Dlaczego skończenie wiele współrzędnych?]
\label{rem:topog-iloczyn-skrzynkowa}
Przy nieskończonym zbiorze indeksów żądanie, aby
\emph{wszystkie} iloczyny $\prod_i U_i$ z otwartymi
$U_i\subseteq X_i$ były otwarte, może dać inną,
\emph{topologię skrzynkową}. Dla $X_n=\mathbb R$ i
$I=\mathbb N$ zbiór $\prod_{n\ge1}(-1/n,1/n)$ jest
otwarty w topologii skrzynkowej, lecz nie w topologii
iloczynu: każde bazowe otoczenie punktu zerowego
pozostawia wszystkie poza skończenie wieloma współrzędne
bez ograniczeń, więc nie może zawierać się w podanym
zbiorze. Dla skończonej liczby czynników obie topologie
są równe.
\end{uwaga}

\section{Suma rozłączna przestrzeni}
\label{sec:topog-suma-rozlaczna}

\begin{definicja}[Suma rozłączna]
\label{def:topog-suma-rozlaczna}
Niech $(X_i,\tau_i)_{i\in I}$ będzie rodziną przestrzeni
topologicznych. \emph{Sumą rozłączną} jest zbiór
$S=\coprod_{i\in I}X_i=\{(i,x):i\in I,\ x\in X_i\}$
z topologią określoną następująco: $W\subseteq S$ jest
otwarty wtedy i tylko wtedy, gdy dla każdego $i\in I$
zbiór $W_i=\{x\in X_i:(i,x)\in W\}$ jest otwarty
w $X_i$. Naturalne włożenia oznaczamy przez
$\jmath_i\colon X_i\to S$, $\jmath_i(x)=(i,x)$.
Dla $I=\varnothing$ otrzymujemy pustą przestrzeń.
\end{definicja}

\begin{stwierdzenie}[Topologia i składniki sumy rozłącznej]
\label{prop:topog-suma-rozlaczna-topologia}
Definicja \ref{def:topog-suma-rozlaczna} wyznacza topologię
na $S$. Dla każdego $i\in I$ odwzorowanie $\jmath_i$
jest zanurzeniem topologicznym, a jego obraz
$\jmath_i(X_i)$ jest zarówno otwarty, jak i domknięty w $S$.
\end{stwierdzenie}
\begin{proof}
Przekroje zbiorów pustego i całego $S$ są odpowiednio
puste i równe $X_i$. Ponadto przekrój sumy dowolnej
rodziny zbiorów jest sumą ich przekrojów, a przekrój
przecięcia dwóch zbiorów jest przecięciem ich przekrojów.
Aksjomaty topologii wynikają z aksjomatów w każdym $X_i$.

Zbiór $\jmath_i(X_i)$ jest otwarty, ponieważ jego przekrój
w składniku $i$ jest całym $X_i$, a w każdym innym jest
pusty. Jego dopełnienie także jest otwarte, zatem obraz
jest domknięty. Jeśli $U\subseteq X_i$ jest otwarty,
to $\jmath_i(U)$ jest otwarty w $S$, gdyż tylko przekrój
w składniku $i$ jest równy $U$, a pozostałe są puste.
Dla otwartego $W\subseteq S$ przeciwobraz
$\jmath_i^{-1}(W)=W_i$ jest otwarty w $X_i$.
Odwzorowanie $\jmath_i\colon X_i\to\jmath_i(X_i)$
jest więc ciągłą i otwartą bijekcją, czyli homeomorfizmem.
\end{proof}

\begin{twierdzenie}[Własność uniwersalna sumy rozłącznej]
\label{thm:topog-suma-uniwersalnosc}
Niech $S=\coprod_{i\in I}X_i$ będzie sumą rozłączną
przestrzeni topologicznych, $Y$ przestrzenią topologiczną,
a $f\colon S\to Y$ dowolnym odwzorowaniem. Wtedy $f$
jest ciągłe wtedy i tylko wtedy, gdy dla każdego $i\in I$
odwzorowanie $f\circ\jmath_i\colon X_i\to Y$ jest ciągłe.
\end{twierdzenie}
\begin{proof}
Jeśli $f$ jest ciągłe, to wszystkie złożenia z ciągłymi
$\jmath_i$ są ciągłe. Odwrotnie, niech wszystkie złożenia
będą ciągłe. Dla otwartego $V\subseteq Y$ oraz każdego
$i\in I$ mamy
$\jmath_i^{-1}(f^{-1}(V))=(f\circ\jmath_i)^{-1}(V)$,
zatem wszystkie przekroje $f^{-1}(V)$ są otwarte.
Z definicji topologii na $S$ zbiór $f^{-1}(V)$ jest otwarty.
\end{proof}

\begin{przyklad}[Dwie kopie tej samej przestrzeni]
\label{ex:topog-dwie-kopie}
Przestrzeń $\mathbb R\coprod\mathbb R$ składa się z dwóch
rozłącznych kopii prostej. Każda kopia jest otwarta i
domknięta. Nie jest to zwykła suma zbiorów
$\mathbb R\cup\mathbb R=\mathbb R$: znaczniki składników
odróżniają punkty o tej samej współrzędnej.
Nie oznacza to, że suma rozłączna nie może być homeomorficzna
z pewną podprzestrzenią prostej; chodzi o różnicę
między samymi konstrukcjami zbiorów.
Rozłączność początkowych zbiorów $X_i$ nie jest wymagana.
\end{przyklad}

\section{Topologie początkowe i końcowe}
\label{sec:topog-topologie-poczatkowe-koncowe}

\begin{definicja}[Topologia początkowa]
\label{def:topog-topologia-poczatkowa}
Niech $X$ będzie zbiorem, $(Y_i,\tau_i)_{i\in I}$ rodziną
przestrzeni topologicznych, a $f_i\colon X\to Y_i$
odwzorowaniami. \emph{Topologia początkowa} wyznaczona
przez rodzinę $(f_i)_{i\in I}$ jest topologią na $X$
generowaną przez podbazę
\[
\mathcal S=\{f_i^{-1}(U):i\in I,\ U\in\tau_i\}.
\]
Jest to najmniejsza topologia na $X$, dla której wszystkie
$f_i$ są ciągłe. Jeśli $I=\varnothing$, jest to
topologia antydyskretna $\{\varnothing,X\}$,
rozumiana jako jednoelementowa, gdy $X=\varnothing$.
\end{definicja}

\begin{twierdzenie}[Własność uniwersalna topologii początkowej]
\label{thm:topog-poczatkowa-uniwersalnosc}
Niech $X$ ma topologię początkową wyznaczoną przez
$f_i\colon X\to Y_i$, $i\in I$, a $Z$ będzie dowolną
przestrzenią topologiczną. Odwzorowanie $g\colon Z\to X$
jest ciągłe wtedy i tylko wtedy, gdy wszystkie złożenia
$f_i\circ g\colon Z\to Y_i$ są ciągłe.
\end{twierdzenie}
\begin{proof}
Jeden kierunek wynika ze składania odwzorowań ciągłych.
Jeśli wszystkie złożenia są ciągłe, to dla $i\in I$
i otwartego $U\subseteq Y_i$ zachodzi
$g^{-1}(f_i^{-1}(U))=(f_i\circ g)^{-1}(U)$,
a prawa strona jest otwarta w $Z$.
Ciągłość $g$ wynika z kryterium przez podbazę.
\end{proof}

\begin{uwaga}[Podprzestrzeń i iloczyn jako konstrukcje początkowe]
Topologia podprzestrzeni na $A\subseteq X$ jest początkowa
względem włożenia $\iota\colon A\to X$. Topologia iloczynu
na $\prod_i X_i$ jest początkowa względem rzutowań $\pi_i$.
Oba przypadki wynikają bezpośrednio z definicji rodzin
podbazowych. Ich własności uniwersalne są więc
szczególnymi przypadkami twierdzenia
\ref{thm:topog-poczatkowa-uniwersalnosc}.
\end{uwaga}

\begin{definicja}[Topologia końcowa]
\label{def:topog-topologia-koncowa}
Niech $(X_i,\tau_i)_{i\in I}$ będzie rodziną przestrzeni
topologicznych, $Y$ zbiorem oraz
$f_i\colon X_i\to Y$ odwzorowaniami. \emph{Topologią końcową}
wyznaczoną przez rodzinę $(f_i)_{i\in I}$ nazywamy
\[
\tau_{\mathrm k}=\{V\subseteq Y:
 f_i^{-1}(V)\in\tau_i\text{ dla każdego }i\in I\}.
\]
Jest to największa topologia na $Y$, dla której wszystkie
$f_i$ są ciągłe. Jeśli $I=\varnothing$, otrzymujemy
topologię dyskretną $\mathcal P(Y)$.
\end{definicja}

\begin{stwierdzenie}[Poprawność definicji topologii końcowej]
\label{prop:topog-koncowa-topologia}
Rodzina $\tau_{\mathrm k}$ z definicji
\ref{def:topog-topologia-koncowa} jest topologią.
Każda topologia na $Y$, dla której wszystkie $f_i$
są ciągłe, jest zawarta w $\tau_{\mathrm k}$.
\end{stwierdzenie}
\begin{proof}
Przeciwobrazy $\varnothing$ i $Y$ są odpowiednio
$\varnothing$ i $X_i$, a więc są otwarte.
Przeciwobrazy zachowują dowolne sumy i skończone
przecięcia zbiorów, co dowodzi aksjomatów topologii.
Jeśli w pewnej topologii na $Y$ wszystkie $f_i$
są ciągłe, każdy jej zbiór otwarty ma otwarte
przeciwobrazy przy wszystkich $f_i$, należy więc
do $\tau_{\mathrm k}$.
\end{proof}

\begin{twierdzenie}[Własność uniwersalna topologii końcowej]
\label{thm:topog-koncowa-uniwersalnosc}
Niech na zbiorze $Y$ zadana będzie topologia końcowa
względem $f_i\colon X_i\to Y$, $i\in I$, i niech
$Z$ będzie przestrzenią topologiczną. Odwzorowanie
$h\colon Y\to Z$ jest ciągłe wtedy i tylko wtedy, gdy
wszystkie złożenia $h\circ f_i\colon X_i\to Z$
są ciągłe.
\end{twierdzenie}
\begin{proof}
Jeśli $h$ jest ciągłe, ciągłość złożeń wynika z ciągłości
$f_i$. Jeśli wszystkie złożenia są ciągłe, to dla dowolnego
otwartego $W\subseteq Z$ wszystkie zbiory
\[
f_i^{-1}(h^{-1}(W))=(h\circ f_i)^{-1}(W)
\]
są otwarte. Na mocy definicji topologii końcowej
$h^{-1}(W)$ jest otwarty w $Y$.
\end{proof}

\begin{uwaga}[Suma rozłączna jako konstrukcja końcowa]
Topologia na $\coprod_i X_i$ jest końcowa względem
włożeń $\jmath_i$. W istocie warunek otwartości
$\jmath_i^{-1}(W)$ jest dokładnie warunkiem otwartości
odpowiedniego przekroju $W_i$ z definicji sumy rozłącznej.
\end{uwaga}

\section{Topologia ilorazowa}
\label{sec:topog-iloraz}

\begin{definicja}[Przestrzeń ilorazowa]
\label{def:topog-przestrzen-ilorazowa}
Niech $(X,\tau)$ będzie przestrzenią topologiczną,
a $\sim$ relacją równoważności na $X$. Niech $X/{\sim}$
oznacza zbiór klas równoważności, a
$q\colon X\to X/{\sim}$, $q(x)=[x]$, kanoniczną surjekcję.
\emph{Topologią ilorazową} jest rodzina
\[
\tau_{/{\sim}}=\{V\subseteq X/{\sim}:q^{-1}(V)\in\tau\}.
\]
Przestrzeń $(X/{\sim},\tau_{/{\sim}})$ nazywamy
\emph{przestrzenią ilorazową}. Topologia ilorazowa jest
topologią końcową względem $q$; w szczególności $q$ jest ciągłe.
\end{definicja}

\begin{definicja}[Odwzorowanie ilorazowe i zbiory nasycone]
\label{def:topog-odwzorowanie-ilorazowe}
Niech $q\colon X\to Y$ będzie ciągłą surjekcją
przestrzeni topologicznych. Nazywamy ją
\emph{odwzorowaniem ilorazowym}, jeśli
$V\subseteq Y$ jest otwarty dokładnie wtedy, gdy
$q^{-1}(V)$ jest otwarty w $X$.
Zbiór $A\subseteq X$ nazywamy \emph{nasyconym względem $q$},
jeśli $q^{-1}(q(A))=A$. Jest to równoważne temu,
że $A$ jest sumą włókien $q^{-1}(y)$.
\end{definicja}

\begin{stwierdzenie}[Zbiory nasycone i otwartość w ilorazie]
\label{prop:topog-iloraz-nasycone}
Niech $q\colon X\to Y$ będzie odwzorowaniem ilorazowym.
Dla dowolnego zbioru $A\subseteq X$ zachodzi:
\begin{enumerate}[label=(\roman*)]
\item $q(A)$ jest otwarty w $Y$ wtedy i tylko wtedy, gdy
      $q^{-1}(q(A))$ jest otwarty w $X$;
\item $q(A)$ jest domknięty w $Y$ wtedy i tylko wtedy, gdy
      $q^{-1}(q(A))$ jest domknięty w $X$.
\end{enumerate}
Jeśli $A$ jest nasycony, warunki po prawej stronach
sprowadzają się do otwartości lub domkniętości $A$.
\end{stwierdzenie}
\begin{proof}
Pierwsza równoważność jest definicją odwzorowania
ilorazowego zastosowaną do zbioru $q(A)$.
Dla drugiej stosujemy tę samą definicję do dopełnienia
$Y\setminus q(A)$ i używamy równości
$q^{-1}(Y\setminus q(A))=X\setminus q^{-1}(q(A))$.
Ostatnie zdanie wynika z definicji nasycenia.
\end{proof}

\begin{twierdzenie}[Własność uniwersalna ilorazu]
\label{thm:topog-iloraz-uniwersalnosc}
Niech $q\colon X\to Y$ będzie odwzorowaniem ilorazowym,
$Z$ dowolną przestrzenią topologiczną, a $h\colon Y\to Z$
odwzorowaniem. Wtedy $h$ jest ciągłe wtedy i tylko wtedy,
gdy $h\circ q\colon X\to Z$ jest ciągłe.
Ponadto jeśli $g\colon X\to Z$ jest stałe na każdym
włóknie $q$, istnieje dokładnie jedno odwzorowanie
$h\colon Y\to Z$ z $g=h\circ q$, a $h$ jest ciągłe
dokładnie wtedy, gdy $g$ jest ciągłe.
\end{twierdzenie}
\begin{proof}
Pierwsza równoważność jest własnością uniwersalną
topologii końcowej względem jednego odwzorowania.
Jeśli $g$ ma stałą wartość na każdym włóknie, określamy
$h(y)=g(x)$, gdzie $x$ jest dowolnym punktem z
$q^{-1}(y)$. Włókno jest niepuste z surjektywności,
a stałość $g$ zapewnia niezależność od wyboru $x$.
Wówczas $h(q(x))=g(x)$ dla wszystkich $x\in X$.
Surjektywność $q$ gwarantuje jednoznaczność $h$.
Równoważność ciągłości wynika z pierwszej części.
\end{proof}

\begin{stwierdzenie}[Wystarczające kryterium ilorazowości]
\label{prop:topog-otwarta-surjekcja-iloraz}
Każda ciągła surjekcja $q\colon X\to Y$, która jest
odwzorowaniem otwartym lub odwzorowaniem domkniętym,
jest odwzorowaniem ilorazowym.
\end{stwierdzenie}
\begin{proof}
Ciągłość daje: dla każdego otwartego $V\subseteq Y$
przeciwobraz $q^{-1}(V)$ jest otwarty w $X$.
Jeśli $q$ jest otwarte i $q^{-1}(V)$ jest otwarty,
to $V=q(q^{-1}(V))$ jest otwarty dzięki surjektywności.
Jeśli $q$ jest domknięte i $q^{-1}(V)$ jest otwarty,
to $q^{-1}(Y\setminus V)=X\setminus q^{-1}(V)$
jest domknięty; zatem $Y\setminus V$ jest domknięty
jako obraz zbioru domkniętego. W obu przypadkach
otrzymujemy równoważność definiującą ilorazowość.
\end{proof}

\begin{przyklad}[Sklejenie dwóch punktów prostej]
\label{ex:topog-iloraz-dwa-punkty}
W $\mathbb R$ z topologią euklidesową przyjmijmy
$x\sim y$ wtedy i tylko wtedy, gdy $x=y$ albo
$\{x,y\}=\{0,1\}$. Iloraz otrzymujemy, utożsamiając
$0$ i $1$, a pozostałe punkty pozostają rozróżnione.
Dla $U=(-1/2,1/2)$ zbiór nasycony jest równy
$q^{-1}(q(U))=U\cup\{1\}$, który nie jest otwarty
w $\mathbb R$. Wobec tego $q(U)$ nie jest otwarty
w ilorazie, mimo że $U$ jest otwarty w $\mathbb R$.
Zatem odwzorowanie ilorazowe nie musi być otwarte.
\end{przyklad}

\begin{przyklad}[Iloraz nie musi być odwzorowaniem domkniętym]
\label{ex:topog-iloraz-niedomkniety}
Na $\mathbb R$ utożsammy ze sobą wszystkie punkty
przedziału $(0,1)$, nie utożsamiając żadnych innych.
Niech $q$ będzie odwzorowaniem do ilorazu.
Zbiór $F=\{1/2\}$ jest domknięty w $\mathbb R$,
lecz $q^{-1}(q(F))=(0,1)$ nie jest domknięty.
Z kryterium zbiorów nasyconych wynika, że $q(F)$
nie jest domknięty. Odwzorowanie ilorazowe nie musi
więc być domknięte.
\end{przyklad}

\begin{przyklad}[Ciągła surjekcja, która nie jest ilorazowa]
\label{ex:topog-surjekcja-nieilorazowa}
Niech $S=\{0,1\}$ i rozważmy odwzorowanie
identycznościowe
\[
\operatorname{id}\colon(S,\mathcal P(S))\longrightarrow
       (S,\{\varnothing,S\}).
\]
Jest ono ciągłą surjekcją, ponieważ przeciwdziedzina
ma topologię antydyskretną. Nie jest ilorazowe:
przeciwobraz $\{0\}$ jest otwarty w dyskretnej dziedzinie,
ale $\{0\}$ nie jest otwarty w przeciwdziedzinie.
\end{przyklad}

\section{Zestawienie własności uniwersalnych}
\label{sec:topog-konstrukcje-podsumowanie}

Wszystkie konstrukcje rozdziału powstają z tego samego
pytania: jakie zbiory uznać za otwarte, aby wskazane
odwzorowania były ciągłe?
Dla topologii początkowej wymagamy ciągłości odwzorowań
\emph{do} danych przestrzeni i wybieramy najmniejszą
taką topologię. Dla topologii końcowej wymagamy ciągłości
odwzorowań \emph{z} danych przestrzeni i wybieramy
największą taką topologię. Iloczyny i podprzestrzenie
są szczególnymi przypadkami pierwszej konstrukcji;
sumy rozłączne i ilorazy --- drugiej.

\begin{uwaga}[Kierunek sprawdzania ciągłości]
W topologii podprzestrzeni i iloczynu sprawdzamy ciągłość
odwzorowań $g\colon Z\to X$ na podstawie ich złożeń
z włożeniami lub rzutowaniami \emph{wychodzącymi z $X$}.
W sumie rozłącznej i topologii ilorazowej sprawdzamy
ciągłość odwzorowań $h\colon Y\to Z$ na podstawie
złożeń z naturalnymi odwzorowaniami \emph{wchodzącymi do $Y$}.
Rozróżnienie to będzie wykorzystywane przy konstrukcjach
bardziej złożonych przestrzeni i badaniu ich niezmienników.
\end{uwaga}
\chapter{Aksjomaty oddzielania}
\label{ch:topog-oddzielanie}

\noindent\textbf{Cel rozdziału.}
W dotychczasowych konstrukcjach nie wymagaliśmy, aby różne punkty
miały różne otoczenia ani by punkty i zbiory domknięte można było
oddzielić zbiorami otwartymi. Tymczasem własności te wpływają na
jednoznaczność granic, zachowanie funkcji ciągłych i konstrukcje ilorazowe.
Wprowadzimy aksjomaty $T_0$, $T_1$, $T_2$, regularność, całkowitą
regularność i normalność. Wykażemy także lemat Urysohna
i twierdzenie Tietzego o przedłużeniu. Każde założenie podamy jawnie: samych pojęć
regularności i normalności nie utożsamiamy odpowiednio z aksjomatami
$T_3$ i $T_4$.

\section{Rozróżnianie punktów: $T_0$, $T_1$ i Hausdorff}
\label{sec:topog-t0-t1-t2}

\begin{definicja}[Aksjomaty $T_0$, $T_1$, $T_2$]
\label{def:topog-t0-t1-t2}
Niech $(X,\tau)$ będzie przestrzenią topologiczną.
\begin{enumerate}[label=(\roman*)]
\item $X$ spełnia \emph{aksjomat $T_0$} (jest przestrzenią
Kołmogorowa), gdy dla każdych dwóch różnych punktów $x,y\in X$
istnieje zbiór otwarty zawierający dokładnie jeden z nich.
\item $X$ spełnia \emph{aksjomat $T_1$}, gdy dla każdych $x\ne y$
istnieją zbiory otwarte $U,V\subseteq X$ takie, że
$x\in U$, $y\notin U$ oraz $y\in V$, $x\notin V$.
\item $X$ spełnia \emph{aksjomat $T_2$} (jest przestrzenią
Hausdorffa), gdy dla każdych $x\ne y$ istnieją rozłączne
zbiory otwarte $U,V\subseteq X$ takie, że $x\in U$ i $y\in V$.
\end{enumerate}
Dla przestrzeni pustej i jednoelementowej wszystkie trzy warunki
są spełnione: nie istnieje para różnych punktów.
\end{definicja}


\begin{figure}[htbp]
\centering
\begin{minipage}[t]{.48\textwidth}
\centering
\textbf{Warunek $T_1$}\\[2pt]
\begin{tikzpicture}[scale=.90]
  \fill[blue!10] (-.7,0) ellipse (1.55 and .92);
  \fill[orange!13] (.7,0) ellipse (1.55 and .92);
  \draw[blue!65!black,thick] (-.7,0) ellipse (1.55 and .92);
  \draw[orange!75!black,thick] (.7,0) ellipse (1.55 and .92);
  \fill (-1.75,0) circle (2.3pt);
  \fill (1.75,0) circle (2.3pt);
  \node[below] at (-1.75,-.06) {$x$};
  \node[below] at (1.75,-.06) {$y$};
  \node[blue!60!black] at (-1.23,.52) {$U$};
  \node[orange!70!black] at (1.23,.52) {$V$};
  \node[font=\small] at (0,-1.22) {$y\notin U,\quad x\notin V$};
\end{tikzpicture}
\end{minipage}\hfill
\begin{minipage}[t]{.48\textwidth}
\centering
\textbf{Warunek $T_2$}\\[2pt]
\begin{tikzpicture}[scale=.90]
  \fill[blue!10] (-1.5,0) ellipse (1.05 and .85);
  \fill[orange!13] (1.5,0) ellipse (1.05 and .85);
  \draw[blue!65!black,thick] (-1.5,0) ellipse (1.05 and .85);
  \draw[orange!75!black,thick] (1.5,0) ellipse (1.05 and .85);
  \fill (-1.5,0) circle (2.3pt);
  \fill (1.5,0) circle (2.3pt);
  \node[below] at (-1.5,-.08) {$x$};
  \node[below] at (1.5,-.08) {$y$};
  \node[blue!60!black] at (-1.5,.48) {$U$};
  \node[orange!70!black] at (1.5,.48) {$V$};
  \node[font=\small] at (0,-1.15) {$U\cap V=\varnothing$};
\end{tikzpicture}
\end{minipage}
\caption{Porównanie aksjomatów $T_1$ i $T_2$. Po lewej
\emph{wybrane} otoczenie $U$ punktu $x$ nie zawiera $y$,
a wybrane otoczenie $V$ punktu $y$ nie zawiera $x$;
otoczenia te mogą się przecinać. Po prawej możemy wybrać
otoczenia rozłączne. Rysunki przedstawiają schematy warunków,
a nie konkretne topologie.}
\label{fig:topog-t1-t2-porownanie}
\end{figure}

\begin{stwierdzenie}[Pierwsze implikacje]
\label{prop:topog-t2-t1-t0}
Dla dowolnej przestrzeni topologicznej $X$ zachodzą implikacje
$T_2\Rightarrow T_1\Rightarrow T_0$.
\end{stwierdzenie}
\begin{proof}
Jeśli $x\in U$, $y\in V$ i $U\cap V=\varnothing$, to
$y\notin U$ i $x\notin V$, co daje $T_1$. Z kolei przy założeniu
$T_1$ jeden ze zbiorów otwartych przewidzianych w definicji
rozróżnia punkty w sensie $T_0$.
\end{proof}

\begin{przyklad}[Przestrzeń Sierpińskiego]
\label{ex:topog-sierpinski-oddzielanie}
Niech $X=\{0,1\}$ i $\tau=\{\varnothing,\{1\},X\}$.
Otwarty zbiór $\{1\}$ rozróżnia $0$ i $1$, więc $X$ jest $T_0$.
Nie ma jednak otwartego zbioru zawierającego $0$ i niezawierającego
$1$, ponieważ jedynym otoczeniem otwartym punktu $0$ jest $X$.
Przestrzeń nie spełnia $T_1$, a tym bardziej $T_2$.
\end{przyklad}

\begin{twierdzenie}[Kryterium za pomocą singletonów]
\label{thm:topog-t1-singletony}
Przestrzeń topologiczna $X$ jest $T_1$ wtedy i tylko wtedy,
gdy każdy zbiór jednoelementowy $\{x\}$, $x\in X$, jest domknięty.
\end{twierdzenie}
\begin{proof}
Załóżmy $T_1$ i ustalmy $x\in X$. Dla każdego $y\ne x$
istnieje otwarty zbiór $U_y$ zawierający $y$ i niezawierający
$x$. Stąd
\[
X\setminus\{x\}=\bigcup\{U\in\tau:x\notin U\},
\]
więc dopełnienie singletonu jest otwarte, a singleton domknięty.
Równość jest prawdziwa również, gdy $X=\{x\}$: obie strony są puste.

Odwrotnie, niech wszystkie singletony będą domknięte i $x\ne y$.
Wówczas otwarte zbiory $X\setminus\{y\}$ oraz
$X\setminus\{x\}$ zawierają odpowiednio $x$ i $y$
i nie zawierają drugiego punktu, co daje $T_1$.
\end{proof}

\begin{wniosek}[Skończone przestrzenie $T_1$]
\label{cor:topog-skonczone-t1}
Skończona przestrzeń topologiczna jest $T_1$ wtedy i tylko wtedy,
gdy ma topologię dyskretną.
\end{wniosek}
\begin{proof}
W przestrzeni $T_1$ każdy singleton jest domknięty, a każdy
podzbiór zbioru skończonego jest skończoną sumą singletonów.
Wszystkie podzbiory są więc domknięte, a zatem także otwarte.
Odwrotnie, w topologii dyskretnej singletony są domknięte.
\end{proof}

\begin{przyklad}[Topologia koskończona]
\label{ex:topog-koskonczona-t1-nie-t2}
Niech $X$ będzie nieskończony z topologią koskończoną:
zbiory otwarte to $\varnothing$ i zbiory o skończonym dopełnieniu.
Każdy singleton jest domknięty, zatem $X$ jest $T_1$.
Jeżeli $U,V$ są niepustymi otwartymi, to
$X\setminus(U\cap V)=(X\setminus U)\cup(X\setminus V)$
jest skończony; ponieważ $X$ jest nieskończony,
$U\cap V\ne\varnothing$. Nie można więc rozdzielić
żadnej pary różnych punktów rozłącznymi zbiorami otwartymi:
$X$ nie jest Hausdorffa.
\end{przyklad}

\begin{stwierdzenie}[Przekątna i aksjomat Hausdorffa]
\label{prop:topog-przekatna-hausdorff}
Niech $X$ będzie przestrzenią topologiczną, a $X\times X$
ma topologię iloczynu. Wówczas $X$ jest Hausdorffa wtedy i tylko
wtedy, gdy przekątna $\Delta_X=\{(x,x):x\in X\}$ jest domknięta
w $X\times X$.
\end{stwierdzenie}
\begin{proof}
Jeśli $X$ jest Hausdorffa i $(x,y)\notin\Delta_X$, to $x\ne y$.
Wybierzmy rozłączne otwarte $U\ni x$ i $V\ni y$.
Prostokąt $U\times V$ jest otwarty, zawiera $(x,y)$ i nie przecina
$\Delta_X$. Takie prostokąty pokrywają dopełnienie przekątnej,
które jest więc otwarte.

Odwrotnie, niech $\Delta_X$ będzie domknięta i $x\ne y$.
Dopełnienie przekątnej jest otwarte i zawiera $(x,y)$.
Z kryterium bazy topologii iloczynu, udowodnionego
w rozdziale 04, istnieją otwarte $U\ni x$,
$V\ni y$ takie, że $(x,y)\in U\times V\subseteq
(X\times X)\setminus\Delta_X$. Gdyby $z\in U\cap V$,
to $(z,z)\in U\times V\cap\Delta_X$, sprzeczność.
Zatem $U\cap V=\varnothing$.
\end{proof}

\begin{twierdzenie}[Zgodność odwzorowań na zbiorze gęstym]
\label{thm:topog-rownosc-na-gestym}
Niech $X$ będzie przestrzenią topologiczną, $Y$ przestrzenią
Hausdorffa, a $f,g\colon X\to Y$ odwzorowaniami ciągłymi.
Wówczas zbiór $E=\{x\in X:f(x)=g(x)\}$ jest domknięty.
Jeśli ponadto $D\subseteq X$ jest gęsty w $X$ i $f=g$ na $D$,
to $f=g$ na całym $X$.
\end{twierdzenie}
\begin{proof}
Odwzorowanie $h=(f,g)\colon X\to Y\times Y$ jest ciągłe:
jego współrzędne są ciągłe, a więc stosujemy własność
uniwersalną iloczynu z rozdziału 04.
Z poprzedniego stwierdzenia $\Delta_Y$ jest domknięta,
zatem $E=h^{-1}(\Delta_Y)$ jest domknięty.
Jeśli $D\subseteq E$ i $\overline D=X$, to
$X=\overline D\subseteq E$, a więc $E=X$.
\end{proof}

\begin{uwaga}[Warunek Hausdorffa ma znaczenie]
Niech $Y$ będzie przestrzenią Sierpińskiego i niech
$X=Y$, $D=\{1\}$. Zbiór $D$ jest gęsty, ponieważ jedynym
niepustym otwartym zbiorem zawierającym $0$ jest $Y$.
Odwzorowanie identycznościowe $f\colon X\to Y$ i odwzorowanie
stałe $g(x)=1$ są ciągłe, zgadzają się na $D$, lecz nie
zgadzają się w punkcie $0$. Pokazuje to konieczność
warunku oddzielania w poprzednim twierdzeniu.
\end{uwaga}

\section{Regularność i aksjomat $T_3$}
\label{sec:topog-regularnosc}

\begin{definicja}[Regularność i aksjomat $T_3$]
\label{def:topog-regularnosc-t3}
Przestrzeń topologiczna $X$ jest \emph{regularna}, jeżeli dla
każdego zbioru domkniętego $F\subseteq X$ i każdego $x\in X\setminus F$
istnieją rozłączne zbiory otwarte $U,V\subseteq X$ takie, że
$x\in U$ i $F\subseteq V$.
Przestrzeń spełnia \emph{aksjomat $T_3$}, gdy jest jednocześnie
regularna i $T_1$.
\end{definicja}

\begin{stwierdzenie}[Kryterium zwężania otoczeń]
\label{prop:topog-regularnosc-zwezanie}
Przestrzeń topologiczna $X$ jest regularna wtedy i tylko wtedy,
gdy dla każdego $x\in X$ i każdego otwartego $W\subseteq X$
zawierającego $x$ istnieje otwarty $U\subseteq X$ taki, że
$x\in U\subseteq\overline U\subseteq W$.
\end{stwierdzenie}
\begin{proof}
Załóżmy regularność i przyjmijmy $F=X\setminus W$.
Ponieważ $F$ jest domknięty i $x\notin F$, istnieją rozłączne
otwarte $U\ni x$ i $V\supseteq F$. Domknięty zbiór
$X\setminus V$ zawiera $U$, stąd
$\overline U\subseteq X\setminus V\subseteq W$.

Odwrotnie, dla domkniętego $F$ i $x\notin F$ zastosujmy warunek
do otwartego $W=X\setminus F$. Otrzymamy
$x\in U\subseteq\overline U\subseteq X\setminus F$.
Zbiory $U$ i $V=X\setminus\overline U$ są otwarte i rozłączne,
przy czym $F\subseteq V$.
\end{proof}

\begin{wniosek}[$T_3$ pociąga za sobą $T_2$]
\label{cor:topog-t3-t2}
Każda przestrzeń $T_3$ jest Hausdorffa.
\end{wniosek}
\begin{proof}
Dla $x\ne y$ singleton $\{y\}$ jest domknięty na mocy
aksjomatu $T_1$ i twierdzenia \ref{thm:topog-t1-singletony}.
Regularność rozdziela punkt $x$ i zbiór $\{y\}$
rozłącznymi otwartymi otoczeniami.
\end{proof}

\begin{figure}[htbp]
\centering
\begin{tikzpicture}[scale=1,>=Stealth]
 \fill[blue!6] (0,0) ellipse (3.1 and 1.65);
 \draw[blue!65!black,thick] (0,0) ellipse (3.1 and 1.65);
 \fill[blue!12] (-.9,0) ellipse (1.1 and 1.04);
 \draw[blue!70!black,thick,dashed] (-.9,0) ellipse (1.1 and 1.04);
 \fill[green!14] (-.9,0) ellipse (.58 and .50);
 \draw[green!45!black,thick] (-.9,0) ellipse (.58 and .50);
 \fill (-.9,0) circle (2.4pt);
 \node[below] at (-.9,-.1) {$x$};
 \node[green!35!black] at (-.9,.31) {$U$};
 \node[blue!65!black] at (-1.5,.7) {$W$};
 \draw[fill=red!18,draw=red!70!black,thick] (1.7,-.2) ellipse (.55 and .64);
 \node at (1.7,-.2) {$F$};
 \node[blue!55!black] at (-2.4,.7) {$X$};
\end{tikzpicture}
\caption{Regularność: wokół $x\notin F$ można dobrać otwarty
$U$ tak, by $\overline U$ nie spotykało domkniętego $F$.
Elipsy przedstawiają schemat otoczeń, a nie ich kanoniczny kształt.}
\label{fig:topog-regularnosc}
\end{figure}

\begin{przyklad}[Przestrzeń Hausdorffa, która nie jest regularna]
\label{ex:topog-k-topologia}
Niech $K=\{1/n:n\in\mathbb N,\ n\ge1\}\subseteq\mathbb R$.
Na $\mathbb R$ rozważmy topologię generowaną przez
zbiory otwarte euklidesowo oraz zbiory $U\setminus K$,
gdzie $U$ jest otwarty euklidesowo. Ta topologia jest
silniejsza od euklidesowej, a więc jest Hausdorffa.

Zbiór $K$ jest domknięty w nowej topologii, gdyż jego dopełnienie
$\mathbb R\setminus K$ ma postać $U\setminus K$ dla
$U=\mathbb R$ i jest otwarte. Załóżmy, że istnieją rozłączne
otwarte $P\ni0$ i $Q\supseteq K$. Każde otoczenie $P$ punktu $0$
zawiera pewien zbiór $(-\varepsilon,\varepsilon)\setminus K$:
taką bazę uzyskuje się przez skończone przecięcia podanych
otwartych zbiorów (zob. kryterium podbazy w rozdziale 01).
Wybierzmy $n$ z $1/n<\varepsilon$.
Ponieważ $Q$ jest otwarty i zawiera $1/n\in K$, zawiera
pewien euklidesowy przedział wokół $1/n$: zbiory
$U\setminus K$ nie zawierają bowiem żadnego punktu $K$.
Przedział ten spotyka $(-\varepsilon,\varepsilon)\setminus K$,
gdyż $K$ jest przeliczalny, a każdy niepusty przedział
jest nieprzeliczalny. Sprzeczność z $P\cap Q=\varnothing$.
Zatem przestrzeń nie jest regularna i nie spełnia $T_3$.
\end{przyklad}

\section{Całkowita regularność i przestrzenie Tychonoffa}
\label{sec:topog-tych}

\begin{definicja}[Całkowita regularność]
\label{def:topog-calkowita-regularnosc}
Przestrzeń topologiczna $X$ jest \emph{całkowicie regularna},
jeśli dla każdego domkniętego $F\subseteq X$ i $x\in X\setminus F$
istnieje ciągła funkcja $f\colon X\to[0,1]$ taka, że
$f(x)=0$ oraz $f(y)=1$ dla każdego $y\in F$.
Przedział $[0,1]$ ma topologię euklidesową.
Przestrzeń \emph{Tychonoffa} (spełniającą $T_{3\frac12}$)
nazywamy przestrzeń całkowicie regularną i $T_1$.
\end{definicja}

\begin{stwierdzenie}[Całkowita regularność implikuje regularność]
\label{prop:topog-calkowita-regularnosc-regularna}
Każda przestrzeń całkowicie regularna jest regularna.
W szczególności każda przestrzeń Tychonoffa spełnia $T_3$.
\end{stwierdzenie}
\begin{proof}
Dla domkniętego $F$ i $x\notin F$ weźmy funkcję $f$
z definicji. Wtedy
$U=f^{-1}([0,1/3))$ i $V=f^{-1}((2/3,1])$
są otwarte, rozłączne, $x\in U$ i $F\subseteq V$.
Otwartość wynika z topologii względnej przedziału $[0,1]$
i ciągłości $f$. Jeśli ponadto zachodzi $T_1$,
otrzymujemy $T_3$.
\end{proof}

\begin{przyklad}[Baza złożona ze zbiorów otwarto-domkniętych]
\label{ex:topog-zero-wymiarowy-tych}
Załóżmy, że przestrzeń $X$ jest $T_1$ i ma bazę złożoną
ze zbiorów jednocześnie otwartych i domkniętych.
Dla domkniętego $F$ i $x\notin F$ otwarty zbiór
$X\setminus F$ zawiera pewien bazowy zbiór
$C$ taki, że $x\in C\subseteq X\setminus F$.
Funkcja $f\colon X\to[0,1]$ przyjmująca $0$ na $C$
i $1$ poza $C$ jest ciągła: przeciwobrazy zbiorów
otwartych w $[0,1]$ są jednym ze zbiorów
$\varnothing,C,X\setminus C,X$. Spełnia także
$f(x)=0$ i $f|_F=1$, więc $X$ jest Tychonoffa.
\end{przyklad}


\begin{przyklad}[Przestrzeń Mysiora: $T_3$ bez aksjomatu Tychonoffa]
\label{ex:topog-mysior-t3-nie-tych}
Zbudujemy przestrzeń Hausdorffa, która jest regularna,
lecz nie jest całkowicie regularna. Konstrukcja pochodzi od
A.~Mysiora (\emph{A regular space which is not completely regular},
\emph{Proc. Amer. Math. Soc.} 81 (1981), 652--653).

Niech $H=\mathbb R\times[0,2)$ i dołączmy nowy punkt
$\infty\notin H$, otrzymując $X=H\cup\{\infty\}$.
Dla $x\in\mathbb R$ określmy zbiór
\[
B_x=\{(x,t):0\le t<2\}\;\cup\;
\{(x+t,t):0<t<2\}.
\]
Topologię na $X$ określamy przez następujące bazy otoczeń:
\begin{enumerate}[label=(\roman*)]
\item każdy punkt $(u,t)$ z $t>0$ jest izolowany;
\item w punkcie $(x,0)$ bazą są zbiory $B_x\setminus E$,
      gdzie $E$ przebiega skończone podzbiory
      $B_x\setminus\{(x,0)\}$;
\item w punkcie $\infty$ bazą są zbiory
      $W_n=\{\infty\}\cup\{(u,t)\in H:u>n\}$, $n\in\mathbb N$.
\end{enumerate}
Jest to rzeczywiście baza topologii. Każde $B_x\setminus E$
poza $(x,0)$ zawiera tylko punkty izolowane; w $W_n$
każdy punkt osi $(x,0)$ ma $x>n$, a wówczas
$B_x\subseteq W_n$, gdyż pierwsza współrzędna
wszystkich punktów $B_x$ jest nie mniejsza niż $x$.
Ponadto $W_m\subseteq W_n$ dla $m\ge n$.
Stąd spełnione są warunki zgodności lokalnych baz.

\emph{Krok 1: przestrzeń jest $T_3$.}
Dla $x\ne z$ przecięcie $B_x\cap B_z$ jest skończone:
ukośne odcinki są równoległe, a każdy ukośny odcinek
może przecinać pionowy odcinek najwyżej w jednym punkcie.
Zatem punkty $(x,0)$ i $(z,0)$ można rozdzielić
rozłącznymi otwartymi zbiorami, usuwając z jednego
z $B_x,B_z$ ich skończone przecięcie.
Punkt izolowany rozdzielamy od punktu osi, usuwając go
w razie potrzeby z bazowego otoczenia tego punktu.
Punkt $\infty$ rozdzielamy od $(x,0)$, wybierając
$n>x+2$, ponieważ wszystkie punkty $B_x$ mają
pierwszą współrzędną mniejszą niż $x+2$.
Od punktu izolowanego oddzielamy $\infty$ przez
odpowiednio duże $W_n$. Tak więc $X$ jest Hausdorffa,
a w szczególności spełnia $T_1$.

Co więcej, każdy zbiór $B_x\setminus E$ jest
\emph{otwarto-domknięty} w $X$.
Jest otwarty z konstrukcji. Punkt izolowany poza nim
ma otwarte otoczenie jednoelementowe, punkt $(z,0)$
dla $z\ne x$ ma otoczenie
$B_z\setminus(B_x\cap B_z)$ rozłączne z $B_x\setminus E$,
a punkt $\infty$ ma takie otoczenie $W_n$ przy $n>x+2$.
Oznacza to otwartość jego dopełnienia.
Zatem każdy punkt $H$ ma bazę otoczeń
otwarto-domkniętych (dla punktów izolowanych
korzystamy dodatkowo z $T_1$).

Pozostało sprawdzić regularność w punkcie $\infty$.
Dla każdego $n\in\mathbb N$ zachodzi
\[
\overline{W_{n+2}}\subseteq W_n.
\]
Istotnie, żaden punkt izolowany spoza $W_n$ nie należy
do domknięcia $W_{n+2}$. Jeśli $(x,0)\notin W_n$,
to $x\le n$, a wszystkie punkty $B_x$ mają
pierwszą współrzędną mniejszą niż $x+2\le n+2$;
stąd $B_x\cap W_{n+2}=\varnothing$.
Wykazaliśmy zatem kryterium zwężania otoczeń
ze stwierdzenia \ref{prop:topog-regularnosc-zwezanie}
we wszystkich punktach $X$.
Przestrzeń jest regularna i $T_1$, czyli spełnia $T_3$.

\emph{Krok 2: przestrzeń nie jest całkowicie regularna.}
Zbiór $F=[0,1]\times\{0\}$ jest domknięty w $X$:
każdy punkt osi poza $F$ ma otoczenie $B_x$ rozłączne
z $F$, punkty nad osią są izolowane, a $W_2\cap F=\varnothing$.
Pokażemy, że dowolna ciągła funkcja
$f\colon X\to[0,1]$, spełniająca $f|_F=1$,
musi również spełniać $f(\infty)=1$.

Najpierw zauważmy, że jeżeli $f(x,0)=1$, to
$f=1$ na zbiorze $B_x\setminus E_x$ dla pewnego
\emph{przeliczalnego} $E_x\subseteq B_x\setminus\{(x,0)\}$.
Istotnie, dla każdego $m\ge1$ ciągłość w $(x,0)$
daje skończony $E_{x,m}$, poza którym
$|f-1|<1/m$ na $B_x$. Wystarczy wziąć
$E_x=\bigcup_{m\ge1}E_{x,m}$.

Dla $k\ge1$ oznaczmy
$A_k=\{x\in(k-1,k]:f(x,0)=1\}$.
Ponieważ $f|_F=1$, zbiór $A_1=(0,1]$ jest nieskończony.
Udowodnimy indukcyjnie, że każdy $A_k$ jest nieskończony.
Wybierzmy parami różne $x_j\in A_k$, $j\ge1$,
oraz odpowiadające im przeliczalne zbiory wyjątkowe $E_{x_j}$.
Zbiór
\[
C=\bigcup_{j\ge1}\{x_j+t:\ 0<t<2,\ (x_j+t,t)\in E_{x_j}\}
\]
jest przeliczalny. Dla każdego
$z\in(k,k+1]\setminus C$ oraz każdego $j$
liczba $t_j=z-x_j$ należy do $(0,2)$,
a punkt $(z,t_j)=(x_j+t_j,t_j)$ należy do
$B_{x_j}\setminus E_{x_j}$. Zatem $f(z,t_j)=1$.
Liczby $t_j$ są parami różne, toteż otrzymaliśmy
nieskończenie wiele punktów pionowego odcinka
$B_z$, w których $f$ jest równa $1$.
Ciągłość w $(z,0)$ wymusza $f(z,0)=1$:
gdyby było inaczej, pewne otoczenie bazowe
$B_z\setminus E$ z $E$ skończonym miałoby
obraz w otoczeniu $f(z,0)$ niezawierającym $1$,
co przeczyłoby istnieniu tych punktów.
W konsekwencji
$(k,k+1]\setminus C\subseteq A_{k+1}$,
więc $A_{k+1}$ jest nieskończony.

W każdym otoczeniu $W_n$ punktu $\infty$ znajduje się
zatem punkt $(z,0)$ z $z\in A_{n+1}$ i $f(z,0)=1$.
Z ciągłości $f$ w $\infty$ wynika $f(\infty)=1$:
gdyby $f(\infty)\ne1$, przeciwobraz pewnego otoczenia
$f(\infty)$ niezawierającego $1$ obejmowałby
jedno z $W_n$, co jest niemożliwe.
Nie istnieje więc ciągła funkcja oddzielająca
$\infty$ od domkniętego $F$ wartościami $0$ i $1$.
Oznacza to, że $X$ nie jest całkowicie regularna,
a zatem $T_3\not\Rightarrow T_{3\frac12}$.
\end{przyklad}

\begin{figure}[htbp]
\centering
\begin{tikzpicture}[>=Stealth,scale=.95]
  \draw[->,thick] (-.9,0)--(4.1,0) node[right] {$u$};
  \draw[->,thick] (0,-.25)--(0,2.6) node[above] {$t$};
  \draw[densely dotted,gray] (-.5,2)--(3.65,2);
  \node[gray,above,font=\footnotesize] at (3.1,2.02) {$t=2$};
  \draw[blue!65!black,very thick] (1,0)--(1,2);
  \draw[orange!75!black,very thick] (1,0)--(3,2);
  \draw[white,fill=white,draw=blue!65!black,thick] (1,2) circle (2.4pt);
  \draw[white,fill=white,draw=orange!75!black,thick] (3,2) circle (2.4pt);
  \fill (1,0) circle (2.5pt);
  \node[below] at (1,0) {$(x,0)$};
  \node[blue!65!black,left,font=\small] at (.9,1.15) {pionowy};
  \node[orange!75!black,below right,font=\small] at (2.15,1.15) {ukośny};
\end{tikzpicture}
\caption{Przestrzeń Mysiora: schemat dwóch odcinków tworzących
$B_x$. Punkt $(x,0)$ należy do $B_x$, a końce na wysokości $2$
są wyłączone. Pozostałe punkty nad osią są izolowane.
Rysunek nie przedstawia otwartych zbiorów euklidesowych.}
\label{fig:topog-mysior-baza}
\end{figure}

\section{Normalność, aksjomat $T_4$ i lemat Urysohna}
\label{sec:topog-normalnosc}

\begin{definicja}[Normalność i aksjomat $T_4$]
\label{def:topog-normalnosc-t4}
Przestrzeń topologiczna $X$ jest \emph{normalna}, jeżeli
dla każdych rozłącznych zbiorów domkniętych $F,G\subseteq X$
istnieją rozłączne zbiory otwarte $U,V\subseteq X$
takie, że $F\subseteq U$ i $G\subseteq V$.
Przestrzeń spełnia \emph{aksjomat $T_4$}, jeśli jest
normalna i $T_1$.
\end{definicja}

\begin{stwierdzenie}[Kryterium normalności]
\label{prop:topog-normalnosc-zwezanie}
Przestrzeń $X$ jest normalna wtedy i tylko wtedy,
gdy dla każdego domkniętego $F\subseteq X$ i otwartego
$W\subseteq X$, przy $F\subseteq W$, istnieje otwarty
$U\subseteq X$ taki, że
$F\subseteq U\subseteq\overline U\subseteq W$.
\end{stwierdzenie}
\begin{proof}
Jeśli $X$ jest normalna, rozdzielmy rozłączne domknięte
zbiory $F$ i $X\setminus W$ otwartymi $U$ i $V$.
Z $U\cap V=\varnothing$ wynika
$\overline U\subseteq X\setminus V\subseteq W$.
Odwrotnie, dla rozłącznych domkniętych $F,G$ weźmy
$W=X\setminus G$ i otwarty $U$ z kryterium.
Wtedy $V=X\setminus\overline U$ jest otwarty, zawiera $G$
i jest rozłączny z $U$.
\end{proof}

\begin{twierdzenie}[Lemat Urysohna]
\label{thm:topog-urysohn}
Niech $X$ będzie przestrzenią normalną, a $F,G\subseteq X$
rozłącznymi zbiorami domkniętymi. Wtedy istnieje
ciągła funkcja $f\colon X\to[0,1]$ taka, że
$f|_F=0$ i $f|_G=1$.
\end{twierdzenie}
\begin{proof}
Niech $D$ będzie zbiorem liczb dwójkowych wymiernych
w $[0,1]$, czyli liczb $k/2^m$ dla $m\ge0$ i
$0\le k\le2^m$. Zbudujemy otwarte zbiory $U_r\subseteq X$
dla $r\in D$ spełniające
\[
F\subseteq U_0,\qquad U_1=X\setminus G,\qquad
\overline{U_r}\subseteq U_s\quad(r<s,\ r,s\in D).
\]
Korzystając z kryterium normalności
(\ref{prop:topog-normalnosc-zwezanie}), wybierzmy najpierw
otwarty $U_0$ taki, że
$F\subseteq U_0\subseteq\overline{U_0}\subseteq X\setminus G=U_1$.
Następnie indukcyjnie, gdy zbiory odpowiadające wszystkim
$k/2^m$ są wybrane, dla każdej sąsiedniej pary
$r=k/2^m$, $s=(k+1)/2^m$ wybierzmy otwarty
$U_{(r+s)/2}$ tak, aby
\[
\overline{U_r}\subseteq U_{(r+s)/2}
\subseteq\overline{U_{(r+s)/2}}\subseteq U_s.
\]
Jest to możliwe, ponieważ $\overline{U_r}$ jest domknięty
i zawiera się w otwartym $U_s$. Zachowujemy stare zbiory
przy przejściu do kolejnego poziomu. Przez indukcję
oraz przechodniość zawierania otrzymujemy wymaganą
relację dla wszystkich $r<s$ w $D$.

Zdefiniujmy
\[
f(x)=\inf\bigl(\{r\in D:x\in U_r\}\cup\{1\}\bigr).
\]
Funkcja jest dobrze określona i przyjmuje wartości w $[0,1]$.
Jeśli $x\in F$, to $x\in U_0$ i $f(x)=0$.
Jeśli $x\in G$, to nie należy on do $U_1$,
a z $U_r\subseteq U_1$ dla $r<1$ wynika
$x\notin U_r$ dla żadnego $r\in D$, więc $f(x)=1$.

Sprawdzimy ciągłość. Dla $0<a\le1$ zachodzi
\[
\{x:f(x)<a\}=\bigcup_{\substack{r\in D\\r<a}}U_r.
\]
Rzeczywiście, $f(x)<a$ oznacza, że w zbiorze, którego
infimum definiuje $f(x)$, występuje element mniejszy
od $a$; liczba $1$ nim nie jest. Odwrotna implikacja
wynika bezpośrednio z definicji infimum.
Dla $0\le a<1$ mamy również
\[
\{x:f(x)>a\}=\bigcup_{\substack{r\in D\\a<r<1}}
\bigl(X\setminus\overline{U_r}\bigr).
\]
Jeżeli $x\notin\overline{U_r}$ przy $r>a$, to
$x\notin U_s$ dla żadnego $s\le r$ (bo
$U_s\subseteq\overline{U_s}\subseteq U_r$ dla $s<r$),
zatem $f(x)\ge r>a$.
Odwrotnie, gdy $f(x)>a$, wybierzmy liczby
$r,s\in D$ takie, że $a<r<s<f(x)$; można je dobrać,
gdyż $D$ jest gęsty w $[0,1]$.
Z $s<f(x)$ wynika $x\notin U_s$, a
$\overline{U_r}\subseteq U_s$ daje
$x\notin\overline{U_r}$. Oba zbiory po prawej stronie
powyższych równości są sumami zbiorów otwartych.
Ponieważ przeciwobrazy przedziałów $[0,a)$ i $(a,1]$
tworzą podbazę topologii $[0,1]$, $f$ jest ciągła.
\end{proof}


\begin{wniosek}[$T_4$ implikuje aksjomat Tychonoffa]
\label{cor:topog-t4-tych}
Każda przestrzeń $T_4$ jest Tychonoffa, a w szczególności $T_3$
i Hausdorffa.
\end{wniosek}
\begin{proof}
Niech $F\subseteq X$ będzie domknięty i $x\notin F$.
Z aksjomatu $T_1$ wynika domkniętość singletonu $\{x\}$.
Lemat Urysohna zastosowany do $\{x\}$ i $F$ daje
ciągłą $f\colon X\to[0,1]$ z $f(x)=0$ i $f|_F=1$.
To całkowita regularność; założenie $T_1$ daje warunek
Tychonoffa. Poprzednie implikacje kończą dowód.
\end{proof}

\begin{przyklad}[Normalność bez regularności]
\label{ex:topog-normalna-nie-regularna}
Przestrzeń Sierpińskiego $X=\{0,1\}$ z topologią
$\{\varnothing,\{1\},X\}$ jest normalna. Jedynymi
domkniętymi podzbiorami są $\varnothing$, $\{0\}$, $X$,
więc rozłączne zbiory domknięte nie mogą być jednocześnie
niepuste. Można je oddzielić otwartymi $\varnothing$
i $X$. Nie jest jednak regularna: punkt $1$ nie należy
do domkniętego $F=\{0\}$, lecz jedynym otwartym zbiorem
zawierającym $F$ jest $X$, który przecina każde otwarte
otoczenie punktu $1$. Tak więc normalność bez $T_1$
nie implikuje regularności.
\end{przyklad}

\section{Twierdzenie Tietzego o przedłużeniu}
\label{sec:topog-tietze}

Lemat Urysohna rozdziela dwa rozłączne zbiory domknięte
za pomocą funkcji ciągłej. Twierdzenie Tietzego pozwala
zachować \emph{dowolne} wartości ciągłej funkcji określonej
na zbiorze domkniętym, a nie tylko wartości $0$ i $1$.

\begin{twierdzenie}[Twierdzenie Tietzego o przedłużeniu]
\label{thm:topog-tietze}
Niech $X$ będzie przestrzenią normalną, $A\subseteq X$
zbiorem domkniętym, a $f\colon A\to\mathbb R$ funkcją ciągłą.
Wówczas istnieje ciągła funkcja $F\colon X\to\mathbb R$
taka, że $F|_A=f$. Ponadto, jeśli dla pewnych
$a,b\in\mathbb R$, $a\le b$, zachodzi $f(A)\subseteq[a,b]$,
to można wybrać przedłużenie $F$ o wartościach w $[a,b]$.
Nie zakładamy aksjomatu $T_1$.
\end{twierdzenie}
\begin{proof}
\emph{Krok 1: przybliżenie funkcji ograniczonej.}
Niech $h\colon A\to\mathbb R$ będzie ciągła i
$|h(x)|\le M$ dla każdego $x\in A$, gdzie $M>0$.
Podzbiory
\[
C_- = \{x\in A:h(x)\le -M/3\},\qquad
C_+ = \{x\in A:h(x)\ge M/3\}
\]
są rozłączne i domknięte w $A$ na mocy ciągłości $h$.
Są także domknięte w $X$, ponieważ $A$ jest
zbiorem domkniętym w $X$ (korzystamy z opisu zbiorów
domkniętych w podprzestrzeni z rozdziału 04).
Z lematu Urysohna \ref{thm:topog-urysohn} istnieje
ciągła $u\colon X\to[0,1]$ taka, że $u|_{C_-}=0$
oraz $u|_{C_+}=1$. Funkcja
$g=(M/3)(2u-1)$ jest ciągła, $|g|\le M/3$ na $X$
i spełnia
\[
|h(x)-g(x)|\le \frac{2M}{3}\qquad(x\in A).
\]
Istotnie, na $C_-$ mamy $g=-M/3$, więc
$-2M/3\le h-g\le 0$; na $C_+$ mamy $g=M/3$,
więc $0\le h-g\le 2M/3$. W pozostałych punktach
$|h(x)|<M/3$, a $|g(x)|\le M/3$, zatem
nierówność wynika z oszacowania
$|h-g|\le|h|+|g|$. Jest to przybliżenie o błędzie
co najwyżej $2M/3$.

\emph{Krok 2: przedłużenie z zachowaniem przedziału $[-1,1]$.}
Załóżmy teraz, że $f\colon A\to[-1,1]$ jest ciągła.
Zdefiniujmy $h_0=f$ oraz $M_n=(2/3)^n$.
Stosując konstrukcję z kroku 1 kolejno do $h_n$
i $M_n$, otrzymujemy ciągłe funkcje
$g_n\colon X\to\mathbb R$ i
$h_{n+1}=h_n-g_n|_A$ takie, że
\[
|h_n|\le M_n\ \text{na }A,\qquad
|g_n|\le\frac{M_n}{3}\ \text{na }X,\qquad
|h_{n+1}|\le M_{n+1}\ \text{na }A.
\]
Indukcję rozpoczyna $|h_0|\le M_0=1$.
Ponieważ
$\sum_{n=0}^{\infty} M_n/3
=\frac13\sum_{n=0}^{\infty}(2/3)^n=1$,
szereg $\sum_{n=0}^{\infty}g_n$ jest jednostajnie
zbieżny na $X$ (kryterium Weierstrassa).
Niech $F=\sum_{n=0}^{\infty}g_n$.
Funkcja $F$ jest ciągła jako granica jednostajna
funkcji ciągłych i $|F|\le1$. Dla każdego $N\ge1$
na $A$ zachodzi, z konstrukcji $h_{n+1}$,
\[
f-\sum_{n=0}^{N-1}g_n|_A=h_N,\qquad
|h_N|\le(2/3)^N.
\]
Po przejściu z $N$ do nieskończoności otrzymujemy
$F|_A=f$. To dowodzi twierdzenia dla funkcji
o wartościach w $[-1,1]$. Jeśli $a=b$, przedłużeniem
jest funkcja stała $a$. Gdy $a<b$, stosujemy właśnie
udowodniony przypadek do funkcji
$(2f-a-b)/(b-a)\colon A\to[-1,1]$,
a następnie wykonujemy odwrotną zmianę afiniczną.
Dostajemy przedłużenie o wartościach w $[a,b]$.

\emph{Krok 3: dowolna funkcja rzeczywista.}
Rozważmy homeomorfizm
$\varphi\colon\mathbb R\to(-1,1)$ zadany wzorem
$\varphi(t)=t/(1+|t|)$, którego odwrotność ma postać
$\varphi^{-1}(s)=s/(1-|s|)$ dla $|s|<1$.
Z kroku 2 otrzymujemy ciągłe przedłużenie
$H\colon X\to[-1,1]$ funkcji $\varphi\circ f$.
Zbiór $C=H^{-1}(\{-1,1\})$ jest domknięty w $X$
i rozłączny z $A$, gdyż
$(\varphi\circ f)(A)\subseteq(-1,1)$.
Lemat Urysohna zastosowany do $C$ i $A$
daje ciągłą funkcję $v\colon X\to[0,1]$ taką,
że $v|_C=0$ i $v|_A=1$.
Wówczas $|vH|<1$ w każdym punkcie: na $C$ mamy
$vH=0$, a poza $C$ zachodzi $|H|<1$ i $|v|\le1$.
Stąd $F=\varphi^{-1}\circ(vH)\colon X\to\mathbb R$
jest ciągła, a na $A$ równa
$\varphi^{-1}\circ\varphi\circ f=f$.
\end{proof}

\begin{przyklad}[Przedłużenie funkcji z rozłącznych odcinków]
\label{ex:topog-tietze-odcinki}
Niech $X=[0,4]$ z topologią euklidesową,
$A=[0,1]\cup[3,4]$ oraz
$f\colon A\to[0,1]$ będzie zadana wzorami
$f(x)=0$ dla $x\in[0,1]$ i $f(x)=1$ dla $x\in[3,4]$.
Zbiór $A$ jest domknięty w $X$, a $f$ jest ciągła
w topologii podprzestrzeni, ponieważ $[0,1]$
i $[3,4]$ są rozłącznymi zbiorami otwartymi w $A$.
Jednym z przedłużeń jest funkcja
\[
F(x)=\begin{cases}
0,&0\le x\le1,\\
(x-1)/2,&1<x<3,\\
1,&3\le x\le4.
\end{cases}
\]
Jest ciągła na $X$, gdyż wzory na przedziałach
są ciągłe i przy punktach łączenia $1$ oraz $3$
przyjmują zgodne wartości.
\end{przyklad}

\begin{figure}[htbp]
\centering
\begin{tikzpicture}[x=1.28cm,y=2cm,>=Stealth]
  \draw[->,thin] (-.16,0)--(4.48,0) node[right] {$x$};
  \draw[->,thin] (0,-.12)--(0,1.28) node[above] {$y$};
  \fill[blue!10] (0,-.10) rectangle (1,1.10);
  \fill[blue!10] (3,-.10) rectangle (4,1.10);
  \draw[blue!75!black,line width=1.5pt] (0,0)--(1,0);
  \draw[blue!75!black,line width=1.5pt] (3,1)--(4,1);
  \draw[orange!80!black,line width=1.2pt,dashed] (1,0)--(3,1);
  \fill[blue!70!black] (1,0) circle (.04);
  \fill[blue!70!black] (3,1) circle (.04);
  \foreach \x in {1,2,3,4} {\draw (\x,.025)--(\x,-.025) node[below,font=\small] {$\x$};}
  \node[left,font=\small] at (0,1) {$1$};
  \node[blue!65!black,font=\small] at (.45,.26) {$f$};
  \node[blue!65!black,font=\small] at (3.5,1.16) {$f$};
  \node[orange!70!black,font=\small] at (2.15,.64) {$F$};
\end{tikzpicture}
\caption{Twierdzenie Tietzego: funkcja ciągła zadana na
domkniętym podzbiorze $A=[0,1]\cup[3,4]$ (odcinki niebieskie)
może zostać przedłużona na całe $X=[0,4]$.
Pomarańczowy odcinek obrazuje jedną z możliwych kontynuacji.}
\label{fig:topog-tietze-odcinki}
\end{figure}


\FloatBarrier

\section{Przestrzenie metryczne i dziedziczenie aksjomatów}
\label{sec:topog-metryczne-dziedziczenie}

\begin{twierdzenie}[Każda przestrzeń metryczna jest $T_4$]
\label{thm:topog-metryczne-t4}
Niech $(X,d)$ będzie przestrzenią metryczną z topologią
generowaną przez kule otwarte. Wtedy $X$ spełnia $T_4$.
Co więcej, dla rozłącznych, niepustych, domkniętych
$F,G\subseteq X$ istnieje ciągła $f\colon X\to[0,1]$
taka, że $f|_F=0$ i $f|_G=1$, dana wzorem
\[
f(x)=\frac{d(x,F)}{d(x,F)+d(x,G)},\qquad
d(x,A)=\inf_{a\in A}d(x,a).
\]
\end{twierdzenie}
\begin{proof}
Dla niepustego $A\subseteq X$ funkcja $x\mapsto d(x,A)$
jest $1$-Lipschitzowska, ponieważ nierówność trójkąta daje
$d(x,a)\le d(x,y)+d(y,a)$ dla każdego $a\in A$;
po wzięciu infimum i zamianie $x,y$ otrzymujemy
$|d(x,A)-d(y,A)|\le d(x,y)$. Jest więc ciągła.
Dla domkniętego $A$ mamy $d(x,A)=0$ dokładnie wtedy,
gdy $x\in A$: przy $d(x,A)=0$ w każdej kuli o środku $x$
znajduje się punkt $A$, co oznacza $x\in\overline A=A$.

Dla rozłącznych niepustych domkniętych $F,G$
nie istnieje $x$ z $d(x,F)=d(x,G)=0$, więc mianownik
we wzorze jest dodatni. Funkcja $f$ jest ciągła,
przyjmuje wartości w $[0,1]$, jest równa $0$ na $F$
i $1$ na $G$. Otwarte i rozłączne zbiory
$U=f^{-1}([0,1/3))$ i $V=f^{-1}((2/3,1])$
zawierają odpowiednio $F$ i $G$, zatem $X$ jest normalna.
Gdy jeden z dwóch zbiorów jest pusty, oddzielają je
$\varnothing$ i $X$.
Ponadto pojedynczy punkt jest domknięty,
gdyż jego dopełnienie jest sumą kul
$B(y,d(x,y)/2)$ dla $y\ne x$. Zatem $X$ jest $T_1$
i spełnia $T_4$.
\end{proof}

\begin{stwierdzenie}[Własności dziedziczne]
\label{prop:topog-oddzielanie-dziedziczenie}
Niech $A$ będzie dowolną podprzestrzenią przestrzeni topologicznej $X$.
Jeżeli $X$ spełnia $T_0$, $T_1$, $T_2$, jest regularna
lub całkowicie regularna, to $A$ ma odpowiednią własność.
W szczególności podprzestrzenie przestrzeni $T_3$
i Tychonoffa zachowują te aksjomaty.
Jeśli $X$ jest normalna i $A\subseteq X$ jest domknięty,
to $A$ jest normalna. W konsekwencji domknięte
podprzestrzenie przestrzeni $T_4$ są $T_4$.
\end{stwierdzenie}
\begin{proof}
Dla różnych $x,y\in A$ rozróżniające lub rozłączne
otwarte zbiory w $X$ po przecięciu z $A$ zachowują
wymagane własności. To dowodzi dziedziczenia $T_0,T_1,T_2$.

Niech $F\subseteq A$ będzie domknięty w $A$ i $x\in A\setminus F$.
Z kryterium domkniętości względnej
z rozdziału 04 mamy
$F=A\cap C$ dla pewnego domkniętego $C\subseteq X$.
Wtedy $x\notin C$. Jeśli $X$ jest regularna, rozdzielmy
$x$ i $C$ otwartymi zbiorami $U,V$ w $X$; ich przecięcia
z $A$ oddzielają $x$ i $F$ w $A$. Jeśli $X$ jest całkowicie
regularna, wybierzmy ciągłą $f\colon X\to[0,1]$
z $f(x)=0$, $f|_C=1$. Ograniczenie $f|_A$ jest ciągłe
i oddziela $x$ od $F$.

Jeśli $A$ jest domknięty w $X$, a $F,G$ są rozłącznymi
zbiorami domkniętymi w $A$, to są domknięte także w $X$:
każdy z nich jest przecięciem dwóch domkniętych zbiorów
$A$ i pewnego podzbioru domkniętego $X$. Normalność $X$
pozwala je oddzielić rozłącznymi otwartymi $U,V\subseteq X$,
a $U\cap A,V\cap A$ spełniają warunek normalności w $A$.
\end{proof}

\begin{uwaga}[Ilorazy nie zachowują warunku Hausdorffa]
\label{rem:topog-iloraz-nie-t2}
Rozważmy na $X=\mathbb R$ z topologią euklidesową
surjekcję $q\colon X\to Y=\{0,1\}$ określoną przez
$q(x)=1$ dla $x>0$ i $q(x)=0$ dla $x\le0$.
W topologii ilorazowej na $Y$ zbiór $\{1\}$ jest otwarty,
ponieważ $q^{-1}(\{1\})=(0,\infty)$, lecz $\{0\}$
nie jest otwarty, gdyż $q^{-1}(\{0\})=(-\infty,0]$
nie jest otwarty. Tak więc topologią ilorazową jest
$\{\varnothing,\{1\},Y\}$, czyli topologia Sierpińskiego.
Przestrzeń $X$ jest Hausdorffa, a jej iloraz $Y$ nie jest
nawet $T_1$. Sama ciągła surjekcja nie zachowuje zatem
aksjomatów oddzielania.
\end{uwaga}

\section{Przykład graniczny: płaszczyzna Sorgenfreya}
\label{sec:topog-sorgenfrey-normalnosc}

\begin{twierdzenie}[Iloczyn przestrzeni Tychonoffa nie musi być normalny]
\label{thm:topog-sorgenfrey-plane}
Niech $S$ oznacza prostą rzeczywistą z topologią Sorgenfreya,
której bazę stanowią przedziały $[a,b)$, $a<b$.
Wówczas $S$ jest Tychonoffa, a iloczyn $S\times S$
jest Tychonoffa, lecz nie jest normalny.
\end{twierdzenie}
\begin{proof}
Każdy przedział bazowy $[a,b)$ jest zarazem domknięty w $S$:
jego dopełnienie $(-\infty,a)\cup[b,\infty)$ jest otwarte
w topologii Sorgenfreya. Przestrzeń $S$ jest Hausdorffa,
ponieważ jej topologia jest silniejsza od euklidesowej.
Zatem jest $T_1$ i ma bazę złożoną ze zbiorów
otwarto-domkniętych, czyli jest Tychonoffa
na mocy przykładu \ref{ex:topog-zero-wymiarowy-tych}.
Prostokąty $[a,b)\times[c,d)$ tworzą bazę iloczynu
i także są otwarto-domknięte, gdyż dopełnienie
$C\times D$ dla otwarto-domkniętych $C,D$ jest
$((S\setminus C)\times S)\cup(S\times(S\setminus D))$,
a więc otwarte. Iloczyn jest Hausdorffa jako iloczyn
dwóch przestrzeni Hausdorffa; stosujemy argument
rozdzielania par różniących się co najmniej jedną współrzędną.
Jest więc Tychonoffa.

Pokażemy, że $S\times S$ nie jest normalny.
Zbiór $Q=\mathbb Q\times\mathbb Q$ jest przeliczalny
i gęsty w $S\times S$: każdy niepusty otwarty
prostokąt bazowy zawiera punkt o obu współrzędnych
wymiernych. Natomiast
$D=\{(t,-t):t\in\mathbb R\}$ jest domkniętym
dyskretnym podzbiorem $S\times S$.
Dyskretność wynika z faktu, że prostokąt
$[t,t+\varepsilon)\times[-t,-t+\varepsilon)$ przecina
$D$ dokładnie w punkcie $(t,-t)$: dla $(s,-s)$
należącego do prostokąta jednocześnie $s\ge t$
i $s\le t$. Zbiór $D$ jest domknięty: jeśli
$(a,b)\notin D$, to $a+b\ne0$; wybierając
$\varepsilon>0$ tak małe, że $2\varepsilon<|a+b|$,
widzimy, że na prostokącie
$[a,a+\varepsilon)\times[b,b+\varepsilon)$
suma współrzędnych nie może być zerem.

Załóżmy dla sprzeczności, że $S\times S$ jest normalny.
Dla dowolnego podzbioru $A\subseteq D$ zarówno $A$,
jak i $D\setminus A$ są domknięte w $S\times S$:
każdy podzbiór przestrzeni dyskretnej $D$ jest w niej
domknięty, a $D$ jest domknięty w przestrzeni otaczającej.
Z normalności wybierzmy rozłączne zbiory otwarte
$U_A,V_A$ takie, że $A\subseteq U_A$
i $D\setminus A\subseteq V_A$.

Odwzorowanie $A\mapsto U_A\cap Q$ z $\mathcal P(D)$
do $\mathcal P(Q)$ jest różnowartościowe.
Istotnie, jeśli $A\ne B$, po ewentualnej zamianie
nazw przyjmijmy $(t,-t)\in A\setminus B$.
Wtedy $(t,-t)\in U_A\cap V_B$, zatem
$U_A\cap V_B$ jest niepusty i otwarty.
Gęstość $Q$ daje punkt
$q\in Q\cap U_A\cap V_B$. Ponieważ $U_B\cap V_B=\varnothing$,
mamy $q\in U_A\cap Q$, lecz $q\notin U_B\cap Q$.
Stąd $|\mathcal P(D)|\le|\mathcal P(Q)|$.
Jednak $|D|=|\mathbb R|=2^{\aleph_0}$, podczas gdy
$|Q|=\aleph_0$. Otrzymujemy sprzeczną z twierdzeniem Cantora
nierówność $2^{2^{\aleph_0}}\le2^{\aleph_0}$.
Zatem $S\times S$ nie jest normalny.
\end{proof}

\section{Zestawienie implikacji i zakresu założeń}
\label{sec:topog-oddzielanie-implikacje}

\begin{stwierdzenie}[Hierarchia aksjomatów]
\label{prop:topog-oddzielanie-hierarchia}
W klasie wszystkich przestrzeni topologicznych zachodzi łańcuch
\[
T_4\Longrightarrow T_{3\frac12}\Longrightarrow T_3
\Longrightarrow T_2\Longrightarrow T_1\Longrightarrow T_0.
\]
Wszystkie przestrzenie metryzowalne spełniają $T_4$.
\end{stwierdzenie}
\begin{proof}
Pierwsza implikacja wynika z lematu Urysohna
(\ref{cor:topog-t4-tych}), druga ze stwierdzenia
\ref{prop:topog-calkowita-regularnosc-regularna},
trzecia z wniosku \ref{cor:topog-t3-t2}, a pozostałe
dwie ze stwierdzenia \ref{prop:topog-t2-t1-t0}.
Twierdzenie \ref{thm:topog-metryczne-t4} daje ostatnią tezę.
\end{proof}

\begin{uwaga}[Ścisłość hierarchii i konwencje]
\label{rem:topog-hierarchia-ostra}
Każda z pięciu implikacji w powyższym łańcuchu jest ścisła.
Kontrprzykłady znajdują się w tym rozdziale:
\begin{enumerate}[label=(\roman*)]
\item $T_0\not\Rightarrow T_1$: przestrzeń Sierpińskiego
      (przykład \ref{ex:topog-sierpinski-oddzielanie});
\item $T_1\not\Rightarrow T_2$: nieskończona przestrzeń
      koskończona (przykład \ref{ex:topog-koskonczona-t1-nie-t2});
\item $T_2\not\Rightarrow T_3$: $K$-topologia
      (przykład \ref{ex:topog-k-topologia});
\item $T_3\not\Rightarrow T_{3\frac12}$: przestrzeń
      Mysiora (przykład \ref{ex:topog-mysior-t3-nie-tych});
\item $T_{3\frac12}\not\Rightarrow T_4$:
      płaszczyzna Sorgenfreya
      (twierdzenie \ref{thm:topog-sorgenfrey-plane}).
\end{enumerate}
W literaturze spotyka się odmienne konwencje numeracji:
niektórzy autorzy w definicję regularności lub normalności
włączają już $T_1$. W tym skrypcie pojęcia te definiujemy
bez $T_1$, a aksjomaty $T_3$, $T_{3\frac12}$, $T_4$
wymagają dodatkowo $T_1$ zgodnie z definicjami powyżej.
\end{uwaga}

\begin{figure}[htbp]
\centering
\begin{tikzpicture}[>=Stealth, every node/.style={font=\small}]
\foreach \name/\pos in {$T_4$/0,$T_{3\frac12}$/1,$T_3$/2,$T_2$/3,$T_1$/4,$T_0$/5} {
  \node[draw=blue!55!black,fill=blue!6,rounded corners=2pt,
        minimum width=1.35cm,minimum height=.65cm] (n\pos) at ({-4.65+1.86*\pos},0) {\name};
}
\foreach \a/\b in {0/1,1/2,2/3,3/4,4/5}{\draw[->,thick] (n\a)--(n\b);}
\end{tikzpicture}
\caption{Implikacje aksjomatów oddzielania przy ustalonej konwencji
obejmującej $T_1$ w definicjach $T_3$, $T_{3\frac12}$ i $T_4$.}
\label{fig:topog-oddzielanie-hierarchia}
\end{figure}

\chapter{Zbieżność topologiczna}
\label{ch:topog-zbieznosc}

\noindent\textbf{Cel rozdziału.}
Zbieżność opisuje, w jaki sposób punkty przestrzeni mogą zbliżać się
w sensie jej topologii. W przestrzeniach metrycznych zwykle posługujemy
się ciągami, ale w dowolnej przestrzeni topologicznej same ciągi
nie muszą wykrywać domknięcia ani ciągłości. Zaczniemy od zbieżności
ciągów i jej związku z pierwszym aksjomatem przeliczalności.
Następnie wprowadzimy
sieci (nazywane też ciągami uogólnionymi), które przywracają
pełne kryteria domknięcia i ciągłości, a na końcu filtry,
będące równoważnym językiem opisu granic.
W każdym wyniku jawnie zaznaczymy konieczne założenia.

\section{Zbieżność ciągów i jej podstawowe własności}
\label{sec:topog-zbieznosc-ciagi}

\begin{definicja}[Granica ciągu]
\label{def:topog-granica-ciagu}
Niech $(X,\tau)$ będzie przestrzenią topologiczną, a
$(x_n)_{n\ge1}$ ciągiem punktów $X$. Mówimy, że
\emph{ciąg $(x_n)$ jest zbieżny do $x\in X$}, i zapisujemy
$x_n\to x$, jeżeli dla każdego otoczenia $U$ punktu $x$
istnieje $N\ge1$ takie, że $x_n\in U$ dla wszystkich $n\ge N$.
Wystarczy sprawdzać otoczenia otwarte, ponieważ każde otoczenie
zawiera otwarte otoczenie punktu $x$.
\end{definicja}

\begin{stwierdzenie}[Zależność od bazy otoczeń]
\label{prop:topog-ciag-baza}
Niech $\mathcal B_x$ będzie bazą otoczeń punktu $x$ w przestrzeni
$(X,\tau)$. Ciąg $(x_n)$ jest zbieżny do $x$ wtedy i tylko wtedy,
gdy dla każdego $B\in\mathcal B_x$ wszystkie wyrazy $(x_n)$,
poza skończoną liczbą, należą do $B$.
\end{stwierdzenie}
\begin{proof}
Konieczność wynika bezpośrednio z definicji granicy. Odwrotnie,
niech $U$ będzie otoczeniem $x$. Z definicji bazy istnieje
$B\in\mathcal B_x$ zawarty w $U$. Od pewnego indeksu
$x_n\in B\subseteq U$, co dowodzi zbieżności.
\end{proof}

\begin{stwierdzenie}[Kryterium metryczne]
\label{prop:topog-ciag-metryka}
Jeżeli topologia przestrzeni $X$ pochodzi od metryki $d$,
to $x_n\to x$ jest równoważne warunkowi:
\[
\forall\varepsilon>0\ \exists N\ge1\ \forall n\ge N:
\quad d(x_n,x)<\varepsilon.
\]
\end{stwierdzenie}
\begin{proof}
Kule $B_d(x,\varepsilon)$, $\varepsilon>0$, tworzą bazę otoczeń
$x$. Zastosowanie stwierdzenia \ref{prop:topog-ciag-baza}
daje dokładnie zapisany warunek.
\end{proof}

\begin{przyklad}[Ciąg liczb odwrotnych]
\label{ex:topog-ciag-odwrotnosci}
W prostej euklidesowej $x_n=1/n$ jest zbieżny do $0$.
Istotnie, dla $\varepsilon>0$ wybieramy liczbę naturalną
$N>1/\varepsilon$. Wówczas dla $n\ge N$ zachodzi
$|1/n-0|=1/n\le1/N<\varepsilon$.
Zauważmy, że żaden wyraz ciągu nie jest równy granicy.
\end{przyklad}

\begin{figure}[htbp]
\centering
\begin{tikzpicture}[x=1.3cm,y=1cm,>=Stealth]
  \fill[blue!9] (-1.20,-.48) rectangle (1.20,.48);
  \draw[blue!70!black,dashed] (-1.20,-.62)--(-1.20,.65);
  \draw[blue!70!black,dashed] (1.20,-.62)--(1.20,.65);
  \draw[->,thick] (-2,0)--(3.4,0);
  \node[below] at (-1.20,-.58) {$-\varepsilon$};
  \node[below] at (1.20,-.58) {$\varepsilon$};
  \fill (0,0) circle (2pt) node[below left] {$0$};
  \fill[orange!85!black] (3,0) circle (2.5pt) node[above] {$x_1$};
  \fill[orange!85!black] (1.5,0) circle (2.5pt) node[above right] {$x_2$};
  \fill[blue!70!black] (1,0) circle (2.5pt) node[above] {$x_3$};
  \fill[blue!70!black] (.60,0) circle (2.5pt) node[above] {$x_5$};
  \fill[blue!70!black] (.30,0) circle (2.5pt);
  \fill[blue!70!black] (.15,0) circle (2.5pt);
  \node[above,blue!70!black] at (0,.62) {$B(0,\varepsilon)$};
  \node[below,font=\scriptsize] at (.25,-1.05) {wyrazy dalsze};
\end{tikzpicture}
\caption{Schemat zbieżności $x_n=1/n$ do $0$: dla ustalonej
kuli wokół granicy wszystkie dostatecznie późne wyrazy
znajdują się wewnątrz. Na rysunku wybrano
$\varepsilon=2/5$ po przeskalowaniu poziomej osi.}
\label{fig:topog-ciag-kula}
\end{figure}

\begin{stwierdzenie}[Podciągi i obrazy ciągłe]
\label{prop:topog-ciag-podciag-ciaglosc}
Niech $(X,\tau_X)$ i $(Y,\tau_Y)$ będą przestrzeniami topologicznymi.
\begin{enumerate}[label=(\roman*)]
\item Jeżeli $x_n\to x$, a $n_1<n_2<\cdots$,
      to $x_{n_k}\to x$.
\item Jeżeli odwzorowanie $f\colon X\to Y$ jest ciągłe w $x$
      oraz $x_n\to x$, to $f(x_n)\to f(x)$.
\end{enumerate}
\end{stwierdzenie}
\begin{proof}
(i) Każde otoczenie $U$ punktu $x$ zawiera wszystkie $x_n$
dla $n\ge N$. Ponieważ $n_k\ge k$, mamy $n_k\ge N$
dla $k\ge N$, a więc $x_{n_k}\in U$.
(ii) Niech $V$ będzie otwartym otoczeniem $f(x)$. Z ciągłości
w $x$ istnieje otwarte $U\ni x$ takie, że $f(U)\subseteq V$.
Od pewnego indeksu $x_n\in U$, zatem $f(x_n)\in V$.
\end{proof}

\begin{twierdzenie}[Jednoznaczność granicy w przestrzeniach Hausdorffa]
\label{thm:topog-ciag-hausdorff}
Jeżeli $X$ jest przestrzenią Hausdorffa, to każdy ciąg w $X$
ma co najwyżej jedną granicę.
\end{twierdzenie}
\begin{proof}
Załóżmy, że $x_n\to x$ oraz $x_n\to y$, gdzie $x\ne y$.
Istnieją rozłączne otwarte otoczenia $U$ punktu $x$ i $V$
punktu $y$. Dla dostatecznie dużego $n$ wyraz $x_n$ należy
jednocześnie do $U$ i $V$, co przeczy $U\cap V=\varnothing$.
\end{proof}

\begin{przyklad}[Brak jednoznaczności]
\label{ex:topog-ciag-wiele-granic}
W dwupunktowej przestrzeni antydyskretnej
$X=\{a,b\}$ jedynym otwartym otoczeniem dowolnego punktu
jest $X$. Dlatego każdy ciąg, w tym ciąg stały $x_n=a$,
zbiega zarówno do $a$, jak i do $b$. Nie można więc
opuszczać aksjomatu Hausdorffa w
 twierdzeniu \ref{thm:topog-ciag-hausdorff}.
\end{przyklad}

\section{Co mogą, a czego nie mogą wykryć ciągi}
\label{sec:topog-ciagi-ograniczenia}

\begin{definicja}[Domknięcie ciągowe]
\label{def:topog-domkniecie-ciagowe}
Dla podzbioru $A\subseteq X$ definiujemy
\emph{zbiór granic ciągów o wyrazach w $A$}:
\[
\operatorname{scl}(A)=
\{x\in X:\ \exists(x_n)_{n\ge1}\in A^{\mathbb N},\ x_n\to x\}.
\]
Zbiór $A$ nazywamy \emph{ciągowo domkniętym}, gdy
$\operatorname{scl}(A)=A$.
\end{definicja}

\begin{stwierdzenie}[Relacja z domknięciem]
\label{prop:topog-ciagi-domkniecie}
Dla każdego $A\subseteq X$ zachodzi
$A\subseteq\operatorname{scl}(A)\subseteq\overline A$.
W szczególności każdy zbiór domknięty jest ciągowo domknięty.
\end{stwierdzenie}
\begin{proof}
Jeżeli $a\in A$, to ciąg stały $x_n=a$ zbiega do $a$,
stąd pierwsze zawieranie. Jeżeli $x_n\in A$ i $x_n\to x$,
to każde otoczenie $U$ punktu $x$ zawiera pewne $x_n$,
a więc $U\cap A\ne\varnothing$. Z kryterium domknięcia
wynika $x\in\overline A$. Jeśli $A$ jest domknięty,
$\overline A=A$, więc oba zawierania stają się równościami.
\end{proof}

\begin{definicja}[Pierwszy aksjomat przeliczalności]
\label{def:topog-pierwsza-przeliczalnosc}
Mówimy, że przestrzeń topologiczna $X$ \emph{spełnia pierwszy
aksjomat przeliczalności}, jeżeli każdy punkt $x\in X$
ma co najwyżej przeliczalną bazę otoczeń.
Warunek jest automatycznie spełniony dla $X=\varnothing$.
\end{definicja}

\begin{przyklad}[Przestrzenie metryczne]
\label{ex:topog-metryczna-pierwsza-przeliczalnosc}
W każdej przestrzeni metrycznej baza
$\{B_d(x,1/n):n\ge1\}$ jest przeliczalna.
Jeżeli $U$ jest otwartym otoczeniem $x$, istnieje $r>0$
z $B_d(x,r)\subseteq U$; wybór $n$ takiego, że $1/n<r$,
daje $B_d(x,1/n)\subseteq U$. Każda przestrzeń metryczna
spełnia więc pierwszy aksjomat przeliczalności.
\end{przyklad}

\begin{twierdzenie}[Ciągi wykrywają domknięcie przy pierwszym aksjomacie przeliczalności]
\label{thm:topog-pierwsza-przeliczalnosc-domkniecie}
Niech $X$ będzie przestrzenią topologiczną spełniającą
pierwszy aksjomat przeliczalności,
a $A\subseteq X$. Wówczas
\[
 x\in\overline A
 \quad\Longleftrightarrow\quad
 \text{istnieje ciąg }(a_n)_{n\ge1}\text{ w }A
 \text{ taki, że }a_n\to x.
\]
W szczególności $A$ jest domknięty wtedy i tylko wtedy,
gdy jest ciągowo domknięty.
\end{twierdzenie}
\begin{proof}
Implikację od prawej do lewej wykazaliśmy w
stwierdzeniu \ref{prop:topog-ciagi-domkniecie}.
Niech teraz $x\in\overline A$ i niech
$\{B_n:n\ge1\}$ będzie bazą otoczeń $x$
(w przypadku bazy skończonej numerujemy jej elementy
z powtórzeniami). Zastępując $B_n$ przez $W_n=B_1\cap\cdots\cap B_n$,
otrzymujemy malejącą bazę otoczeń punktu $x$.
Każdy $W_n$ jest otoczeniem $x$ i przecina $A$,
ponieważ $x\in\overline A$.
Wybieramy $a_n\in A\cap W_n$. Dla dowolnego otoczenia
$U$ punktu $x$ istnieje $k$ z $B_k\subseteq U$.
Gdy $n\ge k$, mamy $a_n\in W_n\subseteq B_k\subseteq U$.
Zatem $a_n\to x$.
Wreszcie, zbiór $A$ jest domknięty dokładnie wtedy,
gdy $\overline A=A$, a na mocy dowiedzionej równoważności
$\overline A=\operatorname{scl}(A)$.
\end{proof}

\begin{twierdzenie}[Ciągowe kryterium ciągłości]
\label{thm:topog-ciagowe-kryterium-ciaglosci}
Niech $X$ będzie przestrzenią spełniającą
pierwszy aksjomat przeliczalności,
$Y$ dowolną przestrzenią topologiczną, $f\colon X\to Y$
i $x\in X$. Odwzorowanie $f$ jest ciągłe w $x$
wtedy i tylko wtedy, gdy dla każdego ciągu $x_n\to x$
zachodzi $f(x_n)\to f(x)$.
\end{twierdzenie}
\begin{proof}
Implikacja wprost wynika ze stwierdzenia
\ref{prop:topog-ciag-podciag-ciaglosc}.
Załóżmy warunek ciągowy i przypuśćmy, że $f$
nie jest ciągłe w $x$. Istnieje otwarte otoczenie
$V$ punktu $f(x)$ takie, że żadne otoczenie $U$
punktu $x$ nie jest zawarte w $f^{-1}(V)$.
Niech $W_n$ będą malejącą bazą otoczeń $x$,
jak w dowodzie twierdzenia
\ref{thm:topog-pierwsza-przeliczalnosc-domkniecie}.
Dla każdego $n$ możemy wybrać
$x_n\in W_n\setminus f^{-1}(V)$.
Wówczas $x_n\to x$, lecz wszystkie $f(x_n)$ leżą poza
$V$, zatem $f(x_n)$ nie może zbiegać do $f(x)$.
Sprzeczność dowodzi ciągłości $f$ w $x$.
\end{proof}

\begin{przyklad}[Topologia koprzeliczalna: ciągi nie wystarczają]
\label{ex:topog-ciagi-koprzeliczalna}
Niech $X=\R$ z topologią koprzeliczalną:
zbiór jest otwarty wtedy i tylko wtedy, gdy jest pusty
lub jego dopełnienie jest co najwyżej przeliczalne.
\begin{enumerate}[label=(\roman*)]
\item Jeżeli $x_n\to x$, to $(x_n)$ jest od pewnego miejsca
stały, równy $x$. Istotnie, zbiór
$C=\{x_n:n\ge1,\ x_n\ne x\}$ jest przeliczalny i nie
zawiera $x$. Otoczenie otwarte $X\setminus C$ punktu $x$
musi zawierać wszystkie wyrazy dostatecznie późne,
a jedyną możliwą ich wartością jest $x$.
Odwrotnie, każdy ciąg ostatecznie stały równy $x$
zbiega do $x$.
\item Każdy podzbiór $A\subseteq X$ jest ciągowo domknięty:
każdy ciąg o wyrazach w $A$ i granicy $x$ jest ostatecznie
stały równy $x\in A$.
Jednak nie każdy podzbiór jest domknięty. Dla przykładu
$A=(0,\infty)$ ma domknięcie $X$, ponieważ przecina
każdy niepusty zbiór otwarty (ten ma przeliczalne dopełnienie,
a $A$ jest nieprzeliczalny).
\item Identyczność $f\colon(X,\tau_{\mathrm{koprz}})
\to(X,\tau_{\mathrm{dysk}})$ zachowuje granice wszystkich
ciągów, bo są one ostatecznie stałe. Nie jest jednak
ciągła: przeciwobraz jednoelementowego zbioru otwartego
w topologii dyskretnej nie jest otwarty w topologii
koprzeliczalnej.
\item Granica dowolnego zbieżnego ciągu jest jednoznaczna,
chociaż $X$ nie jest przestrzenią Hausdorffa:
dwa niepuste zbiory otwarte zawsze się przecinają,
gdyż suma dwóch zbiorów przeliczalnych nie wyczerpuje $X$.
\end{enumerate}
Przestrzeń ta nie spełnia pierwszego aksjomatu przeliczalności:
gdyby była, z twierdzenia
\ref{thm:topog-pierwsza-przeliczalnosc-domkniecie}
wynikałoby, że ciągowo domknięty zbiór $A$ jest domknięty,
wbrew punktowi (ii).
Przykład pokazuje więc, że bez pierwszego aksjomatu przeliczalności
ciągi nie dają ogólnego kryterium domknięcia ani ciągłości.
\end{przyklad}

\section{Sieci i zbieżność Moore'a--Smitha}
\label{sec:topog-sieci}

\begin{definicja}[Zbiór skierowany]
\label{def:topog-zbior-skierowany}
\emph{Zbiorem skierowanym} nazywamy niepusty zbiór $D$
wyposażony w relację $\preceq$ zwrotną, przechodnią
i spełniającą warunek: dla każdych $d,e\in D$ istnieje
$c\in D$ takie, że $d\preceq c$ oraz $e\preceq c$.
Nie wymagamy antysymetrii relacji.
\end{definicja}

\begin{definicja}[Sieć i jej granica]
\label{def:topog-siec-granica}
Niech $(D,\preceq)$ będzie zbiorem skierowanym,
a $X$ przestrzenią topologiczną. \emph{Siecią w $X$}, czyli \emph{ciągiem uogólnionym},
nazywamy rodzinę $(x_d)_{d\in D}$ punktów $X$.
Sieć jest \emph{zbieżna do $x\in X$}, gdy dla każdego
otoczenia $U$ punktu $x$ istnieje $d_0\in D$ takie,
że $x_d\in U$ dla wszystkich $d\succeq d_0$.
Zapisujemy wtedy $x_d\to x$.
\end{definicja}

\begin{przyklad}[Ciąg jako szczególna sieć]
\label{ex:topog-ciag-jako-siec}
Zbiór $\mathbb N=\{1,2,\ldots\}$ z naturalnym porządkiem
jest skierowany. Definicja granicy sieci o indeksach
w $\mathbb N$ jest dokładnie definicją granicy ciągu.
\end{przyklad}

\begin{przyklad}[Indeksowanie skończonymi zbiorami]
\label{ex:topog-siec-skonczone}
Niech $I$ będzie dowolnym niepustym zbiorem.
Rodzina $\mathcal D=\{F\subseteq I:F\text{ jest skończony}\}$
z porządkiem przez zawieranie $\subseteq$ jest skierowana:
dla $F,G$ wspólnym ograniczeniem górnym jest $F\cup G$.
Dzięki temu sieci można indeksować skończonymi
zbiorami nawet wtedy, gdy $I$ jest nieprzeliczalny.
\end{przyklad}

\begin{figure}[htbp]
\centering
\begin{tikzpicture}[>=Stealth, node distance=1.3cm,
  every node/.style={draw,rounded corners=2pt,fill=blue!5,inner sep=5pt}]
  \node (empty) at (0,0) {$\varnothing$};
  \node (a) at (-2,1.25) {$\{a\}$};
  \node (b) at (2,1.25) {$\{b\}$};
  \node (ab) at (0,2.6) {$\{a,b\}$};
  \node (abc) at (0,3.9) {$\{a,b,c\}$};
  \draw[->] (empty)--(a);
  \draw[->] (empty)--(b);
  \draw[->] (a)--(ab);
  \draw[->] (b)--(ab);
  \draw[->] (ab)--(abc);
\end{tikzpicture}
\caption{Fragment zbioru skierowanego skończonych podzbiorów
$I$: dwa indeksy $\{a\}$ i $\{b\}$ mają wspólne ograniczenie
górne $\{a,b\}$. Nie przedstawiono wszystkich indeksów
ani wszystkich strzałek relacji zawierania.}
\label{fig:topog-skierowanie}
\end{figure}

\begin{stwierdzenie}[Obraz sieci zbieżnej]
\label{prop:topog-siec-ciaglosc-jeden-kierunek}
Jeżeli $f\colon X\to Y$ jest ciągłe w $x\in X$
oraz sieć $(x_d)_{d\in D}$ zbiega do $x$, to
$f(x_d)\to f(x)$ w $Y$.
\end{stwierdzenie}
\begin{proof}
Dla dowolnego otoczenia $V$ punktu $f(x)$
z ciągłości $f$ w $x$ wynika, że $f^{-1}(V)$ jest
otoczeniem $x$. Stąd istnieje $d_0$ takie, że
dla $d\succeq d_0$ zachodzi
$x_d\in f^{-1}(V)$, czyli $f(x_d)\in V$.
\end{proof}

\begin{twierdzenie}[Domknięcie jest wykrywane przez sieci]
\label{thm:topog-siec-domkniecie}
Niech $X$ będzie dowolną przestrzenią topologiczną,
$A\subseteq X$, $x\in X$. Wtedy
\[
 x\in\overline A
 \quad\Longleftrightarrow\quad
 \text{istnieje sieć $(a_d)_{d\in D}$ o wyrazach w $A$,
 zbieżna do $x$}.
\]
W szczególności $A$ jest domknięty wtedy i tylko wtedy,
gdy granica każdej zbieżnej sieci o wyrazach w $A$
należy do $A$.
\end{twierdzenie}
\begin{proof}
Jeżeli sieć $(a_d)$ o wyrazach w $A$ zbiega do $x$,
to każde otoczenie $U$ punktu $x$ zawiera jakiś
$a_d\in A$, więc $x\in\overline A$.

Odwrotnie, załóżmy $x\in\overline A$. Utwórzmy
\[
 D=\{(U,a):\ U\text{ jest otwartym otoczeniem }x,
                         \ a\in U\cap A\}.
\]
Zbiór $D$ jest niepusty: $X$ jest otwartym otoczeniem
$x$ i $A\ne\varnothing$ (ponieważ $x\in\overline A$).
Określamy $(U,a)\preceq(V,b)$, gdy $V\subseteq U$.
Relacja jest zwrotna i przechodnia. Jeśli
$(U,a),(V,b)\in D$, to otoczenie $U\cap V$ przecina $A$;
wybieramy $c\in U\cap V\cap A$ i otrzymujemy
$(U\cap V,c)$, które jest ograniczeniem górnym obu indeksów.
Zatem $D$ jest skierowany.
Sieć $a_{(U,a)}=a$ ma wartości w $A$.
Dla dowolnego otwartego otoczenia $W$ punktu $x$
wybieramy $a\in W\cap A$ i indeks $d_0=(W,a)$.
Gdy $(V,b)\succeq d_0$, mamy $V\subseteq W$
i $a_{(V,b)}=b\in V\subseteq W$. Sieć zbiega do $x$.

Ostatnia równoważność wynika z wykazanej charakterystyki:
$A$ jest domknięty dokładnie wtedy, gdy
$\overline A=A$.
\end{proof}

\begin{uwaga}[Unikanie niejawnego wyboru]
\label{rem:topog-siec-bez-wyboru}
Indeksy $(U,a)$ obejmują \emph{wszystkie} dopuszczalne
pary, dlatego do skonstruowania sieci nie musimy
jednocześnie wybrać po jednym punkcie z każdej rodziny
$U\cap A$. Podobne indeksowanie wykorzystamy w dowodzie
jednoznaczności granic.
\end{uwaga}

\begin{twierdzenie}[Ciągłość przez zachowanie granic sieci]
\label{thm:topog-siec-kryterium-ciaglosci}
Niech $X,Y$ będą dowolnymi przestrzeniami topologicznymi,
$f\colon X\to Y$ oraz $x\in X$. Wtedy $f$ jest ciągłe w $x$
dokładnie wtedy, gdy dla każdej sieci $(x_d)$ zbieżnej
do $x$ sieć $(f(x_d))$ zbiega do $f(x)$.
\end{twierdzenie}
\begin{proof}
Konieczność wynika ze stwierdzenia
\ref{prop:topog-siec-ciaglosc-jeden-kierunek}.
Jeśli $f$ nie jest ciągłe w $x$, istnieje otwarty
zbiór $V\ni f(x)$ taki, że $f^{-1}(V)$ nie jest
otoczeniem $x$. Wtedy każde otwarte otoczenie $U$
punktu $x$ przecina zbiór
$A=X\setminus f^{-1}(V)$, czyli $x\in\overline A$.
Z twierdzenia \ref{thm:topog-siec-domkniecie} istnieje
sieć $(a_d)$ o wyrazach w $A$ zbieżna do $x$.
Jednak $f(a_d)\notin V$ dla każdego $d$, więc sieć
$f(a_d)$ nie zbiega do $f(x)$: przeczy temu otoczenie $V$.
Dowodzi to implikacji odwrotnej.
\end{proof}

\begin{twierdzenie}[Jednoznaczność granic sieci i aksjomat Hausdorffa]
\label{thm:topog-sieci-hausdorff}
Dla przestrzeni topologicznej $X$ równoważne są:
\begin{enumerate}[label=(\roman*)]
\item $X$ jest przestrzenią Hausdorffa;
\item każda sieć w $X$ ma co najwyżej jedną granicę.
\end{enumerate}
\end{twierdzenie}
\begin{proof}
Załóżmy (i) i niech $(x_d)$ zbiega do różnych
punktów $x,y\in X$. Wybieramy rozłączne otwarte
otoczenia $U\ni x$ i $V\ni y$. Istnieją indeksy
$d_U,d_V$ takie, że odpowiednie ogony sieci są
zawarte w $U$ i $V$. Z kierunkowości istnieje
$d_0\succeq d_U,d_V$, a wtedy $x_{d_0}\in U\cap V$,
co jest sprzeczne z rozłącznością.

Niech teraz $X$ nie będzie Hausdorffa. Istnieją różne
$x,y\in X$ takie, że każde otwarte otoczenia
$U\ni x$ i $V\ni y$ przecinają się.
Rozważmy wszystkie trójki
\[
 D=\{(U,V,z):\ U\ni x,\ V\ni y\text{ są otwarte},
                                 \ z\in U\cap V\}.
\]
Porządek określamy przez
$(U,V,z)\preceq(U',V',z')$, jeśli
$U'\subseteq U$ i $V'\subseteq V$.
Indeksy mają wspólne ograniczenia górne: dla dwóch
trójek bierzemy otoczenia $U\cap U'$ i $V\cap V'$
oraz dowolny punkt ich niepustego przecięcia.
Sieć $z_{(U,V,z)}=z$ zbiega zarówno do $x$, jak i $y$.
Istotnie, dla otwartego otoczenia $W\ni x$
możemy rozpocząć od trójki $(W,X,z)$,
a wszystkie jej późniejsze indeksy mają pierwszą
współrzędną zawartą w $W$, więc odpowiadające im
punkty należą do $W$. Analogicznie dla $y$,
zaczynając od $(X,W',z)$.
Zatem (ii) nie zachodzi, co kończy dowód.
\end{proof}

\begin{uwaga}
W przeciwieństwie do sieci, jednoznaczność granic
\emph{samych ciągów} nie charakteryzuje aksjomatu
Hausdorffa: kontrprzykład daje topologia koprzeliczalna
z przykładu \ref{ex:topog-ciagi-koprzeliczalna}.
\end{uwaga}

\pagebreak[4]
\begin{przyklad}[Sieć zbieżna, której nie można zastąpić ciągiem]
\label{ex:topog-siec-koprzeliczalna-bez-ciagu}
Rozważmy $X=\R$ z topologią koprzeliczalną,
$A=(0,\infty)$ i punkt $0\notin A$.
Jak pokazaliśmy w przykładzie
\ref{ex:topog-ciagi-koprzeliczalna},
$0\in\overline A$, lecz żaden ciąg o wyrazach w $A$
nie zbiega do $0$: każdy ciąg zbieżny do $0$ musiałby
być od pewnego miejsca stale równy $0$.
Skonstruujemy natomiast konkretną sieć w $A$ zbieżną do $0$.

Niech $D$ będzie zbiorem wszystkich par $(C,a)$, w których
$C\subseteq A$ jest co najwyżej przeliczalny,
a $a\in A\setminus C$. Określamy
$(C,a)\preceq(E,b)$, gdy $C\subseteq E$.
Zbiór $D$ jest niepusty, gdyż $(\varnothing,a)\in D$
dla każdego $a\in A$. Relacja jest zwrotna i przechodnia.
Dla dowolnych $(C,a),(E,b)\in D$ suma $C\cup E$
jest przeliczalna, a $A$ jest nieprzeliczalny,
więc istnieje $c\in A\setminus(C\cup E)$.
Wówczas $(C\cup E,c)$ jest wspólnym ograniczeniem
górnym tych indeksów. Zatem $D$ jest skierowany.
Nie wymagamy tu antysymetrii relacji (zob.
definicję \ref{def:topog-zbior-skierowany}).

Niech $x_{(C,a)}=a$. Zbiór wartości tej sieci jest
dokładnie równy $A$, gdyż $(\varnothing,a)\in D$ dla
każdego $a\in A$. Dla dowolnego otwartego otoczenia
$U$ punktu $0$ zbiór $C=A\setminus U$ jest co najwyżej
przeliczalny. Wybieramy $a\in A\setminus C$
i indeks $d_0=(C,a)$. Gdy $(E,b)\succeq d_0$,
to $C\subseteq E$, zatem
$x_{(E,b)}=b\in A\setminus E\subseteq A\setminus C\subseteq U$.
To dowodzi $x_d\to0$, mimo że żadna z wartości sieci
nie jest równa $0$ i żaden ciąg w $A$ nie zbiega do $0$.
Konstrukcja nie wymaga wyboru jednego punktu dla każdego
zbioru $C$: wszystkie dopuszczalne pary stanowią indeksy.
\end{przyklad}

\section{Podsieci, punkty skupienia sieci i iloczyny}
\label{sec:topog-podsieci}

\begin{definicja}[Podsieć]
\label{def:topog-podsiec}
Niech $(x_d)_{d\in D}$ będzie siecią, a $E$ innym
zbiorem skierowanym. Sieć $(x_{\phi(e)})_{e\in E}$
nazywamy \emph{podsiecią} $(x_d)$, jeśli
$\phi\colon E\to D$ spełnia warunek:
dla każdego $d_0\in D$ istnieje $e_0\in E$ takie,
że $\phi(e)\succeq d_0$ dla wszystkich $e\succeq e_0$.
Nie wymagamy, aby $\phi$ była iniekcją ani
monotoniczna; wymagamy jedynie wskazanej ostatecznej
współkońcowości.
\end{definicja}

\begin{stwierdzenie}[Granice podsieci]
\label{prop:topog-podsiec-granica}
Każda podsieć sieci zbieżnej do $x$ jest zbieżna do $x$.
\end{stwierdzenie}
\begin{proof}
Niech $U$ będzie otoczeniem $x$ i niech
$x_d\in U$ dla wszystkich $d\succeq d_0$.
Z warunku na $\phi$ istnieje $e_0$ takie,
że $\phi(e)\succeq d_0$ dla wszystkich $e\succeq e_0$.
Stąd $x_{\phi(e)}\in U$ dla wszystkich $e\succeq e_0$.
\end{proof}

\begin{definicja}[Punkt skupienia sieci]
\label{def:topog-punkt-skupienia-sieci}
Punkt $x\in X$ nazywamy \emph{punktem skupienia sieci}
$(x_d)_{d\in D}$, jeśli dla każdego otoczenia $U$
punktu $x$ i każdego $d_0\in D$ istnieje
$d\succeq d_0$ takie, że $x_d\in U$.
\end{definicja}

\begin{twierdzenie}[Punkt skupienia a zbieżna podsieć]
\label{thm:topog-punkt-skupienia-podsiec}
Niech $(x_d)_{d\in D}$ będzie siecią w przestrzeni
$X$, a $x\in X$. Wtedy $x$ jest punktem skupienia
sieci dokładnie wtedy, gdy istnieje jej podsieć
zbieżna do $x$.
\end{twierdzenie}
\begin{proof}
Jeżeli podsieć $(x_{\phi(e)})$ zbiega do $x$,
to dla otoczenia $U$ istnieje indeks $e_U$,
od którego wszystkie jej wyrazy leżą w $U$.
Dla zadanego $d_0\in D$ istnieje też $e_d$,
od którego $\phi(e)\succeq d_0$.
Wybierając wspólne ograniczenie górne dla $e_U,e_d$,
otrzymujemy $d=\phi(e)\succeq d_0$ z $x_d\in U$.

Przypuśćmy odwrotnie, że $x$ jest punktem skupienia.
Niech
\[
 E=\{(d,U):d\in D,\ U\text{ jest otwartym otoczeniem }x,
                                        \ x_d\in U\}.
\]
Zbiór $E$ jest niepusty, gdyż dla $U=X$
i dowolnego $d\in D$ mamy $(d,X)\in E$.
Określmy $(d,U)\preceq(e,V)$, gdy $d\preceq e$
oraz $V\subseteq U$. Jest to relacja zwrotna i przechodnia.
Dla $(d,U),(e,V)\in E$ bierzemy $d_0\succeq d,e$.
Z własności punktu skupienia istnieje $c\succeq d_0$
takie, że $x_c\in U\cap V$.
Indeks $(c,U\cap V)$ jest wspólnym ograniczeniem górnym,
a więc $E$ jest skierowany.
Definiujemy $\phi(d,U)=d$. Dla każdego $d_0\in D$
wybieramy $c\succeq d_0$ i zauważamy, że
$(c,X)\in E$. Dla wszystkich późniejszych
$(d,U)\in E$ zachodzi $d\succeq c\succeq d_0$.
Zatem $\phi$ spełnia warunek podsieci.
Wreszcie, dla otwartego otoczenia $W\ni x$
wybieramy $d$ z $x_d\in W$ (istnieje z własności
punktu skupienia) i indeks $(d,W)\in E$.
Dla każdego późniejszego $(e,V)$ mamy $V\subseteq W$,
a ponieważ $x_e\in V$, zachodzi $x_e\in W$.
Podsieć $(x_{\phi(d,U)})$ zbiega do $x$.
\end{proof}

\begin{twierdzenie}[Zbieżność w iloczynie]
\label{thm:topog-siec-iloczyn}
Niech $\{X_i:i\in I\}$ będzie rodziną przestrzeni
topologicznych, a $X=\prod_{i\in I}X_i$ jej iloczynem
z topologią produktową. Niech $(x_d)_{d\in D}$
będzie siecią w $X$, a $x\in X$. Wówczas
\[
 x_d\to x\text{ w }X
 \quad\Longleftrightarrow\quad
 \pi_i(x_d)\to\pi_i(x)\text{ w }X_i
 \text{ dla każdego }i\in I.
\]
\end{twierdzenie}
\begin{proof}
Implikacja wprost wynika z ciągłości rzutowań
$\pi_i\colon X\to X_i$, wykazanej przy konstrukcji
 topologii produktowej.
Odwrotnie, niech wszystkie sieci współrzędnych
zbiegają do właściwych granic. Wystarczy sprawdzić
zbieżność na bazie otoczeń punktu $x$.
Jej elementami są zbiory postaci
$B=\bigcap_{i\in F}\pi_i^{-1}(U_i)$,
gdzie $F\subseteq I$ jest skończony, a $U_i$
jest otwartym otoczeniem $\pi_i(x)$.
Dla każdego $i\in F$ istnieje $d_i\in D$ takie,
że $\pi_i(x_d)\in U_i$ dla wszystkich $d\succeq d_i$.
Z kierunkowości $D$ i skończoności $F$ istnieje
wspólne ograniczenie górne $d_0$ tych indeksów
(dla $F=\varnothing$ wystarczy dowolny $d_0\in D$).
Wtedy dla $d\succeq d_0$ zachodzi
$x_d\in B$. To dowodzi zbieżności w iloczynie.
\end{proof}

\section{Filtry i ich zbieżność}
\label{sec:topog-filtry}

\begin{definicja}[Filtr właściwy]
\label{def:topog-filtr}
Niech $X$ będzie zbiorem. \emph{Filtrem (właściwym)}
na $X$ nazywamy rodzinę $\mathcal F\subseteq\mathcal P(X)$
spełniającą następujące warunki:
\begin{enumerate}[label=(\roman*)]
\item $X\in\mathcal F$ i $\varnothing\notin\mathcal F$;
\item jeżeli $F,G\in\mathcal F$, to $F\cap G\in\mathcal F$;
\item jeśli $F\in\mathcal F$ i $F\subseteq B\subseteq X$,
      to $B\in\mathcal F$.
\end{enumerate}
\emph{Bazą filtru} nazywamy niepustą rodzinę
$\mathcal B\subseteq\mathcal P(X)\setminus\{\varnothing\}$
taką, że dla $B,C\in\mathcal B$ istnieje
$D\in\mathcal B$ z $D\subseteq B\cap C$.
Generuje ona filtr $\mathcal F_{\mathcal B}$ wszystkich
zbiorów zawierających co najmniej jeden element
$\mathcal B$.
\end{definicja}

\begin{proof}[Uzasadnienie konstrukcji filtru z bazy]
Rodzina $\mathcal F_{\mathcal B}$ zawiera $X$,
a nie zawiera $\varnothing$, ponieważ elementy bazy
są niepuste. Jeśli $F$ zawiera $B\in\mathcal B$, a $G$
zawiera $C\in\mathcal B$, to $F\cap G$ zawiera
odpowiedni $D\subseteq B\cap C$, więc należy
do $\mathcal F_{\mathcal B}$. Własność rozszerzania
zbiorów jest wprost wpisana w definicję.
\end{proof}

\begin{definicja}[Granica filtru]
\label{def:topog-filtr-granica}
Niech $X$ będzie przestrzenią topologiczną,
$\mathcal F$ filtrem na $X$ i $x\in X$.
Mówimy, że $\mathcal F$ \emph{zbiega do $x$}, gdy
każde otoczenie punktu $x$ należy do $\mathcal F$.
Zapisujemy $\mathcal F\to x$.
\end{definicja}

\begin{przyklad}[Filtr otoczeń]
\label{ex:topog-filtr-otoczen}
Dla każdego $x\in X$ rodzina wszystkich otoczeń
punktu $x$ jest filtrem na $X$: jest niepusta,
nie zawiera $\varnothing$, jest zamknięta na skończone
przecięcia i na rozszerzanie zbiorów.
Filtr ten zbiega do $x$ z samej definicji.
\end{przyklad}

\begin{twierdzenie}[Domknięcie przez filtry]
\label{thm:topog-filtr-domkniecie}
Dla dowolnej przestrzeni topologicznej $X$,
zbioru $A\subseteq X$ i punktu $x\in X$ warunek
$x\in\overline A$ jest równoważny istnieniu
filtru $\mathcal F$ na $X$ takiego, że
$A\in\mathcal F$ i $\mathcal F\to x$.
\end{twierdzenie}
\begin{proof}
Załóżmy $x\in\overline A$. Niech $\mathcal B$
będzie rodziną zbiorów $A\cap U$, gdzie $U$
przebiega wszystkie otwarte otoczenia $x$.
Każdy element $\mathcal B$ jest niepusty z
charakterystyki domknięcia. Dla $A\cap U$ i $A\cap V$
element $A\cap(U\cap V)$ należy do $\mathcal B$
i jest ich przecięciem. Zatem $\mathcal B$ jest bazą
filtru $\mathcal F$. Ponieważ $A=A\cap X\in\mathcal B$,
mamy $A\in\mathcal F$. Każde otwarte otoczenie
$U$ zawiera $A\cap U\in\mathcal B$, więc należy
do $\mathcal F$; dowolne otoczenie $x$ zawiera
otwarte otoczenie, zatem również należy do filtru.

Odwrotnie, jeżeli $A\in\mathcal F$ i
$\mathcal F\to x$, to dla każdego otwartego
otoczenia $U$ punktu $x$ mamy $U\in\mathcal F$.
Stąd $A\cap U\in\mathcal F$, a ponieważ filtr
jest właściwy, $A\cap U\ne\varnothing$.
Oznacza to $x\in\overline A$.
\end{proof}

\begin{twierdzenie}[Sieci i filtry opisują tę samą zbieżność]
\label{thm:topog-sieci-filtry}
Niech $X$ będzie przestrzenią topologiczną.
\begin{enumerate}[label=(\roman*)]
\item Dla każdej sieci $(x_d)_{d\in D}$ w $X$
rodzina wszystkich zbiorów zawierających
pewien ogon $\{x_d:d\succeq d_0\}$
jest filtrem $\mathcal F_{(x_d)}$.
Sieć zbiega do $x$ wtedy i tylko wtedy,
gdy $\mathcal F_{(x_d)}\to x$.
\item Dla każdego filtru $\mathcal F$ na $X$
istnieje sieć $(x_d)_{d\in D}$ taka, że dla każdego
$x\in X$ zachodzi
$\mathcal F\to x\iff x_d\to x$.
\end{enumerate}
\end{twierdzenie}
\begin{proof}
(i) Ogony są niepuste, bo dla każdego $d_0$
zawierają $x_{d_0}$. Przecięcie dwóch ogonów
zawiera ogon o indeksie będącym wspólnym
ograniczeniem górnym ich indeksów początkowych.
Stanowią zatem bazę filtru. Otoczenie $U$ należy
do generowanego filtru dokładnie wtedy, gdy zawiera
pewien ogon sieci. Jest to definicja zbieżności sieci.

(ii) Niech
\[
 D=\{(F,a):F\in\mathcal F,\ a\in F\}.
\]
Rodzina $D$ jest niepusta, ponieważ $X\in\mathcal F$
i $X\ne\varnothing$ (właściwy filtr na zbiorze pustym
nie istnieje). Porządkujemy ją przez
$(F,a)\preceq(G,b)\iff G\subseteq F$.
Dla dwóch indeksów wybieramy dowolny punkt
$c\in F\cap G$, który istnieje, gdyż
$F\cap G\in\mathcal F$. Wtedy $(F\cap G,c)$
jest wspólnym ograniczeniem górnym.
Definiujemy sieć $x_{(F,a)}=a$.

Jeśli $\mathcal F\to x$, to dla każdego otwartego
otoczenia $U$ punktu $x$ mamy $U\in\mathcal F$.
Wybieramy $a\in U$ i indeks $(U,a)$. Każdy późniejszy
$(G,b)$ spełnia $G\subseteq U$, więc $x_{(G,b)}=b\in U$.
Zatem sieć zbiega do $x$.

Na odwrót, niech $(x_{(F,a)})$ zbiega do $x$.
Dla każdego otwartego otoczenia $U$ punktu $x$
istnieje $(F,a)$ takie, że
$x_{(G,b)}\in U$ dla wszystkich $(G,b)\succeq(F,a)$.
Dla dowolnego $b\in F$ indeks $(F,b)$ należy do $D$
i spełnia $(F,b)\succeq(F,a)$, zatem $b\in U$.
Wynika stąd $F\subseteq U$, a więc $U\in\mathcal F$
przez własność rozszerzania zbiorów. Każde inne
otoczenie punktu $x$ zawiera otwarte, więc również
należy do filtru. Otrzymujemy $\mathcal F\to x$.
\end{proof}

\begin{wniosek}[Ciągłość i jednoznaczność granic filtrów]
\label{cor:topog-filtr-ciaglosc-hausdorff}
Niech $X,Y$ będą przestrzeniami topologicznymi.
\begin{enumerate}[label=(\roman*)]
\item Odwzorowanie $f\colon X\to Y$ jest ciągłe
w $x\in X$ dokładnie wtedy, gdy dla każdego filtru
$\mathcal F\to x$ filtr generowany przez zbiory
$f(F)$, $F\in\mathcal F$, zbiega do $f(x)$.
\item Przestrzeń $X$ jest Hausdorffa dokładnie wtedy,
gdy każdy filtr na $X$ ma co najwyżej jedną granicę.
\end{enumerate}
\end{wniosek}
\begin{proof}
Dla (i) najpierw sprawdźmy, że rodzina
$\{f(F):F\in\mathcal F\}$ jest bazą filtru na $Y$.
Każdy jej element jest niepusty, ponieważ
$\varnothing\notin\mathcal F$. Jeśli $F,G\in\mathcal F$,
to $F\cap G\in\mathcal F$ oraz
$f(F\cap G)\subseteq f(F)\cap f(G)$.
Spełnione są więc warunki bazy filtru z
 definicji \ref{def:topog-filtr}.
Ponadto $V\subseteq Y$ należy do filtru obrazowego
dokładnie wtedy, gdy
$f^{-1}(V)\in\mathcal F$: jeśli $f(F)\subseteq V$
dla pewnego $F\in\mathcal F$, to
$F\subseteq f^{-1}(V)$ i filtr zawiera przeciwobraz;
odwrotnie z $f^{-1}(V)\in\mathcal F$ wynika
$f(f^{-1}(V))\subseteq V$.
Jeżeli $f$ jest ciągłe w $x$ i $\mathcal F\to x$,
to dla każdego otwartego otoczenia $V$ punktu $f(x)$
zbiór $f^{-1}(V)$ jest otoczeniem $x$, zatem należy
do $\mathcal F$. Z wykazanej równoważności
$V$ należy do filtru obrazowego, co dowodzi jego
zbieżności do $f(x)$.

Odwrotnie, przypuśćmy, że $f$ zachowuje granice
wszystkich filtrów, ale nie jest ciągłe w $x$.
Z twierdzenia \ref{thm:topog-siec-kryterium-ciaglosci}
istnieje sieć $(x_d)$ zbieżna do $x$, której obraz
$(f(x_d))$ nie zbiega do $f(x)$. Filtr ogonów
$\mathcal F_{(x_d)}$ zbiega do $x$ na mocy
twierdzenia \ref{thm:topog-sieci-filtry}(i).
Ponieważ ogony sieci obrazowej są dokładnie
obrazami ogonów sieci wyjściowej, ich filtr
jest filtrem obrazowym $f_*(\mathcal F_{(x_d)})$.
Nie zbiega on do $f(x)$, co przeczy założeniu.

Dla (ii) jeśli $X$ jest Hausdorffa, dwa różne punkty
$x,y$ mają rozłączne otwarte otoczenia $U,V$.
Filtr zbieżny do obu musiałby zawierać
$U\cap V=\varnothing$, co jest niemożliwe.
Jeśli $X$ nie jest Hausdorffa, to z dowodu
twierdzenia \ref{thm:topog-sieci-hausdorff}
istnieje sieć zbieżna do dwóch różnych punktów.
Jej filtr ogonów z punktu (i) twierdzenia
\ref{thm:topog-sieci-filtry} ma te same granice.
\end{proof}

\section{Podsumowanie: trzy języki zbieżności}
\label{sec:topog-zbieznosc-zestawienie}

\begin{center}
\begin{tabular}{@{}p{.22\textwidth}p{.31\textwidth}p{.34\textwidth}@{}}
\hline
\textbf{Narzędzie} & \textbf{Co wykrywa zawsze} &
\textbf{Kiedy wystarcza}\ \\ \hline
Ciągi & Granice i konieczny warunek ciągłości &
Dla domknięcia i ciągłości, jeśli przestrzeń
spełnia pierwszy aksjomat przeliczalności.\\
Sieci & Domknięcie, ciągłość, aksjomat Hausdorffa &
W każdej przestrzeni topologicznej.\\
Filtry & Domknięcie, ciągłość, aksjomat Hausdorffa &
W każdej przestrzeni topologicznej.\\
\hline
\end{tabular}
\end{center}

\begin{uwaga}[Powiązania z dalszymi rozdziałami]
\label{rem:topog-zbieznosc-dalsze}
Sieci i filtry pozwolą później sformułować
charakterystyki zwartości oraz zbadać dokładniej
przestrzenie bez przeliczalnej bazy otoczeń.
Warunki przeliczalności i aksjomaty metryzowalności
będą analizowane odrębnie; w szczególności
spełnienia pierwszego aksjomatu przeliczalności nie
utożsamiamy tutaj z metryzowalnością.
\end{uwaga}

\chapter{Zwartość i jej charakterystyki}
\label{ch:topog-zwartosc}

\noindent\textbf{Cel rozdziału.}
Zwartość jest własnością pozwalającą zastąpić dowolne pokrycie
otwarte pokryciem skończonym. Jej siła ujawnia się w badaniu
odwzorowań ciągłych (w szczególności osiągania kresów
przez funkcje rzeczywiste), rozdzielania zbiorów
i dowolnych iloczynów.
Po definicji i podstawowych przykładach wyprowadzimy równoważne
kryteria przez rodziny zbiorów domkniętych, sieci i ultrafiltry.
Udowodnimy twierdzenie Tychonoffa o zwartości iloczynów oraz
uzgodnione wcześniej twierdzenie o zanurzaniu przestrzeni
Tychonoffa w kostkę $[0,1]^I$. Tam, gdzie wykorzystamy
aksjomat wyboru, zaznaczymy to jawnie.

\section{Definicja zwartości i pierwsze przykłady}
\label{sec:topog-zwartosc-definicja}

\begin{definicja}[Pokrycie otwarte i zwartość]
\label{def:topog-zwartosc}
Niech $(X,\tau)$ będzie przestrzenią topologiczną.
\emph{Pokryciem otwartym} $X$ nazywamy rodzinę
$(U_i)_{i\in I}$ zbiorów z $\tau$ taką, że
$X=\bigcup_{i\in I}U_i$.
\emph{Podpokryciem skończonym} tej rodziny jest
podrodzina $(U_i)_{i\in J}$, gdzie $J\subseteq I$ jest
skończony i $X=\bigcup_{i\in J}U_i$.
Przestrzeń $X$ jest \emph{zwarta}, jeżeli każde jej
pokrycie otwarte ma podpokrycie skończone.
Podzbiór $K\subseteq X$ jest \emph{zwarty w $X$}, gdy
przestrzeń $K$ z topologią podprzestrzeni jest zwarta.
Nie zakładamy żadnego aksjomatu oddzielania.
\end{definicja}

\begin{stwierdzenie}[Zwartość podzbioru w języku przestrzeni otaczającej]
\label{prop:topog-zwartosc-wzgledna}
Niech $K\subseteq X$. Wówczas $K$ jest zwarty dokładnie
wtedy, gdy dla każdej rodziny $(U_i)_{i\in I}$ zbiorów
otwartych w $X$ takiej, że $K\subseteq\bigcup_{i\in I}U_i$,
istnieje skończony $J\subseteq I$ spełniający
$K\subseteq\bigcup_{i\in J}U_i$.
\end{stwierdzenie}
\begin{proof}
Jeśli $K$ jest zwarty, rodzina $(K\cap U_i)$ pokrywa
$K$ zbiorami otwartymi w podprzestrzeni, więc ma
podpokrycie skończone. Odwrotnie, każde pokrycie $K$
zbiorami otwartymi w $K$ ma postać $(K\cap U_i)$
z pewnymi $U_i$ otwartymi w $X$, zgodnie z definicją
topologii podprzestrzeni (rozdział~04). Zastosowanie
warunku dla $U_i$ daje skończone podpokrycie w $K$.
\end{proof}

\begin{przyklad}[Zbiory skończone i przestrzeń pusta]
\label{ex:topog-zwartosc-skonczone}
Każda przestrzeń topologiczna o skończonym zbiorze
punktów jest zwarta. Z dowolnego pokrycia wybieramy
po jednym zbiorze zawierającym każdy punkt; otrzymana
rodzina ma co najwyżej tyle elementów, ile punktów
przestrzeni. Dla $X=\varnothing$ pokrycie puste
ma puste podpokrycie skończone, które pokrywa $X$.
Także każdy skończony podzbiór dowolnej przestrzeni
jest zwarty, niezależnie od jej topologii.
\end{przyklad}

\begin{przyklad}[Niezwartość prostej i przedziału otwartego]
\label{ex:topog-zwartosc-przeciwwklady}
Prosta $\R$ z topologią euklidesową nie jest zwarta:
rodzina $\{(-n,n):n\ge1\}$ pokrywa $\R$, lecz żadna
jej podrodzina skończona tego nie robi, ponieważ
zawiera się w $(-N,N)$ dla największego wybranego $N$.
Podobnie pokrycie $\{(1/n,1):n\ge2\}$ przedziału
$(0,1)$ nie ma podpokrycia skończonego: skończona
rodzina nie pokrywa punktów dostatecznie bliskich $0$.
\end{przyklad}

\begin{figure}[htbp]
\centering
\begin{tikzpicture}[x=9cm,y=.72cm,>=Stealth]
  \draw[->,thick] (-.04,0)--(1.11,0);
  \draw[gray,dashed] (0,-.2)--(0,4.1);
  \draw[gray,dashed] (1,-.2)--(1,4.1);
  \node[below] at (0,-.05) {$0$};
  \node[below] at (1,-.05) {$1$};
  \foreach \a/\y/\n in {.5/1/2,.3333/1.75/3,.25/2.5/4,.1667/3.25/6} {
    \draw[blue!75!black,line width=1.1pt] (\a,\y)--(1,\y);
    \draw[blue!75!black,fill=white,line width=.8pt] (\a,\y) circle (.035cm);
    \draw[blue!75!black,fill=white,line width=.8pt] (1,\y) circle (.035cm);
    \node[anchor=east] at (.015,\y) {$U_{\n}$};
  }
  \node[align=center,blue!55!black,font=\small] at (.61,4.05)
    {Coraz większe zbiory pokrywają $(0,1)$,\\ ale żaden skończony wybór nie wystarcza};
\end{tikzpicture}
\caption{Pierwsze zbiory z pokrycia $U_n=(1/n,1)$, $n\ge2$.
W każdej skończonej podrodzinie pozostaje niepokryty
przedział przy punkcie $0$. Puste końce oznaczają
brak krańców przedziałów w $U_n$.}
\label{fig:topog-zwartosc-pokrycie-otwarte}
\end{figure}

\begin{twierdzenie}[Zwartość przedziału domkniętego]
\label{thm:topog-zwartosc-odcinek}
Dla każdych $a,b\in\R$ z $a\le b$ przedział $[a,b]$
jest zwarty w topologii euklidesowej.
\end{twierdzenie}
\begin{proof}
Przy $a=b$ stosujemy przykład
\ref{ex:topog-zwartosc-skonczone}. Niech $a<b$
i niech $(U_i)_{i\in I}$ będzie pokryciem
otwartym przedziału $[a,b]$ zbiorami otwartymi
w $\R$ (wystarczy dzięki stwierdzeniu
\ref{prop:topog-zwartosc-wzgledna}).
Niech $S$ oznacza zbiór $t\in[a,b]$, dla których
$[a,t]$ pokrywa skończona liczba zbiorów $U_i$.
Zbiór $S$ jest niepusty, ponieważ istnieje $i$
z $a\in U_i$, a wtedy $a\in S$.
Jest też ograniczony z góry przez $b$; niech $c=\sup S$.

Wybierzmy $i_0$ z $c\in U_{i_0}$. Ponieważ $U_{i_0}$
jest otwarty, istnieje $\delta>0$ takie, że
$(c-\delta,c+\delta)\subseteq U_{i_0}$.
Z definicji supremum istnieje $t\in S$ z
$t>c-\delta/2$ (gdy $c=a$, możemy wziąć $t=a$).
Wtedy $[a,t]$ ma pokrycie skończone, a dodanie
$U_{i_0}$ pokrywa także odcinek
$[t,\min\{b,c+\delta/2\}]$.
Gdyby $c<b$, punkt $\min\{b,c+\delta/2\}$
byłby większy od $c$ i należał do $S$, wbrew
definicji supremum. Stąd $c=b$. Ponownie,
ponieważ $t> b-\delta/2$ i $t\in S$, dołączenie
$U_{i_0}$ do skończonego pokrycia $[a,t]$
pokrywa całe $[a,b]$. Dowód jest zakończony.
\end{proof}

\section{Zbiory domknięte, obrazy ciągłe i oddzielanie}
\label{sec:topog-zwartosc-podstawowe-wlasnosci}

\begin{twierdzenie}[Kryterium skończonych przecięć]
\label{thm:topog-zwartosc-fip}
Przestrzeń topologiczna $X$ jest zwarta wtedy i tylko
wtedy, gdy każda \emph{niepusta} rodzina $(F_i)_{i\in I}$
zbiorów domkniętych w $X$, dla której każde
przecięcie skończenie wielu jej elementów jest
niepuste, spełnia $\bigcap_{i\in I}F_i\ne\varnothing$.
Warunek dotyczący skończonych przecięć nazywamy
\emph{własnością skończonych przecięć}.
\end{twierdzenie}
\begin{proof}
Niech $X$ będzie zwarta i załóżmy, że
$\bigcap_i F_i=\varnothing$. Dopełnienia $X\setminus F_i$
tworzą pokrycie otwarte $X$. Istnieje skończona
podrodzina pokrywająca $X$, a z praw de Morgana
przecięcie odpowiednich $F_i$ jest puste,
sprzecznie z założeniem.

Odwrotnie, załóżmy kryterium i niech $(U_i)_{i\in I}$
będzie pokryciem otwartym bez podpokrycia skończonego.
Jeśli $I$ byłby pusty, to $X=\varnothing$, a pusta
podrodzina stanowiłaby podpokrycie skończone.
Zatem $I\ne\varnothing$. Dla $F_i=X\setminus U_i$
każde skończone przecięcie jest niepuste,
gdyż żadne skończenie wiele $U_i$ nie pokrywa $X$.
Z założenia $\bigcap_i F_i\ne\varnothing$, co przeczy
temu, że $(U_i)$ pokrywa $X$.
\end{proof}

\begin{twierdzenie}[Trwałość zwartości]
\label{thm:topog-zwartosc-trwalosc}
Niech $X,Y$ będą przestrzeniami topologicznymi.
\begin{enumerate}[label=(\roman*)]
\item Jeśli $X$ jest zwarta i $F\subseteq X$ jest
domknięty, to $F$ jest zwarty.
\item Jeśli $K\subseteq X$ jest zwarty i
$f\colon X\to Y$ jest ciągłe, to $f(K)$ jest zwarty.
W szczególności ciągły obraz przestrzeni zwartej
jest zwarty.
\end{enumerate}
\end{twierdzenie}
\begin{proof}
Dla (i) niech $(V_i)$ będzie otwartym pokryciem
$F$ w topologii względnej; zapiszmy
$V_i=F\cap U_i$, gdzie $U_i$ są otwarte w $X$.
Rodzina $(U_i)$ uzupełniona o otwarty zbiór
$X\setminus F$ pokrywa $X$. Skończone podpokrycie
$X$ po ograniczeniu do $F$ daje skończone
podpokrycie $(V_i)$ zbioru $F$.

Dla (ii) niech $(W_i)$ będzie otwartym pokryciem
$f(K)$ w $Y$. Ciągłość odwzorowania $f|_K$
sprawia, że $(K\cap f^{-1}(W_i))$ jest otwartym
pokryciem $K$. Wybór skończonego podpokrycia
przenosi się przez $f$ na skończone pokrycie $f(K)$.
\end{proof}

\begin{twierdzenie}[Zwarte podzbiory przestrzeni Hausdorffa]
\label{thm:topog-zwartosc-hausdorff-domkniete}
Jeśli $X$ jest przestrzenią Hausdorffa, to każdy
zwarty podzbiór $K\subseteq X$ jest domknięty w $X$.
\end{twierdzenie}
\begin{proof}
Dla $K=\varnothing$ teza jest oczywista.
Niech $x\in X\setminus K$. Rozważmy wszystkie pary
$(U,V)$ rozłącznych zbiorów otwartych w $X$,
dla których $x\in U$. Zbiory $V$ występujące
w tych parach pokrywają $K$: dla każdego $y\in K$
aksjomat Hausdorffa zapewnia taką parę z $y\in V$.
Ze zwartości $K$ wybieramy skończenie wiele par
$(U_1,V_1),\ldots,(U_m,V_m)$, dla których
$K\subseteq V_1\cup\cdots\cup V_m$.
Zbiór $U=U_1\cap\cdots\cap U_m$ jest
otwartym otoczeniem $x$ i nie przecina $K$:
każdy punkt $K$ należy do pewnego $V_j$,
rozłącznego z $U\subseteq U_j$.
Oznacza to, że $X\setminus K$ jest otwarty,
a więc $K$ domknięty. Konstrukcja nie wymaga
jednoczesnego wyboru otoczeń dla wszystkich $y\in K$.
\end{proof}

\begin{przyklad}[Dlaczego założenie Hausdorffa jest potrzebne]
\label{ex:topog-zwartosc-niedomkniety}
W przestrzeni Sierpińskiego
$X=\{0,1\}$ z topologią
$\{\varnothing,\{1\},X\}$ podzbiór $\{1\}$
jest zwarty, ponieważ jest skończony,
lecz nie jest domknięty: jego dopełnienie $\{0\}$
nie jest otwarte. Nie można więc bez założenia
Hausdorffa wnioskować o domkniętości zbioru zwartego.
\end{przyklad}

\begin{twierdzenie}[Twierdzenie Weierstrassa o osiąganiu kresów]
\label{thm:topog-weierstrass-kresy}
Niech $K$ będzie \emph{niepustą} zwartą przestrzenią
topologiczną, a $f\colon K\to\R$ funkcją ciągłą.
Wówczas $f$ jest ograniczona oraz istnieją punkty
$x_{\min},x_{\max}\in K$ takie, że
\[
 f(x_{\min})=\inf_{x\in K} f(x),\qquad
 f(x_{\max})=\sup_{x\in K} f(x).
\]
W szczególności funkcja $f$ osiąga minimum i maksimum.
Nie zakładamy, że $K$ jest przestrzenią Hausdorffa.
\end{twierdzenie}
\begin{proof}
Na mocy twierdzenia \ref{thm:topog-zwartosc-trwalosc}(ii)
zbiór $A=f(K)$ jest zwarty w $\R$; jest niepusty,
ponieważ $K\ne\varnothing$. Rodzina przedziałów
$((-n,n))_{n\ge1}$ jest otwartym pokryciem $\R$,
a więc również $A$. Ze zwartości $A$ istnieje
$n\ge1$ takie, że $A\subseteq(-n,n)$: jeśli wybrano
skończenie wiele przedziałów, wystarczy wziąć
największy z ich indeksów. Funkcja $f$ jest zatem
ograniczona.

Z zupełności porządku liczb rzeczywistych istnieją
skończone liczby $m=\inf A$ oraz $M=\sup A$.
Ponieważ $\R$ jest przestrzenią Hausdorffa, zbiór
$A$ jest domknięty na mocy twierdzenia
\ref{thm:topog-zwartosc-hausdorff-domkniete}.
Wykażemy, że $M\in A$. Dla każdego $\varepsilon>0$
z definicji supremum istnieje $a\in A$ spełniające
$M-\varepsilon<a\le M$. Stąd każde otwarte otoczenie
$M$ przecina $A$, czyli $M\in\overline A=A$.
Analogicznie dla każdego $\varepsilon>0$ istnieje
$a\in A$ takie, że $m\le a<m+\varepsilon$,
więc $m\in\overline A=A$. Z definicji obrazu
istnieją $x_{\min},x_{\max}\in K$ z
$f(x_{\min})=m$ i $f(x_{\max})=M$.
\end{proof}

\begin{przyklad}[Rola zwartości i ciągłości w twierdzeniu Weierstrassa]
\label{ex:topog-weierstrass-zalozenia}
Funkcja $f\colon(0,1)\to\R$, $f(x)=x$, jest ciągła
i ograniczona, lecz nie osiąga ani swojego infimum $0$,
ani supremum $1$. Jej dziedzina nie jest zwarta.

Z kolei funkcja $g\colon[0,1]\to\R$ określona przez
$g(x)=x$ dla $x<1$ oraz $g(1)=0$ nie jest ciągła
w punkcie $1$. Mimo że $[0,1]$ jest zwarty,
$\sup g([0,1])=1$ nie jest wartością funkcji $g$.
Oba przykłady pokazują, że założeń twierdzenia
nie można po prostu pominąć.
\end{przyklad}

\begin{twierdzenie}[Zwartość i normalność]
\label{thm:topog-zwartosc-hausdorff-normalna}
Każda zwarta przestrzeń Hausdorffa jest $T_4$,
a więc spełnia aksjomat Tychonoffa.
\end{twierdzenie}
\begin{proof}
Przestrzeń Hausdorffa jest $T_1$ (rozdział~05).
Wystarczy dowieść normalności. Niech $F,G\subseteq X$
będą rozłącznymi zbiorami domkniętymi.
Na mocy twierdzenia \ref{thm:topog-zwartosc-trwalosc}
oba są zwarte. Jeśli któryś z nich jest pusty,
możemy użyć $\varnothing$ i $X$ jako rozłącznych
otwartych otoczeń $F$ i $G$. W dalszej części
zakładamy więc, że $F\ne\varnothing$ i $G\ne\varnothing$.

Ustalmy $x\in F$ i rozważmy wszystkie pary
$(U,V)$ rozłącznych zbiorów otwartych takich,
że $x\in U$. Z aksjomatu Hausdorffa ich drugie
składowe $V$ tworzą pokrycie otwarte $G$.
Ze zwartości $G$ wybieramy skończenie wiele par
$(U_1,V_1),\ldots,(U_m,V_m)$, dla których
$G\subseteq V_1\cup\cdots\cup V_m$.
Zbiory $U_x=\bigcap_{j=1}^m U_j$ oraz
$V_x=\bigcup_{j=1}^m V_j$ są otwarte i rozłączne,
$x\in U_x$ oraz $G\subseteq V_x$.

Zatem dla każdego $x\in F$ istnieje para
rozłącznych zbiorów otwartych $(U,V)$ spełniająca
$x\in U$ i $G\subseteq V$. Rozważmy rodzinę
\emph{wszystkich} takich par: ich pierwsze składowe
pokrywają $F$. Zwartość $F$ daje skończenie wiele
par $(U'_1,V'_1),\ldots,(U'_n,V'_n)$ takich,
że $F\subseteq\bigcup_{k=1}^n U'_k$.
Wtedy $U=\bigcup_{k=1}^n U'_k$ i
$V=\bigcap_{k=1}^n V'_k$ są rozłącznymi
otwartymi zbiorami zawierającymi odpowiednio $F,G$.
Otrzymaliśmy normalność; połączenie z $T_1$ daje $T_4$.
Dzięki rozważaniu wszystkich dopuszczalnych par
nie wybieraliśmy jednocześnie otoczeń z nieskończonych
rodzin.
Implikację $T_4\Rightarrow T_{3\frac12}$
udowodniono w rozdziale~05.
\end{proof}

\begin{wniosek}[Ciągła bijekcja z przestrzeni zwartej]
\label{cor:topog-zwartosc-bijekcja-homeomorfizm}
Niech $X$ będzie przestrzenią zwartą, $Y$ przestrzenią
Hausdorffa, a $f\colon X\to Y$ ciągłą bijekcją.
Wówczas $f$ jest homeomorfizmem.
\end{wniosek}
\begin{proof}
Każdy domknięty $F\subseteq X$ jest zwarty na mocy
twierdzenia \ref{thm:topog-zwartosc-trwalosc}.
Jego obraz $f(F)$ jest zwarty jako obraz ciągły,
a więc domknięty w $Y$ na mocy twierdzenia
\ref{thm:topog-zwartosc-hausdorff-domkniete}.
Zatem $f$ jest odwzorowaniem domkniętym.
Bijektywność sprawia, że dla dowolnego domkniętego
$F\subseteq X$ zachodzi $(f^{-1})^{-1}(F)=f(F)$;
stąd odwrotność jest ciągła, więc $f$ jest
homeomorfizmem.
\end{proof}

\section{Zwartość przez sieci i filtry}
\label{sec:topog-zwartosc-zbieznosc}

\begin{definicja}[Punkt skupienia sieci i filtru]
\label{def:topog-zwartosc-punkt-skupienia}
Niech $(x_d)_{d\in D}$ będzie siecią w $X$.
Punkt $x\in X$ nazywamy \emph{punktem skupienia sieci},
jeżeli dla każdego otwartego otoczenia $U$ punktu $x$
i każdego $d_0\in D$ istnieje $d\succeq d_0$
z $x_d\in U$.
Jeśli $\mathcal F$ jest filtrem na $X$, punkt $x$
jest \emph{punktem skupienia filtru}, gdy
$F\cap U\ne\varnothing$ dla każdego
$F\in\mathcal F$ i każdego otwartego otoczenia $U$
punktu $x$. Jest to równoważne
$x\in\bigcap_{F\in\mathcal F}\overline F$,
ponieważ $x\in\overline F$ dokładnie wtedy,
gdy każde otoczenie $x$ przecina $F$.
\end{definicja}

\begin{stwierdzenie}[Punkt skupienia i podsieć zbieżna]
\label{prop:topog-zwartosc-podsiec}
Punkt $x$ jest punktem skupienia sieci $(x_d)_{d\in D}$
wtedy i tylko wtedy, gdy istnieje podsieć $(x_{\varphi(e)})_{e\in E}$
zbieżna do $x$. Przyjmujemy definicję podsieci
z rozdziału~06: odwzorowanie indeksów
$\varphi\colon E\to D$ spełnia warunek ostatecznej
współkońcowości (nie wymagamy jego monotoniczności).
\end{stwierdzenie}
\begin{proof}
Jeżeli podsieć jest zbieżna do $x$, to dla danego
otwartego otoczenia $U$ punktu $x$ od pewnego
$e_0$ mamy $x_{\varphi(e)}\in U$ dla wszystkich
$e\succeq e_0$. Dla dowolnego $d_0\in D$
na mocy kofinalności ogonowej istnieje $e_1$,
po którym $\varphi(e)\succeq d_0$.
Wspólne ograniczenie górne $e_0,e_1$ daje
$e$ z $\varphi(e)\succeq d_0$ i
$x_{\varphi(e)}\in U$, więc $x$ jest punktem skupienia.

Odwrotnie, niech $x$ będzie punktem skupienia.
Oznaczmy przez $\mathcal N_x^o$ rodzinę otwartych
otoczeń $x$ i zbudujmy zbiór
\[
E=\{(d,U,t):d,t\in D,\ d\preceq t,\
U\in\mathcal N_x^o,\ x_t\in U\}.
\]
Z warunku skupienia wynika, że $E\ne\varnothing$.
Definiujemy
$(d,U,t)\preceq_E(d',U',t')$ wtedy i tylko wtedy,
gdy $d\preceq d'$, $U'\subseteq U$ oraz $t\preceq t'$.
Relacja jest zwrotna i przechodnia. Dla dwóch
indeksów wybierzmy $d'$ większy od ich pierwszych
i trzecich współrzędnych; ze skupienia wybierzmy
$t'\succeq d'$ z $x_{t'}\in U_1\cap U_2$.
Indeks $(d',U_1\cap U_2,t')$ jest wtedy wspólnym
ograniczeniem górnym, więc $E$ jest skierowany.

Niech $\varphi(d,U,t)=t$. Monotoniczność wynika
z definicji $\preceq_E$. Dla dowolnego $d_0\in D$
wybierzmy $t_0\succeq d_0$ z
$x_{t_0}\in X$ (możemy wziąć $t_0=d_0$);
indeks $(d_0,X,t_0)\in E$ ma tę własność,
że powyżej niego $\varphi(d,U,t)=t\succeq t_0\succeq d_0$.
Zatem $\varphi$ jest kofinalne w sensie ogonów.
Wreszcie dla otwartego otoczenia $V$ punktu $x$
wybierzmy indeks $(d,V,t)\in E$ ze skupienia.
Wszystkie późniejsze indeksy $(d',U',t')$
spełniają $U'\subseteq V$, zatem
$x_{\varphi(d',U',t')}=x_{t'}\in U'\subseteq V$.
Podsieć zbiega do $x$.
\end{proof}

\begin{twierdzenie}[Kryterium zwartości przez sieci]
\label{thm:topog-zwartosc-sieci}
Przestrzeń topologiczna $X$ jest zwarta wtedy i tylko
wtedy, gdy każda sieć w $X$ ma punkt skupienia.
Równoważnie: każda sieć ma podsieć zbieżną.
\end{twierdzenie}
\begin{proof}
Niech $X$ będzie zwarta i $(x_d)_{d\in D}$
niech będzie siecią w $X$. Przypuśćmy, że nie ma
punktu skupienia. Wówczas rodzina
\[
\mathcal V=\{U\subseteq X:\ U\text{ otwarty oraz }\
\exists d_U\in D\ \forall d\succeq d_U:\ x_d\notin U\}
\]
pokrywa $X$: dla każdego punktu z negacji warunku
skupienia otrzymujemy pewien zbiór należący do
$\mathcal V$, który ten punkt zawiera.
Zwartość daje skończenie wiele
$U_1,\ldots,U_n\in\mathcal V$ pokrywających $X$.
Dla każdego $j$ istnieje $d_j$, po którym sieć
omija $U_j$. Z kierowalności $D$ wybieramy
$d_0\succeq d_1,\ldots,d_n$.
Wtedy $x_{d_0}$ nie należy do żadnego $U_j$,
wbrew temu, że zbiory te pokrywają $X$.
Dla $X=\varnothing$ nie ma żadnych sieci;
implikacja jest wtedy spełniona automatycznie.

Odwrotnie, niech $(U_i)_{i\in I}$ będzie pokryciem
otwartym bez podpokrycia skończonego.
Zamiast wybierać jednocześnie punkty z wielu
niepustych zbiorów, zbudujemy sieć bez użycia
aksjomatu wyboru. Niech
\[
D=\{(J,x):J\subseteq I\text{ skończony},\
 x\in X\setminus\bigcup_{i\in J}U_i\}.
\]
Zbiór $D$ jest niepusty, ponieważ para
$(\varnothing,x)$ istnieje dla pewnego $x\in X$;
gdyby $X=\varnothing$, pokrycie puste byłoby
podpokryciem skończonym.
Ustalmy $(J,x)\preceq(K,y)$ dokładnie wtedy,
gdy $J\subseteq K$. Jest to relacja zwrotna
i przechodnia. Dwie pary mają wspólne ograniczenie
górne: skończona suma $J\cup K$ nie pokrywa
$X$, więc istnieje $z\in X\setminus
\bigcup_{i\in J\cup K}U_i$ i można wziąć
$(J\cup K,z)$. Tak więc $D$ jest skierowany.
Zdefiniujmy sieć $x_{(J,x)}=x$.
Dla dowolnego $z\in X$ wybierzmy $i\in I$
z $z\in U_i$. Ponieważ żadne skończone
zbiory $U_i$ nie pokrywają $X$, istnieje
$(\{i\},x)\in D$. Powyżej tego indeksu
wszystkie wyrazy sieci leżą poza $U_i$;
zatem $z$ nie jest punktem skupienia.
Otrzymujemy sprzeczność i wnioskujemy zwartość.
Równoważność z podsieciami wynika ze stwierdzenia
\ref{prop:topog-zwartosc-podsiec}.
\end{proof}

\begin{definicja}[Ultrafiltr]
\label{def:topog-zwartosc-ultrafiltr}
\emph{Ultrafiltrem} na $X$ nazywamy filtr właściwy
$\mathcal U$, maksymalny względem zawierania wśród
wszystkich filtrów właściwych na $X$.
Filtr właściwy rozumiemy zgodnie z definicją filtru
z rozdziału~06 (w szczególności nie zawiera zbioru pustego). Na zbiorze pustym filtr właściwy
nie istnieje.
\end{definicja}

\begin{lemat}[Kryterium ultrafiltru]
\label{lem:topog-zwartosc-ultrafiltr-alternatywa}
Filtr właściwy $\mathcal U$ na zbiorze $X$ jest
ultrafiltrem wtedy i tylko wtedy, gdy dla każdego
$B\subseteq X$ dokładnie jeden ze zbiorów
$B$ i $X\setminus B$ należy do $\mathcal U$.
\end{lemat}
\begin{proof}
Oba zbiory nie mogą należeć do właściwego filtru,
bo ich przecięcie jest puste. Jeśli $\mathcal U$
jest maksymalny i $B\notin\mathcal U$, to
musi istnieć $F\in\mathcal U$ z $F\cap B=\varnothing$.
W przeciwnym razie rodzina $\mathcal U\cup\{B\}$
miałaby własność skończonych przecięć i generowałaby
filtr właściwy ściśle większy od $\mathcal U$,
co przeczyłoby maksymalności.
Wtedy $F\subseteq X\setminus B$, zatem
$X\setminus B\in\mathcal U$ z domknięcia filtru
na nadzbiory.

Odwrotnie, jeśli filtr wybiera dla każdego $B$
jedną ze stron podziału $X=B\cup(X\setminus B)$,
to żaden większy właściwy filtr nie może zawierać
$B\notin\mathcal U$, bo zawierałby też
$X\setminus B\in\mathcal U$ i stąd zbiór pusty.
To dowodzi maksymalności.
\end{proof}

\begin{lemat}[Przedłużenie filtru do ultrafiltru]
\label{lem:topog-zwartosc-ultrafiltr-zorn}
Każdy filtr właściwy na zbiorze $X$ jest zawarty
w pewnym ultrafiltrze na $X$. W dowodzie korzystamy
z lematu Zorna, a zatem z zasady wyboru.
\end{lemat}
\begin{proof}
Rozważmy uporządkowaną przez zawieranie rodzinę
wszystkich filtrów właściwych zawierających ustalony
filtr $\mathcal F$. Jest niepusta, gdyż zawiera
$\mathcal F$. Jeśli $\mathcal C$ jest łańcuchem,
to suma rodzin $\mathcal G=\bigcup_{\mathcal H\in\mathcal C}
\mathcal H$ jest filtrem właściwym: dla dwóch zbiorów
należących do $\mathcal G$ ich filtry z łańcucha
są porównywalne, zatem większy z nich zawiera oba
i ich przecięcie; domknięcie na nadzbiory
jest równie bezpośrednie, a $\varnothing$ nie należy
do żadnego z filtrów. Dla pustego łańcucha
górnym ograniczeniem jest $\mathcal F$.
Lemat Zorna daje element maksymalny, czyli ultrafiltr.
\end{proof}

\begin{twierdzenie}[Ultrafiltry charakteryzują zwartość]
\label{thm:topog-zwartosc-ultrafiltry}
Przestrzeń topologiczna $X$ jest zwarta wtedy i tylko
wtedy, gdy każdy ultrafiltr na $X$ jest zbieżny
do przynajmniej jednego punktu $X$.
\end{twierdzenie}
\begin{proof}
Niech $X$ będzie zwarta i $\mathcal U$
ultrafiltrem na $X$. Rodzina
$(\overline F)_{F\in\mathcal U}$ jest niepusta
i ma własność skończonych przecięć:
$F_1\cap\cdots\cap F_n\in\mathcal U$ jest
niepusty i zawiera się w każdym $\overline{F_j}$.
Na mocy twierdzenia \ref{thm:topog-zwartosc-fip}
istnieje $x\in\bigcap_{F\in\mathcal U}\overline F$.
Jeżeli otwarte otoczenie $V$ punktu $x$ nie należałoby
do $\mathcal U$, to z lematu
\ref{lem:topog-zwartosc-ultrafiltr-alternatywa}
$X\setminus V\in\mathcal U$.
Mielibyśmy wtedy $x\in\overline{X\setminus V}$,
co jest niemożliwe, ponieważ otwarty zbiór $V$
nie przecina $X\setminus V$.
Zatem każde otwarte otoczenie $x$ należy do
$\mathcal U$, czyli $\mathcal U\to x$.

Odwrotnie, załóżmy, że $X$ nie jest zwarta.
Istnieje pokrycie otwarte $(U_i)_{i\in I}$
bez podpokrycia skończonego. Dopełnienia
$F_i=X\setminus U_i$ mają własność skończonych
przecięć, więc generują filtr właściwy:
jego elementami są nadzbiory niepustych skończonych
przecięć $F_i$. Lemat
\ref{lem:topog-zwartosc-ultrafiltr-zorn}
daje ultrafiltr $\mathcal U$ zawierający każdy $F_i$.
Gdyby $\mathcal U\to x$, wybralibyśmy $i$ z
$x\in U_i$; wtedy również $U_i\in\mathcal U$,
a $F_i=X\setminus U_i\in\mathcal U$, sprzeczność.
Zatem pewien ultrafiltr nie zbiega, co kończy dowód.
\end{proof}

\begin{uwaga}[Punkt skupienia filtru i granica ultrafiltru]
\label{rem:topog-zwartosc-filtr-zbieznosc}
Każdy filtr na przestrzeni zwartej ma punkt skupienia:
rodzina $(\overline F)_{F\in\mathcal F}$ ma własność
skończonych przecięć, bo przecięcie skończonej liczby
elementów filtru jest niepuste. Po zastosowaniu
kryterium \ref{thm:topog-zwartosc-fip} uzyskujemy
$x\in\bigcap_{F\in\mathcal F}\overline F$.
Nie musi to być granica całego filtru; dla ultrafiltru
punkt skupienia jest jednak granicą na mocy dowodu
poprzedniego twierdzenia. W przestrzeni Hausdorffa
każdy filtr ma co najwyżej jedną granicę, jak
udowodniono w rozdziale~06.
\end{uwaga}

\section{Zwartość iloczynów — twierdzenie Tychonoffa}
\label{sec:topog-zwartosc-iloczyny}

\begin{twierdzenie}[Twierdzenie Tychonoffa o iloczynach]
\label{thm:topog-tychonoff-iloczyny}
Niech $I$ będzie dowolnym zbiorem, a
$((X_i,\tau_i))_{i\in I}$ rodziną przestrzeni
zwartych. Wówczas iloczyn
$X=\prod_{i\in I}X_i$ z topologią produktową
jest zwarty. W dowodzie dla dowolnej rodziny
wykorzystujemy aksjomat wyboru.
\end{twierdzenie}
\begin{proof}
Jeśli $I=\varnothing$, iloczyn jest przestrzenią
jednopunktową i jest zwarty. Jeśli któryś czynnik
jest pusty, iloczyn też jest pusty i jest zwarty.
Rozważmy więc rodzinę niepustych przestrzeni.
Aksjomat wyboru zapewnia, że iloczyn $X$ jest
niepusty.

Niech $\mathcal U$ będzie dowolnym ultrafiltrem
na $X$ i $p_i\colon X\to X_i$ rzutowaniem.
Definiujemy filtr obrazowy
\[
\mathcal U_i=\{B\subseteq X_i:
 p_i^{-1}(B)\in\mathcal U\}.
\]
Jest on ultrafiltrem: dla każdego $B\subseteq X_i$
przeciwobrazy $B$ oraz $X_i\setminus B$ są
wzajemnymi dopełnieniami w $X$, więc na mocy lematu
\ref{lem:topog-zwartosc-ultrafiltr-alternatywa}
dokładnie jeden z nich należy do $\mathcal U$.
Pozostałe własności filtru wynikają z własności
przeciwobrazów i tego, że $p_i^{-1}(X_i)=X$.

Zwartość $X_i$ i twierdzenie
\ref{thm:topog-zwartosc-ultrafiltry} dają punkt
$x_i\in X_i$, do którego zbiega $\mathcal U_i$.
Wybierzmy jednocześnie takie punkty dla wszystkich
$i\in I$ (korzystamy tu z aksjomatu wyboru).
Otrzymujemy $x=(x_i)_{i\in I}\in X$.
Dla bazowego otoczenia $x$ w topologii produktowej
\[
W=\bigcap_{j\in J}p_j^{-1}(V_j),\qquad J\subseteq I
\text{ skończony},\quad x_j\in V_j\in\tau_j,
\]
każdy $V_j$ należy do $\mathcal U_j$, a zatem
$p_j^{-1}(V_j)\in\mathcal U$.
Filtr jest domknięty na skończone przecięcia,
więc $W\in\mathcal U$.
Ponieważ zbiory $W$ tworzą bazę otoczeń $x$,
ultrafiltr $\mathcal U$ zbiega do $x$.
Na mocy twierdzenia \ref{thm:topog-zwartosc-ultrafiltry}
iloczyn $X$ jest zwarty.
\end{proof}

\begin{stwierdzenie}[Przypadek skończonego iloczynu]
\label{prop:topog-zwartosc-iloczyn-skonczony}
Dla $n\ge1$ iloczyn przestrzeni zwartych
$X_1\times\cdots\times X_n$ jest zwarty.
Jeżeli ponadto wszystkie te przestrzenie są
Hausdorffa, ich iloczyn jest Hausdorffa.
\end{stwierdzenie}
\begin{proof}
Zwartość wynika z twierdzenia
\ref{thm:topog-tychonoff-iloczyny}.
Jeśli dwa punkty iloczynu są różne, różnią się
w pewnej współrzędnej $i$. Dla tej współrzędnej
wybieramy rozłączne otwarte otoczenia odpowiednich
punktów w $X_i$. Ich przeciwobrazy przez $p_i$
są rozłącznymi otwartymi otoczeniami punktów
w iloczynie, co dowodzi aksjomatu Hausdorffa.
\end{proof}

\begin{twierdzenie}[Twierdzenie Heinego--Borela]
\label{thm:topog-heine-borel}
Dla każdego $n\ge1$ podzbiór $K\subseteq\R^n$
z topologią euklidesową jest zwarty wtedy i tylko
wtedy, gdy jest domknięty i ograniczony.
\end{twierdzenie}
\begin{proof}
Załóżmy, że $K$ jest zwarty.
Przestrzeń $\R^n$ jest Hausdorffa, więc
$K$ jest domknięty na mocy twierdzenia
\ref{thm:topog-zwartosc-hausdorff-domkniete}.
Zbiory $(-m,m)^n$, $m\ge1$, stanowią pokrycie
$\R^n$, więc pewien z nich zawiera $K$:
z każdego skończonego podpokrycia możemy wybrać
największe $m$. To znaczy, że $K$ jest ograniczony.

Odwrotnie, niech $K$ będzie domknięty i ograniczony.
Istnieje $M>0$ z $K\subseteq[-M,M]^n$.
Każdy przedział $[-M,M]$ jest zwarty na mocy
\ref{thm:topog-zwartosc-odcinek}, a jego
$n$-krotny iloczyn jest zwarty na mocy
\ref{prop:topog-zwartosc-iloczyn-skonczony}.
W topologii euklidesowej $\R^n$ topologia
produktowa współrzędnych jest tą samą topologią:
dla metryki maksimum $d_\infty$ i euklidesowej
$d_2$ zachodzi $d_\infty(x,y)\le d_2(x,y)
\le\sqrt n\,d_\infty(x,y)$, a zatem kule
obu metryk są wzajemnie porównywalne i wyznaczają
tę samą rodzinę zbiorów otwartych. Kule metryki
maksimum są właśnie produktami otwartych
przedziałów o wspólnym promieniu.
Zbiór $K$ jest domknięty w $[-M,M]^n$, zatem
jest zwarty na mocy
\ref{thm:topog-zwartosc-trwalosc}(i).
\end{proof}

\begin{figure}[htbp]
\centering
\begin{tikzpicture}[scale=.87,>=Stealth]
\begin{scope}[shift={(0,0)}]
 \fill[blue!10] (0,0) rectangle (3.15,2.5);
 \draw[blue!70!black,thick] (0,0) rectangle (3.15,2.5);
 \foreach \x in {.63,1.26,1.89,2.52} {
   \draw[blue!32] (\x,0)--(\x,2.5);
 }
 \foreach \y in {.5,1,1.5,2} {
   \draw[blue!32] (0,\y)--(3.15,\y);
 }
 \node[fill=white,inner sep=2pt,font=\small] at (1.57,1.25)
   {$[-M,M]\times[-M,M]$};
 \node[font=\small,align=center] at (1.57,-.48)
   {iloczyn dwóch\\zwartych odcinków};
\end{scope}
\begin{scope}[shift={(5,0)}]
 \fill[blue!10] (0,0) rectangle (3.15,2.5);
 \draw[blue!70!black,thick] (0,0) rectangle (3.15,2.5);
 \filldraw[orange!30,draw=orange!85!black,thick]
   (1.52,1.25) ellipse (0.92 and .73);
 \node[font=\small] at (1.52,1.25) {$K$};
 \node[font=\small,blue!70!black] at (2.32,2.23) {$[-M,M]^2$};
 \node[font=\small,align=center] at (1.57,-.48)
   {domknięty podzbiór\\zwartego prostokąta};
\end{scope}
\end{tikzpicture}
\caption{Dwa kroki dowodu twierdzenia Heinego--Borela:
zwartość prostokąta jako iloczynu zwartych odcinków
oraz trwałość zwartości po przejściu do zbioru domkniętego.
Rysunek ma charakter schematyczny.}
\label{fig:topog-heine-borel}
\end{figure}
\FloatBarrier

\section{Zanurzanie przestrzeni Tychonoffa w kostkę}
\label{sec:topog-zwartosc-kostka}

\begin{twierdzenie}[Zanurzanie w kostkę Tychonoffa]
\label{thm:topog-zwartosc-zanurzenie-kostka}
Dla przestrzeni topologicznej $X$ następujące
warunki są równoważne:
\begin{enumerate}[label=(\roman*)]
\item $X$ spełnia aksjomat Tychonoffa, czyli jest
całkowicie regularna i $T_1$;
\item istnieje zbiór indeksów $I$ oraz zanurzenie
topologiczne $e\colon X\to[0,1]^I$.
\end{enumerate}
Można wziąć $I=C(X,[0,1])$, zbiór wszystkich
ciągłych funkcji $X\to[0,1]$, oraz
$e(x)=(f(x))_{f\in I}$.
\end{twierdzenie}
\begin{proof}
Załóżmy (i). Zbiór $I=C(X,[0,1])$ jest zbiorem,
ponieważ jest podzbiorem zbioru wszystkich funkcji
$[0,1]^X$. Definiujemy
$e(x)=(f(x))_{f\in I}$.
Każda współrzędna $p_f\circ e=f$ jest ciągła,
więc $e$ jest ciągłe na mocy własności uniwersalnej
topologii produktowej z rozdziału~04.

Jeżeli $x\ne y$, to $\{y\}$ jest domknięty,
ponieważ $X$ spełnia $T_1$. Całkowita regularność
daje $f\in I$ z $f(x)=0$ i $f(y)=1$.
Zatem $e(x)\ne e(y)$, czyli $e$ jest różnowartościowe.

Pokażemy, że $e^{-1}\colon e(X)\to X$ jest ciągłe.
Niech $U\subseteq X$ będzie otwarty, a $x\in U$.
Zbiór $F=X\setminus U$ jest domknięty i nie zawiera
$x$, więc całkowita regularność daje $f_x\in I$
z $f_x(x)=0$ i $f_x(F)=\{1\}$, o ile $F\ne\varnothing$;
jeśli $F=\varnothing$, wybieramy po prostu funkcję
stałą równą $0$, co wystarcza w poniższym argumencie.
Zbiór
\[
W_x=p_{f_x}^{-1}([0,1/2))\cap e(X)
\]
jest otwarty w $e(X)$, ponieważ $[0,1/2)$ jest otwarty
w $[0,1]$ i rzutowanie $p_{f_x}$ jest ciągłe.
Zawiera $e(x)$ i spełnia $W_x\subseteq e(U)$:
dla $y\in F$ mamy
$f_x(y)=1$, więc $e(y)\notin W_x$.
Ponieważ $x\in U$ było dowolne, każdy punkt
$e(U)$ ma otwarte w $e(X)$ otoczenie zawarte
w $e(U)$. Stąd $e(U)$ jest otwarty w $e(X)$.
To oznacza ciągłość odwrotności i dowodzi,
że $e$ jest zanurzeniem topologicznym.

Odwrotnie, załóżmy (ii). Przestrzeń $[0,1]$ jest
Hausdorffa i całkowicie regularna: dla punktu
$t$ oraz domkniętego $F\subseteq[0,1]$ z $t\notin F$
można, gdy $F\ne\varnothing$, użyć funkcji
$g(s)=|s-t|/(|s-t|+d(s,F))$, gdzie
$d(s,F)=\inf_{z\in F}|s-z|$.
Funkcja $s\mapsto d(s,F)$ jest ciągła:
dla każdego $z\in F$ nierówność trójkąta daje
$|s-z|\le |s-u|+|u-z|$. Biorąc infimum po
$z\in F$ i zamieniając $s$ z $u$, otrzymujemy
$|d(s,F)-d(u,F)|\le |s-u|$.
Mianownik jest dodatni dla $s\ne t$ przez
$|s-t|>0$, a dla $s=t$ przez $d(t,F)>0$,
ponieważ $F$ jest domknięty i $t\notin F$.
Zatem $g$ jest ciągła, $g(t)=0$ i $g|_F=1$.
Dla $F=\varnothing$ wystarczy $g\equiv0$.

Pokażmy bez powoływania się na nieudowodnioną
regułę o iloczynach, że kostka $Q=[0,1]^I$
jest całkowicie regularna. Niech $q\in Q$ i
$F\subseteq Q$ będzie domknięty, $q\notin F$.
Istnieje bazowe otoczenie
$B=\bigcap_{j\in J}p_j^{-1}(V_j)$ punktu $q$
zawarte w $Q\setminus F$, gdzie $J\subseteq I$
jest skończony, a $V_j\subseteq[0,1]$ są otwarte
i zawierają $q_j$.
Dla każdego $j\in J$ wybierzmy ciągłe
$g_j\colon[0,1]\to[0,1]$ z $g_j(q_j)=0$
oraz $g_j=1$ na $[0,1]\setminus V_j$.
Istnienie wynika z całkowitej regularności odcinka;
przy $V_j=[0,1]$ wystarczy $g_j\equiv0$.
Jeśli $J\ne\varnothing$, zdefiniujmy
$G(z)=\max_{j\in J}g_j(z_j)$.
Każda funkcja $g_j\circ p_j$ jest ciągła,
a maksimum dwóch funkcji ciągłych jest ciągłe
na mocy tożsamości
$\max\{a,b\}=(a+b+|a-b|)/2$.
Przez indukcję otrzymujemy ciągłość $G$.
Ponadto $G(q)=0$ oraz $G|_F=1$,
gdyż każdy punkt poza $B$ narusza przynajmniej
jeden z warunków $z_j\in V_j$.
Jeśli $J=\varnothing$, to $B=Q$ oraz
$F=\varnothing$; wówczas weźmy $G\equiv0$.
Ponadto $Q$ jest $T_1$: dla $q\ne r$ istnieje
współrzędna $i$ z $q_i\ne r_i$, a rozłączne
otoczenia w $[0,1]$ dają rozłączne otoczenia
w $Q$. Zatem $Q$ jest Tychonoffa.
Z twierdzenia o dziedziczeniu z rozdziału~05
każda podprzestrzeń $Q$, w szczególności $e(X)$,
jest Tychonoffa. Homeomorfizm
$X\cong e(X)$ przenosi tę własność na $X$.
\end{proof}

\begin{wniosek}[Zwartość i kostkowe uzwarcenie]
\label{cor:topog-zwartosc-uzwarcenie}
Niech $X$ będzie przestrzenią Tychonoffa i niech
$e\colon X\to Q=[0,1]^I$ będzie zanurzeniem
z poprzedniego twierdzenia. Wówczas
$K=\overline{e(X)}^{\,Q}$ jest zwartą przestrzenią
Hausdorffa, a $e(X)$ jest gęsty w $K$.
Jeśli ponadto $X$ jest zwarta, to $e(X)$ jest
domknięty w $Q$ i $K=e(X)$.
\end{wniosek}
\begin{proof}
Na mocy twierdzenia \ref{thm:topog-tychonoff-iloczyny}
kostka $Q$ jest zwarta. Jest Hausdorffa,
jak wykazano w dowodzie twierdzenia
\ref{thm:topog-zwartosc-zanurzenie-kostka}.
Zbiór $K$ jest domknięty w $Q$, więc zwarty
na mocy twierdzenia \ref{thm:topog-zwartosc-trwalosc}.
Jako podprzestrzeń $Q$ jest Hausdorffa.
Gęstość $e(X)$ w $K$ wynika z definicji domknięcia.
Jeśli $X$ jest zwarta, obraz $e(X)$ jest zwarty
jako obraz ciągły, a więc domknięty w
przestrzeni Hausdorffa $Q$. Stąd $K=e(X)$.
\end{proof}

\begin{przyklad}[Prosty model zanurzenia i domknięcia]
\label{ex:topog-zwartosc-uzwarcenie-odcinka}
Funkcja $e\colon\R\to(0,1)$ dana wzorem
$e(t)=\tfrac12+\tfrac{1}{\pi}\arctan t$
jest homeomorfizmem $\R$ na $(0,1)$:
jest ciągła i ściśle rosnąca, jej granice przy
$t\to\pm\infty$ są równe odpowiednio $0$ i $1$,
a odwrotność $t=\tan(\pi(s-1/2))$ jest ciągła
na $(0,1)$. Jako odwzorowanie $\R\to[0,1]$ funkcja $e$ jest
więc zanurzeniem o niedomkniętym obrazie $(0,1)$.
Jego domknięcie $[0,1]$ jest zwarte,
mimo że $\R$ nie jest zwarta.
\end{przyklad}

\section{Podsumowanie i uwagi o założeniach}
\label{sec:topog-zwartosc-podsumowanie}

\begin{center}
\begin{tabular}{@{}p{.33\textwidth}p{.51\textwidth}@{}}
\hline
\textbf{Charakterystyka} & \textbf{Zakres obowiązywania} \\ \hline
Pokrycia otwarte i przecięcia zbiorów domkniętych &
Wszystkie przestrzenie topologiczne.\\
Punkt skupienia każdej sieci; zbieżna podsieć &
Wszystkie przestrzenie topologiczne.\\
Zbieżność każdego ultrafiltru &
Wszystkie przestrzenie topologiczne; dowód przedłużania
filtru używa lematu Zorna.\\
Osiąganie kresów (Weierstrass) &
Każda ciągła funkcja rzeczywista na niepustej
przestrzeni zwartej (bez aksjomatów oddzielania).\\
Zwartość dowolnego iloczynu &
Dla topologii produktowej; w ogólnym przypadku
stosujemy aksjomat wyboru.\\
Domkniętość każdego zwartego podzbioru &
Jeśli przestrzeń otaczająca jest Hausdorffa.\\
Zanurzenie w $[0,1]^I$ &
Dokładnie dla przestrzeni Tychonoffa ($T_{3\frac12}$).\\
\hline
\end{tabular}
\end{center}

\begin{uwaga}[Granice kryteriów ciągowych]
\label{rem:topog-zwartosc-ciagowe}
W ogólnej przestrzeni topologicznej nie zastępujemy
w definicji zwartości sieci dowolnymi ciągami.
Charakterystyka przez sieci i ultrafiltry jest
bezwarunkowa co do aksjomatów oddzielania,
podczas gdy równoważność zwartości ze zwartością
ciągową wymaga dodatkowych hipotez i zostanie
omówiona w odpowiednim późniejszym rozdziale.
\end{uwaga}

\chapter{Spójność i spójność drogowa}
\label{ch:topog-spojnosc}

\noindent\textbf{Cel rozdziału.}
Zwartość, badana w rozdziale~07, pozwala wybierać skończone
podpokrycia. Spójność wyraża inne zjawisko: niemożność
podziału przestrzeni na dwa niepuste, rozłączne zbiory otwarte.
Zbadamy jej charakteryzacje, zachowanie przy odwzorowaniach
ciągłych i iloczynach oraz składowe spójne. Następnie
odróżnimy spójność od spójności drogowej i omówimy
spójność lokalną. We wszystkich stwierdzeniach, jeśli
nie zaznaczono inaczej, nie zakładamy aksjomatów oddzielania.

\section{Spójność i jej charakteryzacje}
\label{sec:topog-spojnosc-definicja}

\begin{definicja}[Rozdzielenie i przestrzeń spójna]
\label{def:topog-spojnosc}
Niech $(X,\tau)$ będzie przestrzenią topologiczną.
\emph{Rozdzieleniem} $X$ nazywamy parę $(U,V)$ niepustych,
rozłącznych zbiorów otwartych w $X$ takich, że $X=U\cup V$.
Przestrzeń $X$ jest \emph{spójna}, jeżeli nie istnieje
jej rozdzielenie. Podzbiór $A\subseteq X$ jest spójny,
gdy jest spójny z topologią podprzestrzeni.
Pusta przestrzeń jest zatem spójna.
\end{definicja}

\begin{twierdzenie}[Charakteryzacje spójności]
\label{thm:topog-spojnosc-charakteryzacje}
Dla przestrzeni topologicznej $X$ następujące warunki
są równoważne:
\begin{enumerate}[label=\textup{(\roman*)}]
\item $X$ jest spójna;
\item jedynymi zbiorami równocześnie otwartymi i domkniętymi
w $X$ są $\varnothing$ i $X$;
\item każde ciągłe odwzorowanie $f:X\to\{0,1\}$,
gdzie $\{0,1\}$ ma topologię dyskretną, jest stałe.
\end{enumerate}
\end{twierdzenie}
\begin{proof}
Jeśli istnieje rozdzielenie $X=U\cup V$, to $U$
jest niepusty i właściwy, otwarty, a także domknięty,
bo $X\setminus U=V$ jest otwarty. Odwrotnie, dla
niepustego właściwego zbioru otwarto-domkniętego $U$
para $(U,X\setminus U)$ jest rozdzieleniem. Dowodzi
to równoważności (i) i (ii).

Jeśli $U$ jest niepustym właściwym zbiorem
otwarto-domkniętym, funkcja $f:X\to\{0,1\}$
równa $0$ na $U$ i $1$ na jego dopełnieniu jest
ciągła: przeciwobrazy obu singletonów są otwarte.
Nie jest stała. I odwrotnie, jeśli $f$ jest ciągła
i niestała, oba przeciwobrazy $f^{-1}(\{0\})$,
$f^{-1}(\{1\})$ są niepuste oraz otwarte,
a wzajemnie się dopełniają. Uzyskaliśmy
rozdzielenie $X$. Dla $X=\varnothing$ wszystkie
trzy warunki zachodzą.
\end{proof}

\begin{przyklad}[Przestrzenie dyskretne i niedyskretne]
\label{ex:topog-spojnosc-pierwsze}
Przestrzeń z topologią antydyskretną
$\{\varnothing,X\}$ jest spójna: nie ma niepustego
właściwego zbioru otwartego. Z kolei przestrzeń
dyskretna o co najmniej dwóch punktach nie jest
spójna, ponieważ każdy singleton jest
otwarto-domknięty. Jednopunktowa przestrzeń
jest spójna niezależnie od topologii.
\end{przyklad}

\begin{przyklad}[Spójność nie oznacza aksjomatu Hausdorffa]
\label{ex:topog-spojnosc-sierpinski}
Przestrzeń Sierpińskiego $X=\{0,1\}$ z topologią
$\{\varnothing,\{1\},X\}$ jest spójna: jej jedyne
zbiory otwarto-domknięte to $\varnothing$ i $X$.
Jednocześnie nie jest przestrzenią $T_1$, a tym
bardziej Hausdorffa (por. rozdział~05).
\end{przyklad}

\begin{twierdzenie}[Spójność przedziałów]
\label{thm:topog-spojnosc-przedzialy}
Każdy przedział $I\subseteq\R$, z topologią
indukowaną przez topologię euklidesową, jest
spójny. Odwrotnie, każdy spójny podzbiór
$A\subseteq\R$ jest przedziałem, przy czym
dopuszczamy przedział pusty i jednopunktowy.
\end{twierdzenie}
\begin{proof}
Najpierw niech $a<b$. Przypuśćmy, że $[a,b]$
ma rozdzielenie $U\cup V$. Wybierzmy $u\in U$
i $v\in V$; po ewentualnej zamianie nazw
możemy przyjąć $u<v$. W zbiorze $[u,v]$
rozważmy $S=U\cap[u,v]$ i liczbę $c=\sup S$.
Istnienie kresu wynika z niepustości $S$ i
ograniczenia go z góry przez $v$.
Zbiór $U$ jest względnie domknięty w $[a,b]$,
bo $V=[a,b]\setminus U$ jest względnie otwarty.
Wobec tego $c\in U$: istotnie, z definicji
kresu istnieją elementy $S$ dowolnie bliskie $c$,
a $U$ zawiera swoje punkty domknięcia w $[a,b]$.
Jeżeli $c<v$, względna otwartość $U$ daje
punkt $s\in U$ z $c<s<v$, wbrew definicji $c$.
Jeżeli $c=v$, otrzymujemy $v\in U\cap V$, również
sprzeczność. Zatem $[a,b]$ jest spójny.
Przedział jednopunktowy lub pusty jest spójny
bezpośrednio z definicji.

Niech teraz $I$ będzie dowolnym przedziałem.
Gdyby $I=U\cup V$ było rozdzieleniem,
wybralibyśmy $u\in U$, $v\in V$; ponieważ
$I$ jest przedziałem, cały odcinek $[u,v]$
(z końcami w dowolnej kolejności) zawierałby
się w $I$. Ograniczenia $U\cap[u,v]$ i
$V\cap[u,v]$ rozdzielałyby ten odcinek,
wbrew pierwszej części dowodu.

Odwrotnie, jeśli $A\subseteq\R$ nie jest
przedziałem, istnieją $a,b\in A$ oraz
$c\notin A$ takie, że $a<c<b$.
Zbiory $A\cap(-\infty,c)$ i $A\cap(c,\infty)$
są wówczas niepuste, rozłączne, otwarte
w $A$ i pokrywają $A$. Zatem $A$ nie jest
spójny.
\end{proof}

\begin{twierdzenie}[Zachowanie spójności przez odwzorowania]
\label{thm:topog-spojnosc-obraz-ciagly}
Niech $X,Y$ będą przestrzeniami topologicznymi,
$f:X\to Y$ odwzorowaniem ciągłym, a $A\subseteq X$
spójnym podzbiorem. Wówczas $f(A)$ jest spójny.
W szczególności spójność jest zachowywana
przez homeomorfizmy oraz odwzorowania ilorazowe
(przy obrazie całej przestrzeni).
\end{twierdzenie}
\begin{proof}
Ograniczenie $f|_A:A\to f(A)$ jest ciągłe
i surjektywne; ciągłość wynika z definicji
topologii podprzestrzeni. Gdyby $f(A)$
miało rozdzielenie $U\cup V$, przeciwobrazy
$(f|_A)^{-1}(U)$ i $(f|_A)^{-1}(V)$
byłyby rozłączne, otwarte w $A$, niepuste
dzięki surjektywności i pokrywałyby $A$.
To sprzeczne ze spójnością $A$.
Homeomorfizm jest ciągły i surjektywny
na swój obraz, a odwzorowanie ilorazowe
jest, z definicji, ciągłe i surjektywne.
\end{proof}

\begin{wniosek}[Własność wartości pośredniej]
\label{cor:topog-spojnosc-wartosci-posredniej}
Niech $X$ będzie spójną przestrzenią topologiczną,
$f:X\to\R$ funkcją ciągłą, a $x,y\in X$.
Wtedy dla każdej liczby $t$ leżącej pomiędzy
$f(x)$ i $f(y)$ istnieje $z\in X$ takie,
że $f(z)=t$.
\end{wniosek}
\begin{proof}
Z twierdzenia~\ref{thm:topog-spojnosc-obraz-ciagly}
zbiór $f(X)$ jest spójny, a z twierdzenia
\ref{thm:topog-spojnosc-przedzialy} wynika,
że jest przedziałem. Zawiera $f(x)$ oraz
$f(y)$, więc zawiera także każdą liczbę
pomiędzy nimi.
\end{proof}

\begin{wniosek}[Zwartość, spójność i osiąganie wartości]
\label{cor:topog-spojnosc-zwarty-przedzial}
Jeśli $X$ jest niepustą, zwartą i spójną
przestrzenią topologiczną, a $f:X\to\R$
jest ciągła, to istnieją $m,M\in\R$, $m\le M$,
takie, że $f(X)=[m,M]$. Oba kresy są
osiągane przez $f$.
\end{wniosek}
\begin{proof}
Z twierdzenia Weierstrassa z rozdziału~07
wynika istnienie i osiąganie najmniejszej
i największej wartości $m,M$ funkcji $f$.
Na mocy poprzedniego wniosku każda liczba
pomiędzy nimi należy do $f(X)$; z definicji
kresów nic spoza $[m,M]$ nie należy do obrazu.
\end{proof}

\section{Operacje zachowujące spójność}
\label{sec:topog-spojnosc-operacje}

\begin{twierdzenie}[Suma zbiorów spójnych o części wspólnej]
\label{thm:topog-spojnosc-suma}
Niech $X$ będzie przestrzenią topologiczną,
a $(A_i)_{i\in I}$ rodziną niepustych spójnych
podzbiorów $X$. Jeżeli $\bigcap_{i\in I}A_i\ne
\varnothing$, to $\bigcup_{i\in I}A_i$ jest spójny.
Dla $I=\varnothing$ suma jest pusta i zatem spójna.
\end{twierdzenie}
\begin{proof}
Załóżmy $I\ne\varnothing$ i wybierzmy
$p\in\bigcap_{i\in I}A_i$. Gdyby suma
$A=\bigcup_i A_i$ miała rozdzielenie $U\cup V$,
punkt $p$ należałby do jednego z tych zbiorów,
powiedzmy do $U$. Dla każdego $i$ ograniczenia
$U\cap A_i$ i $V\cap A_i$ są otwarte w $A_i$;
ponieważ $A_i$ jest spójny i $p\in U\cap A_i$,
mamy $A_i\subseteq U$. Zatem $A\subseteq U$,
co przeczy niepustości $V$.
\end{proof}

\begin{figure}[htbp]
\centering
\begin{tikzpicture}[scale=1]
  \fill[blue!17] (-1,0) circle (1.2);
  \fill[orange!20] (1,0) circle (1.2);
  \draw[blue!60!black,thick] (-1,0) circle (1.2);
  \draw[orange!65!black,thick] (1,0) circle (1.2);
  \fill (0,0) circle (1.6pt);
  \node[above] at (0,.05) {$p$};
  \node at (-1.15,0) {$A_1$};
  \node at (1.15,0) {$A_2$};
\end{tikzpicture}
\caption{Dwa zbiory spójne $A_1,A_2$ o niepustej
części wspólnej (z zaznaczonym punktem $p$).
Ich suma jest spójna. Rysunek ilustruje
twierdzenie, nie stanowi jego dowodu.}
\label{fig:topog-spojnosc-suma}
\end{figure}

\begin{twierdzenie}[Domknięcie zbioru spójnego]
\label{thm:topog-spojnosc-domkniecie}
Niech $A\subseteq X$ będzie spójny i niech
$A\subseteq B\subseteq\overline A$, gdzie
domknięcie jest liczone w $X$. Wówczas $B$
jest spójny. W szczególności domknięcie
zbioru spójnego jest spójne.
\end{twierdzenie}
\begin{proof}
Załóżmy, że $B=U\cup V$ jest rozdzieleniem.
Zbiory $U\cap A$, $V\cap A$ są otwarte w $A$
i rozłączne. Co najmniej jeden z nich
jest pusty, ponieważ $A$ jest spójny.
Załóżmy $A\subseteq U$. Wtedy $V$ jest
niepustym zbiorem otwartym w $B$, więc
$V=B\cap W$ dla pewnego otwartego $W\subseteq X$.
Wybierzmy $v\in V$. Ponieważ $v\in\overline A$
i $v\in W$, zachodzi $W\cap A\ne\varnothing$,
a z drugiej strony $W\cap A=V\cap A=\varnothing$,
gdyż $A\subseteq B$. Otrzymujemy sprzeczność.
Przypadek $A\subseteq V$ jest symetryczny.
\end{proof}

\begin{przyklad}[Przecięcie zbiorów spójnych nie musi być spójne]
\label{ex:topog-spojnosc-przeciecie}
Okrąg $S^1=\{(x,y)\in\R^2:x^2+y^2=1\}$
jest spójny: jest ciągłym obrazem odcinka
$[0,2\pi]$ przez $t\mapsto(\cos t,\sin t)$.
Oś pozioma $L=\R\times\{0\}$ jest spójna,
ponieważ jest homeomorficzna z $\R$.
Jednak $S^1\cap L=\{(-1,0),(1,0)\}$
nie jest spójny, gdyż oba singletony są
otwarte w podprzestrzeni dwupunktowej.
\end{przyklad}

\begin{twierdzenie}[Spójność iloczynów]
\label{thm:topog-spojnosc-iloczyn}
Niech $(X_i)_{i\in I}$ będzie rodziną
przestrzeni spójnych. Wówczas iloczyn
$\prod_{i\in I}X_i$ z topologią produktową
jest spójny, także gdy jest pusty.
\end{twierdzenie}
\begin{proof}
Najpierw rozważmy dwie niepuste przestrzenie
spójne $X,Y$. Wybierzmy $y_0\in Y$.
Zbiór $B=X\times\{y_0\}$ jest spójny,
bo jest homeomorficzny z $X$.
Dla każdego $x\in X$ zbiór
$A_x=\{x\}\times Y$ jest spójny,
bo jest homeomorficzny z $Y$, a
$A_x\cap B=\{(x,y_0)\}\ne\varnothing$.
Z twierdzenia~\ref{thm:topog-spojnosc-suma}
wynika, że $A_x\cup B$ jest spójny.
Wszystkie zbiory $A_x\cup B$ zawierają $B$,
więc ich suma $X\times Y$ jest spójna.
Indukcyjnie otrzymujemy spójność dowolnego
skończonego iloczynu niepustych przestrzeni
spójnych; iloczyn pusty jest jednopunktowy.

Niech teraz $X=\prod_{i\in I}X_i$.
Jeżeli $X=\varnothing$, jest spójny z definicji.
W przeciwnym razie ustalmy $p=(p_i)_{i\in I}\in X$.
Dla każdego skończonego $J\subseteq I$
niech $Y_J\subseteq X$ składa się z punktów
$x$, dla których $x_i=p_i$ przy $i\notin J$.
Przestrzeń $Y_J$ jest homeomorficzna ze
skończonym iloczynem $\prod_{i\in J}X_i$,
a więc spójna. Wszystkie $Y_J$ zawierają $p$.
Na mocy twierdzenia o sumie zbiorów spójnych
zbiór $D=\bigcup_{J\subseteq I,\,J\,\text{skończony}}Y_J$
jest spójny.

Wykażmy, że $D$ jest gęsty w $X$.
Niech $W$ będzie niepustym otwartym zbiorem
w $X$. Z definicji topologii produktowej
istnieje punkt $q\in W$ oraz podstawowy
zbiór otwarty $O$ taki, że $q\in O\subseteq W$
i $O=\prod_{i\in I}O_i$, gdzie $O_i=X_i$
poza skończonym zbiorem $J$. Zdefiniujmy
$x_i=q_i$ dla $i\in J$ oraz $x_i=p_i$
dla $i\notin J$. Wówczas $x\in Y_J\cap O$,
więc $W\cap D\ne\varnothing$.
Oznacza to $\overline D=X$, a twierdzenie
\ref{thm:topog-spojnosc-domkniecie}
daje spójność $X$. Nie wybieraliśmy
niezależnie punktów z całej rodziny czynników:
w przypadku niepustego iloczynu jeden punkt
$p$ został już ustalony.
\end{proof}

\begin{uwaga}[Spójność i zwartość iloczynów]
\label{rem:topog-spojnosc-tychonoff}
Twierdzenie~\ref{thm:topog-spojnosc-iloczyn}
nie wymaga twierdzenia Tychonoffa ani
aksjomatu wyboru w dowodzie, gdy iloczyn
jest zadany jako niepusty. Inaczej niż przy
zwartości, wystarczają spójność skończonych
iloczynów i domknięć. Odrębną kwestią jest
wnioskowanie o \emph{niepustości} iloczynu
samych niepustych czynników; w ogólności
wymaga to zasady wyboru.
\end{uwaga}

\section{Składowe spójne i całkowita niespójność}
\label{sec:topog-spojnosc-skladowe}

\begin{definicja}[Składowa spójna]
\label{def:topog-spojnosc-skladowa}
Dla przestrzeni $X$ i punktu $x\in X$
\emph{składową spójną} $C(x)$ nazywamy sumę
wszystkich spójnych podzbiorów $X$
zawierających $x$.
Przestrzeń $X$ jest \emph{całkowicie niespójna},
jeśli wszystkie jej składowe spójne są
jednopunktowe (dla $X=\varnothing$ warunek
jest spełniony).
\end{definicja}

\begin{twierdzenie}[Własności składowych spójnych]
\label{thm:topog-spojnosc-skladowe}
Niech $X$ będzie przestrzenią topologiczną.
Każda składowa spójna jest spójna, domknięta
i maksymalna względem zawierania wśród
spójnych podzbiorów $X$. Dwie składowe
spójne są równe albo rozłączne; w konsekwencji
składowe spójne tworzą rozkład $X$.
Nie muszą być otwarte.
\end{twierdzenie}
\begin{proof}
Dla $x\in X$ rodzina definiująca $C(x)$
jest niepusta, bo zawiera $\{x\}$,
a wszystkie jej elementy mają wspólny punkt
$x$. Twierdzenie~\ref{thm:topog-spojnosc-suma}
daje spójność $C(x)$. Jeśli spójny zbiór
$B$ zawiera $C(x)$, to zawiera również $x$,
więc jest jednym ze zbiorów uczestniczących
w definicji $C(x)$; stąd $B\subseteq C(x)$,
co dowodzi maksymalności.
Domknięcie $\overline{C(x)}$ jest spójne
na mocy twierdzenia
\ref{thm:topog-spojnosc-domkniecie}.
Zawiera $C(x)$, więc z maksymalności
$\overline{C(x)}=C(x)$.

Jeśli $C(x)\cap C(y)\ne\varnothing$, suma
tych dwóch zbiorów jest spójna, a zawiera
$x$ oraz $y$. Z maksymalności obu składowych
wynika $C(x)=C(y)$. Każdy punkt należy do
swojej składowej, więc składowe pokrywają $X$.
\end{proof}

\begin{przyklad}[Składowe nie muszą być otwarte]
\label{ex:topog-spojnosc-Q}
W $\Q$ z topologią podprzestrzeni z $\R$
każdy spójny podzbiór ma co najwyżej jeden
punkt. Istotnie, jeśli zawierałby różne
liczby wymierne $p<q$, to dowolna liczba
niewymierna $r\in(p,q)$ wyznaczałaby
jego rozdzielenie przez przecięcia z
$(-\infty,r)$ i $(r,\infty)$.
Zatem składowymi spójnymi $\Q$ są singletony.
Nie są one otwarte w $\Q$, bo każde
względnie otwarte otoczenie liczby wymiernej
zawiera inne liczby wymierne.
Przestrzeń $\Q$ jest więc całkowicie
niespójna, ale nie dyskretna.
\end{przyklad}

\begin{przyklad}[Zbiór Cantora]
\label{ex:topog-spojnosc-cantor}
Niech $\mathcal C\subseteq[0,1]$ będzie klasycznym
zbiorem Cantora. Jego punkty można jednoznacznie
zapisać jako $x=\sum_{n\ge1}a_n3^{-n}$,
gdzie $a_n\in\{0,2\}$. Dla różnych takich
ciągów pierwsza różna cyfra na pozycji $k$
powoduje różnicę wartości co najmniej
$2\cdot3^{-k}-\sum_{n>k}2\cdot3^{-n}=3^{-k}>0$.
Zatem zapis jest jednoznaczny.

Dla ustalonego skończonego ciągu cyfr
$a_1,\ldots,a_m$ podzbiór punktów o tych
pierwszych $m$ cyfrach jest otwarty
i domknięty w $\mathcal C$. Wynika to z tego,
że punktom odpowiadają ciągi cyfr, a suma
pozostałych wyrazów jest co najwyżej
$\sum_{n>m}2\cdot3^{-n}=3^{-m}$:
zbiory o różnych prefiksach długości $m$
leżą w rozłącznych przedziałach domkniętych
z dodatnimi przerwami i tworzą skończony
rozkład $\mathcal C$. Dwa różne punkty
$\mathcal C$ mają różne cyfry na pewnej pozycji,
więc takie względnie otwarto-domknięte
zbiory je rozdzielają. Każdy spójny podzbiór
$\mathcal C$ zawierający dwa różne punkty
miałby wówczas rozdzielenie. Zatem
$\mathcal C$ jest całkowicie niespójny.

\begin{samepage}
Zbiór Cantora jest zarazem zwarty. Aby to
uzasadnić, dla $m\ge1$ niech
\[
K_m=\bigcup_{(a_1,\ldots,a_m)\in\{0,2\}^m}
\left[\sum_{j=1}^m a_j3^{-j},\,
\sum_{j=1}^m a_j3^{-j}+3^{-m}\right].
\]
\end{samepage}

Każdy $K_m$ jest skończoną sumą
przedziałów domkniętych, a więc jest
domknięty w $\R$. Z reprezentacji cyframi
wynika $\mathcal C\subseteq\bigcap_m K_m$.
Dla inkluzji odwrotnej zauważmy, że zbiory
przedziałów w $K_m$ są zagnieżdżone:
przedział odpowiadający pierwszym $m+1$
cyfrom zawiera się w przedziale wyznaczonym
pierwszymi $m$ cyframi. Dla
$x\in\bigcap_m K_m$ jego położenie
w jednoznacznie wyznaczonym przedziale
$K_m$ dostarcza więc zgodnych ze sobą
prefiksów cyfr $a_1,\ldots,a_m$.
Ponieważ długości tych przedziałów to $3^{-m}$,
otrzymujemy $x=\sum_{n\ge1}a_n3^{-n}$.
Wobec tego $\mathcal C=\bigcap_m K_m$
jest domknięty. Zawiera się w zwartym
przedziale $[0,1]$, zatem jest zwarty
na mocy twierdzenia z rozdziału~07.
\end{przyklad}

\section{Drogi i spójność drogowa}
\label{sec:topog-spojnosc-drogi}

\begin{definicja}[Droga i spójność drogowa]
\label{def:topog-spojnosc-droga}
Niech $X$ będzie przestrzenią topologiczną.
\emph{Drogą} od $x$ do $y$ nazywamy ciągłe
odwzorowanie $\gamma:[0,1]\to X$ takie,
że $\gamma(0)=x$ i $\gamma(1)=y$.
Przestrzeń $X$ jest \emph{spójna drogowo},
jeśli dla każdych $x,y\in X$ istnieje
droga od $x$ do $y$.
\end{definicja}

\begin{uwaga}[Droga a łuk]
\label{rem:topog-spojnosc-luk}
W literaturze nazwy \emph{spójność drogowa}
i \emph{spójność łukowa} bywają używane
zamiennie. Aby uniknąć niejednoznaczności,
tu słowem \emph{droga} określamy dowolne
odwzorowanie ciągłe $[0,1]\to X$;
nie wymagamy jego różnowartościowości.
\emph{Łukiem} będziemy nazywać odwzorowanie
homeomorficzne $[0,1]$ na swój obraz.
Twierdzenia niniejszej sekcji dotyczą dróg.
\end{uwaga}

\begin{lemat}[Odwracanie i składanie dróg]
\label{lem:topog-spojnosc-skladanie-drog}
Jeśli istnieją drogi od $x$ do $y$ oraz
od $y$ do $z$, to istnieje droga od $x$
do $z$. Każda droga od $x$ do $y$ daje
także drogę od $y$ do $x$.
\end{lemat}
\begin{proof}
Jeśli $\gamma$ jest drogą od $x$ do $y$,
to $t\mapsto\gamma(1-t)$ jest drogą od
$y$ do $x$, gdyż $t\mapsto1-t$ jest ciągła.
Jeśli $\delta$ jest drogą od $y$ do $z$,
określmy
\[
\eta(t)=\begin{cases}
\gamma(2t),&0\le t\le\tfrac12,\\
\delta(2t-1),&\tfrac12\le t\le1.
\end{cases}
\]
Obie części są ciągłe na swoich przedziałach,
a w $t=\tfrac12$ ich wartości są równe $y$.
Z lematu o sklejaniu odwzorowań ciągłych
(rozdział~03) wynika ciągłość $\eta$.
Ponadto $\eta(0)=x$, $\eta(1)=z$.
\end{proof}

\begin{twierdzenie}[Spójność drogowa implikuje spójność]
\label{thm:topog-spojnosc-droga-implikuje}
Każda przestrzeń spójna drogowo jest spójna.
Każdy ciągły obraz przestrzeni spójnej
drogowo jest spójny drogowo.
\end{twierdzenie}
\begin{proof}
Niech $X$ będzie spójna drogowo.
Gdyby istniało rozdzielenie $X=U\cup V$,
wybralibyśmy $x\in U$, $y\in V$ i drogę
$\gamma:[0,1]\to X$ od $x$ do $y$.
Odcinek $[0,1]$ jest spójny z twierdzenia
\ref{thm:topog-spojnosc-przedzialy}, a jego
obraz $\gamma([0,1])$ jest spójny z twierdzenia
\ref{thm:topog-spojnosc-obraz-ciagly}.
Przecina jednak oba zbiory $U,V$,
więc ich przecięcia z obrazem rozdzielałyby
ten obraz, co daje sprzeczność.

Jeżeli $f:X\to Y$ jest ciągła i $X$
spójna drogowo, to dla dowolnych
$a,b\in f(X)$ możemy wybrać
$x,y\in X$ z $f(x)=a$, $f(y)=b$.
Droga $\gamma$ od $x$ do $y$ daje
ciągłą drogę $f\circ\gamma$ od $a$
do $b$ w $f(X)$.
\end{proof}

\begin{przyklad}[Zbiory wypukłe]
\label{ex:topog-spojnosc-wypukle}
Każdy niepusty zbiór wypukły $C$ w
przestrzeni unormowanej jest spójny drogowo.
Dla $x,y\in C$ droga
$\gamma(t)=(1-t)x+ty$, $0\le t\le1$,
przyjmuje wartości w $C$ z wypukłości
i jest ciągła, ponieważ działania liniowe
są ciągłe w przestrzeni unormowanej.
W szczególności kule, odcinki i całe
przestrzenie unormowane są spójne drogowo.
\end{przyklad}

\begin{definicja}[Składowa drogowa]
\label{def:topog-spojnosc-skladowa-drogowa}
Dla $x\in X$ \emph{składową drogową} $P(x)$
nazywamy zbiór wszystkich punktów $y\in X$,
które można połączyć drogą z $x$.
\end{definicja}

\begin{stwierdzenie}[Rozkład na składowe drogowe]
\label{prop:topog-spojnosc-skladowe-drogowe}
Dla dowolnej przestrzeni $X$ składowe drogowe
tworzą rozkład $X$. Każda składowa drogowa
jest spójna drogowo i zawiera się w pewnej
składowej spójnej. Nie musi być ani otwarta,
ani domknięta.
\end{stwierdzenie}
\begin{proof}
Zachodzą: zwrotność (droga stała), symetria
(odwrócenie drogi) i przechodniość
(składanie dróg), na mocy lematu
\ref{lem:topog-spojnosc-skladanie-drog}.
Jest to zatem relacja równoważności,
a jej klasy stanowią rozkład $X$.
Jeśli $y,z\in P(x)$, możemy połączyć
$y$ z $x$ i $x$ z $z$; złożenie dróg
pokazuje, że $P(x)$ jest spójna drogowo,
więc i spójna. Dlatego $P(x)\subseteq C(x)$.
Ostatnie zdanie zilustrujemy konkretnym
przykładem poniżej.
\end{proof}

\section{Krzywa sinusowa topologa}
\label{sec:topog-spojnosc-sinus}

\begin{przyklad}[Spójność bez spójności drogowej]
\label{ex:topog-spojnosc-sinus}
Rozważmy podzbiór płaszczyzny
\[
G=\{(x,\sin(1/x)):0<x\le1\},\qquad
T=G\cup(\{0\}\times[-1,1]).
\]
Zbiór $G$ jest spójny drogowo: jest obrazem
przedziału $(0,1]$ przez ciągłe odwzorowanie
$x\mapsto(x,\sin(1/x))$. Mamy także
$T=\overline G$ w $\R^2$. Aby wykazać zawieranie
$\overline G\subseteq T$, zauważmy najpierw,
że z $G\subseteq(0,1]\times[-1,1]$ wynika
$\overline G\subseteq[0,1]\times[-1,1]$.
Jeśli $(x,y)\in\overline G$ i $x>0$, to
w przestrzeni metrycznej $\R^2$ istnieje ciąg
$(x_n,\sin(1/x_n))\in G$ zbieżny do $(x,y)$
(kryterium domknięcia przez ciągi w przestrzeniach
pierwszo przeliczalnych z rozdziału~06). Ponieważ $x_n\to x>0$,
z ciągłości funkcji $t\mapsto\sin(1/t)$
w punkcie $x$ wynika
$y=\lim_n\sin(1/x_n)=\sin(1/x)$, czyli
$(x,y)\in G$. Pozostałe punkty domknięcia
mogą więc należeć jedynie do
$\{0\}\times[-1,1]$.

Odwrotnie, $G\subseteq\overline G$, a każdy
punkt $(0,y)$ odcinka pionowego jest punktem
domknięcia $G$: wybierzmy $\theta\in\R$
z $\sin\theta=y$. Dla dostatecznie dużych $n$
mamy $\theta+2\pi n>0$, liczbę
$(\theta+2\pi n)^{-1}\in(0,1]$ oraz
$\bigl((\theta+2\pi n)^{-1},y\bigr)\in G$;
przy tym pierwsza współrzędna dąży do $0$.
Otrzymujemy $\overline G=T$.
Z twierdzenia~\ref{thm:topog-spojnosc-domkniecie}
wynika, że $T$ jest spójny.

Wykażemy, że \emph{nie istnieje droga w $T$}
łącząca punkt odcinka pionowego z punktem $G$.
Przypuśćmy, że taka droga istnieje.
Po ewentualnym odwróceniu parametrów możemy
przyjąć $\gamma(0)\in\{0\}\times[-1,1]$
oraz $\gamma(1)\in G$.
Zapiszmy $\gamma(t)=(u(t),v(t))$; wtedy
$u(0)=0$, $u(1)>0$, a $u(t)\ge0$.
Niech $t_0$ będzie największym punktem
zbioru $\{t\in[0,1]:u(t)=0\}$.
Istnieje on, ponieważ zbiór ten jest
niepusty i domknięty w $[0,1]$, a zatem
zwarty. Ponieważ $u(1)>0$, mamy $t_0<1$.
Dla $t\in(t_0,1]$ zachodzi $u(t)>0$,
a więc $v(t)=\sin(1/u(t))$.

Z ciągłości $v$ w punkcie $t_0$
wybierzmy $\delta\in(0,1-t_0)$ takie,
że $|v(t)-v(t_0)|<\tfrac12$ dla wszystkich
$t\in[t_0,t_0+\delta]$. Z ciągłości $u$
i nierówności $u(t_0)=0<u(t_0+\delta)$
wynika, na mocy własności wartości pośredniej,
że $u([t_0,t_0+\delta])$ zawiera
$[0,u(t_0+\delta)]$. W tym ostatnim
przedziale znajdują się dodatnie liczby
\[
a_n=(\tfrac\pi2+2\pi n)^{-1},\qquad
b_n=(\tfrac{3\pi}2+2\pi n)^{-1}
\]
dla dostatecznie dużego $n$.
Odpowiadają im czasy $s,t\in(t_0,t_0+\delta]$
takie, że $v(s)=1$ i $v(t)=-1$.
Z nierówności $|v(s)-v(t_0)|<1/2$
i $|v(t)-v(t_0)|<1/2$ otrzymujemy
$2=|v(s)-v(t)|<1$, sprzeczność.
Zatem $T$ jest spójny, ale nie jest
spójny drogowo.

Co więcej, $G$ i $\{0\}\times[-1,1]$
są dwiema składowymi drogowymi $T$:
każdy z nich jest spójny drogowo,
a przed chwilą wykluczyliśmy drogi
między nimi. Składowa drogowa $G$
nie jest domknięta, ponieważ $\overline G=T$.
Odcinek pionowy nie jest otwarty w $T$,
ponieważ każde jego otoczenie w $T$
przecina $G$.
\end{przyklad}

\begin{figure}[htbp]
\centering
\begin{tikzpicture}[x=7.5cm,y=1.15cm,>=stealth]
  \draw[->,gray!70] (-.08,0)--(1.1,0) node[below right] {$x$};
  \draw[->,gray!70] (0,-1.3)--(0,1.35) node[above] {$y$};
  \draw[red!70!black,very thick] (0,-1)--(0,1);
  \draw[blue!75!black,thick,samples=280,smooth,domain=.035:1]
    plot (\x,{sin(1/\x r)});
  \fill[red!70!black] (0,0) circle (1.9pt);
  \node[anchor=west,red!70!black,font=\small] at (.035,1.18)
    {odcinek pionowy};
  \node[anchor=west,blue!70!black,font=\small] at (.50,-1.2)
    {$G:\ y=\sin(1/x)$};
\end{tikzpicture}
\caption{Krzywa sinusowa topologa $T$: wykres $G$ wraz
z pionowym odcinkiem w $x=0$. Rysunek przedstawia
skończone przybliżenie coraz gęstszych oscylacji,
których nie da się pokonać drogą ciągłą z odcinka
pionowego do wykresu.}
\label{fig:topog-spojnosc-sinus}
\end{figure}
\FloatBarrier

\section{Spójność lokalna}
\label{sec:topog-spojnosc-lokalna}

\begin{definicja}[Spójność lokalna i drogowa lokalna]
\label{def:topog-spojnosc-lokalna}
Przestrzeń $X$ jest \emph{lokalnie spójna},
jeśli dla każdego $x\in X$ i każdego
otwartego otoczenia $U\subseteq X$ punktu $x$
istnieje otwarty spójny zbiór $V$ taki,
że $x\in V\subseteq U$.
Jest \emph{lokalnie spójna drogowo},
jeśli można dodatkowo wymagać, by $V$
był spójny drogowo. Definicje obejmują
przestrzeń pustą.
\end{definicja}

\begin{twierdzenie}[Charakteryzacja spójności lokalnej]
\label{thm:topog-spojnosc-lokalna-komponenty}
Przestrzeń $X$ jest lokalnie spójna wtedy
i tylko wtedy, gdy każda składowa spójna
każdego otwartego zbioru $U\subseteq X$
jest otwarta w $X$.
Analogicznie, $X$ jest lokalnie spójna
drogowo wtedy i tylko wtedy, gdy wszystkie
składowe drogowe wszystkich otwartych
podzbiorów $X$ są otwarte w $X$.
\end{twierdzenie}
\begin{proof}
Przypuśćmy najpierw, że $X$ jest lokalnie
spójna. Ustalmy otwarty $U\subseteq X$
i jego składową spójną $C$ oraz $x\in C$.
Z definicji istnieje otwarty spójny $V_x$
z $x\in V_x\subseteq U$. Ponieważ
$V_x\cap C$ zawiera $x$, twierdzenie
\ref{thm:topog-spojnosc-suma} pokazuje,
że $V_x\cup C$ jest spójny w $U$.
Maksymalność składowej $C$ w $U$ daje
$V_x\subseteq C$. Zatem
$C=\bigcup_{x\in C}V_x$ jest otwarty w $X$.

Odwrotnie, jeśli składowe spójne każdego
otwartego $U$ są otwarte w $X$, to dla
$x\in U$ składowa spójna $C_U(x)$ jest
otwartym spójnym otoczeniem $x$ zawartym
w $U$. Jest to definicja lokalnej spójności.

W wersji drogowej zastępujemy składowe
spójne składowymi drogowymi i stosujemy
lemat~\ref{lem:topog-spojnosc-skladanie-drog}:
dwa spójne drogowo zbiory o niepustej części
wspólnej mają spójną drogowo sumę, ponieważ
drogi można połączyć we wspólnym punkcie.
Pozostały argument jest identyczny.
\end{proof}

\begin{wniosek}[Spójność i drogi w przestrzeni lokalnie drogowej]
\label{cor:topog-spojnosc-lokalna-drogowa}
Jeśli przestrzeń $X$ jest lokalnie spójna
drogowo, to jej składowe drogowe i spójne
są tymi samymi zbiorami. W szczególności,
spójna przestrzeń lokalnie spójna drogowo
jest spójna drogowo.
\end{wniosek}
\begin{proof}
Z poprzedniego twierdzenia każda składowa
drogowa $P$ przestrzeni $X$ jest otwarta
w $X$. Jej dopełnienie jest sumą innych
składowych drogowych, zatem również jest
otwarte. Tak więc $P$ jest otwarto-domknięta.
Jeśli składowa spójna $C$ zawiera $P$,
a także punkt $y\notin P$, to zbiory
$P\cap C$ i $C\setminus P$ byłyby
niepuste, rozłączne, otwarte w $C$
i pokrywałyby $C$. Przeczy to spójności $C$.
Dlatego $C=P$.
\end{proof}

\begin{wniosek}[Spójność zbiorów otwartych w przestrzeni euklidesowej]
\label{cor:topog-spojnosc-euklidesowe}
Niech $n\ge1$ będzie liczbą całkowitą,
a $U\subseteq\R^n$ zbiorem otwartym.
Wówczas $U$ jest spójny wtedy i tylko wtedy,
gdy jest spójny drogowo.
\end{wniosek}
\begin{proof}
Każdy punkt $x\in U$ ma kulę otwartą
$B(x,r)\subseteq U$, a kula jest wypukła,
więc spójna drogowo na mocy przykładu
\ref{ex:topog-spojnosc-wypukle}.
Wynika stąd lokalna spójność drogowa $U$.
Poprzedni wniosek dowodzi implikacji
ze spójności do spójności drogowej,
a implikacja odwrotna wynika z twierdzenia
\ref{thm:topog-spojnosc-droga-implikuje}.
\end{proof}

\begin{przyklad}[Spójność globalna i lokalna to różne własności]
\label{ex:topog-spojnosc-lokalna-przyklady}
Przestrzeń $\R$ jest lokalnie spójna drogowo:
kule otwarte w $\R$ są przedziałami,
a więc są spójne drogowo. Nie jest to
jednak własność wymuszająca spójność
całej przestrzeni: zbiór
$(-2,-1)\cup(1,2)\subseteq\R$ jest lokalnie
spójny drogowo, lecz nie jest spójny.

Z kolei $\Q$ nie jest lokalnie spójny.
Każdy jego spójny podzbiór jest
jednopunktowy (przykład
\ref{ex:topog-spojnosc-Q}), a żaden
singleton nie jest otwarty w $\Q$.
Nie istnieje więc otwarta spójna baza
otoczeń punktu. W szczególności
całkowita niespójność nie wyklucza
lokalnej spójności (przestrzeń dyskretna
jest lokalnie spójna), ale nie gwarantuje jej.
\end{przyklad}

\section{Podsumowanie i kontrprzykłady}
\label{sec:topog-spojnosc-podsumowanie}

\begin{uwaga}[Implikacje i ich ścisłość]
\label{rem:topog-spojnosc-hierarchia}
Spójność drogowa implikuje spójność,
ale implikacja odwrotna jest fałszywa
(przykład~\ref{ex:topog-spojnosc-sinus}).
Lokalna spójność drogowa implikuje lokalną
spójność, ale nie implikuje spójności
globalnej (przykład
\ref{ex:topog-spojnosc-lokalna-przyklady}).
Spójność i zwartość są niezależnymi
własnościami: $\R$ jest spójna,
ale niezwarta; zbiór dwupunktowy z
topologią dyskretną jest zwarty,
ale niespójny. Zamknięty przedział
$[0,1]$ jest zarazem zwarty i spójny.
Spójność nie wymaga aksjomatu $T_2$
(przykład~\ref{ex:topog-spojnosc-sierpinski}).
\end{uwaga}

\begin{uwaga}[Dalsze związki]
\label{rem:topog-spojnosc-dalsze}
Spójność lokalna stanie się istotna
przy badaniu rozmaitości i przestrzeni
nakryciowych. Spójność drogowa pozwala
konstruować klasy dróg i wprowadzać
grupę podstawową, co będzie rozwijane
w topologii algebraicznej. W niniejszym
rozdziale używaliśmy wyłącznie zwykłej
topologii podprzestrzeni, iloczynu,
odwzorowań ciągłych i zwartości
z poprzednich rozdziałów.
\end{uwaga}

\chapter{Aksjomaty przeliczalności i własności pokrewne}
\label{ch:topog-przeliczalnosc}

\noindent\textbf{Cel rozdziału.}
Baza topologii może być ogromna nawet wtedy, gdy pojedyncze punkty
mają bardzo proste otoczenia. Rozróżnimy warunek lokalny
(pierwszy aksjomat przeliczalności) od globalnego
(drugi aksjomat przeliczalności),
a następnie porównamy je z ośrodkowością i własnością Lindelöfa.
Podamy wszystkie istotne implikacje, przeanalizujemy ograniczenia
ich odwracania oraz wykażemy, że w przestrzeniach metrycznych
trzy z tych własności są równoważne. Nie zakładamy aksjomatów
oddzielania, o ile nie zostaną wyraźnie wymienione.

\section{Pierwszy i drugi aksjomat przeliczalności}
\label{sec:topog-dwie-przeliczalnosci}

\begin{definicja}[Drugi aksjomat przeliczalności]
\label{def:topog-druga-przeliczalnosc}
Niech $(X,\tau)$ będzie przestrzenią topologiczną.
Mówimy, że $X$ \emph{spełnia drugi aksjomat przeliczalności},
jeżeli istnieje co najwyżej przeliczalna rodzina $\mathcal B\subseteq\tau$
taka, że dla każdego $U\in\tau$ i $x\in U$ istnieje
$B\in\mathcal B$ spełniający $x\in B\subseteq U$.
Innymi słowy, $\mathcal B$ jest przeliczalną bazą \emph{całej}
topologii. Dla $X=\varnothing$ można wziąć $\mathcal B=\varnothing$.
\end{definicja}

\begin{uwaga}[Warunek lokalny a globalny]
\label{rem:topog-lokalnie-globalnie-przeliczalnosc}
Pierwszy aksjomat przeliczalności zdefiniowaliśmy w rozdziale~06:
każdy punkt ma co najwyżej przeliczalną bazę otoczeń.
Rodzina tych baz może jednak zmieniać się z punktem i nie musi
istnieć jedna przeliczalna rodzina wystarczająca dla całej przestrzeni.
Bazę otoczeń rozumiemy jak w rozdziale~01: jej elementy są
otoczeniami punktu, niekoniecznie zbiorami otwartymi.
\end{uwaga}

\begin{twierdzenie}[Drugi aksjomat przeliczalności implikuje pierwszy]
\label{thm:topog-druga-pierwsza}
Każda przestrzeń spełniająca drugi aksjomat przeliczalności
spełnia pierwszy aksjomat przeliczalności.
\end{twierdzenie}
\begin{proof}
Niech $\mathcal B$ będzie co najwyżej przeliczalną bazą $X$.
Dla każdego $x\in X$ rodzina
$\mathcal B_x=\{B\in\mathcal B:x\in B\}$
jest co najwyżej przeliczalna. Jeśli $U$ jest dowolnym
otwartym otoczeniem $x$, z definicji bazy istnieje
$B\in\mathcal B_x$ z $B\subseteq U$. Jeżeli $N$ jest
otoczeniem niekoniecznie otwartym, zawiera takie $U$,
więc pewien $B\in\mathcal B_x$ zawiera $x$ i jest zawarty w $N$.
Zatem $\mathcal B_x$ jest bazą otoczeń $x$.
Dla przestrzeni pustej warunek jest spełniony automatycznie.
\end{proof}

\begin{przyklad}[Baza euklidesowa o wymiernych parametrach]
\label{ex:topog-przeliczalna-baza-euklidesowa}
Dla $n\geq1$ rodzina wszystkich kul otwartych
$B(q,r)\subseteq\R^n$ o środku $q\in\Q^n$
i promieniu $r\in\Q$, $r>0$, jest przeliczalna.
Jest bazą topologii euklidesowej: gdy $x\in U$ i $U$ jest otwarty,
istnieje $\varepsilon>0$ z $B(x,\varepsilon)\subseteq U$.
Wybieramy $q\in\Q^n$ tak blisko $x$, aby
$\|x-q\|<\varepsilon/4$, a następnie
$r\in\Q$ takie, że
$\|x-q\|<r<\varepsilon-\|x-q\|$.
Wówczas $x\in B(q,r)$, a z nierówności trójkąta
$B(q,r)\subseteq B(x,\varepsilon)\subseteq U$.
Zatem każda $\R^n$ spełnia drugi aksjomat przeliczalności.
\end{przyklad}

\begin{twierdzenie}[Dziedziczenie przeliczalności]
\label{thm:topog-przeliczalnosc-dziedziczenie}
Niech $A$ będzie podprzestrzenią przestrzeni topologicznej $X$.
Jeżeli $X$ spełnia pierwszy aksjomat przeliczalności, to $A$
również go spełnia. Jeżeli $X$ spełnia drugi aksjomat
przeliczalności, to $A$ również go spełnia.
\end{twierdzenie}
\begin{proof}
Ustalmy $a\in A$. Jeśli $\mathcal N_a$ jest przeliczalną
bazą otoczeń $a$ w $X$, to
$\{N\cap A:N\in\mathcal N_a\}$
jest bazą otoczeń $a$ w $A$. Istotnie, względne otoczenie
$a$ zawiera zbiór postaci $U\cap A$ dla pewnego otwartego
$U\subseteq X$ z $a\in U$; istnieje $N\in\mathcal N_a$
takie, że $a\in N\subseteq U$, więc
$a\in N\cap A\subseteq U\cap A$.
Jeśli $\mathcal B$ jest przeliczalną bazą $X$, to
$\mathcal B|_A=\{B\cap A:B\in\mathcal B\}$
jest przeliczalną bazą topologii podprzestrzeni:
każdy względnie otwarty zbiór ma postać $U\cap A$ i jest
sumą tych $B\cap A$, dla których $B\subseteq U$.
Oba stwierdzenia obejmują również $A=\varnothing$.
\end{proof}

\begin{twierdzenie}[Przeliczalne iloczyny przestrzeni spełniających drugi aksjomat przeliczalności]
\label{thm:topog-druga-iloczyn-przeliczalny}
Niech $I$ będzie zbiorem co najwyżej przeliczalnym,
a $\{X_i:i\in I\}$ rodziną przestrzeni spełniających
drugi aksjomat przeliczalności.
Wówczas iloczyn $\prod_{i\in I}X_i$, wyposażony w topologię
iloczynu, spełnia drugi aksjomat przeliczalności.
\end{twierdzenie}
\begin{proof}
Dla każdego $i\in I$ wybieramy co najwyżej przeliczalną bazę
$\mathcal B_i$ przestrzeni $X_i$. Dla skończonego $F\subseteq I$
i zbiorów $B_i\in\mathcal B_i$ rozważamy cylinder
\[
 C(F,(B_i))=\bigcap_{i\in F}\pi_i^{-1}(B_i).
\]
Dla $F=\varnothing$ jest to cały iloczyn. Zbiory te tworzą
bazę topologii iloczynu, ponieważ każdy podstawowy otwarty
cylinder jest takiego rodzaju po ewentualnym zastąpieniu
otwartych zbiorów w czynnikach elementami odpowiednich baz.
Rodzina skończonych podzbiorów przeliczalnego $I$ jest
przeliczalna: jest sumą po $m\geq0$ obrazów zbiorów
$I^m$ przez odwzorowanie $(i_1,\ldots,i_m)\mapsto
\{i_1,\ldots,i_m\}$. Dla ustalonego skończonego $F$
rodzina kombinacji $(B_i)_{i\in F}$ również jest
co najwyżej przeliczalna, jako skończony iloczyn rodzin
przeliczalnych. Suma przeliczalnie wielu takich rodzin jest przeliczalna.
W przypadku pustego czynnika iloczyn jest pusty,
a dla $I=\varnothing$ jest jednoelementowy; w obu sytuacjach
wniosek zachodzi bezpośrednio.
\end{proof}

\begin{przyklad}[Nieprzeliczalny iloczyn nie musi spełniać pierwszego aksjomatu przeliczalności]
\label{ex:topog-iloczyn-nieprzeliczalny}
Niech $I$ będzie nieprzeliczalny. W przestrzeni
$X=\R^I$ z topologią iloczynu oznaczmy przez $0$
punkt o wszystkich współrzędnych równych zero.
Pokażemy, że $X$ nie spełnia pierwszego aksjomatu
przeliczalności w punkcie $0$.
Załóżmy przeciwnie, że $\{U_n:n\geq1\}$ jest przeliczalną
bazą otoczeń $0$. Dla każdego $n$ istnieje podstawowy
otwarty cylinder $C_n$ z
$0\in C_n\subseteq U_n$, ograniczający tylko skończenie
wiele współrzędnych $F_n\subseteq I$. Zbiór
$F=\bigcup_{n\geq1}F_n$ jest przeliczalny, więc istnieje
$i_0\in I\setminus F$. Zbiór
$V=\pi_{i_0}^{-1}((-1,1))$ jest otwartym otoczeniem $0$.
Dla każdego $n$ cylinder $C_n$ dopuszcza dowolną wartość
współrzędnej $i_0$, ponieważ $i_0\notin F_n$.
Stąd $C_n\not\subseteq V$ i tym bardziej
$U_n\not\subseteq V$, sprzecznie z własnością bazy
otoczeń. Jest to przykład zawodności pierwszego aksjomatu
przeliczalności w dowolnych iloczynach, mimo że każdy
czynnik spełnia nawet drugi aksjomat przeliczalności.
\end{przyklad}

\section{Ośrodkowość i własność Lindelöfa}
\label{sec:topog-osrodkowosc-lindelof}

\begin{definicja}[Ośrodkowość]
\label{def:topog-osrodkowosc}
Przestrzeń topologiczna $X$ jest \emph{ośrodkowa}, jeżeli
istnieje co najwyżej przeliczalny zbiór $D\subseteq X$
gęsty w $X$, to znaczy $\overline D=X$.
Dla przestrzeni pustej można przyjąć $D=\varnothing$.
\end{definicja}

\begin{definicja}[Własność Lindelöfa]
\label{def:topog-lindelof}
Przestrzeń topologiczna $X$ jest \emph{przestrzenią Lindelöfa},
jeśli z każdego pokrycia $X$ zbiorami otwartymi można
wybrać co najwyżej przeliczalne podpokrycie.
Przestrzeń pusta spełnia ten warunek.
\end{definicja}

\begin{twierdzenie}[Dwie konsekwencje drugiego aksjomatu przeliczalności]
\label{thm:topog-druga-osrodkowa-lindelof}
Jeśli przestrzeń topologiczna $X$ spełnia drugi aksjomat
przeliczalności,
to jest ośrodkowa i jest przestrzenią Lindelöfa.
\end{twierdzenie}
\begin{proof}
Niech $\mathcal B$ będzie co najwyżej przeliczalną bazą $X$.
Dla każdego niepustego $B\in\mathcal B$ wybieramy punkt $x_B\in B$.
Zbiór $D=\{x_B:B\in\mathcal B,\ B\ne\varnothing\}$
jest co najwyżej przeliczalny. Jeśli $U\subseteq X$ jest
niepusty i otwarty, zawiera pewien niepusty
$B\in\mathcal B$, a zatem $x_B\in D\cap U$.
Każdy niepusty zbiór otwarty przecina $D$,
co jest równoważne równości $\overline D=X$.

Niech teraz $\mathcal U$ będzie pokryciem otwartym $X$.
Rozważmy rodzinę
$\mathcal B'=\{B\in\mathcal B:\exists U\in\mathcal U,
\ B\subseteq U\}$. Dla każdego $B\in\mathcal B'$
wybieramy jeden zbiór $U_B\in\mathcal U$ z $B\subseteq U_B$.
Rodzina $\{U_B:B\in\mathcal B'\}$ jest co najwyżej
przeliczalna. Pokrywa $X$: dla każdego $x\in X$
istnieje $U\in\mathcal U$ z $x\in U$, a z własności bazy
istnieje $B\in\mathcal B$ z $x\in B\subseteq U$;
wtedy $B\in\mathcal B'$ i $x\in U_B$.
Oba argumenty są poprawne również dla $X=\varnothing$.
\end{proof}

\begin{uwaga}[Aksjomat wyboru]
\label{rem:topog-wybor-przeliczalnosc}
Wybór punktu $x_B$ albo zbioru $U_B$ dla wszystkich
elementów rodziny przeliczalnej korzysta co najwyżej
z przeliczalnej postaci aksjomatu wyboru; w tym skrypcie
pracujemy w zwykłym systemie ZFC. Samo twierdzenie
Tychonoffa z rozdziału~07 wymaga w ogólności silniejszego
narzędzia. Nie utożsamiamy zatem powyższych argumentów
z dowodami zupełnie pozbawionymi wyboru.
\end{uwaga}

\begin{twierdzenie}[Zachowanie ośrodkowości i własności Lindelöfa]
\label{thm:topog-osrodk-lindelof-zachowanie}
Niech $X,Y$ będą przestrzeniami topologicznymi.
\begin{enumerate}[label=(\roman*)]
\item Jeśli $X$ jest ośrodkowa, to każda niepusta otwarta
podprzestrzeń $U\subseteq X$ jest ośrodkowa.
\item Jeśli $f:X\to Y$ jest ciągła i surjektywna, to
ośrodkowość $X$ implikuje ośrodkowość $Y$, a własność
Lindelöfa przestrzeni $X$ implikuje własność Lindelöfa
przestrzeni $Y$.
\item Jeśli $X$ jest przestrzenią Lindelöfa,
to każda domknięta podprzestrzeń $F\subseteq X$ jest
przestrzenią Lindelöfa.
\item Jeśli $X=\bigcup_{n\geq1}X_n$, gdzie każde $X_n$
(jako podprzestrzeń $X$) ma własność Lindelöfa,
to $X$ również ma własność Lindelöfa.
\end{enumerate}
\end{twierdzenie}
\begin{proof}
(i) Niech $D\subseteq X$ będzie przeliczalny i gęsty.
Dla niepustego zbioru otwartego $W$ w $U$ mamy
$W=V\cap U$ dla pewnego otwartego $V\subseteq X$.
Ponieważ $U$ jest otwarty, $W$ jest niepustym zbiorem
otwartym w $X$, a więc $W\cap D\ne\varnothing$.
Zatem $D\cap U$ jest przeliczalny i gęsty w $U$.
Przypadek $U=\varnothing$ jest równie prosty.

(ii) Jeśli $D$ jest gęsty w $X$, to dla każdego niepustego
otwartego $V\subseteq Y$ zbiór $f^{-1}(V)$ jest niepusty
(z surjektywności) i otwarty (z ciągłości). Przecina więc $D$,
a zatem $V$ przecina $f(D)$; zbiór $f(D)$ jest
przeliczalny i gęsty w $Y$. Jeśli $\mathcal V$ jest
pokryciem otwartym $Y$, to
$\{f^{-1}(V):V\in\mathcal V\}$ jest pokryciem otwartym
$X$. Wybieramy przeliczalne podpokrycie
$\{f^{-1}(V_n)\}$; wtedy $\{V_n\}$ pokrywa $Y$
z surjektywności $f$.

(iii) Niech $\mathcal W$ będzie pokryciem $F$ zbiorami
otwartymi w topologii podprzestrzeni. Każde $W\in\mathcal W$
ma postać $W=U_W\cap F$, gdzie $U_W$ jest otwarty w $X$.
Rodzina $\{U_W:W\in\mathcal W\}\cup\{X\setminus F\}$
jest otwartym pokryciem $X$, bo $F$ jest domknięty.
Istnieje jej przeliczalne podpokrycie; po pominięciu
$X\setminus F$ i przecięciu pozostałych zbiorów z $F$
otrzymujemy przeliczalne podpokrycie $F$.

(iv) Niech $\mathcal U$ będzie otwartym pokryciem $X$.
Dla każdego $n$ rodzina
$\{U\cap X_n:U\in\mathcal U\}$ pokrywa $X_n$.
Wybieramy przeliczalną podrodzinę $\mathcal U_n\subseteq\mathcal U$
pokrywającą $X_n$. Suma mnogościowa
$\bigcup_{n\geq1}\mathcal U_n$ jest przeliczalna
i pokrywa $\bigcup_nX_n=X$.
\end{proof}

\begin{przyklad}[Przestrzeń ośrodkowa nie musi spełniać pierwszego aksjomatu przeliczalności]
\label{ex:topog-kosk-onczona-nie-pierwsza}
Niech $X$ będzie nieprzeliczalnym zbiorem z topologią
koskończoną. Dowolny nieskończony przeliczalny
$D\subseteq X$ jest gęsty: każdy niepusty zbiór otwarty
ma skończone dopełnienie, więc przecina $D$.
Przestrzeń jest także zwarta: z danego pokrycia otwartego
wybieramy jeden niepusty zbiór otwarty $U$, a pozostały
skończony zbiór $X\setminus U$ pokrywamy skończenie
wieloma elementami pokrycia.
Jednak $X$ nie spełnia pierwszego aksjomatu przeliczalności.
Istotnie, jeśli $\{N_n:n\geq1\}$ byłaby bazą otoczeń
punktu $x$, wybierzmy dla każdego $n$ otwarty zbiór
$U_n$ z $x\in U_n\subseteq N_n$. Zbiór
$F=\bigcup_n(X\setminus U_n)$ jest przeliczalny.
Wybierzmy $y\in X\setminus(F\cup\{x\})$. Wtedy
otwarte otoczenie $X\setminus\{y\}$ punktu $x$
nie zawiera żadnego $N_n$, gdyż
$y\in U_n\subseteq N_n$ dla każdego $n$.
Otrzymujemy sprzeczność z definicją bazy otoczeń.
W szczególności ośrodkowość, a nawet zwartość i aksjomat
$T_1$, nie zapewniają spełnienia pierwszego aksjomatu przeliczalności.
\end{przyklad}

\section{Przestrzenie metryczne: trzy równoważne własności}
\label{sec:topog-przeliczalnosc-metryczna}

\begin{twierdzenie}[Ośrodkowość, drugi aksjomat przeliczalności i własność Lindelöfa]
\label{thm:topog-metryczne-trzy-warunki}
Niech $(X,d)$ będzie przestrzenią metryczną.
Następujące warunki są równoważne:
\begin{enumerate}[label=(\roman*)]
\item $X$ jest ośrodkowa;
\item $X$ spełnia drugi aksjomat przeliczalności;
\item $X$ jest przestrzenią Lindelöfa.
\end{enumerate}
\end{twierdzenie}
\begin{proof}
Implikacje (ii)$\Rightarrow$(i) i (ii)$\Rightarrow$(iii)
wynikają z twierdzenia~\ref{thm:topog-druga-osrodkowa-lindelof}.

(i)$\Rightarrow$(ii). Niech $D\subseteq X$ będzie
co najwyżej przeliczalny i gęsty. Rodzina
\[
 \mathcal B=\{B_d(a,r):a\in D,\ r\in\Q,\ r>0\}
\]
jest co najwyżej przeliczalna. Wykażemy, że jest bazą.
Niech $x\in U$, gdzie $U$ jest otwarty. Wybierzmy
$\varepsilon>0$ z $B_d(x,\varepsilon)\subseteq U$.
Z gęstości $D$ istnieje $a\in D$ z
$d(x,a)<\varepsilon/4$. Wybieramy liczbę wymierną
$r$ spełniającą
$d(x,a)<r<\varepsilon-d(x,a)$; jest to możliwe,
ponieważ $2d(x,a)<\varepsilon/2<\varepsilon$.
Wtedy $x\in B_d(a,r)$, a jeżeli $y\in B_d(a,r)$,
to z nierówności trójkąta
$d(x,y)\leq d(x,a)+d(a,y)<d(x,a)+r<\varepsilon$.
Zatem $B_d(a,r)\subseteq U$.

(iii)$\Rightarrow$(i). Dla każdego $n\geq1$ rodzina
$\{B_d(x,1/n):x\in X\}$ jest otwartym pokryciem $X$.
Z własności Lindelöfa wybieramy co najwyżej przeliczalne
podpokrycie o środkach ze zbioru $D_n\subseteq X$.
Niech $D=\bigcup_{n\geq1}D_n$. Zbiór $D$ jest
co najwyżej przeliczalny. Dla dowolnego niepustego
zbioru otwartego $U$ wybieramy $x\in U$ i
$\varepsilon>0$ z $B_d(x,\varepsilon)\subseteq U$.
Bierzemy $n$ tak duże, że $1/n<\varepsilon$.
Ponieważ kule o środkach z $D_n$ pokrywają $X$,
istnieje $a\in D_n$ z $d(x,a)<1/n<\varepsilon$.
Wtedy $a\in D\cap U$. Zatem $D$ jest gęsty.
Jeśli $X=\varnothing$, wszystkie trzy własności
zachodzą bezpośrednio; argument można zastosować do
$X\ne\varnothing$.
\end{proof}

\begin{wniosek}[Przykłady metryczne]
\label{cor:topog-osrodkowa-metryczna-baza}
Dla każdego $n\geq1$ przestrzeń $\R^n$ z topologią
euklidesową spełnia drugi aksjomat przeliczalności,
jest ośrodkowa i ma
własność Lindelöfa. Każda jej podprzestrzeń ma te
trzy własności.
\end{wniosek}
\begin{proof}
Rodzina $\Q^n$ jest przeliczalna i gęsta w $\R^n$,
więc stosujemy twierdzenie
\ref{thm:topog-metryczne-trzy-warunki}. Dowolna podprzestrzeń
$\R^n$ jest metryczna, a drugi aksjomat przeliczalności
przechodzi na podprzestrzenie na mocy twierdzenia
\ref{thm:topog-przeliczalnosc-dziedziczenie}.
Ponownie stosując równoważność, otrzymujemy ośrodkowość
i własność Lindelöfa podprzestrzeni.
\end{proof}

\section{Prosta Sorgenfreya: ostrzeżenie przed odwracaniem implikacji}
\label{sec:topog-sorgenfrey-przeliczalnosc}

\begin{definicja}[Prosta Sorgenfreya]
\label{def:topog-sorgenfrey-linia}
Na $\R$ rozważamy topologię mającą za bazę wszystkie
półotwarte przedziały $[a,b)$, gdzie $a<b$.
Otrzymaną przestrzeń oznaczamy przez $S$ i nazywamy
\emph{prostą Sorgenfreya}. Baza ta rzeczywiście wyznacza
topologię: każdy punkt należy do pewnego takiego przedziału,
a jeśli $x\in[a,b)\cap[c,d)$, to
$[x,\min\{b,d\})$ jest zawarty w przecięciu.
\end{definicja}

\begin{figure}[htbp]
\centering
\begin{tikzpicture}[x=1cm,y=1cm]
\draw[-{Latex[length=2mm]},gray!75] (-3.4,.9)--(3.4,.9);
\draw[line width=2.1pt,blue!65!black] (-1.6,.9)--(1.65,.9);
\fill[blue!65!black] (-1.6,.9) circle (2.9pt);
\draw[blue!65!black,fill=white,line width=.95pt] (1.65,.9) circle (2.9pt);
\node[above=3mm] at (0,.9) {$[a,b)$};
\node[below=2mm] at (-1.6,.9) {$a$};
\node[below=2mm] at (1.65,.9) {$b$};
\draw[-{Latex[length=2mm]},gray!75] (-3.4,-.45)--(3.4,-.45);
\draw[line width=2.1pt,orange!80!black] (-1.6,-.45)--(1.65,-.45);
\draw[orange!80!black,fill=white,line width=.95pt] (-1.6,-.45) circle (2.9pt);
\draw[orange!80!black,fill=white,line width=.95pt] (1.65,-.45) circle (2.9pt);
\node[above=3mm] at (0,-.45) {$(a,b)$};
\node[below=2mm] at (-1.6,-.45) {$a$};
\node[below=2mm] at (1.65,-.45) {$b$};
\end{tikzpicture}
\caption{Przedział bazowy topologii Sorgenfreya ma domknięty lewy
koniec i otwarty prawy koniec (u góry). W topologii euklidesowej
przedział bazowy jest otwarty z obu stron (u dołu).}
\label{fig:topog-sorgenfrey-porownanie}
\end{figure}

\begin{twierdzenie}[Cztery własności prostej Sorgenfreya]
\label{thm:topog-sorgenfrey-wlasnosci-przeliczalnosc}
Prosta Sorgenfreya $S$ spełnia pierwszy aksjomat
przeliczalności, jest ośrodkowa i Lindelöfa,
ale nie spełnia drugiego aksjomatu przeliczalności.
\end{twierdzenie}
\begin{proof}
\emph{Pierwszy aksjomat przeliczalności.} Dla każdego $x\in\R$ rodzina
$\{[x,x+1/n):n\geq1\}$ jest przeliczalną bazą otoczeń $x$.
Rzeczywiście, każde otwarte otoczenie $x$ zawiera przedział
bazowy $[a,b)$ z $a\leq x<b$, a wówczas dla dostatecznie
dużego $n$ zachodzi $[x,x+1/n)\subseteq[a,b)$.

\emph{Ośrodkowość.} Zbiór $\Q$ jest przeliczalny i gęsty,
ponieważ każdy niepusty przedział $[a,b)$ zawiera
liczbę wymierną z $(a,b)$.

\emph{Niespełnianie drugiego aksjomatu przeliczalności.} Gdyby
$\mathcal B=\{B_n:n\geq1\}$ była przeliczalną bazą,
dla każdego $x\in\R$ istniałby najmniejszy indeks
\[
 n(x)=\min\{n:x\in B_n\subseteq[x,x+1)\}.
\]
Indeksy $n(x)$ dla różnych $x$ byłyby różne. Istotnie,
przy $x<y$ równość $B_{n(x)}=B_{n(y)}$ byłaby
niemożliwa, bo ten sam zbiór musiałby zawierać $x$,
a jednocześnie być zawarty w $[y,y+1)$,
który nie zawiera $x$. W ten sposób otrzymalibyśmy
iniekcję $\R\to\mathbb N$, sprzeczną z
nieprzeliczalnością $\R$.

\emph{Własność Lindelöfa.} Niech $\mathcal U$ będzie
dowolnym otwartym pokryciem $S$. Rozważmy rodzinę
\[
 \mathcal I=\{[a,b):a<b,\ [a,b)\subseteq U
 \text{ dla pewnego }U\in\mathcal U\}.
\]
Jest to pokrycie $\R$ przedziałami bazowymi.
Zbiór $G=\bigcup_{[a,b)\in\mathcal I}(a,b)$ jest
otwarty w zwykłej topologii euklidesowej. Ponieważ
$\R$ z tą topologią spełnia drugi aksjomat
przeliczalności, a jej podprzestrzeń $G$ również go spełnia
(twierdzenie~\ref{thm:topog-przeliczalnosc-dziedziczenie}),
a więc jest przestrzenią Lindelöfa. Z pokrycia $G$ przedziałami $(a,b)$,
gdzie $[a,b)\in\mathcal I$, możemy wybrać
przeliczalną podrodzinę $(a_n,b_n)$ pokrywającą $G$.
Dla każdego $n$ wybieramy $U_n\in\mathcal U$
z $[a_n,b_n)\subseteq U_n$. Rodzina $\{U_n\}$ pokrywa $G$.

Pozostaje zbiór $F=\R\setminus G$. Dla każdego $x\in F$
istnieje $[x,b)\in\mathcal I$ dla pewnego $b>x$.
Istotnie, dowolny przedział $[a,b)\in\mathcal I$
zawierający $x$ musi mieć $a=x$, gdyż w przypadku
$a<x<b$ mielibyśmy $x\in(a,b)\subseteq G$.
Ustalmy raz na zawsze numerację $\Q=\{q_1,q_2,\ldots\}$.
Dla każdego $x\in F$ istnieje wymierne $q>x$ takie,
że $[x,q)\in\mathcal I$: wystarczy wziąć
$q\in\Q\cap(x,b)$ i zauważyć, że $[x,q)\subseteq[x,b)$.
Niech $q_x$ będzie pierwszą taką liczbą wymierną
w ustalonej numeracji. Wówczas $(x,q_x)\subseteq G$.
Jeśli $x,y\in F$ oraz $x<y$, to $y\notin(x,q_x)$,
a zatem $q_x\leq y<q_y$. Stąd odwzorowanie
$x\mapsto q_x$ jest iniekcją $F\to\Q$, więc $F$
jest co najwyżej przeliczalny. Dla każdego $x\in F$
wybieramy $U_x\in\mathcal U$ zawierający $[x,q_x)$,
co jest możliwe z definicji $\mathcal I$.
Przeliczalna rodzina $\{U_n\}\cup\{U_x:x\in F\}$
pokrywa zarówno $G$, jak i $F$, a zatem całe $S$.
Dowodzi to własności Lindelöfa.
\end{proof}

\begin{wniosek}[Trzy warunki nie zapewniają drugiego aksjomatu przeliczalności]
\label{cor:topog-sorgenfrey-trzy-nie-druga}
Pierwszy aksjomat przeliczalności, ośrodkowość
i własność Lindelöfa, nawet przy założeniu Hausdorffa,
nie implikują łącznie drugiego aksjomatu przeliczalności.
\end{wniosek}
\begin{proof}
Prosta Sorgenfreya spełnia trzy wymienione warunki,
lecz nie spełnia drugiego aksjomatu przeliczalności.
Jest Hausdorffa:
jeśli $x<y$, to przedziały $[x,y)$ i $[y,y+1)$
są rozłącznymi otwartymi otoczeniami odpowiednio $x$ i $y$.
\end{proof}

\begin{przyklad}[Iloczyn przestrzeni Lindelöfa nie musi być przestrzenią Lindelöfa]
\label{ex:topog-sorgenfrey-iloczyn-nie-lindelof}
Płaszczyzna Sorgenfreya $S\times S$ jest iloczynem
przestrzeni Lindelöfa, lecz sama nie ma własności Lindelöfa.
Uzasadnimy to na podstawie zbioru
$D=\{(t,-t):t\in\R\}$, który pojawił się również
w dowodzie twierdzenia o płaszczyźnie Sorgenfreya
w rozdziale~05.
Dla $(x,-x)\in D$ podstawowy prostokąt
$[x,x+\varepsilon)\times[-x,-x+\varepsilon)$
przecina $D$ tylko w $(x,-x)$: jeśli $(t,-t)$ należy
do prostokąta, mamy jednocześnie $t\ge x$ i $t\le x$.
Zatem $D$ jest podprzestrzenią dyskretną.
Jest też domknięty. Jeśli $(u,v)\notin D$, to $u+v\ne0$.
Dla $u+v>0$ każdy podstawowy prostokąt wychodzący
z $(u,v)$ składa się z punktów o dodatniej sumie
współrzędnych; dla $u+v<0$ wystarczy wybrać
$\varepsilon>0$ tak małe, aby $u+v+2\varepsilon<0$.
W obu przypadkach otrzymujemy otwarte otoczenie
nieprzecinające $D$.

Nieprzeliczalna przestrzeń dyskretna nie jest przestrzenią
Lindelöfa, ponieważ pokrycie wszystkimi singletonami
nie ma przeliczalnego podpokrycia. Gdyby $S\times S$
była przestrzenią Lindelöfa, jej domknięta podprzestrzeń
$D$ byłaby Lindelöfa na mocy twierdzenia
\ref{thm:topog-osrodk-lindelof-zachowanie}(iii), sprzeczność.
\end{przyklad}

\section{Które własności przechodzą na podprzestrzenie?}
\label{sec:topog-przeliczalnosc-kontrprzyklady}

\begin{przyklad}[Ośrodkowość nie przechodzi na domknięte podprzestrzenie]
\label{ex:topog-osrodkowosc-nie-dziedziczna}
Płaszczyzna Sorgenfreya $S\times S$ jest ośrodkowa:
$\Q\times\Q$ jest przeliczalny i przecina każdy
niepusty podstawowy prostokąt
$[a,b)\times[c,d)$, ponieważ każdy z obu przedziałów
zawiera liczbę wymierną. Jej domknięta podprzestrzeń
$D=\{(t,-t):t\in\R\}$ jest jednak nieprzeliczalna
i dyskretna, jak wykazaliśmy w przykładzie
\ref{ex:topog-sorgenfrey-iloczyn-nie-lindelof}.
W każdej przestrzeni dyskretnej jedynym gęstym podzbiorem
jest cała przestrzeń, ponieważ każdy singleton jest otwarty.
Zatem $D$ nie jest ośrodkowa.
\end{przyklad}

\begin{przyklad}[Przestrzeń zwarta z otwartą podprzestrzenią
niebędącą przestrzenią Lindelöfa]
\label{ex:topog-lindelof-nie-dziedziczna}
Niech $E$ będzie zbiorem nieprzeliczalnym i niech
$X=E\cup\{\infty\}$, gdzie $\infty\notin E$.
Przyjmijmy za otwarte wszystkie podzbiory $E$ oraz
zbiory zawierające $\infty$, których dopełnienie w $E$
jest skończone. Rodzina ta jest topologią:
sumy zbiorów otwartych są otwarte, a przecięcie
skończenie wielu otoczeń $\infty$ nadal ma skończone
dopełnienie w $E$.

Przestrzeń $X$ jest zwarta. Każde pokrycie otwarte ma
element $U$ zawierający $\infty$; pozostała część
$X\setminus U$ jest skończona, więc wystarczy dobrać
do $U$ skończenie wiele dalszych elementów pokrycia.
Jest też Hausdorffa: różne punkty $e,f\in E$
oddzielają singletony, a $e\in E$ i $\infty$
oddzielają zbiory $\{e\}$ oraz $X\setminus\{e\}$.

Podprzestrzeń $E$ jest otwarta i dyskretna, a ponieważ
jest nieprzeliczalna, nie jest przestrzenią Lindelöfa.
Zatem własność Lindelöfa nie przechodzi na dowolne
otwarte podprzestrzenie, mimo że przechodzi na domknięte.
\end{przyklad}

\begin{samepage}
\begin{przyklad}[Pierwszy aksjomat przeliczalności nie zapewnia ośrodkowości]
\label{ex:topog-dyskretna-nie-osrodkowa}
Każda przestrzeń dyskretna spełnia pierwszy aksjomat
przeliczalności:
dla punktu $x$ jednoelementowa rodzina $\{\{x\}\}$
jest bazą jego otoczeń. Jeśli jednak zbiór punktów $E$
jest nieprzeliczalny, topologia dyskretna na $E$ nie jest
ośrodkowa (gęsty jest jedynie cały $E$), nie jest
przestrzenią Lindelöfa (pokrycie singletonami) ani
nie spełnia drugiego aksjomatu przeliczalności
(każda baza musi zawierać
wszystkie singletony).
\end{przyklad}
\end{samepage}

\section{Zestawienie implikacji i ich ograniczeń}
\label{sec:topog-przeliczalnosc-podsumowanie}

\begin{figure}[htbp]
\centering
\begin{tikzpicture}[
 >=Stealth,
 every node/.style={font=\small,align=center},
 main/.style={draw=blue!65!black,fill=blue!6,rounded corners=2pt,
 text width=3.25cm,minimum height=9mm},
 aux/.style={draw=black!55,fill=black!3,rounded corners=2pt,
 text width=3.25cm,minimum height=9mm}]
\node[main] (second) at (0,2.1) {drugi aksjomat\\przeliczalności};
\node[aux] (first) at (-4.1,.1) {pierwszy aksjomat\\przeliczalności};
\node[aux] (separ) at (0,.1) {ośrodkowość};
\node[aux] (lind) at (4.1,.1) {własność Lindelöfa};
\draw[->,thick] (second.south west) -- (first.north east);
\draw[->,thick] (second.south) -- (separ.north);
\draw[->,thick] (second.south east) -- (lind.north west);
\end{tikzpicture}
\caption{Implikacje prawdziwe dla wszystkich przestrzeni topologicznych.
Strzałki w ogólności nie są odwracalne, nawet przy dodatkowych
warunkach podanych w tekście. Równoważność metryczna dotyczy
wyłącznie drugiego aksjomatu przeliczalności, ośrodkowości
i własności Lindelöfa.}
\label{fig:topog-przeliczalnosc-implikacje}
\end{figure}

\begin{stwierdzenie}[Implikacje bez aksjomatów oddzielania]
\label{prop:topog-przeliczalnosc-hierarchia}
Dla dowolnej przestrzeni topologicznej zachodzą implikacje
\[
\text{drugi aksjomat przeliczalności}\ \Longrightarrow\
\begin{cases}
\text{pierwszy aksjomat przeliczalności},\\
\text{ośrodkowość},\\
\text{własność Lindelöfa}.
\end{cases}
\]
Żadna z implikacji odwrotnych nie zachodzi bez dodatkowych
założeń. Ponadto żadne dwa spośród trzech warunków po prawej
nie są, w ogólności, równoważne.
\end{stwierdzenie}
\begin{proof}
Implikacje udowodniono w twierdzeniach
\ref{thm:topog-druga-pierwsza} oraz
\ref{thm:topog-druga-osrodkowa-lindelof}.
Nieprzeliczalna przestrzeń dyskretna spełnia pierwszy
aksjomat przeliczalności, lecz nie jest ośrodkowa ani Lindelöfa.
Nieprzeliczalna przestrzeń koskończona jest ośrodkowa
i Lindelöfa, ale nie spełnia pierwszego aksjomatu przeliczalności.
Prosta Sorgenfreya spełnia pierwszy aksjomat
przeliczalności, jest ośrodkowa i Lindelöfa, lecz
nie spełnia drugiego aksjomatu przeliczalności.
Natomiast płaszczyzna Sorgenfreya jest ośrodkowa
i spełnia pierwszy aksjomat przeliczalności, lecz nie jest Lindelöfa.
Z drugiej strony, zwarta przestrzeń $E\cup\{\infty\}$
z przykładu~\ref{ex:topog-lindelof-nie-dziedziczna}
nie jest ośrodkowa, ponieważ jej nieprzeliczalna otwarta
podprzestrzeń dyskretna $E$ nie może zawierać gęstego
przeliczalnego zbioru: dowolny gęsty podzbiór całego
$E\cup\{\infty\}$ musi zawierać wszystkie punkty izolowane $E$.
Przykłady te weryfikują w szczególności, że wszystkie
odwrócenia trzech głównych implikacji zawodzą,
a warunki po prawej nie są parami równoważne.
\end{proof}

\begin{uwaga}[Co następuje dalej?]
\label{rem:topog-przeliczalnosc-metryzowalnosc}
W przestrzeniach metrycznych pierwszy aksjomat przeliczalności
jest zawsze spełniony, a trzy pozostałe własności są równoważne.
W ogólnych przestrzeniach rozdział ten pokazuje już,
że nawet łączne spełnienie pierwszego aksjomatu przeliczalności,
ośrodkowości i własności Lindelöfa nie wystarcza do
metryzowalności. W szczególności prosta Sorgenfreya
nie jest metryzowalna: gdyby jej topologię wyznaczała
metryka, z ośrodkowości i twierdzenia
\ref{thm:topog-metryczne-trzy-warunki}
wynikałoby spełnienie drugiego aksjomatu przeliczalności,
co jest niemożliwe.
Pełniejsze kryteria metryzowalności należą do dalszej
części skryptu.
\end{uwaga}
\FloatBarrier

\chapter{Metryzowalność przestrzeni topologicznych}
\label{ch:topog-metryzowalnosc}

\noindent\textbf{Cel rozdziału.}
Topologia przestrzeni metrycznej jest określona przez odległość,
lecz nie każda przestrzeń topologiczna posiada zgodną z topologią
metrykę. Ustalimy, jakie warunki są konieczne, a następnie
udowodnimy \emph{twierdzenie Urysohna o metryzowalności}:
przestrzeń regularna, spełniająca $T_1$ i drugi aksjomat
przeliczalności, jest metryzowalna. Nie ograniczymy się do
samego sformułowania: skonstruujemy rozdzielającą rodzinę funkcji
ciągłych i wyprowadzimy z niej metrykę. Dalej omówimy
metryzowalność podprzestrzeni i przeliczalnych iloczynów,
a także kontrprzykłady dotyczące iloczynów nieprzeliczalnych
i ilorazów.

\section{Metryki, pseudometryki i topologia wyznaczona przez odległość}
\label{sec:topog-metryki-definicje}

\begin{definicja}[Pseudometryka i metryka]
\label{def:topog-pseudometryka-metryka}
Niech $X$ będzie zbiorem. \emph{Pseudometryką} na $X$
nazywamy odwzorowanie $d\colon X\times X\to[0,\infty)$,
które dla wszystkich $x,y,z\in X$ spełnia
\begin{enumerate}[label=(\roman*)]
\item $d(x,x)=0$;
\item $d(x,y)=d(y,x)$;
\item $d(x,z)\le d(x,y)+d(y,z)$.
\end{enumerate}
Jeżeli ponadto $d(x,y)=0$ implikuje $x=y$, to $d$ jest
\emph{metryką}. Dla $x\in X$ i $r>0$ definiujemy kulę
$B_d(x,r)=\{y\in X:d(x,y)<r\}$.
Topologię generowaną przez wszystkie takie kule oznaczamy
$\tau_d$.
\end{definicja}

\begin{stwierdzenie}[Kule tworzą bazę]
\label{prop:topog-baza-kul}
Dla każdej pseudometryki $d$ na zbiorze $X$ rodzina
$\{B_d(x,r):x\in X,\ r>0\}$ jest bazą pewnej topologii
$\tau_d$. Dla każdego $x\in X$ rodzina
$\{B_d(x,1/n):n\in\mathbb N,\ n\ge1\}$ jest bazą otwartych
otoczeń $x$ w $\tau_d$.
\end{stwierdzenie}
\begin{proof}
Każdy punkt $x$ należy do $B_d(x,1)$, więc kule pokrywają
$X$; dla $X=\varnothing$ warunek jest pusty. Niech
$z\in B_d(x,r)\cap B_d(y,s)$. Liczba
$\delta=\min\{r-d(x,z),s-d(y,z)\}$ jest dodatnia.
Dla $w\in B_d(z,\delta)$ nierówność trójkąta daje
$d(x,w)\le d(x,z)+d(z,w)<r$ i analogicznie $d(y,w)<s$.
Stąd $B_d(z,\delta)$ mieści się w przecięciu kul, co jest
kryterium bazy topologii z rozdziału~01.

Jeśli $U\in\tau_d$ i $x\in U$, to istnieje $r>0$ z
$B_d(x,r)\subseteq U$: z bazowości mamy wprawdzie początkowo
$x\in B_d(z,s)\subseteq U$, lecz kula
$B_d(x,s-d(z,x))$ leży w $B_d(z,s)$. Wybierając
$n$ z $1/n<r$, otrzymujemy $B_d(x,1/n)\subseteq U$.
\end{proof}

\begin{definicja}[Metryzowalność]
\label{def:topog-metryzowalnosc}
Przestrzeń topologiczna $(X,\tau)$ jest \emph{metryzowalna},
jeżeli istnieje metryka $d$ na $X$, dla której $\tau=\tau_d$.
Taką metrykę nazywamy \emph{zgodną z topologią}.
Pojęcie nie odnosi się do konkretnego wzoru na odległość,
lecz do istnienia przynajmniej jednej zgodnej metryki.
\end{definicja}

\begin{przyklad}[Różnica między metryką a pseudometryką]
\label{ex:topog-pseudometryka-r2}
Na $\R^2$ funkcja
$p((x,y),(u,v))=|x-u|$ jest pseudometryką: nieujemność,
symetria i nierówność trójkąta wynikają z analogicznych
własności wartości bezwzględnej. Nie jest metryką, ponieważ
$p((0,0),(0,1))=0$. Każda kula $p$ jest pionowym pasem
$\{(u,v):|u-x|<r\}$, więc punkty o jednakowej pierwszej
współrzędnej należą dokładnie do tych samych zbiorów otwartych
$\tau_p$. W szczególności $\tau_p$ nie spełnia $T_0$.
Natomiast odległość euklidesowa na $\R^2$ jest metryką
i wyznacza zwykłą topologię euklidesową.
\end{przyklad}

\begin{stwierdzenie}[Warunki konieczne metryzowalności]
\label{prop:topog-metryka-t3-pierwszy}
Każda przestrzeń metryzowalna spełnia aksjomat $T_3$
oraz pierwszy aksjomat przeliczalności.
\end{stwierdzenie}
\begin{proof}
Niech $d$ będzie zgodną metryką. Pierwszy aksjomat
przeliczalności wynika z poprzedniego stwierdzenia o kulach.
Dla $x\ne y$ kula $B_d(x,d(x,y)/2)$ nie zawiera $y$;
ponieważ wszystkie kule są otwarte, singletony są domknięte:
$X\setminus\{y\}=\bigcup_{x\ne y} B_d(x,d(x,y)/2)$.
Otrzymujemy $T_1$.

Aby wykazać regularność, ustalmy domknięty $F\subseteq X$
i $x\notin F$. Zbiór $X\setminus F$ jest otwarty,
więc istnieje $r>0$ takie, że $B_d(x,r)\cap F=\varnothing$.
Pokażemy, że $\overline{B_d(x,r/3)}\subseteq B_d(x,2r/3)$.
Jeśli $y$ należy do lewego zbioru, to dla dowolnego
$\varepsilon>0$ istnieje $z\in B_d(x,r/3)$
z $d(y,z)<\varepsilon$; stąd
$d(x,y)<r/3+\varepsilon$.
Przechodząc z $\varepsilon\downarrow0$, otrzymujemy
$d(x,y)\le r/3<2r/3$. Zatem otwarte zbiory
$B_d(x,r/3)$ i $X\setminus\overline{B_d(x,r/3)}$
rozdzielają punkt $x$ i domknięty $F$. Jest to regularność,
a wraz z $T_1$ daje $T_3$.
\end{proof}

\begin{uwaga}[Ograniczenie warunków koniecznych]
Przestrzeń metryzowalna nie musi spełniać drugiego aksjomatu
przeliczalności. Niech $X$ będzie nieprzeliczalnym zbiorem,
a $d$ jego metryką dyskretną: $d(x,y)=0$ dla $x=y$
i $d(x,y)=1$ dla $x\ne y$. Dla każdej bazy $\mathcal A$
topologii $\tau_d$ i każdego $x\in X$, ponieważ
$\{x\}$ jest otwarty, istnieje $A\in\mathcal A$ z
$x\in A\subseteq\{x\}$. Wtedy $A=\{x\}$, zatem
każda baza zawiera wszystkie singletony i jest
nieprzeliczalna.
\end{uwaga}

\section{Metryki zgodne i podprzestrzenie}
\label{sec:topog-metryki-zgodne}

\begin{stwierdzenie}[Metryki o tych samych zbiorach otwartych]
\label{prop:topog-min-metryka}
Niech $d$ będzie metryką na zbiorze $X$ i niech
$\widehat d(x,y)=\min\{1,d(x,y)\}$.
Wtedy $\widehat d$ również jest metryką,
$\tau_{\widehat d}=\tau_d$ oraz $\widehat d\le1$.
\end{stwierdzenie}
\begin{proof}
Dodatniość i symetria są oczywiste. Dla $a,b\ge0$
zachodzi
$\min\{1,a+b\}\le\min\{1,a\}+\min\{1,b\}$:
jeżeli $a\ge1$ lub $b\ge1$, prawa strona jest co najmniej
$1$; w przeciwnym razie prawa strona wynosi $a+b$.
Stosując tę nierówność i nierówność trójkąta dla $d$,
uzyskujemy ją także dla $\widehat d$.
Dla każdego $0<r\le1$ i każdego $x$ zachodzi
$B_{\widehat d}(x,r)=B_d(x,r)$.
Kule o promieniach $1/n$, $n\ge2$, są bazami otoczeń
punktów w obu topologiach, więc topologie są równe.
\end{proof}

\begin{twierdzenie}[Dziedziczenie metryzowalności]
\label{thm:topog-metryzowalnosc-podprzestrzenie}
Niech $(X,\tau)$ będzie przestrzenią metryzowalną,
a $A\subseteq X$ jej podprzestrzenią. Wtedy topologia
podprzestrzeni na $A$ jest metryzowalna.
\end{twierdzenie}
\begin{proof}
Wybierzmy zgodną metrykę $d$ na $X$ i rozważmy jej
ograniczenie $d_A=d|_{A\times A}$. Wszystkie aksjomaty
metryki przechodzą na podzbiór. Dla $a\in A$ i $r>0$
mamy $B_{d_A}(a,r)=A\cap B_d(a,r)$.
Jeżeli $U\subseteq A$ jest otwarty w topologii podprzestrzeni,
ma postać $A\cap V$ dla otwartego $V\subseteq X$;
każdy jego punkt leży w odpowiedniej kuli $B_{d_A}(a,r)$.
Odwrotnie, każda taka kula jest względnie otwarta.
Zatem $\tau_{d_A}=\tau|_A$.
\end{proof}

\begin{stwierdzenie}[Przenoszenie metryki przez homeomorfizm]
\label{prop:topog-metryzowalnosc-homeo}
Jeżeli $h\colon X\to Y$ jest homeomorfizmem przestrzeni
topologicznych i $Y$ ma zgodną metrykę $d$, to
$d_h(x,x')=d(h(x),h(x'))$ jest metryką zgodną z topologią $X$.
\end{stwierdzenie}
\begin{proof}
Z różnowartościowości $h$ wynika dodatnia określoność,
a pozostałe aksjomaty dziedziczą się z $d$.
Dla $x\in X$ zachodzi
$B_{d_h}(x,r)=h^{-1}(B_d(h(x),r))$.
Ponieważ $h$ jest homeomorfizmem, przeciwobrazy kul
stanowią bazę topologii $X$.
\end{proof}

\section{Od regularności do normalności}
\label{sec:topog-t3-drugi-normalna}

Kluczowy etap konstrukcji metryki polega na otrzymaniu
\emph{przeliczalnie wielu} funkcji ciągłych rozdzielających
punkty i zbiory domknięte. Posłużymy się lematem Urysohna
udowodnionym w rozdziale~05. Wcześniej musimy uzasadnić
normalność, ponieważ regularność sama jej nie gwarantuje.

\begin{lemat}[Regularność i własność Lindelöfa dają normalność]
\label{lem:topog-t3-lindelof-normalna}
Niech $X$ będzie przestrzenią regularną spełniającą
własność Lindelöfa. Wtedy $X$ jest normalna. Jeżeli
ponadto $X$ spełnia $T_1$, to spełnia $T_4$.
\end{lemat}
\begin{proof}
Niech $F,G\subseteq X$ będą rozłącznymi zbiorami
domkniętymi. Jeżeli jeden z nich jest pusty, wystarczy
para $\varnothing,X$ ustawiona we właściwej kolejności.
Załóżmy więc, że oba są niepuste.

Dla każdego $x\in F$ regularność oraz kryterium zwężania
otoczeń z rozdziału~05 zapewniają otwarty $U_x\ni x$ taki,
że $\overline{U_x}\cap G=\varnothing$. Ponieważ zamknięta
podprzestrzeń przestrzeni Lindelöfa również spełnia
własność Lindelöfa (rozdział~09), możemy wybrać
ciąg $U_1,U_2,\ldots$ tych zbiorów pokrywający $F$.
Jeżeli podpokrycie jest skończone, powtarzamy jego elementy.
Analogicznie otrzymujemy otwarte $V_1,V_2,\ldots$ pokrywające
$G$ i spełniające $\overline{V_n}\cap F=\varnothing$.

Definiujemy zbiory otwarte
\[
U=\bigcup_{n\ge1}\left(U_n\setminus\bigcup_{k=1}^n
\overline{V_k}\right),\qquad
V=\bigcup_{n\ge1}\left(V_n\setminus\bigcup_{k=1}^n
\overline{U_k}\right).
\]
Każda różnica występująca po prawej jest otwarta,
bo odejmujemy skończoną sumę zbiorów domkniętych.
Dla $x\in F$ wybieramy $n$ z $x\in U_n$;
ponieważ $F\cap\overline{V_k}=\varnothing$ dla każdego $k$,
$x$ należy do $n$-tego składnika sumy definiującej $U$.
Stąd $F\subseteq U$. Symetrycznie $G\subseteq V$.

Pokażemy, że $U\cap V=\varnothing$. Gdyby $z\in U\cap V$,
istniałyby $n,m$ takie, że
$z\in U_n\setminus\bigcup_{k\le n}\overline{V_k}$
i $z\in V_m\setminus\bigcup_{k\le m}\overline{U_k}$.
Jeżeli $n\le m$, to $z\in U_n\subseteq\overline{U_n}$,
ale drugi warunek wyklucza przynależność do
$\overline{U_n}$. Jeżeli $m\le n$, otrzymujemy
analogiczną sprzeczność, korzystając z
$z\in V_m\subseteq\overline{V_m}$.
Tak więc $F$ i $G$ mają rozłączne otwarte otoczenia.
Otrzymaliśmy normalność, a wraz z $T_1$ — aksjomat $T_4$.
\end{proof}

\begin{wniosek}[Drugi aksjomat przeliczalności i $T_3$]
\label{cor:topog-drugi-t3-t4}
Każda przestrzeń spełniająca $T_3$ i drugi aksjomat
przeliczalności spełnia $T_4$.
\end{wniosek}
\begin{proof}
Z rozdziału~09 wynika, że drugi aksjomat przeliczalności
implikuje własność Lindelöfa. Aksjomat $T_3$ oznacza
regularność wraz z $T_1$. Stosujemy poprzedni lemat.
\end{proof}

\section{Twierdzenie Urysohna o metryzowalności}
\label{sec:topog-urysohn-metryzowalnosc}

\begin{twierdzenie}[Twierdzenie Urysohna o metryzowalności]
\label{thm:topog-urysohn-metryzowalnosc}
Niech $X$ będzie przestrzenią topologiczną spełniającą
aksjomat $T_3$ i drugi aksjomat przeliczalności.
Wówczas istnieje metryka $d$ na $X$ wyznaczająca
dokładnie jego wyjściową topologię.
\end{twierdzenie}
\begin{proof}
\textbf{Przypadki co najwyżej jednoelementowe.}
Jeśli $X=\varnothing$ lub $X$ ma jeden punkt, to
jedyna funkcja $d\colon X\times X\to[0,\infty)$
o wartości $0$ jest metryką, a wyznaczona przez nią
topologia jest jedyną możliwą topologią na $X$.
Dalej zakładamy, że $X$ ma co najmniej dwa punkty.

\textbf{Krok 1: przeliczalna rodzina par otwartych.}
Niech $\mathcal B$ będzie co najwyżej przeliczalną
bazą topologii $X$; można ją uznać za niepustą.
Rozważmy rodzinę
\[
\mathcal P=\{(B,C)\in\mathcal B\times\mathcal B:
\overline B\subseteq C\}.
\]
Rodzina $\mathcal P$ jest co najwyżej przeliczalna,
ponieważ jest podzbiorem iloczynu dwóch zbiorów
co najwyżej przeliczalnych.

Wykażemy, że dla każdego $x\in X$ i każdego
otwartego $O\ni x$ istnieje $(B,C)\in\mathcal P$ takie, że
\[
 x\in B\subseteq\overline B\subseteq C\subseteq O.
\]
Z bazowości $\mathcal B$ wybieramy $C\in\mathcal B$
z $x\in C\subseteq O$. Regularność i kryterium zwężania
otoczeń z rozdziału~05 dają otwarty $W$ z
$x\in W\subseteq\overline W\subseteq C$.
Ponownie z bazowości wybieramy $B\in\mathcal B$
z $x\in B\subseteq W$. Wtedy
$\overline B\subseteq\overline W\subseteq C$,
więc $(B,C)\in\mathcal P$. Ponieważ możemy
przyjąć $O=X\setminus\{y\}$ dla $x\ne y$ (aksjomat
$T_1$), rodzina $\mathcal P$ jest niepusta.

\textbf{Krok 2: funkcje oddzielające.}
Na mocy poprzedniego wniosku $X$ jest $T_4$.
Dla każdej pary $(B,C)\in\mathcal P$ zbiory
$\overline B$ i $X\setminus C$ są domknięte i rozłączne.
Lemat Urysohna z rozdziału~05 dostarcza zatem
funkcji ciągłej
$u_{B,C}\colon X\to[0,1]$ takiej, że
$u_{B,C}=0$ na $\overline B$ i $u_{B,C}=1$
na $X\setminus C$.
Wybieramy takie funkcje dla wszystkich par
z $\mathcal P$. Jest ich co najwyżej przeliczalnie wiele;
możemy ułożyć je w ciąg $(f_n)_{n\ge1}$,
powtarzając elementy, jeśli rodzina jest skończona.
Każda funkcja $f_n$ jest ciągła i przyjmuje wartości w $[0,1]$.

\textbf{Krok 3: konstrukcja i aksjomaty metryki.}
Dla $x,y\in X$ definiujemy
\[
 d(x,y)=\sum_{n=1}^{\infty}2^{-n}|f_n(x)-f_n(y)|.
\]
Szereg jest zbieżny, ponieważ jego wyrazy są
nieujemne i nie większe od $2^{-n}$, a
$\sum_{n\ge1}2^{-n}=1$. Otrzymujemy więc
$0\le d(x,y)\le1$, symetrię i $d(x,x)=0$.
Dla $x,y,z\in X$ nierówność trójkąta dla wartości
bezwzględnej, zastosowana do każdego składnika,
daje dla wszystkich $N$
\[
\sum_{n=1}^{N}2^{-n}|f_n(x)-f_n(z)|
\le\sum_{n=1}^{N}2^{-n}|f_n(x)-f_n(y)|
+\sum_{n=1}^{N}2^{-n}|f_n(y)-f_n(z)|.
\]
Po przejściu do granicy otrzymujemy
$d(x,z)\le d(x,y)+d(y,z)$.

Pozostała dodatnia określoność. Dla $x\ne y$
wybieramy w kroku 1 parę $(B,C)$ z
$x\in B$ i $C\subseteq X\setminus\{y\}$.
Odpowiednia funkcja $u_{B,C}=f_n$ spełnia
$f_n(x)=0$ i $f_n(y)=1$. Wobec tego
$d(x,y)\ge2^{-n}>0$. Zatem $d$ jest metryką.

\textbf{Krok 4: każda kula metryczna jest otwarta w $X$.}
Ustalmy $x\in X$ i $\varepsilon>0$. Wybierzmy $N$
tak duże, aby $\sum_{n>N}2^{-n}=2^{-N}<\varepsilon/2$.
Dla każdego $n\le N$, z ciągłości $f_n$ w punkcie $x$
zbiór
$W_n=\{y\in X:|f_n(y)-f_n(x)|<\varepsilon/2\}$
jest otwarty i zawiera $x$. Zbiór
$W=\bigcap_{n=1}^{N}W_n$ jest otwartym otoczeniem $x$.
Dla $y\in W$ mamy
\[
 d(x,y)\le\sum_{n=1}^{N}2^{-n}|f_n(x)-f_n(y)|
+\sum_{n>N}2^{-n}<\varepsilon/2+\varepsilon/2
=\varepsilon.
\]
A zatem $x\in W\subseteq B_d(x,\varepsilon)$.
Wykażmy teraz otwartość całej kuli, a nie tylko istnienie
otwartego otoczenia jej środka. Jeśli
$y\in B_d(x,\varepsilon)$, to
$\delta=\varepsilon-d(x,y)>0$. Powtarzając poprzednią
konstrukcję z punktem $y$ i promieniem $\delta$,
znajdujemy otwarty zbiór $W_y$ z
$y\in W_y\subseteq B_d(y,\delta)$.
Z nierówności trójkąta
$B_d(y,\delta)\subseteq B_d(x,\varepsilon)$.
Każdy punkt kuli ma więc zawarte w niej otwarte
otoczenie, czyli kula jest otwarta w początkowej
topologii. W konsekwencji $\tau_d\subseteq\tau$.

\pagebreak[4]
\textbf{Krok 5: każdy zbiór otwarty w $X$ jest otwarty
w metryce $d$.}
Niech $O\in\tau$ i $x\in O$. W kroku 1 znaleźliśmy
$(B,C)\in\mathcal P$ z $x\in B\subseteq C\subseteq O$.
Dla pewnego $n$ funkcja $f_n$ równa jest $0$ na $B$
i $1$ na $X\setminus C$. Jeśli $y\in X\setminus O$,
to $y\in X\setminus C$ oraz
$d(x,y)\ge2^{-n}|f_n(x)-f_n(y)|=2^{-n}$.
Wynika stąd
$B_d(x,2^{-n})\subseteq O$.
Ponieważ $x\in O$ był dowolny, $O\in\tau_d$,
czyli $\tau\subseteq\tau_d$.

Ostatecznie $\tau=\tau_d$, co kończy dowód.
\end{proof}

\begin{figure}[htbp]
\centering
\begin{tikzpicture}[>=Stealth, node distance=8mm,
 node/.style={draw,rounded corners=2pt,fill=blue!6,align=center,
 font=\small,text width=10.4cm,inner sep=6pt}]
\node[node] (a) {Baza $\mathcal B$ i regularność: dla $x\in O$
 dobieramy $B,C\in\mathcal B$ z $x\in B\subseteq\overline B\subseteq C\subseteq O$};
\node[node,below=of a] (b) {Normalność i lemat Urysohna:
 $f_{B,C}|_{\overline B}=0$, $f_{B,C}|_{X\setminus C}=1$};
\node[node,below=of b] (c) {$\displaystyle
 d(x,y)=\sum_{n\ge1}2^{-n}|f_n(x)-f_n(y)|$};
\node[node,below=of c] (d) {Metryka rozdziela punkty, a
 $\tau_d=\tau$};
\draw[->,thick,blue!65!black] (a)--(b);
\draw[->,thick,blue!65!black] (b)--(c);
\draw[->,thick,blue!65!black] (c)--(d);
\end{tikzpicture}
\caption{Etapy konstrukcji metryki w twierdzeniu Urysohna.
Przeliczalność bazy zapewnia, że do sumy wystarczy
ciąg funkcji $f_n$.}
\label{fig:topog-urysohn-metryka-schemat}
\end{figure}

\begin{wniosek}[Kryterium metryzowalności w klasie z drugim aksjomatem przeliczalności]
\label{cor:topog-drugi-metryzowalna-iff-t3}
Niech $X$ spełnia drugi aksjomat przeliczalności.
Wtedy $X$ jest metryzowalna wtedy i tylko wtedy, gdy
spełnia $T_3$.
\end{wniosek}
\begin{proof}
Jeżeli $X$ jest metryzowalna, to spełnia $T_3$
na mocy stwierdzenia \ref{prop:topog-metryka-t3-pierwszy}.
Odwrotna implikacja jest właśnie twierdzeniem
\ref{thm:topog-urysohn-metryzowalnosc}.
\end{proof}

\section{Kostka przeliczalna i przeliczalne iloczyny}
\label{sec:topog-metryzowalnosc-iloczyny}

\begin{twierdzenie}[Metryka na kostce przeliczalnej]
\label{thm:topog-metryka-kostka-przeliczalna}
Na $K=[0,1]^{\mathbb N}$ funkcja
\[
\rho(a,b)=\sum_{n=1}^{\infty}2^{-n}|a_n-b_n|,
\qquad a=(a_n),\ b=(b_n)\in K,
\]
jest metryką. Topologia $\tau_\rho$ jest dokładnie
topologią iloczynu na $K$.
\end{twierdzenie}
\begin{proof}
Zbieżność szeregu oraz aksjomaty metryki wynikają jak
w kroku 3 dowodu twierdzenia Urysohna, przyjmując
$f_n(a)=a_n$. Każde $a\ne b$ różni się na pewnej
współrzędnej $n$, więc $\rho(a,b)\ge2^{-n}|a_n-b_n|>0$.

Ustalmy $a\in K$ i $\varepsilon>0$.
Dobierzmy $N$ z $2^{-N}<\varepsilon/2$.
W topologii iloczynu jest otwarte otoczenie
\[
W=\{b\in K:|a_n-b_n|<\varepsilon/2
\text{ dla }1\le n\le N\}.
\]
Dla $b\in W$ suma pierwszych $N$ wyrazów definiujących
$\rho(a,b)$ jest mniejsza od $\varepsilon/2$,
a pozostałe wyrazy sumują się do co najwyżej $2^{-N}$.
Zatem $a\in W\subseteq B_\rho(a,\varepsilon)$.
Dla dowolnego $b\in B_\rho(a,\varepsilon)$
ponawiamy konstrukcję z promieniem
$\delta=\varepsilon-\rho(a,b)>0$ i środkiem $b$.
Otrzymujemy otwarty w topologii iloczynu zbiór
$W_b$ z $b\in W_b\subseteq B_\rho(b,\delta)$,
a nierówność trójkąta daje
$B_\rho(b,\delta)\subseteq B_\rho(a,\varepsilon)$.
Kula jest więc otwarta w topologii iloczynu,
czyli $\tau_\rho\subseteq\tau_{\mathrm{ilocz}}$.

Odwrotnie, niech $O$ będzie bazowym otwartym
otoczeniem $a$ w topologii iloczynu. Istnieje skończony
$F\subseteq\mathbb N$ i względnie otwarte
$O_n\subseteq[0,1]$, $n\in F$, takie, że
$a_n\in O_n$ i
$\{b\in K:b_n\in O_n\text{ dla }n\in F\}\subseteq O$.
Dla każdego $n\in F$ wybieramy $\delta_n>0$,
aby $\{t\in[0,1]:|t-a_n|<\delta_n\}\subseteq O_n$.
Jeśli $F\ne\varnothing$, przyjmijmy
$\delta=\min_{n\in F}2^{-n}\delta_n>0$.
Z $\rho(a,b)<\delta$ otrzymujemy, dla każdego $n\in F$,
$2^{-n}|a_n-b_n|\le\rho(a,b)<\delta\le2^{-n}\delta_n$.
Stąd $b_n\in O_n$, czyli $B_\rho(a,\delta)\subseteq O$.
Gdy $F=\varnothing$, możemy wziąć dowolne $\delta>0$,
ponieważ bazowy cylinder jest wtedy całym $K$.
Wynika stąd $\tau_{\mathrm{ilocz}}\subseteq\tau_\rho$.
\end{proof}

\begin{wniosek}[Reprezentacja w kostce Hilberta]
\label{cor:topog-urysohn-kostka-hilberta}
Każda przestrzeń $X$ spełniająca $T_3$ i drugi aksjomat
przeliczalności jest homeomorficzna z podprzestrzenią
kostki $[0,1]^{\mathbb N}$.
\end{wniosek}
\begin{proof}
Jeśli $X$ jest pusta, jej obraz pusty jest podprzestrzenią
kostki; dla przestrzeni jednoelementowej wystarcza odwzorowanie
stałe do ustalonego punktu kostki, które jest homeomorfizmem
na swój obraz.

Dla pozostałych przestrzeni wybierzmy funkcje $f_n$
z dowodu twierdzenia \ref{thm:topog-urysohn-metryzowalnosc}
i połóżmy
$\Phi(x)=(f_1(x),f_2(x),\ldots)$.
Wówczas dla wszystkich $x,y\in X$
\[
\rho(\Phi(x),\Phi(y))
=\sum_{n\ge1}2^{-n}|f_n(x)-f_n(y)|=d(x,y).
\]
Równość ta i metryczna charakteryzacja topologii
wykazana wyżej dowodzą, że $\Phi$ jest izometrią
z $(X,d)$ na $\Phi(X)$ z metryką ograniczoną z $\rho$.
Izometria jest różnowartościowa i zachowuje kule,
więc jest homeomorfizmem na swój obraz.
\end{proof}

\begin{twierdzenie}[Przeliczalne iloczyny przestrzeni metryzowalnych]
\label{thm:topog-przeliczalny-iloczyn-metryzowalny}
Niech $(X_n,\tau_n)$, $n\ge1$, będą przestrzeniami
metryzowalnymi. Wtedy iloczyn
$X=\prod_{n\ge1}X_n$, z topologią iloczynu,
jest metryzowalny. Konkretną zgodną metrykę można
skonstruować z dowolnych metryk $d_n$ zgodnych z $\tau_n$:
\[
D(x,y)=\sum_{n=1}^{\infty}
2^{-n}\min\{1,d_n(x_n,y_n)\}.
\]
\end{twierdzenie}
\begin{proof}
Jeżeli którykolwiek $X_n$ jest pusty, iloczyn jest pusty
i twierdzenie jest natychmiastowe. Zakładamy dalej,
że wszystkie czynniki są niepuste.
Niech $\widehat d_n=\min\{1,d_n\}$.
Poprzednie stwierdzenie
\ref{prop:topog-min-metryka} pokazuje, że każda
$\widehat d_n$ jest zgodną metryką ograniczoną przez $1$.

Szereg definiujący $D$ zbiega, ponieważ jego wyrazy
są ograniczone przez $2^{-n}$. Symetria, zerowanie
na przekątnej oraz nierówność trójkąta wynikają ze
zsumowania odpowiednich własności $\widehat d_n$.
Jeśli $x\ne y$, istnieje $k$ z $x_k\ne y_k$,
więc $D(x,y)\ge2^{-k}\widehat d_k(x_k,y_k)>0$.
Stąd $D$ jest metryką.

Ustalmy $x\in X$ i $\varepsilon>0$.
Wybierzmy $N$ tak, aby $2^{-N}<\varepsilon/2$.
Cylinder
\[
W=\{y\in X:\widehat d_n(x_n,y_n)<\varepsilon/2
\text{ dla }1\le n\le N\}
\]
jest otwartym otoczeniem $x$ w topologii iloczynu,
a dla $y\in W$ mamy $D(x,y)<\varepsilon$,
jak w dowodzie twierdzenia o kostce przeliczalnej.
W ten sposób każde $x\in X$ i każdy promień
$\varepsilon>0$ mają otwarte otoczenie $W$
zawarte w $B_D(x,\varepsilon)$. Aby stwierdzić,
że cała kula jest otwarta, dla dowolnego
$z\in B_D(x,\varepsilon)$ stosujemy ten sam argument
z promieniem $\varepsilon-D(x,z)$ i środkiem $z$.
Nierówność trójkąta zapewnia
$B_D(z,\varepsilon-D(x,z))\subseteq B_D(x,\varepsilon)$.
Stąd $\tau_D\subseteq\tau_{\mathrm{ilocz}}$.

Dla odwrotnej inkluzji rozważmy bazowy cylinder
$O=\prod_{n\ge1}O_n$ zawierający $x$,
gdzie $O_n$ jest otwarty w $X_n$, a $O_n=X_n$
poza pewnym skończonym zbiorem indeksów $F$.
Dla każdego $n\in F$ istnieje $r_n>0$ z
$B_{d_n}(x_n,r_n)\subseteq O_n$.
Przyjmijmy $\delta_n=\min\{1/2,r_n\}$.
Jeśli $D(x,y)<\delta=\min_{n\in F}2^{-n}\delta_n$,
to $\widehat d_n(x_n,y_n)<\delta_n<1$ i w konsekwencji
$d_n(x_n,y_n)=\widehat d_n(x_n,y_n)<r_n$ dla $n\in F$.
Tak więc $B_D(x,\delta)\subseteq O$.
Jeżeli $F=\varnothing$, cylinder jest całym $X$
i można przyjąć dowolny promień. Otrzymujemy
$\tau_{\mathrm{ilocz}}\subseteq\tau_D$, a więc
równość obu topologii.
\end{proof}

\begin{wniosek}[Zwarte przestrzenie Hausdorffa]
\label{cor:topog-zwarta-hausdorff-metryzowalna}
Niech $X$ będzie zwartą przestrzenią Hausdorffa.
Wtedy $X$ jest metryzowalna wtedy i tylko wtedy, gdy
spełnia drugi aksjomat przeliczalności.
\end{wniosek}
\begin{proof}
Jeżeli $X$ spełnia drugi aksjomat przeliczalności,
to na mocy rozdziału~07 zwarta przestrzeń Hausdorffa
jest normalna i $T_1$, a więc spełnia $T_3$.
Stosujemy twierdzenie Urysohna o metryzowalności.

Załóżmy teraz, że $X$ jest metryzowalna i wybierzmy
zgodną metrykę $d$. Dla każdego $n\ge1$ rodzina
$\{B_d(x,1/n):x\in X\}$ jest otwartym pokryciem $X$.
Ze zwartości możemy wybrać skończoną rodzinę
punktów $F_n\subseteq X$ taką, że
$X=\bigcup_{a\in F_n}B_d(a,1/n)$.
Jeśli $X$ jest pusta, te rodziny mogą być puste.
Zbiór $D=\bigcup_{n\ge1}F_n$ jest co najwyżej
przeliczalny. Jest gęsty w $X$: dla każdego
$x\in X$ i dowolnej kuli $B_d(x,r)$ wybierzmy
$n$ z $2/n<r$ oraz $a\in F_n$ z $d(x,a)<1/n<r$.
Zatem $D\cap B_d(x,r)\ne\varnothing$.
W rozdziale~09 udowodniliśmy, że ośrodkowość
przestrzeni metrycznej implikuje drugi aksjomat
przeliczalności. Stosujemy ten wynik do $(X,d)$.
\end{proof}

\section{Kontrprzykłady i granice twierdzeń}
\label{sec:topog-metryzowalnosc-kontrprzyklady}

\begin{przyklad}[Prosta Sorgenfreya: warunki konieczne nie wystarczają]
\label{ex:topog-sorgenfrey-niemetryzowalna}
Niech $S$ oznacza prostą Sorgenfreya z bazą
$\{[a,b):a<b\}$, omawianą w rozdziale~09.
Przestrzeń $S$ spełnia pierwszy aksjomat przeliczalności,
jest ośrodkowa i ma własność Lindelöfa, ale nie spełnia
drugiego aksjomatu przeliczalności; wszystkie te własności
zostały udowodnione w rozdziale~09.

Ponadto $S$ jest regularna i $T_1$: baza składa się
ze zbiorów zarazem otwartych i domkniętych.
Rzeczywiście dopełnienie $[a,b)$ jest sumą
$(-\infty,a)$ i $[b,\infty)$, a oba zbiory są otwarte
w topologii Sorgenfreya: pierwszy jest sumą bazowych
przedziałów $[x,a)$, $x<a$, a drugi sumą przedziałów
$[b,c)$ dla $c>b$ (w szczególności $[b,b+1)$,
$[b,b+2)$ itd. pokrywają $[b,\infty)$).
Dla dowolnego otwartego $O\ni x$ możemy dobrać
bazowy $B$ z $x\in B\subseteq O$, a
$\overline B=B\subseteq O$, więc kryterium
zwężania otoczeń z rozdziału~05 daje regularność.
Singletony są domknięte, ponieważ zwykła topologia
euklidesowa, spełniająca $T_1$, jest słabsza od
Sorgenfreya.

Gdyby $S$ była metryzowalna, jej ośrodkowość
implikowałaby drugi aksjomat przeliczalności,
na mocy twierdzenia o przestrzeniach metrycznych
z rozdziału~09. Prowadzi to do sprzeczności.
Zatem $S$ nie jest metryzowalna,
mimo spełniania $T_3$, pierwszego aksjomatu
przeliczalności i własności Lindelöfa.
\end{przyklad}

\begin{przyklad}[Nieprzeliczalna kostka]
\label{ex:topog-kostka-nieprzeliczalna}
Niech $I$ będzie zbiorem nieprzeliczalnym i
$K=[0,1]^I$ ma topologię iloczynu.
Na mocy twierdzenia Tychonoffa z rozdziału~07
$K$ jest zwarta; jest także Hausdorffa, gdyż każdy
czynnik $[0,1]$ jest Hausdorffa, a różniące się punkty
iloczynu można rozdzielić otoczeniami w jednej
współrzędnej.

Pokażemy, że $K$ nie spełnia nawet pierwszego
aksjomatu przeliczalności. Niech $0=(0)_{i\in I}$.
Przypuśćmy, że $(N_n)_{n\ge1}$ jest przeliczalną
bazą otoczeń $0$. Dla każdego $n$ wybierzmy
bazowy otwarty cylinder $W_n$ taki, że
$0\in W_n\subseteq N_n$. Istnieje skończony
zbiór współrzędnych $F_n\subseteq I$, poza którym
$W_n$ nie nakłada żadnych ograniczeń.
Suma $F=\bigcup_{n\ge1}F_n$ jest przeliczalna,
więc możemy wybrać $j\in I\setminus F$.

Zbiór $O_j=\{x\in K:x_j<1/2\}$ jest otwartym
otoczeniem $0$ w topologii iloczynu.
Istniałoby $n$ z $N_n\subseteq O_j$,
a wówczas $W_n\subseteq O_j$.
Ponieważ jednak $j\notin F_n$, punkt $z\in K$
zdefiniowany przez $z_j=1$ i $z_i=0$ dla $i\ne j$
należy do $W_n$, ale nie należy do $O_j$.
To sprzeczność.

Skoro każda przestrzeń metryzowalna spełnia
pierwszy aksjomat przeliczalności, $K$ nie jest
metryzowalna. W ten sposób otrzymujemy zwartą
przestrzeń Hausdorffa, która nie jest metryzowalna.
\end{przyklad}

\clearpage
\begin{przyklad}[Iloraz przestrzeni metryzowalnej nie musi być metryzowalny]
\label{ex:topog-iloraz-dwa-poczatki}
Weźmy rozłączną sumę dwóch kopii prostej rzeczywistej,
$Y=\R\times\{0,1\}$, z topologią sumy rozłącznej.
Jest metryzowalna: na $Y$ można zadać metrykę
\[
\delta((x,i),(y,j))=
\begin{cases}
\min\{1,|x-y|\},&i=j,\\
2,&i\ne j.
\end{cases}
\]
Nierówność trójkąta wynika z nierówności trójkąta
w każdej kopii oraz faktu, że każda odległość między
różnymi kopiami wynosi $2$, podczas gdy odległość
wewnątrz kopii nie przekracza $1$.
Topologia tej metryki jest sumą rozłączną: kule
o promieniu mniejszym od $1$ zawierają punkty tylko
jednej kopii i wyznaczają w niej zwykłą topologię.

Utożsammy $(t,0)$ z $(t,1)$ dla każdego $t\ne0$,
ale pozostawmy rozróżnione początki $(0,0)$ i $(0,1)$.
Niech $q\colon Y\to Z$ będzie odwzorowaniem ilorazowym.
Oznaczmy $o_0=q(0,0)$ i $o_1=q(0,1)$; mamy $o_0\ne o_1$.
Jeśli $U,V\subseteq Z$ są otwartymi otoczeniami
odpowiednio $o_0$ i $o_1$, to przeciwobrazy
$q^{-1}(U)$ i $q^{-1}(V)$ są otwarte w $Y$,
zawierając odpowiednio $(0,0)$ i $(0,1)$.
Istnieje więc $\varepsilon>0$ takie, że
$(t,0)\in q^{-1}(U)$ i $(t,1)\in q^{-1}(V)$
dla każdego $0<t<\varepsilon$.
Ponieważ $q(t,0)=q(t,1)$ dla $t\ne0$, otrzymujemy
$U\cap V\ne\varnothing$. Zatem dwóch początków
nie można rozdzielić rozłącznymi zbiorami otwartymi:
$Z$ nie jest Hausdorffa, a więc nie jest metryzowalna.
\end{przyklad}

\section{Zestawienie implikacji i zakresu założeń}
\label{sec:topog-metryzowalnosc-zestawienie}

\begin{stwierdzenie}[Podsumowanie kryteriów]
\label{prop:topog-metryzowalnosc-zestawienie}
Dla każdej przestrzeni topologicznej $X$ zachodzą następujące
stwierdzenia.
\begin{enumerate}[label=(\roman*)]
\item Jeśli $X$ jest metryzowalna, to spełnia $T_3$
i pierwszy aksjomat przeliczalności.
\item Jeśli $X$ spełnia $T_3$ i drugi aksjomat przeliczalności,
to jest metryzowalna.
\item Każda podprzestrzeń metryzowalnej przestrzeni
jest metryzowalna; to samo dotyczy przeliczalnego
iloczynu przestrzeni metryzowalnych.
\item Jeżeli $X$ jest zwarta i Hausdorffa, to
metryzowalność jest równoważna spełnianiu
drugiego aksjomatu przeliczalności.
\end{enumerate}
Żadnego z warunków $T_3$, pierwszego aksjomatu
przeliczalności, ośrodkowości i własności Lindelöfa
nie można, nawet po połączeniu wszystkich czterech,
użyć jako zamiennika drugiego aksjomatu przeliczalności
w twierdzeniu Urysohna.
\end{stwierdzenie}
\begin{proof}
Punkty (i)--(iv) są dokładnie wynikami
\ref{prop:topog-metryka-t3-pierwszy},
\ref{thm:topog-urysohn-metryzowalnosc},
\ref{thm:topog-metryzowalnosc-podprzestrzenie},
\ref{thm:topog-przeliczalny-iloczyn-metryzowalny}
i \ref{cor:topog-zwarta-hausdorff-metryzowalna}.
Ostatnią uwagę uzasadnia prosta Sorgenfreya z przykładu
\ref{ex:topog-sorgenfrey-niemetryzowalna}.
\end{proof}

\begin{uwaga}[Granice zakresu rozdziału]
Rozważaliśmy metryzowalność topologiczną,
nie zaś zupełność konkretnej metryki. Ta sama topologia
może pochodzić od wielu metryk o różnych własnościach
metrycznych. Na przykład na $\R$ zarówno
$d(x,y)=|x-y|$, jak i
$d'(x,y)=|\arctan(x)-\arctan(y)|$
indukują topologię euklidesową. Dla $d'$ wynika to
z faktu, że $\arctan\colon\R\to(-\pi/2,\pi/2)$
jest homeomorfizmem: jest ciągła, ściśle rosnąca,
a funkcja odwrotna $\tan$ jest ciągła na podanym
przedziale. Metryka $d'$ jest więc metryką przeniesioną
przez homeomorfizm na podprzestrzeń prostej metrycznej,
zgodnie ze stwierdzeniem
\ref{prop:topog-metryzowalnosc-homeo}.
Nie jest jednak zupełna: ciąg $x_n=n$ jest ciągiem
Cauchy'ego w metryce $d'$ (bo $\arctan(n)$ jest zbieżny),
natomiast nie ma granicy w $\R$ w żadnej z tych
równoważnych topologii. Metryzowalności nie wolno więc
utożsamiać ze zupełną metryzowalnością.
\end{uwaga}

\chapter{Zupełność przestrzeni metrycznych i twierdzenie Baire'a}
\label{ch:topog-baire}

\noindent\textbf{Cel rozdziału.}
Metryzowalność przestrzeni topologicznej oznacza istnienie metryki
zgodnej z jej topologią (rozdział~10). Zupełność zależy natomiast
od konkretnej metryki. Przedstawimy pojęcie ciągu Cauchy'ego,
odróżnimy zupełność od zupełnej metryzowalności i udowodnimy
twierdzenie Baire'a w dwóch podstawowych wersjach:
dla zupełnych przestrzeni metrycznych oraz lokalnie zwartych
przestrzeni Hausdorffa. Zastosowania dotyczą kategorii Baire'a.

\section{Ciągi Cauchy'ego i zupełność}
\label{sec:topog-baire-zupelnosc}

\begin{definicja}[Ciąg Cauchy'ego i przestrzeń zupełna]
\label{def:topog-baire-ciag-cauchyego}
Niech $(X,d)$ będzie przestrzenią metryczną.
Ciąg $(x_n)_{n\ge1}$ w $X$ jest \emph{ciągiem Cauchy'ego},
jeżeli dla każdego $\varepsilon>0$ istnieje $N\ge1$ takie,
że $d(x_m,x_n)<\varepsilon$ dla wszystkich $m,n\ge N$.
Przestrzeń $(X,d)$ jest \emph{zupełna}, jeżeli każdy
jej ciąg Cauchy'ego jest zbieżny w $X$.
\end{definicja}

\begin{stwierdzenie}[Ciąg zbieżny jest ciągiem Cauchy'ego]
\label{prop:topog-baire-zbiezny-cauchy}
W każdej przestrzeni metrycznej $(X,d)$ ciąg zbieżny
jest ciągiem Cauchy'ego i ma dokładnie jedną granicę.
\end{stwierdzenie}
\begin{proof}
Jeżeli $x_n\to x$, to dla $\varepsilon>0$ dobieramy $N$
z $d(x_n,x)<\varepsilon/2$ dla $n\ge N$.
Dla $m,n\ge N$ nierówność trójkąta daje
$d(x_m,x_n)\le d(x_m,x)+d(x,x_n)<\varepsilon$.
Jeśli jednocześnie $x_n\to y$, to
$d(x,y)\le d(x,x_n)+d(x_n,y)\to0$, skąd $x=y$.
\end{proof}

\begin{twierdzenie}[Zupełność podprzestrzeni]
\label{thm:topog-baire-podprzestrzen-zupelna}
Niech $(X,d)$ będzie zupełną przestrzenią metryczną
i $A\subseteq X$. Przestrzeń $(A,d|_{A\times A})$
jest zupełna wtedy i tylko wtedy, gdy $A$ jest domknięty w $X$.
\end{twierdzenie}
\begin{proof}
Jeśli $A$ jest domknięty, każdy ciąg Cauchy'ego
$(a_n)$ w $A$ jest też ciągiem Cauchy'ego w $X$.
Zbiega zatem do $x\in X$, a ponieważ $a_n\in A$,
mamy $x\in\overline A=A$.

Odwrotnie, załóżmy, że $A$ jest zupełna, i weźmy
$x\in\overline A$. Dla każdego $n$ znajdziemy
$a_n\in A$ z $d(a_n,x)<1/n$. Wtedy $a_n\to x$,
więc $(a_n)$ jest ciągiem Cauchy'ego w $A$.
Z zupełności otrzymujemy $a_n\to a$ dla pewnego
$a\in A$. Jednoznaczność granicy w $X$ daje
$x=a\in A$, zatem $A$ jest domknięty.
\end{proof}

\begin{przyklad}[Przestrzenie zupełne i niezupełne]
\label{ex:topog-baire-podstawowe-przyklady}
Prosta $(\R,|\cdot|)$ jest zupełna. Rzeczywiście,
każdy ciąg Cauchy'ego w $\R$ jest ograniczony
(od pewnego miejsca jego wyrazy są w odległości
mniejszej niż $1$ od ustalonego wyrazu).
Na mocy twierdzenia Bolzana--Weierstrassa
ma więc podciąg zbieżny. Jeśli $x_{n_k}\to x$,
to dla dowolnego $\varepsilon>0$ wybieramy $N$
z $|x_m-x_n|<\varepsilon/2$ dla $m,n\ge N$,
a następnie $k$ z $n_k\ge N$ i
$|x_{n_k}-x|<\varepsilon/2$. Dla $n\ge N$
otrzymujemy $|x_n-x|<\varepsilon$, czyli
zbiega cały ciąg.

Przedział $[0,1]$ jest zupełny jako domknięty
podzbiór $\R$. Przedział $(0,1)$ ze zwykłą
metryką nie jest zupełny, gdyż $1/(n+1)$
jest ciągiem Cauchy'ego bez granicy w $(0,1)$.
Podobnie $\Q$ z metryką odziedziczoną z $\R$
nie jest zupełna: wymierne przybliżenia liczby
$\sqrt2$ tworzą ciąg Cauchy'ego bez granicy w $\Q$.
\end{przyklad}

\begin{definicja}[Zupełna metryzowalność]
\label{def:topog-baire-zupelna-metryzowalnosc}
Przestrzeń topologiczna $X$ jest
\emph{zupełnie metryzowalna}, jeżeli istnieje
co najmniej jedna metryka zupełna na $X$
zgodna z jego topologią.
\end{definicja}

\begin{przyklad}[Zupełność metryki nie jest własnością topologiczną]
\label{ex:topog-baire-metryka-przedzial}
Homeomorfizm $h\colon(0,1)\to\R$,
$h(t)=\tan(\pi(t-\tfrac12))$, przenosi na $(0,1)$
metrykę $d_h(s,t)=|h(s)-h(t)|$. Z rozdziału~10
wiemy, że jest ona zgodna ze zwykłą topologią.
Jest zupełna, bo $h$ jest izometrią na zupełną
przestrzeń $(\R,|\cdot|)$. Tymczasem zwykła
metryka $|s-t|$ na $(0,1)$ nie jest zupełna.
A zatem zupełność metryki nie jest niezmiennikiem
topologicznym. Natomiast zupełna metryzowalność
jest niezmiennikiem topologicznym: jeśli
$g\colon X\to Y$ jest homeomorfizmem, a $e$ jest
zgodną metryką zupełną na $Y$, to metryka
$e_g(x,x')=e(g(x),g(x'))$ jest zgodna z topologią
$X$ (rozdział~10), a $g$ jest izometrią
$(X,e_g)$ na $(Y,e)$, więc $(X,e_g)$ jest zupełna.
\end{przyklad}

\section{Nigdziegęstość i kategorie zbiorów}
\label{sec:topog-baire-kategorie}

\begin{definicja}[Pierwsza kategoria, rezydualność, przestrzeń Baire'a]
\label{def:topog-baire-kategoria}
Niech $X$ będzie przestrzenią topologiczną.
Zbiór $A\subseteq X$ jest \emph{nigdziegęsty},
jeżeli $\operatorname{int}_X(\overline A^{\,X})
=\varnothing$. Zbiór $M\subseteq X$ jest
\emph{pierwszej kategorii} (albo \emph{mizerny}),
jeśli istnieją nigdziegęste zbiory $A_n\subseteq X$,
$n\ge1$, takie, że $M\subseteq\bigcup_{n\ge1}A_n$.
Zbiór $R\subseteq X$ jest \emph{rezydualny},
jeśli $X\setminus R$ jest pierwszej kategorii.
Przestrzeń $X$ jest \emph{przestrzenią Baire'a},
jeżeli każde przeliczalne przecięcie
zbiorów otwartych i gęstych w $X$ jest gęste w $X$.
\end{definicja}

\begin{uwaga}[Kategoria i miara]
Zbiór pierwszej kategorii może być gęsty,
a nigdziegęsty nie musi być przeliczalny.
Własność pierwszej kategorii nie oznacza
zerowej miary; kategoria i miara są
odrębnymi sposobami badania wielkości zbioru.
\end{uwaga}

\begin{stwierdzenie}[Równoważne warunki Baire'a]
\label{prop:topog-baire-rownowaznosci}
Dla każdej przestrzeni topologicznej $X$
równoważne są:
\begin{enumerate}[label=(\roman*)]
\item $X$ jest przestrzenią Baire'a;
\item każdy zbiór pierwszej kategorii w $X$
ma puste wnętrze w $X$;
\item żaden niepusty zbiór otwarty w $X$
nie jest pierwszej kategorii w $X$.
\end{enumerate}
\end{stwierdzenie}
\begin{proof}
Jeśli $A_n$ są nigdziegęste, to
$G_n=X\setminus\overline{A_n}^{\,X}$
są otwarte i gęste. Z (i) wynika, że
$\bigcap_nG_n$ jest gęste, więc jego
dopełnienie $\bigcup_n\overline{A_n}^{\,X}$
ma puste wnętrze. Każdy zbiór pierwszej
kategorii zawiera się w takiej sumie:
otrzymujemy (ii). Implikacja (ii)$\Rightarrow$(iii)
jest natychmiastowa.

Załóżmy (iii). Dla dowolnych zbiorów
otwartych i gęstych $G_n$ dopełnienia
$F_n=X\setminus G_n$ są nigdziegęste.
Gdyby $\bigcap_nG_n$ nie było gęste,
istniałby niepusty zbiór otwarty
$U\subseteq X\setminus\bigcap_nG_n
=\bigcup_nF_n$. Oznaczałoby to, że
$U$ jest pierwszej kategorii w $X$,
wbrew (iii).
\end{proof}

\section{Twierdzenie Baire'a w przestrzeniach zupełnych}
\label{sec:topog-baire-twierdzenie-metryczne}

\begin{twierdzenie}[Twierdzenie Baire'a o kategorii --- wersja metryczna]
\label{thm:topog-baire-metryczne}
Jeśli $(X,d)$ jest zupełną przestrzenią metryczną,
a $(G_n)_{n\ge1}$ jest ciągiem zbiorów otwartych
i gęstych w $X$, to dla każdego niepustego
otwartego $U\subseteq X$ zachodzi
$U\cap\bigcap_{n\ge1}G_n\ne\varnothing$.
W szczególności $X$ jest przestrzenią Baire'a.
\end{twierdzenie}
\begin{proof}
Zbiór $U\cap G_1$ jest niepusty i otwarty.
Wybierzmy $x_1\in U\cap G_1$ oraz
$r_1\in(0,2^{-1})$ tak mały, że
$K_1=\{y\in X:d(x_1,y)\le r_1\}
\subseteq U\cap G_1$.

Załóżmy, że mamy już niepustą domkniętą kulę
$K_n=\{y:d(x_n,y)\le r_n\}$, gdzie
$0<r_n<2^{-n}$. Gęstość $G_{n+1}$
zapewnia istnienie punktu
$x_{n+1}\in B_d(x_n,r_n)\cap G_{n+1}$.
Z otwartości tego przecięcia możemy
wybrać $0<r_{n+1}<2^{-(n+1)}$ tak, aby
\[
K_{n+1}:=\{y:d(x_{n+1},y)\le r_{n+1}\}
\subseteq B_d(x_n,r_n)\cap G_{n+1}.
\]
Stąd $K_1\supseteq K_2\supseteq\cdots$
oraz $K_n\subseteq G_n$ dla wszystkich $n$.

Dla $m>n$ mamy $x_m\in K_n$, skąd
$d(x_m,x_n)\le r_n<2^{-n}$. Ciąg $(x_n)$
jest zatem ciągiem Cauchy'ego i na mocy
zupełności zbiega do pewnego $x\in X$.
Przy ustalonym $n$ wszystkie $x_m$ dla
$m\ge n$ należą do domkniętego $K_n$,
więc również granica $x$ należy do $K_n$.
Otrzymujemy
$x\in\bigcap_nK_n\subseteq
U\cap\bigcap_nG_n$, co kończy dowód.
\end{proof}

\begin{figure}[htbp]
\centering
\begin{tikzpicture}
\draw[blue!70!black,thick] (0,0) circle (1.9);
\draw[teal!75!black,thick] (.6,0) circle (1.0);
\draw[orange!80!black,thick] (.96,.08) circle (.43);
\fill (1.03,.09) circle (1.4pt);
\node[blue!70!black] at (-1.45,1.5) {$K_1$};
\node[teal!75!black] at (.22,.7) {$K_2$};
\node[orange!80!black] at (1.05,.33) {$K_3$};
\node[anchor=west] at (2.05,.3) {$K_{n+1}\subseteq K_n$};
\node[anchor=west] at (2.05,-.22) {$r_n<2^{-n}$};
\end{tikzpicture}
\caption{Schemat konstrukcji zstępujących kul domkniętych.
Rysunek ilustruje inkluzje, a nie rzeczywiste zanurzenie
dowolnej przestrzeni metrycznej w płaszczyźnie.}
\label{fig:topog-baire-kule}
\end{figure}

\begin{wniosek}[Zupełna metryzowalność a własność Baire'a]
\label{cor:topog-baire-zupelnie-metryzowalne}
Każda przestrzeń zupełnie metryzowalna
jest przestrzenią Baire'a.
\end{wniosek}
\begin{proof}
Wybieramy zgodną metrykę zupełną i stosujemy
twierdzenie~\ref{thm:topog-baire-metryczne}.
Otwartość i gęstość zależą tylko od topologii.
\end{proof}

\section{Twierdzenie Baire'a w przestrzeniach lokalnie zwartych}
\label{sec:topog-baire-lokalnie-zwarte}

\begin{definicja}[Zwartość lokalna]
\label{def:topog-baire-lokalna-zwartosc}
Przestrzeń topologiczna $X$ jest
\emph{lokalnie zwarta}, jeżeli każdy
$x\in X$ ma zwarte otoczenie:
istnieje zwarty $K\subseteq X$
taki, że $x\in\operatorname{int}_X K$.
\end{definicja}

\begin{lemat}[Otoczenia o zwartym domknięciu]
\label{lem:topog-baire-zwarte-domkniecie}
Niech $X$ będzie lokalnie zwartą przestrzenią Hausdorffa.
Dla każdego $x\in X$ i każdego otwartego otoczenia
$V$ punktu $x$ istnieje otwarty $W\ni x$
o zwartym domknięciu $\overline W^{\,X}\subseteq V$.
\end{lemat}
\begin{proof}
Z lokalnej zwartości wybieramy zwarty
$K\subseteq X$ z $x\in\operatorname{int}_X K$.
Przestrzeń $K$ jest zwarta i Hausdorffa,
więc normalna (rozdział~07). Zbiór
$O=V\cap\operatorname{int}_X K$ jest
otwartym otoczeniem $x$ w $K$.
Kryterium zwężania otoczeń w przestrzeni
regularnej (rozdział~05) daje zbiór $H$
otwarty w $K$ spełniający
$x\in H\subseteq\overline H^{\,K}\subseteq O$.
Ponieważ $H\subseteq\operatorname{int}_X K$,
jest on także otwarty w $X$: jeśli
$H=T\cap K$ dla otwartego $T\subseteq X$,
to $H=T\cap\operatorname{int}_X K$.
Przyjmujemy $W=H$. Zwartość $K$
i własność Hausdorffa zapewniają domkniętość
$K$ w $X$ (rozdział~07), toteż
$\overline W^{\,X}=\overline H^{\,K}\subseteq V$.
Domknięty podzbiór zwartego $K$
jest zwarty, co kończy dowód.
\end{proof}

\begin{twierdzenie}[Twierdzenie Baire'a --- wersja lokalnie zwarta]
\label{thm:topog-baire-lokalnie-zwarte}
Każda lokalnie zwarta przestrzeń Hausdorffa
jest przestrzenią Baire'a.
\end{twierdzenie}
\begin{proof}
Niech $U$ będzie niepustym zbiorem otwartym,
a $G_n$ --- otwartymi zbiorami gęstymi.
Wybieramy $x_1\in U\cap G_1$.
Lemat~\ref{lem:topog-baire-zwarte-domkniecie}
daje niepusty otwarty $W_1$ o niepustym
zwartym domknięciu $K_1\subseteq U\cap G_1$.

Indukcyjnie, mając niepusty otwarty $W_n$
o zwartym domknięciu $K_n$, wybieramy
$x_{n+1}\in W_n\cap G_{n+1}$
i na mocy lematu otrzymujemy niepusty otwarty
$W_{n+1}$, którego zwarte domknięcie
$K_{n+1}$ zawiera się w $W_n\cap G_{n+1}$.
Zbiory $K_n$ są więc niepuste, zstępujące,
domknięte w $X$ (gdyż $X$ jest Hausdorffa)
i wszystkie zawierają się w zwartym $K_1$.

Gdyby $\bigcap_{n\ge1}K_n=\varnothing$,
to otwarte w $K_1$ zbiory $K_1\setminus K_n$
tworzyłyby pokrycie $K_1$. Ze zwartości
istniałoby podpokrycie skończone,
a ponieważ $(K_n)$ jest zstępujący,
oznaczałoby to $K_N=\varnothing$ dla
pewnego $N$. Jest to sprzeczność.
Zatem $\bigcap_nK_n$ jest niepuste
i zawiera się w $U\cap\bigcap_nG_n$.
\end{proof}

\section{Konsekwencje i kontrprzykłady}
\label{sec:topog-baire-przyklady}

\begin{przyklad}[Zbiór wymiernych jest mizerny i gęsty]
\label{ex:topog-baire-wymierne}
Każdy singleton $\{q\}$ jest nigdziegęsty
w $\R$, ponieważ jest domknięty i ma puste wnętrze.
Ponieważ $\Q$ jest przeliczalny, jest pierwszej
kategorii w $\R$, choć jednocześnie jest gęsty.
Natomiast $\R\setminus\Q$ jest rezydualny
i gęsty w $\R$.
\end{przyklad}

\begin{twierdzenie}[Zbiór wymiernych nie jest zbiorem $G_\delta$]
\label{thm:topog-baire-wymierne-nie-gdelta}
Zbioru $\Q$ nie można przedstawić jako
przeliczalnego przecięcia zbiorów otwartych w $\R$.
\end{twierdzenie}
\begin{proof}
Załóżmy, że $\Q=\bigcap_{n\ge1}U_n$,
gdzie $U_n$ są otwarte w $\R$.
Ponieważ $\Q$ jest gęsty, wszystkie $U_n$
są gęste. Wypiszmy $\Q=\{q_n:n\ge1\}$.
Każdy $V_n=\R\setminus\{q_n\}$ jest
otwarty i gęsty. Z twierdzenia Baire'a
dla zupełnej prostej $\R$ wynika,
że przecięcie $\bigcap_n(U_n\cap V_n)$
jest gęste i w szczególności niepuste.
Tymczasem jest równe
$\Q\cap(\R\setminus\Q)=\varnothing$.
Sprzeczność.
\end{proof}

\begin{przyklad}[Metryzowalność nie wystarcza]
\label{ex:topog-baire-q-nie-baire}
Przestrzeń $\Q$ z topologią odziedziczoną
z $\R$ jest metryzowalna, ale nie jest Baire'a.
Dla każdego $q\in\Q$ singleton $\{q\}$
jest domknięty i ma puste wnętrze w $\Q$,
więc $\Q$ jest pierwszej kategorii w sobie.
Jest zarazem niepustym otwartym zbiorem
w sobie, co przeczy
stwierdzeniu~\ref{prop:topog-baire-rownowaznosci}(iii).
\end{przyklad}

\begin{stwierdzenie}[Otwarte podprzestrzenie przestrzeni Baire'a]
\label{prop:topog-baire-otwarte-podprzestrzenie}
Jeżeli $X$ jest przestrzenią Baire'a,
a $U\subseteq X$ jest otwarty,
to $U$ z topologią podprzestrzeni
jest przestrzenią Baire'a.
\end{stwierdzenie}
\begin{proof}
Niech $V\subseteq U$ będzie niepusty i otwarty
w $U$, a $H_n\subseteq U$ będą otwarte i gęste w $U$.
Ponieważ $U$ jest otwarty w $X$, także $V$
i $H_n$ są otwarte w $X$. Definiujemy
$G_n=H_n\cup(X\setminus\overline U^{\,X})$.
Zbiory $G_n$ są otwarte w $X$.
Ponieważ $H_n$ jest gęsty w $U$,
jego domknięcie w $X$ zawiera $\overline U^{\,X}$.
Wobec tego $G_n$ jest gęsty w $X$.
Własność Baire'a przestrzeni $X$ daje
$x\in V\cap\bigcap_nG_n$.
Ponieważ $V\subseteq U$ i $G_n\cap U=H_n$,
mamy $x\in V\cap\bigcap_nH_n$.
\end{proof}

\begin{uwaga}[Własność Baire'a a zupełność]
Przedział $(0,1)$ z metryką euklidesową
nie jest zupełny, ale jest przestrzenią Baire'a
jako otwarta podprzestrzeń $\R$.
Jest on również zupełnie metryzowalny
na mocy przykładu~\ref{ex:topog-baire-metryka-przedzial}.
Wersja metryczna twierdzenia Baire'a
jest więc implikacją, a nie charakteryzacją
zupełności wybranej metryki.
\end{uwaga}

\begin{wniosek}[Przeliczalna suma zbiorów domkniętych]
\label{cor:topog-baire-domkniete}
Niech $X$ będzie niepustą przestrzenią Baire'a.
Jeżeli $X=\bigcup_{n\ge1}F_n$ dla domkniętych
$F_n\subseteq X$, to co najmniej jeden
z $F_n$ ma niepuste wnętrze w $X$.
\end{wniosek}
\begin{proof}
Gdyby każde $F_n$ miało puste wnętrze,
wszystkie byłyby nigdziegęste,
a $X$ byłaby pierwszej kategorii w sobie.
Przeczyłoby to
stwierdzeniu~\ref{prop:topog-baire-rownowaznosci}(iii)
dla niepustego otwartego zbioru $X$.
\end{proof}

\begin{uwaga}[Związek z analizą funkcjonalną]
Ostatni wniosek to podstawowy argument
w twierdzeniach Banacha--Steinhausa,
o odwzorowaniu otwartym oraz o wykresie domkniętym.
Ich pełne dowody znajdują się w skrypcie
\emph{Analiza funkcjonalna}.
\end{uwaga}

\chapter{Przestrzenie jednostajne i ciągłość jednostajna}
\label{ch:topog-przestrzenie-jednostajne}

\noindent\textbf{Cel rozdziału.}
Metryka pozwala porównywać odległości punktów niezależnie od miejsca,
w którym znajdują się one w przestrzeni. Sama topologia określa
jedynie lokalne otoczenia poszczególnych punktów. Struktura jednostajna
uogólnia metryczną bliskość par punktów, nie wymagając liczbowego
pomiaru odległości. Wprowadzimy topologię wyznaczoną przez taką
strukturę, ciągłość jednostajną, konstrukcje podprzestrzeni i iloczynów
oraz ciągi, sieci i filtry Cauchy'ego. Rozdział~13 będzie poświęcony
zupełności jednostajnej i uzupełnieniom.

\section{Relacje i struktury jednostajne}
\label{sec:topog-jednostajne-definicje}

\begin{definicja}[Odwracanie i składanie relacji]
\label{def:topog-jednostajne-relacje}
Niech $X$ będzie zbiorem i $E,F\subseteq X\times X$.
Definiujemy przekątną $\Delta_X=\{(x,x):x\in X\}$,
relację odwrotną $E^{-1}=\{(y,x):(x,y)\in E\}$,
przekrój $E[x]=\{y\in X:(x,y)\in E\}$ oraz złożenie
\[
 E\circ F=\{(x,z)\in X\times X:
  \exists y\in X\ ((x,y)\in E\ \land\ (y,z)\in F)\}.
\]
Relację $E$ nazywamy \emph{symetryczną}, gdy $E=E^{-1}$.
\end{definicja}

\begin{definicja}[Struktura jednostajna]
\label{def:topog-jednostajne-struktura}
\emph{Strukturą jednostajną} na zbiorze $X$ nazywamy
niepustą rodzinę $\mathcal U$ podzbiorów $X\times X$,
której elementy nazywamy \emph{otoczeniami jednostajnymi
przekątnej}, spełniającą warunki:
\begin{enumerate}[label=(\roman*)]
\item $\Delta_X\subseteq E$ dla każdego $E\in\mathcal U$;
\item jeżeli $E\in\mathcal U$ i $E\subseteq H\subseteq X\times X$,
      to $H\in\mathcal U$;
\item jeżeli $E,F\in\mathcal U$, to $E\cap F\in\mathcal U$;
\item jeżeli $E\in\mathcal U$, to $E^{-1}\in\mathcal U$;
\item dla każdego $E\in\mathcal U$ istnieje $F\in\mathcal U$
      takie, że $F\circ F\subseteq E$.
\end{enumerate}
Parę $(X,\mathcal U)$ nazywamy \emph{przestrzenią jednostajną}.
Nie wymagamy aksjomatu Hausdorffa.
\end{definicja}

\begin{uwaga}[Przestrzeń pusta]
Jeśli $X=\varnothing$, jedyną strukturą jednostajną jest
$\{\varnothing\}$. W szczególności nie wymagamy, aby każdy
element $\mathcal U$ był niepusty. Dla $X\ne\varnothing$
niepustość wynika z zawierania przekątnej.
\end{uwaga}

\begin{lemat}[Symetryczne otoczenie o małym złożeniu]
\label{lem:topog-jednostajne-symetryczne}
Dla każdego $E\in\mathcal U$ istnieje symetryczne
$F\in\mathcal U$ takie, że $F\subseteq E$ oraz
$F\circ F\subseteq E$.
\end{lemat}
\begin{proof}
Z aksjomatu (v) wybieramy $V\in\mathcal U$
z $V\circ V\subseteq E$. Kładziemy
$F=V\cap V^{-1}\cap E\cap E^{-1}$.
Z aksjomatów (iii) i (iv) wynika, że $F\in\mathcal U$;
z definicji $F=F^{-1}$ oraz $F\subseteq V\cap E$.
Wobec tego $F\circ F\subseteq V\circ V\subseteq E$.
\end{proof}

\begin{przyklad}[Dwie skrajne struktury]
\label{ex:topog-jednostajne-skrajne}
Na dowolnym zbiorze $X$ struktura
$\mathcal U_{\mathrm{ind}}=\{X\times X\}$ jest
\emph{niedyskretna}. Strukturę
$\mathcal U_{\mathrm{dysk}}=
\{E\subseteq X\times X:\Delta_X\subseteq E\}$
nazywamy \emph{dyskretną}. W drugim przypadku
aksjomat (v) spełnia $F=\Delta_X$, ponieważ
$\Delta_X\circ\Delta_X=\Delta_X$.
Obie konstrukcje działają także dla $X=\varnothing$.
\end{przyklad}

\section{Topologia wyznaczona przez strukturę jednostajną}
\label{sec:topog-jednostajne-topologia}

\begin{twierdzenie}[Topologia indukowana]
\label{thm:topog-jednostajne-topologia}
Niech $(X,\mathcal U)$ będzie przestrzenią jednostajną.
Rodzina
\[
 \tau_{\mathcal U}=
 \{O\subseteq X:\ \forall x\in O\ \exists E\in\mathcal U
                   \ (E[x]\subseteq O)\}
\]
jest topologią na $X$. Dla każdego $x\in X$
rodzina $\{E[x]:E\in\mathcal U\}$ stanowi bazę
otoczeń $x$ (niekoniecznie otwartych).
\end{twierdzenie}
\begin{proof}
Zbiór pusty spełnia podany warunek w sposób pusty.
Ponieważ $\mathcal U$ jest niepusta, dla $O=X$
można wziąć dowolne $E\in\mathcal U$.
Jeżeli $x$ należy do sumy zbiorów należących do
$\tau_{\mathcal U}$, to leży w jednym z nich i
odpowiedni przekrój $E[x]$ mieści się w całej sumie.
Jeśli $x\in O_1\cap O_2$, wybieramy
$E_i[x]\subseteq O_i$ dla $E_i\in\mathcal U$;
wtedy $(E_1\cap E_2)[x]\subseteq O_1\cap O_2$.
To dowodzi aksjomatów topologii.

Pokażemy, że każdy $E[x]$ jest otoczeniem $x$.
Niech
\[
 W=\{y\in X:\exists D\in\mathcal U\
                         (D[y]\subseteq E[x])\}.
\]
Ponieważ $E[x]\subseteq E[x]$, mamy $x\in W$,
a z $\Delta_X\subseteq D$ wynika $W\subseteq E[x]$.
Jeśli $y\in W$, wybieramy $D\in\mathcal U$
z $D[y]\subseteq E[x]$ oraz $F\in\mathcal U$
z $F\circ F\subseteq D$. Dla $z\in F[y]$ i
$w\in F[z]$ mamy $(y,w)\in F\circ F\subseteq D$,
więc $w\in E[x]$. Zatem $F[z]\subseteq E[x]$,
co oznacza $z\in W$. Dowiedliśmy $F[y]\subseteq W$,
a więc $W\in\tau_{\mathcal U}$.
Otrzymujemy $x\in W\subseteq E[x]$.

Odwrotnie, jeśli $V$ jest otoczeniem $x$,
to istnieje otwarty $O$ z $x\in O\subseteq V$.
Z definicji $\tau_{\mathcal U}$ można wybrać
$E\in\mathcal U$ z $E[x]\subseteq O\subseteq V$.
Jest to żądana własność bazy otoczeń.
\end{proof}

\begin{twierdzenie}[Kryterium Hausdorffa]
\label{thm:topog-jednostajne-hausdorff}
Dla przestrzeni jednostajnej $(X,\mathcal U)$
równoważne są warunki:
\begin{enumerate}[label=(\roman*)]
\item topologia $\tau_{\mathcal U}$ spełnia aksjomat Hausdorffa;
\item $\bigcap_{E\in\mathcal U}E=\Delta_X$.
\end{enumerate}
Strukturę spełniającą (ii) nazywamy \emph{rozdzielającą}.
\end{twierdzenie}
\begin{proof}
Jeśli topologia spełnia aksjomat Hausdorffa,
to dla $x\ne y$ istnieje otwarty $O$ z
$x\in O$ i $y\notin O$. Z definicji
topologii wybieramy $E\in\mathcal U$ z
$E[x]\subseteq O$, a zatem $(x,y)\notin E$.
Każda relacja $E\in\mathcal U$ zawiera przekątną,
stąd zachodzi (ii).

Załóżmy (ii) i ustalmy $x\ne y$. Wybieramy
$E\in\mathcal U$ z $(x,y)\notin E$ oraz
symetryczne $F\in\mathcal U$ z
$F\circ F\subseteq E$,
korzystając z lematu~\ref{lem:topog-jednostajne-symetryczne}.
Gdyby $z\in F[x]\cap F[y]$, z symetrii mielibyśmy
$(x,z),(z,y)\in F$, a więc $(x,y)\in E$,
co jest sprzecznością. Zbiory $F[x]$ i $F[y]$
są zatem rozłącznymi otoczeniami odpowiednich punktów.
Z twierdzenia~\ref{thm:topog-jednostajne-topologia}
każde zawiera otwarte otoczenie swojego punktu,
co dowodzi warunku Hausdorffa.
\end{proof}

\begin{wniosek}[Topologie skrajnych struktur]
\label{cor:topog-jednostajne-skrajne-topologie}
Topologia indukowana przez $\mathcal U_{\mathrm{ind}}$
jest niedyskretna, a przez $\mathcal U_{\mathrm{dysk}}$
dyskretna.
\end{wniosek}
\begin{proof}
W pierwszym przypadku $E[x]=X$ dla jedynej
relacji $E=X\times X$. Niepusty zbiór otwarty
musi więc być całym $X$. W drugim
$\Delta_X[x]=\{x\}$; dla każdego $O\subseteq X$
i $x\in O$ przekrój ten zawiera się w $O$.
\end{proof}

\section{Struktury wyznaczone przez metryki i pseudometryki}
\label{sec:topog-jednostajne-metryczne}

\begin{twierdzenie}[Struktura jednostajna pseudometryki]
\label{thm:topog-jednostajne-pseudometryka}
Niech $p$ będzie pseudometryką na zbiorze $X$.
Dla $\varepsilon>0$ oznaczmy
$E_\varepsilon=\{(x,y)\in X\times X:p(x,y)<\varepsilon\}$.
Wtedy
\[
 \mathcal U_p=\{E\subseteq X\times X:
      \exists\varepsilon>0\ (E_\varepsilon\subseteq E)\}
\]
jest strukturą jednostajną, a $\tau_{\mathcal U_p}$
jest topologią wyznaczoną przez $p$
w rozdziale~10. Struktura $\mathcal U_p$
jest rozdzielająca wtedy i tylko wtedy, gdy
$p$ jest metryką.
\end{twierdzenie}
\begin{proof}
Każdy $E_\varepsilon$ zawiera przekątną
i jest symetryczny. Domykanie w górę wynika
z definicji. Dla $E_\varepsilon,E_\delta$
zachodzi $E_{\min\{\varepsilon,\delta\}}
\subseteq E_\varepsilon\cap E_\delta$,
co daje domknięcie na przecięcia. Wreszcie
$E_{\varepsilon/2}\circ E_{\varepsilon/2}
\subseteq E_\varepsilon$ z nierówności trójkąta.
Jeśli $E$ zawiera $E_\varepsilon$,
możemy więc wybrać $F=E_{\varepsilon/2}$
do aksjomatu złożenia.
Dla każdego $x$ mamy
$E_\varepsilon[x]=B_p(x,\varepsilon)$.
Są to bazy otoczeń w obu topologiach,
stąd $\tau_{\mathcal U_p}=\tau_p$.
Przecięcie wszystkich $E_\varepsilon$
jest zbiorem par $(x,y)$ o $p(x,y)=0$.
Z twierdzenia~\ref{thm:topog-jednostajne-hausdorff}
otrzymujemy ostatnią równoważność.
\end{proof}

\begin{figure}[htbp]
\centering
\begin{tikzpicture}[x=0.8cm,y=0.8cm,>=stealth]
\fill[blue!14] (0,0)--(.8,0)--(5,4.2)--(5,5)--(4.2,5)--(0,.8)--cycle;
\draw[gray!70] (0,0) rectangle (5,5);
\draw[very thick,blue!75!black] (0,0)--(5,5);
\draw[dashed,blue!65!black] (0,.8)--(4.2,5);
\draw[dashed,blue!65!black] (.8,0)--(5,4.2);
\draw[->] (0,0)--(5.6,0) node[below right] {$x$};
\draw[->] (0,0)--(0,5.6) node[above left] {$y$};
\node[blue!70!black,anchor=south east,font=\small] at (4.7,4.9)
 {$\Delta_X$};
\node[anchor=west,font=\small] at (5.65,3.35)
 {$E_\varepsilon=\{(x,y):|x-y|<\varepsilon\}$};
\draw[->,blue!65!black] (5.55,3.3)--(3.3,2.9);
\end{tikzpicture}
\caption{Schemat otoczenia jednostajnego przekątnej dla metryki
$|x-y|$ na przedziale $X=[0,5]$ i $\varepsilon=0{,}8$.
Zacieniowany pas zawiera przekątną. Linie przerywane
nie należą do relacji $E_\varepsilon$.}
\label{fig:topog-jednostajne-metryczny-pas}
\end{figure}

\begin{przyklad}[Ta sama topologia, różne struktury jednostajne]
\label{ex:topog-jednostajne-dwie-metryki}
Na $\R$ rozważmy metryki
$d(x,y)=|x-y|$ oraz
$e(x,y)=|\arctan x-\arctan y|$.
Ponieważ $\arctan$ jest homeomorfizmem
$\R$ na $(-\pi/2,\pi/2)$, obie metryki
wyznaczają zwykłą topologię prostej,
jak wyjaśniono w rozdziale~10.
Nie wyznaczają jednak tej samej struktury
jednostajnej. Istotnie, dla $x=n$ i $y=n+2$
mamy $d(x,y)=2$, podczas gdy
$e(n,n+2)\to0$. Dla żadnego $\delta>0$
nie zachodzi więc
$\{(x,y):e(x,y)<\delta\}
 \subseteq\{(x,y):d(x,y)<1\}$.
A zatem relacja $\{d<1\}$ należy do
$\mathcal U_d$, lecz nie do $\mathcal U_e$.
Ta różnica będzie ważna przy badaniu
zupełności w rozdziale~13.
\end{przyklad}

\section{Podprzestrzenie i iloczyny}
\label{sec:topog-jednostajne-konstrukcje}

\begin{twierdzenie}[Struktura jednostajna podprzestrzeni]
\label{thm:topog-jednostajne-podprzestrzen}
Jeśli $(X,\mathcal U)$ jest przestrzenią jednostajną,
a $A\subseteq X$, to rodzina
\[
 \mathcal U|_A=\{H\subseteq A\times A:
       \exists E\in\mathcal U\
          (E\cap(A\times A)\subseteq H)\}
\]
jest strukturą jednostajną na $A$.
Indukuje ona topologię podprzestrzeni
$(A,\tau_{\mathcal U}|_A)$.
\end{twierdzenie}
\begin{proof}
Przecięcia relacji z $A\times A$ zawierają
$\Delta_A$, zachowują odwracanie i przecięcia.
Jeśli $F\circ F\subseteq E$, to
$(F\cap(A\times A))\circ(F\cap(A\times A))
\subseteq E\cap(A\times A)$.
Daje to wszystkie aksjomaty struktury jednostajnej.

Przekroje bazowych relacji na $A$ mają postać
$(E\cap(A\times A))[a]=E[a]\cap A$.
Z twierdzenia~\ref{thm:topog-jednostajne-topologia}
są one otoczeniami $a$ w topologii względnej.
Każde względnie otwarte otoczenie $a$ ma
postać $A\cap O$, gdzie $O$ jest otwarty
w $X$; wybierając $E[a]\subseteq O$,
otrzymujemy bazowy przekrój zawarty w
$A\cap O$. Bazy otoczeń obu topologii są
więc równoważne.
\end{proof}

\begin{twierdzenie}[Produkt przestrzeni jednostajnych]
\label{thm:topog-jednostajne-iloczyn}
Niech $(X_i,\mathcal U_i)$, $i\in I$, będzie
rodziną przestrzeni jednostajnych i
$X=\prod_{i\in I}X_i$. Dla skończonego
$F\subseteq I$ oraz $E_i\in\mathcal U_i$,
$i\in F$, definiujemy
\[
 E(F,(E_i))=\{(x,y)\in X\times X:
     (x_i,y_i)\in E_i\text{ dla każdego }i\in F\}.
\]
Rodzina wszystkich nadzbiorów tych relacji
jest strukturą jednostajną $\mathcal U_\Pi$.
Jej topologia jest topologią iloczynu topologii
$\tau_{\mathcal U_i}$. Każde rzutowanie
$\pi_i:X\to X_i$ jest jednostajnie ciągłe
w sensie definicji podanej w następnej sekcji.
\end{twierdzenie}
\begin{proof}
Każda bazowa relacja zawiera przekątną.
Dla dwóch relacji bazowych warunek z
ich przecięcia uzyskujemy na sumie skończonych
zbiorów współrzędnych, biorąc przecięcie
odpowiednich relacji tam, gdzie indeksy się
powtarzają. Odwracanie wykonujemy
współrzędna po współrzędnej. Do warunku
złożenia wybieramy dla każdego $i\in F$
relację $H_i$ z $H_i\circ H_i\subseteq E_i$;
wówczas
$E(F,(H_i))\circ E(F,(H_i))
\subseteq E(F,(E_i))$.
Aksjomaty domykania w górę są zawarte
w samej definicji $\mathcal U_\Pi$.

Ustalmy $x\in X$. Przekroje bazowych relacji
mają postać
$\bigcap_{i\in F}\pi_i^{-1}(E_i[x_i])$.
Na mocy twierdzenia~\ref{thm:topog-jednostajne-topologia}
każdy $E_i[x_i]$ jest otoczeniem $x_i$;
zatem taki przekrój jest otoczeniem
$x$ w topologii iloczynu. Odwrotnie,
dla każdego bazowego otwartego cylindra
$\bigcap_{i\in F}\pi_i^{-1}(O_i)$
zawierającego $x$ istnieją
$E_i\in\mathcal U_i$ z $E_i[x_i]\subseteq O_i$.
Odpowiedni przekrój $E(F,(E_i))[x]$
mieści się w cylindrze. Obie topologie
mają więc takie same otoczenia.

Ostatnie twierdzenie wynika wprost z tego,
że dla dowolnego $E_i\in\mathcal U_i$
relacja $E(\{i\},(E_i))$ należy do
$\mathcal U_\Pi$ i jej pary mają
$i$-te współrzędne należące do $E_i$.
\end{proof}

\begin{uwaga}[Produkt pustej rodziny]
Dla $I=\varnothing$ iloczyn zbiorów
jest przestrzenią jednoelementową,
a jedyną bazową relacją w konstrukcji
jest $X\times X=\Delta_X$.
Jeśli któryś czynnik jest pusty,
cały iloczyn jest pusty i otrzymujemy
jedyną strukturę jednostajną na zbiorze pustym.
\end{uwaga}

\section{Odwzorowania jednostajnie ciągłe}
\label{sec:topog-jednostajne-ciaglosc}

\begin{definicja}[Ciągłość jednostajna]
\label{def:topog-jednostajne-ciaglosc}
Niech $(X,\mathcal U)$ i $(Y,\mathcal V)$
będą przestrzeniami jednostajnymi.
Odwzorowanie $f:X\to Y$ jest
\emph{jednostajnie ciągłe}, jeżeli
dla każdego $V\in\mathcal V$ istnieje
$E\in\mathcal U$ takie, że
$(x,x')\in E$ pociąga za sobą
$(f(x),f(x'))\in V$ dla wszystkich $x,x'\in X$.
Dla struktur metrycznych odpowiada to
zwykłemu warunkowi $\varepsilon$--$\delta$,
w którym $\delta$ nie zależy od punktu $x$.
\end{definicja}

\begin{twierdzenie}[Ciągłość i składanie]
\label{thm:topog-jednostajne-ciaglosc-wlasnosci}
Niech $X,Y,Z$ będą przestrzeniami jednostajnymi.
Każde odwzorowanie jednostajnie ciągłe
$f:X\to Y$ jest ciągłe w topologiach
indukowanych. Jeśli ponadto
$g:Y\to Z$ jest jednostajnie ciągłe,
to $g\circ f:X\to Z$ jest jednostajnie ciągłe.
\end{twierdzenie}
\begin{proof}
Ustalmy $x\in X$ oraz otwarte otoczenie
$O$ punktu $f(x)$ w $Y$. Istnieje
$V\in\mathcal V$ z $V[f(x)]\subseteq O$.
Wybieramy $E\in\mathcal U$ na mocy
ciągłości jednostajnej. Wówczas
$f(E[x])\subseteq V[f(x)]\subseteq O$.
Ponieważ $E[x]$ jest otoczeniem $x$,
ciągłość $f$ w punkcie $x$ wynika z
kryterium otoczeń.

Dla $W$ należącego do struktury na $Z$
wybieramy $V$ na $Y$, które $g$ przenosi
w $W$, a następnie $E$ na $X$, które
$f$ przenosi w $V$. Dla $(x,x')\in E$
para $(g(f(x)),g(f(x')))$ należy do $W$,
co jest dokładnie definicją.
\end{proof}

\begin{przyklad}[Odwzorowanie ciągłe, ale nie jednostajnie ciągłe]
\label{ex:topog-jednostajne-kwadrat}
Funkcja $f:\R\to\R$, $f(x)=x^2$, jest
ciągła w zwykłej topologii.
Nie jest jednostajnie ciągła dla
struktur wyznaczonych metryką euklidesową.
Rzeczywiście, weźmy $x_n=n$ i
$y_n=n+1/n$. Wtedy $|x_n-y_n|=1/n\to0$,
lecz $|f(y_n)-f(x_n)|=2+1/n^2\ge2$.
Dla $\varepsilon=1$ żadne $\delta>0$
nie spełnia warunku ciągłości jednostajnej.
\end{przyklad}

\begin{lemat}[Zwartość metryczna implikuje zwartość ciągową]
\label{lem:topog-jednostajne-zwarta-ciagowa}
Niech $(K,d)$ będzie zwartą przestrzenią metryczną.
Wówczas każdy ciąg $(x_n)_{n\ge1}$ w $K$ ma
podciąg zbieżny do pewnego punktu $x\in K$.
\end{lemat}
\begin{proof}
Dla każdego $N\ge1$ niech
$F_N=\overline{\{x_n:n\ge N\}}^{\,K}$.
Są to niepuste zbiory domknięte w $K$,
a ponadto $F_{N+1}\subseteq F_N$.
Każda skończona rodzina tych zbiorów ma
przecięcie równe jej najmniejszemu
elementowi w sensie inkluzji, a więc
niepuste. Z kryterium skończonych przecięć
dla zwartości, udowodnionego w
twierdzeniu~\ref{thm:topog-zwartosc-fip},
istnieje $x\in\bigcap_{N\ge1}F_N$.

Konstruujemy ściśle rosnący ciąg indeksów.
Przyjmijmy $n_0=0$. Jeśli wybrano
$n_{k-1}$, to $x\in F_{n_{k-1}+1}$.
Z definicji domknięcia punkt $x$
należy do domknięcia ogona
$\{x_n:n\ge n_{k-1}+1\}$, zatem kula
$B_d(x,1/k)$ zawiera pewien wyraz
$x_{n_k}$ tego ogona. Wybieramy
$n_k\ge n_{k-1}+1$ tak, aby
$d(x_{n_k},x)<1/k$. W rezultacie
$n_k>n_{k-1}$ i $x_{n_k}\to x$.
\end{proof}

\begin{twierdzenie}[Twierdzenie Heinego--Cantora]
\label{thm:topog-jednostajne-heine-cantor}
Niech $(K,d)$ będzie zwartą przestrzenią
metryczną, a $(Y,e)$ przestrzenią metryczną.
Każde ciągłe odwzorowanie $f:K\to Y$
jest jednostajnie ciągłe.
\end{twierdzenie}
\begin{proof}
Przypuśćmy przeciwnie, że istnieje
$\varepsilon_0>0$ takie, że dla każdego
$n\ge1$ można wybrać $x_n,y_n\in K$
z $d(x_n,y_n)<1/n$ oraz
$e(f(x_n),f(y_n))\ge\varepsilon_0$.
Na mocy lematu~\ref{lem:topog-jednostajne-zwarta-ciagowa}
ciąg $(x_n)$ ma podciąg $(x_{n_k})$
zbieżny do $x\in K$.
Z nierówności
$d(y_{n_k},x)\le d(y_{n_k},x_{n_k})
  +d(x_{n_k},x)$
wynika, że również $y_{n_k}\to x$.
Ciągłość $f$ w $x$ daje
$f(x_{n_k})\to f(x)$ i $f(y_{n_k})\to f(x)$.
Nierówność trójkąta w $Y$ prowadzi do
$e(f(x_{n_k}),f(y_{n_k}))\to0$,
wbrew wyborowi $\varepsilon_0$.
\end{proof}

\section{Ciągi, sieci i filtry Cauchy'ego}
\label{sec:topog-jednostajne-cauchy}

\begin{definicja}[Sieć Cauchy'ego]
\label{def:topog-jednostajne-siec-cauchyego}
Niech $(X,\mathcal U)$ będzie przestrzenią
jednostajną, a $D$ niepustym zbiorem skierowanym.
Sieć $(x_\alpha)_{\alpha\in D}$ jest
\emph{siecią Cauchy'ego}, jeśli dla każdego
$E\in\mathcal U$ istnieje $\alpha_0\in D$
takie, że $(x_\alpha,x_\beta)\in E$
dla wszystkich $\alpha,\beta\ge\alpha_0$.
Jeśli $D=\mathbb N$, mówimy o \emph{ciągu Cauchy'ego}.
\end{definicja}

\begin{definicja}[Filtr Cauchy'ego]
\label{def:topog-jednostajne-filtr-cauchyego}
Filtr $\mathcal F$ na $X$ nazywamy
\emph{filtrem Cauchy'ego}, jeżeli
dla każdego $E\in\mathcal U$ istnieje
$A\in\mathcal F$ takie, że $A\times A\subseteq E$.
\end{definicja}

\begin{twierdzenie}[Sieć i filtr ogonów]
\label{thm:topog-jednostajne-ogony}
Niech $(x_\alpha)_{\alpha\in D}$ będzie
siecią w przestrzeni jednostajnej $(X,\mathcal U)$.
Niech $\mathcal F_x$ będzie filtrem
generowanym przez zbiory ogonowe
$A_{\alpha_0}=\{x_\alpha:\alpha\ge\alpha_0\}$.
Wtedy sieć jest Cauchy'ego wtedy i tylko wtedy,
gdy $\mathcal F_x$ jest filtrem Cauchy'ego.
\end{twierdzenie}
\begin{proof}
Zbiory $A_{\alpha_0}$ są niepuste.
Skierowanie $D$ zapewnia, że dwa ogony
zawierają wspólny mniejszy ogon;
tworzą więc bazę filtra.
Jeśli sieć jest Cauchy'ego, dla każdego
$E\in\mathcal U$ pewien ogon spełnia
$A_{\alpha_0}\times A_{\alpha_0}\subseteq E$.
Jest on elementem $\mathcal F_x$.

Odwrotnie, niech dla $E\in\mathcal U$
istnieje $A\in\mathcal F_x$ z
$A\times A\subseteq E$. Z definicji
filtra generowanego przez ogony istnieje
$\alpha_0$ takie, że $A_{\alpha_0}\subseteq A$.
Zatem dla $\alpha,\beta\ge\alpha_0$
para $(x_\alpha,x_\beta)$ należy do $E$.
\end{proof}

\begin{twierdzenie}[Zbieżność implikuje warunek Cauchy'ego]
\label{thm:topog-jednostajne-zbiezna-cauchy}
W każdej przestrzeni jednostajnej
każda sieć zbieżna w topologii
$\tau_{\mathcal U}$ jest siecią Cauchy'ego.
\end{twierdzenie}
\begin{proof}
Załóżmy $x_\alpha\to x$. Ustalmy
$E\in\mathcal U$. Na mocy
lematu~\ref{lem:topog-jednostajne-symetryczne}
istnieje symetryczne $F\in\mathcal U$
z $F\circ F\subseteq E$.
Z twierdzenia~\ref{thm:topog-jednostajne-topologia}
zbiór $F[x]$ jest otoczeniem $x$,
więc dla wszystkich $\alpha\ge\alpha_0$
mamy $x_\alpha\in F[x]$.
Jeśli $\alpha,\beta\ge\alpha_0$, to
$(x_\alpha,x),(x,x_\beta)\in F$
dzięki symetrii. Stąd
$(x_\alpha,x_\beta)\in F\circ F\subseteq E$.
\end{proof}

\begin{stwierdzenie}[Ciągłość jednostajna zachowuje sieci Cauchy'ego]
\label{prop:topog-jednostajne-f-cauchy}
Jeżeli $f:(X,\mathcal U)\to(Y,\mathcal V)$
jest jednostajnie ciągłe, a
$(x_\alpha)$ jest siecią Cauchy'ego w $X$,
to $(f(x_\alpha))$ jest siecią Cauchy'ego w $Y$.
\end{stwierdzenie}
\begin{proof}
Dla $V\in\mathcal V$ wybieramy
$E\in\mathcal U$ takie, że
$(f\times f)(E)\subseteq V$.
Od pewnego indeksu wszystkie pary
$(x_\alpha,x_\beta)$ należą do $E$;
ich obrazy należą więc do $V$.
\end{proof}

\begin{twierdzenie}[Sieci Cauchy'ego w produkcie]
\label{thm:topog-jednostajne-produkt-cauchy}
Niech $X=\prod_{i\in I}X_i$ ma strukturę
jednostajną produktu z
twierdzenia~\ref{thm:topog-jednostajne-iloczyn}.
Sieć $(x_\alpha)$ w $X$ jest Cauchy'ego
wtedy i tylko wtedy, gdy dla każdego $i\in I$
sieć współrzędnych $(\pi_i(x_\alpha))$
jest Cauchy'ego w $X_i$.
\end{twierdzenie}
\begin{proof}
Jeśli $(x_\alpha)$ jest Cauchy'ego,
rzutowania są jednostajnie ciągłe,
więc ich sieci obrazów są Cauchy'ego
na mocy poprzedniego stwierdzenia.
Odwrotnie, ustalmy bazowe otoczenie
$E(F,(E_i))$ produktu, gdzie $F$ jest
skończony. Dla każdego $i\in F$ istnieje
indeks $\alpha_i$, od którego wszystkie
pary odpowiednich współrzędnych leżą
w $E_i$. Skierowanie zbioru indeksów
pozwala znaleźć wspólną górną granicę
skończonej rodziny $(\alpha_i)_{i\in F}$.
Od tego indeksu wszystkie pary
$(x_\alpha,x_\beta)$ należą do
$E(F,(E_i))$. Gdy $F=\varnothing$,
warunek jest automatycznie spełniony.
Każda relacja produktu zawiera bazową,
co kończy dowód.
\end{proof}

\section{Kontrprzykłady i rozróżnienia}
\label{sec:topog-jednostajne-kontrprzyklady}

\begin{przyklad}[Sieć Cauchy'ego bez granicy]
\label{ex:topog-jednostajne-cauchy-bez-granicy}
Ciąg $x_n=1/(n+1)$ w $(0,1)$ ze
strukturą indukowaną przez metrykę
euklidesową jest Cauchy'ego, gdyż
jest Cauchy'ego w $\R$. Nie ma jednak
granicy w $(0,1)$: jedyną granicą
w $\R$ jest $0$. Zatem odwrotność
twierdzenia~\ref{thm:topog-jednostajne-zbiezna-cauchy}
jest fałszywa.
\end{przyklad}

\begin{przyklad}[Zbieżność bez jednoznacznej granicy]
\label{ex:topog-jednostajne-niehausdorff}
Jeśli $X$ ma co najmniej dwa punkty
i strukturę niedyskretną, każda sieć
w $X$ zbiega do każdego punktu.
Istotnie, jedynymi otwartymi zbiorami
są $\varnothing$ i $X$. Każda sieć
jest również Cauchy'ego, ponieważ
jedyną relacją jednostajną jest $X\times X$.
Nie można zatem oczekiwać jednoznaczności
granicy bez założeń oddzielania.
\end{przyklad}

\begin{przyklad}[Topologia nie określa ciągów Cauchy'ego]
\label{ex:topog-jednostajne-cauchy-dwie-metryki}
W przykładzie~\ref{ex:topog-jednostajne-dwie-metryki}
metryki $d$ i $e$ mają tę samą topologię na $\R$,
lecz ciąg $x_n=n$ jest Cauchy'ego dla $e$,
ponieważ ciąg $\arctan(n)$ zbiega w $\R$.
Nie jest Cauchy'ego dla $d$, gdyż
$d(n,2n)=n$ dla $n\ge1$. Nie ma również
granicy w zwykłej topologii $\R$.
Zbieżność jest więc pojęciem topologicznym,
ale warunek Cauchy'ego zależy od struktury
jednostajnej.
\end{przyklad}

\begin{uwaga}[Przejście do rozdziału następnego]
W przestrzeniach jednostajnych zupełność
będziemy formułować za pomocą filtrów
lub sieci Cauchy'ego. Rozdział~13
wyjaśni, jak te pojęcia wiążą się
z zupełnością metryczną z rozdziału~11,
i skonstruuje uzupełnienia.
\end{uwaga}

\chapter{Zupełność jednostajna i uzupełnienia}
\label{ch:topog-uzupelnienia}

\noindent\textbf{Cel rozdziału.}
W rozdziale~12 wprowadziliśmy sieci i filtry Cauchy'ego,
a w rozdziale~11 zupełność przestrzeni metrycznych.
Wyjaśnimy ich związek, skonstruujemy uzupełnienie
przestrzeni jednostajnej rozdzielającej oraz
udowodnimy jednoznaczność i własność uniwersalną tej konstrukcji.
Oddzielimy zupełność ustalonej struktury jednostajnej
od zupełnej metryzowalności przestrzeni topologicznej.

\section{Zupełność w języku filtrów i sieci}
\label{sec:topog-uzupelnienia-definicje}

\begin{definicja}[Zbieżność filtra i zupełność jednostajna]
\label{def:topog-uzupelnienia-zupelnosc}
Niech $(X,\mathcal U)$ będzie przestrzenią jednostajną.
Filtr $\mathcal F$ na $X$ jest \emph{zbieżny do $x\in X$},
gdy każde otoczenie $x$ w topologii $\tau_{\mathcal U}$
należy do $\mathcal F$.
Przestrzeń $(X,\mathcal U)$ jest \emph{zupełna},
jeżeli każdy filtr Cauchy'ego na $X$ jest zbieżny
do co najmniej jednego punktu $X$.
\end{definicja}

\begin{uwaga}[Rozdzielanie i przestrzeń pusta]
Bez założenia Hausdorffa granica filtra nie musi być
jednoznaczna. Na przestrzeni pustej nie ma filtrów
właściwych, zatem jest ona zupełna.
\end{uwaga}

\begin{twierdzenie}[Równoważność filtrów i sieci Cauchy'ego]
\label{thm:topog-uzupelnienia-filtry-sieci}
Przestrzeń jednostajna $(X,\mathcal U)$ jest zupełna
wtedy i tylko wtedy, gdy każda jej sieć Cauchy'ego
jest zbieżna w topologii $\tau_{\mathcal U}$.
\end{twierdzenie}
\begin{proof}
Załóżmy, że wszystkie filtry Cauchy'ego są zbieżne.
Jeśli $(x_\alpha)_{\alpha\in D}$ jest siecią Cauchy'ego,
to jej filtr ogonów jest Cauchy'ego na mocy
twierdzenia o filtrze ogonów z rozdziału~12.
Niech zbiega do $x$. Dla każdego otoczenia $O$
punktu $x$ istnieje ogon zawarty w $O$,
więc $x_\alpha\to x$.

Załóżmy teraz, że każda sieć Cauchy'ego jest zbieżna.
Niech $\mathcal F$ będzie filtrem Cauchy'ego.
Rozważmy zbiór skierowany
$D=\{(A,a):A\in\mathcal F,\ a\in A\}$,
z porządkiem $(A,a)\preceq(B,b)$ wtedy i tylko wtedy,
gdy $B\subseteq A$. Jest on niepusty, ponieważ
filtry są właściwe, i skierowany: dla
$(A,a),(B,b)$ wybieramy $c\in A\cap B$
oraz $(A\cap B,c)$. Definiujemy sieć
$x_{(A,a)}=a$.
Dla dowolnego $E\in\mathcal U$ istnieje
$A_0\in\mathcal F$ z $A_0\times A_0\subseteq E$.
Od indeksu $(A_0,a_0)$ wszystkie wyrazy
sieci leżą w $A_0$, więc sieć jest Cauchy'ego.
Niech zbiega do $x$.
Jeżeli $O$ jest otoczeniem $x$, to istnieje
$(A,a)\in D$ takie, że wszystkie późniejsze
wyrazy leżą w $O$. Dla każdego $b\in A$
indeks $(A,b)$ jest późniejszy od $(A,a)$,
a zatem $b\in O$. Stąd $A\subseteq O$
i $O\in\mathcal F$. Filtr $\mathcal F$ zbiega do $x$.
\end{proof}

\begin{stwierdzenie}[Zupełność metryczna i jednostajna]
\label{prop:topog-uzupelnienia-metryczna}
Niech $(X,d)$ będzie przestrzenią metryczną,
a $\mathcal U_d$ strukturą jednostajną wyznaczoną
przez $d$. Wtedy $(X,d)$ jest zupełna wtedy
i tylko wtedy, gdy $(X,\mathcal U_d)$ jest zupełna.
\end{stwierdzenie}
\begin{proof}
Jeżeli $(X,\mathcal U_d)$ jest zupełna, to
każdy ciąg Cauchy'ego jest siecią Cauchy'ego
w $\mathcal U_d$, zatem zbiega w topologii metrycznej.

Odwrotnie, niech $(X,d)$ będzie zupełna, a
$\mathcal F$ filtrem Cauchy'ego. Dla każdego
$n\ge1$ wybieramy $A_n\in\mathcal F$ takie, że
$d(a,b)<2^{-n}$ dla wszystkich $a,b\in A_n$.
Zastępujemy $A_n$ przez $\bigcap_{k=1}^n A_k$;
wówczas $A_{n+1}\subseteq A_n$ i $A_n\ne\varnothing$.
Wybierzmy $x_n\in A_n$. Gdy $m\ge n$,
punkty $x_m,x_n$ należą do $A_n$,
więc $d(x_m,x_n)<2^{-n}$.
Z zupełności $x_n\to x\in X$.
Dla $\varepsilon>0$ wybieramy $n$ tak,
że $2^{-n}<\varepsilon/2$ i
$d(x_n,x)<\varepsilon/2$.
Wtedy $A_n\subseteq B_d(x,\varepsilon)$,
więc kula ta należy do $\mathcal F$.
To dowodzi zbieżności filtra do $x$.
\end{proof}

\begin{przyklad}[Jednakowa topologia, inna zupełność]
\label{ex:topog-uzupelnienia-arctan}
Na $\R$ metryki $d(x,y)=|x-y|$ oraz
$e(x,y)=|\arctan x-\arctan y|$ wyznaczają tę
samą topologię (rozdział~12).
Przestrzeń $(\R,d)$ jest zupełna, lecz
$(\R,e)$ nie jest: ciąg $x_n=n$ jest
Cauchy'ego dla $e$, natomiast nie ma granicy
w zwykłej topologii $\R$.
Zupełność dotyczy struktury jednostajnej,
a nie samej topologii.
\end{przyklad}

\section{Uzupełnienie przestrzeni metrycznej}
\label{sec:topog-uzupelnienia-metryczne}

\begin{definicja}[Uzupełnienie metryczne]
\label{def:topog-uzupelnienia-metryczne}
\emph{Uzupełnieniem metrycznym} przestrzeni
metrycznej $(X,d)$ nazywamy parę
$(\widehat X,\iota)$, gdzie
$(\widehat X,\widehat d)$ jest zupełną przestrzenią
metryczną, a $\iota:(X,d)\to(\widehat X,\widehat d)$
jest izometrią o gęstym obrazie.
\end{definicja}

\begin{twierdzenie}[Konstrukcja uzupełnienia metrycznego]
\label{thm:topog-uzupelnienia-metryczne}
Każda przestrzeń metryczna $(X,d)$ ma uzupełnienie
metryczne. Jeżeli $X\ne\varnothing$,
można je zbudować z klas ciągów Cauchy'ego w $X$.
\end{twierdzenie}
\begin{proof}
Dla $X=\varnothing$ przyjmujemy $\widehat X=\varnothing$
i jedyne możliwe odwzorowanie $\iota$.
Niech dalej $X\ne\varnothing$. Oznaczmy przez
$\mathcal C(X)$ zbiór ciągów Cauchy'ego w $X$
i połóżmy $(x_n)\sim(y_n)$, gdy
$d(x_n,y_n)\to0$. Zwrotność i symetria są
bezpośrednie, a przechodniość wynika z
$d(x_n,z_n)\le d(x_n,y_n)+d(y_n,z_n)$.
Definiujemy $\widehat X=\mathcal C(X)/{\sim}$.

Dla ciągów Cauchy'ego $(x_n),(y_n)$
ciąg liczb $d(x_n,y_n)$ jest Cauchy'ego, gdyż
\[
 |d(x_n,y_n)-d(x_m,y_m)|
 \le d(x_n,x_m)+d(y_n,y_m).
\]
Istnieje zatem granica; definiujemy
$\widehat d([(x_n)],[(y_n)])=\lim_nd(x_n,y_n)$.
Zmiana przedstawicieli nie zmienia granicy,
ponieważ różnica odległości jest nie większa
niż suma odległości między odpowiednimi
przedstawicielami. Symetria i nierówność
trójkąta przechodzą na granicę.
Jeżeli $\widehat d=0$, ciągi są równoważne,
więc $\widehat d$ jest metryką.

Odwzorowanie $\iota(x)=[(x,x,\ldots)]$
jest izometrią. Dla $z=[(x_n)]$ mamy
\[
 \widehat d(\iota(x_n),z)=\lim_{m\to\infty}d(x_n,x_m)
 \longrightarrow0.
\]
Ostatnie przejście wynika wprost z warunku
Cauchy'ego: dla dużego $n$ wszystkie
$d(x_n,x_m)$ z $m$ dostatecznie dużym
są mniejsze od zadanego $\varepsilon$.
Stąd $\iota(X)$ jest gęsty w $\widehat X$.

Niech $(z_k)$ będzie ciągiem Cauchy'ego
w $(\widehat X,\widehat d)$.
Z gęstości wybieramy $x_k\in X$ z
$\widehat d(z_k,\iota(x_k))<2^{-k}$.
Wtedy
\[
 d(x_k,x_l)\le
 2^{-k}+\widehat d(z_k,z_l)+2^{-l},
\]
więc $(x_k)$ jest ciągiem Cauchy'ego.
Niech $z=[(x_k)]$. Z udowodnionej gęstości
$\iota(x_k)\to z$; nierówność
$\widehat d(z_k,z)\le
\widehat d(z_k,\iota(x_k))+
\widehat d(\iota(x_k),z)$
daje $z_k\to z$. Uzupełnienie jest zupełne.
\end{proof}

\begin{przyklad}[Uzupełnienie liczb wymiernych]
\label{ex:topog-uzupelnienia-q}
Zbiór $\Q$ ze zwykłą metryką jest gęsty w
zupełnej przestrzeni $\R$. Dlatego
$(\R,|\cdot|)$ jest uzupełnieniem metrycznym
$(\Q,|\cdot|)$. Klasa ciągu wymiernych
przybliżeń liczby $\sqrt2$ odpowiada
punktowi $\sqrt2\in\R$.
\end{przyklad}

\section{Filtry Cauchy'ego i minimalność}
\label{sec:topog-uzupelnienia-filtry}

\begin{definicja}[Filtry bliskie]
\label{def:topog-uzupelnienia-bliskie}
Niech $(X,\mathcal U)$ będzie przestrzenią
jednostajną. Dwa filtry Cauchy'ego
$\mathcal F,\mathcal G$ na $X$ nazywamy
\emph{bliskimi}, i piszemy $\mathcal F\sim\mathcal G$,
gdy dla każdego $E\in\mathcal U$
istnieją $A\in\mathcal F$ i $B\in\mathcal G$
takie, że $A\times B\subseteq E$.
\end{definicja}

\begin{lemat}[Relacja bliskości jest równoważnością]
\label{lem:topog-uzupelnienia-rownowaznosc}
Relacja $\sim$ jest relacją równoważności
na zbiorze filtrów Cauchy'ego w $X$.
\end{lemat}
\begin{proof}
Zwrotność wynika z warunku Cauchy'ego:
dla każdego $E$ istnieje $A\in\mathcal F$
z $A\times A\subseteq E$.
Symetria: dla $E$ wybieramy $E\cap E^{-1}$
i zamieniamy role obu filtrów.

Dla przechodniości wybierzmy $E\in\mathcal U$
i symetryczne $D\in\mathcal U$ z
$D\circ D\subseteq E$. Jeżeli
$\mathcal F\sim\mathcal G$ oraz
$\mathcal G\sim\mathcal H$, to istnieją
$A\in\mathcal F$, $B,B'\in\mathcal G$,
$C\in\mathcal H$ z $A\times B\subseteq D$
i $B'\times C\subseteq D$.
Z właściwości filtra wybieramy
$b\in B\cap B'$. Dla wszystkich
$a\in A,c\in C$ mamy $(a,b),(b,c)\in D$,
więc $(a,c)\in D\circ D\subseteq E$.
Stąd $\mathcal F\sim\mathcal H$.
\end{proof}

\begin{definicja}[Filtr Cauchy'ego minimalny]
\label{def:topog-uzupelnienia-minimalny}
Filtr Cauchy'ego $\mathcal F$ jest
\emph{minimalny}, gdy nie istnieje
ściśle mniejszy w sensie inkluzji
filtr Cauchy'ego $\mathcal G\subsetneq\mathcal F$.
\end{definicja}

\begin{lemat}[Minimalny przedstawiciel]
\label{lem:topog-uzupelnienia-minimalizacja}
Każdy filtr Cauchy'ego $\mathcal F$ ma
dokładnie jeden minimalny filtr Cauchy'ego
bliski $\mathcal F$. Jest nim filtr
$\mathcal F^\circ$ generowany przez zbiory
$E[A]=\bigcup_{a\in A}E[a]$, gdzie
$E\in\mathcal U$ i $A\in\mathcal F$.
\end{lemat}
\begin{proof}
Zbiory $E[A]$ są niepuste. Ich rodzina
jest bazą filtra, ponieważ
$(E\cap D)[A\cap B]\subseteq E[A]\cap D[B]$.
Ponadto $A\subseteq E[A]$, zatem
$\mathcal F^\circ\subseteq\mathcal F$.

Ustalmy $V\in\mathcal U$. Wybierzmy
symetryczne $D\in\mathcal U$ z
$D\circ D\circ D\subseteq V$:
najpierw wybieramy otoczenie o kwadracie
zawartym w $V$, potem symetryczne otoczenie
o kwadracie zawartym w poprzednim.
Z warunku Cauchy'ego istnieje
$A\in\mathcal F$ z $A\times A\subseteq D$.
Jeżeli $u,v\in D[A]$, wybieramy
$a,b\in A$ z $(a,u),(b,v)\in D$.
Z symetrii $(u,a),(a,b),(b,v)\in D$,
stąd $(u,v)\in D^3\subseteq V$.
Zatem $\mathcal F^\circ$ jest Cauchy'ego.

Aby wykazać bliskość $\mathcal F^\circ$
i $\mathcal F$, dla $V$ wybieramy
symetryczne $D$ z $D^2\subseteq V$
oraz $A\in\mathcal F$ z $A\times A\subseteq D$.
Wtedy $D[A]\in\mathcal F^\circ$
i $D[A]\times A\subseteq D^2\subseteq V$.

Niech teraz $\mathcal G$ będzie dowolnym
filtrem Cauchy'ego bliskim $\mathcal F$.
Pokażemy $\mathcal F^\circ\subseteq\mathcal G$.
Dla bazowego $E[A]$, gdzie
$A\in\mathcal F$, wybieramy symetryczne
$D\subseteq E$. Bliskość zapewnia
$A'\in\mathcal F$ i $B\in\mathcal G$
z $A'\times B\subseteq D$.
Biorąc $a\in A\cap A'$ otrzymujemy
$B\subseteq D[a]\subseteq E[A]$,
więc $E[A]\in\mathcal G$.
Wobec tego $\mathcal F^\circ$
zawiera się w każdym bliskim filtrze Cauchy'ego.

W szczególności, jeśli $\mathcal H$
jest filtrem Cauchy'ego zawartym w
$\mathcal F^\circ$, to jest bliski
$\mathcal F^\circ$: dla $V$ weźmy
$A\in\mathcal H$ z $A\times A\subseteq V$;
ponieważ $\mathcal H\subseteq\mathcal F^\circ$,
mamy również $A\in\mathcal F^\circ$.
Z już wykazanej minimalności względem
filtrów bliskich wynika
$\mathcal F^\circ\subseteq\mathcal H$.
Stąd równość, czyli minimalność.
Jeśli dwa minimalne filtry są bliskie,
oba muszą się zawierać jeden w drugim,
więc są równe.
\end{proof}

\section{Konstrukcja uzupełnienia jednostajnego}
\label{sec:topog-uzupelnienia-konstrukcja}

\begin{twierdzenie}[Istnienie uzupełnienia jednostajnego]
\label{thm:topog-uzupelnienia-istnienie}
Niech $(X,\mathcal U)$ będzie rozdzielającą
przestrzenią jednostajną. Istnieje zupełna,
rozdzielająca przestrzeń jednostajna
$(\widehat X,\widehat{\mathcal U})$
oraz jednostajnie ciągłe odwzorowanie
$\iota:X\to\widehat X$ będące
izomorfizmem jednostajnym na swój gęsty obraz.
\end{twierdzenie}
\begin{proof}
Niech $\widehat X$ będzie zbiorem wszystkich
minimalnych filtrów Cauchy'ego na $X$.
Dla $E\in\mathcal U$ definiujemy relację
\[
 \widehat E=\{(\mathcal F,\mathcal G)\in
 \widehat X^2:\exists A\in\mathcal F,\ B\in\mathcal G,\
 A\times B\subseteq E\}.
\]
Relacje $\widehat E$ tworzą bazę struktury jednostajnej.
Istotnie, każda zawiera przekątną przez warunek
Cauchy'ego. Dla przecięć: wybierzmy symetryczne
$D$ z $D^3\subseteq E\cap F$.
Jeśli $(\mathcal A,\mathcal B)\in\widehat D$,
to należy także do $\widehat E\cap\widehat F$.
Odwrotność $\widehat E$ jest równa
$\widehat{E^{-1}}$. Gdy $D$ jest symetryczne
i $D^2\subseteq E$, mamy
$\widehat D\circ\widehat D\subseteq\widehat E$:
z $A\times B\subseteq D$ oraz
$B'\times C\subseteq D$ i
$b\in B\cap B'$ otrzymujemy
$A\times C\subseteq D^2\subseteq E$.
Domykanie relacji w górę daje
$\widehat{\mathcal U}$.

Jeśli $(\mathcal F,\mathcal G)$ należy do
wszystkich $\widehat E$, to
$\mathcal F\sim\mathcal G$, a więc
$\mathcal F=\mathcal G$ na mocy
lematu~\ref{lem:topog-uzupelnienia-minimalizacja}.
Struktura jest rozdzielająca.

Dla $x\in X$ filtr otoczeń $\mathcal N_x$
w $\tau_{\mathcal U}$ jest filtrem Cauchy'ego:
dla każdego $E$ wybieramy symetryczne $D$
z $D^2\subseteq E$, a
$D[x]\times D[x]\subseteq E$.
Ponadto jest minimalny: dla dowolnego
$A\in\mathcal N_x$ wybieramy $E$ z
$E[x]\subseteq A$. Ponieważ
$\{x\}\in\mathcal N_x$ nie musi zachodzić,
zamiast tego z twierdzenia o topologii
indukowanej bierzemy mniejsze $D[x]$
takie, że $D^2[x]\subseteq A$.
Wówczas $D[D[x]]\subseteq A$, więc
$A\in\mathcal N_x^\circ$ i
$\mathcal N_x^\circ=\mathcal N_x$.
Kładziemy $\iota(x)=\mathcal N_x$.
Ponieważ $X$ jest rozdzielająca,
różne punkty mają różne filtry otoczeń.

Pokażemy zgodność struktur na obrazie.
Jeżeli $(\iota(x),\iota(y))\in\widehat E$,
to istnieją otoczenia $A$ punktu $x$
i $B$ punktu $y$ z $A\times B\subseteq E$.
W szczególności $(x,y)\in E$.
Odwrotnie, jeśli $(x,y)\in D$
dla symetrycznego $D$ z $D^3\subseteq E$,
to $D[x]\times D[y]\subseteq E$;
obie sekcje są otoczeniami punktów,
więc $(\iota(x),\iota(y))\in\widehat E$.
To wykazuje jednostajną ciągłość
$\iota$ i jej odwrotności na obraz.

Gęstość obrazu: niech
$\mathcal F\in\widehat X$ i $E\in\mathcal U$.
Wybieramy symetryczne $D$ z $D^2\subseteq E$
oraz niepusty $A\in\mathcal F$
z $A\times A\subseteq D$.
Dla dowolnego $x\in A$ mamy
$A\times D[x]\subseteq D^2\subseteq E$,
a $D[x]\in\mathcal N_x$.
Zatem $(\mathcal F,\iota(x))\in\widehat E$.
Każde otoczenie $\mathcal F$ w $\widehat X$
spotyka $\iota(X)$.

Pozostaje zupełność. Niech $\mathcal H$
będzie filtrem Cauchy'ego na $\widehat X$.
Z gęstości $\iota(X)$ zbudujemy filtr
$\mathcal G$ na $X$ o bazie złożonej ze zbiorów
\[
 B(H,E)=\{x\in X:\exists \mathcal F\in H,\
       (\iota(x),\mathcal F)\in\widehat E\},
 \qquad H\in\mathcal H,\quad E\in\mathcal U.
\]
Są one niepuste: dla $\mathcal F\in H$
gęstość $\iota(X)$ zapewnia punkt
$\iota(x)\in\widehat E[\mathcal F]$.
Mamy
$B(H\cap H',E\cap E')\subseteq
 B(H,E)\cap B(H',E')$,
więc jest to baza właściwego filtra.

Dla $V\in\mathcal U$ wybieramy symetryczne
$D$ z $D^3\subseteq V$ i
$H\in\mathcal H$ o $H\times H\subseteq\widehat D$.
Jeżeli $x,y\in B(H,D)$, wybieramy
$\mathcal F,\mathcal G\in H$ z
$(\iota(x),\mathcal F),(\iota(y),\mathcal G)
\in\widehat D$.
Wówczas
$(\iota(x),\iota(y))\in\widehat{D^3}$,
a zgodność struktury na obrazie daje
$(x,y)\in D^3\subseteq V$.
Tak więc $\mathcal G$ jest filtrem Cauchy'ego.

Niech $\mathcal F_0=\mathcal G^\circ$
będzie jego minimalnym bliskim filtrem.
Pokażemy, że $\mathcal H$ zbiega do
$\mathcal F_0$. Dla symetrycznego
$E\in\mathcal U$ wybieramy $D$ z $D^4\subseteq E$.
Istnieją $A\in\mathcal G$ i $C\in\mathcal F_0$
z $A\times C\subseteq D$.
Ponadto możemy wybrać bazowy
$B(H,D)\subseteq A$.
Dla $\mathcal K\in H$ wybieramy
$x\in B(H,D)$ oraz
$\mathcal L\in H$ z
$(\iota(x),\mathcal L)\in\widehat D$.
Ponieważ $\mathcal H$ jest Cauchy'ego,
po zmniejszeniu $H$ na początku konstrukcji
możemy wymagać
$H\times H\subseteq\widehat D$,
więc $(\mathcal L,\mathcal K)\in\widehat D$.
Z $x\in A$ i $A\times C\subseteq D$
oraz definicji filtrów otoczeń wynika
$(\iota(x),\mathcal F_0)\in\widehat{D^2}$.
Przez złożenie otrzymujemy
$(\mathcal K,\mathcal F_0)\in\widehat{D^4}
\subseteq\widehat E$.
Zatem $H\subseteq\widehat E[\mathcal F_0]$;
ponieważ $\widehat E[\mathcal F_0]$ jest
otoczeniem, $\mathcal H\to\mathcal F_0$.
\end{proof}

\begin{uwaga}[Przestrzenie nierozdzielające]
Jeżeli $\mathcal U$ nie jest rozdzielająca,
to relacja $x\approx y$ dana warunkiem
$(x,y)\in\bigcap_{E\in\mathcal U}E$
jest relacją równoważności. Na ilorazie
otrzymuje się rozdzielającą strukturę
jednostajną, którą można uzupełnić.
Nie jest możliwe zanurzenie dowolnej
przestrzeni nierozdzielającej w rozdzielającą
przestrzeń jednostajną bez uprzedniego
utożsamienia punktów nierozróżnialnych.
\end{uwaga}

\section{Własność uniwersalna i jednoznaczność}
\label{sec:topog-uzupelnienia-uniwersalnosc}

\begin{lemat}[Rozszerzanie przez gęstość]
\label{lem:topog-uzupelnienia-rozszerzenie}
Niech $A$ będzie gęstą podprzestrzenią
jednostajną przestrzeni $Z$, a $Y$ zupełną
rozdzielającą przestrzenią jednostajną.
Każde jednostajnie ciągłe
$f:A\to Y$ ma dokładnie jedno
jednostajnie ciągłe rozszerzenie
$\overline f:Z\to Y$.
\end{lemat}
\begin{proof}
Dla $z\in Z$ oznaczmy przez $\mathcal F_z$
filtr na $A$ generowany przez
$A\cap E[z]$, gdzie $E$ przebiega
otoczenia jednostajne w $Z$.
Każdy taki zbiór jest niepusty,
gdyż $E[z]$ zawiera otwarte otoczenie
$z$, a $A$ jest gęsty. Filtr
$\mathcal F_z$ jest Cauchy'ego:
dla danego $V$ wybieramy symetryczne
$D$ z $D^2\subseteq V$;
wówczas $(D[z]\cap A)^2\subseteq V$.
Jednostajna ciągłość $f$ oznacza,
że filtr obrazów $f_*(\mathcal F_z)$
jest Cauchy'ego w $Y$. Z zupełności
zbiega do punktu $\overline f(z)$;
punkt ten jest jedyny, gdyż $Y$
jest Hausdorffa. Dla $z\in A$
filtr $\mathcal F_z$ zbiega do $z$,
więc $\overline f(z)=f(z)$.

Dowiedziemy jednostajnej ciągłości.
Ustalmy symetryczne $V$ w $Y$ i wybierzmy
symetryczne $W$ z $W^3\subseteq V$.
Z jednostajnej ciągłości $f$ istnieje
symetryczne $D$ w $Z$ takie, że
$(a,b)\in D\cap(A\times A)$ pociąga
$(f(a),f(b))\in W$.
Wybierzmy symetryczne $E$ z $E^3\subseteq D$.
Jeżeli $(z,z')\in E$, to dla
$a\in A\cap E[z]$ i $b\in A\cap E[z']$
mamy $(a,b)\in E^3\subseteq D$,
więc $(f(a),f(b))\in W$.
Przechodząc do granic filtrów w
topologii produktu $Y\times Y$,
otrzymujemy
$(\overline f(z),\overline f(z'))\in
\overline W^{\,Y\times Y}$.
Możemy na początku wybrać symetryczne $W$
takie, że $\overline W^{\,Y\times Y}\subseteq V$:
wystarczy najpierw wziąć symetryczne $T$
z $T^3\subseteq V$.
Jeśli $(y,y')\in\overline T^{\,Y\times Y}$,
to otwarte otoczenia punktów $y,y'$
zawarte odpowiednio w $T[y]$ i $T[y']$
wyznaczają otoczenie $(y,y')$ spotykające $T$.
Istnieją więc $(a,b)\in T$ z
$(y,a),(b,y')\in T$, co daje
$(y,y')\in T^3\subseteq V$.
Możemy zatem przyjąć $W=T$,
co dowodzi jednostajnej ciągłości.

Jednoznaczność: dwa ciągłe rozszerzenia
zgadzają się na gęstym $A$, więc
są równe na $Z$ dzięki Hausdorffowości $Y$.
\end{proof}

\begin{twierdzenie}[Własność uniwersalna uzupełnienia]
\label{thm:topog-uzupelnienia-uniwersalnosc}
Niech $\iota:X\to\widehat X$ będzie
uzupełnieniem rozdzielającej przestrzeni
jednostajnej $X$, skonstruowanym
w twierdzeniu~\ref{thm:topog-uzupelnienia-istnienie}.
Dla każdej zupełnej rozdzielającej
przestrzeni jednostajnej $Y$
i każdego jednostajnie ciągłego
$f:X\to Y$ istnieje dokładnie jedno
jednostajnie ciągłe $\widehat f:\widehat X\to Y$
takie, że $\widehat f\circ\iota=f$.
Ponadto dowolne dwa uzupełnienia $X$
są jednoznacznie izomorficzne jednostajnie
przez izomorfizm zachowujący obraz $X$.
\end{twierdzenie}
\begin{proof}
Zastosujmy lemat~\ref{lem:topog-uzupelnienia-rozszerzenie}
do gęstej podprzestrzeni $\iota(X)$
w $\widehat X$ i odwzorowania
$f\circ\iota^{-1}$. Otrzymujemy
żądane rozszerzenie.

Dla dwóch uzupełnień $X_1,X_2$
własność uniwersalna daje jednostajnie
ciągłe odwzorowania
$T:X_1\to X_2$ i $S:X_2\to X_1$
zgodne z odpowiednimi zanurzeniami.
Złożenia $S\circ T$ i $T\circ S$
są identycznościami na gęstych
obrazach $X$; z jednoznaczności
rozszerzenia są identycznościami
na $X_1$ i $X_2$. Tak więc
$T$ jest izomorfizmem jednostajnym.
\end{proof}

\section{Domknięte podprzestrzenie i przykłady}
\label{sec:topog-uzupelnienia-przyklady}

\begin{twierdzenie}[Zupełność podprzestrzeni]
\label{thm:topog-uzupelnienia-podprzestrzenie}
Niech $X$ będzie zupełną rozdzielającą
przestrzenią jednostajną i $A\subseteq X$
z odziedziczoną strukturą jednostajną.
Wówczas $A$ jest zupełna wtedy i tylko wtedy,
gdy jest domknięta w $X$.
\end{twierdzenie}
\begin{proof}
Jeżeli $A$ jest domknięta, to każdy filtr
Cauchy'ego na $A$, traktowany przez
obrazy zbiorów jako filtr na $X$,
jest Cauchy'ego. Ma granicę $x\in X$.
Ponieważ filtr ma bazę złożoną ze
zbiorów zawartych w $A$, granica należy
do $\overline A=A$.
Stąd filtr zbiega w $A$.

Odwrotnie, niech $A$ będzie zupełna
i $x\in\overline A$. Filtr na $A$
generowany przez $A\cap E[x]$ jest
Cauchy'ego, jak w dowodzie
lematu~\ref{lem:topog-uzupelnienia-rozszerzenie}.
Zbiega więc do pewnego $a\in A$.
Jako filtr na $X$ zbiega zarazem do $x$.
Hausdorffowość $X$ daje $a=x$;
stąd $A$ jest domknięta.
\end{proof}

\begin{przyklad}[Dwa uzupełnienia tej samej topologii]
\label{ex:topog-uzupelnienia-dwie-metryki}
Uzupełnienie $(\R,|\cdot|)$ jest równe
$\R$ z tą samą metryką. Z kolei dla
$e(x,y)=|\arctan x-\arctan y|$
izometria $\arctan$ identyfikuje $\R$
z otwartym przedziałem
$(-\pi/2,\pi/2)$, którego uzupełnieniem
jest $[-\pi/2,\pi/2]$.
Uzupełnienia te nie są nawet
homeomorficzne: pierwsze jest niezwartą
przestrzenią, a drugie zwartą.
\end{przyklad}

\begin{przyklad}[Produkt i zupełność]
\label{ex:topog-uzupelnienia-produkt}
Jeżeli każda $(X_i,\mathcal U_i)$ jest
zupełna, to produkt $\prod_{i\in I}X_i$
ze strukturą produktu z rozdziału~12
jest zupełny. Istotnie, filtr Cauchy'ego
$\mathcal F$ w produkcie przenosi się
przez każde rzutowanie na filtr Cauchy'ego
w $X_i$. Wybierając jego granicę $x_i$
dla każdego $i$ (dla pustego produktu
otrzymujemy jedyny punkt), dostajemy
$x=(x_i)$, do którego zbiega
$\mathcal F$: każde otwarte otoczenie
$x$ zawiera cylinder oparty na skończenie
wielu otoczeniach $x_i$, a przeciwobrazy
tych otoczeń należą do $\mathcal F$.
Gdy którykolwiek czynnik jest pusty,
produkt jest pusty i zupełny w sensie
definicji~\ref{def:topog-uzupelnienia-zupelnosc}.
\end{przyklad}

\begin{uwaga}[Kierunek dalszej pracy]
Każda zwarta przestrzeń jednostajna jest
zupełna. Rzeczywiście, filtr Cauchy'ego
$\mathcal F$ w zwartej przestrzeni $X$
ma punkt skupienia $x$: rodzina
$\{\overline A:A\in\mathcal F\}$ ma własność
skończonych przecięć, więc jej przecięcie
jest niepuste. Ustalmy $E\in\mathcal U$
i symetryczne $D$ z $D^2\subseteq E$.
Istnieje $A\in\mathcal F$ z
$A\times A\subseteq D$. Ponieważ
$x\in\overline A$ i $D[x]$ jest otoczeniem
$x$, istnieje $a\in A\cap D[x]$.
Dla każdego $b\in A$ mamy
$(x,a),(a,b)\in D$, a zatem
$b\in E[x]$. Stąd $E[x]\in\mathcal F$,
czyli $\mathcal F\to x$.
Konstrukcja uzupełnienia
nie jest jednak uzwarceniem:
uzupełnienie nie musi być zwarte,
a jego funkcją jest domknięcie braków
w filtrach Cauchy'ego, nie dodanie granic
wszystkim filtrom. Powiązania z
uzwarceniami rozważymy oddzielnie.
\end{uwaga}

\chapter{Zwartość lokalna}
\label{ch:topog-zwartosc-lokalna}

\noindent\textbf{Cel rozdziału.}
Zwartość lokalna pozwala w otoczeniu każdego punktu
pracować ze zbiorami zwartymi, chociaż cała przestrzeń
nie musi być zwarta. Zaczniemy od definicji i podstawowego
lematu o otoczeniach mających zwarte domknięcia,
a następnie udowodnimy lokalnie zwartą wersję
twierdzenia Baire'a. Wersja metryczna tego twierdzenia
oraz pojęcie przestrzeni Baire'a zostały omówione
w rozdziale~11.

\section{Zwartość lokalna i twierdzenie Baire'a}
\label{sec:topog-baire-lokalnie-zwarte}

\begin{definicja}[Zwartość lokalna]
\label{def:topog-baire-lokalna-zwartosc}
Przestrzeń topologiczna $X$ jest
\emph{lokalnie zwarta}, jeżeli każdy
$x\in X$ ma zwarte otoczenie:
istnieje zwarty $K\subseteq X$
taki, że $x\in\operatorname{int}_X K$.
\end{definicja}

\begin{lemat}[Otoczenia o zwartym domknięciu]
\label{lem:topog-baire-zwarte-domkniecie}
Niech $X$ będzie lokalnie zwartą przestrzenią Hausdorffa.
Dla każdego $x\in X$ i każdego otwartego otoczenia
$V$ punktu $x$ istnieje otwarty $W\ni x$
o zwartym domknięciu $\overline W^{\,X}\subseteq V$.
\end{lemat}
\begin{proof}
Z lokalnej zwartości wybieramy zwarty
$K\subseteq X$ z $x\in\operatorname{int}_X K$.
Przestrzeń $K$ jest zwarta i Hausdorffa,
więc normalna (rozdział~07). Zbiór
$O=V\cap\operatorname{int}_X K$ jest
otwartym otoczeniem $x$ w $K$.
Kryterium zwężania otoczeń w przestrzeni
regularnej (rozdział~05) daje zbiór $H$
otwarty w $K$ spełniający
$x\in H\subseteq\overline H^{\,K}\subseteq O$.
Ponieważ $H\subseteq\operatorname{int}_X K$,
jest on także otwarty w $X$: jeśli
$H=T\cap K$ dla otwartego $T\subseteq X$,
to $H=T\cap\operatorname{int}_X K$.
Przyjmujemy $W=H$. Zwartość $K$
i własność Hausdorffa zapewniają domkniętość
$K$ w $X$ (rozdział~07), toteż
$\overline W^{\,X}=\overline H^{\,K}\subseteq V$.
Domknięty podzbiór zwartego $K$
jest zwarty, co kończy dowód.
\end{proof}

\begin{twierdzenie}[Twierdzenie Baire'a --- wersja lokalnie zwarta]
\label{thm:topog-baire-lokalnie-zwarte}
Każda lokalnie zwarta przestrzeń Hausdorffa
jest przestrzenią Baire'a.
\end{twierdzenie}
\begin{proof}
Niech $U$ będzie niepustym zbiorem otwartym,
a $G_n$ --- otwartymi zbiorami gęstymi.
Wybieramy $x_1\in U\cap G_1$.
Lemat~\ref{lem:topog-baire-zwarte-domkniecie}
daje niepusty otwarty $W_1$ o niepustym
zwartym domknięciu $K_1\subseteq U\cap G_1$.

Indukcyjnie, mając niepusty otwarty $W_n$
o zwartym domknięciu $K_n$, wybieramy
$x_{n+1}\in W_n\cap G_{n+1}$
i na mocy lematu otrzymujemy niepusty otwarty
$W_{n+1}$, którego zwarte domknięcie
$K_{n+1}$ zawiera się w $W_n\cap G_{n+1}$.
Zbiory $K_n$ są więc niepuste, zstępujące,
domknięte w $X$ (gdyż $X$ jest Hausdorffa)
i wszystkie zawierają się w zwartym $K_1$.

Gdyby $\bigcap_{n\ge1}K_n=\varnothing$,
to otwarte w $K_1$ zbiory $K_1\setminus K_n$
tworzyłyby pokrycie $K_1$. Ze zwartości
istniałoby podpokrycie skończone,
a ponieważ $(K_n)$ jest zstępujący,
oznaczałoby to $K_N=\varnothing$ dla
pewnego $N$. Jest to sprzeczność.
Zatem $\bigcap_nK_n$ jest niepuste
i zawiera się w $U\cap\bigcap_nG_n$.
\end{proof}


\chapter{Parazwartość i podziały jedności}
\label{ch:topog-parazwartosc}

\noindent\textbf{Cel rozdziału.}
Własność Lindelöfa pozwala wybrać przeliczalne podpokrycie,
ale nie kontroluje liczby zbiorów spotykających małe otoczenie
danego punktu. Parazwartość wymaga lokalnie skończonych
uszczegółowień pokryć. Powiążemy ją z podziałami jedności
i skonstruujemy takie podziały dla regularnych przestrzeni
Lindelöfa, w tym dla ważnych klas przestrzeni metryzowalnych
oraz lokalnie zwartych. Korzystamy wyłącznie z wyników
rozdziałów~5, 9, 10 i 14.

\section{Pokrycia lokalnie skończone}
\label{sec:topog-parazwartosc-pokrycia}

\begin{definicja}[Rodzina lokalnie skończona]
\label{def:topog-parazwartosc-lokalnie-skonczona}
Niech $X$ będzie przestrzenią topologiczną, a
$(A_i)_{i\in I}$ rodziną jej podzbiorów.
Rodzina jest \emph{lokalnie skończona}, jeżeli
dla każdego $x\in X$ istnieje otwarte otoczenie
$O_x$ takie, że zbiór
$\{i\in I:O_x\cap A_i\ne\varnothing\}$
jest skończony. Powtórzenia indeksowanych
zbiorów są liczone osobno.
\end{definicja}

\begin{definicja}[Uszczegółowienie i parazwartość]
\label{def:topog-parazwartosc-definicja}
Niech $\mathcal U$ będzie pokryciem otwartym
przestrzeni $X$. Pokrycie $\mathcal V$ jest
\emph{uszczegółowieniem} $\mathcal U$, gdy
każdy $V\in\mathcal V$ zawiera się w pewnym
$U\in\mathcal U$. Przestrzeń $X$ nazywamy
\emph{parazwartą}, gdy każde jej pokrycie
otwarte ma lokalnie skończone uszczegółowienie
otwarte.
Nie zakładamy tu automatycznie
aksjomatu Hausdorffa; w twierdzeniach
wymagających go podajemy go jawnie.
\end{definicja}

\begin{stwierdzenie}[Domknięcia rodzin lokalnie skończonych]
\label{prop:topog-parazwartosc-domkniecia}
Niech $(A_i)_{i\in I}$ będzie lokalnie skończoną
rodziną podzbiorów przestrzeni topologicznej $X$.
Wówczas $(\overline{A_i}^{\,X})_{i\in I}$
jest lokalnie skończona, a ponadto
\[
\overline{\bigcup_{i\in I} A_i}^{\,X}
 =\bigcup_{i\in I}\overline{A_i}^{\,X}.
\]
\end{stwierdzenie}
\begin{proof}
Niech $x\in X$ i niech otwarty $O\ni x$
przecina tylko zbiory $A_i$ o indeksach
należących do skończonego $F\subseteq I$.
Jeżeli $i\notin F$, to $O\cap A_i=\varnothing$;
ponieważ $O$ jest otwarty, również
$O\cap\overline{A_i}^{\,X}=\varnothing$.
Zatem rodzina domknięć jest lokalnie skończona.

Inkluzja od prawej do lewej w podanej
równości wynika z monotoniczności domknięcia.
Dla odwrotnej ustalmy
$x\notin\bigcup_i\overline{A_i}^{\,X}$.
Wybierzmy otwarty $O\ni x$ przecinający
jedynie $A_i$, $i\in F$, gdzie $F$ jest
skończony. Zbiór
$O'=O\setminus\bigcup_{i\in F}
\overline{A_i}^{\,X}$
jest otwartym otoczeniem $x$ i nie
przecina żadnego $A_i$.
Stąd $x\notin\overline{\bigcup_i A_i}^{\,X}$.
\end{proof}

\begin{przyklad}[Lokalna skończoność nie oznacza skończoności]
\label{ex:topog-parazwartosc-przedzialy}
W $\R$ rodzina
$A_n=(n-\frac13,n+\frac13)$,
$n\in\mathbb Z$, jest nieskończona
i lokalnie skończona.
Istotnie, dla $x\in\R$ wybieramy
$\varepsilon\in(0,1)$.
Jeśli $(x-\varepsilon,x+\varepsilon)$
spotyka $A_n$, to
$|n-x|<\varepsilon+\frac13$,
co zachodzi tylko dla skończenie
wielu liczb całkowitych $n$.
W przeciwieństwie do tego rodzina
$(-1/n,1/n)$, $n\ge1$, nie jest
lokalnie skończona w punkcie $0$:
każde otoczenie $0$ przecina
wszystkie jej elementy.
\end{przyklad}

\section{Podziały jedności}
\label{sec:topog-parazwartosc-jednosc}

\begin{definicja}[Nośnik i podział jedności]
\label{def:topog-parazwartosc-podzial}
Niech $X$ będzie przestrzenią topologiczną.
Nośnikiem funkcji $\varphi:X\to\R$ nazywamy
$\operatorname{supp}\varphi
=\overline{\{x\in X:\varphi(x)\ne0\}}^{\,X}$.
\emph{Podziałem jedności} na $X$ nazywamy
rodzinę ciągłych funkcji
$\varphi_i:X\to[0,1]$, $i\in I$, dla której:
\begin{enumerate}[label=(\roman*)]
\item $(\operatorname{supp}\varphi_i)_{i\in I}$
      jest lokalnie skończona;
\item $\sum_{i\in I}\varphi_i(x)=1$
      dla każdego $x\in X$.
\end{enumerate}
Suma jest punktowo skończona na mocy (i).
Podział jest \emph{podporządkowany}
pokryciu otwartemu $\mathcal U$,
jeśli dla każdego $i$ istnieje
$U_i\in\mathcal U$ z
$\operatorname{supp}\varphi_i\subseteq U_i$.
\end{definicja}

\begin{stwierdzenie}[Podział jedności daje uszczegółowienie]
\label{prop:topog-parazwartosc-jednosc-pokrycie}
Jeżeli każde pokrycie otwarte przestrzeni
topologicznej $X$ ma podporządkowany
mu podział jedności, to $X$ jest parazwarta.
\end{stwierdzenie}
\begin{proof}
Niech $\mathcal U$ będzie pokryciem otwartym
i $(\varphi_i)$ podporządkowanym mu
podziałem jedności. Zbiory
$V_i=\{x:\varphi_i(x)>0\}$ są otwarte
na mocy ciągłości. Pokrywają $X$,
ponieważ $\sum_i\varphi_i(x)=1$
wyklucza jednoczesne zerowanie
wszystkich składników w punkcie $x$.
Każdy $V_i$ zawiera się w nośniku
$\varphi_i$, a przez to w pewnym
elemencie $\mathcal U$.
Rodzina $(V_i)$ jest lokalnie skończona,
gdyż każda jej część zawiera się
w odpowiadającym nośniku.
\end{proof}

\begin{lemat}[Funkcje odcinające w przestrzeni normalnej]
\label{lem:topog-parazwartosc-odciecia}
Niech $X$ będzie normalną przestrzenią
Hausdorffa, a $F\subseteq V\subseteq X$,
gdzie $F$ jest domknięty, a $V$ otwarty.
Istnieje funkcja ciągła $f:X\to[0,1]$
taka, że $f|_F=1$ oraz
$\operatorname{supp}f\subseteq V$,
o ile istnieje otwarty $W$ spełniający
$F\subseteq W\subseteq\overline W^{\,X}
\subseteq V$.
\end{lemat}
\begin{proof}
Zbiory $F$ i $X\setminus W$ są
rozłączne i domknięte. Lemat Urysohna
z rozdziału~5 daje ciągłą funkcję
$f:X\to[0,1]$ równą $1$ na $F$
i $0$ na $X\setminus W$.
Zatem $\{f\ne0\}\subseteq W$,
a po domknięciu
$\operatorname{supp}f
\subseteq\overline W^{\,X}\subseteq V$.
\end{proof}

\section{Konstrukcja przeliczalnego podziału jedności}
\label{sec:topog-parazwartosc-konstrukcja}

\begin{twierdzenie}[Podział jedności w regularnej przestrzeni Lindelöfa]
\label{thm:topog-parazwartosc-lindelof}
Niech $X$ będzie regularną przestrzenią
Hausdorffa o własności Lindelöfa,
a $\mathcal U$ dowolnym pokryciem
otwartym $X$. Istnieje co najwyżej
przeliczalny, lokalnie skończony
podział jedności na $X$
podporządkowany $\mathcal U$.
W konsekwencji $X$ jest parazwarta.
\end{twierdzenie}
\begin{proof}
Dla każdego $x\in X$ wybieramy
$U_x\in\mathcal U$ z $x\in U_x$.
Z regularności i kryterium zwężania
otoczeń z rozdziału~5 otrzymujemy
otwarte zbiory $W_x,V_x$ takie, że
\[
x\in W_x\subseteq\overline{W_x}^{\,X}
\subseteq V_x\subseteq
\overline{V_x}^{\,X}\subseteq U_x.
\]
Rodzina $(W_x)_{x\in X}$ pokrywa $X$.
Z własności Lindelöfa wybieramy
przeliczalne podpokrycie
$(W_n)_{n\ge1}$ z odpowiadającymi
zbiorami $V_n$ i $U_n$.
Jeśli wybrane podpokrycie jest skończone,
dopowiadamy puste $W_n,V_n$ i dowolne
$U_n\in\mathcal U$ dla pozostałych
indeksów; przy $X=\varnothing$ twierdzenie
jest bezpośrednie.

Z rozdziału~10 wiemy, że przestrzeń
regularna i Lindelöfa jest normalna.
W związku z tym
lemat~\ref{lem:topog-parazwartosc-odciecia}
zastosowany do
$F_n=\overline{W_n}^{\,X}$,
otwartego zbioru $U_n$ i otwartego
zbioru pośredniego $V_n$ daje funkcje ciągłe
$f_n:X\to[0,1]$ takie, że
$f_n|_{F_n}=1$ i
$\operatorname{supp}f_n\subseteq
\overline{V_n}^{\,X}\subseteq U_n$.
Definiujemy kolejno
\[
\varphi_1=f_1,\qquad
\varphi_n=f_n\prod_{j=1}^{n-1}(1-f_j),
\quad n\ge2.
\]
Każda $\varphi_n$ jest ciągła,
nieujemna i nie przekracza $1$.
Ponadto $\{x:\varphi_n(x)\ne0\}
\subseteq\{x:f_n(x)\ne0\}$,
więc
$\operatorname{supp}\varphi_n
\subseteq\operatorname{supp}f_n
\subseteq U_n$.

Sprawdzimy lokalną skończoność nośników,
a nie tylko punktową skończoność sumy.
Ustalmy $x\in X$.
Ponieważ $(W_n)$ pokrywa $X$,
istnieje $N$ z $x\in W_N$.
Na całym otwartym otoczeniu $W_N$
funkcja $f_N$ jest równa $1$.
Dla $n>N$ we wzorze na $\varphi_n$
występuje czynnik $1-f_N$,
więc $\varphi_n$ zeruje się na $W_N$.
Wobec otwartości $W_N$ również
$\operatorname{supp}\varphi_n
\cap W_N=\varnothing$ dla $n>N$.
Otoczenie $W_N$ spotyka zatem
co najwyżej $N$ nośników.

Dla każdego $m\ge1$ zachodzi
tożsamość teleskopowa
\[
\sum_{n=1}^{m}\varphi_n(x)
=1-\prod_{n=1}^{m}(1-f_n(x)).
\]
Można ją sprawdzić przez indukcję,
odejmując kolejne iloczyny.
Jeżeli $x\in W_N$, to $f_N(x)=1$,
a więc dla wszystkich $m\ge N$
prawa strona jest równa $1$.
To dowodzi warunku sumy.
Otrzymaliśmy żądany podział jedności;
parazwartość wynika ze
stwierdzenia~\ref{prop:topog-parazwartosc-jednosc-pokrycie}.
\end{proof}

\begin{wniosek}[Dwie klasy przestrzeni parazwartych]
\label{cor:topog-parazwartosc-klasy}
Każda przestrzeń metryzowalna
spełniająca drugi aksjomat przeliczalności
oraz każda lokalnie zwarta, $\sigma$-zwarta
przestrzeń Hausdorffa jest parazwarta
i ma podział jedności podporządkowany
dowolnemu pokryciu otwartemu.
\end{wniosek}
\begin{proof}
W pierwszym przypadku przestrzeń jest
regularna i Hausdorffa jako metryzowalna,
a z rozdziału~9 spełnia własność Lindelöfa.
W drugim przypadku lokalna zwartość
i Hausdorffowość zapewniają regularność:
lemat o otoczeniach mających zwarte
domknięcia z rozdziału~14 daje
dla dowolnego otwartego otoczenia
punktu mniejsze otwarte otoczenie
o domknięciu zawartym w wyjściowym.
$\sigma$-zwartość implikuje własność
Lindelöfa: każde pokrycie otwarte
ma skończone podpokrycie każdego
ze zwartych zbiorów tworzących
przeliczalną sumę $X$; suma tych
skończonych rodzin jest przeliczalna.
W obu przypadkach stosujemy
twierdzenie~\ref{thm:topog-parazwartosc-lindelof}.
\end{proof}

\begin{przyklad}[Jawny podział na prostej]
\label{ex:topog-parazwartosc-prosta}
Dla $n\in\mathbb Z$ i $x\in\R$ połóżmy
$h_n(x)=\max\{0,1-|x-n|\}$.
Są to funkcje ciągłe o nośnikach
$[n-1,n+1]$. Rodzina nośników
jest lokalnie skończona.
Jeśli $x\in[n,n+1]$, to
$h_n(x)=n+1-x$ i
$h_{n+1}(x)=x-n$,
a wszystkie pozostałe $h_k(x)$
są równe zeru. Stąd
$\sum_{n\in\mathbb Z}h_n(x)=1$.
Otrzymujemy podział jedności
podporządkowany pokryciu
$((n-\tfrac32,n+\tfrac32))_{n\in\mathbb Z}$.
Warto zauważyć, że w punktach całkowitych
jeden składnik ma wartość $1$,
a wszystkie pozostałe wartość $0$.
\end{przyklad}

\begin{przyklad}[Nieprzeliczalna przestrzeń dyskretna]
\label{ex:topog-parazwartosc-dyskretna}
Niech $D$ będzie dowolną, także
nieprzeliczalną przestrzenią dyskretną.
Jest ona parazwarta: z danego
pokrycia otwartego wybieramy dla każdego
$x\in D$ singleton $\{x\}$.
Rodzina singletonów jest lokalnie skończona.
Funkcje $\chi_{\{x\}}:D\to[0,1]$
są ciągłe i tworzą podporządkowany
temu pokryciu podział jedności.
Przykład pokazuje, że parazwartość
nie implikuje własności Lindelöfa.
\end{przyklad}

\section{Twierdzenie Stone'a}
\label{sec:topog-parazwartosc-stone}

\begin{twierdzenie}[Twierdzenie Stone'a]
\label{thm:topog-parazwartosc-stone}
Niech $(X,d)$ będzie przestrzenią metryczną.
Wówczas $X$ jest parazwarta.
\end{twierdzenie}
\begin{proof}
Jeżeli $X=\varnothing$, to pusta rodzina
jest żądanym uszczegółowieniem.
Niech zatem $X\ne\varnothing$.
Korzystając z twierdzenia o dobrym uporządkowaniu,
zapiszmy pokrycie jako
$\mathcal U=(U_\alpha)_{\alpha<\kappa}$,
gdzie $\kappa$ jest liczbą porządkową.
Dla $x\in X$ istnieje najmniejszy
indeks $\alpha(x)<\kappa$ taki, że
$x\in U_{\alpha(x)}$.
Piszemy $r_n=2^{-n}$ dla $n\ge1$.

Konstruujemy kolejno, dla $n=1,2,\ldots$,
otwarte zbiory $V_{\alpha,n}$,
$\alpha<\kappa$.
Po zdefiniowaniu wszystkich poziomów
o indeksach mniejszych od $n$ określamy
zbiór dopuszczalnych środków
\[
E_{\alpha,n}=
\left\{z\in X:
\begin{array}{l}
\alpha(z)=\alpha,\\
z\notin\displaystyle\bigcup_{j<n}
            \bigcup_{\beta<\kappa}V_{\beta,j},\\
B_d(z,3r_n)\subseteq U_\alpha
\end{array}\right\}
\]
i przyjmujemy
\[
V_{\alpha,n}=\bigcup_{z\in E_{\alpha,n}}B_d(z,r_n).
\]
Dla $n=1$ suma po $j<n$ jest pusta.
Każdy $V_{\alpha,n}$ jest otwarty,
a z trzeciego warunku
$V_{\alpha,n}\subseteq U_\alpha$.

\emph{Pokrywanie.}
Niech $x\in X$ i $\alpha=\alpha(x)$.
Ponieważ $U_\alpha$ jest otwarty,
istnieje $n\ge1$ z
$B_d(x,3r_n)\subseteq U_\alpha$.
Jeżeli $x$ należy już do któregoś
$V_{\beta,j}$ z $j<n$, to jest pokryty.
W przeciwnym razie $x\in E_{\alpha,n}$,
a stąd $x\in V_{\alpha,n}$.
Zatem rodzina wszystkich
$V_{\alpha,n}$ pokrywa $X$.

\emph{Rozdzielenie na ustalonym poziomie.}
Ustalmy $i\ge1$ i dwa różne indeksy
$\beta<\gamma<\kappa$.
Jeżeli
$p\in V_{\beta,i}$ i
$q\in V_{\gamma,i}$,
to istnieją środki
$y\in E_{\beta,i}$,
$z\in E_{\gamma,i}$
z $d(p,y)<r_i$ i $d(q,z)<r_i$.
Ponieważ $\alpha(z)=\gamma>\beta$,
punkt $z$ nie należy do $U_\beta$.
Jednocześnie
$B_d(y,3r_i)\subseteq U_\beta$,
zatem $d(y,z)\ge3r_i$.
Nierówność trójkąta daje
\[
d(p,q)\ge d(y,z)-d(p,y)-d(q,z)>r_i.
\]
W konsekwencji otwarta kula
o średnicy nie większej niż $r_i$
może spotkać najwyżej jeden
z różnych zbiorów $V_{\beta,i}$.
Istotnie, gdyby spotykała dwa,
odległość wybranych w nich punktów
byłaby mniejsza od jej średnicy,
co przeczyłoby powyższemu oszacowaniu.

\emph{Lokalna skończoność wszystkich poziomów.}
Ustalmy $x\in X$. Z wykazanego
pokrywania wybierzmy $\alpha$ i $n$
takie, że $x\in V_{\alpha,n}$.
Z otwartości tego zbioru istnieje
$\rho>0$ z
$B_d(x,\rho)\subseteq V_{\alpha,n}$.
Wybierzmy liczbę całkowitą $M>n$
tak dużą, że $2r_M<\rho$,
i połóżmy $O=B_d(x,r_M)$.

Jeśli $i\ge M$ i
$z\in E_{\beta,i}$, to
drugi warunek definicji środka
wyklucza $z\in V_{\alpha,n}$,
ponieważ $n<i$.
W szczególności $d(x,z)\ge\rho$.
Dla $p\in B_d(z,r_i)$ mamy
\[
d(x,p)\ge d(x,z)-d(z,p)
>\rho-r_i\ge\rho-r_M>r_M.
\]
Zatem $O\cap V_{\beta,i}=\varnothing$
dla wszystkich $i\ge M$ i
$\beta<\kappa$.

Dla każdego z pozostałych poziomów
$1\le i<M$ zachodzi
$r_i\ge 2r_M$, a więc
średnica kuli $O$ jest nie większa
niż $r_i$. Na mocy wykazanego
rozdzielenia $O$ spotyka
co najwyżej jeden zbiór
$V_{\beta,i}$ na każdym takim poziomie.
W sumie $O$ spotyka najwyżej
$M-1$ zbiorów indeksowanej
rodziny $(V_{\beta,i})_{\beta<\kappa,\ i\ge1}$.
Dowodzi to lokalnej skończoności.

Otrzymana rodzina jest otwartym
pokryciem, jest lokalnie skończona
i uszczegóławia $\mathcal U$.
Ponieważ pokrycie $\mathcal U$ było
dowolne, przestrzeń $X$ jest parazwarta.
\end{proof}

\begin{uwaga}[Znaczenie założeń]
Dowód nie wymaga zupełności, ośrodkowości,
zwartości lokalnej ani żadnego aksjomatu
przeliczalności. Wykorzystuje
dobry porządek indeksów pokrycia,
zatem w standardowym ujęciu odwołuje
się do aksjomatu wyboru. W odróżnieniu
od wcześniejszej konstrukcji dla
przestrzeni Lindelöfa uzyskano tutaj
od razu lokalnie skończone pokrycie
otwarte, a nie podział jedności.
\end{uwaga}


\section{Twierdzenie Nagaty--Smirnowa o metryzowalności}
\label{sec:topog-parazwartosc-nagata-smirnow}

\begin{definicja}[Baza będąca przeliczalną sumą rodzin lokalnie skończonych]
\label{def:topog-parazwartosc-sigma-lokalnie-skonczona-baza}
Niech $X$ będzie przestrzenią topologiczną.
Rodzinę $\mathcal B$ jej otwartych podzbiorów
nazywamy \emph{bazą będącą przeliczalną sumą
rodzin lokalnie skończonych}, jeżeli jest bazą
topologii i istnieją rodziny
$\mathcal B_n$, $n\ge1$, lokalnie skończone
w $X$, takie, że
$\mathcal B=\bigcup_{n\ge1}\mathcal B_n$.
Nie zakładamy, że poszczególne
$\mathcal B_n$ są przeliczalne.
\end{definicja}

\begin{twierdzenie}[Twierdzenie Nagaty--Smirnowa]
\label{thm:topog-parazwartosc-nagata-smirnow}
Niech $X$ będzie przestrzenią topologiczną.
Następujące warunki są równoważne:
\begin{enumerate}[label=(\roman*)]
\item $X$ jest metryzowalna;
\item $X$ jest regularną przestrzenią $T_1$
      i ma bazę będącą przeliczalną sumą
      rodzin lokalnie skończonych.
\end{enumerate}
\end{twierdzenie}
\begin{proof}
\emph{Konieczność.}
Załóżmy, że topologię $X$ wyznacza
metryka $d$. Przestrzeń metryczna
jest regularna i spełnia aksjomat $T_1$
(rozdział~10). Dla każdego $n\ge1$
rozpatrzmy pokrycie otwarte
$\mathcal U_n=\{B_d(x,2^{-n}):x\in X\}$.
Na mocy twierdzenia Stone'a,
udowodnionego powyżej, istnieje
lokalnie skończone otwarte pokrycie
$\mathcal V_n$ uszczegóławiające
$\mathcal U_n$.

Pokażemy, że
$\mathcal B=\bigcup_{n\ge1}\mathcal V_n$
jest bazą topologii.
Niech $x\in O$, gdzie $O$ jest otwarty.
Wybierzmy $\varepsilon>0$
z $B_d(x,\varepsilon)\subseteq O$
oraz $n$ z $2^{1-n}<\varepsilon$.
Ponieważ $\mathcal V_n$ pokrywa $X$,
istnieje $V\in\mathcal V_n$ z $x\in V$.
Z uszczegółowienia otrzymujemy
$y\in X$ takie, że
$V\subseteq B_d(y,2^{-n})$.
Wówczas $d(x,y)<2^{-n}$,
więc dla każdego $z\in V$
mamy $d(x,z)<2^{1-n}<\varepsilon$.
Zatem $x\in V\subseteq O$.
Każda $\mathcal V_n$ jest lokalnie
skończona, co dowodzi (ii).

\emph{Wystarczalność.}
Przyjmijmy (ii) i zapiszmy
$\mathcal B=\bigcup_{n\ge1}\mathcal B_n$,
gdzie każda $\mathcal B_n$ jest lokalnie
skończona. Najpierw przygotujemy
dwa fakty potrzebne do zbudowania metryki.

\emph{Krok 1: normalność.}
Niech $A,C\subseteq X$ będą rozłączne
i domknięte. Dla każdego $n$ określmy
\[
\mathcal P_n=
\{B\in\mathcal B_n:
\overline B^{\,X}\cap C=\varnothing\},
\qquad
P_n=\bigcup_{B\in\mathcal P_n}B,
\]
a analogicznie
\[
\mathcal Q_n=
\{B\in\mathcal B_n:
\overline B^{\,X}\cap A=\varnothing\},
\qquad
Q_n=\bigcup_{B\in\mathcal Q_n}B.
\]
Każdy $P_n,Q_n$ jest otwarty.
Lokalna skończoność $\mathcal B_n$
oraz stwierdzenie
\ref{prop:topog-parazwartosc-domkniecia}
dają
$\overline{P_n}^{\,X}
=\bigcup_{B\in\mathcal P_n}\overline B^{\,X}$,
stąd
$\overline{P_n}^{\,X}\cap C=\varnothing$.
Podobnie
$\overline{Q_n}^{\,X}\cap A=\varnothing$.
Z regularności i własności bazy
wynika, że $(P_n)$ pokrywa $A$,
a $(Q_n)$ pokrywa $C$:
dla $a\in A$ można wybrać
otwarty $H\ni a$ z
$\overline H^{\,X}\subseteq X\setminus C$,
a następnie $B\in\mathcal B$
z $a\in B\subseteq H$;
wtedy $B\in\mathcal P_n$ dla pewnego $n$.
Dla punktów $C$ postępujemy analogicznie.

Połóżmy
\[
P=\bigcup_{n\ge1}
 \left(P_n\setminus
       \bigcup_{j\le n}\overline{Q_j}^{\,X}\right),
\qquad
Q=\bigcup_{n\ge1}
 \left(Q_n\setminus
       \bigcup_{j\le n}\overline{P_j}^{\,X}\right).
\]
Oba zbiory są otwarte, gdyż każdorazowo
odejmujemy skończoną sumę zbiorów
domkniętych od otwartego.
Mamy $A\subseteq P$, ponieważ
$A\cap\overline{Q_j}^{\,X}=\varnothing$
dla każdego $j$, i podobnie
$C\subseteq Q$.
Zbiory $P$ i $Q$ są rozłączne:
gdyby $x$ należał do składnika
$P_n$ oraz składnika $Q_m$,
to dla $m\le n$ mielibyśmy
$x\notin\overline{Q_m}^{\,X}$,
wbrew $x\in Q_m$, a dla $n\le m$
otrzymalibyśmy analogiczną
sprzeczność z $x\in P_n$.
Zatem $X$ jest normalna.

\emph{Krok 2: funkcje dodatnie dokładnie
na zadanym zbiorze otwartym.}
Niech $G\subseteq X$ będzie otwarty.
Dla każdego $n$ określmy
\[
F_n(G)=
\bigcup\{\overline B^{\,X}:
 B\in\mathcal B_n,\
 \overline B^{\,X}\subseteq G\}.
\]
Jest to zbiór domknięty:
rodzina występujących tu $B$
jest lokalnie skończona, więc
jej domknięcia mają domkniętą sumę,
zgodnie z wcześniejszym stwierdzeniem.
Ponadto $F_n(G)\subseteq G$.
Z regularności i własności bazy
wynika
$G=\bigcup_{n\ge1}F_n(G)$,
gdyż dla $x\in G$ można wybrać
$B\in\mathcal B$ z
$x\in B\subseteq\overline B^{\,X}
\subseteq G$.

Z kroku 1 i lematu Urysohna
(rozdział~5) dla każdego $n$
istnieje funkcja ciągła
$g_n:X\to[0,1]$ taka, że
$g_n=1$ na $F_n(G)$ i
$g_n=0$ na $X\setminus G$.
Szereg
\[
g_G(x)=\sum_{n\ge1}2^{-n}g_n(x)
\]
jest jednostajnie zbieżny,
ponieważ $0\le g_n\le1$ i
$\sum_{n\ge1}2^{-n}=1$.
Wynika stąd ciągłość $g_G$.
Dla $x\notin G$ wszystkie składniki
są zerowe; dla $x\in G$
co najmniej jeden składnik jest
równy $2^{-n}$.
Zatem $g_G(x)>0$ dokładnie dla $x\in G$.

\emph{Krok 3: konstrukcja metryki.}
Dla każdego $n\ge1$ i
$B\in\mathcal B_n$ wybierzmy
funkcję $g_B$ z kroku 2 i połóżmy
$f_{n,B}=2^{-n}g_B$.
Wtedy $f_{n,B}:X\to[0,2^{-n}]$
jest ciągła oraz
\[
f_{n,B}(x)>0
\ \Longleftrightarrow\ x\in B.
\]
Niech
$J=\{(n,B):n\ge1,\ B\in\mathcal B_n\}$.
Zbiór indeksów $J$ nie musi być
przeliczalny. Definiujemy
$F:X\to\ell^\infty(J)$ wzorem
$F(x)=(f_{n,B}(x))_{(n,B)\in J}$.
Jest to poprawnie określone, ponieważ
$0\le f_{n,B}(x)\le1$.

Pokażemy ciągłość $F$ względem
normy supremum. Ustalmy $x\in X$
i $\varepsilon>0$.
Wybierzmy $N$ z
$2^{-N}<\varepsilon/2$.
Dla $n\ge N$ i dowolnego $B$
różnica
$|f_{n,B}(x)-f_{n,B}(y)|$
jest mniejsza niż $\varepsilon/2$
dla każdego $y\in X$.
Dla każdego z skończenie wielu $n<N$
lokalna skończoność $\mathcal B_n$
daje otwarte otoczenie $O_n\ni x$
przecinające jedynie skończenie
wiele zbiorów z $\mathcal B_n$.
Poza tymi zbiorami wszystkie
$f_{n,B}$ zerują się na $O_n$
oraz w punkcie $x$.
Dla pozostałych, skończenie
wielu funkcji ciągłość w $x$
pozwala dodatkowo zwęzić otoczenie,
tak by
$|f_{n,B}(y)-f_{n,B}(x)|
<\varepsilon/2$.
Przecięcie skończenie wielu
takich otoczeń dla $n<N$
jest otwartym otoczeniem $O$ punktu $x$,
a dla $y\in O$ zachodzi
\[
\|F(y)-F(x)\|_\infty<\varepsilon.
\]
Stąd $F$ jest ciągła.

Ponieważ $X$ spełnia aksjomat $T_1$,
dla różnych $x,y\in X$ istnieje
otwarty zbiór $G$ taki, że
$x\in G$, $y\notin G$.
Wybierzmy $B\in\mathcal B_n$
z $x\in B\subseteq G$.
Wtedy $f_{n,B}(x)>0$ i
$f_{n,B}(y)=0$, zatem $F(x)\ne F(y)$.

Pozostaje sprawdzić ciągłość
odwzorowania odwrotnego
$F^{-1}:F(X)\to X$.
Niech $G\subseteq X$ będzie otwarty
i $x\in G$. Wybierzmy
$B\in\mathcal B_n$ z
$x\in B\subseteq G$ i połóżmy
$\delta=f_{n,B}(x)/2>0$.
Jeśli $y\in X$ i
$\|F(y)-F(x)\|_\infty<\delta$,
to $f_{n,B}(y)>0$, a więc
$y\in B\subseteq G$.
To dowodzi, że obraz
$F(G)$ jest otwarty w $F(X)$.
Zatem $F$ jest homeomorfizmem
$X$ na podprzestrzeń metrycznej
przestrzeni $\ell^\infty(J)$.
Metrykę zgodną na $X$ otrzymujemy
jako
\[
d(x,y)=
\sup_{(n,B)\in J}
|f_{n,B}(x)-f_{n,B}(y)|.
\]
To kończy dowód równoważności.
\end{proof}

\begin{uwaga}[Porównanie z kryterium Urysohna]
\label{rem:topog-parazwartosc-nagata-urysohn}
Twierdzenie Urysohna z rozdziału~10
wymaga regularności, aksjomatu $T_1$
i drugiego aksjomatu przeliczalności.
Baza przeliczalna jest automatycznie
przeliczalną sumą rodzin
jednoelementowych, zatem twierdzenie
Nagaty--Smirnowa obejmuje ten wynik.
Odwrotna implikacja nie zachodzi:
nieprzeliczalna przestrzeń dyskretna
ma bazę singletonów, która sama jest
lokalnie skończona, ale nie spełnia
drugiego aksjomatu przeliczalności.
Główną różnicą w dowodzie jest
użycie przestrzeni $\ell^\infty(J)$
dla dowolnego, także nieprzeliczalnego
zbioru indeksów $J$, z wagami
$2^{-n}$ kontrolującymi ciągłość.
\end{uwaga}

\chapter{Funkcje ciągłe i wprowadzenie do uzwarceń}
\label{ch:topog-funkcje-uzwarcenia}

\noindent\textbf{Cel rozdziału.}
W rozdziale~5 zdefiniowaliśmy całkowitą regularność,
a w rozdziale~7 udowodniliśmy twierdzenie o zanurzaniu
przestrzeni Tychonoffa w kostkę $[0,1]^I$.
Nie powtarzamy tamtych dowodów. Wyjaśnimy teraz,
w jakim sensie topologia jest wyznaczona przez funkcje
ciągłe, czym jest uzwarcenie i kiedy dwa uzwarcenia
uznajemy za równoważne. Własności uniwersalne
konkretnych uzwarceń będą przedmiotem rozdziału~17.

\section{Topologia wyznaczona przez funkcje ciągłe}
\label{sec:topog-uzwarcenia-funkcje}

\begin{definicja}[Topologia początkowa rodziny funkcji]
\label{def:topog-uzwarcenia-topologia-poczatkowa}
Niech $X$ będzie zbiorem, a
$(f_i\colon X\to Y_i)_{i\in I}$ rodziną
odwzorowań do przestrzeni topologicznych $Y_i$.
\emph{Topologią początkową} tej rodziny
nazywamy najmniejszą topologię na $X$,
względem której wszystkie $f_i$ są ciągłe.
Tworzą ją jako podbazę zbiory
$f_i^{-1}(V)$, gdzie $i\in I$,
a $V$ jest otwarty w $Y_i$.
\end{definicja}
\begin{proof}
Przecięcia skończone zbiorów wskazanej
podbazy tworzą bazę pewnej topologii
(uwzględniamy puste przecięcie równe $X$).
W tej topologii każda $f_i$ jest ciągła
z definicji ciągłości przez przeciwobrazy.
Każda inna topologia zapewniająca
ciągłość wszystkich $f_i$ musi zawierać
wszystkie zbiory podbazy, a więc również
otwarte zbiory przez nie generowane.
\end{proof}

\begin{stwierdzenie}[Kryterium ciągłości]
\label{prop:topog-uzwarcenia-poczatkowa-ciaglosc}
Niech $\tau$ będzie topologią początkową
rodziny $(f_i\colon X\to Y_i)_{i\in I}$.
Dla dowolnej przestrzeni topologicznej $Z$
odwzorowanie $g\colon Z\to(X,\tau)$
jest ciągłe wtedy i tylko wtedy,
gdy wszystkie $f_i\circ g$ są ciągłe.
\end{stwierdzenie}
\begin{proof}
Jeśli $g$ jest ciągłe, to ciągłość złożeń
wynika z ciągłości wszystkich $f_i$.
Odwrotnie, przeciwobraz podbazowego zbioru
$f_i^{-1}(V)$ przez $g$ jest równy
$(f_i\circ g)^{-1}(V)$ i jest otwarty.
Przeciwobrazy skończonych przecięć
i dowolnych sum tych zbiorów też są otwarte.
Zatem $g$ jest ciągłe.
\end{proof}

\begin{twierdzenie}[Funkcyjny opis przestrzeni Tychonoffa]
\label{thm:topog-uzwarcenia-funkcyjny-opis}
Niech $X$ będzie przestrzenią topologiczną
spełniającą aksjomat $T_1$. Następujące
warunki są równoważne:
\begin{enumerate}[label=(\roman*)]
\item $X$ jest całkowicie regularna;
\item topologia $X$ jest topologią początkową
rodziny $C(X,[0,1])$ wszystkich ciągłych
funkcji $X\to[0,1]$;
\item funkcje z $C(X,[0,1])$ oddzielają punkty,
a dla każdego otwartego $U\subseteq X$
i każdego $x\in U$ istnieje $f\in C(X,[0,1])$
taka, że $f(x)=1$ i $f|_{X\setminus U}=0$.
\end{enumerate}
\end{twierdzenie}
\begin{proof}
Niech (i) zachodzi. Dla każdego
otwartego $U$ i $x\in U$ całkowita
regularność zastosowana do
$F=X\setminus U$ daje funkcję
$h\in C(X,[0,1])$ równą $0$ w $x$
i $1$ na $F$. Funkcja $f=1-h$
spełnia warunek z (iii).
Dla różnych $x,y$ zastosowanie tego
samego argumentu do $U=X\setminus\{y\}$
jest dozwolone przez $T_1$
i daje $f(x)=1$, $f(y)=0$.
Stąd (iii).

Warunek (iii) implikuje (ii).
Każda funkcja z $C(X,[0,1])$
jest ciągła w wyjściowej topologii,
zatem topologia początkowa jest
od niej nie silniejsza. Z drugiej strony
dla dowolnego otwartego $U$ i $x\in U$
wybieramy $f$ jak w (iii); wtedy
$x\in f^{-1}((1/2,1])\subseteq U$,
a zbiór pośrodku jest otwarty
w topologii początkowej.
Zatem obie topologie są identyczne.

Załóżmy (ii). Jeśli $F\subseteq X$
jest domknięty i $x\notin F$,
to $U=X\setminus F$ jest otwarty.
Z opisu przez podbazę istnieje
skończona rodzina funkcji
$f_1,\ldots,f_m\in C(X,[0,1])$
i otwarte $V_1,\ldots,V_m\subseteq[0,1]$
takie, że
\[
x\in W=\bigcap_{j=1}^{m}f_j^{-1}(V_j)
\subseteq U.
\]
Jeżeli $F=\varnothing$, funkcja stale
równa $0$ spełnia żądanie. W przeciwnym
razie $m\ge1$. Dla każdego $j$
istnieje ciągła funkcja
$h_j\colon[0,1]\to[0,1]$
równa $0$ w $f_j(x)$ i równa $1$
na $[0,1]\setminus V_j$;
jej istnienie wynika z całkowitej
regularności odcinka, wykazanej
w rozdziale~7.
Wówczas
$h=\max_{1\le j\le m}(h_j\circ f_j)$
jest ciągła, $h(x)=0$ i
$h|_{X\setminus W}=1$.
Ponieważ $F\subseteq X\setminus W$,
mamy $h|_F=1$. Jest to (i).
\end{proof}

\begin{przyklad}[Nie wystarcza samo oddzielanie punktów]
\label{ex:topog-uzwarcenia-samo-oddzielanie}
Niech $X=\R$ ma topologię dyskretną,
a w rodzinie funkcji pozostawmy
tylko jedno odwzorowanie
$f(x)=\frac12+\frac1\pi\arctan x$.
Funkcja ta jest ciągła i różnowartościowa,
więc oddziela punkty. Nie wyznacza jednak
topologii dyskretnej: jej topologia
początkowa jest zwykłą topologią
euklidesową na $\R$, ponieważ $f$
jest homeomorfizmem na $(0,1)$.
Rozdzielanie punktów przez
\emph{wybraną} rodzinę funkcji
nie gwarantuje więc odtworzenia
wszystkich zbiorów otwartych.
\end{przyklad}

\section{Definicja uzwarcenia i jego równoważności}
\label{sec:topog-uzwarcenia-definicje}

\begin{definicja}[Uzwarcenie]
\label{def:topog-uzwarcenia-uzwarcenie}
Niech $X$ będzie przestrzenią topologiczną.
\emph{Uzwarceniem Hausdorffa} przestrzeni $X$
nazywamy parę $(K,e)$, gdzie $K$ jest
zwartą przestrzenią Hausdorffa, a
$e\colon X\to K$ jest zanurzeniem
topologicznym o gęstym obrazie.
Przy utożsamieniu $X$ z $e(X)$
mówimy, że $X$ jest gęstą
podprzestrzenią $K$.
\end{definicja}

\begin{twierdzenie}[Istnienie uzwarcenia]
\label{thm:topog-uzwarcenia-istnienie}
Przestrzeń topologiczna $X$ ma uzwarcenie
Hausdorffa wtedy i tylko wtedy, gdy
jest przestrzenią Tychonoffa.
\end{twierdzenie}
\begin{proof}
Jeśli $(K,e)$ jest uzwarceniem,
to $K$ jest przestrzenią Tychonoffa:
zwarta przestrzeń Hausdorffa jest
normalna, a lemat Urysohna z rozdziału~5
zapewnia całkowitą regularność.
Zatem jej podprzestrzeń $e(X)$
jest Tychonoffa na mocy wyników
rozdziału~5, a więc także $X$.

Odwrotnie, dla przestrzeni Tychonoffa
stosujemy twierdzenie
o zanurzaniu w kostkę z rozdziału~7.
Daje ono zanurzenie
$e\colon X\to[0,1]^I$.
Domknięcie
$K=\overline{e(X)}^{\, [0,1]^I}$
jest zwarte na mocy twierdzenia
Tychonoffa i domkniętości
podzbioru zwartej przestrzeni;
jest również Hausdorffa,
a $e(X)$ jest w nim gęsty.
To jest dokładnie wniosek
o kostkowym uzwarceniu
z rozdziału~7.
\end{proof}

\begin{definicja}[Równoważność uzwarceń]
\label{def:topog-uzwarcenia-rownowaznosc}
Dwa uzwarcenia $(K_1,e_1)$
i $(K_2,e_2)$ tej samej przestrzeni $X$
są \emph{równoważne}, jeśli istnieje
homeomorfizm $H\colon K_1\to K_2$
spełniający $H\circ e_1=e_2$.
Wymóg zgodności z zanurzeniami
jest częścią definicji.
\end{definicja}

\begin{stwierdzenie}[Jednoznaczność homeomorfizmu]
\label{prop:topog-uzwarcenia-homeomorfizm-jednoznaczny}
Niech $(K_1,e_1)$ i $(K_2,e_2)$
będą uzwarceniami Hausdorffa przestrzeni $X$.
Istnieje co najwyżej jedno
odwzorowanie ciągłe $H\colon K_1\to K_2$
takie, że $H\circ e_1=e_2$.
W szczególności homeomorfizm
wykazujący ich równoważność
jest jednoznaczny.
\end{stwierdzenie}
\begin{proof}
Jeśli $H_1,H_2$ mają podaną własność,
to są równe na gęstym $e_1(X)$.
Zbiór punktów, w których
$H_1=H_2$, jest domknięty:
jest przeciwobrazem przekątnej
domkniętej w $K_2\times K_2$
(przestrzeń $K_2$ jest Hausdorffa).
Zawiera gęsty $e_1(X)$,
więc jest równy $K_1$.
\end{proof}

\begin{stwierdzenie}[Uzwarcenie przestrzeni zwartej]
\label{prop:topog-uzwarcenia-zwarta}
Jeśli $X$ jest zwartą przestrzenią
Hausdorffa, to każde jej uzwarcenie
Hausdorffa $(K,e)$ spełnia $e(X)=K$.
Tym samym każde takie uzwarcenie
jest równoważne $(X,\mathrm{id}_X)$.
\end{stwierdzenie}
\begin{proof}
Obraz $e(X)$ jest zwarty jako obraz
ciągły zwartej przestrzeni,
a więc domknięty w przestrzeni
Hausdorffa $K$. Jest zarazem gęsty,
więc $e(X)=K$. Zanurzenie
$e\colon X\to K$ jest zatem
homeomorfizmem i zapewnia równoważność.
\end{proof}

\section{Porządek uzwarceń}
\label{sec:topog-uzwarcenia-porzadek}

\begin{definicja}[Porównywanie uzwarceń]
\label{def:topog-uzwarcenia-porzadek}
Niech $(K_1,e_1)$ i $(K_2,e_2)$ będą
uzwarceniami Hausdorffa przestrzeni $X$.
Mówimy, że $(K_1,e_1)$ jest
\emph{nie mniejsze} od $(K_2,e_2)$,
i piszemy $(K_2,e_2)\preccurlyeq(K_1,e_1)$,
jeżeli istnieje ciągłe
$p\colon K_1\to K_2$ spełniające
$p\circ e_1=e_2$.
\end{definicja}

\begin{stwierdzenie}[Porządek na klasach uzwarceń]
\label{prop:topog-uzwarcenia-porzadek-klasy}
Niech $X$ będzie przestrzenią Tychonoffa.
Jeśli $(K_2,e_2)\preccurlyeq(K_1,e_1)$,
to odwzorowanie $p\colon K_1\to K_2$
spełniające $p\circ e_1=e_2$
jest wyznaczone jednoznacznie
i jest surjektywne.
Ponadto relacja $\preccurlyeq$
jest zwrotna i przechodnia;
jeżeli zachodzi ona w obu kierunkach,
to uzwarcenia są równoważne.
Wobec tego wyznacza porządek częściowy
na klasach równoważności uzwarceń.
\end{stwierdzenie}
\begin{proof}
Jednoznaczność $p$ wynika ze
stwierdzenia o jednoznaczności
odwzorowań ciągłych zgodnych
z zanurzeniami z poprzedniej sekcji.
Zbiór $p(K_1)$ jest zwarty jako obraz
ciągły zwartego $K_1$, a więc
domknięty w przestrzeni Hausdorffa $K_2$.
Zawiera $p(e_1(X))=e_2(X)$,
czyli gęsty podzbiór $K_2$.
Stąd $p(K_1)=K_2$.

Zwrotność wynika z możliwości
wyboru $p=\mathrm{id}_{K}$.
Jeżeli $p\colon K_1\to K_2$
oraz $q\colon K_2\to K_3$
są ciągłe i spełniają odpowiednio
$p\circ e_1=e_2$ i $q\circ e_2=e_3$,
to $(q\circ p)\circ e_1=e_3$,
co dowodzi przechodniości.

Załóżmy teraz istnienie
ciągłych $p\colon K_1\to K_2$
i $q\colon K_2\to K_1$ zgodnych
z zanurzeniami w obie strony.
Wtedy
$(q\circ p)\circ e_1=e_1$
oraz $(p\circ q)\circ e_2=e_2$.
Oba złożenia są zgodne z identycznościami
na gęstych obrazach $e_1(X)$ i $e_2(X)$.
Z wcześniej udowodnionej jednoznaczności
wynika $q\circ p=\mathrm{id}_{K_1}$
i $p\circ q=\mathrm{id}_{K_2}$.
Zatem $p$ jest homeomorfizmem,
a uzwarcenia są równoważne.

Na klasach równoważności definicja
nie zależy od wyboru reprezentantów:
przy zamianie uzwarceń na równoważne
wystarczy złożyć $p$ z odpowiednimi
homeomorfizmami zgodnymi z zanurzeniami.
Dowiedzione zwrotność, przechodniość
i antysymetria na tych klasach
dają porządek częściowy.
\end{proof}

\section{Uzwarcenia nierównoważne: przykład}
\label{sec:topog-uzwarcenia-przyklady}

\begin{przyklad}[Dwa uzwarcenia przedziału otwartego]
\label{ex:topog-uzwarcenia-dwa-przedzialu}
Niech $X=(0,1)$ z topologią euklidesową.
Pierwszym uzwarceniem jest
$K_1=[0,1]$ z zanurzeniem
$e_1(t)=t$.

Drugim jest okrąg jednostkowy
$K_2=S^1=\{z\in\mathbb C:|z|=1\}$
z zanurzeniem
$e_2(t)=\exp(2\pi it)$.
Jest ono różnowartościowe na $(0,1)$:
jeśli $e_2(s)=e_2(t)$, to $s-t$
jest liczbą całkowitą z $(-1,1)$,
zatem $s=t$.
Jest homeomorfizmem na obraz
$S^1\setminus\{1\}$:
odwrotnością jest ciągły argument
o wartościach w $(0,2\pi)$,
podzielony przez $2\pi$.
Obraz jest gęsty w $S^1$,
więc $(K_2,e_2)$ jest uzwarceniem.

Te uzwarcenia nie są równoważne.
Gdyby homeomorfizm
$H\colon[0,1]\to S^1$ spełniał
$H(t)=e^{2\pi it}$ dla $t\in(0,1)$,
to z ciągłości w $0$ i $1$
mielibyśmy $H(0)=H(1)=1$.
Jest to sprzeczne z różnowartościowością
homeomorfizmu.
W ten sposób końce przedziału
mogą zostać dodane jako dwa
różne punkty albo jako jeden punkt.
\end{przyklad}

\begin{uwaga}[Zakres kolejnego rozdziału]
W rozdziale~17 zbadamy dokładniej
uzwarcenie Aleksandrowa oraz
uzwarcenie Čecha--Stone'a, w tym ich
założenia i własności uniwersalne.
Nie identyfikujemy na tym etapie
dowolnego uzwarcenia z tymi konstrukcjami:
przykład powyżej pokazuje, że
same zwartość i gęstość obrazu
nie wyznaczają uzwarcenia jednoznacznie.
\end{uwaga}

\chapter{Uzwarcenia Aleksandrowa i Čecha--Stone'a}
\label{ch:topog-uzwarcenia-kanoniczne}

\noindent\textbf{Cel rozdziału.}
W rozdziale~16 ustaliliśmy, czym jest uzwarcenie
Hausdorffa, kiedy ono istnieje i jak porównywać
uzwarcenia przez odwzorowania zgodne z zanurzeniami.
Zbudujemy teraz dwa kanoniczne uzwarcenia,
które rozwiązują różne zagadnienia:
uzwarcenie Aleksandrowa dodaje jeden punkt
do niezwartej przestrzeni lokalnie zwartej,
a uzwarcenie Čecha--Stone'a umożliwia
przedłużanie wszystkich odwzorowań
ciągłych do zwartych przestrzeni Hausdorffa.
Dla każdego z nich podamy dokładne założenia
i udowodnimy odpowiednią własność uniwersalną.

\section{Uzwarcenie Aleksandrowa}
\label{sec:topog-uzwarcenia-aleksandrow}

\begin{definicja}[Konstrukcja Aleksandrowa]
\label{def:topog-uzwarcenia-aleksandrow}
Niech $X$ będzie niezwartą przestrzenią
Hausdorffa, a $\infty$ nowym punktem,
nienależącym do $X$. Na zbiorze
$X^\ast=X\cup\{\infty\}$ rozważmy topologię,
której zbiory otwarte są następujące:
\begin{enumerate}[label=(\roman*)]
\item wszystkie zbiory otwarte w $X$;
\item zbiory $X^\ast\setminus C$,
      gdzie $C\subseteq X$ jest zwarty.
\end{enumerate}
Parę $(X^\ast,\iota)$, gdzie
$\iota\colon X\hookrightarrow X^\ast$
jest włączeniem, nazywamy
\emph{konstrukcją Aleksandrowa}.
\end{definicja}

\begin{lemat}[Topologia i zwartość konstrukcji]
\label{lem:topog-uzwarcenia-aleksandrow-topologia}
Niech $X$ będzie niezwartą przestrzenią
Hausdorffa. Powyższe rodziny tworzą
topologię na $X^\ast$,
względem której $X^\ast$ jest zwarte,
$X$ jest otwartą podprzestrzenią
z pierwotną topologią, a $X$ jest
gęste w $X^\ast$.
\end{lemat}
\begin{proof}
Zbiory $\varnothing$ i $X^\ast$ są otwarte:
pierwszy należy do (i), drugi otrzymujemy
w (ii) dla $C=\varnothing$.
Przecięcie dwóch zbiorów typu (ii)
ma postać $X^\ast\setminus(C_1\cup C_2)$,
gdzie $C_1\cup C_2$ jest zwarty.
Jeżeli $U\subseteq X$ jest otwarty,
a $C\subseteq X$ zwarty, to
$U\cap(X^\ast\setminus C)=U\setminus C$
jest otwarty w $X$, ponieważ zwarty
podzbiór przestrzeni Hausdorffa
jest domknięty.
Dowolna suma zbiorów typu (ii),
jeśli rodzina nie jest pusta,
jest dopełnieniem przecięcia
zbiorów zwartych w $X$;
przecięcie to jest zwarte,
gdyż stanowi domknięty podzbiór
jednego ustalonego zbioru zwartego.
Suma zbiorów typu (i) jest otwarta w $X$.
Suma zbioru typu (ii) z dowolnym otwartym
$U\subseteq X$ jest dopełnieniem
$C\setminus U$, które jest zwarte jako
domknięty podzbiór zwartego $C$.
Otrzymaliśmy aksjomaty topologii.

Zbiory otwarte w $X^\ast$, które
nie zawierają $\infty$, są dokładnie
zbiorami otwartymi w $X$.
Stąd $\iota$ jest zanurzeniem,
a $X=X^\ast\setminus\{\infty\}$
jest otwarty w $X^\ast$.

Niech $\mathcal V$ będzie dowolnym
pokryciem otwartym $X^\ast$.
Wybierzmy $V_\infty\in\mathcal V$
zawierający $\infty$.
Istnieje zwarty $C\subseteq X$ taki, że
$V_\infty=X^\ast\setminus C$.
Zbiory pozostałego pokrycia
indukują pokrycie otwarte $C$;
zwartość $C$ daje skończone
podpokrycie. Wraz z $V_\infty$
pokrywa ono $X^\ast$.

Gdyby jakiś otwarty zbiór
zawierający $\infty$ nie przecinał $X$,
miałby postać $\{\infty\}=X^\ast\setminus X$,
co oznaczałoby zwartość $X$.
Przeczy to założeniu. Dlatego
$\infty\in\overline X^{\,X^\ast}$
i $X$ jest gęste.
\end{proof}

\begin{twierdzenie}[Kryterium uzwarcenia Aleksandrowa]
\label{thm:topog-uzwarcenia-aleksandrow-kryterium}
Niech $X$ będzie niezwartą przestrzenią
Hausdorffa. Konstrukcja Aleksandrowa
$(X^\ast,\iota)$ jest uzwarceniem
Hausdorffa przestrzeni $X$
wtedy i tylko wtedy, gdy $X$
jest lokalnie zwarta.
\end{twierdzenie}
\begin{proof}
Załóżmy lokalną zwartość $X$.
Z poprzedniego lematu mamy
zwartość $X^\ast$ i gęstość $X$.
Dwa różne punkty $x,y\in X$
mają rozłączne otoczenia otwarte w $X$,
a więc również w $X^\ast$.
Niech $x\in X$.
Z lematu o zwartych domknięciach
z rozdziału~14 otrzymujemy otwarty
$U\ni x$ o zwartym domknięciu
$C=\overline U^{\,X}$.
Zbiory $U$ i $X^\ast\setminus C$
są rozłącznymi otoczeniami
odpowiednio $x$ i $\infty$.
Stąd $X^\ast$ jest Hausdorffa.

Odwrotnie, jeśli $X^\ast$ jest
Hausdorffa, to jako przestrzeń zwarta
Hausdorffa jest lokalnie zwarta.
Podprzestrzeń otwarta $X$
jest lokalnie zwarta,
co udowodniono w rozdziale~14.
\end{proof}

\begin{uwaga}[Dlaczego wymagamy niezwartości]
Jeżeli $X$ jest zwarta, to przy
tej samej definicji punkt $\infty$
jest izolowany, ponieważ
$\{\infty\}=X^\ast\setminus X$
jest otwarty. Wówczas $X$
nie jest gęsty w $X^\ast$,
więc konstrukcja nie jest
uzwarceniem w sensie rozdziału~16.
Przestrzeń zwarta Hausdorffa
ma już uzwarcenie $(X,\mathrm{id}_X)$.
\end{uwaga}

\section{Własność uniwersalna i minimalność}
\label{sec:topog-uzwarcenia-aleksandrow-uniwersalne}

\begin{definicja}[Odwzorowanie właściwe]
\label{def:topog-uzwarcenia-wlasciwe}
Niech $X,Y$ będą przestrzeniami
topologicznymi. Odwzorowanie ciągłe
$f\colon X\to Y$ nazywamy
\emph{właściwym}, jeżeli
$f^{-1}(C)$ jest zwarty w $X$
dla każdego zwartego $C\subseteq Y$.
\end{definicja}

\begin{twierdzenie}[Przedłużanie odwzorowań właściwych]
\label{thm:topog-uzwarcenia-aleksandrow-proper}
Niech $X,Y$ będą niezwartymi,
lokalnie zwartymi przestrzeniami Hausdorffa,
a $f\colon X\to Y$ odwzorowaniem ciągłym.
Istnieje ciągłe odwzorowanie
$f^\ast\colon X^\ast\to Y^\ast$
spełniające
$f^\ast|_X=f$ oraz
$f^\ast(\infty_X)=\infty_Y$
wtedy i tylko wtedy, gdy $f$
jest właściwe.
Przy tych warunkach $f^\ast$
jest wyznaczone jednoznacznie.
\end{twierdzenie}
\begin{proof}
Wymagane wartości jednoznacznie
wyznaczają odwzorowanie zbiorów.
Załóżmy, że $f$ jest właściwe.
Dla otwartego $V\subseteq Y$
przeciwobraz $(f^\ast)^{-1}(V)
=f^{-1}(V)$ jest otwarty w $X$,
więc również w $X^\ast$.
Jeśli otwarty zbiór w $Y^\ast$
zawiera $\infty_Y$, ma postać
$Y^\ast\setminus C$ dla pewnego
zwartego $C\subseteq Y$.
Jego przeciwobraz jest równy
$X^\ast\setminus f^{-1}(C)$,
otwarty z definicji konstrukcji,
bo $f^{-1}(C)$ jest zwarty.
Stąd ciągłość $f^\ast$.

Odwrotnie, niech $f^\ast$
będzie ciągłe i spełnia
podane warunki. Dla zwartego
$C\subseteq Y$ zbiór $C$ jest
domknięty w Hausdorffa $Y^\ast$.
Zatem $(f^\ast)^{-1}(C)$
jest domknięty w zwartym $X^\ast$,
a więc zwarty. Nie zawiera
$\infty_X$, ponieważ
$f^\ast(\infty_X)=\infty_Y\notin C$,
i jest równy $f^{-1}(C)$.
Jako podzbiór otwartej podprzestrzeni
$X$ jest również zwarty w $X$.
\end{proof}

\begin{twierdzenie}[Minimalność uzwarcenia Aleksandrowa]
\label{thm:topog-uzwarcenia-aleksandrow-minimalne}
Niech $X$ będzie niezwartą,
lokalnie zwartą przestrzenią Hausdorffa,
a $(K,e)$ dowolnym jej uzwarceniem
Hausdorffa. Istnieje dokładnie jedno
ciągłe odwzorowanie
$p\colon K\to X^\ast$ takie, że
$p\circ e=\iota$.
Jest ono surjektywne
i przenosi cały zbiór
$K\setminus e(X)$ na punkt $\infty$.
W notacji porządku rozdziału~16
oznacza to $(X^\ast,\iota)\preccurlyeq(K,e)$.
\end{twierdzenie}
\begin{proof}
Pokażemy najpierw, że $e(X)$
jest otwarty w $K$.
Ustalmy $x\in X$ i wybierzmy
otwarty zbiór $U\subseteq X$
taki, że $x\in U$ i
$C=\overline U^{\,X}$ jest zwarty.
Istnieje otwarty $O\subseteq K$
z $O\cap e(X)=e(U)$.
Ponieważ $e(X)$ jest gęsty w $K$,
zbiór $O\cap e(X)$ jest gęsty w $O$,
więc
$O\subseteq\overline{e(U)}^{\,K}$.
Zbiór $e(C)$ jest zwarty, a zatem
domknięty w Hausdorffa $K$;
zawiera $e(U)$, skąd
$O\subseteq e(C)\subseteq e(X)$.
Każdy $e(x)$ ma więc
otwarte w $K$ otoczenie zawarte
w $e(X)$, co dowodzi otwartości.

Zbiór $F=K\setminus e(X)$ jest
domknięty w $K$, a zatem zwarty.
Jest niepusty: w przeciwnym razie
$X$ byłaby homeomorficzna ze zwartym
$K$, wbrew założeniu.
Definiujemy
$p(e(x))=x$ dla $x\in X$
oraz $p(z)=\infty$ dla $z\in F$.
Jeśli $V\subseteq X$ jest otwarty,
to $p^{-1}(V)=e(V)$ jest otwarty w $K$,
ponieważ $e(X)$ jest otwarty.
Jeżeli $X^\ast\setminus C$ jest
otoczeniem $\infty$ dla zwartego
$C\subseteq X$, to
$p^{-1}(X^\ast\setminus C)
=K\setminus e(C)$ jest otwarty:
$e(C)$ jest zwarty, więc domknięty w $K$.
To dowodzi ciągłości $p$.
Odwzorowanie jest surjektywne,
gdyż obejmuje każdy punkt $X$
oraz $\infty$ (zbiór $F$ jest niepusty).
Jednoznaczność wynika z gęstości
$e(X)$ i Hausdorffowości $X^\ast$,
zgodnie z rozdziałem~16.
\end{proof}

\begin{wniosek}[Jednoznaczność uzwarcenia jednopunktowego]
\label{cor:topog-uzwarcenia-aleksandrow-jednoznacznosc}
Niech $X$ będzie niezwartą, lokalnie zwartą
przestrzenią Hausdorffa. Jeżeli $(K,e)$
jest jej uzwarceniem Hausdorffa i
$K\setminus e(X)$ składa się z dokładnie
jednego punktu, to $(K,e)$ jest równoważne
uzwarceniu Aleksandrowa $(X^\ast,\iota)$.
\end{wniosek}
\begin{proof}
Z poprzedniego twierdzenia istnieje
ciągłe surjektywne
$p\colon K\to X^\ast$ zgodne
z zanurzeniami. Jest ono różnowartościowe
na $e(X)$, a jedyny punkt
$K\setminus e(X)$ przechodzi na $\infty$;
zatem $p$ jest bijekcją.
Ciągła bijekcja ze zwartej
przestrzeni $K$ na przestrzeń
Hausdorffa $X^\ast$ jest
homeomorfizmem. Spełnia
$p\circ e=\iota$,
więc uzwarcenia są równoważne.
\end{proof}

\begin{przyklad}[Prosta i okrąg]
\label{ex:topog-uzwarcenia-aleksandrow-prosta}
Dla $X=\R$ funkcja
$e(t)=\exp(i(2\arctan t+\pi))$
jest homeomorfizmem $\R$
na $S^1\setminus\{1\}$.
Jej obraz jest gęsty w $S^1$.
Otrzymane uzwarcenie jest
równoważne uzwarceniu Aleksandrowa:
odwzorowanie $e$ przedłużamy,
przypisując $\infty$ wartość $1$.
Oznaczmy to przedłużenie przez
$\widetilde e\colon X^\ast\to S^1$.
Wykażemy jego ciągłość w $\infty$.
Ponieważ
$\lim_{t\to-\infty}2\arctan t=-\pi$
oraz
$\lim_{t\to+\infty}2\arctan t=\pi$,
mamy
$\lim_{t\to\pm\infty}e(t)=1$.
Dla każdego otwartego otoczenia
$V\subseteq S^1$ punktu $1$
istnieje $M>0$ takie, że
$e(t)\in V$ dla $|t|>M$.
Zatem otoczenie
$X^\ast\setminus[-M,M]$ punktu
$\infty$ jest zawarte w
$\widetilde e^{-1}(V)$.
Dowodzi to ciągłości w $\infty$;
na otwartym $X$ ciągłość
wynika z ciągłości $e$.
Przedłużenie jest ciągłą bijekcją
ze zwartej $X^\ast$ do Hausdorffa $S^1$,
więc jest homeomorfizmem.
\end{przyklad}

\begin{lemat}[Całkowita regularność przestrzeni lokalnie zwartych]
\label{lem:topog-uzwarcenia-lokalna-tych}
Każda lokalnie zwarta przestrzeń
Hausdorffa jest przestrzenią Tychonoffa.
\end{lemat}
\begin{proof}
Niech $X$ będzie lokalnie zwarta
i Hausdorffa, $F\subseteq X$ domknięty,
a $x\in X\setminus F$.
Lemat o zwartych domknięciach
z rozdziału~14 daje otwarty
$W\subseteq X$ taki, że
$x\in W$, a
$C=\overline W^{\,X}$ jest zwarty
i zawarty w $X\setminus F$.
Przestrzeń $C$ jest zwarta i Hausdorffa,
więc normalna.
Zbiory $\{x\}$ i $C\setminus W$
są domknięte i rozłączne w $C$.
Lemat Urysohna zastosowany w $C$
daje ciągłą $g\colon C\to[0,1]$
z $g(x)=1$ i $g|_{C\setminus W}=0$.
Definiujemy $h\colon X\to[0,1]$
przez $h|_C=g$ i
$h|_{X\setminus W}=0$.
Jest to definicja zgodna na części
wspólnej $C\setminus W$.
Oba zbiory $C$ i $X\setminus W$
są domknięte i pokrywają $X$,
dlatego lemat o sklejaniu funkcji
ciągłych daje ciągłość $h$.
Mamy $h(x)=1$ i $h|_F=0$,
gdyż $F\cap C=\varnothing$.
Funkcja $1-h$ jest więc równa $0$
w $x$ i $1$ na $F$, co dowodzi
całkowitej regularności.
Aksjomat $T_1$ wynika
z Hausdorffowości.
\end{proof}

\section{Uzwarcenie Čecha--Stone'a}
\label{sec:topog-uzwarcenia-cech-stone}

\begin{definicja}[Konstrukcja Čecha--Stone'a]
\label{def:topog-uzwarcenia-cech-stone}
Niech $X$ będzie przestrzenią Tychonoffa.
Oznaczmy $I=C(X,[0,1])$.
Z rozdziału~7 znamy zanurzenie
$e_X\colon X\to[0,1]^I$ dane wzorem
$e_X(x)=(f(x))_{f\in I}$.
Definiujemy
\[
\beta X=\overline{e_X(X)}^{\, [0,1]^I}.
\]
Parę $(\beta X,e_X)$ nazywamy
\emph{uzwarceniem Čecha--Stone'a}
przestrzeni $X$.
\end{definicja}

\begin{twierdzenie}[Uniwersalna własność uzwarcenia Čecha--Stone'a]
\label{thm:topog-uzwarcenia-cech-stone-uniwersalne}
Niech $X$ będzie przestrzenią Tychonoffa,
$L$ zwartą przestrzenią Hausdorffa,
a $f\colon X\to L$ odwzorowaniem ciągłym.
Istnieje dokładnie jedno
odwzorowanie ciągłe
$\widehat f\colon\beta X\to L$
takie, że $\widehat f\circ e_X=f$.
\end{twierdzenie}
\begin{proof}
Z rozdziału~7 wiemy, że
$e_L\colon L\to[0,1]^J$,
gdzie $J=C(L,[0,1])$,
jest zanurzeniem.
Ponieważ $L$ jest zwarta,
$e_L(L)$ jest zwartym, a zatem
domkniętym podzbiorem kostki
Hausdorffa $[0,1]^J$.

Każde $g\in J$ wyznacza
funkcję $g\circ f\in I$.
Definiujemy
$\Phi\colon\beta X\to[0,1]^J$
współrzędnymi
\[
\Phi(z)_g=z_{g\circ f},\qquad g\in J.
\]
Każda współrzędna jest ograniczeniem
rzutowania z kostki $[0,1]^I$,
więc $\Phi$ jest ciągła z własności
uniwersalnej topologii produktowej.
Dla $x\in X$ i $g\in J$ mamy
\[
\Phi(e_X(x))_g=(g\circ f)(x)
=e_L(f(x))_g.
\]
Stąd $\Phi(e_X(X))\subseteq e_L(L)$.
Gęstość $e_X(X)$ w $\beta X$
i domkniętość $e_L(L)$ dają
$\Phi(\beta X)\subseteq e_L(L)$.
Możemy zatem zdefiniować
$\widehat f=e_L^{-1}\circ\Phi$.
Jest ono ciągłe i spełnia
$\widehat f\circ e_X=f$.
Jeśli dwa ciągłe odwzorowania
$\beta X\to L$ są równe
na gęstym $e_X(X)$,
to są równe wszędzie,
ponieważ $L$ jest Hausdorffa.
Otrzymujemy jednoznaczność.
\end{proof}

\begin{wniosek}[Przedłużanie funkcji ograniczonych]
\label{cor:topog-uzwarcenia-cech-stone-funkcje}
Niech $X$ będzie przestrzenią Tychonoffa.
Każda ograniczona funkcja ciągła
$f\colon X\to\R$ ma dokładnie jedno
ciągłe przedłużenie
$\widehat f\colon\beta X\to\R$.
\end{wniosek}
\begin{proof}
Jeśli $f(X)\subseteq[a,b]$
dla liczb rzeczywistych $a\le b$,
to stosujemy poprzednie twierdzenie
do zwartego przedziału $[a,b]$,
a następnie składamy przedłużenie
z włączeniem $[a,b]\hookrightarrow\R$.
Jednoznaczność wynika z gęstości
$e_X(X)$ w $\beta X$.
\end{proof}

\begin{twierdzenie}[Maksymalność uzwarcenia Čecha--Stone'a]
\label{thm:topog-uzwarcenia-cech-stone-maksymalne}
Niech $X$ będzie przestrzenią Tychonoffa,
a $(K,e)$ dowolnym uzwarceniem
Hausdorffa $X$. Istnieje dokładnie jedno
ciągłe odwzorowanie
$p\colon\beta X\to K$
z $p\circ e_X=e$.
Jest ono surjektywne.
Uzwarcenie $(K,e)$ jest równoważne
$(\beta X,e_X)$ wtedy i tylko wtedy,
gdy każda ciągła funkcja
$f\colon X\to[0,1]$ ma ciągłe
przedłużenie do $K$ zgodne z $e$.
\end{twierdzenie}
\begin{proof}
Istnienie i jednoznaczność $p$
wynikają z własności uniwersalnej
dla odwzorowania $e\colon X\to K$.
Surjektywność: obraz $p(\beta X)$
jest zwarty, a więc domknięty
w Hausdorffa $K$, i zawiera
gęsty $e(X)$.

Jeśli uzwarcenia są równoważne,
to każda współrzędna
$f\colon X\to[0,1]$
przedłuża się do $\beta X$,
a homeomorfizm $\beta X\to K$
pozwala przenieść ją na $K$.

Odwrotnie, załóżmy istnienie
przedłużenia $F_f\colon K\to[0,1]$
dla każdego $f\in I=C(X,[0,1])$,
z $F_f\circ e=f$.
Niech $\pi_f\colon[0,1]^I\to[0,1]$
oznacza rzutowanie na współrzędną
indeksowaną funkcją $f\in I$.
Z własności jednoznaczności
odwzorowań ciągłych na gęstym obrazie
wynika
\[
F_f\circ p=\pi_f|_{\beta X}
\qquad(f\in I),
\]
ponieważ obie strony mają wartości
$f(x)$ w punktach $e_X(x)$.
Jeśli $p(z)=p(w)$, to
$z_f=F_f(p(z))=F_f(p(w))=w_f$
dla wszystkich $f\in I$.
Jako punkty kostki $[0,1]^I$
mamy $z=w$, więc $p$ jest iniektywne.
Ciągła bijekcja ze zwartej $\beta X$
na Hausdorffa $K$ jest homeomorfizmem.
Jest zgodna z zanurzeniami,
zatem uzwarcenia są równoważne.
\end{proof}

\section{Porównanie konstrukcji i przykłady}
\label{sec:topog-uzwarcenia-porownanie}

\begin{wniosek}[Dwa krańce porządku uzwarceń]
\label{cor:topog-uzwarcenia-krance}
Niech $X$ będzie niezwartą,
lokalnie zwartą przestrzenią Hausdorffa.
Dla każdego uzwarcenia Hausdorffa $(K,e)$
zachodzą nierówności w porządku
z rozdziału~16:
\[
(X^\ast,\iota)\preccurlyeq(K,e)
\preccurlyeq(\beta X,e_X).
\]
\end{wniosek}
\begin{proof}
Pierwsza nierówność wynika
z minimalności uzwarcenia Aleksandrowa,
druga z maksymalności uzwarcenia
Čecha--Stone'a. Lokalna zwartość i Hausdorffowość
zapewniają aksjomat Tychonoffa
na mocy
lematu~\ref{lem:topog-uzwarcenia-lokalna-tych}.
\end{proof}

\begin{przyklad}[Uzwarcenie dyskretnej przestrzeni liczb naturalnych]
\label{ex:topog-uzwarcenia-beta-n}
Niech $X=\mathbb N$ z topologią
dyskretną. Na mocy wniosku
\ref{cor:topog-uzwarcenia-aleksandrow-jednoznacznosc}
uzwarcenie Aleksandrowa
jest równoważne uzwarceniu
$\{0\}\cup\{1/n:n\in\mathbb N\}$
z rozdziału~16, ponieważ
to ostatnie dodaje do $\mathbb N$
dokładnie jeden punkt.
Każda funkcja $f\colon\mathbb N\to[0,1]$
jest ciągła, ale przedłuża się
do tego uzwarcenia dokładnie wtedy,
gdy ciąg $(f(n))$ jest zbieżny.

W szczególności funkcja
$f(n)=0$ dla parzystych $n$
oraz $f(n)=1$ dla nieparzystych $n$
nie ma takiego przedłużenia.
Ma natomiast ciągłe przedłużenie
do $\beta\mathbb N$
na mocy twierdzenia
\ref{thm:topog-uzwarcenia-cech-stone-uniwersalne}.
Z kryterium równoważności wynika,
że uzwarcenia Aleksandrowa i Čecha--Stone'a
przestrzeni $\mathbb N$ nie są równoważne.
\end{przyklad}

\begin{uwaga}[Założenia i nazewnictwo]
Uzwarcenie Aleksandrowa w postaci
uzwarcenia Hausdorffa wymagającego
dokładnie jednego nowego punktu
dotyczy niezwartych, lokalnie zwartych
przestrzeni Hausdorffa.
Uzwarcenie Čecha--Stone'a istnieje
dla dowolnych przestrzeni Tychonoffa
i jest największe w porządku uzwarceń.
Te własności są różne:
pierwsza konstrukcja kontroluje
odwzorowania właściwe, druga
wszystkie ciągłe odwzorowania
do zwartych przestrzeni Hausdorffa.
\end{uwaga}

\chapter{Przestrzenie odwzorowań i topologia zwarto-otwarta}
\label{ch:topog-przestrzenie-odwzorowan}

\noindent\textbf{Cel rozdziału.}
Wcześniej rozpatrywaliśmy ciągłość pojedynczego
odwzorowania $X\to Y$.
Teraz sam zbiór wszystkich takich odwzorowań
uczynimy przestrzenią topologiczną.
Pozwoli to badać rodziny odwzorowań, ich granice,
odwzorowanie ewaluacji i zależność funkcji
od parametrów. Szczególną rolę będzie odgrywać
\emph{zbieżność jednostajna na zbiorach zwartych}.
Zwartość lokalna z rozdziału~14 będzie potrzebna
dopiero przy ewaluacji i prawie wykładniczym;
sama zgodność dwóch podstawowych topologii
na zbiorze funkcji ciągłych nie wymaga
tego założenia.

\section{Topologia zwarto-otwarta}
\label{sec:topog-odwzorowania-zwarto-otwarta}

\begin{definicja}[Przestrzeń odwzorowań ciągłych]
\label{def:topog-odwzorowania-Ck}
Niech $X,Y$ będą przestrzeniami topologicznymi.
Oznaczmy przez $C(X,Y)$ zbiór wszystkich
odwzorowań ciągłych $f\colon X\to Y$.
Dla każdego zwartego $K\subseteq X$
i każdego otwartego $U\subseteq Y$
definiujemy
\[
[K,U]=\{f\in C(X,Y):f(K)\subseteq U\}.
\]
\emph{Topologią zwarto-otwartą}
nazywamy topologię na $C(X,Y)$
generowaną przez rodzinę wszystkich
zbiorów $[K,U]$ jako podbazę.
Przestrzeń $C(X,Y)$ z tą topologią
oznaczamy przez $C_k(X,Y)$.
\end{definicja}

\begin{uwaga}[Dlaczego podbaza?]
Zbiory $[K,U]$ nie muszą same
stanowić bazy: warunek jednoczesnego
spełnienia $f(K_1)\subseteq U_1$
i $f(K_2)\subseteq U_2$ wymaga
przecięcia dwóch zbiorów podbazowych.
Bazę tworzą zatem ich przecięcia
skończone, a topologię --- dowolne sumy
takich przecięć. Dla $K=\varnothing$
mamy $[K,U]=C(X,Y)$.
\end{uwaga}

\begin{stwierdzenie}[Oddzielanie punktów i zbieżność punktowa]
\label{prop:topog-odwzorowania-oddzielanie}
Niech $X,Y$ będą przestrzeniami topologicznymi.
Topologia zwarto-otwarta na $C(X,Y)$
jest co najmniej tak silna jak topologia
zbieżności punktowej, czyli topologia
podprzestrzeni odziedziczona z iloczynu $Y^X$.
Jeżeli $Y$ jest Hausdorffa,
to $C_k(X,Y)$ jest Hausdorffa.
\end{stwierdzenie}
\begin{proof}
Dla $x\in X$ singleton $\{x\}$ jest
zwarty bez dodatkowych założeń na $X$.
Przeciwobraz zbioru otwartego
$U\subseteq Y$ przez ewaluację
$f\mapsto f(x)$ ma postać
$[\{x\},U]$, a więc jest otwarty
w topologii zwarto-otwartej.
Ponieważ topologia punktowa
jest początkowa względem wszystkich
tych ewaluacji, jest nie silniejsza.

Niech teraz $f,g\in C(X,Y)$ będą różne.
Istnieje $x\in X$ z $f(x)\ne g(x)$.
Wybierzmy rozłączne otwarte
$U,V\subseteq Y$ zawierające
odpowiednio $f(x)$ i $g(x)$.
Wtedy $[\{x\},U]$ oraz $[\{x\},V]$
są rozłącznymi otoczeniami $f$ i $g$.
\end{proof}

\begin{stwierdzenie}[Składanie z ustalonym odwzorowaniem]
\label{prop:topog-odwzorowania-skladanie}
Niech $X,X',Y,Z$ będą przestrzeniami
topologicznymi.
Jeśli $u\colon X'\to X$ i
$v\colon Y\to Z$ są ciągłe, to
odwzorowania
\[
u^\ast\colon C_k(X,Y)\to C_k(X',Y),
\quad u^\ast(f)=f\circ u,
\]
oraz
\[
v_\ast\colon C_k(X,Y)\to C_k(X,Z),
\quad v_\ast(f)=v\circ f
\]
są ciągłe.
\end{stwierdzenie}
\begin{proof}
Dla zwartego $K'\subseteq X'$
i otwartego $U\subseteq Y$ mamy
\[
(u^\ast)^{-1}([K',U])=[u(K'),U].
\]
Zbiór $u(K')$ jest zwarty
jako obraz ciągły zwartego $K'$,
więc prawa strona jest otwarta.
Dla zwartego $K\subseteq X$
i otwartego $V\subseteq Z$ mamy
\[
(v_\ast)^{-1}([K,V])
=[K,v^{-1}(V)].
\]
Drugi zbiór jest otwarty,
ponieważ $v^{-1}(V)$ jest otwarty.
Ciągłość obu odwzorowań wynika
z kryterium podbazowego.
\end{proof}

\section{Zbieżność jednostajna na zbiorach zwartych}
\label{sec:topog-odwzorowania-jednostajna-zwarte}

\begin{definicja}[Otoczenia jednostajne na zbiorze zwartym]
\label{def:topog-odwzorowania-jednostajna-zwarte}
Niech $X$ będzie przestrzenią topologiczną,
a $(Y,d)$ przestrzenią metryczną.
Dla $f\in C(X,Y)$, zwartego
$K\subseteq X$ i $\varepsilon>0$ określamy
\[
B_K(f,\varepsilon)=
\{g\in C(X,Y):
 d(f(x),g(x))<\varepsilon
 \text{ dla każdego }x\in K\}.
\]
Rodzina tych zbiorów, przy ustalonym $f$
i wszystkich zwartych $K$ oraz
$\varepsilon>0$, tworzy bazę
otoczeń pewnej topologii na $C(X,Y)$.
Nazywamy ją \emph{topologią zbieżności
jednostajnej na zbiorach zwartych}.
Dla $K=\varnothing$ przyjmujemy
$B_K(f,\varepsilon)=C(X,Y)$.
\end{definicja}
\begin{proof}
Każdy taki zbiór zawiera $f$.
Przecięcie dwóch otoczeń
$B_{K_1}(f,\varepsilon_1)$ i
$B_{K_2}(f,\varepsilon_2)$ zawiera
$B_{K_1\cup K_2}(f,\min\{\varepsilon_1,
\varepsilon_2\})$, gdyż skończona suma
zbiorów zwartych jest zwarta.
Jeśli $g\in B_K(f,\varepsilon)$
i $K\ne\varnothing$, to ciągła
funkcja $x\mapsto d(f(x),g(x))$
osiąga na zwartym $K$ maksimum
$m<\varepsilon$. Dla
$\delta=\varepsilon-m>0$
mamy
$B_K(g,\delta)\subseteq
B_K(f,\varepsilon)$
na mocy nierówności trójkąta.
Dla $K=\varnothing$ własność
jest oczywista. Spełnione są
więc aksjomaty bazy otoczeń.
\end{proof}

\begin{twierdzenie}[Równoważność dwóch topologii]
\label{thm:topog-odwzorowania-zwarto-otwarta-jednostajna}
Niech $X$ będzie dowolną przestrzenią
topologiczną, a $(Y,d)$ przestrzenią
metryczną. Na $C(X,Y)$ topologia
zwarto-otwarta jest równa topologii
zbieżności jednostajnej na zbiorach
zwartych. Nie jest wymagane, aby $X$
była Hausdorffa ani lokalnie zwarta.
\end{twierdzenie}
\begin{proof}
Najpierw niech $f\in[K,U]$,
gdzie $K\subseteq X$ jest zwarty,
a $U\subseteq Y$ jest otwarty.
Jeśli $K=\varnothing$ lub $U=Y$,
to $[K,U]=C(X,Y)$.
W pozostałym przypadku
$D=Y\setminus U$ jest niepusty
i domknięty, a $f(K)$ jest
zawartym w $U$ zbiorem zwartym.
Funkcja
$y\mapsto d(y,D)=\inf_{z\in D}d(y,z)$
jest ciągła, a dla każdego
$y\in f(K)$ ma wartość dodatnią.
Osiąga więc na $f(K)$ dodatnie
minimum $\eta>0$.
Jeśli $g\in B_K(f,\eta/2)$,
to $g(x)\notin D$ dla każdego
$x\in K$, a więc $g(K)\subseteq U$.
Zatem
$B_K(f,\eta/2)\subseteq[K,U]$.
Każdy zbiór podbazowy topologii
zwarto-otwartej jest otwarty
w topologii zbieżności jednostajnej
na zbiorach zwartych.

Odwrotnie, niech
$g\in B_K(f,\varepsilon)$.
Jeśli $K=\varnothing$, nie ma
czego dowodzić. Zwartość $K$
i ciągłość $x\mapsto d(f(x),g(x))$
dają liczbę
$m=\max_{x\in K}d(f(x),g(x))
<\varepsilon$.
Połóżmy $\delta=\varepsilon-m>0$.
Zwartość $g(K)$ zapewnia
istnienie punktów
$y_1,\ldots,y_s\in g(K)$
takich, że kule
$B_d(y_j,\delta/4)$ pokrywają
$g(K)$. Definiujemy zwarte zbiory
\[
K_j=K\cap
g^{-1}(\overline B_d(y_j,\delta/4)),
\qquad j=1,\ldots,s.
\]
Są zwarte, ponieważ są domkniętymi
podzbiorami zwartego $K$,
a ponadto $\bigcup_{j=1}^sK_j=K$.
Rozpatrzmy otwarte w topologii
zwarto-otwartej otoczenie
\[
\mathcal W=
\bigcap_{j=1}^s
[K_j,B_d(y_j,\delta/2)].
\]
Jest to rzeczywiście otoczenie
funkcji $g$: jeśli $x\in K_j$,
to $d(g(x),y_j)\le\delta/4<\delta/2$.
Dla dowolnego $h\in\mathcal W$
i $x\in K_j$ otrzymujemy
\[
\begin{split}
d(f(x),h(x))
&\le d(f(x),g(x))
+d(g(x),y_j)+d(y_j,h(x))\\
&<m+\delta/4+\delta/2
<\varepsilon.
\end{split}
\]
Każdy $x\in K$ należy do
pewnego $K_j$, więc
$\mathcal W\subseteq B_K(f,\varepsilon)$.
Dla każdego $g$ z tego otoczenia
znaleźliśmy mniejsze otoczenie
zwarto-otwarte; stąd
$B_K(f,\varepsilon)$ jest otwarty
w topologii zwarto-otwartej.
\end{proof}

\begin{wniosek}[Charakterystyka zbieżności sieci]
\label{cor:topog-odwzorowania-sieci}
Niech $X$ będzie przestrzenią topologiczną,
$(Y,d)$ przestrzenią metryczną,
a $(f_\lambda)_{\lambda\in\Lambda}$
siecią w $C(X,Y)$.
Sieć ta jest zbieżna do $f\in C(X,Y)$
w topologii zwarto-otwartej
wtedy i tylko wtedy, gdy dla każdego
zwartego $K\subseteq X$ i każdej
$\varepsilon>0$ istnieje
$\lambda_0\in\Lambda$ takie, że
dla $\lambda\ge\lambda_0$
oraz wszystkich $x\in K$ mamy
$d(f_\lambda(x),f(x))<\varepsilon$.
\end{wniosek}
\begin{proof}
Na mocy poprzedniego twierdzenia
zbiory $B_K(f,\varepsilon)$ stanowią
bazę otoczeń punktu $f$.
Definicja zbieżności sieci daje
dokładnie sformułowany warunek.
\end{proof}

\begin{przyklad}[Zbieżność na zwartych zbiorach nie jest globalnie jednostajna]
\label{ex:topog-odwzorowania-gorki-uciekajace}
Niech $X=Y=\R$ ze zwykłymi
topologiami i metryką na $Y$.
Dla $n\ge1$ określamy
\[
f_n(x)=\max\{1-|x-n|,0\}.
\]
Każda $f_n$ jest ciągła oraz
$\sup_{x\in\R}|f_n(x)|=1$.
Jeżeli $K\subseteq\R$ jest zwarty,
to jest ograniczony; dla
dostatecznie dużego $n$
mamy $f_n|_K=0$.
Zatem $f_n\to0$ jednostajnie
na każdym zbiorze zwartym,
lecz nie jednostajnie na całej $\R$.
\end{przyklad}

\section{Ewaluacja i prawo wykładnicze}
\label{sec:topog-odwzorowania-ewaluacja}

\begin{twierdzenie}[Ciągłość ewaluacji]
\label{thm:topog-odwzorowania-ewaluacja}
Niech $X$ będzie lokalnie zwartą
przestrzenią Hausdorffa, a $Y$
dowolną przestrzenią topologiczną.
Odwzorowanie ewaluacji
\[
\operatorname{ev}\colon
C_k(X,Y)\times X\to Y,\qquad
\operatorname{ev}(f,x)=f(x),
\]
jest ciągłe.
\end{twierdzenie}
\begin{proof}
Niech $(f_0,x_0)\in C_k(X,Y)\times X$
i niech $V\subseteq Y$ będzie otwartym
otoczeniem $f_0(x_0)$.
Zbiór $f_0^{-1}(V)$ jest otwartym
otoczeniem $x_0$ w $X$.
Z lematu o zwartych domknięciach
z rozdziału~14 istnieje
otwarty zbiór $W$ i zwarty zbiór $K$
takie, że
\[
x_0\in W\subseteq K
\subseteq f_0^{-1}(V).
\]
Wówczas $f_0\in[K,V]$,
a zbiór $[K,V]\times W$
jest otwartym otoczeniem $(f_0,x_0)$.
Dla dowolnego $(g,x)$ z tego
otoczenia mamy $x\in W\subseteq K$
i $g(K)\subseteq V$, więc
$\operatorname{ev}(g,x)=g(x)\in V$.
To wykazuje ciągłość w $(f_0,x_0)$.
\end{proof}

\begin{twierdzenie}[Prawo wykładnicze: funkcje z parametrem]
\label{thm:topog-odwzorowania-wykladnicze}
Niech $X$ będzie lokalnie zwartą
przestrzenią Hausdorffa, a $Z,Y$
dowolnymi przestrzeniami topologicznymi.
Niech $F\colon Z\times X\to Y$
będzie dowolnym odwzorowaniem.
Dla $z\in Z$ oznaczmy jego przekrój
przez $F_z\colon X\to Y$,
$F_z(x)=F(z,x)$.
Wówczas $F$ jest ciągłe
wtedy i tylko wtedy, gdy
każdy przekrój $F_z$ jest ciągły
oraz przyporządkowanie
$\widehat F\colon Z\to C_k(X,Y)$,
$\widehat F(z)=F_z$,
jest ciągłe.
W ten sposób przyporządkowanie
$F\mapsto\widehat F$ daje
bijekcję pomiędzy zbiorami
odwzorowań ciągłych
$C(Z\times X,Y)$ i
$C(Z,C_k(X,Y))$.
\end{twierdzenie}
\begin{proof}
Załóżmy, że $F$ jest ciągłe.
Dla każdego $z\in Z$
odwzorowanie $x\mapsto F(z,x)$
jest ciągłe jako złożenie
$F$ z ciągłym odwzorowaniem
$x\mapsto(z,x)$.
Niech teraz $[K,U]$ będzie
podbazowym zbiorem otwartym
w $C_k(X,Y)$ i niech
$z_0\in\widehat F^{-1}([K,U])$.
Wtedy $F(\{z_0\}\times K)\subseteq U$,
czyli otwarty zbiór $F^{-1}(U)$
zawiera $\{z_0\}\times K$.
Dla każdego $x\in K$ wybieramy
otwarte $O_x\ni z_0$ i $V_x\ni x$
z $O_x\times V_x\subseteq F^{-1}(U)$.
Zwartość $K$ daje skończone
$x_1,\ldots,x_m\in K$ takie, że
$K\subseteq\bigcup_{j=1}^mV_{x_j}$.
Wtedy $O=\bigcap_{j=1}^mO_{x_j}$
jest otwartym otoczeniem $z_0$
spełniającym $O\times K\subseteq
F^{-1}(U)$.
Jeśli $K=\varnothing$, bierzemy $O=Z$.
Zatem $O\subseteq
\widehat F^{-1}([K,U])$.
Wszystkie przeciwobrazy
zbiorów podbazowych są otwarte,
więc $\widehat F$ jest ciągłe.

Odwrotnie, załóżmy, że
$\widehat F\colon Z\to C_k(X,Y)$
jest ciągłe. Wówczas
\[
F=\operatorname{ev}\circ
(\widehat F\times\mathrm{id}_X).
\]
Pierwsze odwzorowanie w złożeniu
jest ciągłe jako iloczyn
odwzorowań ciągłych, a ewaluacja
jest ciągła na mocy poprzedniego
twierdzenia. Zatem $F$ jest ciągłe.
Obie operacje przyporządkowujące
sobie $F$ i $\widehat F$ są
wzajemnie odwrotne, co daje bijekcję.
\end{proof}

\begin{przyklad}[Dlaczego zwartość lokalna w ewaluacji nie jest zbędna]
\label{ex:topog-odwzorowania-ewaluacja-q}
Niech $X=\Q$ z topologią odziedziczoną
z $\R$, a $Y=\R$. Przestrzeń $\Q$
nie jest lokalnie zwarta i żaden
jej zwarty podzbiór nie ma
niepustego wnętrza w $\Q$
(patrz rozdział~14).
Wykażemy, że
$\operatorname{ev}\colon C_k(\Q,\R)
\times\Q\to\R$
nie jest ciągła w punkcie
$(0,0)$, gdzie $0$ oznacza
funkcję stale równą zero.

Niech $\mathcal N$ będzie dowolnym
otoczeniem zerowej funkcji w $C_k(\Q,\R)$,
a $V$ dowolnym otwartym otoczeniem
$0$ w $\Q$. Wybieramy bazowe
otoczenie
$\bigcap_{j=1}^m[K_j,U_j]
\subseteq\mathcal N$ funkcji zerowej,
gdzie $K_j\subseteq\Q$ są zwarte,
a $0\in U_j\subseteq\R$.
Zbiór $K=\bigcup_{j=1}^mK_j$
jest zwarty i ma puste wnętrze
w $\Q$; dla $m=0$ przyjmujemy
$K=\varnothing$. Istnieje więc
$r\in V\setminus K$.
Jeżeli $K\ne\varnothing$, funkcja
\[
g(q)=\min\left\{1,
 \frac{d(q,K)}{d(r,K)}\right\},
 \qquad q\in\Q,
\]
jest ciągła, zerowa na $K$
i spełnia $g(r)=1$.
Ponieważ $K$ jest domknięty
w przestrzeni metrycznej $\Q$
i $r\notin K$, mianownik jest dodatni.
Dla $K=\varnothing$ wybieramy
funkcję $g\equiv1$.
W obu przypadkach $g\in\mathcal N$,
ale $g(r)=1$. Stąd żadne
otoczenie punktu $(0,0)$ nie jest
odwzorowywane przez ewaluację
do $(-1/2,1/2)$, co przeczy
ciągłości w tym punkcie.
\end{przyklad}

\section{Metryzacja i zupełność przestrzeni odwzorowań}
\label{sec:topog-odwzorowania-metryzacja}

\begin{twierdzenie}[Metryka z wyczerpania zwartego]
\label{thm:topog-odwzorowania-metryka-wyczerpanie}
Niech $X$ będzie lokalnie zwartą,
$\sigma$-zwartą przestrzenią Hausdorffa,
a $(Y,d)$ przestrzenią metryczną.
Wybierzmy wyczerpanie zwarte
$(K_n)_{n\ge1}$ przestrzeni $X$
z rozdziału~14, tak że
$K_n\subseteq\operatorname{int}_XK_{n+1}$
oraz $X=\bigcup_{n\ge1}\operatorname{int}_XK_n$.
Dla $f,g\in C(X,Y)$ połóżmy
\[
\rho_n(f,g)=
\sup_{x\in K_n}\min\{1,d(f(x),g(x))\},
\qquad
D(f,g)=\sum_{n=1}^\infty
2^{-n}\rho_n(f,g).
\]
Przy $K_n=\varnothing$ ustalamy
$\rho_n(f,g)=0$.
Wówczas $D$ jest metryką zgodną
z topologią zwarto-otwartą na $C(X,Y)$.
\end{twierdzenie}
\begin{proof}
Każda $\rho_n$ jest nieujemna,
symetryczna i nie większa od $1$.
Nierówność trójkąta wynika
z nierówności
$\min\{1,a+b\}\le
\min\{1,a\}+\min\{1,b\}$
dla $a,b\ge0$, stosowanej punktowo,
a następnie z przejścia do supremum.
Dlatego szereg definiujący $D$
jest zbieżny i spełnia nierówność
trójkąta. Jeżeli $D(f,g)=0$,
to $\rho_n(f,g)=0$ dla każdego $n$.
Ponieważ zbiory $K_n$ pokrywają $X$,
otrzymujemy $f=g$.
Zatem $D$ jest metryką.

Dla każdego zwartego $K\subseteq X$
wyczerpanie z rozdziału~14 daje
$N$ z $K\subseteq K_N$.
W konsekwencji zbieżność
$\rho_N(f_\lambda,f)\to0$
wymusza jednostajną zbieżność
$f_\lambda$ do $f$ na $K$:
dla dowolnego $0<\varepsilon<1$
nierówność $\rho_N<\varepsilon$
oznacza $d(f_\lambda(x),f(x))
<\varepsilon$ dla wszystkich $x\in K$.
Odwrotnie, zbieżność jednostajna
na każdym $K_n$ implikuje
$\rho_n(f_\lambda,f)\to0$
dla każdego $n$.
Dla $\varepsilon>0$ wybieramy
$N$ tak, by
$\sum_{n>N}2^{-n}<\varepsilon/2$.
Następnie zbieżność skończenie
wielu $\rho_n$, $n\le N$,
pozwala zapewnić
$\sum_{n=1}^{N}2^{-n}\rho_n
<\varepsilon/2$.
Zatem $D(f_\lambda,f)<\varepsilon$.
Oba rozumowania obowiązują
dla sieci, nie tylko dla ciągów.
Na mocy twierdzenia
\ref{thm:topog-odwzorowania-zwarto-otwarta-jednostajna}
metryka $D$ wyznacza
topologię zwarto-otwartą.
\end{proof}

\begin{twierdzenie}[Zupełność przestrzeni odwzorowań]
\label{thm:topog-odwzorowania-zupelnosc}
Niech $X$ będzie lokalnie zwartą,
$\sigma$-zwartą przestrzenią Hausdorffa,
a $(Y,d)$ zupełną przestrzenią metryczną.
Dla dowolnego wyczerpania zwartego
$(K_n)$ metryka $D$ z poprzedniego
twierdzenia jest zupełna na $C(X,Y)$.
\end{twierdzenie}
\begin{proof}
Niech $(f_j)_{j\ge1}$ będzie ciągiem
Cauchy'ego w metryce $D$.
Dla każdego $n$ i $0<\varepsilon<1$
istnieje $j_0$ takie, że
\[
D(f_j,f_k)<2^{-n}\varepsilon
\qquad(j,k\ge j_0).
\]
Ponieważ $2^{-n}\rho_n(f_j,f_k)
\le D(f_j,f_k)$, otrzymujemy
$\rho_n(f_j,f_k)<\varepsilon$,
a więc
$d(f_j(x),f_k(x))<\varepsilon$
dla wszystkich $x\in K_n$.
Zatem $(f_j|_{K_n})$ jest
jednostajnie ciągiem Cauchy'ego.
Dla każdego $x\in X$ ciąg
$(f_j(x))$ jest więc ciągiem
Cauchy'ego w $Y$ i ma granicę
$f(x)\in Y$.

Wykażemy ilościowo jednostajną
zbieżność na $K_n$. Niech
$\eta>0$ będzie dowolna.
W rozumowaniu powyżej wybieramy
$0<\varepsilon<\min\{1,\eta\}$
i odpowiadający mu indeks $j_0$.
Dla $j\ge j_0$, $k\ge j_0$
oraz $x\in K_n$ mamy
$d(f_j(x),f_k(x))<\varepsilon$.
Przechodząc przy ustalonych $j,x$
do granicy $k\to\infty$,
otrzymujemy
$d(f_j(x),f(x))\le\varepsilon<\eta$.
Wybór $j_0$ nie zależy od $x\in K_n$,
więc $f_j\to f$ jednostajnie na $K_n$.
Każdy $x\in X$ należy do
$\operatorname{int}_X K_n$
dla pewnego $n$, więc $f$
jest ciągła w otoczeniu $x$
jako jednostajna granica funkcji
ciągłych (o wartościach w przestrzeni
metrycznej). W konsekwencji
$f\in C(X,Y)$.
Zbieżność jednostajna na wszystkich
$K_n$ implikuje $D(f_j,f)\to0$
na mocy poprzedniego twierdzenia.
Dowodzi to zupełności.
\end{proof}

\begin{przyklad}[Nieprzeliczalna przestrzeń dyskretna]
\label{ex:topog-odwzorowania-nieprzeliczalna}
Niech $X$ będzie nieprzeliczalną
przestrzenią dyskretną, a $Y=\R$.
Zwarte podzbiory $X$ są dokładnie
zbiorami skończonymi. Wobec tego
$C_k(X,\R)=\R^X$ z topologią
produktową, czyli topologią
zbieżności punktowej.

Pokażemy, że przestrzeń ta
nie spełnia pierwszego aksjomatu
przeliczalności w funkcji zerowej.
Gdyby $(U_n)_{n\ge1}$ było
przeliczalną bazą otoczeń,
dla każdego $n$ wybralibyśmy
podstawowe otoczenie produktowe
$V_n\subseteq U_n$ funkcji zerowej.
Zbiór współrzędnych ograniczanych
przez $V_n$ jest skończony;
oznaczmy go przez $F_n$.
Ponieważ $X$ jest nieprzeliczalna,
istnieje $x_0\in X\setminus
\bigcup_{n\ge1}F_n$.
Otoczenie
$W=\{f\in\R^X:|f(x_0)|<1\}$
funkcji zerowej nie zawiera
żadnego $V_n$: dla każdego $n$
można wybrać funkcję równą $0$
na $F_n$ i równą $2$ w $x_0$.
Zatem żadne $U_n$ nie zawiera się
w $W$, bo $V_n\subseteq U_n$.
Otrzymujemy sprzeczność.
W szczególności $C_k(X,\R)$
nie jest metryzowalna.
Pokazuje to, dlaczego założenia
o istnieniu przeliczalnego
wyczerpania zwartego są istotne.
\end{przyklad}

\begin{uwaga}[Związki z dalszą teorią]
Topologia zwarto-otwarta jest
naturalna przy badaniu przestrzeni
homotopii, odwzorowań i rodzin
parametryzowanych.
Zbieżność jednostajna na zbiorach
zwartych łączy topologię ogólną
z analizą funkcjonalną i teorią
przestrzeni funkcji.
W rozdziale~19 przejdziemy do
przestrzeni polskich i zbiorów
borelowskich, korzystając m.in.
z metryzacji i zupełności
udowodnionych tutaj.
\end{uwaga}

\chapter{Przestrzenie polskie i zbiory borelowskie}
\label{ch:topog-polskie-borelowskie}

\noindent\textbf{Cel rozdziału.}
Zupełność metryczna z rozdziału~11 i metryzowalność
z rozdziału~10 prowadzą do pojęcia przestrzeni
polskiej. Zbadamy jej trwałość przy podstawowych
konstrukcjach oraz scharakteryzujemy te podprzestrzenie
przestrzeni polskich, które same są polskie.
Następnie zbudujemy $\sigma$-algebrę zbiorów
borelowskich oraz jej hierarchię.
Od rozdziału~20 oddzielimy bardziej zaawansowane
zagadnienia opisowej teorii mnogości, w szczególności
zbiory analityczne i twierdzenia o ich strukturze.

\section{Przestrzenie polskie: definicja i przykłady}
\label{sec:topog-polskie-definicja}

\begin{definicja}[Przestrzeń polska]
\label{def:topog-polskie}
Przestrzeń topologiczną $X$ nazywamy
\emph{przestrzenią polską}, jeśli istnieje
metryka $d$ wyznaczająca jej topologię,
taka, że $(X,d)$ jest zupełna,
oraz $X$ ma przeliczalny zbiór gęsty.
Nie wymagamy, aby każda metryka
zgodna z topologią $X$ była zupełna.
\end{definicja}

\begin{stwierdzenie}[Drugi aksjomat przeliczalności]
\label{prop:topog-polskie-drugi-aksjomat}
Każda przestrzeń polska spełnia
drugi aksjomat przeliczalności
i jest przestrzenią Hausdorffa.
\end{stwierdzenie}
\begin{proof}
Wybierzmy zgodną metrykę $d$
i przeliczalny zbiór gęsty $D\subseteq X$.
Rodzina
$\{B_d(a,r):a\in D,\ r\in\mathbb Q,\ r>0\}$
jest przeliczalna. Jest bazą:
jeśli $x\in U$ i $U$ jest otwarty,
istnieje $\varepsilon>0$ z
$B_d(x,\varepsilon)\subseteq U$.
Wybierzmy $a\in D$ z
$d(x,a)<\varepsilon/2$
oraz wymierne $r$ z
$d(x,a)<r<\varepsilon-d(x,a)$.
Wówczas
$x\in B_d(a,r)\subseteq U$.
Hausdorffowość wynika z metryzowalności.
\end{proof}

\begin{przyklad}[Znane przestrzenie polskie]
\label{ex:topog-polskie-podstawowe}
Przestrzeń $\mathbb R^m$ ze zwykłą
metryką euklidesową jest zupełna,
a $\mathbb Q^m$ jest gęstym zbiorem
przeliczalnym. Każda jej domknięta
podprzestrzeń, z metryką ograniczoną,
jest również zupełna i ośrodkowa.

Przedział $(0,1)$ jest polski,
chociaż metryka euklidesowa
ograniczona do $(0,1)$ nie jest
zupełna. Zgodną metrykę zupełną
można otrzymać przez zanurzenie
$t\mapsto\log(t/(1-t))$
w prostą: odwzorowanie to jest
homeomorfizmem $(0,1)$ na $\mathbb R$,
a więc metryka
\[
d_\ast(s,t)=
\left|\log\frac{s}{1-s}
-\log\frac{t}{1-t}\right|
\]
jest zgodna i zupełna.
\end{przyklad}

\begin{przyklad}[Przestrzenie Cantora i Baire'a]
\label{ex:topog-polskie-cantor-baire}
Niech $A$ oznacza odpowiednio
$\{0,1\}$ lub $\mathbb N$,
w obu przypadkach z topologią dyskretną.
Na zbiorze $A^{\mathbb N}$ ciągów
określmy metrykę
\[
d_A(x,y)=
\begin{cases}
0,&x=y,\\
2^{-k},&k=\min\{n\ge1:x_n\ne y_n\}.
\end{cases}
\]
Jest to ultrametryka: jeżeli
$x_k=y_k$ i $y_k=z_k$
dla każdego $k<m$, to również
$x_k=z_k$ dla tych indeksów;
stąd
$d_A(x,z)\le\max\{d_A(x,y),d_A(y,z)\}$.
Kule o odpowiednio małym promieniu
wyznaczają dokładnie warunek
zgodności skończonej liczby
początkowych współrzędnych,
zatem $d_A$ wyznacza topologię
produktową.

Przestrzeń ta jest zupełna.
Jeżeli $(x^{(j)})_{j\ge1}$ jest
ciągiem Cauchy'ego, to dla każdego
$n$ współrzędne
$x_1^{(j)},\ldots,x_n^{(j)}$
są ostatecznie stałe:
wystarczy wykorzystać oszacowanie
$d_A(x^{(j)},x^{(l)})<2^{-n}$.
Dla każdej współrzędnej $k$
jej ostateczna wartość jest
jednoznacznie określona.
Wartości te definiują
$x\in A^{\mathbb N}$.
Dla ustalonego $n$ wyrazy
ciągu od pewnego miejsca
zgadzają się z $x$
na pierwszych $n$ współrzędnych,
zatem ich odległość od $x$
jest nie większa niż
$2^{-(n+1)}$.
Ponieważ $n$ jest dowolne,
otrzymujemy zbieżność do $x$.

Ciągi, których wszystkie współrzędne
od pewnego miejsca są równe
ustalonemu $a_0\in A$,
tworzą zbiór przeliczalny i gęsty.
Dla $A=\{0,1\}$ otrzymujemy
\emph{przestrzeń Cantora}
$2^{\mathbb N}$, a dla
$A=\mathbb N$ ---
\emph{przestrzeń Baire'a}
$\mathbb N^{\mathbb N}$.
Obie są polskie; pierwsza jest
zwarta na mocy twierdzenia Tychonoffa,
a druga nie jest zwarta, ponieważ
jej obraz przez rzutowanie na
pierwszą współrzędną jest
nieskończoną dyskretną przestrzenią
$\mathbb N$, która nie jest zwarta.
\end{przyklad}

\begin{twierdzenie}[Przeliczalne iloczyny przestrzeni polskich]
\label{thm:topog-polskie-iloczyn}
Niech $(X_n)_{n\ge1}$ będzie ciągiem
przestrzeni polskich.
Wtedy przestrzeń
$X=\prod_{n\ge1}X_n$
z topologią produktową jest polska.
\end{twierdzenie}
\begin{proof}
Jeśli któryś czynnik jest pusty,
to $X=\varnothing$, a przestrzeń
pusta jest polska. Załóżmy więc,
że wszystkie $X_n$ są niepuste.
Wybierzmy na $X_n$ zgodną
zupełną metrykę $d_n$ i
przeliczalny zbiór gęsty $D_n$.
Dla $x=(x_n)$ oraz $y=(y_n)$
określmy
\[
d(x,y)=\sum_{n=1}^{\infty}
2^{-n}\min\{1,d_n(x_n,y_n)\}.
\]
Nierówność trójkąta wynika z
$\min\{1,a+b\}\le
\min\{1,a\}+\min\{1,b\}$,
a dodatniość metryki z faktu,
że różne punkty mają różną
co najmniej jedną współrzędną.
Szereg jest zbieżny, gdyż każdy
składnik jest ograniczony przez $2^{-n}$.

Topologia metryki $d$ jest produktowa.
Zbieżność w $d$ wymusza zbieżność
każdej współrzędnej, ponieważ
$2^{-n}\min\{1,d_n(x_n,y_n)\}\le d(x,y)$.
W drugą stronę skończenie wiele
współrzędnych można kontrolować
równocześnie, a ogon szeregu
$\sum_{n>N}2^{-n}$ uczynić
dowolnie małym. To dowodzi
zgodności topologii.

Jeśli $(x^{(j)})$ jest ciągiem
Cauchy'ego w $d$, każda
$(x_n^{(j)})$ jest ciągiem
Cauchy'ego w $d_n$, więc
ma granicę $x_n\in X_n$.
Oszacowanie ogona szeregu
pokazuje następnie, że
$x^{(j)}\to(x_n)$ w $d$.
Zatem $d$ jest zupełna.

Wybierzmy punkt bazowy
$a=(a_n)\in X$.
Zbiór ciągów różniących się od
$a$ tylko na skończenie wielu
współrzędnych, przyjmujących
tam wartości z odpowiednich $D_n$,
jest przeliczalny jako
przeliczalna suma skończonych
iloczynów zbiorów przeliczalnych.
Jest gęsty, ponieważ bazowy
zbiór otwarty topologii produktowej
ogranicza tylko skończenie wiele
współrzędnych. To kończy dowód.
\end{proof}

\section{Podprzestrzenie polskie i zbiory typu \texorpdfstring{$G_\delta$}{G-delta}}
\label{sec:topog-polskie-gdelta}

\begin{definicja}[Zbiory typu F-sigma i G-delta]
\label{def:topog-polskie-fsigma-gdelta}
Niech $X$ będzie przestrzenią
topologiczną. Zbiór $A\subseteq X$
jest typu $F_\sigma$, jeżeli
istnieją zbiory domknięte
$F_n\subseteq X$ takie, że
$A=\bigcup_{n\ge1}F_n$.
Jest typu $G_\delta$, jeżeli
istnieją otwarte $U_n\subseteq X$
takie, że $A=\bigcap_{n\ge1}U_n$.
\end{definicja}

\begin{lemat}[Zbiory otwarte i domknięte w przestrzeni metrycznej]
\label{lem:topog-polskie-otwarte-fsigma}
Niech $(X,d)$ będzie przestrzenią
metryczną. Każdy jej zbiór otwarty
jest typu $F_\sigma$, a każdy
zbiór domknięty jest typu $G_\delta$.
\end{lemat}
\begin{proof}
Niech $U\subseteq X$ będzie otwarty
oraz $F=X\setminus U$.
Jeśli $F=\varnothing$, mamy
$U=X$, który jest domknięty.
W przeciwnym przypadku
dla $n\ge1$ połóżmy
$F_n=\{x\in X:d(x,F)\ge1/n\}$.
Funkcja $x\mapsto d(x,F)$
jest ciągła, zatem $F_n$ jest
domknięty. Ponieważ dla
$x\in U$ istnieje kula
$B_d(x,\varepsilon)\subseteq U$,
mamy $d(x,F)>0$ i punkt $x$
należy do pewnego $F_n$.
Jeżeli $x\in F$, to
$d(x,F)=0$ i $x$ nie należy
do żadnego $F_n$. Stąd
$U=\bigcup_{n\ge1}F_n$.
Dla domkniętego $C$ stosujemy
pierwszą część do otwartego
$X\setminus C$ i bierzemy
dopełnienia.
\end{proof}

\begin{twierdzenie}[Charakterystyka polskich podprzestrzeni]
\label{thm:topog-polskie-gdelta}
Niech $X$ będzie przestrzenią polską
i $A\subseteq X$. Podprzestrzeń
$A$ jest polska wtedy i tylko wtedy,
gdy $A$ jest zbiorem typu $G_\delta$
w $X$.
\end{twierdzenie}
\begin{proof}
Niech najpierw
$A=\bigcap_{n\ge1}U_n$,
gdzie $U_n$ są otwarte w $X$.
Wybierzmy zupełną zgodną metrykę
$d$ na $X$. Gdy $A=\varnothing$,
twierdzenie jest oczywiste.
Dla każdego $n$ z
$X\setminus U_n\ne\varnothing$
określamy ciągłą na $A$ funkcję
$h_n(x)=1/d(x,X\setminus U_n)$.
Gdy $U_n=X$, przyjmujemy $h_n=0$.
Mianowniki są dodatnie dla
$x\in A$, ponieważ $U_n$ jest
otwarty i $x\in U_n$.
Na $A$ określmy
\[
\rho(x,y)=d(x,y)+
\sum_{n\ge1}2^{-n}
\min\{1,|h_n(x)-h_n(y)|\}.
\]
Jest to metryka: składniki
są pseudometrykami, a pierwszy
jest metryką. Wyznacza
topologię podprzestrzeni $A$.
Istotnie, $\rho\ge d$, zatem
zbieżność w $\rho$ implikuje
zbieżność w $d$; odwrotnie,
ciągłość każdego $h_n$ i
jednostajnie mały ogon szeregu
zapewniają zbieżność w $\rho$
przy zbieżności w $d$:
dla ustalonego $\varepsilon>0$
wybieramy $N$ z
$\sum_{n>N}2^{-n}<\varepsilon/2$,
a następnie kontrolujemy
skończenie wiele różnic
$|h_n(x_j)-h_n(x)|$, $n\le N$.
W przestrzeni metrycznej
zgodność zbiegających ciągów
do wszystkich punktów wystarcza
do równości topologii.

Udowodnimy zupełność.
Niech $(x_j)$ będzie ciągiem
Cauchy'ego w $\rho$.
Wówczas jest ciągiem Cauchy'ego
w $d$, a więc $x_j\to x$
dla pewnego $x\in X$.
Dla każdego ustalonego $n$
ciąg liczb $(h_n(x_j))$
jest ciągiem Cauchy'ego.
Istotnie, dla
$0<\varepsilon<1$ wybieramy
$j_0$ tak, by
$\rho(x_j,x_k)<2^{-n}\varepsilon$
dla wszystkich $j,k\ge j_0$.
Wtedy
$\min\{1,|h_n(x_j)-h_n(x_k)|\}
<\varepsilon<1$,
a więc
$|h_n(x_j)-h_n(x_k)|<\varepsilon$.
To wystarcza również dla
większych wartości $\varepsilon$.
Jeśli $x\notin U_n$,
to $d(x_j,X\setminus U_n)
\le d(x_j,x)\to0$, przez co
$h_n(x_j)\to+\infty$,
wbrew własności Cauchy'ego
(gdy $U_n=X$, taka sytuacja
nie występuje).
Zatem $x\in\bigcap_n U_n=A$.
Dla każdego $n$ ciągłość
$h_n$ na $A$ daje
$h_n(x_j)\to h_n(x)$,
ponieważ $x_j\to x$ w $d$
i wykazaliśmy $x\in A$.
Dla zadanego $\varepsilon>0$
dobieramy $N$ z
$\sum_{n>N}2^{-n}<\varepsilon/3$,
a następnie $j_0$ tak, by
$d(x_j,x)<\varepsilon/3$
oraz
$\sum_{n=1}^N2^{-n}
\min\{1,|h_n(x_j)-h_n(x)|\}
<\varepsilon/3$
dla $j\ge j_0$.
Wówczas $\rho(x_j,x)<\varepsilon$.
Zatem $(A,\rho)$ jest zupełna.
Podprzestrzeń przestrzeni ośrodkowej
metrycznej jest ośrodkowa,
ponieważ spełnia drugi aksjomat
przeliczalności jako podprzestrzeń
przestrzeni z przeliczalną bazą,
a wybór po jednym punkcie z
każdego niepustego zbioru tej bazy
daje przeliczalny zbiór gęsty.
Stąd $A$ jest polska.

Odwrotnie, załóżmy, że
$A$ jest polska.
Wybierzmy zgodną metrykę $d$
na $X$ i zgodną zupełną
metrykę $\sigma$ na $A$.
Dla każdego $a\in A$ i $n\ge1$
istnieje $r_n(a)>0$ takie, że
\[
r_n(a)<2^{-n},\qquad
\operatorname{diam}_{\sigma}
\bigl(B_d(a,2r_n(a))\cap A\bigr)
<2^{-n}.
\]
Istnienie wynika z ciągłości
identyczności
$(A,d|_A)\to(A,\sigma)$:
można wybrać tak małą kulę,
aby wszystkie jej punkty z $A$
leżały w kuli
$B_\sigma(a,2^{-n-2})$.
Niech
\[
V_n=\bigcup_{a\in A}
 B_d(a,r_n(a)).
\]
Są to zbiory otwarte w $X$
i $A\subseteq\bigcap_n V_n$.

Pokażemy inkluzję odwrotną.
Ustalmy $x\in\bigcap_nV_n$.
Dla każdego $n$ wybierzmy
$a_n\in A$ z
$d(x,a_n)<r_n(a_n)$.
Ponieważ $r_n(a_n)<2^{-n}$,
mamy $a_n\to x$ w metryce $d$.
Ustalmy $n$ oraz połóżmy
$\delta_n=r_n(a_n)-d(x,a_n)>0$.
Dla wystarczająco dużych $m$
mamy $r_m(a_m)<\delta_n$,
a więc
\[
d(a_m,a_n)
\le d(a_m,x)+d(x,a_n)
<r_m(a_m)+d(x,a_n)
<r_n(a_n).
\]
Wszystkie dostatecznie późne
$a_m$ leżą zatem w
$B_d(a_n,2r_n(a_n))\cap A$,
którego średnica w metryce
$\sigma$ jest mniejsza niż
$2^{-n}$. Ustalmy teraz dowolne
$\varepsilon>0$ i wybierzmy
$n$ takie, że $2^{-n}<\varepsilon$.
Z poprzedniego oszacowania
istnieje $m_0$ takie, że
dla wszystkich $m\ge m_0$
mamy
$a_m\in B_d(a_n,2r_n(a_n))\cap A$.
Zatem dla dowolnych
$m,l\ge m_0$ zachodzi
$\sigma(a_m,a_l)<2^{-n}<\varepsilon$.
To dowodzi, że $(a_m)$
jest ciągiem Cauchy'ego w $\sigma$. Z zupełności
otrzymujemy granicę $a\in A$.
Metryki $\sigma$ i $d|_A$
wyznaczają tę samą topologię,
więc $a_m\to a$ również w $d$.
Jednoznaczność granicy w
metrycznej $X$ i zbieżność
$a_m\to x$ dają $x=a\in A$.
Zatem $A=\bigcap_nV_n$.
\end{proof}

\begin{przyklad}[Liczby wymierne i niewymierne]
\label{ex:topog-polskie-wymierne}
Zbiór $\mathbb Q\subseteq\mathbb R$
jest typu $F_\sigma$, ponieważ
jest przeliczalną sumą singletonów
domkniętych. Nie jest typu
$G_\delta$ w $\mathbb R$.
Rzeczywiście, jako przestrzeń
z topologią względną jest
przeliczalna, a żaden jej
singleton nie ma niepustego
wnętrza. Gdyby $\mathbb Q$
była polska, twierdzenie Baire'a
z rozdziału~11 zabraniałoby
przedstawienia jej jako
przeliczalnej sumy zbiorów
nigdziegęstych. Stąd
$\mathbb Q$ nie jest polska
i z poprzedniego twierdzenia
nie jest typu $G_\delta$.

Zbiór liczb niewymiernych
$\mathbb R\setminus\mathbb Q$
jest typu $G_\delta$, a więc
jest przestrzenią polską.
Nie jest typu $F_\sigma$:
w przeciwnym razie jego
dopełnienie $\mathbb Q$
byłoby typu $G_\delta$.
\end{przyklad}

\section{Zbiory borelowskie i poziomy hierarchii}
\label{sec:topog-polskie-borel}

\begin{definicja}[Sigma-algebra borelowska]
\label{def:topog-polskie-borel}
Niech $X$ będzie przestrzenią
topologiczną. \emph{Sigma-algebrą
borelowską} $\mathcal B(X)$
nazywamy najmniejszą
$\sigma$-algebrę podzbiorów $X$
zawierającą wszystkie zbiory
otwarte. Jej elementy nazywamy
\emph{zbiorami borelowskimi}.
\end{definicja}

\begin{stwierdzenie}[Zbiory borelowskie w podprzestrzeni]
\label{prop:topog-polskie-borel-slad}
Niech $X$ będzie dowolną
przestrzenią topologiczną,
a $A\subseteq X$ podprzestrzenią.
Wtedy
\[
\mathcal B(A)=
\{A\cap B:B\in\mathcal B(X)\}.
\]
\end{stwierdzenie}
\begin{proof}
Rodzina po prawej stronie
jest $\sigma$-algebrą na $A$,
ponieważ dopełnienia względne
i przeliczalne sumy zachowują
postać $A\cap B$.
Zawiera wszystkie otwarte w $A$
zbiory $A\cap U$, gdzie $U$
jest otwarty w $X$.
Dlatego $\mathcal B(A)$ jest
w niej zawarta.

Odwrotnie, rodzina
$\mathcal C=\{B\subseteq X:
A\cap B\in\mathcal B(A)\}$
jest $\sigma$-algebrą na $X$
i zawiera wszystkie otwarte
zbiory $U\subseteq X$,
ponieważ $A\cap U$ jest
otwarty w $A$.
Stąd $\mathcal B(X)\subseteq
\mathcal C$ i każde $A\cap B$
z $B\in\mathcal B(X)$
należy do $\mathcal B(A)$.
\end{proof}

\begin{stwierdzenie}[Przeliczalne generatory sigma-algebry borelowskiej]
\label{prop:topog-polskie-borel-przeliczalna-baza}
Niech $X$ będzie przestrzenią
topologiczną spełniającą
drugi aksjomat przeliczalności
i niech $\mathcal U$ będzie
jej bazą przeliczalną.
Wówczas
\[
\mathcal B(X)=\sigma(\mathcal U),
\]
gdzie po prawej stronie oznaczamy
$\sigma$-algebrę generowaną
przez rodzinę $\mathcal U$.
W szczególności $\mathcal B(\mathbb R)$
jest generowana przez
przedziały otwarte $(a,b)$
o końcach wymiernych $a<b$.
\end{stwierdzenie}
\begin{proof}
Ponieważ wszystkie elementy
$\mathcal U$ są otwarte,
mamy $\sigma(\mathcal U)
\subseteq\mathcal B(X)$.
Każdy zbiór otwarty w $X$
jest sumą elementów $\mathcal U$;
rodzina $\mathcal U$ jest
przeliczalna, więc suma ta
jest co najwyżej przeliczalna.
Zatem każdy zbiór otwarty
należy do $\sigma(\mathcal U)$,
co daje inkluzję odwrotną.
Przedziały o końcach wymiernych
tworzą przeliczalną bazę
zwykłej topologii na $\mathbb R$.
\end{proof}

\begin{definicja}[Odwzorowanie borelowskie]
\label{def:topog-polskie-odwzorowanie-borelowskie}
Niech $X,Y$ będą przestrzeniami
topologicznymi. Odwzorowanie
$f\colon X\to Y$ nazywamy
\emph{borelowskim}, jeżeli
$f^{-1}(B)\in\mathcal B(X)$
dla każdego $B\in\mathcal B(Y)$.
\end{definicja}

\begin{stwierdzenie}[Kryterium borelowskości]
\label{prop:topog-polskie-borelowskie-kryterium}
Niech $X,Y$ będą przestrzeniami
topologicznymi, a
$f\colon X\to Y$ odwzorowaniem.
Jest ono borelowskie wtedy
i tylko wtedy, gdy
$f^{-1}(U)\in\mathcal B(X)$
dla każdego otwartego
$U\subseteq Y$.
Każde odwzorowanie ciągłe
jest borelowskie, a złożenie
dwóch odwzorowań borelowskich
jest borelowskie.
\end{stwierdzenie}
\begin{proof}
Konieczność jest oczywista.
Dla wystarczalności rozważmy
rodzinę
$\mathcal C=\{B\subseteq Y:
f^{-1}(B)\in\mathcal B(X)\}$.
Przeciwobrazy zachowują
dopełnienia i przeliczalne
sumy, więc $\mathcal C$ jest
$\sigma$-algebrą. Jeśli
zawiera wszystkie otwarte
zbiory $Y$, to zawiera
$\mathcal B(Y)$. To daje
borelowskość $f$.
Gdy $f$ jest ciągłe,
przeciwobrazy otwartych
są otwarte, więc borelowskie.
Jeśli $f\colon X\to Y$
i $g\colon Y\to Z$ są
borelowskie, to dla
$B\in\mathcal B(Z)$
mamy
$(g\circ f)^{-1}(B)
=f^{-1}(g^{-1}(B))
\in\mathcal B(X)$.
\end{proof}

\begin{definicja}[Skończone poziomy hierarchii borelowskiej]
\label{def:topog-polskie-hierarchia-skonczona}
Niech $X$ będzie przestrzenią
metryzowalną. Definiujemy
$\boldsymbol\Sigma^0_1(X)$
jako rodzinę wszystkich
zbiorów otwartych,
a $\boldsymbol\Pi^0_1(X)$
jako rodzinę wszystkich
zbiorów domkniętych.
Dla $n\ge1$ określamy
\[
\begin{aligned}
\boldsymbol\Sigma^0_{n+1}(X)
&=\left\{\bigcup_{k\ge1}A_k:
A_k\in\boldsymbol\Pi^0_n(X)\right\},\\
\boldsymbol\Pi^0_{n+1}(X)
&=\left\{\bigcap_{k\ge1}A_k:
A_k\in\boldsymbol\Sigma^0_n(X)\right\}.
\end{aligned}
\]
Są to kolejno rodziny
zbiorów otwartych i domkniętych,
zbiorów typu $F_\sigma$
i $G_\delta$, a następnie
ich dalszych przeliczalnych
sum i przecięć.
\end{definicja}

\begin{stwierdzenie}[Dualność i operacje]
\label{prop:topog-polskie-hierarchia-wlasnosci}
Niech $X$ będzie przestrzenią
metryzowalną, a $n\ge1$.
Wówczas
\[
A\in\boldsymbol\Sigma^0_n(X)
\ \Longleftrightarrow\
X\setminus A\in\boldsymbol\Pi^0_n(X).
\]
Ponadto rodzina
$\boldsymbol\Sigma^0_n(X)$
jest zamknięta na przeliczalne
sumy, a rodzina
$\boldsymbol\Pi^0_n(X)$
na przeliczalne przecięcia.
Dla każdego $n\ge1$
zachodzą inkluzje
\[
\boldsymbol\Sigma^0_n(X)
\cup\boldsymbol\Pi^0_n(X)
\subseteq
\boldsymbol\Sigma^0_{n+1}(X)
\cap\boldsymbol\Pi^0_{n+1}(X).
\]
\end{stwierdzenie}
\begin{proof}
Dualność dla $n=1$
jest definicyjna.
Jeżeli zachodzi na poziomie $n$,
to prawa de Morgana zamieniają
przeliczalną sumę zbiorów
z $\boldsymbol\Pi^0_n$
w przeliczalne przecięcie
ich dopełnień z
$\boldsymbol\Sigma^0_n$,
co daje dualność na poziomie
$n+1$.

Zamknięcie $\boldsymbol\Sigma^0_1$
na przeliczalne sumy
wynika z aksjomatów topologii.
Na poziomach $n+1$ sumę
przeliczalnie wielu przeliczalnych
sum zbiorów z
$\boldsymbol\Pi^0_n$
można zapisać jako jedną
przeliczalną sumę, indeksując
składniki przez $\mathbb N^2$.
Zamknięcie rodziny
$\boldsymbol\Pi^0_n$ na
przeliczalne przecięcia
otrzymujemy przez dopełnienia.

Lemat
\ref{lem:topog-polskie-otwarte-fsigma}
daje
$\boldsymbol\Sigma^0_1
\subseteq\boldsymbol\Sigma^0_2$
i
$\boldsymbol\Pi^0_1
\subseteq\boldsymbol\Pi^0_2$.
Ponadto każdy zbiór z
$\boldsymbol\Sigma^0_1$ jest
jednoelementowym przeliczalnym
przecięciem zbiorów otwartych,
zatem należy do
$\boldsymbol\Pi^0_2$;
analogicznie
$\boldsymbol\Pi^0_1
\subseteq\boldsymbol\Sigma^0_2$.
Przyjmijmy indukcyjnie obie
inkluzje własnych poziomów
$n\to n+1$.
Jeżeli zbiór jest przeliczalną
sumą zbiorów z
$\boldsymbol\Pi^0_n$, to
każdy składnik należy także do
$\boldsymbol\Pi^0_{n+1}$,
co daje zawieranie
$\boldsymbol\Sigma^0_{n+1}
\subseteq\boldsymbol\Sigma^0_{n+2}$.
Druga inkluzja wynika przez
dopełnienia; inkluzje krzyżowe
na każdym poziomie wynikają
z zapisania zbioru jako
przeliczalnej sumy lub
przecięcia jednego zbioru.
To kończy indukcję.
\end{proof}

\begin{uwaga}[Poziomy skończone nie wyczerpują hierarchii]
Rodziny
$\boldsymbol\Sigma^0_n$
i $\boldsymbol\Pi^0_n$,
$n\in\mathbb N$, opisują
jedynie skończone poziomy
hierarchii borelowskiej.
W ogólności pojawiają się
również zbiory borelowskie
o wyższych, przeliczalnych
rangach. Szczegółowe twierdzenia
o ścisłości hierarchii wymagają
dalszych narzędzi opisowej teorii
mnogości.
Nie należy utożsamiać całej
$\sigma$-algebry borelowskiej
z sumą jej poziomów skończonych.
\end{uwaga}

\section{Pełna hierarchia borelowska}
\label{sec:topog-polskie-hierarchia-pozaskonczona}

\begin{definicja}[Hierarchia borelowska indeksowana liczbami porządkowymi]
\label{def:topog-polskie-hierarchia-porzadkowa}
Niech $X$ będzie przestrzenią
metryzowalną. Dla każdej
niezerowej przeliczalnej
liczby porządkowej
$\alpha<\omega_1$
definiujemy przez indukcję
pozaskończoną rodziny
$\boldsymbol\Sigma^0_\alpha(X)$
i $\boldsymbol\Pi^0_\alpha(X)$.
Dla $\alpha=1$ są to
odpowiednio zbiory otwarte
i domknięte. Dla $\alpha>1$
przyjmujemy
\[
\begin{aligned}
\boldsymbol\Sigma^0_\alpha(X)
&=\left\{\bigcup_{n\ge1}A_n:
\text{ dla każdego }n\
A_n\in\boldsymbol\Pi^0_{\beta_n}(X)
\text{ dla pewnego }1\le\beta_n<\alpha
\right\},\\
\boldsymbol\Pi^0_\alpha(X)
&=\{X\setminus A:
A\in\boldsymbol\Sigma^0_\alpha(X)\}.
\end{aligned}
\]
Symbol $\omega_1$ oznacza
pierwszą nieprzeliczalną
liczbę porządkową.
Na poziomach skończonych
ta definicja zgadza się
z definicją poprzednią:
w przestrzeni metryzowalnej
rodziny $\boldsymbol\Pi^0_n$
są rosnące względem $n$,
więc suma przeliczalna
zbiorów z wcześniejszych
poziomów $\boldsymbol\Pi^0_\beta$,
$\beta<n+1$, jest taką samą
sumą zbiorów z
$\boldsymbol\Pi^0_n$.
\end{definicja}

\begin{twierdzenie}[Wyczerpanie sigma-algebry borelowskiej]
\label{thm:topog-polskie-borel-wyczerpanie}
Niech $X$ będzie przestrzenią
metryzowalną. Wówczas
\[
\mathcal B(X)=
\bigcup_{1\le\alpha<\omega_1}
\boldsymbol\Sigma^0_\alpha(X)
=
\bigcup_{1\le\alpha<\omega_1}
\boldsymbol\Pi^0_\alpha(X).
\]
\end{twierdzenie}
\begin{proof}
Każdy poziom hierarchii
jest zawarty w $\mathcal B(X)$,
co wynika przez indukcję
pozaskończoną z definicji
$\sigma$-algebry.

Oznaczmy przez $\mathcal C$
sumę poziomów
$\boldsymbol\Sigma^0_\alpha$.
Zawiera ona wszystkie zbiory
otwarte i zbiór $X$.
Jeżeli
$A\in\boldsymbol\Sigma^0_\alpha$,
to $X\setminus A$ należy do
$\boldsymbol\Pi^0_\alpha$,
a więc do
$\boldsymbol\Sigma^0_{\alpha+1}$:
można przedstawić go jako
przeliczalną sumę kopii
samego siebie, a
$\alpha<\alpha+1$.
Zatem $\mathcal C$ jest
zamknięta na dopełnienia.

Niech teraz
$A_n\in\mathcal C$,
a $\alpha_n<\omega_1$
będzie rangą taką, że
$A_n\in\boldsymbol\Sigma^0_{\alpha_n}$.
Na mocy poprzedniego argumentu
dopełnienie $A_n$ należy do
$\boldsymbol\Sigma^0_{\alpha_n+1}$,
ponieważ należy do
$\boldsymbol\Pi^0_{\alpha_n}$.
Po ponownym wzięciu dopełnienia
otrzymujemy
$A_n\in\boldsymbol\Pi^0_{\alpha_n+1}$.
Niech $\gamma$ będzie przeliczalną
liczbą porządkową większą od
każdej $\alpha_n+1$
(np. $\gamma=\sup_n(\alpha_n+2)$).
Wtedy każdy $A_n$ należy
do któregoś
$\boldsymbol\Pi^0_{\beta_n}$
z $\beta_n<\gamma$, a zatem
$\bigcup_nA_n$ należy do
$\boldsymbol\Sigma^0_\gamma$.
Ponieważ supremum przeliczalnie
wielu przeliczalnych liczb
porządkowych jest przeliczalne,
mamy $\gamma<\omega_1$.
Dowodzi to, że $\mathcal C$
jest $\sigma$-algebrą.
Jako $\sigma$-algebra
zawierająca wszystkie otwarte
musi zawierać $\mathcal B(X)$,
a więc jest jej równa.
Równość z sumą poziomów
$\boldsymbol\Pi^0_\alpha$
wynika teraz z zamknięcia
$\mathcal B(X)$ na dopełnienia
i z dualności poziomów.
\end{proof}

\begin{uwaga}[Zakres dalszego rozdziału]
W rozdziale~20 zbadamy zbiory
analityczne, obrazy ciągłe
przestrzeni polskich, własności
regularności i dalsze twierdzenia
opisowej teorii mnogości.
Hierarchia borelowska będzie
punktem odniesienia dla tych
konstrukcji; nie zastępujemy
jej samymi poziomami skończonymi.
\end{uwaga}

\chapter{Zbiory analityczne i operacja Suslina}
\label{ch:topog-analityczne-suslin}

\noindent\textbf{Cel rozdziału.}
W rozdziale~\ref{ch:topog-polskie-borelowskie}
poznaliśmy przestrzenie polskie i zbiory borelowskie.
Teraz zbadamy zbiory powstające jako obrazy ciągłe
przestrzeni polskich i rzuty zbiorów domkniętych.
Wykażemy równoważność tych opisów, wprowadzimy
operację Suslina i skonstruujemy zbiór analityczny,
który nie jest borelowski.

Nie powtarzamy rozdziału~30 skryptu
\emph{Logika matematyczna i teoria mnogości},
poświęconego grom nieskończonym, aksjomatowi
determinacji, klasom projektywnym i redukcjom Wadge'a.
Podajemy natomiast pełne dowody topologicznych
faktów, które tam występują jedynie jako
wprowadzenie lub wyniki przywołane bez dowodu.
W całym rozdziale pracujemy w zwykłym
ujęciu teorii mnogości ZFC.

\section{Przestrzeń Baire'a jako przestrzeń parametryczna}
\label{sec:topog20-baire}

Przez \(\mathcal N=\mathbb N^{\mathbb N}\)
oznaczamy przestrzeń Baire'a z topologią
produktową, opisaną w
Przykładzie~\ref{ex:topog-polskie-cantor-baire}.
Indeksy współrzędnych zaczynamy od \(1\).
Zbiór wszystkich skończonych ciągów
liczb naturalnych oznaczamy
\(\mathbb N^{<\mathbb N}\).
Dla \(s=(s_1,\ldots,s_k)\) definiujemy
\[
 [s]=\{x\in\mathcal N:
             x_i=s_i\text{ dla }1\le i\le k\}.
\]
Przyjmujemy \([\varnothing]=\mathcal N\).
Zbiory \([s]\) są otwarte i domknięte,
a rodzina wszystkich \([s]\)
jest przeliczalną bazą topologii
\(\mathcal N\).
Dla \(x\in\mathcal N\) zapisujemy
\(x|k=(x_1,\ldots,x_k)\).

\begin{lemat}[Parametryzacja przestrzeni polskiej]
\label{lem:topog20-surjection}
Niech \(X\) będzie niepustą przestrzenią
polską. Istnieje ciągłe odwzorowanie
surjektywne \(f\colon\mathcal N\to X\).
\end{lemat}
\begin{proof}
Wybierzmy na \(X\) zgodną zupełną
metrykę \(d\) oraz przeliczalny
gęsty zbiór \(D\).
Zbudujemy rekurencyjnie niepuste
domknięte zbiory \(F_s\subseteq X\),
\(s\in\mathbb N^{<\mathbb N}\),
spełniające
\[
 F_{\varnothing}=X,\qquad
 F_{s^\frown k}\subseteq F_s,\qquad
 F_s=\bigcup_{k\in\mathbb N}F_{s^\frown k},
\]
oraz
\(\operatorname{diam}_d(F_s)\le2^{-|s|}\)
dla \(|s|\ge1\).

Załóżmy, że zdefiniowano
niepusty domknięty \(F_s\)
i niech \(l=|s|\).
Wszystkie niepuste zbiory postaci
\[
 F_s\cap\overline B_d(a,2^{-l-3}),
 \qquad a\in D,
\]
są domknięte, zawarte w \(F_s\)
i mają średnicę nie większą
niż \(2^{-l-2}\).
Pokrywają \(F_s\): dla \(x\in F_s\)
z gęstości \(D\) można wybrać
\(a\in D\) z \(d(x,a)<2^{-l-3}\).
Tę niepustą przeliczalną rodzinę
enumerujemy jako
\((F_{s^\frown k})_{k\in\mathbb N}\);
jeśli jest skończona,
powtarzamy jej elementy.
To kończy konstrukcję indukcyjną.

Dla \(\alpha\in\mathcal N\)
zbiory \(F_{\alpha|k}\)
tworzą malejącą rodzinę
niepustych domkniętych
podzbiorów zupełnej przestrzeni
metrycznej, a ich średnice
dążą do zera.
Wybierzmy \(x_k\in F_{\alpha|k}\).
Dla \(p,q\ge k\) punkty
\(x_p,x_q\) należą do
\(F_{\alpha|k}\), zatem
\(d(x_p,x_q)\le2^{-k}\).
Ciąg \((x_k)\) jest więc
ciągiem Cauchy'ego i ma
granicę \(x\in X\).
Domkniętość każdego
\(F_{\alpha|k}\) zapewnia
\(x\in F_{\alpha|k}\).
Średnice pokazują zarazem
jednoznaczność takiego \(x\).
Definiujemy \(f(\alpha)=x\).

Jeżeli \(\alpha,\beta\)
mają ten sam początkowy
fragment długości \(k\), to
\(f(\alpha)\) i \(f(\beta)\)
leżą w \(F_{\alpha|k}\),
więc ich odległość
nie przekracza \(2^{-k}\).
Stąd \(f\) jest ciągłe.

Dla dowolnego \(x\in X\)
wybieramy rekurencyjnie
\(k_1,k_2,\ldots\), aby
\(x\in F_{(k_1,\ldots,k_n)}\)
dla wszystkich \(n\).
Wybór jest możliwy, bo
każdy \(F_s\) jest sumą
swoich następców.
Dla \(\alpha=(k_n)\)
otrzymujemy \(f(\alpha)=x\).
Zatem \(f\) jest surjektywne.
\end{proof}

\begin{uwaga}[Co daje parametryzacja]
\label{rem:topog20-param}
W powyższym lemacie \(X\)
nie musi być zwartą ani
lokalnie zwartą przestrzenią.
Nie twierdzimy również,
że dowolny ciągły obraz
przestrzeni polskiej jest polski.
Właśnie to rozróżnienie
prowadzi do zbiorów analitycznych.
\end{uwaga}

\section{Zbiory analityczne i rzuty zbiorów domkniętych}
\label{sec:topog20-analityczne}

\begin{definicja}[Zbiór analityczny]
\label{def:topog20-analityczny}
Niech \(X\) będzie przestrzenią polską.
Podzbiór \(A\subseteq X\) nazywamy
\emph{analitycznym}, jeśli
\(A=\varnothing\) albo istnieje
ciągłe odwzorowanie
\(f\colon\mathcal N\to X\)
takie, że \(A=f(\mathcal N)\).
Dopełnienie zbioru analitycznego
nazywamy \emph{koanalitycznym}.
\end{definicja}

\begin{stwierdzenie}[Inna przestrzeń parametryczna]
\label{prop:topog20-param}
Dla przestrzeni polskiej \(X\)
i niepustego \(A\subseteq X\)
następujące warunki są równoważne:
\begin{enumerate}[label=\textup{(\roman*)}]
\item \(A\) jest analityczny;
\item istnieje niepusta przestrzeń
polska \(Y\) i ciągłe odwzorowanie
\(g\colon Y\to X\) takie, że \(g(Y)=A\).
\end{enumerate}
\end{stwierdzenie}
\begin{proof}
Implikację \textup{(i)} \(\Rightarrow\)
\textup{(ii)} uzyskujemy,
biorąc \(Y=\mathcal N\).
Odwrotnie,
z Lematu~\ref{lem:topog20-surjection}
istnieje ciągła surjekcja
\(h\colon\mathcal N\to Y\).
Wówczas \(g\circ h\)
jest ciągłe i ma obraz \(A\).
\end{proof}

\begin{twierdzenie}[Charakterystyka przez rzut]
\label{thm:topog20-rzut}
Niech \(X\) będzie przestrzenią polską,
a \(A\subseteq X\).
Wtedy \(A\) jest analityczny
dokładnie wtedy, gdy istnieje
zbiór domknięty
\(F\subseteq X\times\mathcal N\)
taki, że
\[
 A=\pi_X(F)
   =\{x\in X:
       \text{ istnieje }u\in\mathcal N
       \text{ z }(x,u)\in F\}.
\]
\end{twierdzenie}
\begin{proof}
Jeżeli \(A=\varnothing\),
wystarczy przyjąć \(F=\varnothing\).
W przeciwnym przypadku,
jeśli \(A=f(\mathcal N)\)
dla ciągłego \(f\),
rozważamy jego graf
\[
 F=\{(f(u),u):u\in\mathcal N\}.
\]
Graf jest domknięty:
jeśli \((x_j,u_j)\to(x,u)\)
i \(x_j=f(u_j)\),
to z ciągłości \(f\)
wynika \(x=f(u)\).
Jego rzutem na \(X\)
jest dokładnie \(A\).

W drugą stronę załóżmy,
że \(A=\pi_X(F)\)
dla domkniętego \(F\).
Produkt \(X\times\mathcal N\)
jest przestrzenią polską na mocy
Twierdzenia~\ref{thm:topog-polskie-iloczyn}
(uzupełniając iloczyn dwoma czynnikami
o jednoelementowe przestrzenie polskie).
Jego domknięta podprzestrzeń
\(F\) jest polska.
Jeśli \(F=\varnothing\), to
\(A=\varnothing\).
W przeciwnym przypadku
Lemat~\ref{lem:topog20-surjection}
daje ciągłą surjekcję
\(h\colon\mathcal N\to F\).
Złożenie \( \pi_X|_F\circ h \)
jest ciągłe i ma obraz \(A\).
\end{proof}

\begin{przyklad}[Dlaczego rzut nie musi być domknięty]
\label{ex:topog20-rzut}
W \(\mathbb R^2\) zbiór
\[
 F=\{(x,y)\in\mathbb R^2:xy=1\}
\]
jest domknięty, bo jest
przeciwobrazem singletonu
\(\{1\}\) przez ciągłe
mnożenie. Tymczasem
\(\pi_1(F)=\mathbb R\setminus\{0\}\)
nie jest domknięty.
Przykład nie daje zbioru
nieborelowskiego: pokazuje tylko,
dlaczego własności domkniętości
nie wolno przenosić przez
dowolne rzutowanie.
\end{przyklad}

\section{Działania na zbiorach analitycznych}
\label{sec:topog20-operacje}

\begin{twierdzenie}[Przeliczalne sumy i przecięcia]
\label{thm:topog20-closure}
Niech \(X\) będzie przestrzenią polską
i niech \((A_n)_{n\ge1}\)
będzie ciągiem jej podzbiorów
analitycznych.
Wówczas zbiory
\(\bigcup_{n\ge1}A_n\) i
\(\bigcap_{n\ge1}A_n\)
są analityczne.
\end{twierdzenie}
\begin{proof}
Z Twierdzenia~\ref{thm:topog20-rzut}
dla każdego \(n\) możemy wybrać
domknięty \(F_n\subseteq
X\times\mathcal N\) taki,
że \(\pi_X(F_n)=A_n\).

Najpierw rozważmy przestrzeń
\(\mathbb N\times\mathcal N\),
gdzie \(\mathbb N\) jest dyskretna.
Jest ona homeomorficzna
z \(\mathcal N\), ponieważ
\((n,u_1,u_2,\ldots)\)
można utożsamić z
\((n,u_1,u_2,\ldots)\in\mathcal N\).
Zbiór
\[
 F=\{(x,(n,u)):
               (x,u)\in F_n\}
 \subseteq X\times(\mathbb N\times\mathcal N)
\]
jest domknięty: jego przekrój
nad każdym otwarto-domkniętym
\(\{n\}\) jest domknięty,
a dopełnienie \(F\) jest
sumą otwartych dopełnień
tych przekrojów.
Jego rzut na \(X\) jest
równy \(\bigcup_n A_n\).
Twierdzenie~\ref{thm:topog20-rzut}
daje analityczność sumy.

Dla przecięcia użyjmy
\(\mathcal N^{\mathbb N}\),
która jest homeomorficzna
z \(\mathcal N\): wystarczy
ustalić bijekcję
\(\mathbb N^2\to\mathbb N\)
i ponumerować wszystkie
współrzędne obu ciągów indeksów.
Definiujemy
\[
 H=\left\{(x,(u_n)_{n\ge1}):
            (x,u_n)\in F_n
            \text{ dla każdego }n\right\}.
\]
Jest to przecięcie
przeciwobrazów domkniętych \(F_n\)
przez ciągłe odwzorowania
\((x,(u_j))\mapsto(x,u_n)\);
jest zatem domknięty.
Jego rzutem jest
\(\bigcap_n A_n\):
w jednym kierunku wynika to
bezpośrednio z definicji,
w drugim dla \(x\in\bigcap_n A_n\)
wybieramy świadków \(u_n\)
z \((x,u_n)\in F_n\).
Ponownie stosujemy
Twierdzenie~\ref{thm:topog20-rzut}.
\end{proof}

\begin{stwierdzenie}[Obrazy i przeciwobrazy]
\label{prop:topog20-maps}
Niech \(X,Y\) będą przestrzeniami
polskimi i \(f\colon X\to Y\)
odwzorowaniem ciągłym.
Wówczas:
\begin{enumerate}[label=\textup{(\roman*)}]
\item jeżeli \(A\subseteq X\)
jest analityczny, to \(f(A)\)
jest analityczny w \(Y\);
\item jeżeli \(B\subseteq Y\)
jest analityczny, to
\(f^{-1}(B)\) jest analityczny w \(X\).
\end{enumerate}
\end{stwierdzenie}
\begin{proof}
Jeśli \(A=\varnothing\),
pierwsza teza jest oczywista.
Dla \(A=g(\mathcal N)\)
mamy \(f(A)=(f\circ g)(\mathcal N)\).

Dla drugiej tezy przedstawmy
\(B=\pi_Y(F)\), gdzie
\(F\subseteq Y\times\mathcal N\)
jest domknięty.
Odwzorowanie
\[
 \Phi\colon X\times\mathcal N
     \longrightarrow Y\times\mathcal N,
 \qquad \Phi(x,u)=(f(x),u)
\]
jest ciągłe.
Zbiór \(\Phi^{-1}(F)\)
jest więc domknięty,
a jego rzut na \(X\)
jest równy \(f^{-1}(B)\).
\end{proof}

\begin{wniosek}[Zbiory borelowskie są analityczne]
\label{cor:topog20-borel}
W każdej przestrzeni polskiej \(X\)
każdy zbiór borelowski \(B\subseteq X\)
jest zarazem analityczny i koanalityczny.
\end{wniosek}
\begin{proof}
Niech \(\mathcal D\)
oznacza rodzinę wszystkich
zbiorów \(B\subseteq X\),
dla których zarówno \(B\),
jak i \(X\setminus B\)
są analityczne.
Na mocy definicji
\(\mathcal D\) jest zamknięta
na dopełnienia.
Twierdzenie~\ref{thm:topog20-closure}
oraz prawa de Morgana
pokazują, że jest zamknięta
na przeliczalne sumy;
jest więc \(\sigma\)-algebrą.

Każdy zbiór otwarty \(U\)
należy do \(\mathcal D\):
zarówno \(U\), jak i
jego dopełnienie są borelowskie
na poziomie pierwszym,
ale analityczność uzasadniamy
bez odwołania do tezy.
W przestrzeni metrycznej
\(U\) jest typu \(F_\sigma\)
na mocy Lematu~
\ref{lem:topog-polskie-otwarte-fsigma}.
Każdy zbiór domknięty,
jako podprzestrzeń polska,
jest analityczny z
Lematu~\ref{lem:topog20-surjection}
(lub jest pusty).
Stąd \(U\), jako przeliczalna
suma takich zbiorów, jest
analityczny. Jego dopełnienie
jest domknięte, więc również
analityczne.
Ponieważ \(\mathcal D\)
zawiera wszystkie otwarte
i jest \(\sigma\)-algebrą,
zawiera całą
\(\mathcal B(X)\).
\end{proof}


\begin{stwierdzenie}[Zgodność z definicją przez rzut borelowski]
\label{prop:topog20-borel-projection}
Niech \(n\ge1\), a \(A\subseteq\mathbb R^n\).
Zbiór \(A\) jest analityczny
w sensie Definicji~\ref{def:topog20-analityczny}
dokładnie wtedy, gdy
istnieje zbiór borelowski
\(B\subseteq\mathbb R^{n+1}\)
taki, że \(A=\pi_{\mathbb R^n}(B)\).
Jest to definicja stosowana
w rozdziale~30 skryptu
\emph{Logika matematyczna i teoria mnogości}.
\end{stwierdzenie}
\begin{proof}
Jeżeli \(B\) jest borelowski
w \(\mathbb R^{n+1}\),
to jest analityczny na mocy
Wniosku~\ref{cor:topog20-borel}.
Rzutowanie
\(\mathbb R^{n+1}\to\mathbb R^n\)
jest ciągłe, zatem
Stwierdzenie~\ref{prop:topog20-maps}
pokazuje, że jego obraz \(A\)
jest analityczny.

Dla implikacji odwrotnej
skonstruujemy homeomorficzne
zanurzenie przestrzeni Baire'a
w borelowski podzbiór prostej.
Niech \(\mathbb N_0=\{0,1,2,\ldots\}\)
i ustalmy bijekcję
\(b\colon\mathbb N\to\mathbb N_0\).
Dla \(\alpha=(\alpha_k)\in\mathcal N\)
zapiszmy nieskończony ciąg binarny
\[
 e(\alpha)=
 0^{b(\alpha_1)}1\,
 0^{b(\alpha_2)}1\,
 0^{b(\alpha_3)}1\cdots,
\]
gdzie \(0^l\) oznacza blok
złożony z \(l\) zer,
a \(0^0\) jest blokiem pustym.
Obrazem \(e\) jest zbiór
\[
 C_\infty=\{u\in2^{\mathbb N}:
             u_k=1\text{ dla
             nieskończenie wielu }k\}.
\]
Jest on typu \(G_\delta\)
w przestrzeni Cantora, ponieważ
\[
 C_\infty=
 \bigcap_{m\ge1}\bigcup_{k\ge m}
                  \{u:u_k=1\}.
\]
Odwzorowanie \(e\) jest bijekcją
na \(C_\infty\).
Jest ciągłe, bo początkowy
fragment jego wartości
zależy tylko od skończenie
wielu współrzędnych \(\alpha\).
Jego odwrotność jest ciągła:
pierwszych \(l\) wyrazów
\(\alpha\) odczytujemy z długości
bloków zer przed pierwszymi
\(l\) jedynkami, a ich wystąpienie
jest rozstrzygnięte przez
skończony fragment ciągu binarnego.

Odwzorowanie
\[
 j\colon2^{\mathbb N}\longrightarrow
           \mathbb R,\qquad
 j(u)=\sum_{k\ge1}\frac{2u_k}{3^k}
\]
jest ciągłe i różnowartościowe:
jeśli dwa ciągi różnią się
po raz pierwszy na pozycji \(k\),
różnica wkładu tego wyrazu
ma moduł \(2/3^k\),
a cały ogon ma moduł
co najwyżej \(1/3^k\).
Ponieważ \(2^{\mathbb N}\)
jest zwarta, \(j\) jest
homeomorfizmem na swój obraz,
który jest domknięty w
\(\mathbb R\).
Zatem \(E=j(C_\infty)\)
jest borelowski w \(\mathbb R\),
a \(i=j\circ e\)
jest homeomorfizmem
\(\mathcal N\) na \(E\).

Niech teraz \(A\) będzie analityczny.
Z Twierdzenia~\ref{thm:topog20-rzut}
mamy domknięty
\(F\subseteq\mathbb R^n\times\mathcal N\)
z \(\pi_{\mathbb R^n}(F)=A\).
Obraz \(F\) przez
\(\operatorname{id}_{\mathbb R^n}\times i\)
jest domknięty w przestrzeni
względnej \(\mathbb R^n\times E\).
Ponieważ \(E\) jest borelowski,
istnieje domknięty
\(H\subseteq\mathbb R^{n+1}\),
dla którego ten obraz
jest równy
\(H\cap(\mathbb R^n\times E)\).
Jest on zatem borelowski
w \(\mathbb R^{n+1}\)
i ma rzut równy \(A\).
\end{proof}

\section{Operacja Suslina}
\label{sec:topog20-suslin}

\begin{definicja}[Operacja Suslina]
\label{def:topog20-suslin}
Niech \(X\) będzie zbiorem,
a \((F_s)_{s\in\mathbb N^{<\mathbb N}}\)
rodziną jego podzbiorów.
\emph{Operacją Suslina}
na tej rodzinie nazywamy zbiór
\[
 \mathcal A(F_s)=
 \bigcup_{\alpha\in\mathcal N}
     \bigcap_{n\ge1}F_{\alpha|n}.
\]
W definicji nie wymagamy
monotoniczności rodziny \(F_s\).
\end{definicja}

\begin{twierdzenie}[Charakterystyka Suslina]
\label{thm:topog20-suslin}
Niech \(X\) będzie przestrzenią polską,
a \(A\subseteq X\).
Następujące warunki są równoważne:
\begin{enumerate}[label=\textup{(\roman*)}]
\item \(A\) jest zbiorem analitycznym;
\item istnieje rodzina domkniętych
zbiorów \(F_s\subseteq X\),
indeksowana wszystkimi
\(s\in\mathbb N^{<\mathbb N}\),
taka, że \(A=\mathcal A(F_s)\).
\end{enumerate}
\end{twierdzenie}
\begin{proof}
Przypuśćmy najpierw, że
\(A=\mathcal A(F_s)\)
dla domkniętych \(F_s\).
Określmy
\[
 H=\{(x,\alpha)\in X\times\mathcal N:
       x\in F_{\alpha|n}
       \text{ dla każdego }n\ge1\}.
\]
Dla ustalonego \(n\) zbiór
\[
 H_n=\bigcup_{s\in\mathbb N^n}
                         (F_s\times[s])
\]
jest domknięty.
Istotnie, cylindry \([s]\)
o długości \(n\) tworzą
otwarto-domknięty podział
\(\mathcal N\), a jego dopełnienie
jest sumą otwartych zbiorów
\((X\setminus F_s)\times[s]\).
Mamy \(H=\bigcap_n H_n\),
zatem \(H\) jest domknięty.
Jego rzut na \(X\)
jest równy \(\mathcal A(F_s)\),
więc z
Twierdzenia~\ref{thm:topog20-rzut}
wynika analityczność \(A\).

Odwrotnie, niech
\(A=\pi_X(H)\) dla domkniętego
\(H\subseteq X\times\mathcal N\).
Wybierzmy zgodną metrykę \(d\)
na \(X\) i dla skończonego \(s\)
połóżmy
\[
 F_s=
 \overline{\pi_X\bigl(
     H\cap(X\times[s])\bigr)}^{\,X}.
\]
Są to zbiory domknięte.
Jeśli \((x,\alpha)\in H\),
to \(x\in F_{\alpha|n}\)
dla każdego \(n\),
więc \(A\subseteq\mathcal A(F_s)\).

Niech teraz
\(x\in\mathcal A(F_s)\).
Istnieje \(\alpha\in\mathcal N\)
takie, że
\(x\in F_{\alpha|n}\)
dla każdego \(n\).
Z definicji domknięcia
możemy dla każdego \(n\)
wybrać \((x_n,u_n)\in H\)
tak, aby
\[
 d(x_n,x)<1/n,\qquad
 u_n\in[\alpha|n].
\]
Zatem \(x_n\to x\),
a \(u_n\to\alpha\)
w topologii produktowej.
Ponieważ \(H\) jest domknięty
w produkcie przestrzeni metryzowalnych,
otrzymujemy \((x,\alpha)\in H\).
W szczególności
\(x\in\pi_X(H)=A\).
Dowodzi to drugiego zawierania.
\end{proof}

\begin{uwaga}[Granica ustalonego zakresu]
\label{rem:topog20-no-meas}
Powyższe twierdzenie daje
topologiczny opis zbiorów analitycznych.
W szczególności \emph{nie} stanowi
jeszcze dowodu ich uniwersalnej
mierzalności ani twierdzenia
o rozdzielaniu zbiorów analitycznych.
W skrypcie \emph{Logika matematyczna
i teoria mnogości}, rozdział~30,
twierdzenie o mierzalności
zbiorów analitycznych jest
jawnie przywołane jako zewnętrzny wynik.
Nie przedstawiamy go tutaj
jako udowodnionego.
\end{uwaga}

\section{Zbiór analityczny nieborelowski}
\label{sec:topog20-nieborelowski}

\begin{lemat}[Uniwersalna rodzina rzutów domkniętych]
\label{lem:topog20-universal}
Istnieje analityczny
\(U\subseteq\mathcal N\times\mathcal N\)
o następującej własności:
dla każdego analitycznego
\(A\subseteq\mathcal N\)
istnieje \(c\in\mathcal N\)
takie, że
\[
 A=U_c=\{x\in\mathcal N:(c,x)\in U\}.
\]
\end{lemat}
\begin{proof}
Zbiór par \((s,t)\) skończonych
ciągów naturalnych tej samej
długości jest przeliczalny.
Ustalmy jego numerację
\(e(s,t)\in\mathbb N\), będącą bijekcją.
Dla \(c\in\mathcal N\) kodujemy
rodzinę par przez warunek
\(c_{e(s,t)}=1\).
Wartość \(1\) oznacza obecność pary,
a wartość \(2\) jej brak.
Obie wartości należą do \(\mathbb N\),
niezależnie od tego, czy \(0\)
zaliczamy do liczb naturalnych.

Zdefiniujmy
\[
 H=\{(c,x,y)\in\mathcal N^3:
       c_{e(x|n,y|n)}=1
       \text{ dla każdego }n\ge1\}.
\]
Jest to zbiór domknięty.
Dla ustalonego \(n\)
warunek na pierwszych
\(n\) współrzędnych \(x,y\)
i jedną współrzędną \(c\)
wyznacza zbiór otwarty
i domknięty; przecięcie
tych zbiorów po \(n\)
jest domknięte.
Niech
\[
 U=\{(c,x):\text{ istnieje }y
                   \text{ z }(c,x,y)\in H\}.
\]
Z Twierdzenia~\ref{thm:topog20-rzut}
\(U\) jest analityczny,
po utożsamieniu
\(\mathcal N^2\cong\mathcal N\).

Weźmy teraz dowolny analityczny
\(A\subseteq\mathcal N\).
Istnieje domknięty
\(F\subseteq\mathcal N^2\)
z \(\pi_1(F)=A\).
Rozważmy rodzinę
\[
 T_F=\{(s,t):
    |s|=|t|,\quad
    F\cap([s]\times[t])
           \ne\varnothing\}.
\]
Dobierzmy kod \(c\) tak,
że \(c_{e(s,t)}=1\) dla par
\((s,t)\in T_F\), natomiast
\(c_{e(s,t)}=2\) dla
wszystkich pozostałych par.

Dla \((x,y)\in F\)
wszystkie prefiksy
\((x|n,y|n)\) należą do \(T_F\).
Odwrotnie, jeśli każdy
taki prefiks należy do \(T_F\),
dla każdego \(n\) wybieramy
\((x_n,y_n)\in
F\cap([x|n]\times[y|n])\).
Wtedy \((x_n,y_n)\to(x,y)\),
a domkniętość \(F\) daje
\((x,y)\in F\).
Zatem
\[
 (c,x,y)\in H
   \ \Longleftrightarrow\ (x,y)\in F,
\]
co po rzutowaniu na \(x\)
daje \(U_c=A\).
\end{proof}

\begin{twierdzenie}[Istnienie nieborelowskiego zbioru analitycznego]
\label{thm:topog20-nonborel}
W przestrzeni Baire'a
\(\mathcal N=\mathbb N^{\mathbb N}\)
istnieje zbiór analityczny,
który nie jest borelowski.
\end{twierdzenie}
\begin{proof}
Niech \(U\) będzie zbiorem
z Lematu~\ref{lem:topog20-universal}.
Określmy
\[
 D=\{x\in\mathcal N:
                 (x,x)\notin U\}.
\]
Przypuśćmy, że \(D\)
jest analityczny.
Z uniwersalności \(U\)
istnieje \(c\in\mathcal N\)
takie, że \(D=U_c\).
Wtedy
\[
 c\in D
 \ \Longleftrightarrow\ (c,c)\notin U
 \ \Longleftrightarrow\ c\notin U_c=D,
\]
co jest sprzecznością.
Zatem \(D\) nie jest analityczny.

Jego dopełnienie
\[
 E=\{x\in\mathcal N:(x,x)\in U\}
\]
jest analityczne:
to przeciwobraz analitycznego
\(U\) przez ciągłe
odwzorowanie diagonalne
\(x\mapsto(x,x)\), zgodnie ze
Stwierdzeniem~\ref{prop:topog20-maps}.
Gdyby \(E\) był borelowski,
jego dopełnienie \(D\)
też byłoby borelowskie,
a więc analityczne na mocy
Wniosku~\ref{cor:topog20-borel}.
Otrzymujemy sprzeczność.
Wobec tego \(E\) jest
analityczny i nieborelowski.
\end{proof}

\begin{wniosek}[Rzut może zwiększać złożoność]
\label{cor:topog20-projection}
Istnieje zbiór domknięty
\(F\subseteq\mathcal N^2\),
którego rzut na pierwszą
współrzędną nie jest borelowski.
\end{wniosek}
\begin{proof}
Zbiór \(E\) otrzymany w
Twierdzeniu~\ref{thm:topog20-nonborel}
jest analityczny i nieborelowski.
Z Twierdzenia~\ref{thm:topog20-rzut}
jest rzutem zbioru domkniętego
\(F\subseteq\mathcal N^2\).
\end{proof}

\section{Podsumowanie i granice materiału}
\label{sec:topog20-summary}

Niech \(X\) będzie przestrzenią polską.
Podzbiór \(X\) jest analityczny
dokładnie wtedy, gdy jest pusty
albo jest ciągłym obrazem
przestrzeni Baire'a.
Równoważnie, jest rzutem
domkniętego podzbioru
\(X\times\mathcal N\) lub wynikiem
operacji Suslina na rodzinie
zbiorów domkniętych.
Są zamknięte na przeliczalne
sumy, przeliczalne przecięcia,
ciągłe obrazy i przeciwobrazy,
a każdy zbiór borelowski
jest analityczny i koanalityczny.
W przestrzeni Baire'a
klasa zbiorów analitycznych
jest ściśle większa od klasy
zbiorów borelowskich, zgodnie z
Twierdzeniem~\ref{thm:topog20-nonborel}.
Nie zachodzi to dla każdej
przestrzeni polskiej:
w przeliczalnej przestrzeni
dyskretnej każdy podzbiór
jest otwarty, a zatem borelowski.

Nie powtarzaliśmy teorii
determinacji i redukcji Wadge'a
z rozdziału~30 skryptu
\emph{Logika matematyczna i teoria mnogości}.
Wyniki dotyczące
uniwersalnej mierzalności,
własności zbioru doskonałego
i rozdzielania zbiorów analitycznych
pozostają osobnymi,
głębszymi zagadnieniami.

\chapter{Topologie na rodzinach podzbiorów}
\label{ch:topog-hiperprzestrzenie}

\noindent\textbf{Cel rozdziału.}
W rozdziale~\ref{ch:topog-przestrzenie-odwzorowan} rozważaliśmy
topologie na zbiorach odwzorowań. Teraz punktami nowej przestrzeni
będą same podzbiory ustalonej przestrzeni \(X\).
Pozwala to badać rodziny zbiorów zwartych, zmieniające się
figury geometryczne i granice zbiorów domkniętych.
Rozróżnimy topologię Vietorisa na zbiorach zwartych
oraz topologię Fella na zbiorach domkniętych.
W przypadku metrycznym wykażemy związek z metryką Hausdorffa,
udowodnimy zupełność odpowiedniej przestrzeni zbiorów zwartych
i wskażemy, kiedy dwie topologie opisują różne rodzaje zbieżności.

\section{Hiperprzestrzeń zbiorów zwartych i topologia Vietorisa}
\label{sec:topog21-vietoris}

\begin{definicja}[Hiperprzestrzeń zwartych podzbiorów]
\label{def:topog21-kx}
Dla przestrzeni topologicznej \(X\) definiujemy
\(\mathcal K(X)=\{K\subseteq X:K\neq\varnothing,\ K
\text{ jest zwarty}\}\).
Dla zbioru otwartego \(U\subseteq X\) wprowadzamy
\[
 U^-=\{K\in\mathcal K(X):K\cap U\neq\varnothing\},
 \qquad
 U^+=\{K\in\mathcal K(X):K\subseteq U\}.
\]
\emph{Topologią Vietorisa} na \(\mathcal K(X)\)
nazywamy topologię generowaną przez wszystkie zbiory
\(U^-\) i \(U^+\), gdzie \(U\) przebiega zbiory otwarte w \(X\).
Zbiory pierwszego rodzaju określają warunki
\emph{trafienia} zbioru otwartego, a drugiego
warunki \emph{zawierania się} w zbiorze otwartym.
\end{definicja}

\begin{uwaga}[Rola zbioru pustego]
\label{rem:topog21-empty}
Zbiór pusty został wyłączony z \(\mathcal K(X)\).
Dzięki temu zarówno odległość Hausdorffa,
jak i warunek trafienia mają zwykłą postać,
bez wartości nieskończonych i dodatkowych wyjątków.
W topologii Fella będziemy natomiast dopuszczać zbiór pusty.
Jeśli \(X=\varnothing\), to \(\mathcal K(X)=\varnothing\).
\end{uwaga}

\begin{stwierdzenie}[Baza topologii Vietorisa]
\label{prop:topog21-vietoris-basis}
Niech \(X\) będzie przestrzenią topologiczną.
Dla każdego \(m\ge1\) i otwartych \(U_1,\ldots,U_m\subseteq X\)
połóżmy
\[
 \langle U_1,\ldots,U_m\rangle
 =\{K\in\mathcal K(X):
 K\subseteq\textstyle\bigcup_{i=1}^{m}U_i,\
 K\cap U_i\neq\varnothing\text{ dla }1\le i\le m\}.
\]
Rodzina wszystkich takich zbiorów jest bazą
topologii Vietorisa.
\end{stwierdzenie}
\begin{proof}
Każdy rozważany zbiór ma postać
\((\bigcup_iU_i)^+\cap\bigcap_iU_i^-\),
a więc jest otwarty.
Aby otrzymać bazę, wystarczy dla dowolnego
\(K\in\mathcal K(X)\) oraz dowolnego
skończonego przecięcia zbiorów podbazowych,
które zawiera \(K\), znaleźć zbiór przedstawionej
postaci zawierający \(K\) i zawarty w tym przecięciu.

Niech warunki górne mają postać
\(K\subseteq V_1,\ldots,K\subseteq V_r\),
a dolne \(K\cap W_j\neq\varnothing\) dla \(1\le j\le s\).
Przyjmijmy \(V=\bigcap_{i=1}^{r}V_i\),
a dla \(r=0\) przyjmijmy \(V=X\).
Wówczas
\(\langle V,V\cap W_1,\ldots,V\cap W_s\rangle\)
zawiera \(K\): warunki trafienia \(V\cap W_j\)
wynikają z \(K\subseteq V\) oraz trafienia \(W_j\).
Każdy element tego zbioru jest zawarty w \(V\)
i trafia wszystkie \(W_j\).
Dotyczy to także \(s=0\), gdy bazowym zbiorem
jest po prostu \(\langle V\rangle=V^+\).
\end{proof}

\begin{stwierdzenie}[Jednopunktowe zbiory jako podprzestrzeń]
\label{prop:topog21-singletons}
Niech \(X\) będzie przestrzenią topologiczną.
Odwzorowanie
\(j\colon X\to\mathcal K(X)\), \(j(x)=\{x\}\),
jest zanurzeniem topologicznym.
Jeżeli ponadto \(X\) jest przestrzenią Hausdorffa,
to \(j(X)\) jest domknięty w \(\mathcal K(X)\).
\end{stwierdzenie}
\begin{proof}
Dla każdego otwartego \(U\subseteq X\) mamy
\(j^{-1}(U^-)=j^{-1}(U^+)=U\).
Stąd \(j\) jest ciągłe.
Jest różnowartościowe, a odwzorowanie odwrotne
na jego obraz jest ciągłe, ponieważ
\(j(U)=j(X)\cap U^-\) jest otwarte w \(j(X)\).
Jeżeli \(X\) jest Hausdorffa, a \(K\in\mathcal K(X)\)
ma przynajmniej dwa punkty \(x\neq y\), to istnieją
rozłączne zbiory otwarte \(U\ni x\), \(V\ni y\).
Zbiór \(U^-\cap V^-\) jest otwartym otoczeniem \(K\),
niezawierającym żadnego singletonu.
Zatem dopełnienie \(j(X)\) jest otwarte.
\end{proof}

\begin{twierdzenie}[Oddzielanie zbiorów zwartych]
\label{thm:topog21-vietoris-separation}
Jeżeli \(X\) jest przestrzenią Hausdorffa,
to \(\mathcal K(X)\) z topologią Vietorisa
również jest przestrzenią Hausdorffa.
\end{twierdzenie}
\begin{proof}
Niech \(A,B\in\mathcal K(X)\) i \(A\neq B\).
Po ewentualnej zamianie ich ról wybierzmy \(x\in A\setminus B\).
Dla każdego \(b\in B\) istnieją rozłączne zbiory otwarte
\(U_b\ni x\) i \(V_b\ni b\).
Zwartość \(B\) pozwala wybrać \(b_1,\ldots,b_m\)
tak, że \(B\subseteq\bigcup_{i=1}^m V_{b_i}\).
Wówczas \(U=\bigcap_{i=1}^m U_{b_i}\) jest otwartym
otoczeniem \(x\), a \(V=\bigcup_{i=1}^mV_{b_i}\)
jest otwartym otoczeniem \(B\), przy czym
\(U\cap V=\varnothing\).
Zbiory \(U^-\) i \(V^+\) są rozłącznymi otwartymi
otoczeniami odpowiednio \(A\) i \(B\)
w \(\mathcal K(X)\).
\end{proof}

\begin{stwierdzenie}[Obrazy zwartych zbiorów]
\label{prop:topog21-induced}
Niech \(X,Y\) będą przestrzeniami topologicznymi,
a \(f\colon X\to Y\) odwzorowaniem ciągłym.
Wówczas odwzorowanie
\[
 \mathcal K(f)\colon\mathcal K(X)\to\mathcal K(Y),
 \qquad K\longmapsto f(K),
\]
jest dobrze określone i ciągłe dla topologii Vietorisa.
\end{stwierdzenie}
\begin{proof}
Obraz niepustego zbioru zwartego przez odwzorowanie
ciągłe jest niepusty i zwarty.
Dla otwartego \(U\subseteq Y\) zachodzą równości
\[
 \mathcal K(f)^{-1}(U^-)=\bigl(f^{-1}(U)\bigr)^-,
 \qquad
 \mathcal K(f)^{-1}(U^+)=\bigl(f^{-1}(U)\bigr)^+.
\]
Oba przeciwobrazy są otwarte, ponieważ \(f^{-1}(U)\)
jest otwarty w \(X\). Wystarczy sprawdzić
ciągłość na podbazie.
\end{proof}

\section{Metryka Hausdorffa i jej topologia}
\label{sec:topog21-hausdorff}

\begin{definicja}[Odległość Hausdorffa]
\label{def:topog21-hausdorff}
Niech \((X,d)\) będzie przestrzenią metryczną.
Dla \(x\in X\) i niepustego \(A\subseteq X\) piszemy
\(d(x,A)=\inf_{a\in A}d(x,a)\).
Dla \(A,B\in\mathcal K(X)\) definiujemy
\[
 d_H(A,B)=
 \max\left\{\sup_{a\in A}d(a,B),
             \sup_{b\in B}d(b,A)\right\}.
\]
Obie liczby pod supremami są skończone,
gdyż \(A\) i \(B\) są niepuste i zwarte,
a więc ograniczone.
\end{definicja}

\begin{twierdzenie}[Metryczność odległości Hausdorffa]
\label{thm:topog21-hausdorff-metric}
Dla każdej przestrzeni metrycznej \((X,d)\)
funkcja \(d_H\) jest metryką na \(\mathcal K(X)\).
\end{twierdzenie}
\begin{proof}
Nieujemność i symetria są bezpośrednie.
Jeżeli \(d_H(A,B)=0\), to dla każdego \(a\in A\)
mamy \(d(a,B)=0\). Ponieważ \(B\) jest zwarty,
jest domknięty w \(X\), stąd \(a\in B\).
Otrzymujemy \(A\subseteq B\), a analogicznie
\(B\subseteq A\), więc \(A=B\).

Niech \(A,B,C\in\mathcal K(X)\).
Dla \(a\in A\) i \(b\in B\) nierówność trójkąta
daje \(d(a,C)\le d(a,b)+d(b,C)\).
Biorąc infimum po \(b\in B\), otrzymujemy
\(d(a,C)\le d(a,B)+\sup_{b\in B}d(b,C)\).
Po wzięciu supremum po \(a\in A\) wynika
\[
 \sup_{a\in A}d(a,C)
 \le d_H(A,B)+d_H(B,C).
\]
Identycznie oszacowujemy \(\sup_{c\in C}d(c,A)\).
Maksimum obu liczb spełnia więc
\(d_H(A,C)\le d_H(A,B)+d_H(B,C)\).
\end{proof}

\begin{przyklad}[Dwa skończone podzbiory prostej]
\label{ex:topog21-points}
Dla \(A=\{0,2\}\) i
\(B=\{\frac15,\frac{17}{10}\}\)
w metryce euklidesowej mamy
\(\sup_{a\in A}d(a,B)=\frac3{10}\) oraz
\(\sup_{b\in B}d(b,A)=\frac3{10}\);
zatem \(d_H(A,B)=\frac3{10}\).
Nie jest to odległość między środkami
ani suma odległości dobranych punktów.
\begin{figure}[htbp]
\centering
\begin{tikzpicture}[x=3.05cm,y=1cm]
 \draw[-{Latex[length=2mm]},thin] (-.13,0)--(2.25,0);
 \foreach \t/\lab in {0/0,.2/{0{,}2},1.7/{1{,}7},2/2}
   {\draw (\t,.04)--(\t,-.04);
    \node[below,font=\small] at (\t,-.08) {$\lab$};}
 \fill[blue!70!black] (0,.52) circle (2.6pt);
 \fill[blue!70!black] (2,.52) circle (2.6pt);
 \fill[orange!85!black] (.2,.28) circle (2.6pt);
 \fill[orange!85!black] (1.7,.28) circle (2.6pt);
 \node[left,font=\small,blue!70!black] at (-.12,.67) {$A$};
 \node[left,font=\small,orange!85!black] at (-.12,.13) {$B$};
\end{tikzpicture}
\caption{Odległość Hausdorffa mierzy najdalszą odległość
punktu jednego zbioru od drugiego zbioru.}
\label{fig:topog21-points}
\end{figure}
\end{przyklad}

\begin{twierdzenie}[Vietoris i metryka Hausdorffa]
\label{thm:topog21-vietoris-hausdorff}
Niech \((X,d)\) będzie przestrzenią metryczną.
Topologia Vietorisa na \(\mathcal K(X)\) jest równa
topologii metryki \(d_H\).
\end{twierdzenie}
\begin{proof}
Najpierw sprawdzimy, że wszystkie zbiory podbazowe
Vietorisa są otwarte w metryce \(d_H\).
Jeżeli \(K\cap U\neq\varnothing\) i \(U\subseteq X\)
jest otwarty, wybieramy \(x\in K\cap U\)
oraz \(r>0\) z \(B_d(x,r)\subseteq U\).
Gdy \(d_H(K,L)<r/2\), istnieje \(y\in L\)
z \(d(x,y)<r/2\), więc \(L\cap U\neq\varnothing\).
Zatem \(U^-\) jest otwarty.

Jeżeli \(K\subseteq U\), to dla \(U=X\) teza jest
oczywista. W przeciwnym razie funkcja
\(x\mapsto d(x,X\setminus U)\) jest ciągła
i dodatnia na zwartym \(K\). Osiąga więc
dodatnie minimum \(r\).
Jeżeli \(d_H(K,L)<r/2\), to każdy \(y\in L\)
ma odległość mniejszą niż \(r/2\) od pewnego
\(x\in K\), zatem \(y\notin X\setminus U\).
Daje to \(L\subseteq U\) i otwartość \(U^+\).

Odwrotnie, niech \(K\in\mathcal K(X)\) i \(r>0\).
Zwartość \(K\) pozwala wybrać \(x_1,\ldots,x_m\in K\)
tak, że \(K\subseteq\bigcup_i B_d(x_i,r/3)\).
Niech \(V=\bigcup_{x\in K}B_d(x,r/2)\).
Rozważmy otoczenie Vietorisa
\[
 \mathcal O=V^+\cap
 \bigcap_{i=1}^{m} B_d(x_i,r/3)^- .
\]
Zbiór \(K\) należy do \(\mathcal O\).
Dla \(L\in\mathcal O\), każdego \(y\in L\)
mamy \(d(y,K)<r/2\), ponieważ \(L\subseteq V\).
Dla dowolnego \(x\in K\) wybieramy \(i\)
z \(d(x,x_i)<r/3\), a następnie
\(y_i\in L\cap B_d(x_i,r/3)\).
Wtedy \(d(x,L)<2r/3\).
Wobec tego \(d_H(K,L)<r\),
czyli \(\mathcal O\subseteq B_{d_H}(K,r)\).
Otrzymujemy drugie zawieranie topologii.
\end{proof}

\begin{wniosek}[Zbieżność w metryce Hausdorffa]
\label{cor:topog21-hausdorff-convergence}
Niech \((X,d)\) będzie przestrzenią metryczną
i \(K_n,K\in\mathcal K(X)\).
Wówczas \(K_n\to K\) w topologii Vietorisa
dokładnie wtedy, gdy dla każdego \(\varepsilon>0\)
od pewnego miejsca zachodzą oba zawierania
\[
 K_n\subseteq\{x:d(x,K)<\varepsilon\},
 \qquad
 K\subseteq\{x:d(x,K_n)<\varepsilon\}.
\]
\end{wniosek}
\begin{proof}
Oba zawierania są równoważne nierówności
\(d_H(K_n,K)\le\varepsilon\) z odpowiednią
zmianą \(\varepsilon\) na mniejszą liczbę dodatnią.
Zastosujmy Twierdzenie~\ref{thm:topog21-vietoris-hausdorff}
i definicję granicy ciągu w przestrzeni metrycznej.
\end{proof}

\section{Zupełność i zwartość hiperprzestrzeni}
\label{sec:topog21-complete}

\begin{twierdzenie}[Zupełność hiperprzestrzeni]
\label{thm:topog21-complete}
Jeżeli \((X,d)\) jest zupełną przestrzenią metryczną,
to \((\mathcal K(X),d_H)\) jest zupełną przestrzenią
metryczną.
\end{twierdzenie}
\begin{proof}
Niech \((K_n)\) będzie ciągiem Cauchy'ego
w \((\mathcal K(X),d_H)\).
Wybierzmy rosnący ciąg indeksów \((n_j)\)
tak, że dla każdego \(j\ge1\)
\[
 d_H(K_{n_j},K_{n_{j+1}})<2^{-j}.
\]
Taki wybór jest możliwy z warunku Cauchy'ego.
Oznaczmy \(A_j=K_{n_j}\)
i \(f_j(x)=d(x,A_j)\) dla \(x\in X\).
Nierówność
\(|d(x,A)-d(x,B)|\le d_H(A,B)\)
pokazuje, że \((f_j)\) jest jednostajnie Cauchy'ego
na \(X\). Zatem istnieje funkcja \(f\colon X\to[0,\infty)\),
do której \(f_j\) zbiega jednostajnie.
Każde \(f_j\) jest \(1\)-lipschitzowskie,
więc \(f\) jest ciągła.
Niech \(K=\{x\in X:f(x)=0\}\).
Zbiór \(K\) jest domknięty.

Pokażemy, że \(K\neq\varnothing\).
Wybieramy \(x_1\in A_1\);
rekurencyjnie można wybrać \(x_{j+1}\in A_{j+1}\)
tak, aby \(d(x_j,x_{j+1})<2^{-j}\).
Wówczas \((x_j)\) jest ciągiem Cauchy'ego,
więc z zupełności \(X\) zbiega do \(x\in X\).
Ponieważ \(d(x,A_j)\le d(x,x_j)\to0\),
a \(f_j(x)\to f(x)\), mamy \(f(x)=0\),
czyli \(x\in K\).

Dla dowolnego \(j\) i \(z\in A_j\)
można analogicznie rozpocząć łańcuch
w punkcie \(z\), otrzymując punkt
\(x_z\in K\) z
\(d(z,x_z)\le\sum_{k=j}^{\infty}2^{-k}=2^{1-j}\).
Zatem \(\sup_{z\in A_j}d(z,K)\le 2^{1-j}\).
Z drugiej strony dla \(x\in K\) mamy
\(d(x,A_j)=|f_j(x)-f(x)|
\le\sup_{y\in X}|f_j(y)-f(y)|\to0\).
Po wykazaniu zwartości \(K\)
oszacowania te pozwolą zastosować
definicję \(d_H\) i otrzymać \(d_H(A_j,K)\to0\).
Udowodnimy więc zwartość \(K\)
bez zakładania jej z góry.

Niech \(\varepsilon>0\).
Wybierzmy \(j\) tak, że
\(\sup_{x\in K}d(x,A_j)<\varepsilon/3\).
Zwarty \(A_j\) ma skończoną
\(\varepsilon/3\)-sieć \(z_1,\ldots,z_m\in A_j\).
Każdy punkt \(x\in K\) leży
w odległości mniejszej niż \(2\varepsilon/3\)
od któregoś \(z_i\).
Zatem \(K\) jest całkowicie ograniczony.
Jako domknięty podzbiór zupełnej przestrzeni
\(X\) jest również zupełny, a więc zwarty.
Dowiedzione oszacowania dają
\(d_H(A_j,K)\to0\).
Cały ciąg Cauchy'ego \((K_n)\)
zbiega wtedy do \(K\), ponieważ ma podciąg
zbieżny do \(K\).
\end{proof}

\begin{twierdzenie}[Zwartość dla zwartej przestrzeni metrycznej]
\label{thm:topog21-compact-metric}
Jeżeli \((X,d)\) jest zwartą przestrzenią metryczną,
to \((\mathcal K(X),d_H)\) jest zwarte.
\end{twierdzenie}
\begin{proof}
Przestrzeń \(X\) jest zupełna,
więc Twierdzenie~\ref{thm:topog21-complete}
zapewnia zupełność \(\mathcal K(X)\).
Wykażemy całkowitą ograniczoność tej hiperprzestrzeni.

Niech \(\varepsilon>0\).
Zwartość \(X\) daje skończoną
\(\varepsilon/3\)-sieć \(S\subseteq X\).
Dla każdego \(K\in\mathcal K(X)\) połóżmy
\(S_K=\{s\in S:d(s,K)<\varepsilon/3\}\).
Zbiór \(S_K\) jest niepusty: każdy \(x\in K\)
ma odległość mniejszą niż \(\varepsilon/3\)
od pewnego punktu \(S\).
Dla \(x\in K\) możemy wybrać \(s\in S_K\)
z \(d(x,s)<\varepsilon/3\); dla \(s\in S_K\)
z definicji \(d(s,K)<\varepsilon/3\).
Zatem \(d_H(K,S_K)\le\varepsilon/3<\varepsilon\).
Ponieważ niepustych podzbiorów skończonego \(S\)
jest skończenie wiele, stanowią one skończoną
\(\varepsilon\)-sieć w \(\mathcal K(X)\).
Zupełność i całkowita ograniczoność
metrycznej \(\mathcal K(X)\) dają jej zwartość.
\end{proof}

\begin{twierdzenie}[Zwartość Vietorisa bez metryki]
\label{thm:topog21-vietoris-compact}
Jeżeli \(X\) jest zwartą przestrzenią Hausdorffa,
to \(\mathcal K(X)\) z topologią Vietorisa
jest zwartą przestrzenią Hausdorffa.
\end{twierdzenie}
\begin{proof}
Własność Hausdorffa wynika z
Twierdzenia~\ref{thm:topog21-vietoris-separation}.
Aby dowieść zwartości, zastosujemy
charakterystykę przez ultrafiltry z
Twierdzenia~\ref{thm:topog-zwartosc-ultrafiltry}.
Niech \(\mathcal U\) będzie ultrafiltrem
na \(\mathcal K(X)\).
Określmy
\[
 K=\{x\in X:\text{ dla każdego otwartego }
  V\ni x\text{ zachodzi }V^-\in\mathcal U\}.
\]
Jeśli \(x\notin K\), istnieje otwarte \(V\ni x\)
z \(V^-\notin\mathcal U\). Każdy punkt \(V\)
również nie należy do \(K\).
Zatem \(K\) jest domknięty w \(X\).

Gdyby \(K\) był pusty, moglibyśmy dla każdego
\(x\in X\) wybrać otwarty \(V_x\ni x\)
z \(V_x^-\notin\mathcal U\).
Zwartość \(X\) pozwala wybrać skończone
podpokrycie \(V_{x_1},\ldots,V_{x_m}\).
Każdy niepusty zbiór zwarty trafia
przynajmniej jedno \(V_{x_i}\), czyli
\(\mathcal K(X)=\bigcup_i V_{x_i}^-\).
Ultrafiltr musiałby zawierać
co najmniej jeden \(V_{x_i}^-\),
co daje sprzeczność.
Stąd \(K\neq\varnothing\), a jako zbiór
domknięty w zwartej \(X\) jest zwarty.

Sprawdzimy, że \(\mathcal U\to K\) w topologii Vietorisa.
Jeśli \(K\cap V\neq\varnothing\) dla otwartego \(V\),
to z definicji \(K\) mamy \(V^-\in\mathcal U\).
Jeśli \(K\subseteq V\), zbiór \(C=X\setminus V\)
jest zwarty. Dla każdego \(x\in C\) wybieramy
otwarty \(W_x\ni x\) z \(W_x^-\notin\mathcal U\).
Z pokrycia \(C\) wybieramy skończenie wiele
\(W_{x_1},\ldots,W_{x_m}\).
Dopełnienie \(V^+\) w \(\mathcal K(X)\)
jest zawarte w \(\bigcup_iW_{x_i}^-\).
Nie może należeć do ultrafiltru, gdyż wtedy
jedno z \(W_{x_i}^-\) również do niego należałoby.
Zatem \(V^+\in\mathcal U\).
Ultrafiltr zawiera każde podbazowe otoczenie
\(K\), a więc każde otoczenie \(K\).
Jest zbieżny do \(K\).
\end{proof}

\section{Topologia Fella na zbiorach domkniętych}
\label{sec:topog21-fell}

\begin{definicja}[Hiperprzestrzeń Fella]
\label{def:topog21-fell}
Niech \(X\) będzie przestrzenią topologiczną.
Przez \(\mathcal F(X)\) oznaczamy rodzinę
\emph{wszystkich} domkniętych podzbiorów \(X\),
włącznie ze zbiorem pustym.
Dla otwartego \(U\subseteq X\) i zwartego
\(C\subseteq X\) definiujemy
\[
 U^-_{\!F}=\{F\in\mathcal F(X):F\cap U\neq\varnothing\},
 \qquad
 C^0=\{F\in\mathcal F(X):F\cap C=\varnothing\}.
\]
\emph{Topologia Fella} jest topologią generowaną
przez wszystkie takie zbiory.
Pierwszy rodzaj oznacza trafienie otwartego zbioru,
drugi zaś ominięcie zbioru zwartego.
\end{definicja}

\begin{uwaga}[Konwencja dotycząca pustego zbioru]
\label{rem:topog21-fell-empty}
Niektóre źródła rozważają wyłącznie niepuste
zbiory domknięte. W tym rozdziale
\(\mathcal F(X)\) zawiera \(\varnothing\);
to założenie jest istotne dla
twierdzenia o zwartości poniżej.
Przy \(X\) niezwartym ciąg niepustych
zbiorów może w topologii Fella zbiegać
do \(\varnothing\).
\end{uwaga}

\begin{lemat}[Kryterium zwartości przez podbazę]
\label{lem:topog21-alexander}
Niech \(Y\) będzie przestrzenią topologiczną,
a \(\mathcal S\) podbazą jej topologii.
Jeżeli każde pokrycie \(Y\) zbiorami z \(\mathcal S\)
ma skończone podpokrycie, to \(Y\) jest zwarta.
\end{lemat}
\begin{proof}
Niech \(\mathcal U\) będzie dowolnym ultrafiltrem na \(Y\).
Gdyby nie był zbieżny do żadnego punktu,
dla każdego \(y\in Y\) istniałoby otwarte
otoczenie \(O_y\ni y\), które nie należy do \(\mathcal U\).
Ponieważ skończone przecięcia elementów podbazy
tworzą bazę, możemy wybrać
\(y\in S_{y,1}\cap\cdots\cap S_{y,m_y}\subseteq O_y\),
gdzie \(S_{y,i}\in\mathcal S\).
Co najmniej jeden z elementów \(S_{y,i}\)
nie należy do ultrafiltru: gdyby wszystkie
należały, zawierałby też ich przecięcie
i jego nadzbiór \(O_y\).
Wybierzmy taki element i oznaczmy go \(S_y\).
Rodzina \((S_y)_{y\in Y}\) pokrywa \(Y\).
Z założenia mamy skończone podpokrycie
\(S_{y_1},\ldots,S_{y_m}\).
Ultrafiltr zawiera co najmniej jeden
element skończonego pokrycia całej przestrzeni,
sprzecznie z wyborem \(S_y\).
Każdy ultrafiltr jest więc zbieżny,
a Twierdzenie~\ref{thm:topog-zwartosc-ultrafiltry}
kończy dowód.
\end{proof}

\begin{twierdzenie}[Zwartość topologii Fella]
\label{thm:topog21-fell-compact}
Dla dowolnej przestrzeni topologicznej \(X\)
przestrzeń \(\mathcal F(X)\) z topologią Fella
jest zwarta.
\end{twierdzenie}
\begin{proof}
Na mocy Lematu~\ref{lem:topog21-alexander}
wystarczy rozważyć pokrycie podbazowe
\[
 \mathcal F(X)=
  \bigcup_{i\in I}(U_i)^-_{\!F}\ \cup\
  \bigcup_{j\in J}C_j^0,
\]
gdzie \(U_i\) są otwarte, a \(C_j\) zwarte,
i pokazać, że ma skończone podpokrycie.
Przypuśćmy przeciwnie, że takiego podpokrycia
nie ma. Połóżmy
\(F=X\setminus\bigcup_{i\in I}U_i\).
Jest to zbiór domknięty, także gdy \(F=\varnothing\),
więc \(F\in\mathcal F(X)\), i żaden
\((U_i)^-_{\!F}\) go nie zawiera.

Dla ustalonego \(j\in J\) i każdego skończonego
\(I_0\subseteq I\) rodzina
\(\{(U_i)^-_{\!F}:i\in I_0\}\cup\{C_j^0\}\)
nie pokrywa \(\mathcal F(X)\).
Istnieje więc domknięty \(E\subseteq X\),
który nie trafia żadnego \(U_i\), \(i\in I_0\),
a jednocześnie trafia \(C_j\).
To znaczy
\[
 C_j\cap\bigcap_{i\in I_0}(X\setminus U_i)
 \neq\varnothing.
\]
Rodzina zbiorów \(C_j\cap(X\setminus U_i)\),
\(i\in I\), ma własność skończonych przecięć.
Ponieważ \(C_j\) jest zwarty, mamy
\(C_j\cap F\neq\varnothing\).
Zatem \(F\notin C_j^0\) dla każdego \(j\in J\).
Zbudowaliśmy element \(\mathcal F(X)\)
nienależący do żadnego zbioru pokrycia,
sprzeczność. Każde pokrycie podbazowe
ma zatem skończone podpokrycie.
\end{proof}

\begin{twierdzenie}[Oddzielanie w topologii Fella]
\label{thm:topog21-fell-hausdorff}
Jeżeli \(X\) jest lokalnie zwartą przestrzenią
Hausdorffa, to \(\mathcal F(X)\) z topologią Fella
jest zwartą przestrzenią Hausdorffa.
\end{twierdzenie}
\begin{proof}
Zwartość wynika z
Twierdzenia~\ref{thm:topog21-fell-compact}.
Niech \(F,G\in\mathcal F(X)\) będą różne.
Po zamianie ról możemy przyjąć, że
istnieje \(x\in F\setminus G\).
Zbiór \(X\setminus G\) jest otwartym
otoczeniem \(x\). Na mocy
Lematu~\ref{lem:topog-baire-zwarte-domkniecie}
istnieje otwarty \(W\ni x\) o zwartym
domknięciu \(C=\overline W^{\,X}\subseteq X\setminus G\).
Wtedy \(W^-_{\!F}\) jest otwartym otoczeniem \(F\),
a \(C^0\) otwartym otoczeniem \(G\).
Są one rozłączne, ponieważ \(W\subseteq C\).
\end{proof}

\begin{stwierdzenie}[Zbieżność Fella przez trafianie i omijanie]
\label{prop:topog21-fell-nets}
Niech \(X\) będzie przestrzenią topologiczną,
\((F_\lambda)_{\lambda\in D}\) siecią
w \(\mathcal F(X)\), a \(F\in\mathcal F(X)\).
Sieć \(F_\lambda\to F\) w topologii Fella
dokładnie wtedy, gdy spełnione są dwa warunki:
\begin{enumerate}[label=\textup{(\roman*)}]
\item dla każdego otwartego \(U\subseteq X\)
z \(F\cap U\neq\varnothing\)
od pewnego indeksu \(F_\lambda\cap U\neq\varnothing\);
\item dla każdego zwartego \(C\subseteq X\)
z \(F\cap C=\varnothing\)
od pewnego indeksu \(F_\lambda\cap C=\varnothing\).
\end{enumerate}
\end{stwierdzenie}
\begin{proof}
Są to dokładnie warunki, by sieć ostatecznie
należała do każdego podbazowego otoczenia \(F\).
Skończone przecięcie takich otoczeń
również jest ostatecznie trafiane:
dla skończonej liczby indeksów progowych
wybieramy wspólne górne ograniczenie
w skierowanym zbiorze \(D\).
Otrzymujemy zbieżność w topologii
generowanej przez podbazę.
\end{proof}

\section{Porównanie topologii Vietorisa i Fella}
\label{sec:topog21-comparison}

\begin{twierdzenie}[Zgodność w zwartej przestrzeni Hausdorffa]
\label{thm:topog21-fell-vietoris}
Jeżeli \(X\) jest zwartą przestrzenią Hausdorffa,
to na \(\mathcal K(X)=\mathcal F(X)\setminus\{\varnothing\}\)
topologia Vietorisa pokrywa się z topologią
indukowaną przez topologię Fella.
\end{twierdzenie}
\begin{proof}
W przestrzeni zwartej Hausdorffa
każdy niepusty zbiór domknięty jest zwarty,
więc obie rozważane rodziny punktów są równe.
Warunki dolne są identyczne.
Dla otwartego \(U\subseteq X\)
jego dopełnienie \(X\setminus U\) jest zwarte, zatem
\[
 U^+=(X\setminus U)^0
          \cap\mathcal K(X).
\]
Odwrotnie, dla dowolnego zwartego \(C\subseteq X\)
własność Hausdorffa zapewnia domkniętość \(C\),
a więc \(X\setminus C\) jest otwarty i
\[
 C^0\cap\mathcal K(X)=(X\setminus C)^+.
\]
Obie podbazy indukują tę samą topologię.
\end{proof}

\begin{przyklad}[Punkty uciekające do nieskończoności]
\label{ex:topog21-escaping}
Niech \(X=\mathbb R\).
W \(\mathcal F(\mathbb R)\) z topologią Fella
zachodzi \(\{n\}\to\varnothing\).
Istotnie, zbiór pusty nie ma otoczeń wymagających
trafienia jakiegokolwiek otwartego \(U\);
dla każdego zwartego \(C\subseteq\mathbb R\)
od pewnego \(n\) mamy \(n\notin C\).
Nie ma natomiast granicy tego ciągu
w \(\mathcal K(\mathbb R)\) z topologią Vietorisa,
bo \(d_H(\{n\},K)\to\infty\)
dla każdego ustalonego \(K\in\mathcal K(\mathbb R)\).

Ponadto \(K_n=\{0,n\}\) zbiega w topologii Fella
do \(K=\{0\}\): każdy otwarty zbiór trafiany
przez \(K\) jest trafiany przez \(K_n\),
a każdy zwarty \(C\) rozłączny z \(K\)
jest ostatecznie omijany przez \(K_n\),
ponieważ \(n\) opuszcza \(C\).
Jednak
\[
 d_H(K_n,K)=n
\]
dla \(n\ge1\), więc \(K_n\) nie zbiega
do \(K\) w topologii Vietorisa.
W szczególności na \(\mathcal K(\mathbb R)\)
topologia Fella indukowana z \(\mathcal F(\mathbb R)\)
jest ściśle słabsza od topologii Vietorisa.
\end{przyklad}

\begin{uwaga}[Dokładna rola zwartości]
\label{rem:topog21-restriction}
Dla dowolnej przestrzeni Hausdorffa \(X\)
ograniczenie podbazy Fella do \(\mathcal K(X)\)
jest zawarte w topologii Vietorisa:
dla zwartego \(C\) zbiór \(C\) jest domknięty,
więc warunek ominięcia \(C\) jest warunkiem
zawierania się w otwartym \(X\setminus C\).
Odwrotne zawieranie może zawieść,
jak pokazuje poprzedni przykład.
Nie wolno utożsamiać topologii Fella
z metryką Hausdorffa na niezwartym \(X\).
\end{uwaga}

\section{Podsumowanie}
\label{sec:topog21-summary}

W topologii Vietorisa punktami są niepuste
zbiory zwarte; wymaga się trafiania wskazanych
otwartych podzbiorów oraz zawierania w otwartym
otoczeniu całego zbioru. Dla przestrzeni metrycznej
topologia ta jest wyznaczana przez metrykę Hausdorffa.
Zupełność \(X\) przechodzi na \(\mathcal K(X)\);
zwartość \(X\) pociąga zwartość \(\mathcal K(X)\)
zarówno w przypadku metrycznym,
jak i w ogólnym przypadku zwartym Hausdorffa.

W topologii Fella punktami są wszystkie
zbiory domknięte, także pusty; wymagamy
trafiania otwartych i omijania zwartych
podzbiorów. Powstała przestrzeń jest zwarta
dla dowolnego \(X\), a Hausdorffa,
gdy \(X\) jest lokalnie zwarta i Hausdorffa.
Dla zwartej przestrzeni Hausdorffa
topologie Fella i Vietorisa zgadzają się
na niepustych zbiorach domkniętych.
Dla niezwartych przestrzeni zbiór może
utracić odległe punkty w granicy Fella,
co nie zachodzi przy zbieżności
w metryce Hausdorffa.
\chapter{Zbieżność Kuratowskiego oraz topologie Wijsmana i Attoucha--Wetsa}
\label{ch:topog22-zbieznosc-zbiorow}

\noindent\textbf{Cel rozdziału.}
W rozdziale~\ref{ch:topog-hiperprzestrzenie} wprowadziliśmy topologie
Vietorisa i Fella oraz metrykę Hausdorffa. Zbadamy teraz inne
sposoby opisywania granic zbiorów domkniętych: przez punkty,
które można przybliżać punktami kolejnych zbiorów, oraz przez funkcje
odległości od tych zbiorów.
Ustalimy dokładnie, kiedy te konstrukcje są równoważne.
Szczególną rolę odegra założenie, że domknięte kule są zwarte.
Nie zakładamy z góry, że rozważane zbiory domknięte są ograniczone.

\section{Granice dolna i górna Kuratowskiego}
\label{sec:topog22-kuratowski}

\begin{definicja}[Granice zbiorów w sensie Kuratowskiego]
\label{def:topog22-kuratowski}
Niech \((X,d)\) będzie przestrzenią metryczną
i niech \((F_n)_{n\ge1}\) będzie ciągiem jej podzbiorów.
Definiujemy \emph{granicę dolną Kuratowskiego}
oraz \emph{granicę górną Kuratowskiego} wzorami
\[
 \operatorname{Li}_{n\to\infty}F_n
 =\bigl\{x\in X:\ \forall\varepsilon>0\
     \exists N\ \forall n\ge N\
     F_n\cap B_d(x,\varepsilon)\neq\varnothing\bigr\},
\]
\[
 \operatorname{Ls}_{n\to\infty}F_n
 =\bigl\{x\in X:\ \forall\varepsilon>0\
     \forall N\ \exists n\ge N\
     F_n\cap B_d(x,\varepsilon)\neq\varnothing\bigr\}.
\]
Jeśli te zbiory są równe pewnemu \(F\subseteq X\),
mówimy, że \(F_n\) \emph{zbiega do \(F\) w sensie
Kuratowskiego}.
W tych definicjach dopuszczamy także zbiory puste.
\end{definicja}

\begin{stwierdzenie}[Opis przez funkcje odległości]
\label{prop:topog22-kuratowski-dist}
Niech \((X,d)\) będzie przestrzenią metryczną,
\((F_n)\) ciągiem podzbiorów \(X\), a
\(d(x,\varnothing)=+\infty\).
Wówczas
\[
 \operatorname{Li}F_n
 =\{x\in X:d(x,F_n)\longrightarrow0\},
 \qquad
 \operatorname{Ls}F_n
 =\{x\in X:\liminf_{n\to\infty}d(x,F_n)=0\}.
\]
Ponadto \(\operatorname{Li}F_n\subseteq\operatorname{Ls}F_n\),
a oba zbiory są domknięte w \(X\).
\end{stwierdzenie}
\begin{proof}
Dla dowolnego \(x\in X\) i \(\varepsilon>0\)
warunek \(F_n\cap B_d(x,\varepsilon)\neq\varnothing\)
jest równoważny \(d(x,F_n)<\varepsilon\).
Podstawiając tę równoważność do kwantyfikatorów
w Definicji~\ref{def:topog22-kuratowski}, otrzymujemy
oba wzory. Pierwszy warunek implikuje drugi.

Wykażemy domkniętość granicy dolnej bez założeń
o domkniętości \(F_n\).
Niech \(x\) należy do domknięcia \(\operatorname{Li}F_n\)
i niech \(\varepsilon>0\).
Wybierzmy \(y\in\operatorname{Li}F_n\)
z \(d(x,y)<\varepsilon/2\).
Od pewnego \(n\) zachodzi \(d(y,F_n)<\varepsilon/2\),
a z nierówności trójkąta
\(d(x,F_n)\le d(x,y)+d(y,F_n)<\varepsilon\).
Zatem \(x\in\operatorname{Li}F_n\).

Dla granicy górnej wybierzmy
\(x\in\overline{\operatorname{Ls}F_n}^{\,X}\).
Dla dowolnych \(\varepsilon>0\) i \(N\)
istnieje \(y\in\operatorname{Ls}F_n\)
z \(d(x,y)<\varepsilon/2\);
dla pewnego \(n\ge N\) mamy
\(d(y,F_n)<\varepsilon/2\).
Stąd \(d(x,F_n)<\varepsilon\).
Spełnione są wszystkie kwantyfikatory
definicji granicy górnej, więc
\(x\in\operatorname{Ls}F_n\).
\end{proof}

\begin{przyklad}[Brak granicy Kuratowskiego]
\label{ex:topog22-oscillating}
Niech \(X=\mathbb R\), \(F_{2n}=\{0\}\),
a \(F_{2n-1}=\{1\}\) dla \(n\ge1\).
Wówczas
\(\operatorname{Li}F_n=\varnothing\) oraz
\(\operatorname{Ls}F_n=\{0,1\}\).
Istotnie, żaden punkt nie ma wszystkich dostatecznie
późnych zbiorów w każdym swoim otoczeniu,
natomiast otoczenia \(0\) i \(1\) trafiają
nieskończenie wiele wyrazów ciągu.
Inne punkty nie należą do granicy górnej,
bo mają otoczenie rozłączne z \(\{0,1\}\).
\end{przyklad}

\begin{przyklad}[Punkt uciekający do nieskończoności]
\label{ex:topog22-escape}
W \(X=\mathbb R\) niech
\(F_n=\{0,n\}\).
Wtedy
\[
 \operatorname{Li}F_n=\operatorname{Ls}F_n=\{0\}.
\]
Dla \(x\neq0\) można bowiem wybrać kulę wokół \(x\)
rozłączną z \(\{0\}\), która od pewnego miejsca
nie zawiera również \(n\).
Nie zachodzi jednak zbieżność w metryce Hausdorffa,
ponieważ \(d_H(F_n,\{0\})=n\).
Przykład pokazuje potrzebę rozróżnienia
granicy lokalnej od globalnego pomiaru odległości zbiorów.

\begin{figure}[htbp]
\centering
\begin{tikzpicture}[x=1.08cm,y=.9cm,>=Latex]
 \draw[->,thin] (-.4,0)--(6.65,0);
 \foreach \x/\t in {0/0,1/1,2/2,3/3,4/4,5/5,6/6}
    {\draw (\x,-.08)--(\x,.08);
     \node[below,font=\scriptsize] at (\x,-.1) {$\t$};}
 \fill[blue!65!black] (0,.42) circle (2.5pt);
 \foreach \x in {2,4,6}
    {\fill[orange!80!black] (\x,.42) circle (2.5pt);}
 \node[above,blue!70!black,font=\small] at (0,.47) {punkt stały};
 \node[above,orange!85!black,font=\small] at (4,.47)
    {punkty oddalające się};
\end{tikzpicture}
\caption{Przy \(F_n=\{0,n\}\) punkt \(0\) pozostaje,
a drugi punkt opuszcza każdy ustalony zbiór ograniczony.}
\label{fig:topog22-escape}
\end{figure}
\end{przyklad}

\section{Topologia Wijsmana}
\label{sec:topog22-wijsman}

\begin{definicja}[Topologia Wijsmana]
\label{def:topog22-wijsman}
Niech \((X,d)\) będzie niepustą przestrzenią metryczną.
Przez \(\mathcal F_*(X)\) oznaczamy zbiór wszystkich
\emph{niepustych domkniętych} podzbiorów \(X\).
Dla \(x\in X\) definiujemy
\[
 \delta_x\colon\mathcal F_*(X)\to[0,\infty),
 \qquad \delta_x(F)=d(x,F).
\]
\emph{Topologią Wijsmana} na \(\mathcal F_*(X)\)
nazywamy najsłabszą topologię, względem której
wszystkie odwzorowania \(\delta_x\) są ciągłe.
W szczególności ciąg \(F_n\) jest w niej zbieżny
do \(F\in\mathcal F_*(X)\) dokładnie wtedy, gdy
\[
 d(x,F_n)\longrightarrow d(x,F)
 \quad\text{dla każdego }x\in X.
\]
\end{definicja}

\begin{uwaga}[Dlaczego wyłączamy zbiór pusty]
\label{rem:topog22-wijsman-empty}
Dla niepustego \(F\) funkcja \(d(\,\cdot\,,F)\)
ma wartości skończone, nawet gdy \(F\) jest nieograniczony.
Dla \(\varnothing\) przyjęliśmy wartość \(+\infty\).
Można osobno rozważać rozszerzoną topologię Wijsmana
z wartościami w \([0,\infty]\), lecz w tym rozdziale
porównania z metryką Attoucha--Wetsa dotyczą
wyłącznie \(\mathcal F_*(X)\).
Zbieżność w sensie Kuratowskiego i topologia Fella
nadal mogą obejmować \(\varnothing\).
\end{uwaga}

\begin{lemat}[Odległość od zbioru jest lipschitzowska]
\label{lem:topog22-distance-lipschitz}
Niech \((X,d)\) będzie przestrzenią metryczną
i \(F\subseteq X\) zbiorem niepustym.
Dla każdych \(x,y\in X\) zachodzi
\[
 |d(x,F)-d(y,F)|\le d(x,y).
\]
Jeśli \(F\) jest domknięty, to
\(F=\{x\in X:d(x,F)=0\}\).
\end{lemat}
\begin{proof}
Dla każdego \(z\in F\) mamy
\(d(x,F)\le d(x,z)\le d(x,y)+d(y,z)\).
Po wzięciu infimum względem \(z\in F\)
otrzymujemy \(d(x,F)\le d(x,y)+d(y,F)\).
Zamieniając \(x\) z \(y\), uzyskujemy pierwszą tezę.
Druga wynika z definicji domknięcia:
\(d(x,F)=0\) dokładnie wtedy, gdy
\(x\in\overline F^{\,X}\).
\end{proof}

\begin{twierdzenie}[Metryzowalność topologii Wijsmana]
\label{thm:topog22-wijsman-metrizable}
Niech \((X,d)\) będzie niepustą ośrodkową
przestrzenią metryczną i
\(D=\{q_1,q_2,\ldots\}\subseteq X\)
ustaloną przeliczalną rodziną gęstą
(dopuszczamy powtórzenia).
Wówczas wzór
\[
 \rho_W(A,B)=\sum_{j=1}^{\infty}2^{-j}
       \min\{1,\ |d(q_j,A)-d(q_j,B)|\},
 \qquad A,B\in\mathcal F_*(X),
\]
określa metrykę na \(\mathcal F_*(X)\),
której topologia jest topologią Wijsmana.
\end{twierdzenie}
\begin{proof}
Nieujemność i symetria są oczywiste,
a nierówność trójkąta wynika z
\(\min\{1,a+b\}\le\min\{1,a\}+\min\{1,b\}\)
dla \(a,b\ge0\).
Jeżeli \(\rho_W(A,B)=0\), to
\(d(q_j,A)=d(q_j,B)\) dla każdego \(j\).
Z Lematu~\ref{lem:topog22-distance-lipschitz}
obie funkcje odległości są ciągłe,
a \(D\) jest gęsty, zatem są równe na całym \(X\).
Ponieważ \(A,B\) są domknięte, ich zbiory zerowe
są równe, więc \(A=B\).

Zbieżność w metryce \(\rho_W\) jest równoważna
zbieżności \(d(q_j,F_n)\to d(q_j,F)\) dla każdego \(j\):
w jedną stronę wynika to z oszacowania
pojedynczego składnika, a w drugą ze zbieżności
ograniczonego szeregu, rozbijanego na
skończenie wiele pierwszych składników i mały ogon.
Dla dowolnego \(x\in X\), \(q_j\in D\)
i \(A,B\in\mathcal F_*(X)\) mamy
\[
 |d(x,A)-d(x,B)|
 \le 2d(x,q_j)+|d(q_j,A)-d(q_j,B)|.
\]
Dzięki gęstości \(D\) zbieżność na \(D\)
pociąga zbieżność na całym \(X\).
Dowodzi to równoważności z topologią Wijsmana.
Ten sam argument z otoczeniami zamiast ciągów
pokazuje równość topologii, a nie tylko
ich zbieżnych ciągów.
\end{proof}

\begin{twierdzenie}[Wijsman implikuje Kuratowskiego]
\label{thm:topog22-wijsman-kuratowski}
Niech \((X,d)\) będzie przestrzenią metryczną
i \(F,F_n\in\mathcal F_*(X)\) dla \(n\ge1\).
Jeżeli \(F_n\to F\) w topologii Wijsmana,
to \(F_n\to F\) w sensie Kuratowskiego.
\end{twierdzenie}
\begin{proof}
Dla każdego \(x\in F\) mamy \(d(x,F)=0\),
a więc \(d(x,F_n)\to0\).
Stwierdzenie~\ref{prop:topog22-kuratowski-dist}
daje \(x\in\operatorname{Li}F_n\).
Jeżeli \(x\notin F\), domkniętość \(F\) zapewnia
\(d(x,F)>0\). Stąd dla dostatecznie dużych \(n\)
zachodzi \(d(x,F_n)>d(x,F)/2\), więc
\(x\notin\operatorname{Ls}F_n\).
Otrzymujemy
\(F\subseteq\operatorname{Li}F_n
\subseteq\operatorname{Ls}F_n\subseteq F\).
\end{proof}

\section{Topologia Attoucha--Wetsa}
\label{sec:topog22-aw}

\begin{definicja}[Metryka i topologia Attoucha--Wetsa]
\label{def:topog22-aw}
Niech \((X,d)\) będzie niepustą przestrzenią
metryczną i niech \(x_0\in X\).
Dla \(m\ge1\) oraz \(A,B\in\mathcal F_*(X)\) połóżmy
\[
 \Delta_m(A,B)=
  \sup_{x\in\overline B_d(x_0,m)}
       |d(x,A)-d(x,B)|,
\]
\[
 \rho_{AW}(A,B)=\sum_{m=1}^{\infty}
          2^{-m}\min\{1,\Delta_m(A,B)\}.
\]
Topologię tej metryki nazywamy
\emph{topologią Attoucha--Wetsa}.
Zbiory \(\overline B_d(x_0,m)\)
są jedynie kulami domkniętymi:
w definicji nie zakładamy ich zwartości.
\end{definicja}

\begin{twierdzenie}[Metryka Attoucha--Wetsa i zbieżność na ograniczonych zbiorach]
\label{thm:topog22-aw-metric}
Niech \((X,d)\) będzie niepustą przestrzenią metryczną,
\(x_0\in X\), a \(\mathcal F_*(X)\) rodziną
niepustych domkniętych podzbiorów.
Funkcja \(\rho_{AW}\) z Definicji~\ref{def:topog22-aw}
jest metryką. Ciąg \((F_n)\subseteq\mathcal F_*(X)\)
zbiega w niej do \(F\in\mathcal F_*(X)\)
dokładnie wtedy, gdy dla każdego ograniczonego
zbioru \(B\subseteq X\) zachodzi
\[
 \sup_{x\in B}|d(x,F_n)-d(x,F)|\longrightarrow0.
\]
Topologia ta nie zależy od wyboru punktu \(x_0\).
\end{twierdzenie}
\begin{proof}
Dla ustalonych \(A,B\in\mathcal F_*(X)\)
oraz \(x\in\overline B_d(x_0,m)\) mamy
\(d(x,A)\le m+d(x_0,A)\) i analogicznie dla \(B\).
Zatem \(\Delta_m(A,B)\) jest skończona.
Nieujemność i symetria są oczywiste;
nierówność trójkąta dla \(\Delta_m\) wynika
z nierówności trójkąta dla wartości bezwzględnej,
a dla \(\rho_{AW}\) z nierówności
\(\min\{1,a+b\}\le\min\{1,a\}+\min\{1,b\}\).

Jeżeli \(\rho_{AW}(A,B)=0\), to wszystkie
\(\Delta_m(A,B)\) są zerowe.
Ponieważ każde \(x\in X\) należy do pewnej
\(\overline B_d(x_0,m)\), funkcje odległości
do \(A\) i do \(B\) są identyczne.
Lemat~\ref{lem:topog22-distance-lipschitz}
daje \(A=B\).

Z definicji szeregu wynika, że
\(\rho_{AW}(F_n,F)\to0\) wtedy i tylko wtedy,
gdy \(\Delta_m(F_n,F)\to0\) dla każdego \(m\).
Każdy ograniczony zbiór \(B\) zawiera się
w pewnej kuli \(\overline B_d(x_0,m)\),
więc otrzymujemy żądaną zbieżność jednostajną.
Implikacja odwrotna wynika bezpośrednio
z zastosowania założenia do każdej takiej kuli.
Ponieważ rodziny kul wokół dwóch różnych
punktów \(x_0,x_1\) są współkońcowe
w rodzinie wszystkich ograniczonych podzbiorów \(X\),
otrzymana topologia nie zależy od punktu bazowego.
\end{proof}

\begin{stwierdzenie}[Attouch--Wets jest silniejsza od Wijsmana]
\label{prop:topog22-aw-wijsman}
Dla każdej niepustej przestrzeni metrycznej \((X,d)\)
topologia Attoucha--Wetsa na \(\mathcal F_*(X)\)
jest co najmniej tak silna jak topologia Wijsmana.
\end{stwierdzenie}
\begin{proof}
Każdy singleton \(\{x\}\subseteq X\) jest ograniczony.
Zbieżność jednostajna funkcji odległości
na wszystkich ograniczonych zbiorach
pociąga zbieżność w każdym punkcie.
Ponadto każda ewaluacja
\(F\mapsto d(x,F)\) jest ciągła względem
topologii Attoucha--Wetsa, ponieważ
\[
 |d(x,A)-d(x,B)|\le\Delta_m(A,B)
\]
dla \(m\) takiego, że \(x\in\overline B_d(x_0,m)\).
Z najsłabszości topologii Wijsmana
wynika wymagane zawieranie topologii.
\end{proof}

\begin{przyklad}[Ścisłość na ograniczonej przestrzeni]
\label{ex:topog22-strict-aw}
Niech \(X=\mathbb N\) z metryką dyskretną:
\(d(k,l)=1\) dla \(k\neq l\).
Jest to przestrzeń metryczna ograniczona,
zupełna i ośrodkowa, ale niezwarta.
Niech \(F_n=X\setminus\{n\}\) i \(F=X\).
Dla każdego ustalonego \(k\in X\)
od pewnego \(n\) zachodzi \(k\in F_n\);
wtedy \(d(k,F_n)=0=d(k,F)\).
Zatem \(F_n\to F\) w topologii Wijsmana.

Jednak \(d(n,F_n)=1\), podczas gdy
\(d(n,F)=0\). Stąd
\[
 \sup_{x\in X}|d(x,F_n)-d(x,F)|=1
 \quad\text{dla każdego }n.
\]
Ponieważ \(X\) jest ograniczony,
nie ma zbieżności w topologii Attoucha--Wetsa.
Obie topologie są więc różne nawet na rodzinie
niepustych zbiorów domkniętych zupełnej,
ośrodkowej przestrzeni metrycznej.
\end{przyklad}

\section{Topologia Fella i właściwe przestrzenie metryczne}
\label{sec:topog22-proper}

\begin{definicja}[Przestrzeń metryczna właściwa]
\label{def:topog22-proper}
Przestrzeń metryczną \((X,d)\) nazywamy
\emph{właściwą}, jeżeli każda domknięta
kula \(\overline B_d(x,r)\), gdzie
\(x\in X\) i \(r\ge0\), jest zwarta.
Wówczas każdy domknięty ograniczony
podzbiór \(X\) jest zwarty.
\end{definicja}

\begin{lemat}[Zbieżność funkcji odległości na zwartym zbiorze]
\label{lem:topog22-equilipschitz}
Niech \((X,d)\) będzie przestrzenią metryczną,
\(K\subseteq X\) zbiorem zwartym,
a \((f_n)\) ciągiem funkcji
\(f_n\colon X\to\mathbb R\)
spełniających
\(|f_n(x)-f_n(y)|\le d(x,y)\) dla wszystkich \(x,y\in X\).
Jeżeli \(f_n(x)\to f(x)\) dla każdego \(x\in X\),
to \(f_n\to f\) jednostajnie na \(K\).
\end{lemat}
\begin{proof}
Przechodząc do granicy w nierówności,
stwierdzamy, że \(f\) również jest
\(1\)-lipschitzowska.
Jeśli \(K=\varnothing\), nie ma czego dowodzić.
Niech \(\varepsilon>0\).
Zwartość \(K\) daje punkty \(x_1,\ldots,x_m\in K\)
takie, że każdy \(x\in K\)
spełnia \(d(x,x_i)<\varepsilon/3\)
dla pewnego \(i\).
Dla dostatecznie dużych \(n\) mamy
\(|f_n(x_i)-f(x_i)|<\varepsilon/3\)
jednocześnie dla wszystkich \(i\).
Wtedy
\[
 |f_n(x)-f(x)|
 \le |f_n(x)-f_n(x_i)|
     +|f_n(x_i)-f(x_i)|+|f(x_i)-f(x)|
 <\varepsilon.
\]
\end{proof}

\begin{twierdzenie}[Równoważność Wijsmana i Attoucha--Wetsa]
\label{thm:topog22-proper-wijsman-aw}
Niech \((X,d)\) będzie niepustą właściwą
przestrzenią metryczną. Wówczas
topologia Wijsmana na \(\mathcal F_*(X)\)
pokrywa się z topologią Attoucha--Wetsa.
\end{twierdzenie}
\begin{proof}
Zawieranie topologii Wijsmana w Attoucha--Wetsa
wynika ze Stwierdzenia~\ref{prop:topog22-aw-wijsman}.
Odwrotnie, dla dowolnej zbieżnej w sensie Wijsmana
sieci \((F_\lambda)\) do \(F\)
funkcje \(d(\,\cdot\,,F_\lambda)\)
są \(1\)-lipschitzowskie i zbieżne punktowo
do \(d(\,\cdot\,,F)\).
Każda kula \(\overline B_d(x_0,m)\) jest zwarta.
Argument skończonej \(\varepsilon\)-sieci
z Lematu~\ref{lem:topog22-equilipschitz}
działa również dla sieci:
dla skończenie wielu punktów bazowych
wybiera się wspólny indeks dominujący.
Zatem zbieżność jest jednostajna na każdej
kuli i na każdym zbiorze ograniczonym.
Twierdzenie~\ref{thm:topog22-aw-metric}
pokazuje, że jest to zbieżność
w topologii Attoucha--Wetsa.
W szczególności identyczność z topologii Wijsmana
do topologii Attoucha--Wetsa jest ciągła.
\end{proof}

\begin{twierdzenie}[Zbieżność Fella i Kuratowskiego]
\label{thm:topog22-fell-kuratowski}
Niech \((X,d)\) będzie lokalnie zwartą
przestrzenią metryczną, a \(F,F_n\subseteq X\)
zbiorami domkniętymi, przy czym dopuszczamy
\(F=\varnothing\) i \(F_n=\varnothing\).
Wówczas \(F_n\to F\) w topologii Fella
z Definicji~\ref{def:topog21-fell}
dokładnie wtedy, gdy
\[
 \operatorname{Li}F_n=\operatorname{Ls}F_n=F.
\]
\end{twierdzenie}
\begin{proof}
Niech najpierw \(F_n\to F\) w topologii Fella.
Dla każdego \(x\in F\) i \(\varepsilon>0\)
otwarty zbiór \(B_d(x,\varepsilon)\) trafia \(F\),
więc na mocy warunku dolnego
zbieżności Fella jest od pewnego miejsca
trafiany przez wszystkie \(F_n\).
Stąd \(F\subseteq\operatorname{Li}F_n\).

Niech \(x\notin F\).
Z lokalnej zwartości i własności Hausdorffa
oraz Lematu~\ref{lem:topog-baire-zwarte-domkniecie}
istnieje otwarte \(W\ni x\)
o zwartym domknięciu \(C\subseteq X\setminus F\).
Warunek ominięcia zwartego \(C\)
zapewnia \(F_n\cap C=\varnothing\)
od pewnego miejsca, a zatem
\(F_n\cap W=\varnothing\) od pewnego miejsca.
Stąd \(x\notin\operatorname{Ls}F_n\).
Otrzymujemy
\(F\subseteq\operatorname{Li}F_n
\subseteq\operatorname{Ls}F_n\subseteq F\).

Odwrotnie, niech obie granice Kuratowskiego
są równe \(F\).
Jeżeli otwarty \(U\) trafia \(F\),
wybierzmy \(x\in F\cap U\) i
\(\varepsilon>0\) tak, aby
\(B_d(x,\varepsilon)\subseteq U\).
Z \(x\in\operatorname{Li}F_n\)
wynika, że \(F_n\) ostatecznie trafiają \(U\).

Niech \(C\subseteq X\) będzie zwarty
i rozłączny z \(F\).
Gdyby nieskończenie wiele \(F_n\) trafiało \(C\),
moglibyśmy wybrać rosnący ciąg indeksów \(n_k\)
oraz \(x_k\in F_{n_k}\cap C\).
Zwartość metryczna \(C\) daje podciąg
\(x_{k_j}\to x\in C\).
Dla każdego \(\varepsilon>0\) i \(N\)
istnieje \(j\) z \(n_{k_j}\ge N\)
oraz \(d(x_{k_j},x)<\varepsilon\).
To oznacza \(x\in\operatorname{Ls}F_n=F\),
wbrew \(C\cap F=\varnothing\).
Zatem \(F_n\) ostatecznie omijają \(C\).
Oba warunki zbieżności Fella są spełnione
na mocy Stwierdzenia~\ref{prop:topog21-fell-nets}.
\end{proof}

\begin{twierdzenie}[Zbieżności są równoważne w przestrzeni właściwej]
\label{thm:topog22-proper-equivalences}
Niech \((X,d)\) będzie niepustą właściwą
przestrzenią metryczną oraz
\(F,F_n\in\mathcal F_*(X)\) dla wszystkich \(n\ge1\).
Następujące warunki są równoważne:
\begin{enumerate}[label=\textup{(\roman*)}]
\item \(F_n\to F\) w sensie Kuratowskiego;
\item \(F_n\to F\) w topologii Fella;
\item \(F_n\to F\) w topologii Wijsmana;
\item \(F_n\to F\) w topologii Attoucha--Wetsa.
\end{enumerate}
\end{twierdzenie}
\begin{proof}
Przestrzeń właściwa jest lokalnie zwarta,
ponieważ kula domknięta o dodatnim promieniu
jest zwartym otoczeniem swojego środka.
Twierdzenie~\ref{thm:topog22-fell-kuratowski}
daje \textup{(i)} \(\Leftrightarrow\) \textup{(ii)}.
Twierdzenia~\ref{thm:topog22-wijsman-kuratowski}
i~\ref{thm:topog22-proper-wijsman-aw}
oraz Stwierdzenie~\ref{prop:topog22-aw-wijsman}
dają implikacje \textup{(iv)} \(\Rightarrow\)
\textup{(iii)} \(\Rightarrow\) \textup{(i)}
oraz \textup{(iii)} \(\Leftrightarrow\) \textup{(iv)}.
Pozostaje dowieść \textup{(i)} \(\Rightarrow\) \textup{(iii)}.

Ustalmy \(x\in X\) i połóżmy \(r=d(x,F)\).
Ponieważ \(F\neq\varnothing\), liczba \(r\) jest skończona.
Dla każdego \(\varepsilon>0\) wybierzmy \(z\in F\)
takie, że \(d(x,z)<r+\varepsilon\).
Z \(z\in\operatorname{Li}F_n\)
wynika \(d(z,F_n)\to0\).
Stąd
\[
 \limsup_{n\to\infty}d(x,F_n)
 \le d(x,z)<r+\varepsilon.
\]
Po przejściu z \(\varepsilon\downarrow0\)
dostajemy \(\limsup_nd(x,F_n)\le r\).

Jeżeli \(r=0\), nieujemność odległości
kończy dowód. Niech \(r>0\).
Gdyby \(\liminf_nd(x,F_n)<r\),
istniałyby \(\eta>0\) i rosnące indeksy \(n_k\)
takie, że
\(d(x,F_{n_k})<r-\eta\).
Wybierzmy \(z_k\in F_{n_k}\)
z \(d(x,z_k)<r-\eta/2\).
Wszystkie \(z_k\) należą do zwartej
kuli \(\overline B_d(x,r)\);
po przejściu do podciągu możemy
przyjąć \(z_k\to z\).
Ponieważ \(n_k\to\infty\)
i \(z_k\in F_{n_k}\),
punkt \(z\) należy do
\(\operatorname{Ls}F_n=F\).
Jednocześnie \(d(x,z)\le r-\eta/2<r\),
wbrew definicji \(r=d(x,F)\).
Stąd \(\liminf_nd(x,F_n)\ge r\)
i ostatecznie \(d(x,F_n)\to d(x,F)\).
Punkt \(x\) był dowolny,
więc otrzymujemy topologię Wijsmana.
\end{proof}

\begin{przyklad}[Lokalna zwartość nie wystarcza dla Wijsmana]
\label{ex:topog22-nonproper}
Niech \(X=\{p,q\}\cup\{y_n:n\ge1\}\)
i niech odległości różnych punktów spełniają
\[
 d(p,q)=2,\qquad
 d(p,y_n)=d(q,y_n)=1,\qquad
 d(y_n,y_m)=1\quad(n\neq m).
\]
Są to odległości metryczne:
jedyny odcinek długości \(2\)
ma końce \(p,q\), a każda droga
łącząca je przez trzeci punkt ma długość \(2\).
Przestrzeń jest dyskretna, przeliczalna,
ograniczona, zupełna i lokalnie zwarta,
lecz nie jest właściwa, ponieważ sama
jest domkniętą kulą promienia \(2\)
i nie jest zwarta.

Niech \(F_n=\{p,y_n\}\), \(F=\{p\}\).
Każdy punkt \(y_k\) należy tylko do jednego \(F_n\),
a \(p\) należy do wszystkich.
Ponieważ topologia \(X\) jest dyskretna,
\(\operatorname{Li}F_n=\operatorname{Ls}F_n=\{p\}\).
Stąd \(F_n\to F\) również w topologii Fella
na mocy Twierdzenia~\ref{thm:topog22-fell-kuratowski}.
Jednak
\[
 d(q,F_n)=1,\qquad d(q,F)=2
 \quad\text{dla każdego }n,
\]
więc nie ma zbieżności Wijsmana.
Pokazuje to, dlaczego w
Twierdzeniu~\ref{thm:topog22-proper-equivalences}
potrzebna jest dodatkowa zwartość kul,
a sama lokalna zwartość nie wystarcza.
\end{przyklad}

\begin{przyklad}[Zbieżność lokalna nie oznacza zbieżności Hausdorffa]
\label{ex:topog22-aw-not-hausdorff}
W przestrzeni właściwej \(X=\mathbb R\)
niech \(F_n=\{0,n\}\) i \(F=\{0\}\).
Dla każdej ustalonej kuli \([-m,m]\)
oraz \(n>2m\) mamy
\(d(x,F_n)=|x|=d(x,F)\) dla \(x\in[-m,m]\).
Zatem \(\Delta_m(F_n,F)=0\) od pewnego miejsca
i \(F_n\to F\) w topologii Attoucha--Wetsa.
Jednocześnie \(d_H(F_n,F)=n\).
W konsekwencji nawet na rodzinie
niepustych zbiorów zwartych w \(\mathbb R\)
topologia Vietorisa z rozdziału~21
jest silniejsza od topologii Attoucha--Wetsa.
\end{przyklad}

\section{Podsumowanie i zakres twierdzeń}
\label{sec:topog22-summary}

Granice Kuratowskiego wymagają jedynie struktury metrycznej
i opisują punkty dające się przybliżać punktami kolejnych zbiorów.
Topologia Wijsmana bada punktową zbieżność funkcji odległości
do niepustych zbiorów domkniętych,
a topologia Attoucha--Wetsa --- ich zbieżność jednostajną
na zbiorach ograniczonych.
Ta ostatnia jest zawsze silniejsza od Wijsmana;
właściwość przestrzeni metrycznej czyni je równymi.
Dla lokalnie zwartych przestrzeni metrycznych
zbieżność ciągów w topologii Fella jest równoważna
zbieżności w sensie Kuratowskiego.
W przestrzeniach właściwych wszystkie cztery
kryteria są równoważne dla ciągów zbiorów
niepustych domkniętych o niepustej granicy.

Nie są to jednak synonimy metryki Hausdorffa:
punkt, który opuszcza każdy zbiór ograniczony,
może zanikać w granicy Kuratowskiego, Fella,
Wijsmana i Attoucha--Wetsa, chociaż odległość
Hausdorffa do granicy rośnie bez ograniczenia.
Zbiory puste dopuszczamy w granicach Kuratowskiego
i w topologii Fella, lecz nie w przyjętych
tu definicjach Wijsmana i Attoucha--Wetsa.
\chapter{Wymiar topologiczny}
\label{ch:topog23-wymiar}

\noindent\textbf{Cel rozdziału.}
W przestrzeni euklidesowej liczba współrzędnych sugeruje
pojęcie wymiaru. Dla dowolnych przestrzeni topologicznych
nie mamy jednak wyróżnionych współrzędnych.
Wprowadzimy trzy niezmienniki oparte wyłącznie na topologii:
wymiar pokryciowy Lebesgue'a, mały wymiar indukcyjny
oraz duży wymiar indukcyjny. Wyjaśnimy ich sens,
podamy ścisłe charakterystyki wymiaru zerowego
i przeprowadzimy samodzielne dowody wymiaru odcinka
oraz górnego oszacowania wymiaru zwartych podzbiorów
przestrzeni euklidesowych. Pokażemy także rolę
nerwu pokrycia i map do kompleksów symplicjalnych.
W głównych twierdzeniach o wymiarze pokryciowym
pracujemy z przestrzeniami zwartymi Hausdorffa,
aby wszystkie argumenty o pokryciach miały jawnie
podane założenia.

\section{Rząd pokrycia i wymiar Lebesgue'a}
\label{sec:topog23-cover}

\begin{definicja}[Pokrycia, rozdrobnienia i rząd]
\label{def:topog23-order}
Niech \(X\) będzie przestrzenią topologiczną.
Rodzinę \(\mathcal U=(U_i)_{i\in I}\) nazywamy
\emph{pokryciem} \(X\), jeśli
\(X=\bigcup_{i\in I}U_i\).
Pokrycie \(\mathcal V=(V_j)_{j\in J}\)
jest \emph{rozdrobnieniem} \(\mathcal U\),
jeżeli dla każdego \(j\in J\) istnieje
\(i\in I\) takie, że \(V_j\subseteq U_i\).
Rodziny traktujemy jako indeksowane,
więc dwa wystąpienia tego samego zbioru
z różnymi indeksami liczymy osobno.

\emph{Rząd pokrycia} \(\mathcal U\)
jest równy co najwyżej \(m\in\mathbb N\),
jeżeli każdy punkt \(x\in X\)
należy do co najwyżej \(m\) zbiorów \(U_i\).
Inaczej: dla każdych różnych indeksów
\(i_0,\ldots,i_m\in I\) zachodzi
\(U_{i_0}\cap\cdots\cap U_{i_m}=\varnothing\).
\end{definicja}

\begin{definicja}[Wymiar pokryciowy Lebesgue'a]
\label{def:topog23-dim}
Niech \(X\) będzie zwartą przestrzenią Hausdorffa.
Przyjmujemy \(\dim\varnothing=-1\).
Dla niepustego \(X\) i \(n\in\{0,1,2,\ldots\}\)
piszemy \(\dim X\le n\), jeżeli
każde skończone pokrycie otwarte \(X\)
ma skończone otwarte rozdrobnienie
rzędu co najwyżej \(n+1\).
Liczbą \(\dim X\) nazywamy najmniejsze
takie \(n\), o ile istnieje;
w przeciwnym razie piszemy \(\dim X=\infty\).
\end{definicja}

\begin{uwaga}[Zakres definicji]
\label{rem:topog23-normality}
Dla przestrzeni niezwartych używa się również
ogólnej definicji wymiaru pokryciowego,
z odpowiednimi warunkami na rozdrobnienia.
Tutaj ograniczenie do zwartych przestrzeni Hausdorffa
zapewnia równoważność sformułowań przez
pokrycia skończone i dowolne pokrycia otwarte:
z dowolnego pokrycia wybiera się skończone
podpokrycie, które następnie rozdrabnia się
do żądanego rzędu.
Nie należy utożsamiać tego pojęcia
z wymiarem Hausdorffa metrycznej przestrzeni.
\end{uwaga}

\begin{stwierdzenie}[Niezmienniczość homeomorficzna]
\label{prop:topog23-homeomorphism}
Niech \(X,Y\) będą zwartymi przestrzeniami Hausdorffa
i \(h\colon X\to Y\) homeomorfizmem.
Wówczas \(\dim X=\dim Y\).
\end{stwierdzenie}
\begin{proof}
Załóżmy, że \(\dim X\le n\).
Dla skończonego otwartego pokrycia
\((U_i)\) przestrzeni \(Y\)
przeciwobrazy \(h^{-1}(U_i)\)
tworzą skończone otwarte pokrycie \(X\).
Wybierzmy jego skończone otwarte rozdrobnienie
\((V_j)\) rzędu co najwyżej \(n+1\).
Rodzina \((h(V_j))\) jest otwartym
pokryciem \(Y\), rozdrabnia \((U_i)\),
a bijektywność \(h\) zachowuje
liczbę zbiorów zawierających punkt.
Stąd \(\dim Y\le n\).
Stosujemy ten sam argument do \(h^{-1}\).
W przypadku wymiaru nieskończonego
równość wynika z obu implikacji.
Przypadek pusty jest oczywisty.
\end{proof}

\begin{lemat}[Zamknięte zmniejszenie skończonego pokrycia]
\label{lem:topog23-closed-shrinking}
Niech \(X\) będzie zwartą przestrzenią Hausdorffa,
a \(U_1,\ldots,U_m\subseteq X\)
skończonym pokryciem otwartym,
gdzie \(m\ge1\).
Istnieją zbiory domknięte \(F_i\subseteq U_i\)
dla \(1\le i\le m\), takie że
\(X=\bigcup_{i=1}^mF_i\).
\end{lemat}
\begin{proof}
Przestrzeń \(X\) jest normalna
na mocy Twierdzenia~\ref{thm:topog-zwartosc-hausdorff-normalna}.
Dla każdego \(x\in X\) wybieramy \(i(x)\)
z \(x\in U_{i(x)}\).
Regularność pozwala wybrać
otwarty \(W_x\ni x\) spełniający
\(\overline{W_x}^{\,X}\subseteq U_{i(x)}\).
Ze zwartości \(X\) istnieją
\(x_1,\ldots,x_k\in X\) takie, że
\(X=\bigcup_{j=1}^k W_{x_j}\).
Dla \(1\le i\le m\) definiujemy
\[
 F_i=\bigcup_{\{j:i(x_j)=i\}}
           \overline{W_{x_j}}^{\,X},
\]
przyjmując pusty zbiór, gdy wskazana rodzina
indeksów jest pusta.
Każdy \(F_i\) jest skończoną sumą
zbiorów domkniętych, więc jest domknięty.
Z konstrukcji \(F_i\subseteq U_i\),
a suma wszystkich \(F_i\) zawiera
sumę zbiorów \(W_{x_j}\), czyli całe \(X\).
\end{proof}

\begin{twierdzenie}[Dziedziczenie przez zbiory domknięte]
\label{thm:topog23-closed-monotone}
Niech \(X\) będzie zwartą przestrzenią Hausdorffa,
a \(F\subseteq X\) zbiorem domkniętym.
Wówczas \(\dim F\le\dim X\)
w znaczeniu: jeśli \(\dim X\le n\)
dla pewnego \(n\ge0\), to \(\dim F\le n\).
\end{twierdzenie}
\begin{proof}
Gdy \(F=\varnothing\), teza wynika z definicji.
Dla niepustego \(F\) weźmy dowolne
skończone otwarte pokrycie
\(U_1,\ldots,U_m\) przestrzeni \(F\).
Istnieją otwarte \(W_i\subseteq X\)
takie, że \(U_i=W_i\cap F\).
Rodzina \(W_1,\ldots,W_m,X\setminus F\)
jest skończonym otwartym pokryciem \(X\).
Wybierzmy otwarte rozdrobnienie
\((V_j)_{j=1}^k\) rzędu co najwyżej \(n+1\).
Zbiory \(V_j\cap F\) niepuste,
po odrzuceniu tych \(V_j\),
które rozdrabniają \(X\setminus F\),
tworzą otwarte pokrycie \(F\).
Każdy z nich zawiera się w jednym z \(U_i\).
Ich rząd nie przekracza rzędu \((V_j)\),
więc \(\dim F\le n\).
\end{proof}

\section{Wymiar zero i bazy zbiorów otwarto-domkniętych}
\label{sec:topog23-zero}

\begin{twierdzenie}[Charakteryzacja wymiaru zerowego]
\label{thm:topog23-zero-clopen}
Niech \(X\) będzie niepustą zwartą przestrzenią Hausdorffa.
Następujące warunki są równoważne:
\begin{enumerate}[label=\textup{(\roman*)}]
\item \(\dim X=0\);
\item każdy punkt ma bazę otoczeń
złożoną ze zbiorów jednocześnie otwartych
i domkniętych w \(X\);
\item każde skończone otwarte pokrycie \(X\)
ma skończone rozdrobnienie będące
pokryciem przez parami rozłączne
zbiory otwarto-domknięte.
\end{enumerate}
\end{twierdzenie}
\begin{proof}
Załóżmy \textup{(i)}.
Niech \(x\in U\subseteq X\), gdzie \(U\) jest otwarty.
Zwartość Hausdorffa implikuje aksjomat \(T_1\),
więc \(\{U,X\setminus\{x\}\}\)
jest skończonym otwartym pokryciem \(X\).
Istnieje skończone otwarte rozdrobnienie
rzędu co najwyżej \(1\).
Ponieważ rozdrobnienie pokrywa \(X\),
jego niepuste elementy są parami rozłączne.
Jest wśród nich dokładnie jeden \(V\ni x\).
Zbiór \(V\) musi zawierać się w \(U\),
gdyż \(X\setminus\{x\}\) nie zawiera \(x\).
Dopełnienie \(V\) jest sumą pozostałych
elementów otwartego pokrycia,
zatem jest otwarte.
Wobec tego \(V\) jest otwarto-domknięty
i \(x\in V\subseteq U\).
Dowiedliśmy \textup{(ii)}.

Załóżmy \textup{(ii)} i rozważmy
skończone otwarte pokrycie \((U_i)_{i=1}^m\).
Dla każdego \(x\in X\) możemy wybrać
otwarto-domknięty \(C_x\) taki, że
\(x\in C_x\subseteq U_{i(x)}\).
Ze zwartości \(X\) wybieramy skończone
podpokrycie \(C_1,\ldots,C_k\).
Definiujemy
\[
 D_1=C_1,\qquad
 D_j=C_j\setminus\bigcup_{l<j}C_l
 \quad (2\le j\le k).
\]
Zbiory \(D_j\) są otwarto-domknięte,
parami rozłączne i pokrywają \(X\);
każdy zawiera się w odpowiednim \(C_j\),
a więc rozdrabnia pierwotne pokrycie.
Możemy pominąć wyrazy puste.
Dostajemy \textup{(iii)}.

Warunek \textup{(iii)} jest wprost
warunkiem \(\dim X\le0\).
Ponieważ \(X\neq\varnothing\),
otrzymujemy \(\dim X=0\).
\end{proof}

\begin{stwierdzenie}[Przestrzenie skończone i dyskretne]
\label{prop:topog23-finite-zero}
Każda niepusta skończona przestrzeń dyskretna
ma wymiar pokryciowy \(0\).
\end{stwierdzenie}
\begin{proof}
Przestrzeń skończona dyskretna jest zwarta
i Hausdorffa. Jej wszystkie singletony
są otwarto-domknięte, a zatem tworzą bazę.
Twierdzenie~\ref{thm:topog23-zero-clopen}
daje tezę.
\end{proof}

\begin{stwierdzenie}[Iloczyny przestrzeni wymiaru zero]
\label{prop:topog23-zero-products}
Niech \((X_i)_{i\in I}\) będzie dowolną
rodziną niepustych zwartych przestrzeni
Hausdorffa o wymiarze pokryciowym \(0\).
Wówczas iloczyn
\(X=\prod_{i\in I}X_i\)
z topologią produktową ma wymiar \(0\).
Dotyczy to także \(I=\varnothing\),
gdy \(X\) jest jednoelementowy.
\end{stwierdzenie}
\begin{proof}
Twierdzenie Tychonoffa
\ref{thm:topog-tychonoff-iloczyny}
zapewnia zwartość \(X\),
a własność Hausdorffa przechodzi na iloczyny.
Ponieważ każdy \(X_i\) ma bazę
zbiorów otwarto-domkniętych,
dla każdego punktu \(x\in X\)
i podstawowego otoczenia
\[
 O=\bigcap_{j\in J}p_j^{-1}(U_j),
 \qquad J\subseteq I\text{ skończony},
\]
możemy wybrać otwarto-domknięte
\(C_j\subseteq U_j\) z \(x_j\in C_j\).
Wówczas
\(C=\bigcap_{j\in J}p_j^{-1}(C_j)\)
jest zbiorem otwarto-domkniętym
z \(x\in C\subseteq O\).
Zbiory te tworzą bazę iloczynu.
Na mocy Twierdzenia~\ref{thm:topog23-zero-clopen}
wymiar \(X\) jest zerowy.
W przypadku \(I=\varnothing\)
mamy przestrzeń jednoelementową,
a więc również wymiar zero.
\end{proof}

\begin{przyklad}[Zbiór Cantora]
\label{ex:topog23-cantor}
Niech
\[
 C=\left\{\sum_{j=1}^{\infty}
          \frac{2\varepsilon_j}{3^j}:
          \varepsilon_j\in\{0,1\}\right\}
 \subseteq[0,1]
\]
będzie klasycznym zbiorem Cantora.
Dla każdej skończonej listy cyfr
\(\varepsilon_1,\ldots,\varepsilon_m\)
zbiór punktów \(C\) o wskazanych
pierwszych \(m\) cyfrach jest
otwarto-domknięty w \(C\).
Istotnie, jest obrazem podstawowego
zbioru cylindrycznego w
\(\{0,1\}^{\mathbb N}\).
Odwzorowanie kodujące cyfry jest ciągłe,
gdyż ogon szeregu ma moduł co najwyżej \(3^{-m}\);
jest też różnowartościowe:
przy pierwszej różnicy na miejscu \(m\)
różnica odpowiadających sum ma moduł
co najmniej
\(2\cdot3^{-m}-\sum_{j>m}2\cdot3^{-j}
=3^{-m}>0\).
Jest zatem homeomorfizmem zwartej przestrzeni
\(\{0,1\}^{\mathbb N}\) na \(C\).
Podstawowe cylindry są otwarto-domknięte,
stąd \(\dim C=0\).
Równoważnie, można zastosować
Stwierdzenie~\ref{prop:topog23-zero-products}
i niezmienniczość homeomorficzną.
\end{przyklad}

\begin{twierdzenie}[Niezerowy wymiar przestrzeni spójnej]
\label{thm:topog23-connected-positive}
Niech \(X\) będzie zwartą przestrzenią Hausdorffa
zawierającą przynajmniej dwa punkty.
Jeśli \(X\) jest spójna, to \(\dim X\ge1\).
\end{twierdzenie}
\begin{proof}
Wybierzmy różne \(a,b\in X\).
Zbiory \(X\setminus\{a\}\) i
\(X\setminus\{b\}\) są otwarte i właściwe,
a ich suma pokrywa \(X\).
Gdyby \(\dim X=0\), na mocy
Twierdzenia~\ref{thm:topog23-zero-clopen}
pokrycie miałoby rozdrobnienie
przez parami rozłączne
zbiory otwarto-domknięte.
Ponieważ \(X\) jest spójna,
każdy taki zbiór jest pusty albo równy \(X\).
Niepuste rozdrobnienie musiałoby
więc składać się z samego \(X\),
ale \(X\) nie jest zawarty w żadnym
z dwóch wyjściowych zbiorów.
Otrzymujemy sprzeczność.
\end{proof}

\section{Wymiary indukcyjne}
\label{sec:topog23-inductive}

\begin{definicja}[Mały wymiar indukcyjny]
\label{def:topog23-ind}
Niech \(X\) będzie zwartą przestrzenią Hausdorffa.
Definiujemy rekurencyjnie
\(\operatorname{ind}\varnothing=-1\).
Dla \(n\ge0\) piszemy
\(\operatorname{ind}X\le n\), jeśli
dla każdego \(x\in X\) i każdego otwartego
otoczenia \(U\ni x\) istnieje otwarty zbiór \(V\)
spełniający
\[
 x\in V\subseteq\overline V^{\,X}\subseteq U,
 \qquad
 \operatorname{ind}(\partial_XV)\le n-1,
\]
gdzie
\(\partial_XV=\overline V^{\,X}\setminus V\).
Mały wymiar indukcyjny
\(\operatorname{ind}X\) jest najmniejszym
takim \(n\), a w przypadku braku
jakiegokolwiek \(n\) wynosi \(\infty\).
\end{definicja}

\begin{definicja}[Duży wymiar indukcyjny]
\label{def:topog23-Ind}
Niech \(X\) będzie zwartą przestrzenią Hausdorffa.
Przyjmujemy \(\operatorname{Ind}\varnothing=-1\).
Dla \(n\ge0\) piszemy
\(\operatorname{Ind}X\le n\), jeżeli
dla każdej pary \(F\subseteq U\subseteq X\),
gdzie \(F\) jest domknięty, a \(U\) otwarty,
istnieje otwarty \(V\subseteq X\) taki, że
\[
 F\subseteq V\subseteq\overline V^{\,X}\subseteq U,
 \qquad
 \operatorname{Ind}(\partial_XV)\le n-1.
\]
Definiujemy \(\operatorname{Ind}X\)
jako najmniejsze odpowiednie \(n\)
lub \(\infty\), jeśli takie \(n\) nie istnieje.
\end{definicja}

\begin{uwaga}[Znaczenie warunku o brzegu]
\label{rem:topog23-boundary-ind}
Wymiar indukcyjny opisuje nie tylko
liczbę zbiorów spotykających się w punkcie,
lecz możliwość wyboru otoczeń,
których brzeg ma mniejszy wymiar.
Dla \(n=0\) warunek
\(\operatorname{ind}(\partial_XV)\le-1\)
oznacza \(\partial_XV=\varnothing\).
Wtedy \(V\) jest otwarto-domknięty.
Wymiar mały dotyczy otoczeń punktów;
wymiar duży --- oddzielania
całych zbiorów domkniętych od dopełnienia
otwartego otoczenia.
\end{uwaga}

\begin{twierdzenie}[Zgodność trzech wymiarów w wymiarze zero]
\label{thm:topog23-three-zero}
Niech \(X\) będzie niepustą zwartą
przestrzenią Hausdorffa.
Wówczas następujące warunki są równoważne:
\[
 \dim X=0,\qquad
 \operatorname{ind}X=0,\qquad
 \operatorname{Ind}X=0.
\]
\end{twierdzenie}
\begin{proof}
Z definicji małego wymiaru indukcyjnego
warunek \(\operatorname{ind}X\le0\)
jest równoważny istnieniu bazy
otwarto-domkniętych otoczeń każdego punktu.
Istotnie, w takim przypadku
\(\partial_XV=\varnothing\), a ponieważ
\(V\) jest otwarty, oznacza to,
że \(V\) jest także domknięty.
Odwrotnie, otwarto-domknięte \(V\)
ma pusty brzeg.
Twierdzenie~\ref{thm:topog23-zero-clopen}
daje równoważność
\(\dim X=0\Longleftrightarrow
 \operatorname{ind}X=0\).

Jeśli \(\operatorname{Ind}X=0\),
stosujemy definicję dużego wymiaru
do domkniętego singletonu \(F=\{x\}\)
i dowolnego otwartego \(U\ni x\).
Otrzymujemy otwarto-domknięty zbiór \(V\)
z \(x\in V\subseteq U\);
stąd \(\operatorname{ind}X=0\).

Załóżmy wreszcie \(\dim X=0\).
Niech \(F\subseteq U\subseteq X\),
gdzie \(F\) jest domknięty, a \(U\) otwarty.
Zwartość \(F\) oraz baza zbiorów
otwarto-domkniętych pozwalają wybrać
skończenie wiele takich zbiorów
\(C_1,\ldots,C_k\subseteq U\),
które pokrywają \(F\).
Ich suma \(V=\bigcup_{j=1}^k C_j\)
jest otwarto-domknięta,
a \(F\subseteq V\subseteq U\).
W szczególności \(\partial_XV=\varnothing\),
więc \(\operatorname{Ind}X\le0\).
Niepustość \(X\) wyklucza wymiar \(-1\),
co kończy dowód.
\end{proof}

\section{Lemat o liczbie Lebesgue'a i wymiar odcinka}
\label{sec:topog23-interval}

\begin{lemat}[Liczba Lebesgue'a pokrycia]
\label{lem:topog23-lebesgue-number}
Niech \((K,d)\) będzie niepustą
zwartą przestrzenią metryczną,
a \((U_i)_{i=1}^m\) jej skończonym
otwartym pokryciem.
Istnieje \(\delta>0\) takie,
że każdy niepusty podzbiór \(A\subseteq K\)
o średnicy mniejszej niż \(\delta\)
jest zawarty w pewnym \(U_i\).
\end{lemat}
\begin{proof}
Dla każdego \(x\in K\) wybieramy \(i(x)\)
oraz \(\varepsilon_x>0\) tak, że
\(B_d(x,2\varepsilon_x)\subseteq U_{i(x)}\).
Kule \(B_d(x,\varepsilon_x)\) pokrywają \(K\).
Ze zwartości istnieją punkty
\(x_1,\ldots,x_k\) takie, że
\(K=\bigcup_{j=1}^k B_d(x_j,\varepsilon_{x_j})\).
Przyjmijmy
\(\delta=\min_{1\le j\le k}\varepsilon_{x_j}>0\).

Jeśli \(A\neq\varnothing\) i
\(\operatorname{diam} A<\delta\),
wybierzmy \(a\in A\) i indeks \(j\)
z \(a\in B_d(x_j,\varepsilon_{x_j})\).
Dla każdego \(y\in A\)
\[
 d(y,x_j)\le d(y,a)+d(a,x_j)
 <\delta+\varepsilon_{x_j}
 \le2\varepsilon_{x_j}.
\]
Zatem \(A\subseteq U_{i(x_j)}\).
\end{proof}

\begin{twierdzenie}[Wymiar zwartych podzbiorów prostej]
\label{thm:topog23-real-upper}
Niech \(K\subseteq\mathbb R\) będzie zwartym
podzbiorem prostej z jej zwykłą topologią.
Wówczas \(\dim K\le1\),
przy konwencji \(\dim\varnothing=-1\).
\end{twierdzenie}
\begin{proof}
Dla pustego \(K\) teza jest oczywista.
Załóżmy \(K\neq\varnothing\)
i weźmy skończone otwarte pokrycie \((U_i)\)
przestrzeni \(K\).
Niech \(\delta>0\) będzie liczbą
z Lematu~\ref{lem:topog23-lebesgue-number}
dla metryki euklidesowej.
Wybierzmy \(h>0\) z \(2h<\delta\).
Dla każdego \(k\in\mathbb Z\) definiujemy
\[
 V_k=K\cap(kh-h,kh+h).
\]
Zbiory \(V_k\) są otwarte w \(K\) i pokrywają \(K\):
każda liczba rzeczywista leży w odległości
co najwyżej \(h/2\) od pewnej wielokrotności \(h\).
Ponieważ \(K\) jest ograniczony,
tylko skończenie wiele \(V_k\)
jest niepustych.

Średnica każdego \(V_k\)
jest mniejsza lub równa \(2h<\delta\),
więc niepuste \(V_k\) zawiera się w pewnym \(U_i\).
Ponadto każdy punkt należy
do co najwyżej dwóch zbiorów \(V_k\):
gdyby \(x\) należał do \(V_{k_1}\)
i \(V_{k_3}\) dla \(k_3\ge k_1+2\),
musielibyśmy mieć jednocześnie
\(|x-k_1h|<h\) i
\(|x-k_3h|<h\), co przez nierówność trójkąta
dawałoby \((k_3-k_1)h<2h\),
sprzeczność.
Pokrycie \((V_k)\) jest więc
skończonym otwartym rozdrobnieniem
rzędu co najwyżej \(2\).
To dowodzi \(\dim K\le1\).
\end{proof}

\begin{twierdzenie}[Wymiar odcinka]
\label{thm:topog23-interval}
Dla dowolnych \(a,b\in\mathbb R\) z \(a<b\)
zachodzi
\[
 \dim[a,b]=1.
\]
\end{twierdzenie}
\begin{proof}
Twierdzenie~\ref{thm:topog23-real-upper}
daje \(\dim[a,b]\le1\).
Odcinek jest zwartą przestrzenią Hausdorffa,
spójną i zawiera więcej niż jeden punkt.
Twierdzenie~\ref{thm:topog23-connected-positive}
wyklucza wymiar zero.
Otrzymujemy równość.
\end{proof}

\begin{wniosek}[Zwarte spójne podzbiory prostej]
\label{cor:topog23-connected-real}
Niech \(K\subseteq\mathbb R\) będzie niepusty,
zwarty i spójny.
Jeżeli \(K\) ma jeden punkt,
to \(\dim K=0\).
Jeżeli \(K\) zawiera przynajmniej dwa punkty,
to \(\dim K=1\).
\end{wniosek}
\begin{proof}
Jednopunktowa przestrzeń jest dyskretna,
więc ma wymiar zero.
Jeśli \(K\) ma co najmniej dwa punkty,
jego spójność i Twierdzenie
\ref{thm:topog23-connected-positive}
dają \(\dim K\ge1\).
Z Twierdzenia~\ref{thm:topog23-real-upper}
mamy \(\dim K\le1\).
\end{proof}

\begin{figure}[htbp]
\centering
\begin{tikzpicture}[x=1.02cm,y=.7cm]
 \draw[->] (-.5,0)--(7.4,0);
 \foreach \k in {0,...,6}{
  \draw[very thin,gray] (\k,-.12)--(\k,.14);
  \draw[blue!65!black,very thick] (\k-.87,.55)--(\k+.87,.55);
  \fill[blue!65!black] (\k-.87,.55) circle (1.4pt);
  \fill[blue!65!black] (\k+.87,.55) circle (1.4pt);
 }
 \node[anchor=west,font=\small] at (0,-.56)
 {Punkty prostej należą do najwyżej dwóch kolejnych otoczeń.};
\end{tikzpicture}
\caption{Idea rozdrobnienia pokrycia zwartego podzbioru
prostej do rzędu co najwyżej \(2\).}
\label{fig:topog23-line-cover}
\end{figure}

\section{Nerw pokrycia i odwzorowanie barycentryczne}
\label{sec:topog23-nerve}

\begin{definicja}[Nerw skończonego pokrycia]
\label{def:topog23-nerve}
Niech \(\mathcal U=(U_i)_{i=1}^{m}\)
będzie skończonym otwartym pokryciem \(X\).
\emph{Nerwem} \(\mathcal U\) nazywamy
abstrakcyjny kompleks symplicjalny,
którego wierzchołki odpowiadają indeksom
\(1,\ldots,m\), a niepusty zbiór indeksów
\(J\subseteq\{1,\ldots,m\}\)
wyznacza sympleks wtedy i tylko wtedy, gdy
\(\bigcap_{j\in J}U_j\neq\varnothing\).
Realizacja geometryczna nerwu
\(\lvert N(\mathcal U)\rvert\)
jest sumą odpowiednich ścian
standardowego sympleksu w \(\mathbb R^m\).
Jeśli rząd pokrycia jest co najwyżej \(n+1\),
realizacja nerwu ma wymiar symplicjalny
co najwyżej \(n\).
\end{definicja}

\begin{twierdzenie}[Odwzorowanie kanoniczne do nerwu]
\label{thm:topog23-canonical-nerve}
Niech \((X,d)\) będzie niepustą
zwartą przestrzenią metryczną,
a \(U_1,\ldots,U_m\) skończonym otwartym
pokryciem \(X\).
Istnieje ciągłe odwzorowanie
\(p\colon X\to\lvert N(\mathcal U)\rvert\)
o następującej własności:
dla każdego \(x\in X\) współrzędna
barycentryczna \(p_i(x)\) jest dodatnia
tylko wtedy, gdy \(x\in U_i\).
Jeśli rząd pokrycia jest co najwyżej \(n+1\),
obraz \(p(X)\) zawiera się w \(n\)-szkielecie nerwu.
\end{twierdzenie}
\begin{proof}
Dla każdego \(i\) określmy
ciągłą funkcję \(\psi_i\colon X\to[0,\infty)\)
w następujący sposób:
\[
 \psi_i(x)=
 \begin{cases}
   1,&U_i=X,\\
   d(x,X\setminus U_i),&U_i\neq X.
 \end{cases}
\]
Jeśli \(U_i\neq X\), dopełnienie
\(X\setminus U_i\) jest niepuste i domknięte,
więc funkcja odległości jest ciągła.
W obu przypadkach
\(\psi_i(x)>0\) dokładnie wtedy, gdy \(x\in U_i\).
Ponieważ \((U_i)\) pokrywa \(X\),
suma \(S(x)=\sum_i\psi_i(x)\) jest dodatnia
dla każdego \(x\in X\).
Przyjmijmy
\[
 p_i(x)=\frac{\psi_i(x)}{S(x)},
 \qquad
 p(x)=\sum_{i=1}^m p_i(x)e_i\in\mathbb R^m.
\]
Wszystkie współrzędne są ciągłe,
nieujemne i sumują się do \(1\).
Dodatnie współrzędne występują dokładnie
dla tych \(i\), dla których \(x\in U_i\).
Ich indeksy tworzą sympleks nerwu,
gdyż ich odpowiednie zbiory
mają wspólny punkt \(x\).
Stąd \(p(X)\subseteq\lvert N(\mathcal U)\rvert\).

Jeżeli rząd pokrycia nie przekracza \(n+1\),
dla każdego \(x\) co najwyżej \(n+1\)
współrzędnych \(p_i(x)\) jest dodatnich.
Tym samym \(p(x)\) leży na ścianie
o wymiarze nie większym niż \(n\).
\end{proof}

\begin{uwaga}[Związek z topologią algebraiczną]
\label{rem:topog23-nerve-overlap}
Nerw pokrycia występuje również
w topologii algebraicznej, gdzie
jego realizację bada się za pomocą
homologii, homotopii i twierdzenia o nerwie.
Tutaj potrzebujemy wyłącznie
konstrukcji odwzorowania barycentrycznego
i oszacowania liczby niezerowych współrzędnych.
Nie powtarzamy teorii homologii nerwu.
\end{uwaga}

\section{Górne oszacowanie wymiaru w przestrzeni euklidesowej}
\label{sec:topog23-euclidean}

\begin{lemat}[Drobna triangulacja przestrzeni euklidesowej]
\label{lem:topog23-grid-triangulation}
Niech \(n\ge1\) i \(h>0\).
Istnieje lokalnie skończona triangulacja
\(\mathbb R^n\), której wierzchołkami
są punkty siatki \(h\mathbb Z^n\),
a wszystkie sympleksy leżą w sześcianach
\[
 \prod_{j=1}^{n}[k_jh,(k_j+1)h],
 \qquad (k_1,\ldots,k_n)\in\mathbb Z^n.
\]
W tej triangulacji otwarte gwiazdy wierzchołków
tworzą pokrycie \(\mathbb R^n\)
rzędu co najwyżej \(n+1\).
Gwiazda wierzchołka \(v\)
zawiera się w \(v+(-h,h)^n\),
zatem ma średnicę nie większą niż \(2h\sqrt n\).
\end{lemat}
\begin{proof}
Rozważmy najpierw jednostkowy sześcian
\([0,1]^n\).
Dla każdej permutacji \(\sigma\)
zbioru \(\{1,\ldots,n\}\)
weźmy sympleks o wierzchołkach
\[
 0,\quad e_{\sigma(1)},\quad
 e_{\sigma(1)}+e_{\sigma(2)},\quad\ldots,\quad
 e_{\sigma(1)}+\cdots+e_{\sigma(n)}.
\]
Sympleks ten składa się z punktów,
których współrzędne w kolejności
\(\sigma(1),\ldots,\sigma(n)\)
są nierosnące.
Każdy punkt sześcianu spełnia
takie nierówności dla pewnej permutacji,
więc sympleksy pokrywają sześcian.
Dwa z nich przecinają się w ścianie
wyznaczonej przez równości
odpowiednich sąsiednich współrzędnych.
Te same podziały są zgodne na wspólnych
ścianach sąsiednich sześcianów.

Przesuwając i skalując konstrukcję
na sześciany siatki \(h\mathbb Z^n\),
otrzymujemy żądaną triangulację.
Jest lokalnie skończona, ponieważ
każdy zwarty zbiór spotyka tylko
skończenie wiele sześcianów siatki.
Dla wierzchołka \(v\)
jego \emph{otwarta gwiazda}
to suma wnętrz względnych sympleksów
mających \(v\) jako wierzchołek.
Równoważnie, jest to zbiór punktów,
dla których globalna, kawałkami afiniczna
współrzędna barycentryczna przy \(v\)
jest dodatnia. Współrzędna ta jest ciągła
dzięki zgodności triangulacji na ścianach,
więc gwiazda jest otwarta w \(\mathbb R^n\).
Każdy punkt należy do pewnej gwiazdy,
bo suma współrzędnych barycentrycznych
w odpowiednim sympleksie wynosi \(1\).

Dla punktu należącego do względnego wnętrza
sympleksu \(\sigma\)
dodatnie współrzędne barycentryczne
dotyczą wyłącznie jego wierzchołków.
Każdy sympleks ma co najwyżej \(n+1\)
wierzchołków, a więc punkt należy
do co najwyżej \(n+1\) otwartych gwiazd.
Każdy sympleks zawierający wierzchołek \(v\)
leży w jednym z sześcianów siatki
mających \(v\) jako wierzchołek.
Stąd gwiazda \(v\) zawiera się w
\(v+(-h,h)^n\), co daje oszacowanie średnicy.
\end{proof}

\begin{twierdzenie}[Górne ograniczenie wymiaru euklidesowego]
\label{thm:topog23-euclidean-upper}
Niech \(n\ge1\) oraz
\(K\subseteq\mathbb R^n\)
będzie zwartym podzbiorem.
Wówczas \(\dim K\le n\).
\end{twierdzenie}
\begin{proof}
Dla \(K=\varnothing\) teza wynika
z konwencji \(\dim\varnothing=-1\).
Dla niepustego \(K\)
weźmy dowolne skończone otwarte pokrycie
\((U_i)_{i=1}^m\) przestrzeni \(K\).
Niech \(\delta>0\) będzie jego liczbą Lebesgue'a
z Lematu~\ref{lem:topog23-lebesgue-number}.
Wybierzmy \(h>0\) takie, że
\(2h\sqrt n<\delta\),
i rozważmy triangulację z
Lematu~\ref{lem:topog23-grid-triangulation}.
Dla wierzchołków \(v\in h\mathbb Z^n\)
niech \(S_v\) będzie otwartą gwiazdą \(v\).

Rodzina \((K\cap S_v)\)
pokrywa \(K\) zbiorami otwartymi w \(K\)
i ma rząd co najwyżej \(n+1\).
Niepusty zbiór \(K\cap S_v\)
ma średnicę co najwyżej
\(2h\sqrt n<\delta\),
więc zawiera się w pewnym \(U_i\).
Ponieważ \(K\) jest ograniczony,
tylko skończenie wiele gwiazd spotyka \(K\):
jeśli \(S_v\cap K\neq\varnothing\),
to \(v\) leży w odległości
co najwyżej \(h\sqrt n\)
od pewnego punktu \(K\),
a w ustalonym ograniczonym
podzbiorze \(\mathbb R^n\)
jest tylko skończenie wiele punktów
siatki \(h\mathbb Z^n\).
Po odrzuceniu pustych przecięć
otrzymujemy skończone otwarte rozdrobnienie
wyjściowego pokrycia rzędu co najwyżej \(n+1\).
Zatem \(\dim K\le n\).
\end{proof}

\begin{uwaga}[Granice udowodnionych wyników]
\label{rem:topog23-euclidean-lower}
Powyższe twierdzenie dowodzi nierówności
\(\dim K\le n\), a nie równości dla wszystkich
zwartych \(K\subseteq\mathbb R^n\):
na przykład zbiór Cantora ma wymiar zero,
mimo że jest podzbiorem prostej.
Dla sześcianu \([0,1]^n\)
klasyczna równość \(\dim[0,1]^n=n\)
wymaga również dolnego oszacowania.
Wymaga ono dodatkowego twierdzenia
o pokryciach sześcianu lub równoważnego
argumentu topologicznego,
którego nie utożsamiamy
z udowodnioną tutaj nierównością górną.
Taki dowód powinien zostać umieszczony
w osobnej kontynuacji teorii wymiaru,
po sprawdzeniu odpowiednich wyników
w skrypcie topologii algebraicznej.
\end{uwaga}

\section{Obrazy ciągłe nie zachowują wymiaru}
\label{sec:topog23-images}

\begin{twierdzenie}[Obraz zbioru Cantora może mieć większy wymiar]
\label{thm:topog23-cantor-surjection}
Istnieje ciągłe odwzorowanie
\(f\colon C\to[0,1]\) zbioru Cantora
na cały odcinek.
W szczególności
\[
 \dim C=0,\qquad \dim f(C)=1.
\]
Zatem wymiar pokryciowy nie musi
maleć przy dowolnym ciągłym
odwzorowaniu surjektywnym.
\end{twierdzenie}
\begin{proof}
Z Przykładu~\ref{ex:topog23-cantor}
mamy homeomorfizm
\[
 c\colon\{0,1\}^{\mathbb N}\to C,\qquad
 c((\varepsilon_j))=\sum_{j=1}^{\infty}
                       2\varepsilon_j3^{-j}.
\]
Określmy odwzorowanie
\[
 b\colon\{0,1\}^{\mathbb N}\to[0,1],
 \qquad b((\varepsilon_j))=
             \sum_{j=1}^{\infty}\varepsilon_j2^{-j}.
\]
Jest ono ciągłe: suma pierwszych \(m\)
wyrazów zależy od skończenie wielu
współrzędnych, a ogon szeregu
ma moduł co najwyżej \(2^{-m}\),
jednostajnie względem wszystkich ciągów.
Jest także surjektywne,
ponieważ każda liczba z \([0,1]\)
ma co najmniej jedno rozwinięcie dwójkowe
o cyfrach \(0,1\); w przypadku \(1\)
wybieramy ciąg samych jedynek.
Wystarczy zatem przyjąć
\(f=b\circ c^{-1}\).
Ciągłość i surjektywność wynikają
z własności \(b\) oraz \(c\).
Równości wymiarów wynikają
z Przykładu~\ref{ex:topog23-cantor}
i Twierdzenia~\ref{thm:topog23-interval}.
\end{proof}

\section{Podsumowanie}
\label{sec:topog23-summary}

Wymiar pokryciowy mierzy możliwość
rozdrobnienia pokrycia otwartego tak,
aby w jednym punkcie spotykała się
ograniczona liczba jego elementów.
Dla zwartych przestrzeni Hausdorffa
wymiar zero jest równoważny istnieniu
bazy zbiorów otwarto-domkniętych.
W tym szczególnym wymiarze
otrzymaliśmy również zgodność
z małym i dużym wymiarem indukcyjnym.

Dowiedliśmy, że zbiór Cantora ma wymiar zero,
każdy nietrywialny zwarty zbiór
spójny ma wymiar co najmniej jeden,
odcinek ma wymiar jeden,
a każdy zwarty podzbiór \(\mathbb R^n\)
ma wymiar co najwyżej \(n\).
Nerw pokrycia wiąże ograniczenie rzędu
pokrycia z wymiarem kompleksu symplicjalnego.
Wskazaliśmy zarazem granicę obecnego
opracowania: niższe oszacowania wymiaru
sześcianów wyższych wymiarów
i pełne porównanie trzech wymiarów
wymagają dodatkowych wyników;
nie deklarujemy ich dowodów
na podstawie samych definicji.
\chapter{Grupy topologiczne i przestrzenie jednorodne}
\label{ch:topog24-grupy}

\noindent\textbf{Cel rozdziału.}
Działania algebraiczne mogą współistnieć z topologią na tym samym
zbiorze. W grupie topologicznej mnożenie i odwracanie są ciągłe,
a otoczenia elementu neutralnego opisują całą strukturę lokalną.
Zbadamy wynikające stąd własności oddzielania, podgrupy, przestrzenie
warstw, ilorazy grupowe, składowe spójne i struktury jednostajne.
Definicji grup oraz ich podstawowych własności algebraicznych nie
powtarzamy ze skryptu \emph{Algebra abstrakcyjna};
konstrukcja i jednoznaczność miary Haara na grupach lokalnie
zwartych należą do rozdziału 23 skryptu \emph{Teoria miary}.

\section{Grupa topologiczna i działania przez translacje}
\label{sec:topog24-definition}

\begin{definicja}[Grupa topologiczna Hausdorffa]
\label{def:topog24-group}
Niech \(G\) będzie grupą z elementem neutralnym \(e\)
i topologią \(\tau\). Nazywamy \((G,\tau)\)
\emph{grupą topologiczną Hausdorffa}, jeżeli \(G\)
jest przestrzenią Hausdorffa i odwzorowania
\[
 m\colon G\times G\to G,\quad m(x,y)=xy,
 \qquad
 \iota\colon G\to G,\quad \iota(x)=x^{-1}
\]
są ciągłe, przy czym \(G\times G\) ma topologię
produktową. Dalej określenie \emph{grupa topologiczna}
oznacza grupę topologiczną Hausdorffa, chyba że
wyraźnie zaznaczono inaczej.
Dla \(A,B\subseteq G\) przyjmujemy
\(AB=\{ab:a\in A,\ b\in B\}\),
\(A^{-1}=\{a^{-1}:a\in A\}\).
\end{definicja}

\begin{uwaga}[Dlaczego oddzielamy oba warunki ciągłości?]
\label{rem:topog24-continuity}
Sama ciągłość mnożenia nie musi implikować
ciągłości odwracania. Grupa z topologią, dla której
ciągłe jest jedynie mnożenie, jest nazywana
\emph{grupą paratopologiczną}.
Ponadto aksjomat Hausdorffa nie wynika z samej
ciągłości obu działań: na dowolnej grupie
topologia antydyskretna czyni oba odwzorowania
ciągłymi, lecz dla grupy wieloelementowej
nie jest Hausdorffa. Stąd jawna konwencja
w Definicji~\ref{def:topog24-group}.
\end{uwaga}

\begin{twierdzenie}[Jedno kryterium ciągłości działań]
\label{thm:topog24-difference}
Niech \(G\) będzie grupą z dowolną topologią.
Następujące warunki są równoważne:
\begin{enumerate}[label=\textup{(\roman*)}]
\item mnożenie \(G\times G\to G\) i odwracanie
      \(G\to G\) są ciągłe;
\item odwzorowanie \(D\colon G\times G\to G\)
      zadane wzorem \(D(x,y)=xy^{-1}\) jest ciągłe.
\end{enumerate}
W tej równoważności nie potrzeba aksjomatu Hausdorffa.
\end{twierdzenie}
\begin{proof}
Jeśli oba działania są ciągłe, to
\(D=m\circ(\mathrm{id}_G\times\iota)\)
jest ciągłe. Odwrotnie, przyjmując
\(\iota(y)=D(e,y)\), widzimy, że odwracanie
jest ciągłe jako złożenie \(y\mapsto(e,y)\)
z \(D\). Wówczas
\(m(x,y)=D(x,\iota(y))\),
więc mnożenie jest również ciągłe.
\end{proof}

\begin{stwierdzenie}[Translacje i sprzężenia są homeomorfizmami]
\label{prop:topog24-translations}
Niech \(G\) będzie grupą topologiczną,
a \(g\in G\). Odwzorowania
\(L_g(x)=gx\), \(R_g(x)=xg\) i
\(C_g(x)=gxg^{-1}\) są homeomorfizmami \(G\).
Odwracanie \(x\mapsto x^{-1}\) także
jest homeomorfizmem.
\end{stwierdzenie}
\begin{proof}
Każda translacja jest ciągła jako
ograniczenie ciągłego mnożenia do
ustalonego czynnika. Odwrotnościami
\(L_g,R_g,C_g\) są odpowiednio
\(L_{g^{-1}},R_{g^{-1}},C_{g^{-1}}\);
wszystkie są ciągłe. Odwracanie
jest ciągłą inwolucją, a więc
jest swoim ciągłym odwzorowaniem odwrotnym.
\end{proof}

\begin{wniosek}[Jednorodność]
\label{cor:topog24-homogeneous}
Dla dowolnych \(x,y\in G\) istnieje
homeomorfizm \(h\colon G\to G\) z \(h(x)=y\).
W szczególności każda własność lokalna
zachowywana przez homeomorfizmy,
np. zwartość lokalna czy spełnianie
pierwszego aksjomatu przeliczalności,
zachodzi w każdym punkcie \(G\) dokładnie wtedy,
gdy zachodzi w \(e\).
\end{wniosek}
\begin{proof}
Wystarczy wziąć \(h=L_{yx^{-1}}\).
Wtedy \(h(x)=yx^{-1}x=y\).
Każde otoczenie \(x\) jest
przesunięciem pewnego otoczenia \(e\),
więc zachowanie rozważanych własności
wynika z niezmienniczości homeomorficznej.
\end{proof}

\section{Otoczenia elementu neutralnego i regularność}
\label{sec:topog24-neighborhoods}

\begin{lemat}[Symetryczne otoczenie o małym iloczynie]
\label{lem:topog24-small-neighborhoods}
Niech \(G\) będzie grupą topologiczną.
Dla każdego otwartego otoczenia
\(U\subseteq G\) punktu \(e\) istnieje
otwarte otoczenie \(V\) punktu \(e\)
takie, że \(V=V^{-1}\) oraz \(VV\subseteq U\).
Co więcej, dla każdej liczby naturalnej
\(k\ge1\) można dobrać takie \(V\)
z \(V^k\subseteq U\).
\end{lemat}
\begin{proof}
Z ciągłości \(m\) w \((e,e)\)
wynika istnienie otwartych otoczeń
\(W_1,W_2\) elementu \(e\), takich że
\(W_1W_2\subseteq U\).
Zbiór
\(V=W_1\cap W_2\cap W_1^{-1}\cap W_2^{-1}\)
jest otwarty, symetryczny i zawiera \(e\).
Ponieważ \(V\subseteq W_1\cap W_2\),
otrzymujemy \(VV\subseteq U\).

Dla \(k\ge1\) wybieramy kolejno
symetryczne otwarte otoczenia
\(U=V_0\supseteq V_1\supseteq\cdots\supseteq V_r\)
z \(V_{j+1}^2\subseteq V_j\).
Wybierając \(2^r\ge k\),
otrzymujemy \(V_r^k\subseteq
V_r^{2^r}\subseteq U\), ponieważ \(e\in V_r\).
\end{proof}

\begin{lemat}[Zamykanie małych otoczeń]
\label{lem:topog24-closure}
Niech \(G\) będzie grupą topologiczną
i \(V\) otwartym, symetrycznym otoczeniem \(e\).
Wówczas \(\overline V^{\,G}\subseteq VV\).
Dla każdego otwartego otoczenia \(U\ni e\)
istnieje więc otwarte, symetryczne \(V\ni e\)
z \(\overline V^{\,G}\subseteq U\).
\end{lemat}
\begin{proof}
Niech \(x\in\overline V^{\,G}\).
Zbiór \(xV\) jest otwartym otoczeniem \(x\),
więc \(xV\cap V\neq\varnothing\).
Istnieją zatem \(v,w\in V\)
z \(xv=w\), czyli \(x=wv^{-1}\in VV\),
ponieważ \(V=V^{-1}\).
Jeśli teraz \(U\ni e\) jest otwarty,
wybieramy \(V\) z poprzedniego lematu
tak, aby \(VV\subseteq U\).
Wtedy \(\overline V^{\,G}\subseteq U\).
\end{proof}

\begin{twierdzenie}[Regularność grup topologicznych]
\label{thm:topog24-regular}
Każda grupa topologiczna Hausdorffa
jest przestrzenią regularną spełniającą \(T_1\).
Dokładniej, dla każdego \(x\in G\)
i każdego otwartego otoczenia \(O\ni x\)
istnieje otwarty zbiór \(W\) taki, że
\[
 x\in W\subseteq\overline W^{\,G}\subseteq O.
\]
\end{twierdzenie}
\begin{proof}
Zbiór \(x^{-1}O\) jest otwartym
otoczeniem \(e\). Na mocy
Lematu~\ref{lem:topog24-closure}
istnieje otwarte \(V\ni e\)
z \(\overline V^{\,G}\subseteq x^{-1}O\).
Bierzemy \(W=xV\).
Ponieważ \(L_x\) jest homeomorfizmem,
\(\overline{xV}^{\,G}=x\overline V^{\,G}\).
Stąd \(x\in W\) i
\(\overline W^{\,G}\subseteq O\).
Aksjomat \(T_1\) wynika
z aksjomatu Hausdorffa.
\end{proof}

\begin{stwierdzenie}[Ciągłość homomorfizmu wystarczy sprawdzać w jedynce]
\label{prop:topog24-hom-continuity}
Niech \(G,H\) będą grupami topologicznymi,
a \(f\colon G\to H\) homomorfizmem grupowym.
Wówczas \(f\) jest ciągły na całym \(G\)
dokładnie wtedy, gdy jest ciągły w \(e_G\).
\end{stwierdzenie}
\begin{proof}
Kierunek konieczny jest oczywisty.
Załóżmy ciągłość w \(e_G\).
Niech \(x\in G\), a \(O\)
będzie otwartym otoczeniem \(f(x)\) w \(H\).
Zbiór \(f(x)^{-1}O\) jest
otwartym otoczeniem \(e_H\).
Istnieje więc otwarte \(V\ni e_G\)
takie, że \(f(V)\subseteq f(x)^{-1}O\).
Wtedy \(xV\) jest otwartym otoczeniem \(x\),
a \(f(xV)=f(x)f(V)\subseteq O\).
To dowodzi ciągłości w każdym \(x\).
\end{proof}

\begin{uwaga}[Ciągłość i jednostajna ciągłość]
\label{rem:topog24-hom-uniform}
Powyższe twierdzenie dotyczy
homomorfizmów, a nie dowolnych funkcji.
Homomorfizmy ciągłe są także
jednostajnie ciągłe względem
naturalnych lewych struktur jednostajnych,
które zdefiniujemy w ostatniej części rozdziału.
Nie wynika z tego, że mnożenie na
\(G\times G\) jest zawsze jednostajnie
ciągłe dla takich struktur.
\end{uwaga}

\section{Podgrupy, domknięcia i centralizatory}
\label{sec:topog24-subgroups}

\begin{twierdzenie}[Podgrupy i ich domknięcia]
\label{thm:topog24-subgroup-closure}
Niech \(G\) będzie grupą topologiczną
oraz \(H\le G\) podgrupą.
Wtedy \(H\) z topologią podprzestrzeni
jest grupą topologiczną,
a domknięcie \(\overline H^{\,G}\)
jest podgrupą domkniętą \(G\).
Jeżeli \(H\) jest normalna w \(G\),
to \(\overline H^{\,G}\) również jest normalna.
\end{twierdzenie}
\begin{proof}
Ograniczenia mnożenia i odwracania
do \(H\times H\) oraz \(H\)
są ciągłe, a podprzestrzeń
przestrzeni Hausdorffa jest Hausdorffa.

Niech \(a,b\in\overline H^{\,G}\).
Dla dowolnego otwartego otoczenia
\(O\ni ab^{-1}\)
ciągłość odwzorowania
\(D(x,y)=xy^{-1}\)
zapewnia otwarte otoczenia
\(U\ni a\), \(V\ni b\),
dla których \(UV^{-1}\subseteq O\).
Z definicji domknięcia
wybieramy \(h\in H\cap U\)
i \(k\in H\cap V\).
Wtedy \(hk^{-1}\in H\cap O\).
Każde otwarte otoczenie \(ab^{-1}\)
spotyka \(H\), więc
\(ab^{-1}\in\overline H^{\,G}\).
Ponieważ \(H\neq\varnothing\),
to kryterium podgrupy
dowodzi, że \(\overline H^{\,G}\le G\).

Jeśli \(H\) jest normalna,
to dla każdego \(g\in G\)
sprzężenie \(C_g\) jest homeomorfizmem,
zatem
\(g\overline H^{\,G}g^{-1}
=\overline{gHg^{-1}}^{\,G}
=\overline H^{\,G}\).
\end{proof}

\begin{stwierdzenie}[Podgrupa otwarta jest domknięta]
\label{prop:topog24-open-closed}
Jeżeli \(H\le G\) jest podgrupą
otwartą w grupie topologicznej \(G\),
to jest również domknięta.
W szczególności jej wszystkie lewe
i prawe warstwy są otwarto-domknięte.
\end{stwierdzenie}
\begin{proof}
Dla każdego \(g\in G\) lewa
warstwa \(gH\) jest otwarta,
bo jest obrazem \(H\)
przy homeomorfizmie \(L_g\).
Warstwy lewe są parami rozłączne
albo równe, a ich suma daje \(G\).
Dopełnienie \(H\) jest więc sumą
pozostałych otwartych warstw,
zatem jest otwarte.
Stąd \(H\) jest domknięta.
Argument dla prawych warstw
jest analogiczny.
\end{proof}

\begin{stwierdzenie}[Centrum i centralizatory]
\label{prop:topog24-centralizers}
Niech \(G\) będzie grupą topologiczną
oraz \(A\subseteq G\).
Centralizator
\[
 C_G(A)=\{x\in G:xa=ax
              \text{ dla każdego }a\in A\}
\]
jest podgrupą domkniętą.
W szczególności centrum
\(Z(G)=C_G(G)\)
jest domkniętą podgrupą normalną.
\end{stwierdzenie}
\begin{proof}
Algebraiczne własności centralizatora
i centrum wynikają z praw grupowych.
Dla ustalonego \(a\in A\)
zbiór \(\{x:xa=ax\}\)
jest domknięty, ponieważ
odwzorowania \(x\mapsto xa\)
oraz \(x\mapsto ax\)
są ciągłe, a przekątna w
\(G\times G\) jest domknięta
dzięki własności Hausdorffa.
Centralizator jest przecięciem
tych domkniętych zbiorów po \(a\in A\).
Centrum jest ponadto normalne,
ponieważ jego elementy komutują
ze wszystkimi elementami \(G\).
\end{proof}

\begin{przyklad}[Podgrupa gęsta i niedomknięta]
\label{ex:topog24-rational-subgroup}
Grupa \((\mathbb R,+)\) z topologią euklidesową
zawiera podgrupę \((\mathbb Q,+)\).
Jest ona gęsta w \(\mathbb R\):
każdy niepusty przedział
zawiera liczbę wymierną.
Nie jest domknięta, gdyż
\(\overline{\mathbb Q}^{\,\mathbb R}=\mathbb R\)
i \(\mathbb Q\neq\mathbb R\).
Iloraz \(\mathbb R/\mathbb Q\)
z topologią ilorazową
nie będzie Hausdorffa;
dokładne kryterium udowodnimy poniżej.
\end{przyklad}

\section{Przestrzenie warstw i grupy ilorazowe}
\label{sec:topog24-quotients}

\begin{definicja}[Przestrzeń jednorodna wyznaczona podgrupą]
\label{def:topog24-cosetspace}
Niech \(G\) będzie grupą topologiczną
i \(H\le G\). Przez \(G/H\)
oznaczamy zbiór lewych warstw
\(gH=\{gh:h\in H\}\),
wyposażony w topologię ilorazową
względem kanonicznej surjekcji
\[
 q\colon G\to G/H,\qquad q(g)=gH.
\]
Nie zakładamy, że \(H\) jest normalna.
Przestrzeń \(G/H\) jest więc
na razie przestrzenią topologiczną,
niekoniecznie grupą.
\end{definicja}

\begin{twierdzenie}[Otwartość odwzorowania na przestrzeń warstw]
\label{thm:topog24-open-quotient}
Dla każdej podgrupy \(H\le G\)
kanoniczne odwzorowanie
\(q\colon G\to G/H\) jest ciągłe,
surjektywne i otwarte.
\end{twierdzenie}
\begin{proof}
Ciągłość i surjektywność
należą do definicji ilorazu.
Niech \(U\subseteq G\) będzie otwarty.
Mamy równość
\[
 q^{-1}(q(U))=UH
  =\bigcup_{h\in H}Uh.
\]
Każdy \(Uh\) jest otwarty,
gdyż prawa translacja przez \(h\)
jest homeomorfizmem.
Zatem \(UH\) jest otwarty w \(G\),
a z definicji topologii ilorazowej
\(q(U)\) jest otwarty w \(G/H\).
\end{proof}

\begin{twierdzenie}[Ciągłe i przechodnie działanie na przestrzeni warstw]
\label{thm:topog24-coset-action}
Niech \(G\) będzie grupą topologiczną,
\(H\le G\) dowolną podgrupą,
a \(G/H\) przestrzenią warstw z topologią ilorazową.
Wówczas wzór
\[
 \alpha\colon G\times G/H\to G/H,\qquad
 \alpha(g,xH)=(gx)H
\]
określa ciągłe działanie grupy \(G\)
na przestrzeni \(G/H\).
Działanie to jest przechodnie,
a stabilizatorem neutralnej warstwy \(H\)
jest dokładnie podgrupa \(H\).
\end{twierdzenie}
\begin{proof}
Jeżeli \(xH=yH\), to istnieje
\(h\in H\) takie, że \(y=xh\);
wobec tego \((gy)H=(gxh)H=(gx)H\).
Wzór na \(\alpha\) nie zależy
zatem od reprezentanta warstwy.
Prawa grupowe dają
\(\alpha(e,xH)=xH\) oraz
\(\alpha(g,\alpha(k,xH))=\alpha(gk,xH)\).

Z Twierdzenia~\ref{thm:topog24-open-quotient}
odwzorowanie \(q\colon G\to G/H\)
jest ciągłą otwartą surjekcją.
Zatem także
\(\mathrm{id}_G\times q\colon
G\times G\to G\times G/H\)
jest ciągłą otwartą surjekcją,
a więc odwzorowaniem ilorazowym.
Zachodzi równość
\[
 \alpha\circ(\mathrm{id}_G\times q)
 =q\circ m_G,
\]
gdzie \(m_G(g,x)=gx\) jest ciągłe.
Własność uniwersalna
odwzorowania ilorazowego
dowodzi ciągłości \(\alpha\).
Przechodniość wynika z równości
\(\alpha(yx^{-1},xH)=yH\).
Warunek \(\alpha(g,H)=H\)
jest równoważny \(gH=H\),
czyli \(g\in H\).
\end{proof}

\begin{twierdzenie}[Hausdorffowość przestrzeni warstw]
\label{thm:topog24-coset-hausdorff}
Niech \(G\) będzie grupą topologiczną,
a \(H\le G\).
Przestrzeń \(G/H\) jest Hausdorffa
wtedy i tylko wtedy, gdy \(H\)
jest domknięta w \(G\).
\end{twierdzenie}
\begin{proof}
Jeśli \(G/H\) jest Hausdorffa,
singleton \(\{H\}\) jest domknięty
w \(G/H\). Jego przeciwobraz
\(q^{-1}(\{H\})=H\)
jest domknięty w \(G\).

Odwrotnie, niech \(H\) będzie domknięta
i wybierzmy różne warstwy
\(xH\neq yH\). Wtedy
\(g=x^{-1}y\notin H\).
Z ciągłości odwzorowania
\[
 (u,v)\longmapsto u^{-1}gv
\]
w punkcie \((e,e)\)
i otwartości \(G\setminus H\)
istnieje otwarte otoczenie \(V\ni e\)
takie, że
\(V^{-1}gV\subseteq G\setminus H\).
Zbiory \(q(xV)\) i \(q(yV)\)
są otwartymi otoczeniami
odpowiednio \(xH\) i \(yH\)
na mocy poprzedniego twierdzenia.
Gdyby miały wspólny punkt,
istniałyby \(v,w\in V\)
i \(h_1,h_2\in H\) z
\(xvh_1=ywh_2\).
Otrzymalibyśmy
\(v^{-1}gw=h_1h_2^{-1}\in H\),
wbrew \(V^{-1}gV\cap H=\varnothing\).
Otoczenia te są więc rozłączne,
co dowodzi własności Hausdorffa.
\end{proof}

\begin{twierdzenie}[Grupa ilorazowa przez podgrupę normalną]
\label{thm:topog24-normal-quotient}
Niech \(G\) będzie grupą topologiczną,
a \(N\triangleleft G\) podgrupą normalną.
Na zbiorze \(G/N\) działania
\[
 (gN)(hN)=ghN,\qquad (gN)^{-1}=g^{-1}N
\]
wyznaczają grupę z topologią ilorazową,
dla której mnożenie i odwracanie
są ciągłe. Ponadto następujące
warunki są równoważne:
\begin{enumerate}[label=\textup{(\roman*)}]
\item \(N\) jest domknięta w \(G\);
\item topologia ilorazowa na \(G/N\) spełnia \(T_1\);
\item \(G/N\) jest grupą topologiczną Hausdorffa.
\end{enumerate}
\end{twierdzenie}
\begin{proof}
Poprawność algebraicznych działań
wynika z normalności \(N\).
Niech \(q\colon G\to G/N\)
będzie odwzorowaniem ilorazowym.
Jest ono otwarte na mocy
Twierdzenia~\ref{thm:topog24-open-quotient}.
Iloczyn \(q\times q\colon
G\times G\to(G/N)\times(G/N)\)
jest ciągłą otwartą surjekcją:
obraz podstawowego prostokąta
\(U\times V\) jest prostokątem
\(q(U)\times q(V)\), otwartym,
a zbiory otwarte produktu są sumami
takich prostokątów. Wobec tego
\(q\times q\) jest odwzorowaniem ilorazowym.

Mnożenie \(m_N\) w ilorazie
spełnia \(m_N\circ(q\times q)=q\circ m_G\).
Prawa strona jest ciągła,
więc własność uniwersalna odwzorowania
ilorazowego daje ciągłość \(m_N\).
Podobnie \(\iota_N\circ q=q\circ\iota_G\),
a ciągłość prawej strony
i ilorazowość \(q\) zapewniają
ciągłość \(\iota_N\).

Z definicji ilorazu
singleton \(\{N\}\) jest domknięty
w \(G/N\) dokładnie wtedy,
gdy \(N=q^{-1}(\{N\})\)
jest domknięta w \(G\).
W każdej grupie z ciągłymi działaniami
domkniętość singletonu neutralnego
implikuje \(T_1\), ponieważ
wszystkie translacje są homeomorfizmami.
Każda taka grupa spełniająca \(T_1\)
jest także Hausdorffa:
dla różnych \(a,b\) wybieramy
otwarte \(U\ni e\) z \(a^{-1}b\notin U\),
następnie symetryczne otwarte \(V\ni e\)
z \(VV\subseteq U\).
Zbiory \(aV,bV\) są rozłączne,
gdyż punkt wspólny dawałby
\(a^{-1}b\in VV^{-1}\subseteq U\).
To dowodzi równoważności
\textup{(i)}--\textup{(iii)}.
Alternatywnie \textup{(i)} implikuje
\textup{(iii)} wprost z
Twierdzenia~\ref{thm:topog24-coset-hausdorff}.
\end{proof}

\begin{przyklad}[Warstwy nie muszą tworzyć grupy]
\label{ex:topog24-nonnormal-s3}
Niech \(G=S_3\) będzie grupą wszystkich permutacji
trójelementowego zbioru z topologią dyskretną
i \(H=\{e,(1\,2)\}\).
Jest to podgrupa domknięta, ale nie jest normalna:
dla \(\sigma=(1\,2\,3)\) mamy
\(\sigma(1\,2)\sigma^{-1}=(2\,3)\notin H\).
Przestrzeń \(G/H\) ma trzy punkty i topologię dyskretną,
gdyż przeciwobrazy wszystkich jej podzbiorów
są otwarte w dyskretnym \(G\).

Nie można jednak definiować mnożenia warstw
wzorem \((gH)(kH)=gkH\):
mamy \(eH=(1\,2)H\), lecz
\[
 e(1\,3)H\neq(1\,2)(1\,3)H.
\]
Gdyby obie warstwy były równe, mielibyśmy
\((1\,3)^{-1}(1\,2)(1\,3)\in H\),
podczas gdy sprzężenie to jest \((2\,3)\notin H\).
Przykład rozróżnia przestrzeń warstw od grupy ilorazowej.
\end{przyklad}

\begin{twierdzenie}[Pierwsze twierdzenie o izomorfizmie topologicznym]
\label{thm:topog24-isomorphism}
Niech \(f\colon G\to H\) będzie
ciągłym, otwartym i surjektywnym
homomorfizmem grup topologicznych.
Wówczas \(N=\ker f\) jest
domkniętą podgrupą normalną \(G\),
a indukowany homomorfizm
\[
 \widetilde f\colon G/N\to H,\qquad
 \widetilde f(gN)=f(g)
\]
jest izomorfizmem grup topologicznych,
tj. homeomorfizmem i izomorfizmem
algebraicznym.
\end{twierdzenie}
\begin{proof}
Jądro jest podgrupą normalną
z własności algebraicznych homomorfizmu.
Jest domknięte jako
\(f^{-1}(\{e_H\})\), gdyż \(H\)
jest Hausdorffa.
Wzór na \(\widetilde f\)
jest dobrze określony: jeśli \(gN=hN\),
to \(h^{-1}g\in N\), więc \(f(g)=f(h)\).
Odwzorowanie \(\widetilde f\)
jest homomorfizmem, jest surjektywne,
a równość \(f(g)=f(h)\) daje
\(h^{-1}g\in N\), więc jest także iniektywne.

Ponieważ \(f=\widetilde f\circ q\)
i \(q\) jest odwzorowaniem ilorazowym,
ciągłość \(f\) implikuje ciągłość
\(\widetilde f\).
Jeśli \(W\subseteq G/N\) jest otwarty,
to \(q^{-1}(W)\) jest otwarty w \(G\),
a
\(\widetilde f(W)=f(q^{-1}(W))\)
jest otwarty w \(H\)
dzięki założeniu otwartości \(f\).
Zatem \(\widetilde f\) jest otwartą
ciągłą bijekcją, czyli homeomorfizmem.
\end{proof}

\begin{uwaga}[Dlaczego otwartość jest istotna?]
\label{rem:topog24-nonopen-map}
Nie każdy ciągły bijektywny
homomorfizm grup topologicznych
jest homeomorfizmem.
Na tej samej grupie addytywnej
\(\mathbb R\) rozważmy dwie topologie:
dyskretną \(\tau_d\)
i zwykłą euklidesową \(\tau_e\).
Identyczność
\((\mathbb R,\tau_d)\to(\mathbb R,\tau_e)\)
jest ciągłym bijektywnym homomorfizmem,
lecz jej odwrotność nie jest ciągła.
Dlatego nie usuwamy hipotezy
otwartości z
Twierdzenia~\ref{thm:topog24-isomorphism}.
\end{uwaga}

\begin{twierdzenie}[Zwartość i zwartość lokalna przestrzeni warstw]
\label{thm:topog24-coset-compact}
Niech \(G\) będzie grupą topologiczną,
a \(H\le G\) podgrupą domkniętą.
\begin{enumerate}[label=\textup{(\roman*)}]
\item Jeśli \(G\) jest zwarta, to \(G/H\)
jest zwartą przestrzenią Hausdorffa.
\item Jeśli \(G\) jest lokalnie zwarta,
to \(G/H\) jest lokalnie zwartą
przestrzenią Hausdorffa.
\end{enumerate}
\end{twierdzenie}
\begin{proof}
Hausdorffowość \(G/H\)
wynika z
Twierdzenia~\ref{thm:topog24-coset-hausdorff}.
W przypadku \textup{(i)}
przestrzeń \(G/H\) jest zwartym
obrazem ciągłym zwartej \(G\).

Załóżmy \textup{(ii)}.
Wybierzmy zwarty zbiór \(K\subseteq G\)
będący otoczeniem \(e\),
czyli \(e\in\operatorname{int}_G K\).
Zbiór \(q(K)\) jest zwarty,
a \(q(\operatorname{int}_G K)\)
jest otwartym otoczeniem
neutralnej warstwy \(H\)
zawartym w \(q(K)\), gdyż \(q\) jest otwarte.
Każda warstwa \(gH\)
ma analogiczne zwarte otoczenie:
lewe działanie \(G\) na \(G/H\),
określone wzorem
\(g'\cdot(xH)=(g'x)H\),
dla ustalonego \(g'\in G\)
jest homeomorfizmem,
ponieważ jego odwrotnością
jest działanie przez \((g')^{-1}\).
Przesuwając \(q(K)\), otrzymujemy
zwartość lokalną w każdym punkcie.
\end{proof}

\begin{przyklad}[Okrąg jako iloraz prostej]
\label{ex:topog24-circle-quotient}
Niech \(G=(\mathbb R,+)\)
i \(H=(\mathbb Z,+)\).
Podgrupa \(\mathbb Z\)
jest domknięta i normalna.
Odwzorowanie
\[
 p\colon\mathbb R\to S^1,\qquad
 p(t)=\exp(2\pi i t)
\]
jest ciągłym surjektywnym
homomorfizmem, którego jądrem
jest \(\mathbb Z\).
Jest także otwarte: obraz
każdego otwartego przedziału
długości mniejszej od \(1\)
jest otwartym łukiem okręgu,
a dowolny otwarty podzbiór prostej
jest sumą takich przedziałów.
Z Twierdzenia~\ref{thm:topog24-isomorphism}
otrzymujemy homeomorfizm grup
\[
 \mathbb R/\mathbb Z\simeq S^1.
\]
Każda warstwa \(t+\mathbb Z\)
jest zbiorem liczb rzeczywistych
różniących się o liczbę całkowitą.
Grupa \(\mathbb R/\mathbb Q\)
nie jest natomiast Hausdorffa,
ponieważ \(\mathbb Q\) jest
podgrupą niedomkniętą.
\end{przyklad}

\begin{figure}[htbp]
\centering
\begin{tikzpicture}[>=Stealth,scale=.85]
 \draw[->] (-4.4,0)--(4.5,0) node[below right] {$\mathbb R$};
 \foreach \k in {-4,-3,...,4} {
   \draw (\k,.10)--(\k,-.10);
 }
 \fill[blue!70!black] (-3,0) circle (2.2pt);
 \fill[blue!70!black] (-2,0) circle (2.2pt);
 \fill[blue!70!black] (-1,0) circle (2.2pt);
 \fill[blue!70!black] (0,0) circle (2.2pt);
 \fill[blue!70!black] (1,0) circle (2.2pt);
 \fill[blue!70!black] (2,0) circle (2.2pt);
 \fill[blue!70!black] (3,0) circle (2.2pt);
 \node[below,font=\small] at (-3,-.13) {$-3$};
 \node[below,font=\small] at (0,-.13) {$0$};
 \node[below,font=\small] at (3,-.13) {$3$};
 \draw[->,thick] (0,-.85)--(0,-1.55)
   node[midway,right,font=\small] {$t\mapsto e^{2\pi it}$};
 \draw[thick,blue!65!black] (0,-3.0) circle (1.05);
 \fill[blue!70!black] (0,-1.95) circle (2.2pt);
 \node[anchor=west,font=\small] at (1.18,-3.02)
 {$S^1\simeq\mathbb R/\mathbb Z$};
 \node[anchor=east,font=\small] at (-1.18,-1.97)
 {$\mathbb Z\mapsto1$};
\end{tikzpicture}
\caption{Wszystkie liczby całkowite tworzą jedną warstwę
\(\mathbb Z\) i są odwzorowywane na punkt \(1\) okręgu.
Przestrzenią warstw jest grupa topologiczna \(S^1\).}
\label{fig:topog24-circle-quotient}
\end{figure}

\section{Składowa spójna neutralności}
\label{sec:topog24-components}

\begin{twierdzenie}[Składowa neutralna jest normalna]
\label{thm:topog24-identity-component}
Niech \(G\) będzie grupą topologiczną.
Składowa spójna \(G^0\subseteq G\)
zawierająca \(e\) jest domkniętą
podgrupą normalną \(G\).
Co więcej, iloraz \(G/G^0\)
jest całkowicie niespójny,
tj. jego jedynymi niepustymi
spójnymi podzbiorami są singletony.
\end{twierdzenie}
\begin{proof}
Z definicji składowej spójnej
oraz wyników o spójności z rozdziału 08
zbiór \(G^0\) jest niepusty,
spójny i domknięty.
Iloczyn \(G^0\times G^0\) jest spójny
jako iloczyn dwóch spójnych przestrzeni.
Ciągły obraz przez \(D(x,y)=xy^{-1}\)
jest zatem spójny, a ponieważ zawiera
\(e\), musi zawierać się w \(G^0\)
przez maksymalność składowej.
W konsekwencji \(G^0\) jest podgrupą.

Dla każdego \(g\in G\)
homeomorfizm sprzężenia \(C_g\)
przenosi spójny zbiór \(G^0\)
na spójny zbiór zawierający \(e\).
Stąd \(gG^0g^{-1}\subseteq G^0\);
stosując tę inkluzję do \(g^{-1}\),
otrzymujemy równość.
Podgrupa \(G^0\) jest więc normalna.

Niech \(q\colon G\to G/G^0\)
będzie odwzorowaniem ilorazowym.
Jest ono otwarte,
a jego włókna są warstwami \(gG^0\),
czyli zbiorami spójnymi,
jako obrazy \(G^0\) przez translacje.
Pokażemy, że dla każdego spójnego
podzbioru \(C\subseteq G/G^0\)
przeciwobraz \(q^{-1}(C)\) jest spójny.
Gdyby istniał rozkład
\(q^{-1}(C)=A\cup B\)
na dwa niepuste rozłączne zbiory
otwarte w podprzestrzeni,
każde spójne włókno musiałoby
leżeć w całości w \(A\) lub w \(B\).
Zatem \(A,B\) byłyby nasycone
w \(q^{-1}(C)\).
Możemy je zapisać
\(A=O_A\cap q^{-1}(C)\),
\(B=O_B\cap q^{-1}(C)\),
gdzie \(O_A,O_B\subseteq G\)
są otwarte. Dzięki nasyceniu
otrzymujemy
\(q(A)=q(O_A)\cap C\)
i \(q(B)=q(O_B)\cap C\).
Są to dwa niepuste rozłączne
otwarte w \(C\) zbiory pokrywające \(C\),
ponieważ \(q\) jest otwarte.
Jest to sprzeczne ze spójnością \(C\).

Jeżeli teraz \(C\) jest spójnym
podzbiorem \(G/G^0\) zawierającym
neutralną warstwę \(G^0\),
to \(q^{-1}(C)\) jest spójny
i zawiera \(e\).
Z maksymalności \(G^0\)
otrzymujemy \(q^{-1}(C)=G^0\),
a więc \(C=\{G^0\}\).
Dowolny spójny podzbiór niepusty
ilorazu można przez translację
przenieść do spójnego podzbioru
zawierającego neutralną warstwę;
dlatego musi być singletonem.
\end{proof}

\begin{przyklad}[Grupy spójne i całkowicie niespójne]
\label{ex:topog24-components}
Grupa \((\mathbb R^n,+)\)
jest spójna drogowo, a więc \(G^0=G\).
Dla grupy dyskretnej \(G\)
każdy singleton jest otwarto-domknięty,
więc \(G^0=\{e\}\).

Jeśli \(I\) jest dowolnym zbiorem,
grupa \(\{0,1\}^I\) z dodawaniem
modulo \(2\) współrzędna po współrzędnej
i topologią produktową jest zwartą
grupą Hausdorffa.
Dla dwóch różnych punktów istnieje
współrzędna, na której się różnią,
a przeciwobrazy singletonów tej
współrzędnej są otwarto-domknięte.
Żaden spójny podzbiór nie zawiera
dwóch różnych punktów.
Zatem grupa jest całkowicie niespójna,
a jej składowa neutralna jest singletonem.
\end{przyklad}

\section{Naturalne struktury jednostajne grupy}
\label{sec:topog24-uniformities}

\begin{definicja}[Lewa i prawa struktura jednostajna]
\label{def:topog24-uniformities}
Niech \(G\) będzie grupą topologiczną,
a \(U\) przebiega otwarte otoczenia \(e\).
Definiujemy relacje na \(G\times G\):
\[
 E_U^{L}=\{(x,y):x^{-1}y\in U\},\qquad
 E_U^{R}=\{(x,y):xy^{-1}\in U\}.
\]
Rodziny \((E_U^L)\) i \((E_U^R)\)
są bazami odpowiednio
\emph{lewej} i \emph{prawej}
struktury jednostajnej na \(G\).
\end{definicja}

\begin{twierdzenie}[Poprawność i zgodność struktur jednostajnych]
\label{thm:topog24-uniformities}
Niech \(G\) będzie grupą topologiczną.
Wówczas:
\begin{enumerate}[label=\textup{(\roman*)}]
\item obie rodziny z
Definicji~\ref{def:topog24-uniformities}
spełniają aksjomaty baz struktur jednostajnych
i indukują wyjściową topologię na \(G\);
\item relacje \(E_U^L\) są niezmiennicze
przy wspólnej lewej translacji obu argumentów,
a \(E_U^R\) przy wspólnej prawej translacji;
\item odwracanie \(x\mapsto x^{-1}\)
jest izomorfizmem jednostajnym
z lewej struktury na prawą;
\item każdy ciągły homomorfizm
\(f\colon G\to H\) jest jednostajnie
ciągły dla obu lewych struktur
oraz, osobno, dla obu prawych struktur.
\end{enumerate}
\end{twierdzenie}
\begin{proof}
Każda relacja \(E_U^L\)
zawiera przekątną, gdyż \(e\in U\).
Dla dwóch otoczeń \(U,V\)
mamy \(E_{U\cap V}^L=E_U^L\cap E_V^L\);
odwrotność relacji \(E_U^L\)
to \(E_{U^{-1}}^L\).
Dla danego \(U\)
Lemat~\ref{lem:topog24-small-neighborhoods}
dostarcza symetrycznego otoczenia \(V\ni e\)
z \(VV\subseteq U\).
Jeśli \((x,y)\in E_V^L\)
i \((y,z)\in E_V^L\),
to \(x^{-1}z=(x^{-1}y)(y^{-1}z)\in VV\subseteq U\).
Zatem \(E_V^L\circ E_V^L\subseteq E_U^L\).
Po domknięciu w górę otrzymujemy
strukturę jednostajną.
Jej otoczenia punktu \(x\)
mają bazę \(E_U^L[x]=xU\),
czyli bazę wyjściowej topologii.
Wszystkie te argumenty stosują się
analogicznie do prawej struktury,
dla której
\(E_U^R[x]=U^{-1}x\).
To dowodzi \textup{(i)}.

Dla \(g,x,y\in G\)
mamy \((gx)^{-1}(gy)=x^{-1}y\)
oraz \((xg)(yg)^{-1}=xy^{-1}\),
co dowodzi \textup{(ii)}.

Dla odwzorowania \(\iota(x)=x^{-1}\)
zachodzi
\[
 (x,y)\in E_U^L
 \quad\Longleftrightarrow\quad
 (x^{-1},y^{-1})\in E_U^R,
\]
ponieważ
\(x^{-1}(y^{-1})^{-1}=x^{-1}y\).
Ta sama równość daje jednostajną
ciągłość \(\iota\) i jej odwrotności,
gdy stosujemy odpowiednie struktury.
Otrzymujemy \textup{(iii)}.

Dla \textup{(iv)} niech \(W\)
będzie otwartym otoczeniem
elementu neutralnego w \(H\).
Ciągłość \(f\) w \(e_G\)
daje otwarte \(U\ni e_G\)
z \(f(U)\subseteq W\).
Gdy \(x^{-1}y\in U\),
to \(f(x)^{-1}f(y)=f(x^{-1}y)\in W\).
To jest jednostajna ciągłość
między lewymi strukturami.
Równanie \(f(x)f(y)^{-1}=f(xy^{-1})\)
daje dowód dla prawych struktur.
\end{proof}

\begin{wniosek}[Grupy przemienne]
\label{cor:topog24-abelian-uniform}
Jeżeli \(G\) jest przemienna,
to lewa i prawa struktura jednostajna
na \(G\) są równe.
\end{wniosek}
\begin{proof}
Dla wszystkich \(x,y\in G\)
zachodzi \(xy^{-1}=(x^{-1}y)^{-1}\).
Zatem
\(E_U^L=E_{U^{-1}}^R\).
Ponieważ odwracanie zamienia
otoczenia \(e\) na otoczenia \(e\),
obie rodziny generują te same relacje.
\end{proof}

\section{Metryzowalność i twierdzenie Birkhoffa--Kakutaniego}
\label{sec:topog24-metrization}

\begin{twierdzenie}[Twierdzenie Birkhoffa--Kakutaniego]
\label{thm:topog24-birkhoff-kakutani}
Niech \(G\) będzie grupą topologiczną Hausdorffa.
Następujące warunki są równoważne:
\begin{enumerate}[label=\textup{(\roman*)}]
\item \(G\) spełnia pierwszy aksjomat przeliczalności;
\item istnieje metryka \(d\) na \(G\), wyznaczająca
topologię \(G\) i lewostronnie niezmiennicza, to znaczy
\[
 d(gx,gy)=d(x,y)
 \quad\text{dla wszystkich }g,x,y\in G.
\]
\end{enumerate}
Można przy tym wybrać metrykę ograniczoną przez \(1\).
\end{twierdzenie}
\begin{proof}
Jeśli zachodzi \textup{(ii)}, kule
\(B_d(e,1/n)\), \(n\ge1\),
tworzą przeliczalną bazę otoczeń \(e\).
Z jednorodności
(Wniosek~\ref{cor:topog24-homogeneous})
wynika pierwszy aksjomat przeliczalności
w każdym punkcie \(G\).

Załóżmy \textup{(i)}.
Wybierzmy przeliczalną bazę
otwartych otoczeń \(B_n\ni e\), \(n\ge1\).
Budujemy indukcyjnie symetryczne
otwarte otoczenia \(U_n\ni e\)
z \(U_0=G\) i
\[
 U_{n+1}^3\subseteq U_n\cap B_{n+1}
 \qquad(n\ge0).
\]
Istnienie wynika z
Lematu~\ref{lem:topog24-small-neighborhoods}
dla \(k=3\). Ponieważ
\(e\in U_{n+1}\), mamy
\(U_{n+1}\subseteq U_{n+1}^3\subseteq U_n\);
rodzina \((U_n)\) jest więc zstępująca
i stanowi bazę otoczeń \(e\).
Z Hausdorffowości \(G\) wynika
\(\bigcap_{n\ge0}U_n=\{e\}\):
dla \(g\neq e\) istnieje otwarte
otoczenie \(O\ni e\) z \(g\notin O\),
a pewien \(U_n\subseteq O\).

Definiujemy funkcję \(\rho\colon G\to[0,1]\)
wzorem
\[
 \rho(g)=\inf\{2^{-n}:g\in U_n,\ n\ge0\}.
\]
Jest ona dobrze określona,
bo \(U_0=G\), oraz
\(\rho(e)=0\),
\(\rho(g)=\rho(g^{-1})\).
Ponieważ \(\bigcap_nU_n=\{e\}\),
dla \(g\neq e\) istnieje największy
indeks \(n(g)\), dla którego \(g\in U_{n(g)}\).
Dlatego \(\rho(g)=2^{-n(g)}>0\).
Ponadto dla dowolnego \(n\ge0\)
\[
 \rho(g)\le 2^{-n}
 \quad\Longleftrightarrow\quad g\in U_n.
\]

\emph{Krok 1: oszacowanie iloczynu.}
Udowodnimy indukcyjnie względem liczby
czynników \(m\ge1\), że dla każdego \(n\ge0\)
i \(g_1,\ldots,g_m\in G\) zachodzi
\[
 \sum_{j=1}^{m}\rho(g_j)<2^{-n}
 \quad\Longrightarrow\quad
 g_1\cdots g_m\in U_n.
 \tag{*}
\]
Dla \(m=1\) mamy
\(\rho(g_1)<2^{-n}\),
a więc \(g_1\in U_{n+1}\subseteq U_n\).
Załóżmy teraz, że wynik zachodzi
dla wszystkich iloczynów krótszych
niż \(m\), i połóżmy
\(S=\sum_{j=1}^m\rho(g_j)<2^{-n}\),
\(h=2^{-(n+1)}\).

Gdy \(S<h\), rozbijmy iloczyn
na pierwszy czynnik i pozostałe \(m-1\).
Suma wag w każdej z obu części
jest mniejsza niż \(h\).
Z założenia indukcyjnego
oba iloczyny należą do \(U_{n+1}\),
stąd cały iloczyn leży w
\(U_{n+1}^2\subseteq U_n\).

Niech teraz \(h\le S<2h\).
Wybierzmy najmniejszy indeks \(j\)
taki, że
\(\rho(g_1)+\cdots+\rho(g_j)\ge h\).
Suma wag czynników poprzedzających \(g_j\)
jest mniejsza niż \(h\).
Suma wag czynników po \(g_j\)
też jest mniejsza niż \(h\),
ponieważ suma wag do \(j\) wynosi
co najmniej \(h\), a \(S<2h\).
Ponadto \(\rho(g_j)<2h=2^{-n}\);
ponieważ dodatnie wartości \(\rho\)
są potęgami \(2^{-k}\),
mamy \(\rho(g_j)\le h\)
(także gdy \(\rho(g_j)=0\)).
Zatem \(g_j\in U_{n+1}\).
Każdy z dwóch pozostałych,
ewentualnie pustych, iloczynów
należy do \(U_{n+1}\)
na mocy założenia indukcyjnego;
pusty iloczyn to \(e\in U_{n+1}\).
Cały iloczyn należy więc do
\(U_{n+1}^3\subseteq U_n\).
Dowodzi to \((*)\).

\emph{Krok 2: konstrukcja metryki.}
Dla \(x,y\in G\) określamy
\[
 d(x,y)=\inf\left\{
 \sum_{j=1}^{m}\rho(x_{j-1}^{-1}x_j):
 m\ge1,\ x_0=x,\ x_m=y,\
 x_1,\ldots,x_{m-1}\in G
 \right\}.
\]
Zbiór wartości, z którego bierzemy infimum,
jest niepusty, ponieważ można wybrać
łańcuch jednoodcinkowy \(x_0=x,x_1=y\).
Stąd \(0\le d(x,y)\le\rho(x^{-1}y)\le1\).
Z symetrii \(\rho\) wynika symetria \(d\),
a z łączenia dwóch skończonych łańcuchów
wynika nierówność trójkąta:
dla \(\varepsilon>0\) wybieramy
łańcuch z \(x\) do \(y\) o koszcie
mniejszym niż \(d(x,y)+\varepsilon/2\)
oraz z \(y\) do \(z\) o koszcie
mniejszym niż \(d(y,z)+\varepsilon/2\).
Ich połączenie ma koszt mniejszy
niż \(d(x,y)+d(y,z)+\varepsilon\);
przechodzimy z \(\varepsilon\downarrow0\).

Sprawdzimy dodatnią określoność.
Jeśli \(d(x,y)<2^{-n}\),
istnieje łańcuch \(x=x_0,\ldots,x_m=y\)
o sumie wag mniejszej niż \(2^{-n}\).
Z \((*)\) otrzymujemy
\[
 x^{-1}y=
 (x_0^{-1}x_1)\cdots(x_{m-1}^{-1}x_m)
 \in U_n.
\]
Jeśli \(x\neq y\), możemy wybrać \(n\)
takie, że \(x^{-1}y\notin U_n\).
Wówczas \(d(x,y)\ge2^{-n}>0\).
Zatem \(d\) jest metryką.
Jest lewostronnie niezmiennicza,
ponieważ wspólna lewa translacja
całego łańcucha nie zmienia różnic
\((gx_{j-1})^{-1}(gx_j)=x_{j-1}^{-1}x_j\).
Po wzięciu infimum otrzymujemy
\(d(gx,gy)=d(x,y)\).

\emph{Krok 3: zgodność topologii.}
Dla każdego \(n\ge0\) zachodzą inkluzje
\[
 U_{n+1}\subseteq
 B_d(e,2^{-n})\subseteq U_n.
\]
Pierwsza wynika z
\(d(e,g)\le\rho(g)\le2^{-(n+1)}<2^{-n}\)
dla \(g\in U_{n+1}\).
Druga wynika z \((*)\) i definicji \(d\).
Ponieważ \((U_n)\) jest bazą otoczeń \(e\),
obie topologie mają tę samą bazę otoczeń \(e\):
dla dowolnego otwartego \(O\ni e\)
pewne \(U_n\subseteq O\),
a wtedy \(B_d(e,2^{-n})\subseteq O\);
z kolei każda kula \(B_d(e,\varepsilon)\)
zawiera \(U_{n+1}\) dla dostatecznie dużego \(n\).
Translacje są homeomorfizmami
w topologii wyjściowej,
a metryka \(d\) jest lewostronnie
niezmiennicza, więc zgoda baz w \(e\)
oznacza zgodność topologii we wszystkich punktach.
\end{proof}

\begin{wniosek}[Metryka prawostronnie niezmiennicza]
\label{cor:topog24-right-metric}
Jeżeli grupa topologiczna \(G\) spełnia
pierwszy aksjomat przeliczalności,
to istnieje także zgodna z jej topologią
metryka prawostronnie niezmiennicza.
\end{wniosek}
\begin{proof}
Niech \(d_L\) będzie metryką
z Twierdzenia~\ref{thm:topog24-birkhoff-kakutani}.
Definiujemy
\(d_R(x,y)=d_L(x^{-1},y^{-1})\).
Ciągłość odwracania i jego odwrotności
daje zgodność topologii \(d_R\) i \(G\),
a dla wszystkich \(x,y,g\in G\)
\[
 d_R(xg,yg)=
 d_L(g^{-1}x^{-1},g^{-1}y^{-1})
 =d_L(x^{-1},y^{-1})
 =d_R(x,y).
\]
\end{proof}

\begin{przyklad}[Dwie przestrzenie produktowe różnej przeliczalności]
\label{ex:topog24-uncountable-cube}
Niech \(I\) będzie nieprzeliczalny.
Zwartą grupę topologiczną
\(G=\{0,1\}^I\) z topologią produktową
poznaliśmy już w
Przykładzie~\ref{ex:topog24-components}.
Nie spełnia ona pierwszego aksjomatu
przeliczalności w \(e=(0)_{i\in I}\).
Istotnie, przypuśćmy, że
\((U_n)_{n\ge1}\) jest przeliczalną
bazą otoczeń \(e\).
Wybierzmy dla każdego \(n\)
podstawowe otoczenie cylindryczne
\(V_n\) z \(e\in V_n\subseteq U_n\).
Każde \(V_n\) ogranicza współrzędne
tylko ze skończonego zbioru \(J_n\subseteq I\).
Zbiór \(J=\bigcup_nJ_n\) jest przeliczalny;
istnieje \(i_0\in I\setminus J\).
Otwarte otoczenie
\[
 O=\{x\in G:x_{i_0}=0\}
\]
nie zawiera żadnego \(V_n\),
bo dla każdego \(n\) możemy
zmienić współrzędną \(i_0\)
na \(1\) bez opuszczania \(V_n\).
Jest to sprzeczne z bazowością \(U_n\),
gdyż z \(U_n\subseteq O\)
wynikałoby także \(V_n\subseteq O\).
Zatem \(G\) nie jest metryzowalna,
ponieważ każda przestrzeń metryczna
spełnia pierwszy aksjomat przeliczalności.

Z drugiej strony, każda grupa dyskretna,
nawet nieprzeliczalna, spełnia pierwszy
aksjomat przeliczalności (baza \(\{e\}\)),
choć nie spełnia drugiego aksjomatu
przeliczalności, jeśli ma nieprzeliczalnie
wiele punktów.
Własności te należy rozróżniać.
\end{przyklad}

\par\smallskip\noindent
\textbf{Podsumowanie.}
Grupy topologiczne są jednorodne, regularne
i opisane lokalnie przez otoczenia \(e\).
Ich podgrupy domknięte wyznaczają
przestrzenie warstw Hausdorffa, a normalne
podgrupy domknięte --- grupy ilorazowe.
Składowa neutralna jest domknięta i normalna,
a iloraz przez nią jest całkowicie niespójny.
Lewa i prawa struktura jednostajna
pozwalają badać ciągłość homomorfizmów
bez wprowadzania miary Haara.
\chapter{Continua i ich struktura}
\label{ch:topog25-continua}

\noindent\textbf{Cel rozdziału.}
Zwartość i spójność pozwalają wyróżnić przestrzenie, które
mogą mieć złożoną strukturę mimo braku rozdzielenia.
W rozdziale~\ref{ch:topog-spojnosc} omówiliśmy już
podstawowe własności spójności, składowe spójne,
spójność drogową i lokalną oraz krzywą sinusową topologa.
Nie powtarzamy tych dowodów.
Zbadamy istnienie minimalnych podcontinuów,
różnicę między continuami rozkładalnymi i nierozkładalnymi,
granice odwrotne i przykłady, w których lokalna spójność
zawodzi.
Skonstruujemy continuum nierozkładalne Knastera
oraz solenoid dyadyczny.
Na koniec wskażemy dokładny zakres związków
z funkcjami Peana.

\section{Continua i podcontinua}
\label{sec:topog25-basic}

\begin{definicja}[Continuum i podcontinuum]
\label{def:topog25-continuum}
\emph{Continuum} nazywamy niepustą, zwartą,
spójną przestrzeń Hausdorffa.
Jeżeli \(X\) jest continuum,
to \emph{podcontinuum} \(X\)
jest niepustym, spójnym, zwartym podzbiorem \(K\subseteq X\)
z topologią podprzestrzeni.
Continuum nazywamy \emph{nietrywialnym},
gdy zawiera co najmniej dwa punkty.
\emph{Continuum metryczne} to continuum,
którego topologia jest metryzowalna.
\end{definicja}

\begin{uwaga}[Rola aksjomatów]
\label{rem:topog25-closed}
W przestrzeni Hausdorffa każdy podzbiór zwarty
jest domknięty. Zatem podcontinua są zbiorami
domkniętymi w \(X\). Odwrotnie, każdy
niepusty, domknięty i spójny podzbiór
continuum \(X\) jest podcontinuum.
Definicja dopuszcza przestrzeń jednopunktową;
w niektórych zastosowaniach zastrzega się
nietrywialność, dlatego będziemy
ją jawnie zaznaczać, kiedy jest potrzebna.
\end{uwaga}

\begin{stwierdzenie}[Obrazy i produkty continuów]
\label{prop:topog25-images-products}
Niech \(X\) będzie continuum, a \(Y\)
przestrzenią Hausdorffa.
Jeżeli \(f\colon X\to Y\) jest ciągłe,
to \(f(X)\) jest continuum z topologią podprzestrzeni.
Ponadto iloczyn dowolnej niepustej rodziny
continuów jest continuum.
\end{stwierdzenie}
\begin{proof}
Obraz niepustego zbioru jest niepusty.
Jest zwarty jako ciągły obraz zwartego \(X\).
Jest spójny na mocy
Twierdzenia~\ref{thm:topog-spojnosc-obraz-ciagly};
jako podprzestrzeń \(Y\) jest Hausdorffa.
Stąd \(f(X)\) jest continuum.
Dla rodziny \((X_i)_{i\in I}\)
każdy czynnik jest niepusty i zwarty,
więc iloczyn jest niepusty i zwarty
na mocy twierdzenia Tychonoffa
(przy użyciu aksjomatu wyboru).
Jest Hausdorffa, a spójność produktu
wynika z
Twierdzenia~\ref{thm:topog-spojnosc-iloczyn}.
\end{proof}

\begin{przyklad}[Proste continua]
\label{ex:topog25-basic}
Nietrywialny odcinek \([a,b]\subseteq\mathbb R\)
jest continuum metrycznym.
Okrąg \(S^1\subseteq\mathbb R^2\)
jest continuum, ponieważ jest obrazem
odcinka przy odwzorowaniu
\(t\mapsto(\cos(2\pi t),\sin(2\pi t))\).
Z kolei zbiór Cantora jest zwarty,
ale nie jest continuum:
jest niespójny i zawiera więcej niż jeden punkt.
\end{przyklad}

\begin{twierdzenie}[Przecięcie zstępującego ciągu podcontinuów]
\label{thm:topog25-nested}
Niech \(X\) będzie zwartą przestrzenią Hausdorffa,
a \((K_n)_{n\ge1}\) ciągiem niepustych
zwartych i spójnych podzbiorów \(X\),
spełniającym \(K_{n+1}\subseteq K_n\).
Wówczas
\[
 K=\bigcap_{n=1}^{\infty}K_n
\]
jest niepustym podcontinuum \(X\).
\end{twierdzenie}
\begin{proof}
Zbiory \(K_n\) są domknięte w \(X\).
Mają własność skończonych przecięć,
bo każda skończona rodzina przecina się
w najmniejszym z rozważanych zbiorów.
Zwartość \(X\) daje \(K\neq\varnothing\).
Przecięcie zbiorów domkniętych
jest domknięte, zatem \(K\)
jest także zwarte.

Załóżmy, że \(K\) nie jest spójny.
Istnieją niepuste rozłączne zbiory
\(A,B\subseteq K\), otwarte i domknięte w \(K\),
takie że \(K=A\cup B\).
Są one zwarte, a więc domknięte w \(X\).
Normalność zwartego Hausdorffa \(X\)
daje rozłączne otwarte \(U,V\subseteq X\)
z \(A\subseteq U\), \(B\subseteq V\).
Wtedy \(K\subseteq U\cup V\).

Twierdzimy, że istnieje \(N\)
takie, że \(K_N\subseteq U\cup V\).
W przeciwnym razie niepuste
zwarte zbiory \(K_n\setminus(U\cup V)\)
tworzyłyby ciąg zstępujący.
Zwartość \(X\) dałaby punkt ich
przecięcia należący jednocześnie
do \(K\) i do \(X\setminus(U\cup V)\),
co przeczyłoby \(K\subseteq U\cup V\).
Zatem taki \(N\) istnieje.

Ponieważ \(A,B\subseteq K\subseteq K_N\),
oba zbiory \(K_N\cap U\) i \(K_N\cap V\)
są niepuste, otwarte w \(K_N\),
rozłączne i pokrywają \(K_N\).
Jest to rozdzielenie \(K_N\),
wbrew jego spójności.
\end{proof}

\begin{lemat}[Przecięcie skierowanej rodziny continuów]
\label{lem:topog25-chain-intersection}
Niech \(X\) będzie continuum
i \(\mathcal C\) niepustą rodziną jego
podcontinuów uporządkowaną liniowo
przez inkluzję.
Wówczas \(\bigcap_{C\in\mathcal C}C\)
jest niepustym podcontinuum \(X\).
\end{lemat}
\begin{proof}
Każdy \(C\in\mathcal C\) jest
niepusty i domknięty w zwartym \(X\).
Dzięki liniowemu uporządkowaniu
dowolne skończenie wiele członów
rodziny ma niepuste przecięcie,
równe najmniejszemu z nich.
Zatem \(K=\bigcap_{C\in\mathcal C}C\)
jest niepusty i zwarty.

Jeśli \(K\) byłby niespójny,
jak w dowodzie
Twierdzenia~\ref{thm:topog25-nested}
znaleźlibyśmy rozłączne otwarte
\(U,V\subseteq X\), których suma
zawiera \(K\) i które spotykają \(K\).
Zwartość \(X\) implikuje, że pewna
skończona podrodzina \(\mathcal C\)
ma przecięcie zawarte w \(U\cup V\):
istotnie, otwarte zbiory
\(X\setminus C\), \(C\in\mathcal C\),
wraz z \(U\cup V\) pokrywają \(X\);
wybieramy skończone podpokrycie.
Ponieważ \(\mathcal C\) jest łańcuchem,
jego przecięcie skończone jest
pewnym \(C_0\in\mathcal C\).
Wobec tego \(C_0\subseteq U\cup V\).
Ponadto \(K\subseteq C_0\)
spotyka zarówno \(U\), jak i \(V\).
Otrzymujemy rozdzielenie \(C_0\),
co jest sprzeczne ze spójnością.
\end{proof}

\section{Continua nierozkładalne i nieprzywiedlne}
\label{sec:topog25-indecomposable}

\begin{definicja}[Rozkładalność]
\label{def:topog25-indecomposable}
Continuum \(X\) jest
\emph{rozkładalne}, jeżeli istnieją
jego dwa właściwe podcontinua \(A,B\subsetneq X\)
takie, że \(X=A\cup B\).
W przeciwnym razie nazywamy \(X\)
\emph{nierozkładalnym}.
Nie należy utożsamiać rozkładalności
z niespójnością: continuum z definicji
jest spójne.
\end{definicja}

\begin{przyklad}[Przykłady rozkładalne]
\label{ex:topog25-decomposable}
Odcinek \([0,1]\) jest rozkładalny, gdyż
\([0,1]=[0,1/2]\cup[1/2,1]\),
a obie części są właściwymi podcontinuami.
Okrąg jest rozkładalny:
można przedstawić go jako sumę
dwóch domkniętych półokręgów.
Continuum jednoelementowe jest
nierozkładalne w przyjętej definicji,
ponieważ nie ma właściwego niepustego
podcontinuum; dlatego ważne przykłady
nierozkładalne będziemy wybierać
nietrywialne.
\end{przyklad}

\begin{stwierdzenie}[Właściwe podcontinua continuum nierozkładalnego]
\label{prop:topog25-nowhere-dense}
Niech \(X\) będzie nietrywialnym continuum,
w którym każde właściwe podcontinuum
ma puste wnętrze w \(X\).
Wówczas \(X\) jest nierozkładalne.
Równoważnie: jeżeli \(X\) jest rozkładalne,
to ma właściwe podcontinuum
o niepustym wnętrzu.
\end{stwierdzenie}
\begin{proof}
Gdyby \(X=A\cup B\) dla właściwych
podcontinuów \(A,B\), to \(B\)
byłby domknięty w \(X\), a zbiór
\(X\setminus B\) byłby niepusty i otwarty.
Z równości \(X=A\cup B\)
wynika \(X\setminus B\subseteq A\).
Zatem \(A\) byłoby właściwym
podcontinuum o niepustym wnętrzu,
wbrew założeniu.
\end{proof}

\begin{definicja}[Continuum nieprzywiedlne między punktami]
\label{def:topog25-irreducible}
Niech \(X\) będzie continuum,
a \(a,b\in X\).
Mówimy, że \(X\) jest
\emph{nieprzywiedlne między \(a\) i \(b\)},
jeśli nie istnieje właściwe
podcontinuum \(K\subsetneq X\)
zawierające jednocześnie \(a\) i \(b\).
\end{definicja}

\begin{twierdzenie}[Istnienie minimalnego continuum między punktami]
\label{thm:topog25-irreducible-existence}
Niech \(X\) będzie continuum,
a \(a,b\in X\).
Istnieje podcontinuum
\(K\subseteq X\), zawierające \(a,b\),
które jest nieprzywiedlne między
tymi punktami.
\end{twierdzenie}
\begin{proof}
Niech \(\mathcal P\) będzie rodziną
wszystkich podcontinuów \(C\subseteq X\)
zawierających \(a,b\).
Jest ona niepusta, ponieważ \(X\in\mathcal P\).
Uporządkujmy ją przez odwrotną inkluzję:
\(C\preceq D\) wtedy i tylko wtedy,
gdy \(D\subseteq C\).

Niech \(\mathcal L\subseteq\mathcal P\)
będzie niepustym łańcuchem.
Na mocy
Lematu~\ref{lem:topog25-chain-intersection}
zbiór
\(K_{\mathcal L}=\bigcap_{C\in\mathcal L}C\)
jest niepustym podcontinuum.
Zawiera \(a,b\) i jest górnym
ograniczeniem łańcucha
w porządku \(\preceq\).
Dla łańcucha pustego można
wybrać górne ograniczenie \(X\).
Lemat Kuratowskiego--Zorna
daje element maksymalny \(K\)
dla \(\preceq\), czyli
element minimalny dla inkluzji.
Gdyby właściwe podcontinuum
\(D\subsetneq K\) zawierało \(a,b\),
to \(D\in\mathcal P\) i \(K\prec D\),
wbrew maksymalności \(K\).
Zatem \(K\) jest nieprzywiedlne
między \(a,b\).
\end{proof}

\begin{uwaga}[Dwa różne pojęcia]
\label{rem:topog25-irred-vs-indec}
Nieprzywiedlność między punktami
i nierozkładalność nie są synonimami.
Odcinek \([0,1]\) jest
nieprzywiedlny między \(0\) i \(1\):
każdy spójny podzbiór prostej
zawierający te punkty zawiera
cały odcinek. Jednocześnie
odcinek jest rozkładalny.
\end{uwaga}

\section{Granice odwrotne przestrzeni zwartych}
\label{sec:topog25-inverse}

\begin{definicja}[Granica odwrotna ciągu]
\label{def:topog25-inverse-limit}
Niech \((X_n)_{n\ge0}\) będzie
ciągiem niepustych przestrzeni
zwartych Hausdorffa, a
\(f_n\colon X_{n+1}\to X_n\)
ciągłymi odwzorowaniami.
\emph{Granicę odwrotną} układu
\((X_n,f_n)\) definiujemy jako
\[
 \varprojlim(X_n,f_n)=
 \left\{(x_0,x_1,\ldots)\in
  \prod_{n=0}^{\infty}X_n:
  f_n(x_{n+1})=x_n\quad(n\ge0)\right\}
\]
z topologią podprzestrzeni
produktu. Ograniczenie \(n\)-tego
rzutowania do granicy oznaczamy
przez \(p_n\).
\end{definicja}

\begin{twierdzenie}[Granica odwrotna continuów]
\label{thm:topog25-inverse-connected}
Niech \((X_n)_{n\ge0}\) będzie
ciągiem continuów, a
\(f_n\colon X_{n+1}\to X_n\)
ciągłymi surjekcjami.
Wówczas granica odwrotna
\(X=\varprojlim(X_n,f_n)\)
jest continuum.
Wszystkie rzutowania
\(p_n\colon X\to X_n\)
są ciągłe i surjektywne.
Jeśli ponadto każde \(X_n\)
jest metryzowalne,
to \(X\) jest continuum metrycznym.
\end{twierdzenie}
\begin{proof}
Produkt \(P=\prod_{n\ge0}X_n\)
jest zwarty Hausdorffa na mocy
twierdzenia Tychonoffa.
Dla każdego \(n\) warunek
\(f_n(x_{n+1})=x_n\)
wyznacza zbiór domknięty w \(P\),
ponieważ jest przeciwobrazem
przekątnej domkniętej w
\(X_n\times X_n\).
Zatem \(X\) jest domknięty
w \(P\), a więc zwarty
i Hausdorffa.

Udowodnimy niepustość
oraz surjektywność \(p_n\).
Dla ustalonego \(n\) i
\(a\in X_n\) możemy wyznaczyć
kolejno \(x_{n+1},x_{n+2},\ldots\)
tak, aby
\(f_j(x_{j+1})=x_j\)
i \(x_n=a\); używamy kolejno
surjektywności \(f_j\)
oraz aksjomatu wyboru zależnego.
Dla \(j<n\) kładziemy
\(x_j=f_j\circ\cdots\circ
f_{n-1}(a)\).
Otrzymujemy punkt \(x\in X\)
z \(p_n(x)=a\).
W szczególności \(X\neq\varnothing\).

Pozostaje spójność.
Załóżmy przeciwnie, że
\(X=A\cup B\), gdzie \(A,B\)
są niepustymi rozłącznymi
zbiorami otwarto-domkniętymi
w \(X\). Oba są zwarte.
Pokażemy, że dla pewnego \(n\)
zwarte zbiory \(p_n(A)\)
i \(p_n(B)\) są rozłączne.
Gdyby dla każdego \(n\)
ich przecięcie było niepuste,
zbiór
\[
 F_n=\{(a,b)\in A\times B:p_n(a)=p_n(b)\}
\]
byłby niepusty i domknięty
w zwartym \(A\times B\).
Wynika z warunków zgodności,
że \(F_{n+1}\subseteq F_n\).
Zwartość dałaby parę
\((a,b)\in\bigcap_{n\ge0}F_n\).
Wtedy \(p_n(a)=p_n(b)\)
dla wszystkich \(n\), więc
\(a=b\) jako ciągi współrzędnych,
wbrew \(A\cap B=\varnothing\).
Dla pewnego \(n\)
mamy więc \(p_n(A)\cap p_n(B)=\varnothing\).

Ponieważ \(p_n\) jest surjektywne
oraz \(X=A\cup B\), otrzymujemy
\[
 X_n=p_n(A)\cup p_n(B).
\]
Oba zbiory po prawej są
niepuste, zwarte, a więc
domknięte w \(X_n\), i rozłączne.
Stanowią rozdzielenie \(X_n\),
sprzeczne ze spójnością \(X_n\).
To dowodzi spójności \(X\).

Jeśli każde \(X_n\) jest
metryzowalne, wybieramy zgodne
ograniczone metryki \(d_n\)
na \(X_n\). Wówczas
\[
 d(x,y)=\sum_{n=0}^{\infty}
  2^{-(n+1)}\min\{1,d_n(x_n,y_n)\}
\]
jest metryką na \(P\),
wyznaczającą topologię produktu.
Jej ograniczenie do \(X\)
metryzuje \(X\).
\end{proof}

\begin{stwierdzenie}[Continuum z rzutami przedziałowymi]
\label{prop:topog25-inverse-intervals}
Niech każde \(X_n=[0,1]\),
a \(f_n\colon[0,1]\to[0,1]\)
będzie ciągłą surjekcją.
Wówczas granica odwrotna
jest niepustym continuum metrycznym,
a wszystkie jej rzutowania
na przedział są surjektywne.
\end{stwierdzenie}
\begin{proof}
Przedział \([0,1]\)
jest nietrywialnym continuum
metrycznym; wystarczy zastosować
Twierdzenie~\ref{thm:topog25-inverse-connected}.
\end{proof}

\section{Continuum nierozkładalne Knastera}
\label{sec:topog25-knaster}

\begin{definicja}[Odwzorowanie namiotowe]
\label{def:topog25-tent}
Niech \(T\colon[0,1]\to[0,1]\)
będzie odwzorowaniem ciągłym danym wzorem
\[
 T(t)=
 \begin{cases}
  2t,&0\le t\le\tfrac12,\\
  2-2t,&\tfrac12\le t\le1.
 \end{cases}
\]
Przez \(K_T\) oznaczamy granicę
odwrotną stałego układu
\([0,1]\xleftarrow{T}[0,1]
\xleftarrow{T}\cdots\).
\end{definicja}

\begin{lemat}[Pełne gałęzie iteracji namiotu]
\label{lem:topog25-tent-branches}
Dla każdego \(r\ge1\) i
\(j\in\{0,\ldots,2^r-1\}\)
ograniczenie \(T^r\) do przedziału
\[
 I_{r,j}=\left[\frac{j}{2^r},
                     \frac{j+1}{2^r}\right]
\]
jest afiniczną bijekcją tego
przedziału na \([0,1]\).
W szczególności, dla każdego
niepustego otwartego podzbioru
\(U\subseteq[0,1]\), zbiór
\((T^r)^{-1}(U)\) spotyka
przedziały \(I_{r,0}\) i
\(I_{r,2^r-1}\).
\end{lemat}
\begin{proof}
Dla \(r=1\) stwierdzenie
wynika z definicji \(T\):
na obu połówkach przedziału
jest ono afiniczne i odwzorowuje
je bijektywnie na \([0,1]\).

Załóżmy tezę dla \(r\).
Na każdym \(I_{r,j}\)
odwzorowanie \(T^r\) jest
afiniczną bijekcją na \([0,1]\).
Przeciwobrazy dwóch przedziałów
\([0,\tfrac12]\) i
\([\tfrac12,1]\)
przy tym ograniczeniu dzielą
\(I_{r,j}\) na dwie równe połowy.
Na każdej z nich \(T^{r+1}=T\circ T^r\)
jest afiniczną bijekcją na \([0,1]\).
Są to przedziały
\(I_{r+1,2j}\) i \(I_{r+1,2j+1}\).
Indukcja dowodzi pierwszej tezy.
Druga wynika z surjektywności
każdego z tych ograniczeń.
\end{proof}

\begin{twierdzenie}[Continuum Knastera jest nierozkładalne]
\label{thm:topog25-knaster}
Continuum metryczne \(K_T\)
z Definicji~\ref{def:topog25-tent}
jest nietrywialne i nierozkładalne.
Co więcej, każde jego właściwe
podcontinuum ma puste wnętrze
w \(K_T\).
\end{twierdzenie}
\begin{proof}
Odwzorowanie \(T\) jest ciągłą
surjekcją \([0,1]\) na siebie.
Z Twierdzenia~\ref{thm:topog25-inverse-connected}
wynika, że \(K_T\) jest
niepustym continuum metrycznym,
a rzutowanie \(p_0\colon K_T\to[0,1]\)
jest surjektywne. Stąd \(K_T\)
ma co najmniej dwa punkty.

Przypuśćmy, że właściwe
podcontinuum \(C\subsetneq K_T\)
ma niepuste wnętrze.
Topologia granicy odwrotnej
ma bazę cylindrów
\(p_k^{-1}(U)\), gdzie
\(k\ge0\), a \(U\subseteq[0,1]\)
jest otwarty: dowolne skończone
warunki na wcześniejsze współrzędne
można przenieść do najwyższego
wśród ich indeksów.
Istnieją więc \(k\ge0\)
i niepusty otwarty \(U\subseteq[0,1]\)
takie, że
\[
 p_k^{-1}(U)\subseteq C.
 \tag{*}
\]

Niech \(j\ge k\) będzie dowolny.
Wybierzmy \(m\ge j+2\)
oraz \(m\ge k+2\)
i oznaczmy \(r=m-k\).
Surjektywność rzutowania \(p_m\)
i warunki granicy odwrotnej
dają
\[
 p_m(p_k^{-1}(U))=(T^r)^{-1}(U).
\]
W szczególności z \((*)\)
otrzymujemy
\((T^r)^{-1}(U)\subseteq p_m(C)\).
Zbiór \(p_m(C)\)
jest spójnym zwartym podzbiorem
przedziału \([0,1]\),
a zatem jest domkniętym przedziałem.
Lemat~\ref{lem:topog25-tent-branches}
pokazuje, że zawiera on
punkt z \([0,2^{-r}]\)
i punkt z \([1-2^{-r},1]\).
Ponieważ \(r\ge2\),
wynika stąd
\[
 [\tfrac14,\tfrac34]\subseteq p_m(C).
\]
Bezpośrednio z definicji \(T\)
obliczamy
\[
 T([\tfrac14,\tfrac34])=[\tfrac12,1],
 \qquad
 T^2([\tfrac14,\tfrac34])=[0,1].
\]
Ponieważ \(m-j\ge2\) oraz
\(T([0,1])=[0,1]\),
otrzymujemy
\[
 p_j(C)=T^{m-j}(p_m(C))=[0,1].
\]
Punkt \(j\ge k\) był dowolny.
Dla \(j<k\) również
\(p_j(C)=T^{k-j}(p_k(C))=[0,1]\).
Zatem wszystkie rzutowania
\(p_j(C)\) są pełnymi przedziałami.

To implikuje \(C=K_T\).
Istotnie, gdyby \(x\in K_T\setminus C\),
otwartość dopełnienia \(C\)
i bazowość cylindrów dały
\(j\) i otwarty \(W\subseteq[0,1]\)
takie, że
\(x\in p_j^{-1}(W)\subseteq K_T\setminus C\).
Wówczas \(p_j(C)\cap W=\varnothing\),
wbrew \(p_j(C)=[0,1]\).
Otrzymaliśmy sprzeczność.

Każde właściwe podcontinuum
\(K_T\) ma więc puste wnętrze.
Stwierdzenie~\ref{prop:topog25-nowhere-dense}
daje nierozkładalność \(K_T\).
\end{proof}

\begin{figure}[htbp]
\centering
\begin{tikzpicture}[x=5.0cm,y=3.0cm]
  \draw[->] (-.05,0)--(1.09,0) node[below] {$t$};
  \draw[->] (0,-.05)--(0,1.13) node[left] {$T(t)$};
  \draw[very thick,blue!65!black]
    (0,0)--(.5,1)--(1,0);
  \draw[densely dashed,gray] (.5,0)--(.5,1);
  \node[below,font=\small] at (0,0) {$0$};
  \node[below,font=\small] at (.5,0) {$\tfrac12$};
  \node[below,font=\small] at (1,0) {$1$};
  \node[left,font=\small] at (0,1) {$1$};
  \node[anchor=west,font=\small,align=left] at (.57,.88)
    {Każda z dwóch gałęzi\\pokrywa cały przedział};
\end{tikzpicture}
\caption{Odwzorowanie namiotowe \(T\); jego kolejne
iteracje mają \(2^r\) pełnych gałęzi afinicznych.
Ta własność jest kluczowa dla nierozkładalności \(K_T\).}
\label{fig:topog25-tent}
\end{figure}

\section{Solenoid dyadyczny a spójność lokalna}
\label{sec:topog25-solenoid}

\begin{definicja}[Solenoid dyadyczny]
\label{def:topog25-solenoid}
Niech \(S^1=\{z\in\mathbb C:|z|=1\}\).
\emph{Solenoidem dyadycznym} nazywamy
granicę odwrotną
\[
 \Sigma_2=\varprojlim
 \bigl(S^1\xleftarrow{z\mapsto z^2}S^1
 \xleftarrow{z\mapsto z^2}\cdots\bigr).
\]
Punkty \(\Sigma_2\) są ciągami
\((z_0,z_1,\ldots)\) liczb z \(S^1\)
spełniającymi \(z_{n+1}^2=z_n\).
\end{definicja}

\begin{stwierdzenie}[Solenoid jest continuum]
\label{prop:topog25-solenoid-continuum}
Solenoid \(\Sigma_2\) jest
nietrywialnym continuum metrycznym.
Działanie przez mnożenie współrzędnych
\[
 (z_n)_{n\ge0}\cdot(w_n)_{n\ge0}
 =(z_nw_n)_{n\ge0}
\]
czyni go ponadto zwartą
przemienną grupą topologiczną Hausdorffa.
\end{stwierdzenie}
\begin{proof}
Okrąg jest nietrywialnym
continuum metrycznym,
a odwzorowanie \(z\mapsto z^2\)
jest ciągłą surjekcją.
Twierdzenie~\ref{thm:topog25-inverse-connected}
daje pierwszą tezę wraz
z nietrywialnością, ponieważ
rzutowanie \(p_0\) jest surjektywne.

Jeżeli \(z_{n+1}^2=z_n\)
i \(w_{n+1}^2=w_n\),
to \((z_{n+1}w_{n+1})^2=z_nw_n\),
więc iloczyn współrzędny
ponownie należy do \(\Sigma_2\).
Elementem neutralnym jest
ciąg samych jedynek,
a odwrotnością ciągu \((z_n)\)
jest \((z_n^{-1})\).
Łączność i przemienność wynikają
z mnożenia na okręgu.
Mnożenie i odwracanie
są ciągłe w topologii produktu,
więc są ciągłe także
na podprzestrzeni \(\Sigma_2\).
\end{proof}

\begin{lemat}[Lokalny model solenoidu]
\label{lem:topog25-solenoid-local}
Niech \(J\subseteq\mathbb R\)
będzie niepustym przedziałem otwartym
o długości mniejszej niż \(1\),
a \(U=\{e^{2\pi it}:t\in J\}\subseteq S^1\).
Wówczas przeciwobraz
\(p_0^{-1}(U)\subseteq\Sigma_2\)
jest homeomorficzny z
\(J\times\{0,1\}^{\mathbb N}\),
gdzie drugi czynnik ma topologię produktową.
\end{lemat}
\begin{proof}
Odwzorowanie
\(t\mapsto e^{2\pi it}\)
jest homeomorfizmem \(J\to U\),
ponieważ \(|J|<1\).
Dla \(t\in J\) oraz
\(\varepsilon=(\varepsilon_1,\varepsilon_2,\ldots)
\in\{0,1\}^{\mathbb N}\)
konstruujemy rekurencyjnie
\(t_0=t\) oraz
\[
 t_{n+1}=\frac{t_n+\varepsilon_{n+1}}{2},
 \qquad z_n=e^{2\pi it_n}.
\]
Wtedy \(z_{n+1}^2=z_n\),
więc \((z_n)\in p_0^{-1}(U)\).
Daje to odwzorowanie
\(\Phi\colon J\times\{0,1\}^{\mathbb N}
\to p_0^{-1}(U)\).
Każda współrzędna \(z_n\)
zależy ciągle od \(t\)
i od skończonej liczby
\(\varepsilon_1,\ldots,\varepsilon_n\),
zatem \(\Phi\) jest ciągłe.

Z danego ciągu \((z_n)\)
odtwarzamy jednoznacznie \(t_0\in J\).
W każdym kroku \(z_{n+1}\)
jest jednym z dwóch różnych
pierwiastków kwadratowych z \(z_n\),
którym odpowiadają dokładnie
dwie wartości \(t_{n+1}\)
w podanym wzorze.
Wybór \(\varepsilon_{n+1}\)
jest więc jednoznaczny
i otrzymujemy odwrotność \(\Phi\).

Aby wykazać ciągłość tej odwrotności,
ustalmy punkt i dowolne \(N\).
Po ograniczeniu \(t_0\)
do małego przedziału wokół
wybranej wartości oba pierwiastki
kwadratowe \(z_0\) można rozróżnić
przez rozłączne otwarte łuki.
To zapewnia lokalną stałość
\(\varepsilon_1\).
Następnie postępujemy analogicznie
dla kolejnych pierwiastków,
uzyskując ciągłość każdej
skończonej listy
\((\varepsilon_1,\ldots,\varepsilon_N)\).
Są to wszystkie podstawowe warunki
topologii produktowej
na \(\{0,1\}^{\mathbb N}\),
więc odwrotność \(\Phi^{-1}\)
jest ciągła.
\end{proof}

\begin{twierdzenie}[Solenoid nie jest lokalnie spójny]
\label{thm:topog25-solenoid-not-locally-connected}
Solenoid dyadyczny \(\Sigma_2\)
jest continuum metrycznym,
które nie jest lokalnie spójne
w żadnym punkcie.
\end{twierdzenie}
\begin{proof}
Niech \(x\in\Sigma_2\).
Wybierzmy krótki otwarty łuk
\(U\subsetneq S^1\) zawierający
pierwszą współrzędną \(p_0(x)\).
Z Lematu~\ref{lem:topog25-solenoid-local}
otwarte otoczenie
\(p_0^{-1}(U)\) punktu \(x\)
jest homeomorficzne
z \(J\times\{0,1\}^{\mathbb N}\).
Przestrzeń Cantora
\(\{0,1\}^{\mathbb N}\)
jest całkowicie niespójna
i nie ma punktów izolowanych.
Dlatego każdy spójny podzbiór
\(J\times\{0,1\}^{\mathbb N}\)
ma stałą drugą współrzędną:
jej obraz przy rzutowaniu
na przestrzeń Cantora jest spójny,
więc jest singletonem.

Przypuśćmy, że \(\Sigma_2\)
jest lokalnie spójny w \(x\).
Wówczas wewnątrz
\(p_0^{-1}(U)\) istnieje
spójne otwarte otoczenie \(W\ni x\).
Jego obraz przez \(\Phi^{-1}\)
musiałby leżeć w
\(J\times\{\varepsilon\}\)
dla pewnego ciągu \(\varepsilon\).
Jednocześnie byłby
niepusty i otwarty
w \(J\times\{0,1\}^{\mathbb N}\).
To niemożliwe, ponieważ
singleton \(\{\varepsilon\}\)
nie jest otwarty w przestrzeni Cantora.
Otrzymujemy sprzeczność.
Punkt \(x\) był dowolny.
\end{proof}

\section{Continua Peana i odwzorowania przedziału}
\label{sec:topog25-peano}

\begin{definicja}[Continuum Peana]
\label{def:topog25-peano}
\emph{Continuum Peana} nazywamy
continuum metryczne,
które jest lokalnie spójne.
Takie przestrzenie obejmują
odcinki, okręgi i skończone
spójne grafy topologiczne.
\end{definicja}

\begin{definicja}[Skończony graf topologiczny]
\label{def:topog25-finite-graph}
\emph{Skończonym grafem topologicznym}
nazywamy realizację geometryczną
skończonego grafu kombinatorycznego,
otrzymaną przez sklejenie
skończenie wielu egzemplarzy \([0,1]\)
w wierzchołkach, przy założeniu,
że wnętrza różnych krawędzi
są rozłączne.
Dopuszczamy graf jednoelementowy
bez krawędzi.
\end{definicja}

\begin{twierdzenie}[Skończone grafy spójne są obrazami przedziału]
\label{thm:topog25-finite-graph-peano}
Niech \(P\) będzie realizacją
geometryczną skończonego,
spójnego grafu kombinatorycznego.
Wówczas \(P\) jest continuum Peana.
Ponadto istnieje ciągła surjekcja
\(\gamma\colon[0,1]\to P\).
\end{twierdzenie}
\begin{proof}
Jeśli graf ma jeden wierzchołek
i nie ma krawędzi, teza jest oczywista:
wybieramy odwzorowanie stałe.
Dalej załóżmy, że graf ma
przynajmniej jedną krawędź.

Realizacja \(P\) jest ilorazem
skończonej sumy rozłącznej
odcinków, zidentyfikowanych
w skończenie wielu końcach.
Jest zwarta, ponieważ
jest ciągłym obrazem
skończonej sumy zwartych odcinków.
Można ją metryzować
metryką długości najkrótszej
drogi wzdłuż krawędzi
(po przypisaniu każdej krawędzi
długości \(1\)).
Dla różnych punktów odległość
jest dodatnia; trójkątna nierówność
wynika ze składania dróg.
Topologia tej metryki jest
topologią realizacji:
we wnętrzu krawędzi ma bazę
małych przedziałów, a w wierzchołku
ma bazę małych gwiazd
złożonych z początkowych odcinków
wszystkich incydentnych krawędzi.
Są to zbiory otwarte i spójne.
Zatem \(P\) jest lokalnie spójna.

Spójność grafu kombinatorycznego
oznacza, że dowolne dwa wierzchołki
łączy skończony ciąg krawędzi.
Dołączając początkowy i końcowy
odcinek krawędzi, otrzymujemy
drogę ciągłą między dowolnymi
punktami \(P\).
Stąd \(P\) jest spójna,
a więc jest continuum Peana.

Aby skonstruować surjekcję
z przedziału, wybierzmy wierzchołek
bazowy \(v\).
Ponieważ graf jest skończony i spójny,
dla każdej krawędzi \(e\)
istnieje skończona droga z \(v\)
do jednego końca \(e\).
Przechodzimy tę drogę,
następnie pokonujemy krawędź \(e\),
po czym wracamy do \(v\)
wzdłuż odwróconej drogi
z drugiego końca \(e\).
Otrzymujemy zamknięty spacer
zaczynający się i kończący w \(v\),
którego obraz zawiera całą krawędź \(e\).
Wykonujemy kolejno takie spacery
dla wszystkich krawędzi grafu.
Jest ich skończenie wiele,
a wszystkie łączą się w \(v\).
Parametryzując kolejne odcinki
spacerów liniowo na odpowiednich
podprzedziałach \([0,1]\),
otrzymujemy ciągłe odwzorowanie
\(\gamma\colon[0,1]\to P\).
Jego obraz zawiera każdą krawędź
oraz wszystkie wierzchołki,
więc jest równy \(P\).
\end{proof}

\begin{twierdzenie}[Istnienie krzywych wypełniających sześcian]
\label{thm:topog25-space-filling}
Dla każdej liczby całkowitej \(n\ge1\)
istnieje ciągła surjekcja
\[
 F_n\colon[0,1]\longrightarrow[0,1]^n.
\]
W szczególności dla \(n=2\)
istnieje krzywa Peana wypełniająca
cały kwadrat.
\end{twierdzenie}
\begin{proof}
Niech \(C\subseteq[0,1]\)
będzie klasycznym zbiorem Cantora.
Przypomnijmy z
Przykładu~\ref{ex:topog23-cantor},
że odwzorowanie
\[
 c\colon\{0,1\}^{\mathbb N}\longrightarrow C,
 \qquad c((\varepsilon_k)_{k\ge1})
 =\sum_{k=1}^{\infty}\frac{2\varepsilon_k}{3^k}
\]
jest homeomorfizmem.
Dla ustalonego \(n\ge1\)
definiujemy na przestrzeni
\(\{0,1\}^{\mathbb N}\)
funkcje
\[
 b_j((\varepsilon_k)_{k\ge1})
 =\sum_{m=1}^{\infty}
 \frac{\varepsilon_{n(m-1)+j}}{2^m},
 \qquad 1\le j\le n.
\]
Każda \(b_j\) jest ciągła:
suma pierwszych \(M\) wyrazów
zależy od skończenie wielu
współrzędnych, a ogon ma moduł
co najwyżej \(2^{-M}\),
jednostajnie względem wszystkich
ciągów \((\varepsilon_k)\).
Wobec tego
\(b=(b_1,\ldots,b_n)\)
jest ciągłym odwzorowaniem
do sześcianu \([0,1]^n\).

Odwzorowanie \(b\) jest surjektywne.
Istotnie, dla dowolnego
\(y=(y_1,\ldots,y_n)\in[0,1]^n\)
wybieramy rozwinięcia dwójkowe
\[
 y_j=\sum_{m=1}^{\infty}
       \frac{\eta_{j,m}}{2^m},
 \qquad \eta_{j,m}\in\{0,1\}.
\]
Dla liczby \(1\) wybieramy
rozwinięcie z samych jedynek.
Przypisując
\(\varepsilon_{n(m-1)+j}=\eta_{j,m}\),
otrzymujemy ciąg, dla którego
\(b((\varepsilon_k))=y\).
Zatem odwzorowanie
\(f=b\circ c^{-1}\colon C\to[0,1]^n\)
jest ciągłe i surjektywne.

Zbiór \(C\) jest domknięty
w normalnej przestrzeni \([0,1]\).
Każda współrzędna \(f_j\colon C\to[0,1]\)
ma więc, na mocy
twierdzenia Tietzego o przedłużeniu
\ref{thm:topog-tietze},
ciągłe przedłużenie
\(F_j\colon[0,1]\to[0,1]\).
Odwzorowanie
\(F_n=(F_1,\ldots,F_n)\colon[0,1]\to[0,1]^n\)
jest ciągłe.
Jego obraz zawiera \(F_n(C)=f(C)=[0,1]^n\),
a jednocześnie jest podzbiorem \([0,1]^n\);
wobec tego jest surjekcją.
\end{proof}

\begin{uwaga}[Zakres twierdzenia Hahna--Mazurkiewicza]
\label{rem:topog25-hahn-mazurkiewicz}
Klasyczne twierdzenie Hahna--Mazurkiewicza
mówi, że niepusta przestrzeń Hausdorffa
jest ciągłym obrazem przedziału \([0,1]\)
wtedy i tylko wtedy, gdy jest
zwarta, spójna, lokalnie spójna
i metryzowalna.
Pełny dowód obu implikacji
wymaga osobnej konstrukcji
kolejnych coraz drobniejszych
parametryzacji; nie wynika
z udowodnionego wyżej
twierdzenia o skończonych grafach.
W niniejszym rozdziale
podaliśmy pełny dowód
dla skończonych grafów topologicznych,
natomiast ogólną równoważność
zaznaczamy jako klasyczny
wynik do szczegółowego opracowania
w dalszej iteracji.
Solenoid z
Twierdzenia~\ref{thm:topog25-solenoid-not-locally-connected}
nie spełnia warunku lokalnej spójności,
mimo że jest continuum metrycznym.
\end{uwaga}

\par\smallskip\noindent\textbf{Podsumowanie.}
Minimalne podcontinuum między dowolnymi
dwoma punktami istnieje dzięki
lematowi Kuratowskiego--Zorna
i twierdzeniu o przecięciu łańcucha.
Granice odwrotne ciągów continuów
z surjektywnymi odwzorowaniami wiążącymi
pozostają continuami.
Dla odwzorowania namiotowego
otrzymaliśmy nietrywialne
continuum nierozkładalne Knastera.
Solenoid dyadyczny jest natomiast
continuum metrycznym i zwartą
grupą topologiczną, lecz
nie jest lokalnie spójny.
Dla skończonych grafów spójnych
i wszystkich sześcianów \\( [0,1]^n \\)
skonstruowaliśmy ciągłe surjekcje
z odcinka, bez powtarzania
podstaw spójności z rozdziału~08.
\chapter{Twierdzenie Hahna--Mazurkiewicza i konstrukcja krzywych Peana}
\label{ch:topog26-hahn-mazurkiewicz}

\noindent\textbf{Cel rozdziału.}
W rozdziale~\ref{ch:topog25-continua} wprowadziliśmy
continua Peana i udowodniliśmy istnienie ciągłych
surjekcji odcinka na skończone spójne grafy
oraz na sześciany \([0,1]^n\).
Pozostała ogólna charakteryzacja przestrzeni,
które są ciągłymi obrazami odcinka.
Udowodnimy ją w obu kierunkach,
ze szczególnym naciskiem na konstrukcję
krzywej surjektywnej i uzasadnienie jej ciągłości.
Nie powtarzamy definicji zwartości ani
podstawowych twierdzeń o spójności
z rozdziałów~07--08.

\section{Małe spójne otoczenia w przestrzeni zwartej}
\label{sec:topog26-uniform-connected}

\begin{lemat}[Jednostajna lokalna spójność]
\label{lem:topog26-uniform-local-connected}
Niech \((X,d)\) będzie niepustą zwartą
przestrzenią metryczną, która jest lokalnie spójna.
Dla każdego \(\varepsilon>0\) istnieje
\(\delta>0\) takie, że dowolne
\(a,b\in X\) spełniające \(d(a,b)<\delta\)
należą do wspólnego spójnego
zbioru otwartego \(V\subseteq X\)
o średnicy mniejszej niż \(\varepsilon\).
\end{lemat}
\begin{proof}
Dla każdego \(x\in X\) z lokalnej spójności
wybieramy otwarty spójny zbiór
\(V_x\) taki, że
\(x\in V_x\subseteq B_d(x,\varepsilon/3)\).
Wówczas \(\operatorname{diam}V_x<\varepsilon\).
Z otwartości \(V_x\) wybieramy \(r_x>0\)
takie, że \(B_d(x,2r_x)\subseteq V_x\).
Zwartość \(X\) daje skończone
podpokrycie
\[
 X=\bigcup_{i=1}^{m}B_d(x_i,r_{x_i}).
\]
Niech \(\delta=\min_{1\le i\le m}r_{x_i}>0\).
Jeżeli \(d(a,b)<\delta\), wybieramy
\(i\) z \(a\in B_d(x_i,r_{x_i})\).
Wtedy
\[
 d(b,x_i)\le d(b,a)+d(a,x_i)
       <\delta+r_{x_i}\le2r_{x_i}.
\]
Oba punkty należą zatem do \(V_{x_i}\),
który jest otwarty, spójny
i ma średnicę mniejszą niż \(\varepsilon\).
\end{proof}

\begin{lemat}[Łańcuchy w zbiorach spójnych]
\label{lem:topog26-chain}
Niech \((Y,d)\) będzie spójną przestrzenią metryczną
i \(\delta>0\). Dla dowolnych \(a,b\in Y\)
istnieją \(m\ge1\) i punkty
\(a=y_0,y_1,\ldots,y_m=b\)
takie, że \(d(y_{j-1},y_j)<\delta\)
dla wszystkich \(j=1,\ldots,m\).
\end{lemat}
\begin{proof}
Ustalmy \(a\in Y\) i oznaczmy przez \(A\)
zbiór wszystkich punktów, które można
połączyć z \(a\) skończonym łańcuchem
o odległościach kolejnych punktów
mniejszych od \(\delta\).
Zbiór \(A\) jest niepusty, gdyż zawiera \(a\).
Jeżeli \(y\in A\), to \(B_d(y,\delta)\subseteq A\):
wystarczy dołączyć kolejny punkt do łańcucha.
Zatem \(A\) jest otwarty.
Jeżeli \(z\in Y\setminus A\), to kula
\(B_d(z,\delta)\) nie spotyka \(A\),
bo w przeciwnym razie \(z\in A\).
Dopełnienie \(A\) jest więc także otwarte.
Ze spójności \(Y\) otrzymujemy \(A=Y\).
\end{proof}

\section{Budowa drogi z coraz drobniejszych łańcuchów}
\label{sec:topog26-paths}

\begin{twierdzenie}[Małe drogi w continuach Peana]
\label{thm:topog26-small-paths}
Niech \((X,d)\) będzie niepustym
continuum metrycznym, lokalnie spójnym.
Dla każdego \(\eta>0\) istnieje
\(\delta>0\) takie, że dla dowolnych
\(a,b\in X\) spełniających \(d(a,b)<\delta\)
istnieje droga \(\gamma\colon[0,1]\to X\)
łącząca \(a\) z \(b\), której obraz
ma średnicę mniejszą niż \(\eta\).
W szczególności \(X\) jest
lokalnie spójna drogowo i spójna drogowo.
\end{twierdzenie}
\begin{proof}
Ustalmy \(\eta>0\).
Wybierzmy dodatni ciąg \((\varepsilon_n)_{n\ge0}\)
z
\(\sum_{n=0}^{\infty}\varepsilon_n<\eta/4\),
na przykład
\(\varepsilon_n=\eta\,2^{-n-4}\).
Z Lematu~\ref{lem:topog26-uniform-local-connected}
dla każdego \(n\) wybieramy
\(\delta_n>0\) takie, że
każde dwa punkty odległe o mniej
niż \(\delta_n\) należą do
wspólnego spójnego zbioru
o średnicy mniejszej niż \(\varepsilon_n\).
Możemy dodatkowo żądać
\(\delta_{n+1}<\delta_n/2\).
Niech \(\delta=\delta_0\).

\emph{Krok 1: konstruowanie wartości na gęstym zbiorze.}
Dla \(a,b\in X\) z \(d(a,b)<\delta_0\)
przypiszmy wartości \(a,b\)
końcom odcinka \([0,1]\).
Załóżmy indukcyjnie, że skonstruowaliśmy
skończony podział
\[
 0=t_{n,0}<t_{n,1}<\cdots<t_{n,N_n}=1
\]
oraz wartości \(x_{n,j}\in X\)
w jego wierzchołkach, zachowujące
\(x_{n,0}=a\), \(x_{n,N_n}=b\),
takie że
\[
 d(x_{n,j-1},x_{n,j})<\delta_n.
 \tag{1}
\]
Dla \(n=0\) przyjmujemy
podział o jednym odcinku i
wartości \(a,b\).

Dla każdego odcinka
\([t_{n,j-1},t_{n,j}]\)
Lemat~\ref{lem:topog26-uniform-local-connected}
dostarcza spójnego zbioru
\(C_{n,j}\subseteq X\),
zawierającego oba końcowe punkty
i mającego średnicę mniejszą
niż \(\varepsilon_n\).
Lemat~\ref{lem:topog26-chain}
pozwala połączyć te punkty
skończonym \(\delta_{n+1}\)-łańcuchem
złożonym z punktów \(C_{n,j}\).
Dzielimy rozpatrywany odcinek
czasu na tyle równych części,
ile łańcuch ma odcinków,
i przypisujemy nowym wierzchołkom
kolejne punkty łańcucha.
Jeżeli wybrany łańcuch ma
tylko jeden odcinek, dopisujemy
na jego początku jeszcze raz
pierwszy punkt. Dzięki temu
każdy odcinek czasu zostaje
podzielony na co najmniej dwie części.
Na wspólnych końcach starych
odcinków zachowujemy poprzednio
przypisane wartości.

W ten sposób otrzymujemy
podział poziomu \(n+1\),
zgodny z poprzednim,
a na każdym jego odcinku
kolejne wartości spełniają (1)
z indeksem \(n+1\).
Długość najdłuższego odcinka
podziału poziomu \(n\)
jest nie większa niż \(2^{-n}\).

Niech \(D\subseteq[0,1]\) będzie
sumą zbiorów wierzchołków
wszystkich podziałów.
Jest to zbiór gęsty:
średnice odcinków podziału dążą do zera.
Ponieważ wartości na starych wierzchołkach
nie są zmieniane, konstrukcja
wyznacza jednoznaczne odwzorowanie
\(\gamma_0\colon D\to X\).

\emph{Krok 2: oszacowanie oscylacji.}
Niech
\(S_n=\sum_{k=n}^{\infty}\varepsilon_k\).
Ustalmy dowolny odcinek \(J\)
podziału na poziomie \(n\).
Wszystkie wartości przypisane
jego wierzchołkom przy kolejnych
podziałach są odległe o mniej
niż \(S_n\) od wartości
na lewym końcu \(J\):
na pierwszym etapie nowy punkt
należy do \(C_{n,j}\),
więc odległość jest mniejsza
niż \(\varepsilon_n\).
W kolejnych etapach
nowo dodany punkt leży w
spójnym zbiorze o średnicy
mniejszej niż odpowiednie
\(\varepsilon_k\), zawierającym
końce jego odcinka rodzicielskiego.
Indukcja daje sumę oszacowań
\(\varepsilon_n+\cdots+\varepsilon_{N-1}\)
na dowolnym poziomie \(N>n\).

W konsekwencji oscylacja \(\gamma_0\)
na \(D\cap J\) jest co najwyżej \(2S_n\).
Ponieważ wartości przypisane
prawemu końcowi \(J\)
też należą do tego zbioru,
każdy punkt \(D\cap J\)
jest od niego odległy
o co najwyżej \(2S_n\).
Dla dwóch sąsiednich odcinków
poziomu \(n\) dostajemy więc
oszacowanie \(4S_n\),
wykorzystując ich wspólny wierzchołek.

Teraz dla dowolnego \(\rho>0\)
wybieramy \(n\) tak, aby \(4S_n<\rho\).
Ponieważ podział poziomu \(n\)
jest skończony, ma dodatnią
najmniejszą długość odcinka \(h_n\).
Jeżeli \(u,v\in D\) i
\(|u-v|<h_n\), to \(u,v\)
należą do tego samego lub
do dwóch sąsiednich odcinków
podziału poziomu \(n\).
Zatem \(d(\gamma_0(u),\gamma_0(v))<\rho\).
Odwzorowanie \(\gamma_0\)
jest jednostajnie ciągłe na \(D\).

\emph{Krok 3: przedłużenie do drogi.}
Zwartość metrycznego \(X\)
implikuje jego zupełność.
Dla każdego \(t\in[0,1]\)
wybieramy ciąg \(t_k\in D\)
z \(t_k\to t\).
Jednostajna ciągłość \(\gamma_0\)
sprawia, że \((\gamma_0(t_k))\)
jest ciągiem Cauchy'ego,
więc ma granicę w \(X\).
Granica nie zależy od wyboru
ciągu \(t_k\), ponieważ dla
dwóch takich ciągów
ich odległość w parametrze
dąży do zera.
Definiujemy \(\gamma(t)\)
jako tę granicę.
Otrzymujemy ciągłe przedłużenie
\(\gamma\colon[0,1]\to X\)
z \(\gamma(0)=a,\gamma(1)=b\).

Stosując oszacowanie z kroku 2
do całego początkowego odcinka
i przechodząc do granicy,
otrzymujemy
\(d(\gamma(t),a)\le S_0\)
dla wszystkich \(t\).
Wobec tego
\(\operatorname{diam}\gamma([0,1])
\le2S_0<\eta\).
Udowodniliśmy główną tezę.

\emph{Krok 4: spójność drogowa lokalna i globalna.}
Niech \(U\subseteq X\) będzie otwarty,
a \(x\in U\).
Wybierzmy \(r>0\)
z \(B_d(x,r)\subseteq U\).
Konstrukcja powyżej z \(\eta<r\)
daje \(\delta_x>0\) takie, że
każdy \(y\in B_d(x,\delta_x)\)
można połączyć z \(x\)
drogą o obrazie zawartym
w \(B_d(x,r)\subseteq U\).
Dlatego składowa drogowa
\(U\) zawierająca \(x\)
zawiera otwarte otoczenie \(x\).
Stosując ten argument do
każdego punktu tej składowej,
widzimy, że jest ona otwarta.
Wszystkie składowe drogowe
wszystkich otwartych podzbiorów
są więc otwarte.
Z charakterystyki
\ref{thm:topog-spojnosc-lokalna-komponenty}
otrzymujemy lokalną spójność drogową.
Ponieważ \(X\) jest spójna,
wniosek~\ref{cor:topog-spojnosc-lokalna-drogowa}
daje spójność drogową \(X\).
\end{proof}

\begin{uwaga}[Dlaczego nie wystarczy spójność lokalna bez zwartości?]
\label{rem:topog26-compactness}
Zwartość była niezbędna do uzyskania
jednego \(\delta\) działającego
dla wszystkich punktów \(X\).
Ponadto jej konsekwencja, zupełność,
pozwoliła przedłużyć funkcję
z gęstego zbioru parametrów.
Sama spójność lokalna w dowolnej
przestrzeni metrycznej nie uzasadnia
powyższej konstrukcji.
\end{uwaga}

\section{Dołączanie krótkich wycieczek do drogi}
\label{sec:topog26-excursions}

\begin{lemat}[Uzupełnienie obrazu drogi o skończony zbiór]
\label{lem:topog26-excursion}
Niech \((X,d)\) będzie continuum Peana,
\(h\colon[0,1]\to X\) drogą,
\(\varepsilon>0\) i \(\delta>0\)
takimi, że każdą parę punktów
odległych o mniej niż \(\delta\)
można połączyć drogą o obrazie
średnicy mniejszej niż \(\varepsilon\).
Niech \(E\subseteq X\)
będzie niepustym skończonym zbiorem,
spełniającym
\[
 \forall y\in E\quad
 \exists t_y\in[0,1]\quad
 d(y,h(t_y))<\delta.
\]
Wówczas istnieje droga
\(g\colon[0,1]\to X\)
taka, że
\[
 g(0)=h(0),\qquad g(1)=h(1),\qquad
 E\subseteq g([0,1]),\qquad
 \sup_{t\in[0,1]}d(g(t),h(t))<2\varepsilon.
\]
\end{lemat}
\begin{proof}
Zbierzmy wszystkie różne wartości
\(t_y\) w skończony zbiór
\(T=\{s_1,\ldots,s_m\}\subseteq[0,1]\).
Dla każdego \(s_i\) wybieramy
domknięty przedział
\(J_i=[a_i,b_i]\subseteq[0,1]\)
z \(a_i<b_i\) i \(s_i\in J_i\),
tak aby przedziały \(J_i\)
były parami rozłączne oraz
\[
 \operatorname{diam}h(J_i)<\varepsilon.
 \tag{2}
\]
Można to zrobić dzięki jednostajnej
ciągłości \(h\) oraz skończoności \(T\).
Gdy \(s_i=0\) lub \(1\),
wybieramy przedział jednostronny.

Dla każdego \(y\in E\) z \(t_y=s_i\)
Twierdzenie~\ref{thm:topog26-small-paths}
daje drogę
\(\alpha_y\colon[0,1]\to X\)
od \(h(s_i)\) do \(y\),
o obrazie średnicy mniejszej
niż \(\varepsilon\).
Połączenie \(\alpha_y\)
z drogą przebieganą w przeciwnym kierunku
jest pętlą zaczynającą się
i kończącą w \(h(s_i)\),
której obraz zawiera \(y\).

Na każdym \(J_i\) budujemy
odcinkami ciągłą drogę
przebiegającą kolejno:
odcinek starej drogi \(h\)
od \(a_i\) do \(s_i\),
wszystkie pętle do punktów
\(y\) z \(t_y=s_i\),
a następnie odcinek \(h\)
od \(s_i\) do \(b_i\).
Parametryzujemy te finitnie
liczne fragmenty liniowo
na kolejnych częściach \(J_i\).
Jeśli \(s_i=a_i\) lub \(s_i=b_i\),
odpowiedni fragment stały można pominąć.
Na dopełnieniu sumy przedziałów \(J_i\)
kładziemy \(g(t)=h(t)\).
Nowe i stare fragmenty zgadzają się
w punktach \(a_i,b_i\),
zatem \(g\) jest ciągła.
Ma wymagane końce i zawiera każdy punkt \(E\).

Dla \(t\in J_i\) punkt \(g(t)\)
jest albo punktem z \(h(J_i)\),
albo punktem pętli do któregoś \(y\).
W obu przypadkach
\(d(g(t),h(s_i))<\varepsilon\),
na mocy (2) i średnicy obrazu
każdej pętli.
Ponieważ również
\(d(h(t),h(s_i))<\varepsilon\),
mamy \(d(g(t),h(t))<2\varepsilon\).
Poza przedziałami \(J_i\)
funkcje są identyczne.
To kończy dowód.
\end{proof}

\begin{figure}[htbp]
\centering
\begin{tikzpicture}[x=1.1cm,y=1.05cm,>=stealth]
  \draw[gray!60,thick] (0,0) -- (1,0.35) -- (2,0.15)
    -- (3,0.5) -- (4,0.2) -- (5,0.7) -- (6,0.45);
  \draw[very thick,blue!70!black]
    (0,0) -- (1,0.35) -- (2,0.15) -- (2.6,0.42)
    .. controls (2.75,1.55) and (3.4,1.35) .. (3.0,0.5)
    .. controls (3.3,1.45) and (3.75,1.05) .. (3.0,0.5)
    -- (4,0.2) -- (5,0.7) -- (6,0.45);
  \fill[blue!70!black] (3,0.5) circle (2pt);
  \fill[blue!70!black] (3.2,1.13) circle (2pt);
  \node[font=\small,anchor=west] at (3.37,1.40)
     {punkt dołączony krótką pętlą};
  \node[font=\small,anchor=west] at (.1,-.40)
     {szara: droga wyjściowa};
  \node[font=\small,blue!70!black,anchor=west] at (3.6,-.40)
     {niebieska: droga po poprawce};
\end{tikzpicture}
\caption{Schemat lokalnej poprawki drogi:
na krótkim przedziale parametru
dołączamy pętlę przechodzącą przez
nowy punkt, zachowując końce całej drogi.
Rysunek ma charakter ilustracyjny.}
\label{fig:topog26-excursion}
\end{figure}

\section{Twierdzenie Hahna--Mazurkiewicza: wystarczalność}
\label{sec:topog26-sufficiency}

\begin{twierdzenie}[Twierdzenie Hahna--Mazurkiewicza, kierunek konstrukcyjny]
\label{thm:topog26-hm-sufficient}
Niech \(X\) będzie niepustym
continuum metrycznym, lokalnie spójnym.
Istnieje ciągła surjekcja
\(f\colon[0,1]\to X\).
\end{twierdzenie}
\begin{proof}
Ustalmy zgodną metrykę \(d\) na \(X\).
Wybierzmy ciąg dodatni
\(\varepsilon_n=2^{-n-3}\)
oraz, na mocy
Twierdzenia~\ref{thm:topog26-small-paths},
ciąg \((\delta_n)_{n\ge0}\) liczb dodatnich
taki, że
\(\delta_{n+1}<\delta_n/2\)
i każdą parę punktów
odległych o mniej niż \(\delta_n\)
można połączyć drogą
o obrazie średnicy mniejszej
niż \(\varepsilon_n\).

\emph{Krok 1: pierwsza droga o obrazie gęstym
w skali \(\delta_0\).}
Zwartość \(X\) daje
skończony zbiór
\(E_0=\{y_1,\ldots,y_m\}\subseteq X\)
taki, że dla każdego \(x\in X\)
istnieje \(y_j\in E_0\)
z \(d(x,y_j)<\delta_0/2\).
Na mocy
Twierdzenia~\ref{thm:topog26-small-paths}
przestrzeń \(X\) jest spójna drogowo.
Możemy więc połączyć kolejno
punkty \(y_1,\ldots,y_m\)
drogami i złożyć je w jedną
ciągłą drogę
\(f_0\colon[0,1]\to X\),
której obraz zawiera \(E_0\).
Gdy \(m=1\), bierzemy
drogę stałą.
Obraz \(f_0([0,1])\)
jest \(\delta_0/2\)-gęsty w \(X\).

\emph{Krok 2: poprawianie drogi.}
Załóżmy, że skonstruowaliśmy
drogę \(f_n\colon[0,1]\to X\),
której obraz jest
\(\delta_n/2\)-gęsty.
Zwartość \(X\) pozwala
wybrać skończony zbiór
\(E_{n+1}\subseteq X\)
będący \(\delta_{n+1}/2\)-siecią,
tzn.
\[
 \forall x\in X\quad
 \exists y\in E_{n+1}\quad
 d(x,y)<\delta_{n+1}/2.
\]
Z gęstości obrazu \(f_n\)
wynika, że dla każdego
\(y\in E_{n+1}\) istnieje
\(t_y\in[0,1]\)
z \(d(y,f_n(t_y))<\delta_n/2<\delta_n\).
Lemat~\ref{lem:topog26-excursion},
z parametrami \(\delta_n\)
i \(\varepsilon_n\),
daje nową drogę
\(f_{n+1}\colon[0,1]\to X\),
której obraz zawiera
\(E_{n+1}\) oraz
\[
 \sup_{t\in[0,1]}
 d(f_{n+1}(t),f_n(t))
 \le2\varepsilon_n.
 \tag{3}
\]
Wobec tego obraz \(f_{n+1}\)
jest \(\delta_{n+1}/2\)-gęsty.
Rekurencja konstruuje
drogi \(f_n\) dla wszystkich \(n\).

\emph{Krok 3: istnienie granicy jednostajnej.}
Dla \(m>n\), stosując nierówność
trójkąta oraz (3), mamy
\[
 \sup_{t\in[0,1]}d(f_m(t),f_n(t))
 \le \sum_{k=n}^{m-1}2\varepsilon_k.
\]
Szereg \(\sum_{k\ge0}2\varepsilon_k\)
jest zbieżny, więc \((f_n)\)
jest ciągiem Cauchy'ego
w metryce zbieżności jednostajnej.
Ponieważ \(X\) jest zupełna,
dla każdego \(t\) istnieje
granica \(f(t)=\lim_n f_n(t)\),
a powyższe oszacowanie
daje zbieżność jednostajną.
Funkcja \(f\colon[0,1]\to X\)
jest ciągła jako granica
jednostajna funkcji ciągłych.

\emph{Krok 4: surjektywność granicy.}
Niech \(x\in X\).
Dla każdego \(n\) istnieje
\(t_n\in[0,1]\) z
\(d(x,f_n(t_n))<\delta_n/2\).
Zatem
\[
 d(x,f(t_n))
 \le \frac{\delta_n}{2}
  +\sup_{t\in[0,1]}d(f_n(t),f(t))
 \le \frac{\delta_n}{2}
  +\sum_{k=n}^{\infty}2\varepsilon_k
 \longrightarrow0.
\]
Ponieważ \(f([0,1])\)
jest zwartym, a więc domkniętym
podzbiorem Hausdorffowego \(X\),
istnienie punktów \(f(t_n)\)
zbiegających do \(x\)
implikuje \(x\in f([0,1])\).
Punkt \(x\) był dowolny,
więc \(f([0,1])=X\).
\end{proof}

\begin{uwaga}[Charakter konstrukcyjny]
\label{rem:topog26-constructive}
Powyższy dowód nie zakłada
istnienia krzywej wypełniającej \(X\)
na żadnym wcześniejszym etapie.
Poziom \(n\) dostarcza drogi
o obrazie gęstym w skali
\(\delta_n/2\), a koszt poprawki
jest kontrolowany przez
\(2\varepsilon_n\).
Suma kosztów jest skończona,
dlatego granica jest ciągła;
gęstość w coraz drobniejszej
skali i domkniętość obrazu
zapewniają surjektywność.
\end{uwaga}

\section{Własności konieczne obrazów odcinka}
\label{sec:topog26-necessity}

\begin{lemat}[Obraz odcinka spełnia drugi aksjomat przeliczalności]
\label{lem:topog26-image-second-countable}
Niech \(X\) będzie przestrzenią
Hausdorffa, a
\(f\colon[0,1]\to X\)
ciągłą surjekcją.
Wówczas \(X\) ma
przeliczalną bazę topologii.
\end{lemat}
\begin{proof}
Niech \(\mathcal B\)
będzie przeliczalną bazą
topologii odcinka \([0,1]\),
na przykład rodziną
przecięć przedziałów
o wymiernych końcach z \([0,1]\).
Rodzina \(\mathcal V\)
wszystkich skończonych
sum zbiorów z \(\mathcal B\)
jest przeliczalna.
Dla \(V\in\mathcal V\)
definiujemy
\[
 W_V=X\setminus f([0,1]\setminus V).
\]
Zbiór \([0,1]\setminus V\)
jest domknięty w zwartym odcinku,
a więc zwarty.
Jego obraz przez \(f\)
jest zwarty, czyli domknięty
w Hausdorffowym \(X\).
Zatem \(W_V\) jest otwarty.

Pokażemy, że
\(\{W_V:V\in\mathcal V\}\)
jest bazą \(X\).
Niech \(x\in O\), gdzie
\(O\subseteq X\) jest otwarty.
Zbiór \(f^{-1}(\{x\})\)
jest niepusty i zwarty:
jest domknięty w \([0,1]\),
ponieważ \(X\) jest Hausdorffa.
Jest zawarty w otwartym zbiorze
\(f^{-1}(O)\).
Dla każdego jego punktu
dobieramy \(B\in\mathcal B\)
zawierający ten punkt
i spełniający \(B\subseteq f^{-1}(O)\).
Ze zwartości włókna
wybieramy skończenie wiele
takich elementów bazy,
których suma \(V\in\mathcal V\)
spełnia
\[
 f^{-1}(\{x\})\subseteq V
 \subseteq f^{-1}(O).
\]
Pierwsza inkluzja daje
\(x\notin f([0,1]\setminus V)\),
a zatem \(x\in W_V\).
Druga daje \(W_V\subseteq O\):
jeśli \(y\notin O\),
każdy jego przeciwobraz
leży poza \(V\), więc
\(y\in f([0,1]\setminus V)\).
Otrzymaliśmy \(x\in W_V\subseteq O\).
Wszystkie zbiory bazy
tworzą rodzinę przeliczalną.
\end{proof}

\begin{lemat}[Ciągły obraz odcinka jest lokalnie spójny]
\label{lem:topog26-image-locally-connected}
Niech \(X\) będzie przestrzenią
Hausdorffa, a
\(f\colon[0,1]\to X\)
ciągłą surjekcją.
Wówczas \(X\) jest lokalnie spójna.
\end{lemat}
\begin{proof}
Ustalmy otwarty \(O\subseteq X\)
i punkt \(x\in O\).
Włókno
\(F=f^{-1}(\{x\})\)
jest niepuste i zwarte
oraz leży w \(f^{-1}(O)\).
Dla każdego \(t\in F\)
istnieje domknięty przedział
\(I_t\subseteq[0,1]\)
taki, że \(t\) leży w jego
wnętrzu względnym,
\(I_t\subseteq f^{-1}(O)\)
oraz \(I_t\cap F\neq\varnothing\).
Dla \(t=0\) lub \(1\)
używamy przedziałów jednostronnych.
Ze zwartości \(F\)
wybieramy skończoną
rodzinę \(I_1,\ldots,I_m\),
których wnętrza względne
pokrywają \(F\).
Niech
\[
 V=\bigcup_{j=1}^{m}\operatorname{int}_{[0,1]}I_j,
 \qquad K=\bigcup_{j=1}^{m}f(I_j).
\]
Każdy \(f(I_j)\) jest
spójny jako ciągły obraz
przedziału, zawiera \(x\),
gdyż \(I_j\) zawiera
odpowiedni punkt włókna \(F\),
i jest zawarty w \(O\).
Suma skończenie wielu
spójnych zbiorów mających
wspólny punkt \(x\)
jest spójna. Zatem
\(K\subseteq O\)
jest spójny i zawiera \(x\).

Zdefiniujmy
\[
 W=X\setminus f([0,1]\setminus V).
\]
Zbiór \(W\) jest otwarty,
ponieważ \([0,1]\setminus V\)
jest zwarty, a \(X\)
jest Hausdorffa.
Ponieważ \(F\subseteq V\),
mamy \(x\in W\).
Z kolei \(W\subseteq K\):
dla \(y\in W\) surjektywność
\(f\) daje \(t\) z \(f(t)=y\);
z definicji \(W\) wszystkie
takie punkty leżą w \(V\),
a więc należą do pewnego \(I_j\).
Otrzymaliśmy otwarte
otoczenie \(W\ni x\)
z \(W\subseteq K\subseteq O\),
gdzie \(K\) jest spójny.

Niech \(C\) będzie składową
spójną \(O\) zawierającą \(x\).
Ponieważ \(K\) jest spójny
i \(x\in K\cap C\),
mamy \(K\subseteq C\)
na mocy maksymalności składowej.
Zatem \(W\subseteq C\).
Powtarzając to rozumowanie
dla każdego punktu \(C\)
i tego samego \(O\),
widzimy, że \(C\)
jest otwarta w \(X\).
Każda składowa spójna
dowolnego otwartego \(O\)
jest więc otwarta.
Twierdzenie~\ref{thm:topog-spojnosc-lokalna-komponenty}
daje lokalną spójność \(X\).
\end{proof}

\begin{twierdzenie}[Twierdzenie Hahna--Mazurkiewicza]
\label{thm:topog26-hahn-mazurkiewicz}
Niech \(X\) będzie niepustą
przestrzenią topologiczną Hausdorffa.
Następujące warunki są równoważne:
\begin{enumerate}[label=\textup{(\roman*)}]
\item istnieje ciągła surjekcja
\(f\colon[0,1]\to X\);
\item \(X\) jest zwarta, spójna,
lokalnie spójna i metryzowalna.
\end{enumerate}
\end{twierdzenie}
\begin{proof}
Jeżeli zachodzi \textup{(ii)},
to \(X\) jest continuum Peana,
a implikację \textup{(ii)}\(\Rightarrow\)\textup{(i)}
udowodniliśmy konstrukcyjnie
w Twierdzeniu~\ref{thm:topog26-hm-sufficient}.

Załóżmy \textup{(i)}.
Jako ciągły obraz
niepustego zwartego
i spójnego odcinka
\(X\) jest niepusta,
zwarta i spójna.
Ponadto
Lemat~\ref{lem:topog26-image-locally-connected}
daje lokalną spójność \(X\).
Na mocy
Lematu~\ref{lem:topog26-image-second-countable}
przestrzeń \(X\)
spełnia drugi aksjomat
przeliczalności.
Każda zwarta przestrzeń
Hausdorffa jest regularna
i spełnia \(T_1\);
zatem twierdzenie Urysohna
o metryzowalności
\ref{thm:topog-urysohn-metryzowalnosc}
zapewnia metryzowalność \(X\).
To dowodzi \textup{(i)}\(\Rightarrow\)\textup{(ii)}.
\end{proof}

\section{Przykłady, kontrprzykłady i konsekwencje}
\label{sec:topog26-examples}

\begin{wniosek}[Continua Peana są spójne drogowo]
\label{cor:topog26-peano-path-connected}
Niech \(X\) będzie continuum metrycznym
i lokalnie spójnym.
Wówczas \(X\) jest
spójna drogowo oraz lokalnie spójna drogowo.
\end{wniosek}
\begin{proof}
Jest to ostatnia część
Twierdzenia~\ref{thm:topog26-small-paths}.
Można także wyprowadzić
samą spójność drogową
z Twierdzenia~\ref{thm:topog26-hahn-mazurkiewicz},
ponieważ ciągły obraz
drogi jest spójny drogowo,
ale nie dawałoby to
jednostajnego oszacowania
średnicy małych dróg.
\end{proof}

\begin{przyklad}[Zwartość i spójność nie wystarczą]
\label{ex:topog26-solenoid}
Solenoid dyadyczny \(\Sigma_2\),
skonstruowany w
Definicji~\ref{def:topog25-solenoid},
jest zwartym, spójnym,
metryzowalnym continuum.
Nie jest jednak lokalnie spójny
na mocy
Twierdzenia~\ref{thm:topog25-solenoid-not-locally-connected}.
Z Twierdzenia~\ref{thm:topog26-hahn-mazurkiewicz}
wynika, że nie istnieje
ciągła surjekcja
\([0,1]\to\Sigma_2\).
Jest to przykład pokazujący
istotność założenia
lokalnej spójności.
\end{przyklad}

\begin{przyklad}[Niespójność jest osobną przeszkodą]
\label{ex:topog26-cantor}
Dyskretny zbiór \(D=\{0,1\}\)
jest zwarty, Hausdorffa, metryzowalny
i lokalnie spójny: każdy singleton
jest otwartym spójnym otoczeniem
swojego punktu. Nie jest jednak
spójny, więc nie może być ciągłym
obrazem całego odcinka \([0,1]\).

Klasyczny zbiór Cantora jest również
zwarty i metryzowalny, ale nie jest
ani spójny, ani lokalnie spójny:
każdy jego niepusty spójny
podzbiór jest singletonem,
a żaden singleton nie jest otwarty.
Można z niego natomiast konstruować
ciągłe surjekcje na sześciany,
zgodnie z
Twierdzeniem~\ref{thm:topog25-space-filling}.
\end{przyklad}

\begin{przyklad}[Okrąg, grafy i sześciany]
\label{ex:topog26-classical}
Okrąg \(S^1\), każdy skończony
spójny graf topologiczny
oraz każdy sześcian \([0,1]^n\)
są continuami Peana.
Dlatego są ciągłymi
obrazami odcinka.
Dla grafów i sześcianów
konstrukcje surjekcji
zostały już podane w
Twierdzeniach
\ref{thm:topog25-finite-graph-peano}
i~\ref{thm:topog25-space-filling};
obecne twierdzenie
obejmuje je jako przypadki
ogólniejszej charakteryzacji.
\end{przyklad}

\begin{uwaga}[Istotność założenia Hausdorffa]
\label{rem:topog26-metrizability}
Bez aksjomatu Hausdorffa
ciągły obraz odcinka
nie musi być metryzowalny.
Istotnie, weźmy
\(Y=\{0,1\}\) z topologią
\(\{\varnothing,\{1\},Y\}\)
i określmy
\(f\colon[0,1]\to Y\) wzorem
\(f(0)=0\), \(f(t)=1\) dla \(t>0\).
Odwzorowanie jest surjektywne i ciągłe,
gdyż przeciwobrazem jedynego
nietrywialnego zbioru otwartego
jest \((0,1]\), otwarty w \([0,1]\).
Jednak \(Y\) nie jest Hausdorffa,
a zatem nie jest metryzowalna.
To uzasadnia jawne założenie
Hausdorffowości w
Twierdzeniu~\ref{thm:topog26-hahn-mazurkiewicz}.
\end{uwaga}

\par\smallskip\noindent\textbf{Podsumowanie.}
Twierdzenie Hahna--Mazurkiewicza
charakteryzuje ciągłe obrazy
odcinka wśród niepustych
przestrzeni Hausdorffa.
Konstrukcja wystarczalności
ma dwa niezależne etapy:
otrzymanie małych dróg
z lokalnej spójności
i zwartości oraz dołączanie
coraz drobniejszych
wycieczek do skończonych sieci.
Jednostajna zbieżność
poprawek daje krzywą ciągłą,
a coraz większa gęstość
jej przybliżeń
gwarantuje surjektywność.
W kierunku koniecznym
zwartość włókien
ciągłego odwzorowania odcinka
daje lokalną spójność
i drugi aksjomat przeliczalności,
a twierdzenie Urysohna
daje metryzowalność.
\end{document}
